\documentclass[10pt]{amsart}
\usepackage{amsmath,amssymb,amsthm,mathtools}
\newtheorem{theorem}{Theorem}[section]
\newtheorem{lemma}[theorem]{Lemma}
\newtheorem{proposition}[theorem]{Proposition}
\newtheorem{corollary}[theorem]{Corollary}
\theoremstyle{definition}
\newtheorem{definition}[theorem]{Definition}
\newtheorem{problem}[theorem]{Problem}
\theoremstyle{remark}
\newtheorem{remark}[theorem]{Remark}
\newcommand{\NA}{\operatorname{NA}}
\newcommand{\supp}{\operatorname{supp}}
\newcommand{\sgn}{\operatorname{sign}}
\newcommand{\spn}{\operatorname{span}}
\newcommand{\dist}{\operatorname{dist}}
\newcommand{\R}{\mathbb R}
\newcommand{\N}{\mathbb N}
\newcommand{\Cset}{\mathcal C}
\newcommand{\Mset}{\mathcal M}
\newcommand{\Cert}{\operatorname{Cert}}
\newcommand{\Li}{\operatorname*{Li}}
\newcommand{\cl}{\operatorname{cl}}
\newcommand{\gap}{\operatorname{gap}}
\newcommand{\Def}{\operatorname{Def}}
\newcommand{\Rec}{\mathcal R}
\newcommand{\zh}{\hat z}
\newcommand{\ip}[2]{\langle #1,#2\rangle}
\newcommand{\lin}{\mathfrak l}
\newcommand{\Exc}{\operatorname{Exc}}
\setlength{\emergencystretch}{3em}
\title[Recovery of mates on Mart\'in norms II]{Recovery of mates on Mart\'in norms II:\\companions, clean sub-windows and a designed operator}
\author{}
\date{Research continuation, October 2026}
\begin{document}
\begin{abstract}
This is the second part of a study of the density of norm-attaining
operators from Mart\'in's renormings $(c_0,p)$ into $\ell_2^2$.  Part~I
reduced density, for a designed operator $T$ and $N$ blocks, to a single
approximation statement (Lemma~Z) about first rows with finite base support
and signature room.  Here we develop tools that work at \emph{nearby} first
rows: companions with prescribed data, the threshold equation of a block,
the cost of moving the normer, a raise-transfer inequality, a uniform
one-sided transfer expansion, and windowed recovery theorems in which the
window data live at nearby first rows whose transfer error (for raise
companions: the Hilbert footprint of the raise and the change of the
normer) is below a fixed multiple of the square of the bottom of the
window.  We then fix one $N$-independent
admissible operator $T_{\mathrm{final}}$, with a diagonal base whose
entries decay doubly exponentially, bounded-gap signature sets, weights
chosen after the targets, and windows split into pigeonhole sub-windows at
which every rate object of the first row is either tiny or robust.  For
$T_{\mathrm{final}}$ and first rows with finite base support we prove that
every first row is recoverable unless, at every clean sub-window of every
large level, a near-exact coherent shift resonance occurs (the residual
(C$^*$)); tiny minors of the exact system are made exactly zero at a cheap
companion, and the remaining shift obstruction is reduced to a pigeonhole
dichotomy.  We prove that (C$^*$) does not occur at constant-sign maximal
contact, that some (C$^*$) rows are recoverable, and that under a lever
hypothesis it can be resolved for a modified design.  For infinite base
support we prove the transfer of peaks to base vectors outside $c_{00}$,
recovery under signature room, under cushion-sparse support swallowing
(with the room and peak conditions (W$^\star$), (H2), (H3-inf)) and for
box-dominated support (with the block conditions (SC$_I$), (W$_M$)), and,
for $T_{\mathrm{final}}$, recovery of first rows whose finitely many bad
carriers swallow the support with any profile (under (W$^\star$), (H2),
(H3-inf)), and of rows with infinitely many bad carriers under a rate
condition.  Lemma~Z, and the density of $\NA((c_0,p_N),\ell_2^2)$ for any
$N$ and of $\NA((c_0,p),\ell_2^2)$ for Mart\'in's norm, remain open; we
list the precise remaining cases.
\end{abstract}
\maketitle
\tableofcontents

\section{Introduction}\label{sec:II-intro}

\subsection{The problem and what Part~I proved}
Mart\'in \cite{Martin} constructed renormings $p$ of $c_0$ whose
norm-attaining functionals form an algebraically trivial set; the
construction starts from the canonical base $q$ of $c_0$ (a smoothing of
the supremum norm in the sense of Debs, Godefroy and Saint-Raymond
\cite{DGS}) and an injective norm-one operator $T$ given by his Lemma~B,
which rests on \cite[Proposition~2.8]{KLMW20}.  We call such an operator
\emph{admissible} [I,~Definition~1.3]; $p_N$ denotes the norm built with
the blocks $I=\{1,\dots,N\}$ and $p$ the norm built with $I=\N$.  The
question studied in Part~I \cite{PartI} and here is whether
$\NA((c_0,p_N),\ell_2^2)$ and $\NA((c_0,p),\ell_2^2)$ are dense.  It is a
test case for Lindenstrauss' property B \cite{Lindenstrauss} of the
two-dimensional Hilbert space (compare \cite{JW,KLMW}): a negative answer
for one $N$ and one admissible $T$ would produce a two-dimensional range
space without property B [I,~Section~9].

An operator into $\ell_2^2$ is a pair $(f,g)$ of functionals.  For a
\emph{first row} $f\in S_{p^*}$ the \emph{mate fibre} is
$\Cset(f)=\{g:(f,g)\text{ is contractive}\}$, and
\[
 \Rec:=\{f\in S_{p^*}:\ (f,g)\in\cl\NA((c_0,p),\ell_2^2)\text{ for every }
 g\in\Cset(f)\}
\]
is the recoverable set [I,~Definition~1.15]; density holds if and only if
$\Rec=S_{p^*}$ [I,~Proposition~1.13].  By [I,~Lemma~1.5], density for
$p_N$ for infinitely many $N$ implies density for $p$.  Part~I proved,
among other things:
\begin{enumerate}
\item[(I-1)] the transport of finite certificates along every sequence
$f_n\to f$ and an averaging theorem over scales [I,~Theorems~3.8
and~4.1];
\item[(I-2)] the resonance obstruction: an admissible $T$ and a first row
at which the intrinsically recoverable mates lie on one line while
$\Cset(f)$ is infinite dimensional [I,~Theorem~6.6];
\item[(I-3)] engineered recovery: at first rows with finite base support,
every mate with $d$-neutral two-piece data of weighted coefficient at most
one is recovered along engineered norm-attaining approximants
[I,~Theorem~7.18, Corollary~7.21], and every block-tame first row is
recoverable [I,~Corollary~7.22];
\item[(I-4)] for a \emph{signature-ladder} operator [I,~Definition~8.1,
Theorem~8.2]: every first row with finite base support and signature room
is recoverable [I,~Theorem~8.18], and density for $p_N$ is equivalent to
\emph{Lemma~Z} [I,~Theorem~8.20, Problem~8.21]: for every $f\in
S_{p_N^*}$, $g\in\Cset(f)$, $\rho<1$ and $\varepsilon>0$ there is $f'$ with
finite base support and signature room \textup{(}$f'\in\mathcal R_0$,
[I,~Definition~8.3]\textup{)}, $p_N^*(f'-f)<\varepsilon$ and
$\dist(\rho g,\Cset(f'))<\varepsilon$;
\item[(I-5)] for that operator, recovery of window-pinned mates, of
sign-mixed room and of some exactly swallowed signature sets
[I,~Theorems~8.31, 8.35 and~8.44].
\end{enumerate}
Lemma~Z was left open for every $N$.

\subsection{Lemma Z for the designed operator}
In Section~\ref{sec:II-Tfinal} we fix one operator $T_{\mathrm{final}}$
(Definition~\ref{def:II-Tfinal}) with base $U_{\mathrm{final}}$
(Definition~\ref{def:II-base}).  It is admissible and satisfies the
conclusion of Mart\'in's Lemma~B (Theorem~\ref{thm:II-lemmaB}), it does
not depend on $N$, and with any choice of one sub-window of each level it
satisfies (T-a)--(T-d) and (P1)--(P3) of [I,~Theorem~8.2] verbatim (these
are the conclusions of [I,~Theorem~8.2], which are all that [I,~Section~8]
uses; Lemma~\ref{lem:II-GW}(b)).  Consequently all results of [I,~Section~8]
hold for $T_{\mathrm{final}}$; in particular:
\begin{itemize}
\item for $T=T_{\mathrm{final}}$ and every $N\ge1$,
$\NA((c_0,p_N),\ell_2^2)$ is dense if and only if Lemma~Z holds for
$T_{\mathrm{final}}$ and $N$ [I,~Theorem~8.20];
\item if Lemma~Z holds for $T_{\mathrm{final}}$ and infinitely many $N$,
then $\NA((c_0,p),\ell_2^2)$ is dense for Mart\'in's norm $p$ built with
$T_{\mathrm{final}}$ (Remark~\ref{rem:II-blocks}, [I,~Lemma~1.5]);
\item Lemma~Z splits into a statement about first rows with finite base
support and an approximation statement (LSC-trunc) for infinite base
support (Proposition~\ref{prop:II-LSC}).
\end{itemize}
Membership of individual first rows in $\Rec$ does not transfer from $p_N$
to $p$: [I,~Lemma~1.5] transfers density, not rows, and the hypotheses of
the row-wise theorems below depend on $N$.

\subsection{Main results}
\emph{Recovery through nearby first rows (Section~\ref{sec:II-companions}).}
The tools of Section~\ref{sec:II-companions} hold for every admissible
operator.  First rows with prescribed data $f[a,z]$ and their companions
(Definition~\ref{def:II-companion}) allow us to move contacts, raise the
base vector and change the normer while keeping control of the forced
data.  The block threshold is the unique root of one explicit equation
(Lemma~\ref{lem:II-T}), with a Lipschitz bound and a quantitative raising
property (Lemma~\ref{lem:II-T2}); the cost of moving the normer is bounded
by the change of the data, including an unweighted term which cannot be
dropped (Lemma~\ref{lem:II-cost}, Remark~\ref{rem:II-cost}); and a same-sign
raise of the base vector obeys the raise-transfer inequality
(Lemma~\ref{lem:II-RT}), in which the $\ell_1$-mass of the raise cancels
against the normalization.  The uniform one-sided transfer expansion
(Lemma~\ref{lem:II-U}) holds along every sequence of first rows
$f_j\to f$ with finite supports $F\subseteq F_j$, with constants
independent of the quantities that the companions change; for arbitrary
supports the corresponding statement is Lemma~\ref{lem:II-Uinf}, under its
flip-sparsity hypothesis.  The \emph{Theorem~E family} then transfers
recovery from nearby first rows to $f$: if window data for the mate live at
first rows $f_j$ whose transfer error (for raise companions: the Hilbert
footprint of the raise and the change of the normer; this is not
$p^*(f_j-f)$, see Remark~\ref{rem:II-RS}(a)) is smaller than a fixed
multiple of the square of the bottom of the $j$-th window, the mate is
recovered
(Theorem~\ref{thm:II-ERT} with a transfer error, for arbitrary base
support; Theorem~\ref{thm:II-E} with exact $d$-neutral two-piece data;
Theorem~\ref{thm:II-Eshift} in Section~\ref{sec:II-Cstar} with shifted
data).  Section~\ref{sec:II-companions} also proves an exactness
certificate for two-piece data (Lemma~\ref{lem:II-cert}), a shift bound
(Lemma~\ref{lem:II-shift}), the rigidity of two-piece data
(Lemma~\ref{lem:II-sandwich}, Corollaries~\ref{cor:II-S1}
and~\ref{cor:II-S2}), a criterion for the failure of lower semicontinuity
of the fibre map along block-tame first rows (Theorem~\ref{thm:II-NL}; its
hypotheses are realized for $T_{\mathrm{final}}$ and $N=1$ at a self-aligned
row in Proposition~\ref{prop:II-NLreal} of Section~\ref{sec:II-Cstar}), and
the violation
tolerance of one engineered stage (Lemma~\ref{lem:II-VT}).

\emph{The designed operator (Section~\ref{sec:II-Tfinal}).}
$T_{\mathrm{final}}$ is a signature-ladder operator with four
modifications: a diagonal base $U_{\mathrm{final}}k_s=\mu^\star_se_s$ with
$\mu^\star_s=2^{-2^{2^s}}$; bounded-gap signature sets
$S_l=\{2^l(2i+1):i\ge1\}$ with an additional allowedness rule; weights
$c_l$ chosen after the target and the signature size of level $l$; and
windows split into pigeonhole sub-windows governed by a design factor
$\mathrm{Des}(l)$ and a finite family of rate objects
(Definitions~\ref{def:II-data}--\ref{def:II-rates}).  The recursion is well
founded and independent of $N$ (Lemma~\ref{lem:II-W}); $T_{\mathrm{final}}$
satisfies Mart\'in's Lemma~B (Theorem~\ref{thm:II-lemmaB}); and for every
first row and every level there is a \emph{clean} sub-window, at which every
rate object of $f$ of that level is either tiny ($\le b(w)$) or robust
($\ge u(w)$) (Definition~\ref{def:II-clean}, Theorem~\ref{thm:II-clean}).

\emph{Finite base support (Sections~\ref{sec:II-exact}
and~\ref{sec:II-Cstar}).}  At a clean sub-window the switching amplitudes
of two-sided decompositions are pinned diagonally (Lemmas~\ref{lem:II-diagpin}
and~\ref{lem:II-pinned}); far pulls and private banks on the diagonal base
tune finitely many carrier values exactly and two-sidedly
(Lemma~\ref{lem:II-TU}); every block has a peak with robust margin whose
bank raises the block threshold by a buffer dominating all drifts
(Lemmas~\ref{lem:II-donorpeak}, \ref{lem:II-donorraise}
and~\ref{lem:II-status}); and one exactifying companion per clean
sub-window carries all exactifications at cost $o(T_{\rm lo}^2)$
(Definition~\ref{def:II-companionw}, Lemma~\ref{lem:II-compcost},
Proposition~\ref{prop:II-transplant}).  This gives \emph{Master
Theorem~II} (Theorem~\ref{thm:II-master2}): for $T_{\mathrm{final}}$, a
first row with finite base support is in $\Rec$ if at infinitely many
levels some clean sub-window satisfies a shift condition (SH$_w$) and a ray
condition (VR$_w$).  The ray condition is removed by the exactification
of all tiny minors at a companion, with a Hoffman constant of the form
$C_f^l\mathrm{Des}(l)^2/u(w)$ (Theorem~\ref{thm:II-B},
Corollary~\ref{cor:II-B1}), using Hoffman's theorem \cite{Hoffman} through
minors (Lemma~\ref{lem:II-hoffman}) and the {\L}ojasiewicz inequality
\cite{BCR} (Lemma~\ref{lem:II-loj}).  The shift condition is reduced by the
shift dichotomy (Theorem~\ref{thm:II-C1}) to a pigeonhole rate.  The
result is \emph{Master Theorem~III$'$} (Theorem~\ref{thm:II-MTIII}) and:

\begin{itemize}
\item[] \emph{For $T=T_{\mathrm{final}}$, every $N\ge1$ and every $f\in
S_{p_N^*}$ with finite base support: if $f\notin\Rec$, then $f$ is a
\textup{(C$^*$)} row} (Corollary~\ref{cor:II-Cstarnec}), i.e.\ at every
clean sub-window of every large level the shift-pinning rate and the shift
cost are both tiny (Definition~\ref{def:II-Cstar}): a near-exact coherent
shift resonance.
\end{itemize}

About (C$^*$) we prove: it does not occur at constant-sign maximal contact
(Proposition~\ref{prop:II-constcontact}); shifted exact data are rigid
(Proposition~\ref{prop:II-C3}) and oscillating profiles force them
(Proposition~\ref{prop:II-C4}); (C$^*$)-type resonances are realized by
self-aligned first rows, which are block-tame and hence recoverable
(Theorem~\ref{thm:II-SA}, Corollary~\ref{cor:II-SArec}); near-threshold
carriers are Baire generic in every fixed-$z$ fibre
(Proposition~\ref{prop:II-Baire}) while exact all-negative data are not
(Proposition~\ref{prop:II-R5}); exact negative data without dead zones are
recovered (Theorem~\ref{thm:II-negexact}); and, for a modification
$D^{U1'}$ of $T_{\mathrm{final}}$ with absorber targets
(Definition~\ref{def:II-DU1}), \emph{Master Theorem~IV$'$}
(Theorem~\ref{thm:II-MTIV}) proves $f\in\Rec$ whenever, at infinitely many
main stages, some clean sub-window carries enough scales whose activity
class has no active shift, or admits a $\varkappa$-neutral lever
(KN$_{w,a}$) (in a mixed class, together with the absence of frustrated
positive carriers).
A lever is needed in this route: a donor raise moves the block scalar
$\varkappa$ by the same order as the raise, and the exact shifted system
depends on a block only through ratios (Corollary~\ref{cor:II-Rkt},
Proposition~\ref{prop:II-Rkt}), so peak pushes alone cannot achieve
exactness and threshold protection independently
(Remark~\ref{rem:II-onelever}; no claim is made that a one-lever argument
cannot succeed).

\emph{Infinite base support (Section~\ref{sec:II-infinite}).}  Transfer
peaks exist at every base vector, also outside $c_{00}$
(Theorem~\ref{thm:II-transferinf}); hence signature room, sign-mixed
room and window-pinned mates are recovered at arbitrary base support
(Theorems~\ref{thm:II-Bstarinf}--\ref{thm:II-cushionroom}), which turns
[I,~Remark~8.46] into a theorem (Theorem~\ref{thm:II-Binf}).  Engineered
recovery at arbitrary base support needs sparsity of flips only for the
one-sided parts of the data (Theorems~\ref{thm:II-raised}
and~\ref{thm:II-flipext}).  Finitely many cushion-sparse swallowed carriers
are harmless for every signature-ladder operator (Theorem~\ref{thm:II-R1},
under (W$^\star$), (H2), (H3-inf)),
and box-dominated support is harmless for every admissible operator
(Theorem~\ref{thm:II-N}, under (SC$_I$) and (W$_M$)).  For
$T_{\mathrm{final}}$ the raise-transfer
inequality removes all profile conditions: support swallowing of any
profile by finitely many bad carriers is harmless
(Theorems~\ref{thm:II-RS}, \ref{thm:II-RSstar} and~\ref{thm:II-NDprime},
under (W$^\star$), (H2), (H3-inf)),
and infinitely many bad carriers are harmless under the explicit rate
condition (W$_{\rm inf}$) (Theorem~\ref{thm:II-RSinf}).

\subsection{What remains open}
\emph{Lemma~Z remains open for $T_{\mathrm{final}}$ and every $N\ge1$,
including $N=1$; so does the density of $\NA((c_0,p_N),\ell_2^2)$ for every
$N$ and of $\NA((c_0,p),\ell_2^2)$ for Mart\'in's norm $p$, for every
admissible operator, including $T_{\mathrm{final}}$.}  By the results above,
for $T_{\mathrm{final}}$ the open part of Lemma~Z consists of:
\begin{enumerate}
\item[(i)] at finite base support, the (C$^*$) rows
(Problem~\ref{prob:II-Cstar}) that are not block-tame and do not lie in
$\mathcal R_0\cup\mathcal R_0^\pm\cup\mathcal R_S$; the precise
sub-problems are status coherence without a lever, mixed activity classes
for $N\ge2$, exact multi-block coupling, and uniform composition with
$d$-consistent engineered approximants (Problems~\ref{prob:II-KN}--\ref{prob:II-S2});
\item[(ii)] at infinite base support, (LSC-trunc) for the first rows not
covered by Section~\ref{sec:II-infinite}: the transport of Master
Theorem~III$'$ (only sketched, Remark~\ref{rem:II-Minf}), the rate
condition (W$_{\rm inf}$) for infinitely many bad carriers, the tuning
regularity modulus, condition (H2) and degenerate swallowing-sign bad
peaks, and the scrambling condition along deep raises
(Problems~\ref{prob:II-Minf}--\ref{prob:II-SCdeep}).
\end{enumerate}
Section~\ref{sec:II-status} states these problems precisely.  Nothing
proved here or in Part~I points to a counterexample; equally, nothing
proved here settles any one of the items (i), (ii).

\subsection{Conventions}
The notation of Part~I is used without comment, in particular the forced
data $(\xi,q_0,a,w,F,z,\zh,e,\nu)$ of a first row [I,~Proposition~1.9],
two-piece data and the coefficient $\Gamma_w$ [I,~Definition~7.7], and the
signature-ladder notation of [I,~Section~8].  References of the form
[I,~Theorem~8.20] are to Part~I \cite{PartI}.  Constants $C_f$ depend on the
first row (and on the mate where said) but not on the level, the
sub-window or the scale.  Statements that are only sketched or heuristic
in our sources appear only inside remarks that say so explicitly
(``sketch'', ``heuristic'', ``numerical evidence''); every lemma,
proposition, theorem and corollary is given with a complete proof.  Where an earlier
version of an argument was found to be wrong, the corrected statement is
given and the failure is recorded in a remark (for instance
Remarks~\ref{rem:II-corrections} and~\ref{rem:II-notMTIV}).


%%%%%%%%%%%%%%%%%%%%%%%%%%%%%%%%%%%%%%%%%%%%%%%%%%%%%%%%%%%%%%%%%%%%%%%%

\section{Companions and transfer}\label{sec:II-companions}

This section collects the tools of Part~II that are valid for every
admissible operator $T$; where a statement uses more (a property of the
design, a diagonal base, one block) this is said in the statement.
Throughout, $T$ is admissible, the block set $I$ is finite and $p$ is the
norm of [I,~Section~1] built with $I$ (so $p=p_N$ when $I=\{1,\dots,N\}$).
Forced data, the sets $F,K,J$, the block data $\zeta_m=R_m^{**}\xi$,
$\sigma_m$, $M_m$, $C_m$, $P_m$, $Q_m$, $\alpha_m$, $\gap_m$, $\mu_{k,m}$,
two-piece data, $\Gamma_w$, $\kappa_w$ and $\Delta d_m$ are those of
[I,~Proposition~1.9, Lemma~1.7, Definition~7.7].  We also write
\[
 \hat\zeta_m:=R_m^{**}\zh=\zeta_m/q_0,\qquad |\hat\zeta_m|_m=\sigma_m/q_0,
 \qquad \psi_m:=R_m^*w_m\in\ell_1 ,
\]
and $\phi_z(x):=|x|-zx$ for $|z|\le1$, $x\in\R$, as in [I,~Lemma~8.23].

The section is organized as follows.  Subsection~\ref{subsec:II-comp}
introduces first rows with prescribed data and companions.
Subsection~\ref{subsec:II-threshold} gives the block threshold as the
unique root of one explicit equation, with quantitative consequences and
the identities behind the block scalar $\varkappa$.
Subsection~\ref{subsec:II-cost} bounds the cost of moving the normer and
proves the raise-transfer lemma.  Subsection~\ref{subsec:II-U} proves the
uniform one-sided transfer expansion along convergent first rows, and
Subsection~\ref{subsec:II-E} the windowed recovery theorems through nearby
first rows.  Subsection~\ref{subsec:II-rigid} studies exact two-piece
data: the exactness certificate, the shift bound, the rigidity of two-piece
data and a criterion for the failure of lower semicontinuity of the fibre
map along block-tame rows.  Subsection~\ref{subsec:II-VT} shows that the engineered approximants
of [I,~Section~7] tolerate small violations of the sign conditions of
two-piece data at a single stage.

\subsection{First rows with prescribed data; companions}\label{subsec:II-comp}

Let $a'\in\ell_1$ with $q^*(a')=1$ and $z'\in B_{\ell_\infty}$ with
$z'_j=\sgn a'_j$ for $j\in\supp a'$.  By [I,~Remark~8.22(c)] there is a
unique $f'\in S_{p^*}$ whose forced data are
\begin{gather*}
 e':=U^*a'/\|U^*a'\|_H,\qquad \zh':=z'+Ue',\qquad
 q'_0:=\Big(1+\sum_{m\in I}|R_m^{**}\zh'|_m\Big)^{-1},\\
 \xi':=q'_0\zh',\qquad w':=J_V(L^{**}\xi'),\qquad f'=a'+L^*w' ;
\end{gather*}
we denote it by $f[a',z']$.  Since each $J_m$ is positively homogeneous of
degree $0$, $w'_m=J_m(R_m^{**}\zh')$ depends on $\zh'$ only.  If $a'_n\to a'$
in $\ell_1$ and $z'_n\to z'$ coordinatewise (with the same constraints),
then $f[a'_n,z'_n]\to f[a',z']$ in norm [I,~Remark~8.22(c)].  The row
$f[a',z']$ attains its norm if and only if $a'\in c_{00}$ and $z'\in c_0$
[I,~Proposition~1.10(c)]; in general it does not.

\begin{definition}[Companions]\label{def:II-companion}
Let $f\in S_{p^*}$ have forced data $(a,z,\dots)$.
\begin{enumerate}
\item[(a)] A \emph{companion} of $f$ is a row $f^\#:=f[a,z^\#]$ with
$z^\#\in B_{\ell_\infty}$ and $z^\#=z$ on $F$.  We put
$\delta:=z^\#-z$; then $e^\#=e$, $\nu^\#=\nu$ and $\zh^\#=\zh+\delta$, and
$\delta$ vanishes on $F$.
\item[(b)] A \emph{raise} is a vector $\Delta a\in\ell_1$ with
$\supp\Delta a\subseteq F$ and $\sgn(a_j)\Delta a_j\ge0$ for $j\in F$.  The
\emph{raise companion} is $f^\#:=f[a^\#,z]$ with $\lambda:=q^*(a+\Delta a)$
and $a^\#:=(a+\Delta a)/\lambda$.
\end{enumerate}
Forced data of a companion carry the superscript $\#$; for instance
$w^\#_m=J_m(R_m^{**}\zh^\#)$, $C^\#_m$, $P^\#_m$, $Q^\#_m$, $\gap^\#_m$,
$d^\#_m(\omega)=\ip{D_mw^\#_m}{D_m\omega}/C^\#_m$.
\end{definition}

The admissibility of the data in (b) follows from $\supp a^\#=F$ and
$\sgn a^\#_j=\sgn a_j=z_j$ on $F$.  Companions in the sense of (a) are the
rows used to \emph{exactify} data (close rooms, flip or raise contacts);
raise companions are used to enlarge the base part on its own support.
In both cases the contact data off the moved coordinates are those of $f$.

\begin{lemma}[$d$-coefficients through the normer]\label{lem:II-dnormer}
Let $f\in S_{p^*}$ and $m\in I$.
\begin{enumerate}
\item[(a)] For every finitely supported $\omega$ with $\supp\omega\subseteq
Q_m$,
\[
 d_m(\omega)=\frac{(R_m^*\omega)(\zh)}{|R_m^{**}\zh|_m}
 =\frac{\ip{\omega}{\hat\zeta_m}}{|\hat\zeta_m|_m}.
\]
\item[(b)] If $f^\#$ is a companion of $f$ \textup{(}Definition~\textup{\ref{def:II-companion}(a))} and $\omega$ is finitely supported
in $Q_m\cap Q^\#_m$, then
\[
 |R_m^{**}\zh^\#|_m\,d^\#_m(\omega)-|R_m^{**}\zh|_m\,d_m(\omega)
 =(R_m^*\omega)(\delta).
\]
\end{enumerate}
\end{lemma}

\begin{proof}
(a) Off $P_m$, [I,~Lemma~1.7] gives $\Phi_m(k)^2w_m(k)/C_m=
\zeta_m(k)/\sigma_m$.  Hence
$d_m(\omega)=\sum_k\Phi_m(k)^2w_m(k)\omega(k)/C_m=\ip{\omega}{\zeta_m}/\sigma_m
=\ip{\omega}{\hat\zeta_m}/|\hat\zeta_m|_m$, and
$\ip{\omega}{R_m^{**}\zh}=(R_m^*\omega)(\zh)$.
(b) Apply (a) at $f$ and at $f^\#$ and subtract; $\zh^\#-\zh=\delta$.
\end{proof}

\begin{corollary}[Moving contacts converts a first-order defect into a $d$-mismatch]\label{cor:II-defect}
Let $f^\#$ be a companion of $f$, let $\omega$ be finitely supported in
$Q_m\cap Q^\#_m$ and put $V:=R_m^*\omega\in\ell_1$.  Suppose that for every
$j\in\supp\delta$ either $V_j=0$ or $z^\#_j=\sgn V_j$.  Then
\[
 |R_m^{**}\zh^\#|_m\,d^\#_m(\omega)-|R_m^{**}\zh|_m\,d_m(\omega)
 =\sum_{j\in\supp\delta}\phi_{z_j}(V_j)\ \ge0 .
\]
\end{corollary}

\begin{proof}
By Lemma~\ref{lem:II-dnormer}(b) the left side is $V(\delta)=\sum_{j\in
\supp\delta}V_j(z^\#_j-z_j)$, and $V_j(z^\#_j-z_j)=|V_j|-z_jV_j=
\phi_{z_j}(V_j)$ when $z^\#_j=\sgn V_j$; both sides vanish when $V_j=0$.
\end{proof}

Applied to a switching $\omega=\omega^--\omega^+$, the corollary says: if
the switching is $d$-neutral at $f$ and a companion makes the coordinates
it moves exact contacts with the sign of $V$, then at the companion the
switching has a $d$-mismatch equal to its first-order defect
$\sum_{j\in\supp\delta}\phi_{z_j}(V_j)$, divided by $|R_m^{**}\zh^\#|_m$.
Only the defect on $\supp\delta$ is converted; nothing is said about
the coordinates that are not moved.

\begin{lemma}[Approximate resonance and $d$-neutrality]\label{lem:II-resonance}
Let $u\in\ell_1$, $\varepsilon\in\{\pm1\}$, and let $f^\#$ be a companion
of $f$ such that $\varepsilon u_jz^\#_j=|u_j|$ for every $j\in\supp
u\setminus F$.  Then
\[
 \varepsilon u(\zh^\#)-\varepsilon u(\zh)=\mathrm{Re}(u):=\sum_{j\notin F}
 \phi_{z_j}(\varepsilon u_j)\ \ge0 .
\]
Consequently, if $u=u_{k,m}$, $k\in Q_m$ with $w_m(k)=0$ \textup{(}a
$d$-neutral strict non-peak of $f$\textup{)}, $\mathrm{Re}(u)>0$, and $k$ is
a strict non-peak of $f^\#$, then
\[
 \varepsilon\,\frac{\Phi_m(k)w^\#_m(k)}{mC^\#_m}=
 \frac{\mathrm{Re}(u)}{|R_m^{**}\zh^\#|_m}>0 ;
\]
in particular $k$ is not $d$-neutral at $f^\#$.
\end{lemma}

\begin{proof}
$\zh^\#-\zh=\delta$ vanishes on $F$, so $\varepsilon u(\zh^\#)-\varepsilon
u(\zh)=\sum_{j\notin F}\varepsilon u_j(z^\#_j-z_j)$, and
$\varepsilon u_j(z^\#_j-z_j)=|u_j|-\varepsilon u_jz_j=\phi_{z_j}(\varepsilon
u_j)$ for $j\in\supp u\setminus F$ (both vanish off $\supp u$).  If
$w_m(k)=0$, then $u(\zh)=0$: off $P_m$, $w_m(k)=C_m\zeta_m(k)/(\Phi_m(k)^2
\sigma_m)$ by [I,~Lemma~1.7].  At the strict non-peak $k$ of $f^\#$ the same
formula gives
\[
 \frac{\Phi_m(k)w^\#_m(k)}{C^\#_m}=\frac{\lambda_{k,m}u(\zh^\#)}{\Phi_m(k)
 |R_m^{**}\zh^\#|_m}=\frac{m\,u(\zh^\#)}{|R_m^{**}\zh^\#|_m},
\]
and $\varepsilon u(\zh^\#)=\mathrm{Re}(u)$.
\end{proof}

\begin{remark}\label{rem:II-resonance}
(a) If all switching carriers of a block are of the kind described in
Lemma~\ref{lem:II-resonance}, they all acquire weights of the same sign at
$f^\#$, and the cone of exact $d$-neutral switchings at $f^\#$ meets the
nonnegative orthant only in $0$.  A Hoffman constant of that cone is then
at least $|R_m^{**}\zh^\#|_m/\min_l\mathrm{Re}(u_l)$ (test with a unit
vector).  Since $\mathrm{Re}(u_l)$ is the room of the signature part plus
the contribution of the target part, this is of order $1/(\text{room})$
only when the target part is comparable with the room.

(b) By Corollary~\ref{cor:II-defect} a companion produces nonnegative
$d$-mismatches, which [I,~Corollary~7.21] accepts.  However, finitely
supported two-piece data at a row with nonzero mismatches satisfy, off the
supports $\Sigma(\omega)$ of the block parts, $b^+-b^-=-W_\Delta=-\sum_m
\Delta d_m\psi_m$ (Lemma~\ref{lem:II-sandwich} below), and $\psi_m$ has
infinitely many peak terms; so the mismatch can be carried only where
$z_jW_\Delta(j)\le0$ at contacts $j\notin\Sigma(\omega)$ and $W_\Delta(j)=0$
at free coordinates $j\notin\Sigma(\omega)$.  This is the content of [I,~Remark~8.45(b)].
\end{remark}

\subsection{The threshold equation}\label{subsec:II-threshold}

In this subsection we fix one block and drop its index: $\Phi$ is a
sequence of positive numbers with $\|\Phi\|_1<1$, $D$ the diagonal
operator with entries $\Phi(k)$, $|\cdot|$ the norm on $\ell_1$ with unit
ball $B_{\ell_1}+D(B_{\ell_2})$ and $N(w)=\|w\|_\infty+\|Dw\|_2$ its dual
norm [I,~Remark~1.1].  For $\zeta\in\ell_1\setminus\{0\}$ let $w=J(\zeta)$,
$M$, $C$, $P$, $Q:=\N\setminus P$ and $\alpha$ be as in [I,~Lemma~1.7], and put
\begin{gather*}
 \bar\vartheta(\zeta):=\frac{|\zeta|M}{C},\qquad
 \mathsf v_k(\zeta):=\frac{|\zeta(k)|}{\Phi(k)^2},\qquad
 \varrho_k(\zeta):=\frac{\mathsf v_k(\zeta)}{\bar\vartheta(\zeta)},\qquad
 \Phi_P^2:=\sum_{k\in P}\Phi(k)^2,\\
 \mathsf A(x;\zeta):=\sum_k\big(|\zeta(k)|-x\Phi(k)^2\big)_+,\qquad
 \mathsf B(x;\zeta):=\sum_k\Phi(k)^2\min\big(x,\mathsf v_k(\zeta)\big)^2,\\
 \Xi(x;\zeta):=\mathsf A(x;\zeta)^2-\mathsf B(x;\zeta)\qquad(x>0).
\end{gather*}
The series converge: $(|\zeta(k)|-x\Phi(k)^2)_+\le|\zeta(k)|$ and
$\Phi(k)^2\min(x,\mathsf v_k)^2\le x\Phi(k)^2\mathsf v_k=x|\zeta(k)|$.  All
these quantities except $\varrho_k$ are positively homogeneous of degree
one in $\zeta$ (and $\varrho_k$, $M$, $C$, $P$ of degree zero).  For the
block vectors of a first row, $\zeta=R_m^{**}\xi$, $\lambda_{k,m}=m\Phi_m(k)$
and $\mathsf v_k(\zeta)\ge\bar\vartheta(\zeta)$ iff $|\xi(u_{k,m})|\ge
\vartheta_m\Phi_m(k)$; thus $\bar\vartheta(\zeta_m)=m\vartheta_m$ with the
threshold constant $\vartheta_m$ of [I,~Section~1], and for
$\hat\zeta_m=R_m^{**}\zh$
\[
 \begin{gathered}
 \mu_{k,m}=\frac{q_0\Phi_m(k)}{m}\big(\mathsf v_k(\hat\zeta_m)-\bar\vartheta
 (\hat\zeta_m)\big)\quad(k\in P_m),\\
 \gap_m(k)=M_m\big(1-\varrho_k(\hat\zeta_m)\big)\quad(k\in Q_m)
 \end{gathered}
\]
(the second identity is the clamp formula in part (b) below).  We call
$\varrho_k$ the \emph{relative position} of $k$: $k$ is a peak iff
$\varrho_k\ge1$, a degenerate peak iff $\varrho_k=1$.

\begin{lemma}[Threshold equation]\label{lem:II-T}
Let $\zeta\in\ell_1\setminus\{0\}$.
\begin{enumerate}
\item[(a)] For every $\zeta'\in\ell_1\setminus\{0\}$ the function
$\Xi(\cdot;\zeta')$ has exactly one zero in $(0,\infty)$; it is positive to
the left and negative to the right of the zero.
\item[(b)] The zero of $\Xi(\cdot;\zeta)$ is $\bar\vartheta(\zeta)$.  Writing
$\bar\vartheta:=\bar\vartheta(\zeta)$ and $\mathsf v_k:=\mathsf v_k(\zeta)$:
\[
 \begin{gathered}
 |\zeta|=\mathsf A(\bar\vartheta;\zeta),\quad |\zeta|^2=\mathsf B(\bar\vartheta;
 \zeta),\quad M=\frac{\bar\vartheta}{\bar\vartheta+|\zeta|},\quad
 C=\frac{|\zeta|}{\bar\vartheta+|\zeta|},\\
 w(k)=\sgn(\zeta(k))\,M\min\Big(1,\frac{\mathsf v_k}{\bar\vartheta}\Big),
 \end{gathered}
\]
$P=\{k:\mathsf v_k\ge\bar\vartheta\}$, and $|\zeta|\,|\alpha(k)|=|\zeta(k)|-
\bar\vartheta\Phi(k)^2$ for $k\in P$; in particular $k\in P$ is a
degenerate peak iff $\mathsf v_k=\bar\vartheta$.
\item[(c)] \textup{(threshold identities)}
\[
 \sum_{k\in P}\Phi(k)^2(\mathsf v_k-\bar\vartheta)=|\zeta|,\qquad
 |\zeta|^2=\bar\vartheta^2\Phi_P^2+\sum_{k\in Q}\mathsf v_k^2\Phi(k)^2 .
\]
\item[(d)] For $\zeta'\in\ell_1\setminus\{0\}$: if $\Xi(\bar\vartheta(\zeta);
\zeta')>0$ then $\bar\vartheta(\zeta')>\bar\vartheta(\zeta)$, and if
$\Xi(\bar\vartheta(\zeta);\zeta')<0$ then $\bar\vartheta(\zeta')<\bar\vartheta
(\zeta)$.
\item[(e)] \textup{(raising moves)} Let $k'\in\N$ and let $\zeta'$ agree
with $\zeta$ except at $k'$, where $\sgn\zeta'(k')=\sgn\zeta(k')$ or
$\zeta'(k')=0$.  Then $\bar\vartheta(\zeta')>\bar\vartheta(\zeta)$ in each of the
cases: \textup{(i)} $k'\in P$ and $|\zeta'(k')|>|\zeta(k')|$;
\textup{(ii)} $k'\notin P$ and $|\zeta'(k')|<|\zeta(k')|$;
\textup{(iii)} $k'$ is a degenerate peak and $|\zeta'(k')|\ne|\zeta(k')|$.
\item[(f)] \textup{(perturbed raising moves)} Let $k'\in\N$,
$\zeta'\in\ell_1\setminus\{0\}$, $\Delta:=\big||\zeta'(k')|-|\zeta(k')|\big|>0$
and $E:=\sum_{k\ne k'}|\zeta'(k)-\zeta(k)|$.  Then $\bar\vartheta(\zeta')>
\bar\vartheta(\zeta)$ if either \textup{(i')} $k'\in P$,
$|\zeta'(k')|>|\zeta(k')|$ and $E<|\zeta|\Delta/(|\zeta|+\bar\vartheta)$; or
\textup{(ii')} $\mathsf v_{k'}>0$, $\mathsf v_{k'}\le\bar\vartheta$,
$|\zeta'(k')|<|\zeta(k')|$ and $E<\mathsf v_{k'}\Delta/(2(|\zeta|+
\bar\vartheta))$.
\end{enumerate}
\end{lemma}

\begin{proof}
(a) Fix $\zeta'$ and write $\mathsf A',\mathsf B',\Xi',\mathsf v'$.  By
dominated convergence $\mathsf A'$ and $\mathsf B'$ are continuous on
$(0,\infty)$; $\mathsf A'$ is nonincreasing and $\mathsf B'$ nondecreasing.
Let $s^*:=\sup_k\mathsf v'_k\in(0,\infty]$.  For $x<s^*$ the set
$\{k:\mathsf v'_k>x\}$ is nonempty; each of its terms of $\mathsf A'$ is
positive and strictly decreasing and each of its terms of $\mathsf B'$ is
strictly increasing near $x$, so $\Xi'$ is strictly decreasing on
$(0,s^*)$.  As $x\to0^+$, $\mathsf A'(x)\to\|\zeta'\|_1>0$ and $0\le\mathsf
B'(x)\le x\|\zeta'\|_1\to0$, so $\Xi'>0$ near $0$.  For $x\ge s^*$ (if
$s^*<\infty$), $\mathsf A'(x)=0$ and $\Xi'(x)=-\mathsf B'(x)<0$.  If
$s^*=\infty$, then $\mathsf A'(x)\to0$ as $x\to\infty$ by dominated
convergence, while $\mathsf B'(x)\ge\mathsf B'(1)>0$ for $x\ge1$, so $\Xi'<0$
for large $x$.  By continuity and strict monotonicity on $(0,s^*)$,
$\Xi'$ has exactly one zero, which lies in $(0,s^*)$, and the sign
pattern is as stated.

(b) Let $\bar\vartheta>0$ be the zero of $\Xi(\cdot;\zeta)$, put $\mathsf
a:=\mathsf A(\bar\vartheta;\zeta)$ and $s_k:=\sgn\zeta(k)$.  If $\mathsf a=0$
then $\mathsf B(\bar\vartheta;\zeta)=0$, which forces $\zeta=0$; so $\mathsf
a>0$, and some $k$ has $\mathsf v_k>\bar\vartheta$.  \emph{Primal}: put
$x_k:=s_k(|\zeta(k)|-\bar\vartheta\Phi(k)^2)_+$ and $\beta_k:=s_k\Phi(k)
\min(\mathsf v_k,\bar\vartheta)$.  Checking the cases $\mathsf v_k\ge
\bar\vartheta$ and $\mathsf v_k<\bar\vartheta$ separately gives
$\zeta=x+D\beta$, and $\|x\|_1=\mathsf a$, $\|\beta\|_2^2=\mathsf
B(\bar\vartheta;\zeta)=\mathsf a^2$.  Hence $\zeta/\mathsf a\in B_{\ell_1}+
D(B_{\ell_2})$, i.e.\ $|\zeta|\le\mathsf a$.  \emph{Dual}: put
$v_k:=s_k\min(\bar\vartheta,\mathsf v_k)$.  Then $\|v\|_\infty=\bar\vartheta$,
$\|Dv\|_2^2=\mathsf B(\bar\vartheta;\zeta)=\mathsf a^2$, so $N(v)=\bar\vartheta+
\mathsf a$, and, using $\sum_{\mathsf v_k\ge\bar\vartheta}|\zeta(k)|=\mathsf
a+\bar\vartheta\sum_{\mathsf v_k\ge\bar\vartheta}\Phi(k)^2$ and $\mathsf
v_k|\zeta(k)|=\Phi(k)^2\mathsf v_k^2$,
\[
 \ip{v}{\zeta}=\bar\vartheta\sum_{\mathsf v_k\ge\bar\vartheta}|\zeta(k)|+
 \sum_{\mathsf v_k<\bar\vartheta}\Phi(k)^2\mathsf v_k^2=\bar\vartheta\mathsf a
 +\mathsf B(\bar\vartheta;\zeta)=\mathsf a(\bar\vartheta+\mathsf a).
\]
So $w':=v/(\bar\vartheta+\mathsf a)$ has $N(w')=1$ and $\ip{w'}{\zeta}=\mathsf
a\ge|\zeta|\ge\ip{w'}{\zeta}$.  Hence $|\zeta|=\mathsf a$, $w'$ norms $\zeta$,
and $w'=J(\zeta)=w$ because $|\cdot|$ is smooth [I,~Remark~1.1(d)].  This
gives the formulas for $|\zeta|$, $|\zeta|^2$, $w$, $M=\|w\|_\infty=
\bar\vartheta/(\bar\vartheta+|\zeta|)$, $C=\|Dw\|_2=|\zeta|/(\bar\vartheta+
|\zeta|)$, and $|\zeta|M/C=\bar\vartheta$; $P=\{|w(k)|=M\}=\{\mathsf v_k\ge
\bar\vartheta\}$; finally $\alpha=\zeta/|\zeta|-D^2w/C=x/|\zeta|$, which is the
last claim.

(c) The first identity is $\mathsf A(\bar\vartheta;\zeta)=|\zeta|$, the second
$\mathsf B(\bar\vartheta;\zeta)=|\zeta|^2$, both read with (b).

(d) By (a) for $\zeta'$: $\Xi(x;\zeta')>0$ iff $x<\bar\vartheta(\zeta')$, and
$\Xi(x;\zeta')<0$ iff $x>\bar\vartheta(\zeta')$.

(e) Put $\bar\vartheta:=\bar\vartheta(\zeta)$; $\Xi(\bar\vartheta;\zeta)=0$.  In
case (i) the $k'$-term of $\mathsf A(\bar\vartheta;\cdot)$ increases by
$|\zeta'(k')|-|\zeta(k')|>0$ and the $k'$-term of $\mathsf B(\bar\vartheta;
\cdot)$ is $\bar\vartheta^2\Phi(k')^2$ before and after; so
$\Xi(\bar\vartheta;\zeta')>0$.  In case (ii), $\mathsf v'_{k'}<\mathsf
v_{k'}<\bar\vartheta$: the $k'$-terms of $\mathsf A$ vanish before and after
and the $k'$-term of $\mathsf B$ decreases; so $\Xi(\bar\vartheta;\zeta')>0$.
Case (iii) is case (i) for an increase and is computed as case (ii) for a
decrease.  Conclude with (d).

(f) For $k\ne k'$ the $k$-term of $\mathsf A(\bar\vartheta;\cdot)$ is
$1$-Lipschitz in $|\zeta(k)|$, and the $k$-term of $\mathsf B(\bar\vartheta;
\cdot)$ is $2\bar\vartheta$-Lipschitz in $|\zeta(k)|$ (since
$|\min(x,a)^2-\min(x,b)^2|\le2x|a-b|$ for $a,b\ge0$).  Write $\mathsf
a:=|\zeta|$.  (i') $\mathsf A(\bar\vartheta;\zeta')\ge\mathsf a+\Delta-E>0$
and $\mathsf B(\bar\vartheta;\zeta')\le\mathsf a^2+2\bar\vartheta E$, so
$\Xi(\bar\vartheta;\zeta')\ge2\mathsf a(\Delta-E)-2\bar\vartheta E>0$ by the
hypothesis on $E$.  (ii') The $k'$-term of $\mathsf A(\bar\vartheta;\cdot)$
vanishes before and after.  Since $\mathsf v'_{k'}<\mathsf v_{k'}\le
\bar\vartheta$, the $k'$-term of $\mathsf B$ drops by $\Phi(k')^2(\mathsf
v_{k'}^2-\mathsf v_{k'}'^2)\ge\Phi(k')^2\mathsf v_{k'}(\mathsf v_{k'}-\mathsf
v'_{k'})=\mathsf v_{k'}\Delta$.  Hence, as $\mathsf A(\bar\vartheta;\zeta')^2
\ge(\mathsf a-E)_+^2\ge\mathsf a^2-2\mathsf aE$,
$\Xi(\bar\vartheta;\zeta')\ge\mathsf v_{k'}\Delta-2(\mathsf a+\bar\vartheta)E>0$.
Conclude with (d).
\end{proof}

Thus moving block mass toward the peaks (more at a peak, less at a strict
non-peak) raises the threshold, and the reverse moves lower it; at a
degenerate peak $\Xi$ has a concave kink and both moves raise it.  The
following quantitative version is used for the threshold buffers of
companions.

\begin{lemma}[Lipschitz bound and quantitative raising]\label{lem:II-T2}
Let $\zeta\in\ell_1\setminus\{0\}$, $\bar\vartheta:=\bar\vartheta(\zeta)$,
$\mathsf a:=|\zeta|$, $\phi:=\|\Phi\|_2^2$; let $k_0$ be a peak with $\mathsf
v_{k_0}>\bar\vartheta$ \textup{(}it exists, since $\|\alpha\|_1=1$\textup{)},
and put $\phi_0:=\Phi(k_0)^2$, $h_0:=\min(1,\mathsf v_{k_0}-\bar\vartheta)$,
$C_1:=(1+2(\bar\vartheta+1)/\mathsf a)/\phi_0$.
\begin{enumerate}
\item[(a)] If $\zeta'\in\ell_1$ and $E:=\|\zeta'-\zeta\|_1$ satisfies
$C_1E<\min(h_0,\bar\vartheta)$, then $\zeta'\ne0$ and
$|\bar\vartheta(\zeta')-\bar\vartheta|\le C_1E$.  If only $C_1E<h_0$, then
$\bar\vartheta(\zeta')\le\bar\vartheta+C_1E$.
\item[(b)] Let $c$ be a peak \textup{(}degenerate or not\textup{)}, $s>0$,
and let $\zeta'\in\ell_1$ satisfy $|\zeta'(c)|=|\zeta(c)|+s$,
$\sgn\zeta'(c)=\sgn\zeta(c)$ and $E:=\sum_{k\ne c}|\zeta'(k)-\zeta(k)|\le
\mathsf as/(8(\mathsf a+\bar\vartheta+1))$.  Then
\[
 \bar\vartheta(\zeta')-\bar\vartheta\ge\min\Big\{1,\frac{\mathsf a s}{8\phi
 (\mathsf a+\bar\vartheta+1)}\Big\},
\]
and $\bar\vartheta(\zeta')-\bar\vartheta\le C_1(s+E)$ if $C_1(s+E)<h_0$.
\end{enumerate}
\end{lemma}

\begin{proof}
For $h\ge0$ and every $\zeta''$: $\mathsf A(x;\cdot)$ changes by at most $E''$
and $\mathsf B(x;\cdot)$ by at most $2xE''$ when $\zeta''$ moves by $E''$ in
$\ell_1$ (as in the proof of Lemma~\ref{lem:II-T}(f)); $\mathsf
B(\cdot;\zeta)$ is nondecreasing; and, by Lemma~\ref{lem:II-T}(b),
$\mathsf A(\bar\vartheta;\zeta)=\mathsf a$, $\mathsf B(\bar\vartheta;\zeta)=\mathsf
a^2$.

(a) $\zeta'\ne0$ because $E<\phi_0h_0\le|\zeta(k_0)|$.  \emph{Upper bound.}
Let $h:=C_1E+\varepsilon\le h_0$ with $\varepsilon>0$.  Since $\mathsf
v_{k_0}\ge\bar\vartheta+h$, the terms of $\mathsf A(\bar\vartheta+h;\zeta)$ with
$\mathsf v_k\ge\bar\vartheta+h$ are those of $\mathsf A(\bar\vartheta;\zeta)$
minus $h\Phi(k)^2$, and the other terms vanish; so $\mathsf
A(\bar\vartheta+h;\zeta)\le\mathsf a-h\phi_0$.  Put $y:=h\phi_0-E$; then
$0<y\le\mathsf a$ (as $C_1\phi_0>1$, and $h\phi_0\le\Phi(k_0)^2(\mathsf
v_{k_0}-\bar\vartheta)$, which is one term of $\mathsf
A(\bar\vartheta;\zeta)$).  Thus $0\le\mathsf A(\bar\vartheta+h;\zeta')\le\mathsf
a-y$ and $\mathsf B(\bar\vartheta+h;\zeta')\ge\mathsf a^2-2(\bar\vartheta+1)E$,
whence
\[
 \Xi(\bar\vartheta+h;\zeta')\le(\mathsf a-y)^2-\mathsf a^2+2(\bar\vartheta+1)E
 \le-y\mathsf a+2(\bar\vartheta+1)E<0,
\]
because $y\mathsf a>(C_1\phi_0-1)E\mathsf a=2(\bar\vartheta+1)E$.  So
$\bar\vartheta(\zeta')<\bar\vartheta+h$ by Lemma~\ref{lem:II-T}(a).
\emph{Lower bound} (when $C_1E<\bar\vartheta$).  Let $h:=C_1E+\varepsilon<
\bar\vartheta$.  The terms of $\mathsf A(\bar\vartheta-h;\zeta)$ with $\mathsf
v_k\ge\bar\vartheta$ exceed those of $\mathsf A(\bar\vartheta;\zeta)$ by
$h\Phi(k)^2$, so $\mathsf A(\bar\vartheta-h;\zeta')\ge\mathsf a+h\phi_0-E\ge0$,
and $\mathsf B(\bar\vartheta-h;\zeta')\le\mathsf a^2+2\bar\vartheta E$.  Hence
$\Xi(\bar\vartheta-h;\zeta')\ge2\mathsf a(h\phi_0-E)-2\bar\vartheta E>0$, since
$h\phi_0>E(1+\bar\vartheta/\mathsf a)$.  So $\bar\vartheta(\zeta')>\bar\vartheta-h$.
Let $\varepsilon\to0$.

(b) \emph{Lower bound.}  Let $h:=\min\{1,\mathsf as/(8\phi(\mathsf
a+\bar\vartheta+1))\}$; then $h\le s/\phi\le s/\Phi(c)^2$ and $h\phi\le s/8$,
and $E\le s/8$.  Since $\mathsf v_c(\zeta')=\mathsf v_c+s/\Phi(c)^2\ge
\bar\vartheta+h$, the $c$-term of $\mathsf A(\bar\vartheta+h;\zeta')$ is the
$c$-term of $\mathsf A(\bar\vartheta;\zeta)$ plus $s-h\Phi(c)^2$, while every
other term of $\mathsf A(\bar\vartheta+h;\zeta)$ is at least the
corresponding term of $\mathsf A(\bar\vartheta;\zeta)$ minus $h\Phi(k)^2$;
so $\mathsf A(\bar\vartheta+h;\zeta')\ge\mathsf a+s-h\phi-E$.  The $c$-term of
$\mathsf B(\bar\vartheta+h;\zeta')$ is $(\bar\vartheta+h)^2\Phi(c)^2$, that of
$\mathsf B(\bar\vartheta;\zeta)$ is $\bar\vartheta^2\Phi(c)^2$, and
$\min(\bar\vartheta+h,x)^2\le\min(\bar\vartheta,x)^2+2\bar\vartheta h+h^2$; with
$h\le1$ this gives $\mathsf B(\bar\vartheta+h;\zeta')\le\mathsf
a^2+(2\bar\vartheta+1)h\phi+2(\bar\vartheta+1)E$.  As $s-h\phi-E\ge3s/4>0$,
\[
 \Xi(\bar\vartheta+h;\zeta')\ge2\mathsf as-2(\mathsf a+\bar\vartheta+1)E
 -2(\mathsf a+\bar\vartheta+1)h\phi\ge2\mathsf as-\tfrac14\mathsf as-
 \tfrac14\mathsf as>0,
\]
so $\bar\vartheta(\zeta')>\bar\vartheta+h$.  \emph{Upper bound}: part (a)
(upper bound) with $\|\zeta'-\zeta\|_1=s+E$.
\end{proof}

\begin{corollary}[Common rescaling under a raise]\label{cor:II-T3}
In the situation of Lemma~\textup{\ref{lem:II-T2}(b)}, for every $k\ne c$,
\[
 \varrho_k(\zeta')=\frac{|\zeta'(k)|}{|\zeta(k)|}\,\varrho_k(\zeta)\,
 \frac{\bar\vartheta(\zeta)}{\bar\vartheta(\zeta')}\qquad(\zeta(k)\ne0);
\]
so every coordinate $k$ with $\zeta'(k)=\zeta(k)$ has
$\varrho_k(\zeta')=\varrho_k(\zeta)/(1+\delta)$ with
$\delta:=\bar\vartheta(\zeta')/\bar\vartheta(\zeta)-1\ge c_{\rm raise}\,s$ for
$s\le8\phi(\mathsf a+\bar\vartheta+1)/\mathsf a$, where $c_{\rm raise}:=
\mathsf a/(8\phi\bar\vartheta(\mathsf a+\bar\vartheta+1))$.  Moreover, for a
first row $f$, a block $m$ and a companion $f^\#$ of $f$ such that
$\hat\zeta^\#_m:=R_m^{**}\zh^\#$ and $\hat\zeta_m$ agree on a set $G\subseteq
Q_m\cap Q^\#_m$, every finitely supported $\omega$ with $\supp\omega\subseteq
G$ satisfies
\[
 d^\#_m(\omega)=\frac{|\hat\zeta_m|_m}{|\hat\zeta^\#_m|_m}\,d_m(\omega):
\]
all $d$-coefficients of untouched strict non-peaks are multiplied by one
common positive factor, and $d$-neutrality of a switching supported in $G$
is preserved.
\end{corollary}

\begin{proof}
$\varrho_k=|\zeta(k)|/(\Phi(k)^2\bar\vartheta)$, and Lemma~\ref{lem:II-T2}(b)
for the bound on $\delta$.  The last claim is Lemma~\ref{lem:II-dnormer}(a)
at $f$ and at $f^\#$: $d^\#_m(\omega)=\ip{\omega}{\hat\zeta^\#_m}/
|\hat\zeta^\#_m|_m$ and $\ip{\omega}{\hat\zeta^\#_m}=\ip{\omega}{\hat\zeta_m}$.
\end{proof}

\subsubsection*{The block scalar $\varkappa$}
For $\zeta\in\ell_1\setminus\{0\}$ and $\Omega\subseteq Q(\zeta)$ put
\[
 \varkappa(\zeta;\Omega):=\frac{|\zeta|\big(|\zeta|+\bar\vartheta\big)-
 \sum_{k\in\Omega}\mathsf v_k^2\Phi(k)^2}{\bar\vartheta}
 =|\zeta|+\bar\vartheta\Phi_P^2+\frac1{\bar\vartheta}\sum_{k\in Q\setminus\Omega}
 \mathsf v_k^2\Phi(k)^2 ,
\]
with $\bar\vartheta=\bar\vartheta(\zeta)$, $\mathsf v_k=\mathsf v_k(\zeta)$; the
two expressions agree by Lemma~\ref{lem:II-T}(c), and $\varkappa\ge|\zeta|>0$.

\begin{lemma}[The data identity]\label{lem:II-kappa}
Let $f'\in S_{p^*}$, $m\in I$, $\Omega_m\subseteq Q'_m$ finite, and let
$\omega^\pm_m$ be supported in $\Omega_m$.  Put $\Delta\omega:=\omega^-_m-
\omega^+_m$, $\Delta_m:=d'_m(\Delta\omega)$ and $\gamma_k:=\lambda_{k,m}
(\Delta\omega(k)-\Delta_mw'_m(k))$ for $k\in\Omega_m$.  Then, with
$\hat\zeta'_m:=R_m^{**}\zh'$,
\[
 \Delta_mM'_m\,\varkappa(\hat\zeta'_m;\Omega_m)=\sum_{k\in\Omega_m}
 u_{k,m}(\zh')\,\gamma_k .
\]
\end{lemma}

\begin{proof}
Write $\zeta:=\hat\zeta'_m$, $\mathsf a:=|\zeta|$, $\bar\vartheta:=
\bar\vartheta(\zeta)$, and drop $m$ and the primes.  For $k\in Q$,
[I,~Lemma~1.7] gives $\Phi(k)^2w(k)/C=\zeta(k)/\mathsf a$ and hence
$\Phi(k)^2w(k)^2/C=C\mathsf v_k^2\Phi(k)^2/\mathsf a^2$.  On $\Omega$,
$\Delta\omega(k)=\gamma_k/\lambda_k+\Delta w(k)$ and
$\zeta(k)/\lambda_k=u_k(\zh)$.  So
\[
 \Delta=d(\Delta\omega)=\sum_{k\in\Omega}\frac{\zeta(k)}{\mathsf a}
 \Big(\frac{\gamma_k}{\lambda_k}+\Delta w(k)\Big)=\frac1{\mathsf a}
 \sum_{k\in\Omega}u_k(\zh)\gamma_k+\Delta\frac C{\mathsf a^2}\sum_{k\in\Omega}
 \mathsf v_k^2\Phi(k)^2 .
\]
Multiply by $\mathsf a^2/C=\mathsf a(\mathsf a+\bar\vartheta)$
(Lemma~\ref{lem:II-T}(b)); since $\mathsf a/C=\mathsf a+\bar\vartheta$,
\[
 \Delta\Big(\mathsf a(\mathsf a+\bar\vartheta)-\sum_{k\in\Omega}\mathsf v_k^2
 \Phi(k)^2\Big)=(\mathsf a+\bar\vartheta)\sum_{k\in\Omega}u_k(\zh)\gamma_k .
\]
Divide by $\mathsf a+\bar\vartheta$ and use $1/(\mathsf a+\bar\vartheta)=
M/\bar\vartheta$.
\end{proof}

By [I,~(1)], for every carrier $(k,m)$ outside $\Omega_m$ the switching
contributes the forced coefficient $-\Delta_m\lambda_{k,m}w'_m(k)$ (see
Lemma~\ref{lem:II-cert}(a) below), and the lemma shows that the only
$f'$-dependent quantities of the block entering the $d$-row of an exact
switching through $\Omega_m$ are the values $u_{k,m}(\zh')$ ($k\in\Omega_m$),
$M'_m$ and the single scalar $\varkappa(\hat\zeta'_m;\Omega_m)$.

\begin{lemma}[Continuity and derivatives of $\varkappa$]\label{lem:II-kappader}
Let $\zeta\in\ell_1\setminus\{0\}$ and $\Omega\subseteq Q(\zeta)$ finite.
\begin{enumerate}
\item[(a)] $\varkappa(\cdot;\Omega)$ is continuous at $\zeta$ on the set of
$\zeta''\ne0$ with $\Omega\subseteq Q(\zeta'')$; in particular it does not
jump when a coordinate outside $\Omega$ crosses the threshold.
\item[(b)] Let $c$ be a coordinate and $\zeta_s$ the vector obtained from
$\zeta$ by replacing $|\zeta(c)|$ by $|\zeta(c)|+s$ \textup{(}same sign\textup{)}, for
$s$ in an open interval $\mathcal I\ni0$ on which $P(\zeta_s)$ is constant
and $\Omega\subseteq Q(\zeta_s)$.  Put $X:=C\sum_{k\in Q\setminus\Omega}
\mathsf v_k^2\Phi(k)^2/(\bar\vartheta^2\Phi_P^2)\ge0$.  Then
$\bar\vartheta$, $|\zeta_s|$ and $\varkappa$ are differentiable on $\mathcal I$,
and at $s=0$:
\begin{enumerate}
\item[(i)] if $c\in P$: $d\bar\vartheta/ds=C/\Phi_P^2$, $d|\zeta_s|/ds=M$,
$d\varkappa/ds=1-X$;
\item[(ii)] if $c\in Q\setminus\Omega$: $d\bar\vartheta/ds=-\varrho_cM/\Phi_P^2$,
$d|\zeta_s|/ds=\varrho_cM$, $d\varkappa/ds=\varrho_c(2+MX/C)$;
\item[(iii)] if $c\in\Omega$: $d\varkappa/ds=\varrho_cMX/C$.
\end{enumerate}
\end{enumerate}
\end{lemma}

\begin{proof}
(a) $|\cdot|$ is a norm, $\mathsf v_k$ ($k\in\Omega$) are continuous, and
$\bar\vartheta$ is continuous by Lemma~\ref{lem:II-T2}(a); use the first
expression of $\varkappa$, which does not involve $P$.  (In the second
expression a coordinate with $\mathsf v_k=\bar\vartheta$ contributes
$\bar\vartheta\Phi(k)^2$ both as a peak and as a strict non-peak.)
(b) With $P$ fixed, Lemma~\ref{lem:II-T}(c) reads
$\sum_P|\zeta_s(k)|-\bar\vartheta\Phi_P^2=|\zeta_s|$ and $|\zeta_s|^2=
\bar\vartheta^2\Phi_P^2+S_Q(s)$, $S_Q(s):=\sum_{k\in Q}\mathsf v_k(\zeta_s)^2
\Phi(k)^2$; eliminating $|\zeta_s|$ gives a quadratic equation for
$\bar\vartheta$ with coefficients affine or quadratic in $s$ and with nonzero
derivative in $\bar\vartheta$ at the root, so the implicit function theorem
gives differentiability.  (i) Differentiating, $ds-\Phi_P^2d\bar\vartheta=
d|\zeta|$ and $|\zeta|\,d|\zeta|=\bar\vartheta\Phi_P^2d\bar\vartheta$; hence
$d\bar\vartheta=|\zeta|ds/(\Phi_P^2(|\zeta|+\bar\vartheta))=Cds/\Phi_P^2$ and
$d|\zeta|=Mds$; with the second expression of $\varkappa$ (the sum over
$Q\setminus\Omega$ unchanged), $d\varkappa=d|\zeta|+\Phi_P^2d\bar\vartheta-
(S_{Q\setminus\Omega}/\bar\vartheta^2)\,d\bar\vartheta=Mds+Cds-XC^{-1}\Phi_P^2\cdot
Cds/\Phi_P^2=(1-X)ds$, where $S_{Q\setminus\Omega}:=\sum_{k\in Q\setminus
\Omega}\mathsf v_k^2\Phi(k)^2=X\bar\vartheta^2\Phi_P^2/C$.  (ii) Now $-\Phi_P^2d\bar\vartheta=d|\zeta|$ and
$|\zeta|d|\zeta|=\bar\vartheta\Phi_P^2d\bar\vartheta+\mathsf v_cds$ (as
$d(\mathsf v_c^2\Phi(c)^2)=2\mathsf v_cds$); hence $d\bar\vartheta=-\mathsf
v_cds/(\Phi_P^2(|\zeta|+\bar\vartheta))=-\varrho_cMds/\Phi_P^2$,
$d|\zeta|=\varrho_cMds$, and $d\varkappa=d|\zeta|+\Phi_P^2d\bar\vartheta+
2\mathsf v_cds/\bar\vartheta-(S_{Q\setminus\Omega}/\bar\vartheta^2)d\bar\vartheta
=\varrho_c(2+MX/C)ds$.  (iii) As in (ii), but the sum over
$Q\setminus\Omega$ does not change.
\end{proof}

\begin{corollary}[Relative positions as functions of ratios]\label{cor:II-Rkt}
Let $\zeta\in\ell_1\setminus\{0\}$, $\Omega\subseteq Q(\zeta)$ finite,
$\varkappa:=\varkappa(\zeta;\Omega)$, and put $\mathsf r_k:=|\zeta(k)|/(\Phi(k)
\varkappa)$ for $k\in\Omega$, $R^2:=\sum_{k\in\Omega}\mathsf r_k^2$,
$a:=|\zeta|/\bar\vartheta$, $\mathsf k:=\varkappa/\bar\vartheta$ and
$s_Q:=\bar\vartheta^{-2}\sum_{k\in Q\setminus\Omega}\mathsf v_k^2\Phi(k)^2$.  Then
\[
 a^2=\Phi_P^2+\mathsf k^2R^2+s_Q,\qquad \mathsf k=a+\Phi_P^2+s_Q,\qquad
 \varrho_k=\mathsf r_k\mathsf k/\Phi(k)\quad(k\in\Omega),
\]
and $R<1$, $a>\mathsf kR$.  In particular, for fixed $(\Phi_P^2,s_Q)$ the
relative positions of the coordinates of $\Omega$ are functions of the ratios
$(\mathsf r_k)_{k\in\Omega}$ alone.
\end{corollary}

\begin{proof}
Divide the second identity of Lemma~\ref{lem:II-T}(c) by $\bar\vartheta^2$,
using $\sum_\Omega\mathsf v_k^2\Phi(k)^2=\varkappa^2R^2$; divide the second
expression of $\varkappa$ by $\bar\vartheta$; and $\varrho_k=\mathsf v_k/\bar\vartheta$ equals
$(|\zeta(k)|/(\Phi(k)\varkappa))\cdot(\varkappa/\bar\vartheta)/\Phi(k)$.
As $P\ne\emptyset$, $\mathsf k>a$; if $R\ge1$ then $a^2\ge\mathsf k^2R^2\ge
\mathsf k^2>a^2$, which is absurd; and $a^2-\mathsf k^2R^2=\Phi_P^2+s_Q>0$.

For the last assertion put $c_0:=\Phi_P^2+s_Q>0$, so that $\mathsf k=a+c_0$
and $a$ is a root of $G(a):=a^2-(a+c_0)^2R^2-c_0$.  The condition $a>\mathsf
kR=(a+c_0)R$ is equivalent (as $R<1$) to $a>c_0R/(1-R)$.  On the half-line
$\{a>c_0R/(1-R)\}$ we have $G'(a)=2(a-(a+c_0)R^2)=2(a-\mathsf kR^2)>0$,
because $a>\mathsf kR\ge\mathsf kR^2$.  Hence $G$ is strictly increasing on
this half-line and has at most one root there.  Since the true value of $a$
lies on it, $a$ is determined by $(R,\Phi_P^2,s_Q)$, and hence so are
$\mathsf k=a+c_0$ and $\varrho_k=\mathsf r_k\mathsf k/\Phi(k)$ for $k\in
\Omega$ (with $\Phi(k)$ fixed).  Compare Proposition~\ref{prop:II-Rkt}(i)
for the smooth dependence.
\end{proof}

For the block vector $\hat\zeta_m=R_m^{**}\zh$ one has $|\zeta(k)|/\Phi(k)=
m|u_{k,m}(\zh)|$, so $\mathsf r_k=m|u_{k,m}(\zh)|/\varkappa$.  A push of a
peak (Lemma~\ref{lem:II-kappader}(b)(i)) with the values on $\Omega$ fixed
rescales all ratios of the block by one factor; if the ratios are restored
afterwards, all relative positions in $\Omega$ return to their old values up
to the change of $\Phi_P^2$ and $s_Q$.  Exactness of shifted data
(through Lemma~\ref{lem:II-kappa}) and the statuses of the switching
carriers therefore cannot be adjusted independently by peak pushes alone.

\subsection{The cost of moving the normer; raise transfer}\label{subsec:II-cost}

For a first row $f'$ write $\delta:=\zh'-\zh\in\ell_\infty$ and
\[
 \begin{gathered}
 \Delta_m(\delta):=\sum_k\lambda_{k,m}|u_{k,m}(\delta)|=\|R_m^{**}\delta\|_1,\\
 c(\delta):=\sum_{m\in I}\Big[\Delta_m\log\frac e{\Delta_m}+\sum_k
 \min\big(\lambda_{k,m},|u_{k,m}(\delta)|\big)\Big],
 \end{gathered}
\]
with $0\log(e/0):=0$.  We use the elementary counting bound
\begin{equation}\label{eq:II-count}
 \sum_k\min(2\lambda_{k,m},x)\le x\big(\log_2(2m/x)+3\big)\qquad(0<x\le2m),
\end{equation}
which holds because $2\lambda_{k,m}\le2m2^{-m-k}$: at most
$\log_2(2m/x)+1$ indices have $2m2^{-m-k}>x$, and the remaining terms form a
geometric series bounded by $2x$.

\begin{lemma}[Cost of moving the normer]\label{lem:II-cost}
Let $f\in S_{p^*}$.  There are $c_f>0$ and $C_f<\infty$, depending only on
$f$ \textup{(}and on $I$, $T$\textup{)}, such that for every $f'\in S_{p^*}$
with $\max_m\Delta_m(\delta)\le c_f$, $\delta:=\zh'-\zh$:
\begin{enumerate}
\item[(i)] $|C'_m-C_m|+|\sigma'_m-\sigma_m|+|q'_0-q_0|\le C_f\max_{m'}
\Delta_{m'}(\delta)$ for every $m$;
\item[(ii)] $\sum_m\|R_m^*(w'_m-w_m)\|_1\le C_fc(\delta)$;
\item[(iii)] $\|D_m(w'_m-w_m)\|_2^2\le C_fc(\delta)^2$ for every $m$;
\item[(iv)] $p^*(L^*(w'-w))\le(1+\|U\|)C_fc(\delta)$ and
$p^*(f'-f)\le(1+\|U\|)\big(\|a'-a\|_1+C_fc(\delta)\big)$.
\end{enumerate}
In particular, for a companion $f^\#$ of $f$ \textup{(}$a^\#=a$\textup{)},
$p^*(f^\#-f)\le(1+\|U\|)C_fc(\delta)$.
\end{lemma}

\begin{proof}
Fix $m$, drop it, and write $A:=|\hat\zeta|$, $A':=|\hat\zeta'|$ with
$\hat\zeta=R^{**}\zh$, $\hat\zeta'=R^{**}\zh'$; $\Delta:=\Delta_m(\delta)=
\|\hat\zeta'-\hat\zeta\|_1$.  Recall $|u_k(\zh)|\le q^*(u_k)q^{**}(\zh)=1$, and
$|A'-A|\le\Delta$ since $|\cdot|\le\|\cdot\|_1$ [I,~Remark~1.1(c)].

\emph{Scalars.}  $1/q_0=1+\sum_mA_m$, so $|1/q'_0-1/q_0|\le\sum_m\Delta_m$ and
$|q'_0-q_0|\le|I|\max_m\Delta_m$; $\sigma=q_0A$ and $A'\le\|\hat\zeta'\|_1\le
\sum_k\lambda_k\le1$ give the bound for $\sigma$.  By Lemma~\ref{lem:II-T}(b)
(applied to $\hat\zeta$; $M$, $C$ are scale invariant),
$C=A/(\bar\vartheta(\hat\zeta)+A)$, and by Lemma~\ref{lem:II-T2}(a)
$|\bar\vartheta(\hat\zeta')-\bar\vartheta(\hat\zeta)|\le C_1\Delta$ as soon as
$C_1\Delta<\min(h_0,\bar\vartheta(\hat\zeta))$ (constants of that lemma for
$\hat\zeta$).  Hence $|C'-C|\le c_3\Delta$ with $c_3$ depending only on $f$,
once $\Delta$ is below a constant depending only on $f$.  This is (i).

\emph{Clamp formula.}  By Lemma~\ref{lem:II-T}(b), $w(k)=\sgn(\hat\zeta(k))
\min(M,C|\hat\zeta(k)|/(A\Phi(k)^2))$ (use $M/\bar\vartheta=C/A$), and
$|\hat\zeta(k)|=m\Phi(k)|u_k(\zh)|$; so with $v_k:=m|u_k(\zh)|/A$,
\[
 \Phi(k)w(k)=\sgn(u_k(\zh))\min\big(\Phi(k)M,Cv_k\big),
\]
and the same at $f'$ with $v'_k:=m|u_k(\zh')|/A'$.  For $\Delta\le A/2$,
$|v'_k-v_k|\le c_1(|u_k(\delta)|+\Delta)$ with $c_1:=2m(1+1/A)/A$.  If
$u_k(\zh)$ and $u_k(\zh')$ have the same sign or one of them vanishes, then
$\Phi(k)|w'(k)-w(k)|\le|C'-C|(\Phi(k)+v_k)+C'|v'_k-v_k|$ (note $M'-M=C-C'$);
if they have opposite signs, then $|u_k(\zh)|,|u_k(\zh')|\le|u_k(\delta)|$
and $\Phi(k)(|w'(k)|+|w(k)|)\le C'v'_k+Cv_k\le3m|u_k(\delta)|/A$.  Since
$v_k\le m/A$, $|w|,|w'|\le1$ and $\lambda_k=m\Phi(k)$, in both cases, with a
constant $c_5$ depending only on $f$,
\[
 \begin{aligned}
 \lambda_k|w'(k)-w(k)|&\le\min\big(2\lambda_k,\ m\Phi(k)|C'-C|+c_5\Delta+
 c_5|u_k(\delta)|\big)\\
 &\le m\Phi(k)|C'-C|+\min(2\lambda_k,c_5\Delta)+
 \min(2\lambda_k,c_5|u_k(\delta)|).
 \end{aligned}
\]

\emph{(ii).}  Sum over $k$, using $\sum_k\Phi(k)\le1$, (i),
\eqref{eq:II-count} with $x=c_5\Delta$ (for $c_5\Delta\le2m$), and
$\min(2\lambda,c_5y)\le\max(2,c_5)\min(\lambda,y)$; and
$\|R^*x\|_1\le\sum_k\lambda_k|x(k)|$ because $\|u_k\|_1\le q^*(u_k)=1$.

\emph{(iii).}  $\|D(w'-w)\|_2^2\le\max_k\big(\Phi(k)|w'(k)-w(k)|\big)\sum_k
\Phi(k)|w'(k)-w(k)|$.  The second factor is $\frac1m\sum_k\lambda_k|w'(k)-
w(k)|\le C c(\delta)$ by (ii).  By the pointwise bounds, the first factor is
at most $c_6\Delta+\min(2\Phi(k),c_6|u_k(\delta)|)\le c_7c(\delta)$, since
$\Delta\le\Delta\log(e/\Delta)$ for $\Delta\le1$ and $\min(2\Phi(k),c_6y)\le
\max(2/m,c_6)\min(\lambda_k,y)$.

\emph{(iv).}  $p^*\le q^*\le(1+\|U\|)\|\cdot\|_1$ on $X^*$ (by
[I,~Lemma~1.6], $B_{q^*}\subseteq B_{p^*}$), $L^*(w'-w)=\sum_mR_m^*(w'_m-w_m)$,
and $f'-f=(a'-a)+L^*(w'-w)$.  Finally choose $c_f$ so small that all the
smallness conditions used above hold for $\max_m\Delta_m\le c_f$.
\end{proof}

\begin{remark}\label{rem:II-cost}
(a) Lemma~\ref{lem:II-cost} is an upper bound.  Its unweighted term,
the sum of the numbers $\min(\lambda_{k,m},|u_{k,m}(\delta)|)$, cannot in
general be replaced by a $\lambda$-weighted one: at a coordinate $k$ that is a strict non-peak of
$f$ and $f'$, the clamp formula gives
\[
 \lambda_{k,m}(w'_m(k)-w_m(k))=m^2\big(C'_mu_{k,m}(\zh')/A'_m-C_mu_{k,m}(\zh)/
 A_m\big)
\]
\textup{(}$A_m:=|R_m^{**}\zh|_m$\textup{)}, which is linear in the
unweighted quantity $u_{k,m}(\delta)$ with coefficient $m^2C_m/A_m$ when the
scalars do not move; the changes of the scalars are $O(\Delta_m)$, i.e.\ of
weighted size.  That the unweighted term is attained is supported by
numerical evidence in our sources (not used).
(b) The lemma uses $f'$ only through $\zh'$: $w'$ and $q'_0$ depend on $\zh'$
alone.  It applies to companions (where $\delta=z^\#-z$) and to raise
companions (where $\delta=U(e^\#-e)$, Lemma~\ref{lem:II-raise}).
\end{remark}

\begin{lemma}[Raise companions]\label{lem:II-raise}
Let $f\in S_{p^*}$, let $\Delta a$ be a raise, $\lambda:=q^*(a+\Delta a)$, and
let $f^\#=f[a^\#,z]$ be the raise companion.  Then:
\begin{enumerate}
\item[(a)] $1+\|\Delta a\|_1-\|U^*\Delta a\|_H\le\lambda\le1+\|\Delta a\|_1+
\|U^*\Delta a\|_H$;
\item[(b)] $\|e^\#-e\|_H\le2\|U^*\Delta a\|_H/\nu$;
\item[(c)] $\zh^\#-\zh=U(e^\#-e)$; hence $|u(\zh^\#)-u(\zh)|\le\|U^*u\|_H\,
\|e^\#-e\|_H$ for $u\in\ell_1$, and $\Delta_m(\zh^\#-\zh)\le m2^{-m}\|e^\#-e\|_H$;
\item[(d)] if $\lambda\ge1$, then $\|a^\#-a\|_1\le2\|\Delta a\|_1+\|U^*\Delta
a\|_H$; and $f^\#\to f$ in norm as $\|\Delta a\|_1\to0$;
\item[(e)] put $\varepsilon_e:=\|e^\#-e\|_H$ and assume $\varepsilon_e\le1$
\textup{(}true when $\|U^*\Delta a\|_H\le\nu/2$, by \textup{(b))}; then one has
$c(\zh^\#-\zh)\le C_I\varepsilon_e\log(e/\varepsilon_e)$, $C_I$ depending only
on $I$; so, if $\max_m\Delta_m(\zh^\#-\zh)\le c_f$, Lemma~\textup{\ref{lem:II-cost}} gives $p^*(L^*(w^\#-w))\le C'_f\varepsilon_e\log(e/
\varepsilon_e)$.
\end{enumerate}
\end{lemma}

\begin{proof}
(a) On $F$ the signs of $a$ and $\Delta a$ agree, so $\|a+\Delta a\|_1=\|a\|_1+
\|\Delta a\|_1$; and $\|U^*a\|-\|U^*\Delta a\|\le\|U^*(a+\Delta a)\|\le
\|U^*a\|+\|U^*\Delta a\|$; use $q^*(a)=\|a\|_1+\|U^*a\|=1$.  (b) $e^\#$ is
the normalization of $U^*(a+\Delta a)$, and $\|x/\|x\|-y/\|y\|\|\le2\|x-y\|/
\|y\|$ with $y=U^*a$, $\|y\|=\nu$.  (c) $z$ is unchanged, so $\zh^\#-\zh=
U(e^\#-e)$ and $u(Uh)=\ip{U^*u}h$; $\|U^*u_{k,m}\|\le q^*(u_{k,m})=1$ and
$\sum_k\lambda_{k,m}\le m2^{-m}$.  (d) $a^\#-a=\Delta a/\lambda+(1/\lambda-1)
a$, $\|a\|_1\le1$ and $0\le1-1/\lambda\le\lambda-1\le\|\Delta a\|_1+\|U^*\Delta
a\|$ by (a).  If $\|\Delta a\|_1\to0$ then $a^\#\to a$ in $\ell_1$ and $z$ is
fixed, so $f^\#\to f$ [I,~Remark~8.22(c)].  (e) By (c), $|u_{k,m}(\delta)|\le
\varepsilon_e$ and $\Delta_m\le m2^{-m}\varepsilon_e$; since $x\log(e/x)$ is
increasing on $(0,1]$ and by \eqref{eq:II-count},
$c(\delta)\le\sum_m[m2^{-m}\varepsilon_e\log(e2^m/(m\varepsilon_e))+
\varepsilon_e(\log_2(2m/\varepsilon_e)+3)]\le C_I\varepsilon_e\log(e/
\varepsilon_e)$.
\end{proof}

\begin{lemma}[Raise transfer]\label{lem:II-RT}
In the situation of Lemma~\textup{\ref{lem:II-raise}} assume $\lambda\ge1$.
Then for every $h\in X^*$ and $r\in\R$
\[
 p^*(f^\#+rh)\le1+\frac{p^*(f+\lambda rh)-1+2\|U^*\Delta a\|_H}\lambda+
 p^*\big(L^*(w^\#-w)\big).
\]
\end{lemma}

\begin{proof}
For $r=0$ the right side is at least $1=p^*(f^\#)$.  Let $r\ne0$.  By
[I,~Lemma~1.6], $f+\lambda rh=A+L^*W$ with $\max(q^*(A),\|W\|_{V^*})=\pi:=
p^*(f+\lambda rh)$.  Since $f=a+L^*w$, $\lambda rh=(A-a)+L^*(W-w)$, and
\[
 f^\#+rh=\frac{A+\Delta a}\lambda+L^*\Big(\big(1-\tfrac1\lambda\big)w+
 \tfrac1\lambda W\Big)+L^*(w^\#-w).
\]
\emph{Base.}  $q^*(A+\Delta a)\le\pi+\|\Delta a\|_1+\|U^*\Delta a\|\le\pi+
(\lambda-1)+2\|U^*\Delta a\|$ by Lemma~\ref{lem:II-raise}(a), so
$q^*((A+\Delta a)/\lambda)\le1+(\pi-1+2\|U^*\Delta a\|)/\lambda$.
\emph{Blocks.}  As $\lambda\ge1$, $N_m((1-1/\lambda)w_m+W_m/\lambda)\le
(1-1/\lambda)+\pi/\lambda=1+(\pi-1)/\lambda$.  By [I,~Lemma~1.6] the
$p^*$-norm of the first two terms is at most the maximum of these bounds;
add $p^*(L^*(w^\#-w))$.
\end{proof}

No smallness of $\|\Delta a\|_1$ is needed: the $\ell_1$-mass of a same-sign
raise on $F$ cancels against the normalization $\lambda-1$, and only the
Hilbert footprint $\|U^*\Delta a\|$ and the normer change remain.

\begin{remark}[Raise with corrections]\label{rem:II-RTcorr}
Suppose that, in addition to the raise $\Delta a$, corrections $\Delta a''$
of arbitrary sign are made on a finite set $F_0\subseteq F\setminus
\supp\Delta a$ with $|\Delta a''_j|\le|a_j|/2$ for $j\in F_0$.  Put
$D:=\Delta a+\Delta a''$, assume $\lambda:=q^*(a+D)\ge1$, and let
$f^\#:=f[(a+D)/\lambda,z]$ with normer $w^\#$.  The condition
$|\Delta a''_j|\le|a_j|/2$ keeps $\supp(a+D)=F$ and the signs of $a$ on
$F$, so $f^\#$ is well defined (Subsection~\ref{subsec:II-comp}).  Then the
proof of Lemma~\ref{lem:II-RT} gives, for every $h\in X^*$ and $r\in\R$,
\[
 p^*(f^\#+rh)\le1+\frac{p^*(f+\lambda rh)-1+2\|U^*D\|_H+2\|\Delta a''\|_1}
 \lambda+p^*\big(L^*(w^\#-w)\big).
\]
Indeed, $\|a+D\|_1\ge\|a\|_1+\|\Delta a\|_1-\|\Delta a''\|_1$ and
$\|U^*(a+D)\|_H\ge\|U^*a\|_H-\|U^*D\|_H$ give $\|\Delta a\|_1\le\lambda-1+
\|\Delta a''\|_1+\|U^*D\|_H$, hence $q^*(A+D)\le\pi+\|D\|_1+\|U^*D\|_H\le
\pi+(\lambda-1)+2\|\Delta a''\|_1+2\|U^*D\|_H$, and the block part of the
proof is unchanged.  The error is that of the whole change $D$; compared
with Lemma~\ref{lem:II-RT} the extra error is at most $2q^*(\Delta a'')/
\lambda$, since $\|U^*D\|_H\le\|U^*\Delta a\|_H+\|U^*\Delta a''\|_H$.
\end{remark}

\begin{corollary}[Mates under a raise]\label{cor:II-RT}
In the situation of Lemma~\textup{\ref{lem:II-RT}}, let $g\in\Cset(f)$ and
$\rho\in(0,1]$.  Then for all $r\in\R$
\[
 p^*(f^\#+r\rho g)\le1+\frac{s(\lambda\rho r)-1}\lambda+\varepsilon',\qquad
 \varepsilon':=2\|U^*\Delta a\|_H+p^*\big(L^*(w^\#-w)\big),
\]
and $(s(\lambda x)-1)/\lambda\le s(x)-1+(\lambda-1)\min(\lambda^2x^2/2,1)$ for
$x\in\R$.
\end{corollary}

\begin{proof}
Lemma~\ref{lem:II-RT} with $h=\rho g$ and $p^*(f+\lambda r\rho g)\le
s(\lambda\rho r)$ (as $g\in\Cset(f)$), and $2\|U^*\Delta a\|/\lambda\le
2\|U^*\Delta a\|$.  For the inequality put $\varphi(\mu):=(s(\mu x)-1)/\mu$;
with $u=\mu x$ and $s'(u)=u/s(u)$ one computes $\varphi'(\mu)=(s(u)-1)/(\mu^2
s(u))$, and $0\le(s(u)-1)/s(u)\le\min(u^2/2,1)$.  For $\mu\in[1,\lambda]$
this gives $\varphi'(\mu)\le\min(\lambda^2x^2/2,1)$; integrate over
$[1,\lambda]$ and use $\varphi(1)=s(x)-1$.
\end{proof}

\subsection{Uniform one-sided transfer along convergent first rows}\label{subsec:II-U}

For a first row $f_j$ with finite base support $F_j:=\supp a_j$ we use the
notions of [I,~Definition~7.7] at $f_j$ (side-$\pm$ admissible pairs,
representation, $d^{(j)}_m$, $\Gamma^{(j)}_w$), with all data computed at
$f_j$; a superscript or subscript $j$ refers to $f_j$.  The following
coordinate kinds are used for the block parts $\omega_m$ of a pair at $f_j$
and a scale $t>0$, given constants $A_2\ge1$ and $\gamma_B\in(0,1]$:
\begin{enumerate}
\item[{[1]}] $k\in Q^{(j)}_m$ and $|\omega_m(k)|\le3\gap^{(j)}_m(k)/t$;
\item[{[2]}] $k\in Q^{(j)}_m$, $\gap^{(j)}_m(k)\ge\gamma_B$ and
$|\omega_m(k)|\le A_2/t$;
\item[{[3]}] \emph{inward coordinates}: $k$ is a degenerate peak or a strict
non-peak of $w_{j,m}$, $w_{j,m}(k)\ne0$, $|\omega_m(k)|\le A_2/t$, and, with
$\varsigma_k:=\sgn w_{j,m}(k)$, $\varsigma_k\omega_m(k)\le0$ for a
side-$+$ pair, resp.\ $\varsigma_k\omega_m(k)\ge0$ for a side-$-$ pair.
\end{enumerate}
For pairs whose block parts live on coordinates of these kinds we put
$d^{(j)}_m(\omega_m):=\ip{D_mw_{j,m}}{D_m\omega_m}/C_{j,m}$, and the pair
is said to represent the functional
\[
 b+\sum_mR_m^*\big(\omega_m-d^{(j)}_m(\omega_m)w_{j,m}\big),
\]
as in [I,~Definition~7.7]; since $\alpha_{j,m}$ vanishes at degenerate peaks and
off $P_{j,m}$, [I,~Lemma~3.2] applies: the block first-order terms vanish.
Kinds [1] and [2] are the coordinates of [I,~Lemma~8.42] (with the constant
$2$ replaced by $3$ and without the lower bound $\gap\ge t^2$, which is used
there only as a convention of the window certificate); kind [3] allows
arbitrarily small gaps, but only in the inward direction.

\begin{lemma}[Uniform one-sided transfer expansion]\label{lem:II-U}
Let $f\in S_{p^*}$ have $F$ finite, and let $f_j\to f$ in $S_{p^*}$ with
$F_j=\supp a_j$ finite and $F\subseteq F_j$.  For every $\varepsilon_{\rm tr}>0$
and $A_0\ge1$ there are $t_1>0$, $j_0\in\N$ and $c_0\in(0,\frac18]$,
depending only on $f$, the sequence $(f_j)$, $\varepsilon_{\rm tr}$ and
$A_0$, with the following property.  Let $A_2\ge1$, $\gamma_B\in(0,1]$, and
put $c_\flat:=c_0\gamma_B/A_2$.  Let $j\ge j_0$, $t\in(0,t_1]$, and let
$(b,\omega)$, $b\in\ell_1$, $\omega_m\in c_{00}$, satisfy:
\begin{enumerate}
\item[(OS1)] $b(\xi_j)=0$, $t\|b\|_1\le A_0$ and $\Gamma^{(j)}_w(b,\omega)\le2$;
\item[(OS2)] $b(i)=0$ for $i\notin F_j\cup K_j$, $z_j(i)b(i)\ge0$ for $i\in
K_j$, and for every $i\in F_j\setminus F$ either $z_j(i)b(i)\ge0$ or
$t|b(i)|\le|a_j(i)|$;
\item[(OS3)] every $k\in\supp\omega_m$ is of kind \textup{[1]}, \textup{[2]}
or \textup{[3]} \textup{(}side $+$\textup{)} at $f_j$, with the given $A_2,
\gamma_B$.
\end{enumerate}
Then $G:=b+\sum_mR_m^*(\omega_m-d^{(j)}_m(\omega_m)w_{j,m})$ satisfies
\[
 p^*(f_j+rG)\le1+\frac{r^2}2\big(\Gamma^{(j)}_w(b,\omega)+\varepsilon_{\rm tr}
 \big)\qquad(0<r\le c_\flat t).
\]
The same holds for side-$-$ pairs \textup{(}in \textup{(OS2)} and \textup{[3]}
the signs reversed\textup{)} and $-c_\flat t\le r<0$.
\end{lemma}

The sequence may be constant: for $f_j=f$ and $F_j=F$ the lemma is
[I,~Lemma~8.42] with the enlarged coordinate kinds.  The essential point is
that $t_1$ and $j_0$ do not depend on $A_2$ and $\gamma_B$, and that $c_\flat$
depends on them only through $\gamma_B/A_2$.  The coordinates of
$F_j\setminus F$ are of two types: \emph{banks}, at which the data have the
sign of $a_j$, and \emph{pulls}, at which the data are small compared with
the mass $|a_j(i)|$; nothing is required of $f$ at these coordinates.

\begin{proof}
We follow the proofs of [I,~Lemmas~8.16 and~8.42] and record where $f_j$,
$A_2$ and $\gamma_B$ enter.  We treat side $+$; side $-$ is the same with
$-r$, $-b$, $-\omega$.

\emph{Step 0: limits and transfer data.}  We may assume $\varepsilon_{\rm tr}\le1$.  Let $a_{\min}:=\min_{i\in F}|a_i|$,
$C_{\min}:=\min_mC_m$, $M_{\min}:=\min_mM_m$, and fix $\gamma_m\in(M_m/2,M_m)$
with $|w_m(k)|\ne\gamma_m$ for all $k$.  By [I,~Proposition~1.11] the forced
data of $f_j$ converge to those of $f$; let $j'_0$ be such that for $j\ge
j'_0$: $q^{(j)}_0\ge q_0/2$, $\sigma_{j,m}\ge\sigma_m/2$, $C_{j,m}\ge C_m/2$,
$M_{j,m}-\gamma_m\ge(M_m-\gamma_m)/2$, $\nu_j\ge\nu/2$ and $\min_{i\in F}
|a_j(i)|\ge a_{\min}/2$.  Put $K_3:=16/\min(q_0,\min_m\sigma_m)$ and
$\eta_1:=\min(\frac14,\varepsilon_{\rm tr}/(20|I|K_3))$, and choose transfer
data $(\gamma_m,k^\varsigma_{*,m},\Lambda^\varsigma_m)$ of $f$ in every block
by [I,~Lemma~7.4] with this $\eta_1$.  By [I,~Lemma~7.5] there is $j_0\ge
j'_0$ such that for $j\ge j_0$ the transfer vectors $y^\varsigma_{j,m}$
built at $f_j$ (with $L_{j,m}:=\{|w_{j,m}|\ge\gamma_m\}$, which contains
$P_{j,m}$, and $k^\varsigma_{*,m}$ a peak of $w_{j,m}$ of value $M_{j,m}$)
satisfy $e^\varsigma_{y,j,m}\ge\frac12$, $\iota^\varsigma_{j,m}\le2\eta_1$,
$\|D_my^\varsigma_{j,m}\|_2\le K_y$ and $|Y^\varsigma_{j,m}|/q^{(j)}_0\le K_Y$,
where $K_y,K_Y$ depend only on $f$ and the $\gamma_m$
([I,~Lemma~7.4(b),(d)]).  Let $j\ge j_0$, $t\le t_1$, $0<r\le c_\flat t$,
and write $f_j+rG=A+\sum_mR_m^*W_m$ with $A:=a_j+rb$ and $W_m:=
(1-rd_m)w_{j,m}+r\omega_m$, $d_m:=d^{(j)}_m(\omega_m)$.  The first-order terms
at $f_j$ vanish: $\lin_b(A)=rb(\zh_j)=0$ by (OS1), and $\lin_m(W_m)=0$ by
[I,~Lemma~3.2] at $f_j$.  Put $A_3:=3+A_2$; then $t\|\omega_m\|_\infty\le
A_3$ (kind [1]: $\gap\le1$) and $|d_m|\le\|D_m\omega_m\|_2\le\|\Phi_m\|_2
\|\omega_m\|_\infty\le A_3/(2t)$.

\emph{Step 1: base.}  On $F$, $|rb(i)|\le c_\flat t\|b\|_1\le c_\flat A_0\le
a_{\min}/2\le|a_j(i)|$ if $c_\flat\le a_{\min}/(2A_0)$: no flip.  On
$F_j\setminus F$: if $z_j(i)b(i)\ge0$ then $a_j(i)$ and $rb(i)$ have the
same sign $z_j(i)=\sgn a_j(i)$; otherwise $|rb(i)|\le c_\flat t|b(i)|\le
c_\flat|a_j(i)|<|a_j(i)|$; in both cases no flip.  Off $F_j$ the summands
$|rb(i)|-z_j(i)rb(i)$ vanish by (OS2).  By [I,~Lemma~7.1(b)] at $f_j$ and the
estimate of $\Psi$ in [I,~Lemma~3.3],
\[
 G_b(A)=\nu_j\Psi\big(rU^*b/\nu_j\big)\le\hat G_b:=\frac{r^2}2h^{(j)}(b)
 (1+K_Ac_\flat),\qquad K_A:=4\|U\|A_0/\nu,
\]
provided $c_\flat\|U\|A_0\le\nu/4$ (so that $\|rU^*b\|/\nu_j\le\frac12$).

\emph{Step 2: blocks.}  If $c_\flat\le1/A_3$ then $|rd_m|\le\frac12$.  We
claim $\|W_m\|_\infty=(1-rd_m)M_{j,m}$.  At a non-degenerate peak $\omega_m=0$
and $|W_m(k)|=(1-rd_m)M_{j,m}$; off $\supp\omega_m$, $|W_m(k)|=(1-rd_m)
|w_{j,m}(k)|$.  At kinds [1], [2], if $c_\flat\le1/6$ and $c_\flat\le
\gamma_B/(2A_2)$, then $|r\omega_m(k)|\le\gap^{(j)}_m(k)/2$ and
$|W_m(k)|\le(1-rd_m)(M_{j,m}-\gap^{(j)}_m(k))+\gap^{(j)}_m(k)/2\le(1-rd_m)
M_{j,m}$.  At kind [3], $\varsigma_kW_m(k)=(1-rd_m)|w_{j,m}(k)|+r\varsigma_k
\omega_m(k)\le(1-rd_m)M_{j,m}$, and $\varsigma_kW_m(k)\ge-c_\flat A_2\ge
-(1-rd_m)M_{j,m}$ if $c_\flat A_2\le M_{\min}/8$.  This proves the claim.
The proof of [I,~Lemma~3.4(d)] uses its radius condition only to obtain
this equality, and the condition $|rd_m|\le C_{j,m}/(2M_{j,m})$, which holds
if $c_\flat\le C_{\min}/(2A_3)$.  Hence
\[
 G^{(j)}_m(W_m)=N_m(W_m)-1\le\hat G_m:=\frac{r^2}2H^{(j)}_m(\omega_m)\big(1+
 2A_3c_\flat/C_{\min}\big).
\]
Consequently $\hat\Gamma:=q^{(j)}_0\hat G_b+\sum_m\sigma_{j,m}\hat G_m\le\frac{r^2}2
\Gamma^{(j)}_w(b,\omega)(1+(K_A+2A_3/C_{\min})c_\flat)$.

\emph{Step 3: rebalancing.}  Since $h^{(j)}(b)\le2/q^{(j)}_0$ and
$H^{(j)}_m(\omega_m)\le2/\sigma_{j,m}$ by (OS1), the numbers $\varepsilon_m:=
(\hat G_m-\hat\Gamma)/e_{y,j,m}$ (with $y^\downarrow$ or $y^\uparrow$
according to the sign) satisfy $|\varepsilon_m|\le K_3r^2$.  Now
$\|W_m-w_{j,m}\|_\infty\le r(\|\omega_m\|_\infty+|d_m|)\le2A_3c_\flat$ and
$\|D_m(W_m-w_{j,m})\|_2\le A_3c_\flat$.  If $2A_3c_\flat\le\min_m(M_m-\gamma_m)
/16$, $A_3c_\flat\le C_{\min}/4$, $K_3t_1^2(2+\Lambda^\varsigma_m)\le\gamma_m/2$,
$K_3t_1^2\le\min_m(M_m-\gamma_m)/8$ and $|I|K_3K_Yt_1^2\le1$, then
[I,~Lemma~7.2] at $f_j$ gives [I,~(9)] and $|\lambda|\le1$, and
[I,~Proposition~7.3] at $f_j$ gives $p^*(f_j+rG)\le1+\hat\Gamma+E$ with
\[
 E\le2|I|K_3\eta_1r^2+2|I|K_3K_Y(1+\|U\|)A_0c_\flat r^2+\frac{8K_3K_yA_3
 c_\flat r^2+2K_3^2K_y^2r^4}{C_{\min}},
\]
because $|q^*(A)-1|+q^*(A-a_j)\le2q^*(rb)\le2(1+\|U\|)c_\flat A_0$.

\emph{Step 4: constants.}  Require in addition
$2(K_A+2A_3/C_{\min})c_\flat\le\varepsilon_{\rm tr}/5$,
$4|I|K_3K_Y(1+\|U\|)A_0c_\flat\le\varepsilon_{\rm tr}/5$,
$16K_3K_yA_3c_\flat/C_{\min}\le\varepsilon_{\rm tr}/5$ and
$4K_3^2K_y^2t_1^2/C_{\min}\le\varepsilon_{\rm tr}/5$; with $4|I|K_3\eta_1\le
\varepsilon_{\rm tr}/5$ and $\Gamma^{(j)}_w\le2$ this gives $\hat\Gamma+E\le
\frac{r^2}2(\Gamma^{(j)}_w(b,\omega)+\varepsilon_{\rm tr})$.  Every condition
on $c_\flat$ above has the form $c_\flat\le\kappa_1$, $c_\flat\le\kappa_1/A_3$,
$c_\flat\le\kappa_1/A_2$ or $c_\flat\le\gamma_B/(2A_2)$ with $\kappa_1>0$
depending only on $f$, $\varepsilon_{\rm tr}$, $A_0$; since $A_3\le4A_2$ and
$\gamma_B\le1\le A_2$, all of them hold for $c_\flat=c_0\gamma_B/A_2$ with
$c_0>0$ small enough.  The conditions on $t_1$ and $j_0$ do not involve
$A_2$ or $\gamma_B$.
\end{proof}

\subsection{Windowed recovery through nearby first rows}\label{subsec:II-E}

The mechanism of [I,~Theorem~8.17] (windowed averaging) and of
[I,~Theorem~8.44] (averaged exact window data, then [I,~Corollary~7.21])
does not require the window data to live at $f$ itself: they may live at a
nearby first row $f_j$, provided $f_j$ is closer to $f$ than the square of
the \emph{bottom} of the window.  We first prove an abstract averaging
theorem that also allows a relative factor $\lambda_j$ in the comparison
of $f_j$ with $f$ (this is what raise companions produce,
Corollary~\ref{cor:II-RT}), and then the version with exact two-piece data.
For a window $(T,n)$, $T\in(0,1]$, $n\in\N$, we write $\mathcal S(T,n):=
\{T2^{1-i}:1\le i\le n\}$; its least element is $2T2^{-n}$.

\begin{theorem}[Windowed recovery with a transfer error]\label{thm:II-ERT}
Let $f\in S_{p^*}$ \textup{(}$F$ arbitrary\textup{)}, $g\in\Cset(f)$,
$\rho\in(0,1)$ and $\eta_0\in(0,2]$ with $\rho^2(1+\eta_0)\le(1+\rho^2)/2$.
Suppose that there are first rows $f_j\in S_{p^*}$, numbers $\lambda_j\ge1$,
$\varepsilon'_j\ge0$, $c_j\in(0,1]$, and windows $(T_j,n_j)$ with $T_j\in(0,1]$,
$n_j\in\N$, together with numbers $K_j\ge0$, such that:
\begin{enumerate}
\item[(a)] $p^*(f_j+r\rho g)\le1+(s(\lambda_j\rho r)-1)/\lambda_j+
\varepsilon'_j$ for all $r\in\R$;
\item[(b)] for every $t\in\mathcal S_j:=\mathcal S(T_j,n_j)$ there is
$g_{j,t}\in X^*$ with $p^*(g-g_{j,t})\le K_jt$ and $p^*(f_j+r\rho g_{j,t})\le
1+\frac{\rho^2r^2}2(1+\eta_0)$ for $0<\rho|r|\le c_jt$;
\item[(c)] $K_jT_j\to0$, $\lambda_j\to1$, $\varepsilon'_j\to0$,
$p^*(f_j-f)\to0$, and $n_j\ge48\rho^2K_j/(c_j(1-\rho^2))$ for all large $j$;
\item[(d)] $\varepsilon'_j\le\theta_*c_j^2(T_j2^{-n_j})^2$ for all large $j$,
where $\theta_*:=(1-\rho^2)/(12\rho^2)$;
\item[(e)] for every $j$: if $\rho\bar g_j\in\Cset(f_j)$, where $\bar
g_j:=\frac1{n_j}\sum_{t\in\mathcal S_j}g_{j,t}$, then $(f_j,\rho\bar g_j)\in
\cl\NA((c_0,p),\ell_2^2)$.
\end{enumerate}
Then $(f,\rho g)\in\cl\NA((c_0,p),\ell_2^2)$.
\end{theorem}

\begin{proof}
Put $r_0:=\sqrt{1-\rho^2}\in(0,1)$ and let $j$ be so large that
(c), (d) are in force and
\[
 (\lambda_j-1)(1+\lambda_j^2)\le\frac{(1-\rho^2)r_0^2}{24},\qquad
 \varepsilon'_j\le\frac{(1-\rho^2)r_0^2}{12},\qquad
 \frac{2\rho K_jT_j}{n_j}\le\frac{(1-\rho^2)r_0}{12}.
\]
Drop the index $j$.  For $t\in\mathcal S$ and $r\in\R$ we have two bounds:
\begin{enumerate}
\item[(A)] if $\rho|r|\le ct$: $p^*(f_j+r\rho g_t)\le1+\frac{\rho^2r^2}2
(1+\eta_0)\le1+\frac{1+\rho^2}4r^2$, by (b) and the choice of $\eta_0$;
\item[(B)] always: $p^*(f_j+r\rho g_t)\le p^*(f_j+r\rho g)+\rho|r|Kt\le
s(\rho r)+(\lambda-1)\min(\lambda^2\rho^2r^2/2,1)+\varepsilon'+\rho|r|Kt$,
by (a), (b) and Corollary~\ref{cor:II-RT}.
\end{enumerate}
By convexity, $p^*(f_j+r\rho\bar g)\le\frac1n\sum_{t\in\mathcal S}p^*(f_j+r
\rho g_t)$.

\emph{Case $|r|\le r_0$.}  Let $\mathcal S_r:=\{t\in\mathcal S:ct<\rho|r|\}$.
If $\mathcal S_r=\emptyset$, every term obeys (A), and $1+\frac{1+\rho^2}4r^2
\le s(r)$ (see below).  Otherwise $\rho|r|>c\min\mathcal S=2cT2^{-n}$, so
$\varepsilon'\le\theta_*c^2(T2^{-n})^2<\theta_*\rho^2r^2/4\le(1-\rho^2)r^2/48$;
moreover $\sum_{t\in\mathcal S_r}t<2\rho|r|/c$ (the scales are dyadic), so
$\frac1n\sum_{t\in\mathcal S_r}\rho|r|Kt\le2\rho^2Kr^2/(cn)\le(1-\rho^2)r^2/
24$; and $(\lambda-1)\lambda^2\rho^2r^2/2\le(1-\rho^2)r^2/48$.  Using (A) off
$\mathcal S_r$ and (B) on $\mathcal S_r$,
\[
 p^*(f_j+r\rho\bar g)\le\max\Big\{1+\frac{1+\rho^2}4r^2,\ s(\rho r)\Big\}+
 \frac{(1-\rho^2)r^2}{12}.
\]
Now $1+\frac{1+\rho^2}4r^2+\frac{1-\rho^2}{12}r^2=1+\frac{2+\rho^2}6r^2\le
1+\frac{r^2}2-\frac{r^4}8\le s(r)$, because $r^2\le1-\rho^2\le
\frac43(1-\rho^2)$ and $\sqrt{1+x}\ge1+\frac x2-\frac{x^2}8$; and
$s(\rho r)+\frac{1-\rho^2}{12}r^2\le s(r)$ by [I,~Lemma~3.6(a)] ($|r|\le1$).

\emph{Case $|r|>r_0$.}  Using (B) for every $t$ and $\frac1n\sum_t t\le2T/n$,
\[
 \begin{aligned}
 p^*(f_j+r\rho\bar g)&\le s(\rho r)+(\lambda-1)+\varepsilon'+\frac{2\rho KT}n|r|\\
 &\le s(\rho r)+\frac{1-\rho^2}{12}\big(r_0^2+r_0^2+r_0|r|\big)\le s(\rho r)+
 \frac{(1-\rho^2)r_0|r|}4 ,
 \end{aligned}
\]
and $s(r)-s(\rho r)\ge(1-\rho^2)\min(r^2,|r|)/3\ge(1-\rho^2)r_0|r|/3$ by
[I,~Lemma~3.6(a)], since $r_0\le1$.

Hence $\rho\bar g_j\in\Cset(f_j)$ for all large $j$, and by (e)
$(f_j,\rho\bar g_j)\in\cl\NA$.  Finally $p^*(\rho\bar g_j-\rho g)\le\frac\rho
{n_j}\sum_tK_jt\le2\rho K_jT_j/n_j\to0$ and $p^*(f_j-f)\to0$, so
$\|(f_j,\rho\bar g_j)-(f,\rho g)\|\le p^*(f_j-f)+p^*(\rho\bar g_j-\rho g)\to0$;
as $\cl\NA$ is closed, $(f,\rho g)\in\cl\NA$.
\end{proof}

Only pieces of the window that are finer than the probing scale are
compared with $g$ at $f_j$, where $g$ is a mate only up to the error
$\varepsilon'_j$; such pieces exist exactly when $\rho|r|\gtrsim c_jT_j2^{-n_j}$,
which is the reason for (d).  With $\lambda_j=1$ and
$\varepsilon'_j:=p^*(f_j-f)$, hypothesis (a) holds by the triangle inequality
($p^*(f_j+r\rho g)\le p^*(f+r\rho g)+\varepsilon'_j$).  For raise companions,
Corollary~\ref{cor:II-RT} gives (a) with $\varepsilon'_j=2\|U^*\Delta a_j\|+
p^*(L^*(w^\#_j-w))$, which involves the Hilbert footprint of the raise and
the normer change but not $\|\Delta a_j\|_1$.

\begin{theorem}[Windowed recovery through nearby first rows]\label{thm:II-E}
Let $f\in S_{p^*}$ have $F$ finite, $g\in\Cset(f)$, $\rho\in(0,1)$, and
$\eta_0\in(0,2]$ with $\rho^2(1+\eta_0)\le(1+\rho^2)/2$.  Let $f_j\to f$ in
$S_{p^*}$ with $F_j=\supp a_j$ finite and $F\subseteq F_j$, and let
$t_1,j_0,c_0$ be given by Lemma~\textup{\ref{lem:II-U}} for
$\varepsilon_{\rm tr}:=\eta_0/2$ and a constant $A_0\ge1$.  Suppose that for
every $j$ there are $A_{2,j}\ge1$, $\gamma_j\in(0,1]$, a window $(T_j,n_j)$
and $K_j\ge1$, and for every $t\in\mathcal S(T_j,n_j)$ $d$-neutral two-piece
data $(b^\pm,\omega^\pm)=(b^\pm_{j,t},\omega^\pm_{j,t})$ at $f_j$
\textup{(}[I,~Definition~7.7] at $f_j$\textup{)} for a functional
$g_{j,t}$, such that, with $c_j:=c_0\gamma_j/A_{2,j}$:
\begin{enumerate}
\item[(E-a)] $b^\pm(\xi_j)=0$, $t\|b^\pm\|_1\le A_0$ and
$\Gamma^{(j)}_w(b^\pm,\omega^\pm)\le1+\eta_0/2$;
\item[(E-b)] every $k\in\supp\omega^\pm_m$ is of kind \textup{[1]},
\textup{[2]} \textup{(}with $A_{2,j},\gamma_j$\textup{)} or \textup{[3]}
\textup{(}for the side of the pair\textup{)} at $f_j$, and for every $i\in
F_j\setminus F$: $z_j(i)b^+(i)\ge0$ or $t|b^+(i)|\le|a_j(i)|$, and
$z_j(i)b^-(i)\le0$ or $t|b^-(i)|\le|a_j(i)|$;
\item[(E-c)] $p^*(g-g_{j,t})\le K_jt$;
\item[(E-d)] $K_jT_j\to0$ and $n_j\ge48\rho^2K_j/(c_j(1-\rho^2))$ for all
large $j$;
\item[(E-e)] $\varepsilon_j:=p^*(f_j-f)\le\theta_*c_j^2(T_j2^{-n_j})^2$ for all
large $j$, $\theta_*=(1-\rho^2)/(12\rho^2)$.
\end{enumerate}
Then $(f,\rho g)\in\cl\NA((c_0,p),\ell_2^2)$.  If the hypotheses hold for
every $\rho<1$, then $(f,g)\in\cl\NA((c_0,p),\ell_2^2)$.
\end{theorem}

\begin{proof}
We verify the hypotheses of Theorem~\ref{thm:II-ERT} with $\lambda_j:=1$,
$\varepsilon'_j:=\varepsilon_j$ and $c_j$ as given.  (a) holds by the
triangle inequality, since $p^*(f+r\rho g)\le s(\rho r)$.  For (b): as
$K_j\ge1$ and $K_jT_j\to0$, $T_j\le t_1$ for large $j$, and $j\ge j_0$.  Both
pairs represent $g_{j,t}$ at $f_j$, are side-admissible at $f_j$, and
satisfy (OS1)--(OS3) of Lemma~\ref{lem:II-U} with $(A_{2,j},\gamma_j)$ (note
$\Gamma^{(j)}_w\le1+\eta_0/2\le2$ and, as two-piece data at $f_j$, they
satisfy the sign conditions of (OS2) off $F_j$).  Lemma~\ref{lem:II-U} applied
to the $+$ pair for $0<\rho r\le c_jt$ and to the $-$ pair for $-c_jt\le\rho
r<0$ gives $p^*(f_j+\rho rg_{j,t})\le1+\frac{\rho^2r^2}2(1+\eta_0/2+\eta_0/2)$.
(c), (d) are (E-d), (E-e), together with $\varepsilon_j\to0$.  For (e): the
averaged data $\bar D_j:=\frac1{n_j}\sum_t(b^\pm_{j,t},\omega^\pm_{j,t})$ are
$d$-neutral two-piece data at $f_j$ for $\bar g_j$: the vanishing off
$F_j\cup K_j$, the sign conditions on $K_j$, the finite supports in
$Q^{(j)}_m$ (coordinates of kind [3] in two-piece data are strict
non-peaks), the two representations and $\Delta d^{(j)}_m=0$ are linear
conditions, preserved under averaging.  As $\Gamma^{(j)}_w$ is convex,
$\kappa_w(\bar D_j)\le1+\eta_0/2$, so the data $\rho\bar D_j$ of $\rho\bar g_j$
have $\kappa_w\le\rho^2(1+\eta_0/2)\le1$.  If $\rho\bar g_j\in\Cset(f_j)$,
[I,~Corollary~7.21] at $f_j$ (which has finite base support $F_j$) gives
$(f_j,\rho\bar g_j)\in\cl\NA$.  Theorem~\ref{thm:II-ERT} applies.  The last
sentence follows because $(f,g)=\lim_{\rho\to1}(f,\rho g)$ and $\cl\NA$ is
closed.
\end{proof}

\begin{remark}[Special cases and comments]\label{rem:II-E}
(a) With $f_j=f$ (so $\varepsilon_j=0$) the theorem is the mechanism of the
proof of [I,~Theorem~8.44]; there (E-d) is provided by the window
arithmetic of the design.
(b) The case $F_j=F$, $A_{2,j}=A_2$, $\gamma_j=\gamma_B$ fixed and
$\varepsilon_j\le\theta_j(T_j2^{-n_j})^2$ with $\theta_j\to0$ is the original
form; window-dependent constants $A_{2,j}$, $\gamma_j$ (hence $c_j\to0$
possibly) are needed when the gaps of the switching carriers at the
companions are only controlled by the window; kind [3] allows used
degenerate peaks of $f$ that have become strict non-peaks of $f_j$ with
arbitrarily small gap, provided the data move them only inward.
(c) The coordinates of $F_j\setminus F$ (banks and pulls) need not be
contacts of $f$: in the proof of Lemma~\ref{lem:II-U} only the sign of
$a_j(i)$ at $f_j$ and the size condition enter.
(d) The approximants $f_j$ need not belong to any recovered class, and no
lower semicontinuity of $\Cset$ along $f_j\to f$ is used:
[I,~Corollary~7.21] is applied at each $f_j$ to the mate $\rho\bar g_j$ of
$f_j$ obtained by averaging.
(e) Windows are arbitrary: any sub-window $(T_j,n_j)$ of a design window can
be used, which improves the exactification thresholds in applications.
\end{remark}

\subsection{Exact two-piece data: certificate, shift bound and rigidity}\label{subsec:II-rigid}

Let $f'\in S_{p^*}$ have finite base support $F'$.  A vector $V\in\ell_1$ is
\emph{$z'$-admissible off $F'$} if $V(j)=0$ for $j\notin F'\cup K'$ and
$z'_jV(j)\ge0$ for $j\in K'$.  For two-piece data $(b^\pm,\omega^\pm)$ at a
row $f$ (with $F$ finite) we use the \emph{data convention}
$\Delta d_m=d_m(\omega^-_m)-d_m(\omega^+_m)$ of [I,~Definition~7.7], and put
\[
 W_\Delta:=\sum_m\Delta d_m\,\psi_m,\qquad
 \Sigma(\omega):=F\cup\bigcup_{m\in I}\ \bigcup_{k\in\supp\omega^+_m\cup
 \supp\omega^-_m}\supp u_{k,m}.
\]

\begin{lemma}[Exactness certificate]\label{lem:II-cert}
Let $f'\in S_{p^*}$ have $F'$ finite, let $\Omega_m\subseteq Q'_m$ be finite,
and let $\omega^\pm_m\in c_{00}$ be supported in $\Omega_m$ \textup{(}$m\in
I$\textup{)}.  Put $\Delta\omega_m:=\omega^-_m-\omega^+_m$,
$\Delta_m:=d'_m(\Delta\omega_m)$ and
\[
 \mathcal V(\Delta\omega):=\sum_mR_m^*\big(\Delta\omega_m-\Delta_mw'_m\big)\in
 \ell_1 .
\]
\begin{enumerate}
\item[(a)] $\mathcal V(\Delta\omega)=\sum_m\sum_k\gamma_{k,m}u_{k,m}$
\textup{(}absolutely convergent in $\ell_1$\textup{)} with
$\gamma_{k,m}:=\lambda_{k,m}(\Delta\omega_m(k)-\Delta_mw'_m(k))$; in
particular every $(k,m)$ with $k\notin\Omega_m$ enters with the forced
coefficient $-\Delta_m\lambda_{k,m}w'_m(k)$, which equals
$-\Delta_m\lambda_{k,m}\varsigma_kM'_m$ at a peak $k$ of sign $\varsigma_k$.
\item[(b)] If $(b^\pm,\omega^\pm)$ are two-piece data at $f'$, then
$b^+-b^-=\mathcal V(\Delta\omega)$; hence $\mathcal V(\Delta\omega)$ is
$z'$-admissible off $F'$.
\item[(c)] Conversely, suppose $\mathcal V:=\mathcal V(\Delta\omega)$ is
$z'$-admissible off $F'$, let $\beta\in\ell_1$ be supported in $F'$ and
$\chi:K'\to[0,1]$, and put $b^+:=\beta+\chi\mathcal V1_{K'}$, $b^-:=b^+-
\mathcal V$.  Then $(b^+,\omega^+)$ and $(b^-,\omega^-)$ are two-piece data at
$f'$ for $g:=b^++\sum_mR_m^*(\omega^+_m-d'_m(\omega^+_m)w'_m)$, with
$\Delta d_m=\Delta_m$.
\end{enumerate}
\end{lemma}

So the existence of exact two-piece data with prescribed block parts is a
property of the single vector $\Delta\omega$: $\mathcal V(\Delta\omega)$ must
be $z'$-admissible off $F'$.  Off $F'$ its restriction is a finite sum over
the carriers of $\Omega$ plus the infinite series
\[
 -\sum_m\Delta_m\sum_{k\notin\Omega_m}\lambda_{k,m}w'_m(k)u_{k,m}.
\]

\begin{proof}
(a) $R_m^*x=\sum_k\lambda_{k,m}x(k)u_{k,m}$ [I,~(1)], and $\sum_{k,m}
\lambda_{k,m}\|u_{k,m}\|_1<\infty$.  (b) Subtracting the two representations
of the functional gives $b^+-b^-=\sum_mR_m^*\big((\omega^-_m-\omega^+_m)-
(d'_m(\omega^-_m)-d'_m(\omega^+_m))w'_m\big)=\mathcal V(\Delta\omega)$.  Off
$F'\cup K'$ both $b^\pm$ vanish, and on $K'$, $z'b^+\ge0\ge z'b^-$, so
$z'(b^+-b^-)\ge0$.  (c) Off $F'$, $b^+=\chi\mathcal V1_{K'}$ vanishes off $K'$
and is $z'$-signed on $K'$; $b^-=\chi\mathcal V-\mathcal V=-(1-\chi)\mathcal V$
on $K'$ and $b^-=0$ at the free coordinates (where $\mathcal V=0$), so $b^-$
is $(-z')$-signed.  The block parts are finitely supported in $Q'_m$.  By
linearity of $d'_m$, $b^-+\sum_mR_m^*(\omega^-_m-d'_m(\omega^-_m)w'_m)=b^+-
\mathcal V+\sum_mR_m^*(\omega^+_m-d'_m(\omega^+_m)w'_m)+\mathcal V=g$.
\end{proof}

\begin{lemma}[Shift bound]\label{lem:II-shift}
Let $f'\in S_{p^*}$ have $F'$ finite and let $(b^\pm,\omega^\pm)$ be
two-piece data at $f'$ with $\Gamma'_w(b^\pm,\omega^\pm)\le2$.  Then for
every $m$, with $\Phi_{P'_m}^2:=\sum_{k\in P'_m}\Phi_m(k)^2>0$,
\[
 |\Delta d_m|\le\Big(\frac{8\,C'^3_m}{\sigma'_mM'^2_m\Phi_{P'_m}^2}\Big)^{1/2}.
\]
\end{lemma}

\begin{proof}
Drop $m$ and the primes.  $\sigma H(\omega^\pm)\le\Gamma_w\le2$, and
$\sqrt H$ is a seminorm (it is $\|P^\perp D\cdot\|_2/\sqrt C$, $P^\perp$
the projection onto $(Dw)^\perp$), so $H(\Delta\omega)\le8/\sigma$ for
$\Delta\omega:=\omega^--\omega^+$, which is supported in a finite set
$\Omega\subseteq Q$.  Put $v:=D(w1_\Omega)$ and $\chi:=\|v\|_2^2/C^2$.  As
$\|Dw\|_2^2=C^2\ge\|v\|_2^2+\|D(w1_P)\|_2^2=\|v\|_2^2+M^2\Phi_P^2$, we have
$1-\chi\ge M^2\Phi_P^2/C^2>0$.  If $v=0$ then $\Delta d=\ip{Dw}{D\Delta\omega}/C
=\ip{v}{D\Delta\omega}/C=0$.  Otherwise write $D\Delta\omega=\mathsf av+y$ with
$y\perp v$; then $\Delta d=\mathsf a\|v\|^2/C=\mathsf aC\chi$ and
$\|P^\perp D\Delta\omega\|^2=\|D\Delta\omega\|^2-\ip{Dw}{D\Delta\omega}^2/C^2\ge
\mathsf a^2C^2\chi(1-\chi)$.  Hence $CH(\Delta\omega)\ge\mathsf a^2C^2\chi(1-\chi)$
and $\Delta d^2=\mathsf a^2C^2\chi^2\le H(\Delta\omega)C\chi/(1-\chi)\le
(8/\sigma)C^3/(M^2\Phi_P^2)$.
\end{proof}

At companions $f^\#$ of $f$ with small cost (Lemma~\ref{lem:II-cost}) the
quantities $C^\#_m,\sigma^\#_m,M^\#_m$ are close to those of $f$, and a fixed
non-degenerate peak of $f$ remains a peak, so the shifts of all data with
$\Gamma^\#_w\le2$ are bounded by a constant depending only on $f$.

\begin{lemma}[Sandwich]\label{lem:II-sandwich}
Let $f\in S_{p^*}$ have $F$ finite and let $h\in X^*$ carry two-piece data
$(b^\pm,\omega^\pm)$.  Put $c^\pm:=\sum_md_m(\omega^\pm_m)\psi_m$.  Then
$b^+-b^-=\sum_mR_m^*(\omega^-_m-\omega^+_m)-W_\Delta$, and for every
$j\notin\Sigma(\omega)$:
\begin{enumerate}
\item[(a)] $h(j)=b^+(j)-c^+(j)=b^-(j)-c^-(j)$;
\item[(b)] if $j$ is a free coordinate \textup{(}$j\notin F\cup K$\textup{)}, then
$b^\pm(j)=0$, $h(j)=-c^+(j)=-c^-(j)$ and $W_\Delta(j)=0$;
\item[(c)] if $j\in K$, then $z_j(h(j)+c^+(j))\ge0\ge z_j(h(j)+c^-(j))$; in
particular $z_jW_\Delta(j)\le0$.
\end{enumerate}
\end{lemma}

\begin{proof}
The identity is the difference of the two representations.  For
$j\notin\Sigma(\omega)$ every $(R_m^*\omega^\pm_m)(j)=\sum_k\lambda_{k,m}
\omega^\pm_m(k)u_{k,m}(j)$ vanishes, and $R_m^*(d\,w_m)=d\,\psi_m$; so the
representations read $h(j)=b^\pm(j)-c^\pm(j)$.  Side admissibility gives
$b^\pm(j)=0$ off $F\cup K$, and $z_jb^+(j)\ge0\ge z_jb^-(j)$ on $K$.  Finally
$c^+-c^-=-W_\Delta$.
\end{proof}

No smallness or dominance hypothesis is needed; the lemma holds at every
coordinate outside the supports of the carriers used by the data.

\begin{corollary}[One block: constancy of the profile]\label{cor:II-S1}
Let $I=\{1\}$, $f\in S_{p^*}$ with $F$ finite, $\psi:=R_1^*w_1$, and let $h$
carry two-piece data, $d^\pm:=d_1(\omega^\pm_1)$, $\Delta:=d^--d^+$.  For
$j\notin\Sigma(\omega)$ with $\psi(j)\ne0$ put $\pi_h(j):=h(j)/\psi(j)$ (the
\emph{profile}); call a contact $j$ \emph{aligned} if $z_j\psi(j)>0$ and
\emph{anti-aligned} if $z_j\psi(j)<0$.  For $j\notin\Sigma(\omega)$ with
$\psi(j)\ne0$:
\begin{enumerate}
\item[(i)] if $j$ is free, then $\Delta=0$ and $\pi_h(j)=-d^+$;
\item[(ii)] if $j$ is an aligned contact, then $-d^+\le\pi_h(j)\le-d^-$
\textup{(}so $\Delta\le0$\textup{)};
\item[(iii)] if $j$ is an anti-aligned contact, then $-d^-\le\pi_h(j)\le-d^+$
\textup{(}so $\Delta\ge0$\textup{)}.
\end{enumerate}
Consequently the oscillation of $\pi_h$ over the aligned contacts outside
$\Sigma(\omega)$ is at most $|\Delta|$; and if outside $\Sigma(\omega)$ there
are an aligned contact and, in addition, a free coordinate or an
anti-aligned contact with $\psi\ne0$, then $\Delta=0$ and $h=-d^+\psi$ at
every coordinate outside $\Sigma(\omega)$.
\end{corollary}

\begin{proof}
Lemma~\ref{lem:II-sandwich} with $c^\pm=d^\pm\psi$.  (i) is (b) ($W_\Delta(j)=
\Delta\psi(j)=0$).  (ii), (iii): divide the inequalities of (c) by $z_j\psi(j)$.
For the last claim, the hypotheses give $\Delta\le0$ and $\Delta\ge0$
(or $\Delta=0$ directly); with $d^+=d^-$, (b) and (c) give $h(j)=-d^+\psi(j)$
at every free coordinate and every contact outside $\Sigma(\omega)$, and
$\Sigma(\omega)\supseteq F$.
\end{proof}

\begin{corollary}[Finitely many contacts: no switching, no shift]\label{cor:II-S2}
If $F$ and $K$ are finite \textup{(}in particular if $f$ attains its norm,
[I,~Proposition~1.10(c)]\textup{)}, then all two-piece data at $f$ have
$b^+=b^-$, $\omega^+=\omega^-$ and $\Delta d_m=0$ for every $m$.
\end{corollary}

\begin{proof}
By Lemma~\ref{lem:II-sandwich}, $v:=b^+-b^-=L^*X$ with $X_m:=(\omega^-_m-
\omega^+_m)-\Delta d_mw_m$, so $v\in Y$; and $v$ is supported in the finite
set $F\cup K$.  By (T-c), $v=0$, and as $L^*$ is injective on $V^*$
[I,~(1)], $\omega^-_m-\omega^+_m=\Delta d_mw_m$ for every $m$.  The left side
vanishes on $P_m\ne\emptyset$, where $|w_m|=M_m>0$; so $\Delta d_m=0$,
$\omega^-=\omega^+$ and $b^+=b^-$.
\end{proof}

At norm-attaining rows exact data cannot carry any shift or switching; the
scale-dependent structure of the mates of such rows lives only in their
decompositions at positive scales (compare [I,~Remark~7.11(a)]).

\begin{proposition}[Following a non-constant profile]\label{prop:II-S3}
Let $I=\{1\}$, $f\in S_{p^*}$ with $F$ finite, $g\in\Cset(f)$, $\rho\in(0,1]$,
and let $s_1,s_2\notin F$ with $\psi(s_1)\psi(s_2)\ne0$ and $\pi_g(s_1)\ne
\pi_g(s_2)$, where $\pi_g:=g/\psi$; put $\kappa_0:=|\psi(s_1)\psi(s_2)|\,
|\pi_g(s_1)-\pi_g(s_2)|>0$.  Let $f_n\to f$ be first rows with finite base
supports $F_n$ at which every mate carries two-piece data \textup{(}for
instance rows satisfying \textup{(BT)}, [I,~Theorem~7.10(c)]\textup{)}; let
$\psi_n:=R_1^*w_{n,1}$, $\Sigma_n:=F_n\cup\bigcup_{k\in Q_{n,1}}\supp u_{k,1}$,
and define aligned and anti-aligned contacts of $f_n$ with $z^{(n)}$ and
$\psi_n$.  Suppose that for infinitely many $n$:
\begin{enumerate}
\item[(1)] $s_1,s_2\notin\Sigma_n$, and
\item[(2)] outside $\Sigma_n$, $f_n$ has an aligned contact and, in addition,
a free coordinate with $\psi_n\ne0$ or an anti-aligned contact.
\end{enumerate}
Then, along these $n$,
\[
 \dist_{\ell_1}(\rho g,\Cset(f_n))\ge2\rho\Big(\kappa_0-\|g\|_\infty
 \big(|\psi_n-\psi|(s_1)+|\psi_n-\psi|(s_2)\big)\Big)\longrightarrow
 2\rho\kappa_0,
\]
and $\dist_{p^*}(\rho g,\Cset(f_n))\ge\frac13\dist_{\ell_1}(\rho g,\Cset(f_n))$;
in particular $\rho g\notin\Li_n\Cset(f_n)$.
\end{proposition}

\begin{proof}
Let $G\in\Cset(f_n)$ with two-piece data at $f_n$; their block parts are
supported in $Q_{n,1}$, so $\Sigma(\omega)\subseteq\Sigma_n$.  By (2) and
Corollary~\ref{cor:II-S1} at $f_n$, $G=-d^+\psi_n$ outside $\Sigma_n$, in
particular at $s_1,s_2$ by (1).  So the functional $\Lambda_n(h):=h(s_1)
\psi_n(s_2)-h(s_2)\psi_n(s_1)$ vanishes at $G$.  Since $\|\psi_n\|_\infty\le
\sum_k\lambda_{k,1}\|u_{k,1}\|_\infty\le\sum_k\lambda_{k,1}\le\frac12$, we
have $|\Lambda_n(h)|\le\frac12\|h\|_1$, so $\|\rho g-G\|_1\ge2|\Lambda_n(\rho
g)|$.  With $\Lambda(g):=g(s_1)\psi(s_2)-g(s_2)\psi(s_1)=\psi(s_1)\psi(s_2)
(\pi_g(s_1)-\pi_g(s_2))$, $|\Lambda(g)|=\kappa_0$ and $|\Lambda_n(g)-\Lambda(g)|
\le\|g\|_\infty(|\psi_n-\psi|(s_1)+|\psi_n-\psi|(s_2))$.  By
[I,~Proposition~1.11], $\psi_n\to\psi$ in $\ell_1$.  Finally, if $h=A+L^*W$
with $\max(q^*(A),\|W\|)=p^*(h)$ [I,~Lemma~1.6], then $\|h\|_1\le q^*(A)+
\sum_{k,m}\lambda_{k,m}\|W\|\le3p^*(h)$, since $\sum_{k,m}\lambda_{k,m}\le
\sum_mm2^{-m}\le2$.
\end{proof}

To recover a mate whose profile is not constant on $\{s_1,s_2\}$,
approximants at which all mates have exact data must therefore, for all large
$n$, put $s_1$ or $s_2$ into $\Sigma_n$ (make it a support coordinate, a
\emph{bank}, or cover it by the support of a strict non-peak), or be aligned
with maximal contact outside $\Sigma_n$ (no free coordinate and no
anti-aligned contact with $\psi_n\ne0$).  The engineered approximants of
[I,~Section~7] use the first alternative (window masses are banks on
contacts); norm-attaining rows can use only the first alternative, since
they have cofinitely many free coordinates.

\begin{theorem}[A criterion for the failure of lower semicontinuity along block-tame rows]\label{thm:II-NL}
Let $I=\{1\}$, let $f\in S_{p^*}$ have $F$ finite, $g\in\Cset(f)$ and
$s_1,s_2\notin F$ with $\psi(s_1)\psi(s_2)\ne0$ and $\pi_g(s_1)\ne\pi_g(s_2)$.
Let $f_n\to f$ be first rows satisfying \textup{(BT)} such that, for every
$n$, $s_1,s_2\notin\Sigma_n$ and, outside $\Sigma_n$, $f_n$ has an aligned
contact and a free coordinate with $\psi_n\ne0$ or an anti-aligned contact.
Then every $f_n$ belongs to $\Rec$, and for every $\rho\in(0,1]$
\[
 \liminf_n\dist_{p^*}(\rho g,\Cset(f_n))\ge\tfrac23\rho\,|\psi(s_1)\psi(s_2)|\,
 |\pi_g(s_1)-\pi_g(s_2)|>0 .
\]
In particular the fibre map $f'\mapsto\Cset(f')$ is not lower semicontinuous
at $f$ along the block-tame rows $f_n$; if $f$ itself satisfies \textup{(BT)},
then nevertheless $f\in\Rec$.
\end{theorem}

\begin{proof}
$f_n\in\Rec$ and (if (BT) holds at $f$) $f\in\Rec$ by [I,~Corollary~7.22];
every mate of $f_n$ carries two-piece data by [I,~Theorem~7.10(c)].  Apply
Proposition~\ref{prop:II-S3} to all $n$.
\end{proof}

\begin{remark}[Realization at self-aligned rows]\label{rem:II-NL}
For $T_{\mathrm{final}}$ and $N=1$ the hypotheses of
Theorem~\ref{thm:II-NL} are realized, with complete proof, in
Proposition~\ref{prop:II-NLreal} below (with $f=f_{\rm SA}$ of
Theorem~\ref{thm:II-SA} and the mates of
Proposition~\ref{prop:II-SAmates}).  \emph{Sketch, not proved here:} the
same construction should work, with $f$ satisfying \textup{(BT)}, for
signature-ladder designs that satisfy the dominance
condition (SF$^*$) and contain a target supported in two coordinates outside
all signature sets (condition (Z0)), with a diagonal base: there $f$ can be
taken to be a \emph{self-aligned} row (every carrier except one exceptional
strict non-peak $k_-$ is an exactly swallowed peak of its own sign, every
coordinate off $F$ is a contact, $Q=\{k_-\}$), $f_n$ is obtained from $f$ by
flipping $z$ on the signature set of one far carrier $o_n$ whose target is
supported in $F$ (which becomes an exactly swallowed anti-type peak, so its
contacts are anti-aligned) and re-assigning the later carriers by a
hysteretic owner rule, and $g$ is a coherent-shift mate of $f$ with a
coordinate-dependent split of its contact part, whose profile on a signature
set of a carrier $l\ne k_-$ is not constant.  The construction of these rows
and the verification of these properties belong to the design-specific part
of our sources and are not reproduced here in this generality.  For
$T_{\mathrm{final}}$ (Proposition~\ref{prop:II-NLreal}), and for any such
design for which the construction is granted, it follows from
Theorem~\ref{thm:II-NL} that membership in $\Rec$ cannot be propagated along an \emph{arbitrary} dense
sequence of block-tame rows, and approximants for mates with oscillating
profiles must either bank the oscillation coordinates or keep every peak of a
shifted block exactly swallowed and aligned outside the data supports.
\end{remark}

\subsection{Violation tolerance of the engineered approximants}\label{subsec:II-VT}

In the proof of [I,~Theorem~7.18] the sign conditions of two-piece data enter
only through the claim on the base excess in Step~3.  We show that a small
violation of these conditions costs a first-order term that is affordable at
a single stage of the construction.  We use the notation of
[I,~Definition~7.12] and of the proof of [I,~Theorem~7.18]: for pairs
$(b^\pm,\omega^\pm)$, $b^\theta:=\frac12(b^++b^-)$, $\omega^\theta:=\frac12
(\omega^++\omega^-)$, $v:=b^+-b^-$; the stage parameters $(N_w,s_1,N'')$; the
row $f'$ with data $a'$, $z'$, $\hat x'$, $c'$; the target $g'=\rho(g''-ca')$;
the vectors $\beta^\theta:=b^\theta1_{[1,N_w]}$,
$\beta^\pm:=b^\pm1_{[1,N_w]}\pm\frac12v1_{(N_w,\infty)}$, $B^\diamond:=
\rho(\beta^\diamond-ca')$; and $V_{>N''}:=\sum_{j\in K,\,j>N''}|v_j|$.

\begin{lemma}[Violation tolerance at a single stage]\label{lem:II-VT}
Let $I$ be finite, $f\in S_{p^*}$ with $F$ finite, $\rho\in(0,1)$, and
$g\in X^*$ with $g(\xi)=0$.  Let $(b^+,\omega^+)$ and $(b^-,\omega^-)$ be
pairs that represent $g$ in the sense of [I,~Definition~7.7]
\textup{(}$\omega^\pm_m\in c_{00}$, $\supp\omega^\pm_m\subseteq Q_m$\textup{)},
not necessarily side-admissible.  For $j\notin F$ put
\[
 \mathrm{viol}(j):=(z_jb^+_j)_-+(z_jb^-_j)_+\ \ (j\in K),\qquad
 \mathrm{viol}(j):=|b^+_j|+|b^-_j|\ \ (j\notin F\cup K),
\]
and $\epsilon:=\sum_{j\notin F}\mathrm{viol}(j)$ \textup{(}two-piece data are
exactly the case $\epsilon=0$\textup{)}.  Assume
\begin{enumerate}
\item[(V1)] $\rho^2\kappa_w<1$, where $\kappa_w:=\max_\pm\Gamma_w(b^\pm,
\omega^\pm)$;
\item[(V2)] $I_-:=\{m:\Delta d_m<0\}$ satisfies \textup{(SC)}
[I,~Definition~7.16] \textup{(}no assumption if $I_-=\emptyset$\textup{)};
\item[(V3)] $b^\theta_j=0$ at every free coordinate $j\notin F\cup K$ with
$\mathrm{viol}(j)>0$.
\end{enumerate}
Choose $\delta$, $T_1$, $K_\sharp$, $\eta_1$, the transfer data and $T_0$ as in
Step~0 of the proof of [I,~Theorem~7.18] for $(f,b^\pm,\omega^\pm,\rho)$, with
$K_\sharp$ replaced by $K_\sharp+1$.  Let $(N_w,s_1,N'')$ be a stage at which
all inequalities used in Steps~1--5 of that proof hold \textup{(}[I,~(11)--(13)],
the persistence of the transfer data, $|\lin'_m|\le K_\sharp|\tau|s_1$,
$\hat G_b,\hat G_m\le K_\sharp\tau^2$ for the exact parts, the bracket bound
$\Gamma_\diamond+\delta/8$ of Step~5, and $\frac12\rho|\tau|V_{>N''}+|\tau|o(s_1)
\le\delta\tau^2/16$ for $|\tau|>s_1$\textup{)}, and at which
\[
 \text{(VT)}\qquad 2\rho\epsilon\le\delta s_1/16 .
\]
Then the engineered approximant $f'$ \textup{(}norm attaining\textup{)} and
$g'$ satisfy
\[
 p^*(f'+\tau g')\le1+\frac{\tau^2}2(1-\delta)\qquad(0<|\tau|\le T_0).
\]
\end{lemma}

\begin{proof}
\emph{Where side admissibility is used.}  Step~0 uses the data only through
$\kappa_w$ (and $\Gamma_\theta\le\kappa_w$ by convexity of $\Gamma_w$, which
does not need admissibility) and through finitely many constants.  Steps~1
and~2 use only the two representations: $g''\to g$, $g'(\hat x')=0$, and the
identities $\beta^\pm-\beta^\theta=\pm\frac12v$ and $v=\sum_mR_m^*((\omega^-_m-
\omega^+_m)-\Delta d_mw_m)$, which holds for any two representations of the
same functional.  Step~3 (Blocks) uses only $\supp\omega^\diamond_m\subseteq
Q_m$ and the radius conditions, and [I,~Lemma~7.13] uses $b^\theta$ only
through the masses $\mathsf m_j=4\rho s_1|b^\theta_j|$ and $\|b^\theta\|_1$.
The sign conditions enter only the claim of Step~3 (Base),
$\Exc'(\tau B^\diamond)\le\frac12\rho|\tau|V_{>N''}$ (and $=0$ for
$\diamond=\theta$).  We replace it by
\begin{equation}\label{eq:II-VTstar}
 \begin{gathered}
 \Exc'(\tau B^\diamond)\le\tfrac12\rho|\tau|V_{>N''}+2\rho|\tau|\epsilon\qquad
 (\diamond=\sgn\tau,\ s_1<|\tau|\le T_0),\\
 \Exc'(\tau B^\theta)=0\qquad(0<|\tau|\le s_1).
 \end{gathered}
\end{equation}

\emph{Proof of \eqref{eq:II-VTstar}.}  By [I,~Lemma~7.1(b)] at $f'$,
$\Exc'(B)=\sum_j(|a'_j+B_j|-|a'_j|-z'_jB_j)$.  On $\supp a'$ the summand is
$2(-\sgn(a'_j)B_j-|a'_j|)_+$, and $a'_j+\tau B^\diamond_j=(1-\tau\rho c)a'_j+
\tau\rho\beta^\diamond_j$ with $1-\tau\rho c\ge\frac12$, so the summand is at
most $2(x)_-$ with $x:=z'_j\tau\rho\beta^\diamond_j$ (recall $\sgn a'_j=z'_j$).
Off $\supp a'$ the summand is $|\tau\rho\beta^\diamond_j|-z'_j\tau\rho
\beta^\diamond_j$.  We go through the coordinates.
\begin{itemize}
\item $j\in F$: no flip by condition $(T_1)$; summand $0$.
\item Window contacts with mass ($j\in K$, $j\le N_w$, $b^\theta_j\ne0$;
$z'_j=z_j$): for $\diamond=\theta$, $|\tau\rho b^\theta_j|\le\mathsf m_j/4\le
|a'_j|/2$ [I,~Lemma~7.13(a)], summand $0$ whatever the sign of $b^\theta_j$;
for $\diamond=+$ ($\tau>0$), $2(x)_-=2\rho\tau(z_jb^+_j)_-$; for $\diamond=-$
($\tau<0$), $2(x)_-=2\rho|\tau|(z_jb^-_j)_+$; both $\le2\rho|\tau|\,
\mathrm{viol}(j)$.
\item Window contacts without mass ($b^\theta_j=0$, so $\beta^\theta_j=0$ and
$b^+_j=-b^-_j$): off $\supp a'$; summand $0$ for $\theta$, and
$2\rho|\tau|(z_jb^+_j)_-$, resp.\ $2\rho|\tau|(z_jb^-_j)_+$, for $\diamond=\pm$.
\item Free $j\le N_w$: off $\supp a'$; for $\diamond=\theta$,
$\beta^\theta_j=b^\theta_j=0$ by (V3) if $\mathrm{viol}(j)>0$, and $b^\pm_j=0$
if $\mathrm{viol}(j)=0$: summand $0$; for $\diamond=\pm$ the summand is at
most $2\rho|\tau||b^\diamond_j|\le2\rho|\tau|\,\mathrm{viol}(j)$.
\item Contacts $j\in(N_w,N'']$: $\beta^\theta_j=0$, and for $\diamond=\sgn\tau$,
$\tau\beta^\diamond_j=\frac12|\tau|v_j$, so the summand is $\rho|\tau|
(z_jv_j)_-$; since $z_jv_j\ge-(z_jb^+_j)_--(z_jb^-_j)_+$, this is at most
$\rho|\tau|\,\mathrm{viol}(j)$.
\item Free $j>N_w$: $\beta^\theta_j=0$ and $|\beta^\pm_j|=|v_j|/2\le
\mathrm{viol}(j)/2$; summand at most $\rho|\tau|\,\mathrm{viol}(j)$.
\item Contacts $j>N''$: $z'_j=0$; summand $\frac12\rho|\tau||v_j|$ for
$\diamond=\pm$ and $0$ for $\theta$ (unchanged; this is the term with
$V_{>N''}$).
\end{itemize}
Free coordinates never belong to $\supp a'=F\cup\{j\in K\cap[1,N_w]:
b^\theta_j\ne0\}$.  Summing gives \eqref{eq:II-VTstar}.

\emph{Constants.}  For $|\tau|>s_1$, (VT) gives $2\rho|\tau|\epsilon\le\delta
|\tau|s_1/16\le\delta\tau^2/16\le\tau^2$; so the bound $\hat G_b\le(K_\sharp+1)
\tau^2$ holds, and Step~4 (rebalancing) is unchanged with $K_\sharp+1$ in
place of $K_\sharp$ (the choice of $\eta_1$ and the conditions $(T_0)$ were
made with $K_\sharp+1$), giving $E\le\delta\tau^2/8$.  In Step~5,
$\hat\Gamma=c'\hat G_b+\sum_m\sigma'_m\hat G_m$ gains at most $c'\cdot2\rho
|\tau|\epsilon\le\delta\tau^2/16$ in the regimes $\diamond=\pm$ and nothing in
the regime $\diamond=\theta$ ($c'\le1$).  Hence, for $0<|\tau|\le T_0$,
\[
 p^*(f'+\tau g')\le1+\frac{\tau^2}2\Big(\rho^2\kappa_w+\frac\delta8+
 \frac\delta8+\frac\delta8+\frac\delta4+\frac\delta8\Big)=1+\frac{\tau^2}2
 \Big(1-2\delta+\frac{3\delta}4\Big)\le1+\frac{\tau^2}2(1-\delta),
\]
since $\rho^2\kappa_w=1-2\delta$.
\end{proof}

\begin{remark}[Scope of Lemma~\ref{lem:II-VT}]\label{rem:II-VT}
(a) Contact violations inside the window need no further repair when (VT)
holds: their cost is the linear term $2\rho|\tau|\mathrm{viol}(j)$, paid only
for $|\tau|>s_1$; in the regime $|\tau|\le s_1$ the window masses protect every
window contact, whatever the sign of $b^\theta_j$.  Free violations inside
the window with $b^\theta_j\ne0$ are not tolerated: they would cost
$2\rho|\tau||b^\theta_j|$ for $|\tau|\le s_1$, which is not $O(\tau^2)$, and a
free coordinate cannot carry a mass (that would change $|z'_j|<1$ into
$1$); this is the reason for (V3).
(b) The lemma concerns a single stage and does not contain Step~6 of the
proof of [I,~Theorem~7.18]: a functional with violated data is in general not
a mate of $f$, so [I,~Lemma~7.6] has to be applied relative to some pair
$(f_0,g_0)$ with $g_0\in\Cset(f_0)$, which requires $T_0$ to be large compared
with $\sqrt{p^*(f'-f_0)}$.  Moreover (VT) bounds $s_1$ from \emph{below},
whereas the lateness conditions bound $s_1$ from above; the lemma is useful
only when $\epsilon$ is small compared with the lateness threshold of the
data.  (Numerical evidence; not used.)  A finite-model computation
indicates that the threshold (VT) is real: below it the bound fails in the
predicted window of linear flip cost.
(c) A uniform use of Lemma~\ref{lem:II-VT} across the scales of a window
meets a further difficulty: with window data whose switching has size $1/t$,
the window masses create a mismatch of order up to $s_1/t^2$ (an upper
bound; that it is attained is generic and supported only by numerical
evidence) between the
$\theta$-decomposition and the $\pm$-decompositions.  \emph{Sketch, not
proved here:} re-tuning at $f'$ the normalized values of the finitely many
coarse switching carriers exactly, so that $d'_m=d_m$ on every switching
vector, would remove this mismatch.
\end{remark}

%%%%%%%%%%%%%%%%%%%%%%%%%%%%%%%%%%%%%%%%%%%%%%%%%%%%%%%%%%%%%%%%%%%%%%%%

\section{The designed operator $T_{\mathrm{final}}$}\label{sec:II-Tfinal}

Mart\'in's construction works with every admissible operator $T$
[I, Remark~1.2, Definition~1.3], and the recovery results of
[I, Section~8] use a designed operator only through a short list
of properties [I, Theorem~8.2 and the remark following it].  In
this section we fix \emph{one} operator $T_{\mathrm{final}}$, together with
its base $U_{\mathrm{final}}$, which carries every design feature used in
the rest of this part, and we prove that it is admissible, that it does
not depend on the number $N$ of blocks, and that it has the window
properties on which the later sections rely.  The operator is a
signature-ladder design in the sense of [I, Definition~8.1] with
four modifications: a diagonal base whose entries decay doubly
exponentially; bounded-gap signature sets and an additional allowedness
rule; weights $c_l$ chosen \emph{after} the target and the signature size
of level $l$, as the minimum of finitely many explicit upper bounds; and
windows split into pigeonhole sub-windows governed by a design factor
$\mathrm{Des}(l)$ and a finite family of rate objects.

Throughout, $\mathfrak j:\N\times\N\to\N$ is a bijection with
$\mathfrak j(k,m)<\mathfrak j(k',m)$ for $k<k'$, and
$(k(l),m(l)):=\mathfrak j^{-1}(l)$.  Then $\mathfrak j(k,m)\ge k$ for all
$(k,m)$, and $k(1)=1$.  The notation of [I, Sections~1, 7 and~8]
is used without comment; in particular, for $f\in S_{p^*}$ we use the
forced data $(\xi,q_0,a,w,F,z,\zh,e,\nu)$ of
[I, Proposition~1.9], $M_m:=\|w_m\|_\infty$,
$C_m:=\|D_mw_m\|_2$, and the threshold constant
$\vartheta_m:=|R_m^{**}\xi|_mM_m/(mC_m)$ of [I, Section~1.3].

\subsection{The base}

\begin{definition}[The base $U_{\mathrm{final}}$]\label{def:II-base}
Let $H:=\ell_2$ with orthonormal basis $(k_s)_{s\ge1}$, put
\[
 \mu^\star_s:=2^{-2^{2^s}}\qquad(s\ge1),
\]
and let $U=U_{\mathrm{final}}:H\to c_0$ be the diagonal operator
$Uk_s:=\mu^\star_se_s$.
\end{definition}

\begin{lemma}[Properties of the base]\label{lem:II-base}
\textup{(a)} $U$ is compact, its range contains $c_{00}$ and is therefore
dense, $U^*e_s^*=\mu^\star_sk_s$, $U^*a=\sum_sa_s\mu^\star_sk_s$ is
injective on $\ell_1$, $\|U\|=\mu^\star_1=\frac1{16}$, and
\[
 q^*(a)=\|a\|_1+\Big(\sum_s(\mu^\star_sa_s)^2\Big)^{1/2}\qquad(a\in\ell_1).
\]
Consequently $\NA((c_0,q),\R)=c_{00}$, and every result of
[I] that is stated for an arbitrary compact operator $U$ with
dense range holds for $U_{\mathrm{final}}$.

\textup{(b)} \textup{(Diagonal identities.)}  Let $b\in\ell_1$ and
$s\notin\supp b$; then $\ip{U^*b}{k_s}=0$.  Let $a\in\ell_1$, let
$J\subseteq\N$ with $J\cap\supp a=\emptyset$ and let $(c_s)_{s\in J}$ be
summable; put $A:=a+\sum_{s\in J}c_se_s^*$ and
$X:=\sum_{s\in J}c_s\mu^\star_sk_s$.  Then $U^*A=U^*a+X$, $X\perp U^*a$,
$\|U^*A\|^2=\|U^*a\|^2+\|X\|^2$, and
$\ip{U^*u}{X}=\sum_{s\in J}c_s(\mu^\star_s)^2u(s)$ for every
$u\in\ell_1$.

\textup{(c)} \textup{(Decay.)}  $(\mu^\star_s)$ is strictly decreasing,
$\mu^\star_{s+1}=(\mu^\star_s)^{2^{2^s}}$, and for every $K>0$ there is
$C_K$ with $\mu^\star_s\le C_K2^{-Ks}$ for all $s$ \textup{(}the decay is
faster than every exponential\textup{)}.

\textup{(d)} \textup{(Depth function.)}  For $x\in(0,\frac12]$ let
$J(x):=\min\{J\ge1:\mu^\star_J\le x^2\}$.  Then
\[
 2^{J(x)}\le2\log_2\big(2\log_2(1/x)\big).
\]
\end{lemma}

\begin{proof}
(a) $U$ is the norm limit of its finite-rank truncations because
$\mu^\star_s\to0$; $Uk_s=\mu^\star_se_s$ gives $c_{00}\subseteq U(H)$.
For $a\in\ell_1$ and $h\in H$, $\ip{a}{Uh}=\sum_sa_s\mu^\star_s\ip{h}{k_s}$,
so $U^*a=\sum_sa_s\mu^\star_sk_s$; this vanishes only for $a=0$.  Since
$\mu^\star$ is decreasing (part (c)), $\|U\|=\sup_s\mu^\star_s=\mu^\star_1
=2^{-4}$.  The formula for $q^*$ is $q^*(a)=\|a\|_1+\|U^*a\|_H$
[I, Section~1.1].  The statement on $\NA((c_0,q),\R)$ is
[I, Proposition~1.10(c)], whose proof uses only that $U$ is
compact with dense range.

(b) $\ip{U^*b}{k_s}=\mu^\star_sb_s=0$.  The vector $X$ is supported (in
the basis $(k_s)$) in $J$ and $U^*a$ in $\supp a$, which is disjoint from
$J$; this gives orthogonality and the Pythagoras identity, and
$\ip{U^*u}{X}=\sum_{s\in J}u(s)\mu^\star_s\,c_s\mu^\star_s$.

(c) $2^{2^{s+1}}=(2^{2^s})^2$, so $\mu^\star_{s+1}=2^{-2^{2^s}\cdot
2^{2^s}}=(\mu^\star_s)^{2^{2^s}}<\mu^\star_s$.  Since $2^{2^s}/s\to\infty$,
$2^{2^s}\ge Ks$ for $s\ge s_K$, and $C_K:=\max_{s<s_K}2^{Ks}$ works.

(d) $\mu^\star_J\le x^2$ is equivalent to $2^{2^J}\ge2\log_2(1/x)$, i.e.
to $2^J\ge\log_2(2\log_2(1/x))=:Y$; note $Y\ge1$ because $x\le\frac12$.
If $J(x)=1$ the claim reads $2\le2Y$.  If $J(x)\ge2$, minimality gives
$2^{J(x)-1}<Y$.
\end{proof}

\begin{remark}[Why a doubly exponential base]\label{rem:II-whybase}
This remark is motivation only; nothing below depends on it.  The
infinite-support arguments of the later sections raise support
coefficients only beyond a depth $J$ and pay a Hilbert footprint of
order $\mu^\star_J$, while the pinning constant of the coordinates left
unraised is of order $2^J$.  A raise of footprint at most $x^2$ forces
$J\ge J(x)$, and the arguments require $2^{J(x)}=o(\log(1/x))$ as
$x\to0$, i.e.\ $\log_2(1/\mu_J)/2^J\to\infty$ for the base entries
$\mu_J$.  By Lemma~\ref{lem:II-base}(d) this holds for $\mu^\star$.  It
fails for $\mu_s=2^{-s^2-1}$ (then $2^{J(x)}$ is of order
$2^{\sqrt{2\log_2(1/x)}}$) and for $\mu_s=2^{-2^s}$ (then $2^{J(x)}$ is of
order $\log_2(1/x)$).  Faster than exponential decay of the base entries
is also what makes the raise-room condition of the earlier support
swallowing arguments hold for every power profile
$|a_s|\approx v_l(s)^{1+b}$; with entries $2^{-s}$ it fails for every
$b\ge2$.
\end{remark}

\subsection{The recursion}

\begin{definition}[Fixed data]\label{def:II-data}
\textup{(i)} \emph{Signature sets.}  $j_0:=1$ and
$S_l:=\{2^l(2i+1):i\ge1\}$ $(l\ge1)$.  The sets $S_l$ are pairwise
disjoint, infinite, consist of even numbers (so $j_0\notin S_l$), have
\emph{bounded gaps} $G_l:=2^{l+1}$ (consecutive elements differ by $G_l$),
and $\min S_l=3\cdot2^l$.  Put $Z_0:=\N\setminus\bigcup_lS_l$, the set of
odd numbers together with the powers $2^r$, $r\ge1$.

\textup{(ii)} \emph{Signatures.}
$h_l:=\sum_{s\in S_l}2^{-s}e_s^*$,
$H_l:=\|h_l\|_1=2^{-3\cdot2^l}/(1-2^{-2^{l+1}})$ and
$\delta^{\max}_l:=\min\{2^{-l},(4(1+\|U\|)H_l)^{-1}\}$.

\textup{(iii)} \emph{Tolerance sequence.}  $\iota_l:=2^{-l}$.

\textup{(iv)} \emph{Targets.}  A sequence $(y^{(i)})_{i\ge1}$ in
$c_{00}\cap S_{q^*}$ which is norm dense in $S_{q^*}$, with
$y^{(1)}:=e_1^*/q^*(e_1^*)$, and which contains
$y^\ast:=(e_3^*-e_5^*)/q^*(e_3^*-e_5^*)$ and
$y^{\ast\ast}:=(e_7^*-e_9^*)/q^*(e_7^*-e_9^*)$.

\textup{(v)} \emph{Schedule.}  $K_{\rm FD}:=\{2^r:r\ge1\}$; write
$\N\setminus K_{\rm FD}=\{k_1<k_2<\cdots\}$; $(i_r)_{r\ge1}$ is a sequence
in which every positive integer occurs infinitely often.
\end{definition}

\begin{definition}[The operator $T_{\mathrm{final}}$]\label{def:II-Tfinal}
The following quantities are defined by recursion on the stage
$l=1,2,\dots$; at stage $l$ all quantities of the stages $l'<l$ are
known, and the steps are carried out in the order listed.

\smallskip\noindent\textup{(1) Target.}  If $k(l)\in K_{\rm FD}$, let
$j_l$ be the least odd integer exceeding $\max\big(\{9\}\cup
\bigcup_{l'<l}\supp y_{l'}\big)$ and $y_l:=e_{j_l}^*/q^*(e_{j_l}^*)$ (a
\emph{fresh-dominated target}).  Otherwise $k(l)=k_r$ for a unique $r$;
call $y\in c_{00}$ \emph{allowed at $l$} if
\begin{enumerate}
\item[(a)] $\supp y\cap S_{l'}=\emptyset$ for every $l'\ge l$, and
\item[(c)] $\supp y\cap S_{l'}\subseteq[1,l]$ for every $l'<l$,
\end{enumerate}
and put $y_l:=y^{(i_r)}$ if $y^{(i_r)}$ is allowed at $l$, and
$y_l:=y^{(1)}$ otherwise.  (Every $y\in c_{00}$ supported in $Z_0$ is
allowed at every $l$.)

\smallskip\noindent\textup{(2) Signature size.}
$\delta_1:=\delta^{\max}_1$.  For $l\ge2$ put
$g_l:=\delta^{\max}_lH_l/(32\cdot3^{|\supp y_l|})$ and let $\delta_l$ be
the largest element of the set $A_l$ of all
$\delta\in[\frac12\delta^{\max}_l,\delta^{\max}_l]$ such that
\[
 \Big|\sum_{j\in\supp y_l}\sigma_jy_l(j)+\epsilon\delta H_l\Big|\ge g_l
 \qquad\text{for all }\sigma\in\{-1,0,1\}^{\supp y_l},\ \epsilon=\pm1 .
 \tag{GM}
\]

\smallskip\noindent\textup{(3) Normalisation.}
$n_l:=q^*(y_l+\delta_lh_l)$, $u_l:=(y_l+\delta_lh_l)/n_l$,
$v_l:=\delta_lh_l/n_l$, $\delta^\circ_l:=\delta_lH_l/n_l$,
$\underline y_l:=\min\{1,\min_{j\in\supp y_l}|y_l(j)|\}$ and
$\Lambda^\circ(l):=\prod_{l''\le l}(1+2/\delta^\circ_{l''})$.

\smallskip\noindent\textup{(4) Weight.}  $c_1:=1$.  For $l\ge2$, $c_l$ is
the minimum of the following finitely many positive numbers:
\begin{enumerate}
\item[(W1)] $c_{l-1}/4$;
\item[(W2)] $b(l-1,M(l-1))^2$ (defined in step (6) of stage $l-1$);
\item[(W3)] $(\delta_lH_l)^2$;
\item[(W4)] $\iota_{l'}2^{-2s}c_{l'}\delta_{l'}/2$, for every $l'<l$ and
every $s\in\supp y_l\cap S_{l'}$;
\item[(W5)] $c_{l-1}\delta_l2^{-\min S_l-l-10}/(1+\|U\|)$;
\item[(W6)] $g_l2^{m(l)+k(l)-l}$;
\item[(W7)] $3\cdot2^{-10}\iota_{l-1}\lambda_{l-1}\delta_{l-1}
2^{-(m(l-1)+k(l-1)+\min S_{l-1})}\underline y_{l-1}/
(n_{l-1}(1+\|U\|))$.
\end{enumerate}
Then $\Phi_l:=2^{-m(l)-k(l)}c_l$ and $\lambda_l:=m(l)\Phi_l$.

\smallskip\noindent\textup{(5) Level data.}
\begin{gather*}
 T(l):=\bigcup_{l''\le l}\supp y_{l''},\qquad
 s_{\max}(l):=\max(T(l)\cup\{1\}),\\
 s_{\rm b}(l):=\max_{l''\le l}\min\big(S_{l''}\cap(s_{\max}(l),\infty)\big),
 \qquad s_{\rm t}(l):=s_{\rm b}(l)+2^{l+2},\\
 \mathfrak m^{\rm nat}_{l''}(l):=\|v_{l''}1_{S_{l''}\setminus([1,l]\cup
 T(l))}\|_1\ (l''\le l),\\
 D_{\rm room}(l):=1+\sum_{l''\le l}\Big(\frac1{\mathfrak m^{\rm nat}_{l''}(l)}
 +\frac1{\Phi_{l''}}\Big),\\
 \delta_{\min}(l):=\min_{l''\le l}\delta_{l''},\qquad G(l):=2^{l+1},\qquad
 B^\star(l):=\frac{2^{s_{\rm t}(l)}}{\delta_{\min}(l)\,
 (\mu^\star_{s_{\rm t}(l)})^2} .
\end{gather*}
Further, the design constants of level $l$ of
Definition~\ref{def:II-constants} and the finite set $\mathcal O(l)$ of
rate objects of level $l$ of Definition~\ref{def:II-rates} are formed, and
$\omega(l):=|\mathcal O(l)|$.  The \emph{design factor} of level $l$ is
\begin{align*}
 \mathrm{Des}(l):={}&\Big[l\,2^{l^3}B^\star(l)\,2^{s_{\max}(l)}2^{G(l)}
 \delta_{\min}(l)^{-1}\Lambda^\circ(l)(1+|T(l)|)\\
 &\qquad\times D_{\rm room}(l)\,H_{\rm conf}(l)
 G^{**}(l)H_{\rm tune}(l)\Big]^6\\
 &\times\mathrm{Lip}(l)\Big(2+\frac1{\delta_{\rm comb}(l)}\Big)
 (r(l)\,c'(l))^l\,C_*(l)\,C_L(l)\,N_L(l)\Big(1+\frac1{\delta_{\rm sh}(l)}
 \Big).
\end{align*}

\smallskip\noindent\textup{(6) Sub-windows.}  $M(l):=\omega(l)+1$.  The
\emph{sub-windows of level $l$} are the pairs $w=(l,i)$,
$1\le i\le M(l)$; all sub-windows are ordered lexicographically, and
$w^-$ (resp.\ $w^+$) denotes the predecessor (resp.\ successor) of $w$.
For $w=(l,1),\dots,(l,M(l))$ in this order:
\begin{gather*}
 u(w):=\tfrac14\ \ (w=(1,1)),\qquad u(w):=b(w^-)\ \ (w\ne(1,1)),\\
 Q(w):=\Big(\frac{4\,\mathrm{Des}(l)}{u(w)}\Big)^{\omega(l)+20},\\
 n(w):=\big\lceil l\,2^{l^3}Q(w)\big\rceil,\qquad
 T_{\rm lo}(w):=2^{-n(w)}T_{\rm hi}(w),\\
 T_{\rm hi}(w):=\min\Big\{T_{\rm lo}(w^-),\frac{2^{-l^3}}{l\,Q(w)}\Big\}
 \ \ (T_{\rm lo}((1,1)^-):=1),\\
 \eta(w):=\frac{T_{\rm lo}(w)^4}{l\,\mathrm{Des}(l)},\qquad
 b(w):=\frac{\min\{\eta(w),(\eta(w)/C_L(l))^{N_L(l)/2}\}}
 {1+\mathrm{Lip}(l)C_*(l)} .
\end{gather*}
The interval $\mathcal W(w):=[T_{\rm lo}(w),T_{\rm hi}(w)]$, which contains
the $n(w)$ dyadic scales $T_{\rm hi}(w)2^{1-i}$ $(1\le i\le n(w))$, is the
\emph{sub-window} $w$, and $\mathrm{Band}(w):=(b(w),u(w))$ is its
\emph{band}.

\smallskip\noindent\textup{(D2)} $T_{\mathrm{final}}e_{k,m}:=
c_{\mathfrak j(k,m)}u_{\mathfrak j(k,m)}$, extended linearly and
continuously to $\ell_1(\N\times\N)$.
\end{definition}

For $l=\mathfrak j(k,m)$ we write $u_{k,m}=u_l$, $\Phi_m(k)=\Phi_l$ and
$\lambda_{k,m}=\lambda_l$, in accordance with [I, Section~1.2]
(this is justified in Theorem~\ref{thm:II-lemmaB}).  Note that
$\Phi_l\le c_l/4$ and $\lambda_l=m(l)2^{-m(l)-k(l)}c_l\le c_l/4$, because
$m2^{-m}\le\frac12$ and $2^{-k}\le\frac12$.

\subsection{Design constants and rate objects}

In the next two definitions $l\ge1$ is fixed and all data of the stages
$\le l$ up to step (5) of Definition~\ref{def:II-Tfinal} are known
($y_{l''},\delta_{l''},n_{l''},u_{l''},c_{l''},\Phi_{l''}$ for $l''\le l$, and
$T(l)$, $s_{\max}(l)$, $s_{\rm t}(l)$, $\mathfrak m^{\rm nat}_{l''}(l)$,
$D_{\rm room}(l)$, $\Lambda^\circ(l)$).  A number is called a \emph{design
number of level $l$} if it is determined by these data and by finitely
many combinatorial choices, and does not depend on $N$.

\begin{definition}[Patterns and design constants]\label{def:II-constants}
\textup{(a)} \emph{Patterns.}  A \emph{pattern of level $l$} is a tuple
$\kappa=(E,P,\epsilon,F',\mathrm{type})$ with $E\subseteq[1,l]$,
$P\subseteq E$, $\epsilon\in\{\pm1\}^E$, $F'\subseteq T(l)$ and
$\mathrm{type}:T(l)\setminus F'\to\{1,-1,0\}$; its \emph{kept set} is
$\mathrm{Kp}(\kappa):=E\setminus P$.  With
$L_j(\tau):=\sum_{l''\in E}\epsilon_{l''}\tau_{l''}u_{l''}(j)$, its
\emph{zero-cost cone} is
\[
 \begin{split}
 C(\kappa):=\{\tau\in\R^E:{}&\ \tau_{l''}\ge0\ (l''\in\mathrm{Kp}(\kappa)),\
 \tau_{l''}=0\ (l''\in P),\\ &\ \mathrm{type}(j)L_j(\tau)\ge0\ (\mathrm{type}(j)
 =\pm1),\ L_j(\tau)=0\ (\mathrm{type}(j)=0)\}.
 \end{split}
\]
$C(\kappa)\subseteq[0,\infty)^E$ is a pointed polyhedral cone; $\mathrm{Ext}
(\kappa)$ denotes the finite set of its extreme rays normalised by
$\|r\|_1=1$.  A \emph{component} of $\kappa$ is a pair $(r,m)$ with
$r\in\mathrm{Ext}(\kappa)$ and $m\in\{m(l''):l''\in\supp r\}$.

\textup{(b)} \emph{Tuning constant.}  For a set $\Sigma$ of components of
a pattern $\kappa$ let $R_\Sigma$ be the matrix with one row
$(r(l'')1_{m(l'')=m})_{l''\le l}$ per component $(r,m)\in\Sigma$, and
$R_\Sigma^+$ its Moore--Penrose inverse.  $H_{\rm tune}(l):=\max\{1,
\max_{\kappa,\Sigma}(\#\Sigma)^{1/2}\|R^+_\Sigma\|_2\}$.

\textup{(c)} \emph{Determinantal data.}  For a pattern $\kappa$ and
$v\in\R^{[1,l]}$ let $A_\kappa(v)$ be the matrix with columns indexed by
$E$ and the following rows: the unit rows $e_{l''}$ $(l''\in E)$; the
\emph{coordinate rows} $(\epsilon_{l''}u_{l''}(j))_{l''\in E}$ for
$j\in T(l)\setminus F'$; and, for every block $m\in\{m(l''):l''\in E\}$,
the \emph{value row} $(v_{l''}1[l''\in\mathrm{Kp}(\kappa),\
m(l'')=m])_{l''\in E}$.  For row and column index sets $J,K$ with
$|J|=|K|\ge1$ put $\pi_{\kappa,J,K}(v):=\det A_\kappa(v)_{J,K}$, a
polynomial in $v$ whose coefficients are design numbers.  Let
$\delta_{\rm comb}(l)$ be the least nonzero value of
$|\pi_{\kappa,J,K}|$ over all $(\kappa,J,K)$ with $J$ free of value rows
($\delta_{\rm comb}(l):=1$ if there is none); let $\mathrm{Lip}(l)\ge1$ be
a common Lipschitz constant, with respect to the maximum norm
$\|\cdot\|_\infty$ on $\R^{[1,l]}$, on $[-2,2]^{[1,l]}$ of all
$\pi_{\kappa,J,K}$ whose $J$ contains a value row (polynomials are Lipschitz
for every norm on a compact cube, so such a constant exists); and let $r(l)$, $c'(l)$ be the maximal
numbers of rows and of columns of the matrices $A_\kappa$.  Finally let
$N_L(l)\ge1$ be the least integer for which there is $C\ge1$ with the
following property, and $C_L(l):=1+\inf\{C\}$ for this integer: for every
subfamily $\mathcal S$ of these polynomials and every
$v\in[-2,2]^{[1,l]}$ with $\max_{\pi\in\mathcal S}|\pi(v)|\le\beta\le1$,
either $\beta\ge1/C$ or there is $v'$ with $\pi(v')=0$ for all
$\pi\in\mathcal S$ and $\|v'-v\|_2\le C\beta^{2/N_L(l)}$.  Such integers
exist by the {\L}ojasiewicz inequality \cite[Corollary~2.6.7]{BCR}, applied to the finitely many sums
of squares $\sum_{\pi\in\mathcal S}\pi^2$ and the distance functions to
their zero sets on the compact set $[-2,2]^{[1,l]}$; and $1+\inf\{C\}$ is
admissible because admissibility is preserved under enlarging $C$.  (If
the common zero set $Z_{\mathcal S}$ of $\mathcal S$ is empty, then
$\sum_{\pi\in\mathcal S}\pi^2\ge c_{\mathcal S}>0$ on the cube, and for $C$
large only the alternative $\beta\ge1/C$ can occur; otherwise the exponent
$2/N_L$ comes from $\dist(\cdot,Z_{\mathcal S})^{N}\le C\sum_{\pi\in
\mathcal S}\pi^2\le C|\mathcal S|\beta^2$.)

\textup{(d)} \emph{Configuration and shift constants.}  Let
$\mathcal R(l)$ be the finite set of pairs $(E,\mathbf r)$ with
$E\subseteq[1,l]$ and $\mathbf r\in\R^E$ either a unit vector $\pm e_{l''}$
$(l''\in E)$ or a coordinate row $(\epsilon_{l''}u_{l''}(j))_{l''\in E}$
$(\epsilon\in\{\pm1\}^E$, $j\in T(l))$.  A \emph{configuration system of
level $l$} is a system of linear conditions $\ip{\mathbf r}{\tau}\le\beta_{\mathbf r}$
or $\ip{\mathbf r}{\tau}=\beta_{\mathbf r}$ in $\tau\in\R^E$, each row $\mathbf r$
taken from $\{\mathbf r:(E,\mathbf r)\in\mathcal R(l)\}$ at most once; its
($\ell_1$-)\emph{Hoffman constant} is the least $H$ such that
$\dist_1(\tau,\mathcal P)\le H\cdot(\text{$\ell_1$-norm of the violations at
$\tau$})$ for all $\tau$ and all right-hand sides $\beta$ for which the
solution set $\mathcal P$ is nonempty; it is finite by Hoffman's theorem
\cite{Hoffman}.  $H_{\rm conf}(l)$ is the maximum of these constants over
the finitely many configuration systems of level $l$ (and $1$).
For a pattern $\kappa$ of level $l$ and box radii $\beta\in(0,\infty)^E$ put
$Z_\kappa(\beta):=\{\tau\in C(\kappa):\tau_{l''}\le\beta_{l''}\
(l''\in\mathrm{Kp}(\kappa))\}$.  The system of linear conditions defining
$Z_\kappa(\beta)$ is a configuration system of level $l$; let $G^{**}(l)$ be
the maximum over all patterns $\kappa$ of level $l$ of its
$\ell_1$-Hoffman constant (and $1$).  It does not depend on $\beta$, and
$G^{**}(l)\le H_{\rm conf}(l)$ (see also Definition~\ref{def:II-pattern}).
Further $C_*(l):=|T(l)|+2$, and $\delta_{\rm sh}(l)$ is the number defined
in \textup{(e)} below.

\textup{(e)} \emph{Combinatorial shift cost.}  Let $\kappa^{\rm sh}=(\kappa,
N_0,R_7,I_{\rm sh},\mathrm{src},G_{\rm pk})$ be an enlarged shift pattern of
level $l$ (the index set of (O9) in Definition~\ref{def:II-rates} below),
$\kappa=(E,P,\epsilon,F',\mathrm{type})$, and let $\mathrm{vs}\in\{\pm1\}^{
G_{\rm pk}}$.  Put $\Pi_m:=\sum_{l''\in G_{\rm pk},\,m(l'')=m}\mathrm{vs}_{l''}
\lambda_{l''}u_{l''}$ for blocks $m\in I_{\rm sh}$, $G_{\rm np}:=E\setminus
G_{\rm pk}$, and, for $\delta\in\R^{I_{\rm sh}}$ and $x\in[0,\infty)^{G_{\rm
np}}$, $L:=\sum_m\delta_m\Pi_m+\sum_{l''\in G_{\rm np}}x_{l''}\epsilon_{l''}
u_{l''}$.  Let
\[
 \begin{split}
 c^{\rm comb}(\delta;\kappa^{\rm sh},\mathrm{vs}):=\inf_{x\ge0}\Big[&\sum_{j\in
 T(l)\setminus F',\,\mathrm{type}(j)=\pm1}2(\mathrm{type}(j)L(j))_-\\
 &+\sum_{j\in T(l)\setminus F',\,\mathrm{type}(j)=0}|L(j)|
 +\sum_{l''\in E}2\mathfrak m^{\rm nat}_{l''}(l)(\epsilon_{l''}
 \mathrm{coef}_{l''})_-\Big],
 \end{split}
\]
where $\mathrm{coef}_{l''}$ is the coefficient of $v_{l''}$ in $L$ on
$S_{l''}\setminus([1,l]\cup T(l))$, namely $\epsilon_{l''}x_{l''}$ for
$l''\in G_{\rm np}$ and $\mathrm{vs}_{l''}\lambda_{l''}\delta_{m(l'')}$ for
$l''\in G_{\rm pk}$ (on this set every $u_{l'''}$ with $l'''\le l$,
$l'''\ne l''$, vanishes, by the disjointness of the signature sets and
Definition~\ref{def:II-Tfinal}(1)).  Let $c^{\rm comb}_*(\kappa^{\rm sh},
\mathrm{vs})$ be the minimum of $c^{\rm comb}(\cdot;\kappa^{\rm sh},
\mathrm{vs})$ over the compact set $\{\delta:\|\delta\|_1=1,\ \delta_m\le0$
if $\mathrm{src}(m)$ is ``$\le$'', $\delta_m\ge0$ if it is ``$\ge$''$\}$
(it is attained: $c^{\rm comb}$ is the minimum of a polyhedral function of
$(\delta,x)$ with design coefficients, and it is continuous in $\delta$).
Then $\delta_{\rm sh}(l)$ is the least positive value of $c^{\rm comb}_*(
\kappa^{\rm sh},\mathrm{vs})$ over all pairs $(\kappa^{\rm sh},\mathrm{vs})$
of level $l$, and $\delta_{\rm sh}(l):=1$ if there is none.  Since there
are finitely many pairs and the coefficients involve only $u_{l''}$
$(l''\le l)$, $\lambda_{l''}$ and $\mathfrak m^{\rm nat}_{l''}(l)$,
$\delta_{\rm sh}(l)$ is a design number of level $l$.
\end{definition}

Only the following properties of the constants of
Definition~\ref{def:II-constants} are used in this section; their precise
role is explained where they are used.
\begin{enumerate}
\item[(K1)] Each of $H_{\rm tune},\mathrm{Lip},r,c',C_L,N_L,H_{\rm conf},
G^{**},C_*$ is a real number $\ge1$, and $\delta_{\rm
comb}(l),\delta_{\rm sh}(l)\in(0,\infty)$.
\item[(K2)] Each is a design number of level $l$: the patterns, cones,
extreme rays and polynomials depend only on the finitely many numbers
$u_{l''}(j)$ $(l''\le l$, $j\in T(l))$, on the blocks $m(l'')$
$(l''\le l)$ and on finitely many sign, contact and status choices; the
Hoffman constants are finite (Hoffman's theorem \cite{Hoffman}), and do
not depend on right-hand sides.
\item[(K3)] None of them depends on $N$: blocks enter only as the blocks
$m(l'')$ of carriers $l''\le l$.
\end{enumerate}

\begin{definition}[Rate objects]\label{def:II-rates}
Let $N\ge1$ and $f\in S_{p_N^*}$ with forced data
$(\xi,q_0,a,w,\allowbreak F,z,\zh,e,\nu)$ ($F$ finite or infinite).  For a carrier
$l''=\mathfrak j(k,m)$ with $m\le N$ put
$\varrho_{l''}(f):=|u_{l''}(\xi)|/(\Phi_{l''}\vartheta_m)$, and, given a sign
$\epsilon_{l''}$, $\mathrm{val}_{l''}(f):=\epsilon_{l''}u_{l''}(\zh)$;
put $\gamma_k(f):=\ip{U^*u_k}{U^*a}$.  The \emph{rate objects of level
$l$} are the following indexed families; the \emph{value} of each object
at $f$ is a number in $[0,\infty]$.  \emph{Convention}: an object that
involves a carrier $l''$ with $m(l'')>N$ has the value $1$.
\begin{enumerate}
\item[(O1)] \emph{Rooms}, indexed by $l''\le l$: with
$S^{\rm nat}:=S_{l''}\setminus(F\cup T(l))$, the value
$\min_{\sigma=\pm1}\sum_{s\in S^{\rm nat}}v_{l''}(s)(1+\sigma z_s)/
\|v_{l''}1_{S^{\rm nat}}\|_1$ ($:=1$ if $S^{\rm nat}=\emptyset$).
\item[(O2)] \emph{Target rooms}, indexed by $j\in T(l)$: $1-|z_j|$
($:=1$ if $j\in F$).
\item[(O3)] \emph{Threshold distances}, indexed by $l''\le l$:
$|\varrho_{l''}(f)-1|$.
\item[(O4)] \emph{Relative positions}, indexed by $l''\le l$:
$\varrho_{l''}(f)$.
\item[(O5)] \emph{Ray components}, indexed by $(\kappa,r,m)$, $(r,m)$ a
component of a pattern $\kappa$: $|D_{r,m}(f)|/\Phi_{\max}(r,m)$, where
$D_{r,m}(f):=\sum_{l''\in\supp r,\ m(l'')=m}r(l'')\epsilon_{l''}
u_{l''}(\zh)$ and $\Phi_{\max}(r,m):=\max\{\Phi_{l''}:l''\in\supp r,\
m(l'')=m\}$ (the signs $\epsilon$ are those of $\kappa$).
\item[(O6)] \emph{Shift costs}, indexed by the shift patterns
$\pi=(I_{\rm up},I_{\rm lo},P^\star,\mathrm{vs},B_f,\epsilon')$ of level
$l$: $I_{\rm up},I_{\rm lo}$ disjoint sets of blocks of carriers $\le l$,
$P^\star\subseteq[1,l]$ with $m(P^\star)\subseteq I_{\rm up}\cup I_{\rm
lo}$, $\mathrm{vs}\in\{\pm1\}^{P^\star}$, $B_f\subseteq[1,l]\setminus
P^\star$, $\epsilon'\in\{\pm1\}^{B_f}$.  With
$\Pi_m:=\sum_{l''\in P^\star,\,m(l'')=m}\mathrm{vs}_{l''}\lambda_{l''}
u_{l''}1_{F^c}$ and $\phi_z(x)=|x|-zx$,
\[
 c(\delta;\pi):=\inf_{x\in[0,\infty)^{B_f}}\sum_{j\notin F}\phi_{z_j}
 \Big(\sum_m\delta_m\Pi_m(j)+\sum_{l''\in B_f}x_{l''}\epsilon'_{l''}
 u_{l''}(j)\Big),
\]
and the value is $c_\pi(f):=\inf\{c(\delta;\pi):\|\delta\|_1=1,\
\delta_m\ge0\ (m\in I_{\rm up}),\ \delta_m\le0\ (m\in I_{\rm lo})\}$
($:=\infty$ if $I_{\rm up}\cup I_{\rm lo}=\emptyset$).
\item[(O7)] \emph{Joint ray objects}, indexed by $(\kappa,\mathcal J)$
with $\emptyset\ne\mathcal J\subseteq\mathrm{Ext}(\kappa)$: the smallest
singular value of the matrix with columns
$\big(D_{r,m}(f)/\Phi_{\max}(r,m)\big)_m$, $r\in\mathcal J$ (rows indexed
by the blocks of carriers $\le l$, entries $0$ for blocks not meeting
$\supp r$; the smallest singular value is $0$ if there are more columns
than rows).
\item[(O8)] \emph{Minors}, indexed by $(\kappa,J,K)$ as in
Definition~\ref{def:II-constants}(c) with $J$ containing a value row:
$|\pi_{\kappa,J,K}((\mathrm{val}_{l''}(f))_{l''\in\mathrm{Kp}(\kappa)})|$.
\item[(O9)] \emph{Shift-pinning rates}, indexed by the enlarged shift
patterns $\kappa^{\rm sh}=(\kappa,N_0,R_7,\allowbreak I_{\rm sh},\mathrm{src},
\allowbreak G_{\rm pk})$: $\kappa=(E,P,\epsilon,F',\mathrm{type})$ a pattern,
$N_0,R_7,G_{\rm pk}\subseteq E$, $I_{\rm sh}$ a set of blocks of carriers
$\le l$, $\mathrm{src}:I_{\rm sh}\to\{\le,\ge,\text{free}\}$.  Let
$\Sigma^{\rm sh}$ be the system in $(\delta,\tau)\in\R^{I(l)}\times\R^E$,
$I(l):=\{m(l''):l''\le l\}$:
(S1) $\tau_{l''}\ge0$ $(l''\in E)$;
(S2) $\tau_{l''}=\epsilon_{l''}\mathrm{vs}_{l''}\lambda_{l''}
\delta_{m(l'')}$ $(l''\in G_{\rm pk})$, $\mathrm{vs}_{l''}:=\sgn
w_{m(l'')}(k(l''))$;
(S3) $\mathrm{type}(j)V_j(\tau)\ge0$ if $\mathrm{type}(j)=\pm1$ and
$V_j(\tau)=0$ if $\mathrm{type}(j)=0$, $j\in T(l)\setminus F'$, where
$V_j(\tau):=\sum_{l''\in E}\epsilon_{l''}\tau_{l''}u_{l''}(j)$;
(S4) $\delta_m+\sum_{l''\in E,\,m(l'')=m}q'_{l''}\tau_{l''}=0$ $(m\in I(l))$,
$q'_{l''}:=\epsilon_{l''}\Phi_{l''}w_m(k(l''))/(mC_m)$ for $l''\notin N_0$
and $q'_{l''}:=0$ for $l''\in N_0$;
(S5) $\delta_m=0$ $(m\notin I_{\rm sh})$;
(S6) $\delta_m\le0$, resp.\ $\delta_m\ge0$, if $\mathrm{src}(m)$ is $\le$,
resp.\ $\ge$;
(S7) $\tau_{l''}+\lambda_{l''}(|w_m(k(l''))|/M_m)\delta_{m(l'')}\le0$
$(l''\in R_7)$.
With $\mathrm{viol}(\delta,\tau)$ the sum of the positive parts of the
violations of the inequalities and of the moduli of the defects of the
equations, the value is
$\rho^{\rm sh}(\kappa^{\rm sh},f):=\inf\{\mathrm{viol}(\delta,\tau):
\|\delta\|_1=1,\ \tau\in\R^E\}$.
\item[(O10)] \emph{Support moduli}, indexed by $s\in[1,s_{\rm t}(l)]$:
$|a_s|$.
\item[(O11)] \emph{Spreads}, indexed by $(l'',s,s')$, $l''\le l$, $s<s'$
among the first three elements of $S_{l''}\cap(s_{\max}(l),\infty)$:
$|1-\mathfrak a_{s'}/\mathfrak a_s|$ with $\mathfrak a_s:=a_s/v_{l''}(s)$ if
$a_s\ne0\ne a_{s'}$, and $1$ otherwise.
\item[(O12)] \emph{Resource minors}, indexed by $(L,K)$ with
$\emptyset\ne L\subseteq[1,l]$, $K\subseteq[1,s_{\rm t}(l)]$, $|K|=|L|$:
$|\det(\Gamma(f)_{L,K})|$, where
$\Gamma(f)_{k,s}:=u_k(s)-\gamma_k(f)a_s/\nu^2$.
\end{enumerate}
$\mathcal O(l)$ is the disjoint union of the index sets of
(O1)--(O12), and $\omega(l):=|\mathcal O(l)|$.
\end{definition}

\begin{lemma}[The rate scheme is a design datum]\label{lem:II-ratescheme}
For every $l$, $\mathcal O(l)$ is a finite set which is determined by the
data of the stages $\le l$ up to step (5), and does not depend on $N$ or
on $f$; $\omega(l)\ge3l+|T(l)|$.  For every $N$ and every
$f\in S_{p_N^*}$, each object of $\mathcal O(l)$ has a well-defined value in
$[0,\infty]$.
\end{lemma}

\begin{proof}
The index sets of (O1)--(O4) are $[1,l]$, $T(l)$, $[1,l]$, $[1,l]$.
Patterns, shift patterns and enlarged shift patterns of level $l$ are
tuples of subsets of $[1,l]$, of $T(l)$ and of the set of blocks
$\{m(l''):l''\le l\}$, together with signs and types: finitely many.
$\mathrm{Ext}(\kappa)$ is finite (Minkowski--Weyl), and the row and column
index sets of $A_\kappa$ are finite.  The index sets of (O10)--(O12) are
subsets of $[1,s_{\rm t}(l)]$ and of $[1,l]$ and triples of
coordinates of the sets $S_{l''}$ determined by $s_{\max}(l)$.  None of
these involves $N$ or $f$.  The bound on $\omega(l)$ counts (O1)--(O4).
The values: $\varrho_{l''}$ is defined because $\Phi_{l''}>0$ and
$\vartheta_m>0$ ($M_m,C_m>0$ by [I, Lemma~1.7]); the
denominator in (O1) is positive when $S^{\rm nat}\ne\emptyset$; the series
in (O6) converge absolutely because $\sum_l\lambda_l\|u_l\|_1<\infty$ and
$|\phi_z(x)|\le2|x|$; $\nu>0$ by [I, Proposition~1.9(c)].  All
other values are finite minima, infima or determinants.
\end{proof}

\begin{remark}[Interpretation of (O10)--(O12)]\label{rem:II-newrates}
The families (O1)--(O9) are the rate objects of the finite-support
theory; (O10)--(O12) are needed only when $F$ is infinite.  If
$s\in S_{l''}\cap F$ lies beyond $s_{\max}(l)$, a change of $a_s$ by
$\Delta$ (with $z$ fixed) changes $\mathrm{val}_k$ at first order by
$((\mu^\star_s)^2/\nu)\epsilon_k\Gamma(f)_{k,s}\Delta$, and for
$s\notin F$ the column $\Gamma(f)_{\cdot,s}=(u_k(s))_k$ is the column of
a private bank.  So the minors (O12) control the regularity of the
first-order tuning resources at depth $\le s_{\rm t}(l)$.  In the
proportional case $a=c_kv_k$ on chosen coordinates $s_k\in S_k$ beyond
$s_{\max}(l)$ $(k\in L)$ one has $\Gamma_{k',s_k}=v_k(s_k)(1[k'=k]-
\gamma_{k'}c_k/\nu^2)$ for $k'\in L$, hence
$\det\Gamma_{L,\{s_k\}}=\prod_{k\in L}v_k(s_k)\,
(1-\sum_{k\in L}c_k\gamma_k/\nu^2)$, the Sherman--Morrison
denominator.  These statements are used and proved where needed; here
only the finiteness of the families matters.
\end{remark}

\subsection{Well-foundedness and admissibility}

\begin{lemma}[The recursion is well founded]\label{lem:II-W}
\textup{(a)} Every quantity of stage $l$ in
Definition~\ref{def:II-Tfinal} is a function of the fixed data of
Definition~\ref{def:II-data}, of quantities of the stages $l'<l$, and of
quantities of earlier steps of stage $l$.

\textup{(b)} All these quantities are finite; $A_l$ is a nonempty compact
set, so $\delta_l$ exists, and $\delta_l\in[\frac12\delta^{\max}_l,
\delta^{\max}_l]$; $n_l\in[\frac34,\frac54]$; $c_l$,
$\mathfrak m^{\rm nat}_{l''}(l)$, $g_l$, $u(w)$, $Q(w)$, $T_{\rm hi}(w)$,
$T_{\rm lo}(w)$, $\eta(w)$, $b(w)$ are positive; $D_{\rm room}(l)$,
$B^\star(l)$, $\mathrm{Des}(l)\in[1,\infty)$.

\textup{(c)} No quantity depends on $N$.
\end{lemma}

\begin{proof}
(a) Step (1) uses $\supp y_{l'}$ $(l'<l)$, the fixed sets $S_{l'}$ and the
schedule; allowedness consists of the support conditions (a), (c) only.
Step (2) uses $y_l$, $H_l$ and $\delta^{\max}_l$.  Step (3) uses $y_l$ and
$\delta_l$.  Step (4) uses $c_{l'}$, $\delta_{l'}$, $\lambda_{l'}$,
$n_{l'}$, $\underline y_{l'}$ $(l'<l)$, the number $b(l-1,M(l-1))$ from
step (6) of stage $l-1$, and $\delta_l,H_l,g_l,\supp y_l$ from steps
(1)--(2).  Step (5) uses data of index $\le l$, including $c_l$ (through
$\Phi_{l''}$ in $D_{\rm room}(l)$); the design constants and the rate
objects are design numbers of level $l$ (properties (K2) and
Lemma~\ref{lem:II-ratescheme}).  Step (6) uses $\mathrm{Des}(l)$,
$\omega(l)$, $C_L(l)$, $N_L(l)$, $\mathrm{Lip}(l)$, $C_*(l)$ and the values
$T_{\rm lo}(w^-)$, $b(w^-)$ of the predecessor.  Note that $\delta_l$ must
be chosen after $y_l$, since both $g_l$ and the inequalities (GM) involve
$y_l$, and that $c_l$ must be chosen after $\delta_l$; the order of the
steps is the only consistent one.

(b) For fixed $(\sigma,\epsilon)$ the set of $\delta$ violating (GM) is
an open interval of length $2g_l/H_l$; there are $2\cdot3^{|\supp y_l|}$
pairs, so the excluded set is open of total length at most
$4\cdot3^{|\supp y_l|}g_l/H_l=\delta^{\max}_l/8$.  Hence $A_l$ is closed,
bounded and of measure at least $\frac12\delta^{\max}_l-\frac18
\delta^{\max}_l>0$, and its largest element exists.  Next,
$q^*(y_l)=1$ and $q^*(\delta_lh_l)\le(1+\|U\|)\delta_lH_l\le\frac14$
because $\delta_l\le\delta^{\max}_l$; so $n_l\in[\frac34,\frac54]$.  The
numbers in (W1)--(W7) are positive: (W2) by induction, (W3), (W5) because
$\delta_l,H_l>0$, (W4) is a finite list ($\supp y_l$ is finite), (W6)
because $g_l>0$, (W7) because $\underline y_{l-1}>0$ and
$\lambda_{l-1}>0$; so $c_l>0$, and $\Phi_l,\lambda_l>0$.  The set
$S_{l''}\setminus([1,l]\cup T(l))$ is infinite and $v_{l''}>0$ on
$S_{l''}$, so $\mathfrak m^{\rm nat}_{l''}(l)>0$.  All factors of
$\mathrm{Des}(l)$ are finite and $\ge1$ by (K1) (note $\delta_{\min}(l)\le
1$, $\mu^\star_{s_{\rm t}(l)}<1$).  Inductively along the sub-windows,
$u(w)>0$, hence $Q(w),n(w)<\infty$, and $T_{\rm hi}(w),T_{\rm lo}(w),
\eta(w),b(w)>0$.

(c) The fixed data and steps (1)--(4) do not involve $N$; the design
constants and the rate scheme are $N$-free by (K3) and
Lemma~\ref{lem:II-ratescheme}; steps (5) and (6) are built from these.
\end{proof}

\begin{remark}[Two corrections of the literal union of earlier designs]
\label{rem:II-corrections}
(a) If the condition (GM) together with $\Phi_l\le g_l2^{-l}$ were imposed
also at $l=1$, it would contradict $c_1=1$: $S_1=\{6,10,14,\dots\}$ gives
$H_1=2^{-6}/(1-2^{-4})=\frac1{60}$, and $\delta^{\max}_1=\frac12$
(because $4(1+\|U\|)H_1<2$); (GM) with $\sigma=0$ requires
$g_1\le\delta_1H_1\le\frac1{120}$, and then $\Phi_1=2^{-m(1)-k(1)}c_1\le
g_1/2$ forces $c_1\le2^{m(1)+k(1)}/240=2^{m(1)+1}/240$ (recall
$k(1)=1$), e.g.\ $c_1\le\frac1{60}$ if $\mathfrak j(1,1)=1$.  Since
$c_l\le c_1/4$ for $l\ge2$, this gives $\|T\|=c_1<1$ whenever
$m(1)\le6$, in conflict with the equality clause of (T-a) and with
``norm one'' in Mart\'in's Lemma B.  This is why
(GM) and (W6) are imposed only for $l\ge2$; their only use concerns
carriers of large index.

(b) A bank factor calibrated to the base $2^{-j}$ (a factor
$8^{s_{\rm b}(l)}$ in place of $B^\star(l)$) does not dominate the bank
efficiency $(\mu^\star_{s'})^2v_c(s')$ of a fast-decaying base: target
supports may contain coordinates of $Z_0$ far beyond everything fixed at
earlier stages, so that $s_{\rm b}(l)$ is not controlled by the other
factors of the design.  (For the base $2^{-s^2-1}$ an explicit choice of
pure fresh targets $e_{s_i}^*/q^*(e_{s_i}^*)$ at a sparse set of first
schedulings makes the uncalibrated inequality fail at infinitely many
levels.  This is stated for orientation only and not used.)  The factor
$B^\star(l)$ removes the problem, see
Proposition~\ref{prop:II-bank}.
\end{remark}

\begin{theorem}[Admissibility; Mart\'in's Lemma B for $T_{\mathrm{final}}$]
\label{thm:II-lemmaB}
Let $T=T_{\mathrm{final}}$, $U=U_{\mathrm{final}}$.
\begin{enumerate}
\item[(a)] $T$ is admissible: \textup{(T-a)} $q^*(Te_{k,m})=
c_{\mathfrak j(k,m)}\le1$, with equality at $(k,m)=\mathfrak j^{-1}(1)$;
\textup{(T-b)} $T$ is injective; \textup{(T-c)}
$Y:=T(\ell_1(\N\times\N))$ satisfies $Y\cap c_{00}=\{0\}$;
\textup{(T-d)} for every $m$ the set $\{u_{\mathfrak j(k,m)}:k\in\N\}$ is
norm dense in $S_{q^*}$.
\item[(b)] \textup{(Conclusion of Mart\'in's Lemma B.)}  $\|T\|=1$, $T$ is
injective, $Y$ is a separable operator range which is dense in
$(\ell_1,q^*)=(c_0,q)^*$, $Y\cap\NA((c_0,q),\R)=\{0\}$, and for every $m$
the normalised vectors $Te_{k,m}/q^*(Te_{k,m})=u_{\mathfrak j(k,m)}$,
$k\in\N$, are dense in $S_{q^*}$, hence in the unit sphere of $Y$.
\item[(c)] For $l=\mathfrak j(k,m)$: $u_{k,m}=u_l$, $\Phi_m(k)=\Phi_l$,
$\lambda_{k,m}=\lambda_l$.  The pair $(y_l,c_l)$ satisfies allowedness
(a) and (b) of [I, Definition~8.1]:
$\supp y_l\cap S_{l'}=\emptyset$ for $l'\ge l$, and
$2c_l\le\iota_{l'}2^{-2s}c_{l'}\delta_{l'}\le2^{-2s}c_{l'}\delta_{l'}$ for
$l'<l$ and $s\in\supp y_l\cap S_{l'}$.
\item[(d)] \textup{(P1)} $y_l\in c_{00}$, $\|y_l\|_1\le1$,
$n_l\in[\frac34,\frac54]$, $\supp y_l\cap S_{l'}=\emptyset$ $(l'\ge l)$;
for $s\in S_l$: $u_l(s)=v_l(s)=\delta_l2^{-s}/n_l$, $u_{l'}(s)=0$
$(l'<l)$ and $u_{l'}(s)=y_{l'}(s)/n_{l'}$ $(l'>l)$.
\textup{(P2)} $c_{l+1}\le c_l/4$, $\sum_{l'>l}c_{l'}\le\frac43c_{l+1}$,
$\sum_{l'>l}\lambda_{l'}\le\frac13c_{l+1}\le\frac13b(l,M(l))^2$.
\end{enumerate}
\end{theorem}

\begin{proof}
(c), (d).  $q^*(Te_{k,m})=c_lq^*(u_l)=c_l$, so
$\Phi_m(k)=2^{-m-k}q^*(Te_{k,m})=\Phi_l$, $\lambda_{k,m}=m\Phi_m(k)=
\lambda_l$, and $u_{k,m}=Te_{k,m}/q^*(Te_{k,m})=u_l$.  The support
statement holds for targets chosen by allowedness (a) and for targets
supported in $Z_0$ (fresh-dominated targets and $y^{(1)}$).  The weight
inequality is (W4), since $\iota_{l'}\le1$.  (P1): $\|y_l\|_1\le q^*(y_l)=
1$ and $n_l\in[\frac34,\frac54]$ (Lemma~\ref{lem:II-W}(b)); for $s\in S_l$
the signatures $h_{l'}$, $l'\ne l$, vanish at $s$ (disjointness), and
$y_{l'}(s)=0$ for $l'\le l$ by the support statement.  (P2): (W1) gives
$c_{l'+1}\le c_{l'}/4$, hence $\sum_{l'>l}c_{l'}\le c_{l+1}\sum_{i\ge0}
4^{-i}=\frac43c_{l+1}$, and $\lambda_{l'}\le c_{l'}/4$ gives the bound on
$\sum\lambda$; finally $c_{l+1}\le b(l,M(l))^2$ by (W2).

(T-a): $c_l\le c_1=1$ by (W1), and $q^*(Te_{\mathfrak j^{-1}(1)})=c_1=1$.

(T-b), (T-c).  Let $x\in\ell_1(\N\times\N)$, indexed by $l$, with
$\mathcal Z:=Tx=\sum_lx_lc_lu_l\in c_{00}$, and suppose $x_{l'}\ne0$.
For $s\in S_{l'}$, by (P1),
\[
 \mathcal Z(s)=\frac{x_{l'}c_{l'}\delta_{l'}2^{-s}}{n_{l'}}+
 \sum_{l>l',\ s\in\supp y_l}\frac{x_lc_ly_l(s)}{n_l}.
\]
If the sum is not empty, let $L(s)$ be its least index; by (c) at stage
$L(s)$, $2c_{L(s)}\le2^{-2s}c_{l'}\delta_{l'}$, and since
$|y_l(s)|/n_l\le\frac43$ the sum is at most
$\frac43\|x\|_\infty\sum_{l\ge L(s)}c_l\le\frac{16}9\|x\|_\infty
c_{L(s)}\le\frac89\|x\|_\infty2^{-2s}c_{l'}\delta_{l'}$.  Hence
$|\mathcal Z(s)|\ge c_{l'}\delta_{l'}2^{-s}(|x_{l'}|/n_{l'}-\frac89
\|x\|_\infty2^{-s})>0$ for all large $s\in S_{l'}$, which contradicts
$\mathcal Z\in c_{00}$.  So $x=0$.  This proves both (T-b) and (T-c).

(T-d).  Fix $m$ and $i$.  The finite set $\supp y^{(i)}$ meets only
finitely many $S_{l'}$, so $y^{(i)}$ satisfies (a) and (c) at every
$l\ge l_i:=1+\max(\supp y^{(i)}\cup\{l':S_{l'}\cap\supp y^{(i)}\ne
\emptyset\})$.  There are infinitely many $r$ with $i_r=i$, and for them
$l=\mathfrak j(k_r,m)\ge k_r\to\infty$; for all such $l\ge l_i$ we have
$k(l)=k_r\notin K_{\rm FD}$ and $y_l=y^{(i)}$.  Along these $l$,
$|n_l-1|\le q^*(\delta_lh_l)\le\frac54\,2^{-l}H_l\to0$ and
$q^*(u_l-y^{(i)})\le|1/n_l-1|+q^*(\delta_lh_l)/n_l\to0$.  Hence every
$y^{(i)}$ is a limit of elements of $\{u_{\mathfrak j(k,m)}:k\in\N\}$, and
the targets are dense in $S_{q^*}$.

(b).  On $\ell_1(\N\times\N)$, $\|T\|=\sup_{(k,m)}q^*(Te_{k,m})=c_1=1$.
$Y$ is the range of a bounded operator on a separable Banach space, i.e.\
a separable operator range.  The closure of the linear space $Y$ contains
the dense subset $\{u_{\mathfrak j(k,1)}\}$ of $S_{q^*}$, hence the unit
ball, so $Y$ is dense.  $\NA((c_0,q),\R)=c_{00}$
(Lemma~\ref{lem:II-base}(a)), so $Y\cap\NA((c_0,q),\R)=\{0\}$ is (T-c).
The last statement is (T-d), since $u_{\mathfrak j(k,m)}\in Y$.
\end{proof}

\begin{remark}[Blocks and Mart\'in's construction]\label{rem:II-blocks}
At the stages $l=\mathfrak j(2^r,m)$, $r\ge1$, the target is a
fresh-dominated target which differs from block to block, so the
literal requirement ``every block follows the same target sequence'' of
[I, Remark~1.2] holds only at the stages with $k(l)\notin K_{\rm
FD}$, where $y_{\mathfrak j(k_r,m)}=y^{(i_r)}$ whenever this target is
allowed.  This is harmless: Mart\'in's proof, and everything below, uses
only the conclusion of Lemma~B, i.e.\ (T-a)--(T-d)
[I, Remark~1.2].  Since $T_{\mathrm{final}}$ does not depend on
$N$, [I, Lemma~1.5] applies to it: density of
$\NA((c_0,p_N),\mathcal F)$ for infinitely many $N$ implies density of
$\NA((c_0,p),\mathcal F)$ for Mart\'in's norm $p$ built with
$T_{\mathrm{final}}$.
\end{remark}

\subsection{Window arithmetic and design features}

\begin{proposition}[Sub-windows]\label{prop:II-subwindows}
For every sub-window $w=(l,i)$ of $T_{\mathrm{final}}$:
\begin{enumerate}
\item[(a)] $u(w)\le\frac14$, $b(w)\le2^{-64/u(w)}<u(w)$ and
$u(w^+)=b(w)$; hence $b$ is strictly decreasing along the sub-windows and
the bands $\mathrm{Band}(w)$ of all sub-windows (of all levels) are
pairwise disjoint.
\item[(b)] \textup{(P3) on sub-windows.}  $l\,2^{l^3}Q(w)T_{\rm hi}(w)\le1$,
$n(w)\ge l\,2^{l^3}Q(w)\ge l\,2^{l^3}\mathrm{Des}(l)\ge l\,2^{l^3}
\Lambda^\circ(l)$, $T_{\rm hi}(w)2^{l^3}\Lambda^\circ(l)\le1/l$, and
$T_{\rm hi}(w^+)\le T_{\rm lo}(w)$ \textup{(}also when $w^+$ is the first
sub-window of level $l+1$\textup{)}.
\item[(c)] \textup{(Fine weights.)}  $\sum_{l'>l}c_{l'}\le\frac43b(w)^2$
and $\sum_{l'>l}\lambda_{l'}\le\frac13b(w)^2\le T_{\rm lo}(w)^8$.
\item[(d)] With $\beta(w):=(1+\mathrm{Lip}(l)C_*(l))b(w)$:
$\beta(w)\le\eta(w)<1/\mathrm{Des}(l)\le1/C_L(l)$ and
$C_L(l)\beta(w)^{2/N_L(l)}\le\eta(w)$.
\item[(e)] $\mathrm{Des}(l)^kT_{\rm lo}(w)\le u(w)/8$ for every $k\le40$.
\item[(f)] $\log_2(1/T_{\rm hi}(w))\le l^3+l+Q(w)$,
$\log_2(1/T_{\rm lo}(w))\le2n(w)$ and
$\log_2(1/b(w))\le11\,N_L(l)\,n(w)\le11\,\mathrm{Des}(l)\,n(w)$.
\end{enumerate}
\end{proposition}

\begin{proof}
(a) $u((1,1))=\frac14$; otherwise $u(w)=b(w^-)\le\eta(w^-)\le
T_{\rm lo}(w^-)^4\le2^{-4}$, because $b\le\eta$ (the denominator of $b$
is $\ge1$ and $\min\{\eta,\cdot\}\le\eta$) and $T_{\rm lo}\le2^{-n}\le\frac12$.
Since $4/u(w)\ge16$, $n(w)\ge Q(w)\ge(4/u(w))^{20}\ge16/u(w)$, hence
$b(w)\le T_{\rm lo}(w)^4\le2^{-4n(w)}\le2^{-64/u(w)}$, and $2^{-64y}<1/y$
for $y=1/u(w)\ge4$.  $u(w^+)=b(w)$ is the definition.  So the endpoints
form one decreasing sequence $u(w_1)>u(w_2)=b(w_1)>u(w_3)=b(w_2)>\cdots$
along the lexicographic order, and $\mathrm{Band}(w_j)=(u(w_{j+1}),
u(w_j))$ are disjoint.

(b) The first two inequalities are the definitions of $T_{\rm hi}(w)$ and
$n(w)$.  $Q(w)\ge\mathrm{Des}(l)$ because $4/u(w)\ge1$ and
$\mathrm{Des}(l)\ge1$; $\mathrm{Des}(l)\ge\Lambda^\circ(l)$ because
$\Lambda^\circ(l)$ is a factor of the bracket and all factors are $\ge1$.
Then $T_{\rm hi}(w)2^{l^3}\Lambda^\circ(l)\le\Lambda^\circ(l)/(lQ(w))\le
1/l$.  $T_{\rm hi}(w^+)\le T_{\rm lo}(w)$ is part of the definition of
$T_{\rm hi}(w^+)$.

(c) By Theorem~\ref{thm:II-lemmaB}(d), the sums are bounded by
$\frac43c_{l+1}$ and $\frac13c_{l+1}$, and $c_{l+1}\le b(l,M(l))^2\le
b(w)^2$ by (W2) and (a).  Finally $b(w)^2\le T_{\rm lo}(w)^8$.

(d) $\beta(w)=\min\{\eta,(\eta/C_L)^{N_L/2}\}\le\eta$, and
$\beta^{2/N_L}\le\eta/C_L$.  $\eta(w)=T_{\rm lo}(w)^4/(l\mathrm{Des}(l))<
1/\mathrm{Des}(l)$ since $T_{\rm lo}(w)<1$, and $C_L(l)\le\mathrm{Des}(l)$
since $C_L(l)$ is one of the factors of $\mathrm{Des}(l)$, all of which are
$\ge1$.

(e) $Q(w)\ge16$, so
\[
 T_{\rm lo}(w)\le2^{-n(w)}\le2^{-Q(w)}\le Q(w)^{-2}=
 \Big(\frac{u(w)}{4\mathrm{Des}(l)}\Big)^{2\omega(l)+40}\le
 \Big(\frac{u(w)}{4\mathrm{Des}(l)}\Big)^{40};
\]
hence $\mathrm{Des}(l)^kT_{\rm lo}(w)\le4^{-40}u(w)^{40}\le u(w)/8$ for
$k\le40$.

(f) $\log_2(1/T_{\rm hi}(w))=\max\{\log_2(1/T_{\rm lo}(w^-)),l^3+\log_2
(lQ(w))\}$.  For $w\ne(1,1)$, $u(w)=b(w^-)\le T_{\rm lo}(w^-)^4$, so
$\log_2(1/T_{\rm lo}(w^-))\le\frac14\log_2(1/u(w))\le\log_2Q(w)$; for
$w=(1,1)$ this term is $0$.  Hence $\log_2(1/T_{\rm hi}(w))\le
l^3+\log_2l+\log_2Q(w)\le l^3+l+Q(w)$.  Since $Q(w)\ge\mathrm{Des}(l)\ge
2^{6l^3}\ge l^3+l$ and $l\,2^{l^3}\ge2$, we get $n(w)\ge2Q(w)\ge Q(w)+l^3+l$,
so $\log_2(1/T_{\rm lo}(w))=n(w)+\log_2(1/T_{\rm hi}(w))\le2n(w)$.  Next
$\log_2(1/\eta(w))=4\log_2(1/T_{\rm lo}(w))+\log_2(l\mathrm{Des}(l))\le
9n(w)$.  With $x:=\log_2(1/\eta(w))$ and $c:=\log_2C_L(l)\le\mathrm{Des}(l)
\le n(w)$,
\begin{align*}
 \log_2\frac1{b(w)}&=\log_2(1+\mathrm{Lip}(l)C_*(l))+\max\Big\{x,
 \frac{N_L(l)}2(x+c)\Big\}\\
 &\le n(w)+N_L(l)(9n(w)+n(w)),
\end{align*}
because $\log_2(1+\mathrm{Lip}(l)C_*(l))\le\mathrm{Des}(l)\le n(w)$.  So
$\log_2(1/b(w))\le11N_L(l)n(w)$, and $N_L(l)\le\mathrm{Des}(l)$.
\end{proof}

\begin{proposition}[Design features]\label{prop:II-features}
$T_{\mathrm{final}}$ and $U_{\mathrm{final}}$ have the following
properties.
\begin{enumerate}
\item[(a)] \textup{(Bounded gaps, rule (c).)}  Every $S_l$ has gaps
$G_l=2^{l+1}$; every scheduled target $y_l$ with $k(l)\notin K_{\rm FD}$
satisfies $\supp y_l\cap S_{l'}\subseteq[1,l]$ for $l'<l$; every
$y\in c_{00}$ is allowed at all large $l$.
\item[(b)] $c_{l+1}\le\min\{c_l/4,\ b(l,M(l))^2,\
(\delta_{l+1}H_{l+1})^2\}$ for every $l$.
\item[(c)] \textup{(SF$_\iota$), (SF$^*$).}  For every $l$,
\[
 (1+\|U\|)\sum_{l''>l}\lambda_{l''}\le2^{-10}\iota_l\lambda_l\delta_l
 2^{-(m(l)+k(l)+\min S_l)}\underline y_l/n_l,
\]
with $(\iota_l)$ decreasing to $0$; and
$c_l\le c_{l-1}\delta_l2^{-\min S_l-l-10}/(1+\|U\|)$ for $l\ge2$.
\item[(d)] \textup{(b$'$).}  $2c_l\le\iota_{l'}2^{-2s}c_{l'}\delta_{l'}$ for
$l'<l$ and $s\in\supp y_l\cap S_{l'}$.
\item[(e)] \textup{(Z0).}  $3,5,7,9$ are distinct elements of $Z_0$, and
$y^\ast$, $y^{\ast\ast}$ belong to the target family.
\item[(f)] \textup{(GM) for $l\ge2$.}  $|\sum_{j\in\supp y_l}\sigma_j
y_l(j)+\epsilon\delta_lH_l|\ge g_l>0$ for all $\sigma\in\{-1,0,1\}^{\supp
y_l}$, $\epsilon=\pm1$, and $\Phi_l\le g_l2^{-l}$.
\item[(g)] \textup{(FD).}  For every block $m$ and every $r\ge1$ the
carrier $l=\mathfrak j(2^r,m)$ has $y_l=e_{j_l}^*/q^*(e_{j_l}^*)$ with
$j_l$ odd, $j_l\notin\bigcup_iS_i$ and $j_l\notin\supp y_{l'}$ for
$l'<l$; hence $u_{l'}(j_l)=0$ for $l'<l$, $u_l(j_l)\ne0$, and
\[
 |y_l(j_l)|>(1+\|U\|)\big(\|y_l\|_1-|y_l(j_l)|+\delta_lH_l\big).
\]
\item[(h)] The base is diagonal with positive, strictly decreasing
entries, and $\|U\|=\frac1{16}\le\frac12$.
\end{enumerate}
\end{proposition}

\begin{proof}
(a) Definition~\ref{def:II-data}(i), allowedness (c), and the proof of
(T-d).  (b) (W1), (W2), (W3) at stage $l+1$.  (c) By (W7) at stage $l+1$
and (W1) at the later stages, $\sum_{l''>l}\lambda_{l''}\le\frac14\sum_{l''
\ge l+1}c_{l''}\le\frac13c_{l+1}\le2^{-10}\iota_l\lambda_l\delta_l
2^{-(m(l)+k(l)+\min S_l)}\underline y_l/(n_l(1+\|U\|))$.  The second
inequality is (W5).  (d) is (W4).  (e) Definition~\ref{def:II-data}(i),
(iv).  (f) The first statement is $\delta_l\in A_l$; $g_l>0$ trivially;
$\Phi_l=2^{-m(l)-k(l)}c_l\le g_l2^{-l}$ by (W6).  (g) $k(l)=2^r\in K_{\rm
FD}$, so step (1) chooses the fresh-dominated target; $j_l$ is odd, hence
in $Z_0$, and exceeds every element of $\supp y_{l'}$, $l'<l$.  For
$l'<l$, $u_{l'}(j_l)=(y_{l'}(j_l)+\delta_{l'}h_{l'}(j_l))/n_{l'}=0$, and
$u_l(j_l)=y_l(j_l)/n_l\ne0$.  Finally $|y_l(j_l)|=1/(1+\mu^\star_{j_l})\ge
\frac{16}{17}$, $\|y_l\|_1-|y_l(j_l)|=0$ and
$(1+\|U\|)\delta_lH_l\le(1+\|U\|)\delta^{\max}_lH_l\le\frac14$.  (h)
Lemma~\ref{lem:II-base}.
\end{proof}

\begin{proposition}[Bank efficiency and depths]\label{prop:II-bank}
For every level $l$:
\begin{enumerate}
\item[(a)] $s_{\max}(l)<s_{\rm b}(l)\le\max\{3\cdot2^l,\ s_{\max}(l)+
2^{l+1}\}$.
\item[(b)] For every $l''\le l$ the first three elements of
$S_{l''}\cap(s_{\max}(l),\infty)$ lie in $(s_{\max}(l),s_{\rm t}(l)]$.
\item[(c)] For every $c\le l$ and every $s'\in S_c$ with $s'\le
s_{\rm t}(l)$:
\[
 \begin{gathered}
 v_c(s')\ge\tfrac45\delta_{\min}(l)2^{-s_{\rm t}(l)},\qquad
 (\mu^\star_{s'})^2v_c(s')\ge\frac4{5B^\star(l)}\ge\mathrm{Des}(l)^{-1/6},\\
 (\mu^\star_{s'}v_c(s'))^{-2}\le\mathrm{Des}(l)^{1/3}.
 \end{gathered}
\]
\item[(d)] $2^{s_{\rm t}(l)+2}/\delta_{\min}(l)\le\mathrm{Des}(l)^{1/6}$ and
$8^{s_{\rm b}(l)}\le\mathrm{Des}(l)^{1/6}$.
\item[(e)] $\mathrm{Des}(l)^{1/6}$ dominates each factor inside the
bracket of $\mathrm{Des}(l)$, and $\mathrm{Des}(l)$ dominates each of
$\mathrm{Lip}(l)$, $1/\delta_{\rm comb}(l)$, $(r(l)c'(l))^l$, $C_*(l)$,
$C_L(l)$, $N_L(l)$, $1/\delta_{\rm sh}(l)$; moreover
$\mathrm{Des}(l)\ge2^{6l^3}$.
\end{enumerate}
\end{proposition}

\begin{proof}
(a) Each $\min(S_{l''}\cap(s_{\max}(l),\infty))$ exceeds $s_{\max}(l)$.  If
$s_{\max}(l)<\min S_{l''}=3\cdot2^{l''}$, this minimum is $3\cdot2^{l''}\le
3\cdot2^l$; otherwise it is at most $s_{\max}(l)+G_{l''}\le s_{\max}(l)+
2^{l+1}$ (bounded gaps).  (The cruder bound $s_{\rm b}(l)\le s_{\max}(l)+
2^{l+1}$ fails whenever $s_{\max}(l)+2^{l+1}<3\cdot2^{l}$, i.e.\ whenever
$s_{\max}(l)<2^l$.)  (b) The first
element is $\le s_{\rm b}(l)$, and the next two are at most $2G_{l''}\le
2^{l+2}$ further.  (c) $n_c\le\frac54$ and $\delta_c\ge\delta_{\min}(l)$
give $v_c(s')=\delta_c2^{-s'}/n_c\ge\frac45\delta_{\min}(l)2^{-s_{\rm t}(l)}$.
Since $\mu^\star$ is decreasing, $(\mu^\star_{s'})^2\ge(\mu^\star_{s_{\rm
t}(l)})^2$, so $(\mu^\star_{s'})^2v_c(s')\ge\frac45/B^\star(l)$, and
$(\mu^\star_{s'}v_c(s'))^{-2}\le\frac{25}{16}(\mu^\star_{s_{\rm t}(l)})^{-2}
4^{s_{\rm t}(l)}\delta_{\min}(l)^{-2}\le\frac{25}{16}B^\star(l)^2$.  Since
$l\,2^{l^3}\ge2$ and $2^{G(l)}\ge4$, $\mathrm{Des}(l)^{1/6}\ge8B^\star(l)
\ge\frac54B^\star(l)$, which gives both remaining inequalities.  (d)
$B^\star(l)\ge2^{s_{\rm t}(l)}/\delta_{\min}(l)$ and
$B^\star(l)\ge(\mu^\star_{s_{\rm b}(l)})^{-2}=2^{2\cdot2^{2^{s_{\rm b}(l)}}}
\ge8^{s_{\rm b}(l)}$, and $\mathrm{Des}(l)^{1/6}\ge8B^\star(l)$.  (e) All
factors are $\ge1$ by (K1) and Lemma~\ref{lem:II-W}(b);
$\mathrm{Des}(l)\ge(l2^{l^3})^6$.
\end{proof}

\begin{remark}[The role of $B^\star(l)$]\label{rem:II-baseindep}
In the diagonal-base arguments of the later sections (private banks,
far pulls, donor raises, exact two-sided tuning, pair and anchor tuning
on support coordinates) the base enters a design constant only through
the bank efficiency $(\mu^\star_{s'})^2v_c(s')$ at coordinates
$s'\le s_{\rm t}(l)$ and through $\|U\|\le\frac12$; the remaining uses of
the base are the identities of Lemma~\ref{lem:II-base}(b) and constants
depending on $f$.  Proposition~\ref{prop:II-bank}(c) turns the bank
efficiency into a power of $\mathrm{Des}(l)$, for instance the donor bank
mass is at most $\mathrm{Des}(l)\Lambda$ when the required threshold raise
is $\Lambda$, and the contraction constant of the tuning fixed point,
which contains $(\mu^\star_{s'}v_c(s'))^{-2}$, is at most
$\mathrm{Des}(l)^{1/3}$ times $f$-dependent constants.  This is the sense
in which the factor $B^\star(l)$ makes the bank and pull machinery
independent of the choice of the diagonal base; each such use is checked
where it occurs.  The depth $s_{\rm t}(l)$ (rather than $s_{\rm b}(l)$)
is needed because pair tuning on support coordinates uses the first three
coordinates of each $S_{l''}$ beyond $s_{\max}(l)$
(Proposition~\ref{prop:II-bank}(b)), and because deep support raises must
be taken beyond $s_{\rm t}(l)$; the pinning constant at that depth,
$2^{s_{\rm t}(l)+2}/\delta_{\min}(l)$, is a design factor by
Proposition~\ref{prop:II-bank}(d).
\end{remark}

\subsection{Clean sub-windows}

\begin{definition}[Clean sub-windows]\label{def:II-clean}
Let $N\ge1$ and $f\in S_{p_N^*}$.  A sub-window $w$ of level $l$ is
\emph{clean for $f$} if no object of $\mathcal O(l)$ has its value at $f$
in $\mathrm{Band}(w)=(b(w),u(w))$.  At a clean sub-window a value is
called \emph{tiny} if it is $\le b(w)$ and \emph{robust} if it is
$\ge u(w)$.
\end{definition}

\begin{theorem}[Clean sub-windows]\label{thm:II-clean}
Let $N\ge1$ and $f\in S_{p_N^*}$, with $F=\supp a$ finite or infinite.
For every level $l$ there is $i\in\{1,\dots,M(l)\}$ such that the
sub-window $(l,i)$ is clean for $f$.  At such a sub-window the value at
$f$ of every rate object of level $l$ is tiny or robust.
\end{theorem}

\begin{proof}
The $M(l)=\omega(l)+1$ bands $\mathrm{Band}(l,i)$ are pairwise disjoint
(Proposition~\ref{prop:II-subwindows}(a)), so each of the $\omega(l)$
values lies in at most one of them, and at least one band contains none of
them.  The second statement is the definition, since every value lies in
$[0,\infty]$ and $[0,\infty]\setminus(b(w),u(w))=[0,b(w)]\cup
[u(w),\infty]$.
\end{proof}

\begin{corollary}[Consequences of cleanness]\label{cor:II-clean}
Let $w$ be a clean sub-window of level $l$ for $f\in S_{p_N^*}$, and
write $b=b(w)$, $u=u(w)$.
\begin{enumerate}
\item[(a)] Every maximum or minimum of finitely many values of objects of
$\mathcal O(l)$ is tiny or robust.
\item[(b)] For every carrier $l''\le l$ with $m(l'')\le N$,
$\varrho_{l''}(f)\in[0,b]\cup[u,1-u]\cup[1-b,1+b]\cup[1+u,\infty)$.
\item[(c)] $b\le2^{-64/u}$; in particular $u/b\ge2^{64}$, and
$\mathrm{Des}(l)^kT_{\rm lo}(w)\le u/8$ for $k\le40$.
\end{enumerate}
\end{corollary}

\begin{proof}
(a) The set $[0,b]\cup[u,\infty]$ is closed under finite maxima and
minima.  (b) By (O3), (O4), $\varrho:=\varrho_{l''}(f)\notin(b,u)$ and
$|\varrho-1|\notin(b,u)$.  If $\varrho\ge u$, then $|\varrho-1|\le b$ or
$|\varrho-1|\ge u$, i.e.\ $\varrho\in[1-b,1+b]$, $\varrho\le1-u$ or $\varrho\ge1+u$.
(c) Proposition~\ref{prop:II-subwindows}(a), (e).
\end{proof}

\begin{remark}[Quantifiers]\label{rem:II-quantifiers}
The order of choices is: the fixed data and $T_{\mathrm{final}}$ (no
dependence on $N$); then $N$ and $f\in S_{p_N^*}$; then, for each level
$l$, a clean sub-window $w_l$ (which depends on $f$); then the
exactifying companion $f^\#_{w_l}$ (Definition~\ref{def:II-companionw}; it
depends on $f$ and $w_l$ only); then the mate $g$ and $\rho<1$; then the
data at each scale $t\in\mathcal W(w_l)$.  The rate
objects are fixed before $f$; only their values depend on $f$.  An
$f$-dependent quantity used in a window argument must therefore be an
$f$-constant (absorbed by Lemma~\ref{lem:II-GW}(a)(GW2)) or the value of
a rate object of the level (tiny or robust at $w_l$).
\end{remark}

\subsection{Generic window use}

\begin{lemma}[Generic window use]\label{lem:II-GW}
\textup{(a)} For every sub-window $w$ of level $l$:
\begin{enumerate}
\item[(GW1)] $\mathcal W(w)$ contains the $n(w)$ dyadic scales
$T_{\rm hi}(w)2^{1-i}$, $1\le i\le n(w)$, and
$T_{\rm hi}(w)2^{l^3}\Lambda^\circ(l)\le1/l$,
$n(w)\ge l\,2^{l^3}\Lambda^\circ(l)$.
\item[(GW2)] Let $C,K\ge1$, and for every $l$ let $w_l$ be a sub-window of
level $l$ and $R_l\ge1$ with $R_l\le K\,l\,2^{l^3}\Lambda^\circ(l)$.  Put
$X_l:=C^{l^2}(C/u(w_l))^{\omega(l)+8}\mathrm{Des}(l)^8R_l$.  Then
$T_{\rm hi}(w_l)X_l\to0$ and $n(w_l)/X_l\to\infty$ as $l\to\infty$.
\item[(GW3)] $\sum_{l'>l}\lambda_{l'}\le T_{\rm lo}(w)^8$; consequently, for
every two-sided decomposition at a scale $t\in\mathcal W(w)$ (in the
sense of [I, Section~8.3]), $\sum_{l'>l}|\Delta\theta_{l'}|\le
6t^7\le6t^2$.
\item[(GW4)] $T_{\rm hi}(w')\le T_{\rm lo}(w)$ for every sub-window $w'$
following $w$ in the lexicographic order.
\end{enumerate}
\textup{(b)} Let $(w_l)_{l\ge1}$ be any choice of one sub-window $w_l$ of
each level $l$, and use $\mathcal W(w_l)$, $T_{\rm hi}(w_l)$,
$T_{\rm lo}(w_l)$, $n(w_l)$ in place of the $l$-th window $\mathcal W(l)$,
$T_{\rm hi}(l)$, $T_{\rm lo}(l)$, $n^w_l$ of [I, Definition~8.1].
Then $T_{\mathrm{final}}$ satisfies (T-a)--(T-d), (P1), (P2) and (P3) of
[I, Theorem~8.2] verbatim; consequently every result of
[I, Section~8] holds for $T_{\mathrm{final}}$, every $N$, and this
choice of windows.
\end{lemma}

\begin{proof}
(GW1) Definition~\ref{def:II-Tfinal}(6) and
Proposition~\ref{prop:II-subwindows}(b).

(GW2) Since $\mathrm{Des}(l)\ge2^{6l^3}\to\infty$, for all large $l$ we
have $C\le4\mathrm{Des}(l)$, so $(C/u)^{\omega+8}\le(4\mathrm{Des}/u)^{\omega
+8}=Q(w_l)(u/(4\mathrm{Des}))^{12}$ (we write $u=u(w_l)$,
$\mathrm{Des}=\mathrm{Des}(l)$, $\omega=\omega(l)$).  Using $T_{\rm hi}
(w_l)\le2^{-l^3}/(lQ(w_l))$, $u\le1$ and $\mathrm{Des}\ge(l2^{l^3}
\Lambda^\circ(l))^6\ge2^{6l^3}\Lambda^\circ(l)$,
\[
 T_{\rm hi}(w_l)X_l\le\frac{C^{l^2}\mathrm{Des}^8R_l}{l\,2^{l^3}
 \mathrm{Des}^{12}}\le\frac{K\,C^{l^2}\Lambda^\circ(l)}{\mathrm{Des}^4}\le
 K\,C^{l^2}2^{-24l^3}\to0 .
\]
Similarly, using $n(w_l)\ge l2^{l^3}Q(w_l)$ and $(4\mathrm{Des}/u)^{12}\ge
\mathrm{Des}^{12}$,
\[
 \frac{n(w_l)}{X_l}\ge\frac{l\,2^{l^3}\mathrm{Des}^{12}}{C^{l^2}
 \mathrm{Des}^8R_l}\ge\frac{\mathrm{Des}^4}{K\,C^{l^2}\Lambda^\circ(l)}\ge
 \frac{2^{24l^3}}{K\,C^{l^2}}\to\infty .
\]

(GW3) The first inequality is Proposition~\ref{prop:II-subwindows}(c).
By the box bound [I, Lemma~8.10], $|\theta^\pm_{l'}|\le3\lambda_{l'}/t$,
hence $|\Delta\theta_{l'}|\le6\lambda_{l'}/t$, so $\sum_{l'>l}|\Delta\theta_{l'}|\le6T_{\rm lo}(w)^8/t
\le6t^7$ for $t\ge T_{\rm lo}(w)$, and $t\le1$.

(GW4) $T_{\rm hi}(w^+)\le T_{\rm lo}(w)$ by the definition of
$T_{\rm hi}(w^+)$, and $T_{\rm lo}(w^+)\le T_{\rm hi}(w^+)$; hence
$T_{\rm hi}$ and $T_{\rm lo}$ are nonincreasing along the order, and
induction gives the claim.

(b) (T-a)--(T-d), (P1) and the first part of (P2) are
Theorem~\ref{thm:II-lemmaB}.  The window part of (P2):
$\sum_{l'>l}c_{l'}\le\frac43c_{l+1}\le2c_{l+1}$ and $c_{l+1}\le
b(w_l)^2\le T_{\rm lo}(w_l)^8\le T_{\rm lo}(w_l)^3$, and
$\sum_{l'>l}\lambda_{l'}\le T_{\rm lo}(w_l)^8\le T_{\rm lo}(w_l)^3$
(Proposition~\ref{prop:II-subwindows}(c); note that $c_{l+1}\le
b(l,M(l))^2\le b(w_l)^2$).  (P3): $T_{\rm hi}(w_l)2^{l^3}\Lambda^\circ(l)
\le1/l$ and $n(w_l)\ge l2^{l^3}\Lambda^\circ(l)$ by (GW1), and
$T_{\rm hi}(w_{l+1})\le T_{\rm lo}(w_l)$ by (GW4).  The targets are
allowed in the sense of [I, Definition~8.1(D1)(a),(b)]
(Theorem~\ref{thm:II-lemmaB}(c)), the vectors $u_l$ have the form
$(y_l+\delta_lh_l)/n_l$ with $q^*(\delta_lh_l)\le\frac14$ and
$\delta_l\to0$, and $Te_{k,m}=c_{\mathfrak j(k,m)}u_{\mathfrak j(k,m)}$.  As
stated after [I, Theorem~8.2], only (T-a)--(T-d), (P1), (P2) and
(P3) are used in [I, Section~8]; the value of $\delta_l$ enters
only through $q^*(\delta_lh_l)\le\frac14$ and $\delta_l\to0$ (both hold
here although $\delta_l$ may differ from $\delta^{\max}_l$ by a factor at
most $2$).  Hence those results hold for $T_{\mathrm{final}}$ with the
windows $\mathcal W(w_l)$.
\end{proof}

\begin{remark}[How the window properties are used]\label{rem:II-GWuse}
In every window argument of [I, Section~8] and of the later
sections of this part, the window enters only through (i) the dyadic
scales of one window, (ii) conditions ``$T_{\rm hi}\,X_l\to0$'' and
``$n/X_l\to\infty$'' along the levels used, where $X_l$ is an
$f$-constant to the power $l^2$, times a product of at most $48$ design
factors of level $l$ (each $\le\mathrm{Des}(l)^{1/6}$, so the product is
$\le\mathrm{Des}(l)^8$), times at most $\omega(l)+8$ factors
$C_f/u(w)$ coming from robust rates, times a room product compared with
$l2^{l^3}\Lambda^\circ(l)$; (iii) the box bound (GW3); and (iv) the
ordering (GW4).  Lemma~\ref{lem:II-GW}(a) provides all four on every
sub-window, so these arguments may use any sub-window of level $l$, in
particular a clean one, as ``the window of level $l$''.  All four
conditions are monotone in $\mathrm{Des}(l)$: enlarging the design factor
only lengthens and deepens the sub-windows.  This is a description of the
proofs, verified for each window theorem where it is proved; it is the
reason why the factors $N_L(l)$ (Proposition~\ref{prop:II-subwindows}(f))
and $B^\star(l)$ could be added to the design factor without affecting
any estimate.
\end{remark}

\begin{proposition}[Window factor for the infinite-support arguments]
\label{prop:II-RSwindow}
For a sub-window $w$ of level $l$ put
$\Omega_K(l):=s_{\max}(l)2^{s_{\max}(l)}/\delta_{\min}(l)$ and
\[
 \Omega(w):=\Omega_K(l)\Big(2+\log_2\big(2+\log_2(1/T_{\rm lo}(w^-))\big)
 \Big)\Big(2+\log_2\big(l2^{l^3}\Lambda^\circ(l)\Omega_K(l)\big)\Big)
\]
\textup{(}$T_{\rm lo}((1,1)^-)=1$\textup{)}.  Then
\begin{enumerate}
\item[(a)] $\Lambda^\circ(l)\Omega(w)\le Q(w)$; hence
$n(w)\ge l2^{l^3}\Lambda^\circ(l)\Omega(w)$ and
$T_{\rm hi}(w)\le2^{-l^3}/(l\,\Omega(w)\Lambda^\circ(l))$;
\item[(b)] $\big(\log_2(1/T_{\rm hi}(w))+\log_2n(w)\big)/n(w)\le
(2l^3+2l+4)/(l2^{l^3})$;
\item[(c)] for every $x\in[T_{\rm lo}(w)^2,\frac12]$,
$2^{J(x)}\le2\log_2(4\log_2(1/T_{\rm lo}(w)))\le2\log_2(8n(w))$, with
$J(\cdot)$ as in Lemma~\ref{lem:II-base}(d).
\end{enumerate}
\end{proposition}

\begin{proof}
(a) Write $\mathrm{Des}=\mathrm{Des}(l)$, $u=u(w)$.  Since $s\le2^s$ and
$\delta_{\min}(l)\le1$, $\Omega_K(l)\le(2^{s_{\max}(l)}/\delta_{\min}(l))^2$,
so $l2^{l^3}\Lambda^\circ(l)\Omega_K(l)\le(l2^{l^3}\Lambda^\circ(l)
2^{s_{\max}(l)}\delta_{\min}(l)^{-1})^2\le\mathrm{Des}^{1/3}$.  The last
factor of $\Omega(w)$ is therefore at most $2+\log_2\mathrm{Des}\le
\mathrm{Des}^{1/2}$, because $\mathrm{Des}\ge64$ and
$y\mapsto y^{1/2}-\log_2y-2$ vanishes at $64$ and increases on
$[64,\infty)$.  For the middle factor: if $w=(1,1)$ it equals $3\le4/u$;
otherwise $T_{\rm lo}(w^-)^4\ge b(w^-)=u$, so with $x:=\log_2(1/u)\ge2$
the middle factor is at most $2+\log_2(2+x/4)\le3+x/4\le2+x\le2^{x+2}=4/u$
(we used $\log_2(2+y)\le1+y$ for $y\ge0$).
Hence $\Lambda^\circ(l)\Omega(w)\le\mathrm{Des}^{1/3}(4/u)
\mathrm{Des}^{1/2}\le4\mathrm{Des}/u\le Q(w)$, and the two consequences
follow from the definitions of $n(w)$ and $T_{\rm hi}(w)$.

(b) By Proposition~\ref{prop:II-subwindows}(f),
$\log_2(1/T_{\rm hi}(w))\le l^3+l+Q(w)$; also $\log_2n(w)\le
\log_2(l2^{l^3}Q(w)+1)\le l^3+l+Q(w)+1$.  Divide by $n(w)\ge l2^{l^3}Q(w)$
and use $Q(w)\ge1$.

(c) $J$ is nonincreasing in $x$, so $2^{J(x)}\le2^{J(T_{\rm lo}(w)^2)}\le
2\log_2(2\log_2(T_{\rm lo}(w)^{-2}))$ by Lemma~\ref{lem:II-base}(d), and
$\log_2(1/T_{\rm lo}(w))\le2n(w)$ by
Proposition~\ref{prop:II-subwindows}(f).
\end{proof}

\begin{remark}[Scope of $T_{\mathrm{final}}$]\label{rem:II-scope}
(a) $T_{\mathrm{final}}$ contains the features of the earlier
signature-ladder designs with sub-windows: the sub-window recursion and
the rate objects (O1)--(O4); the bounded gaps and allowedness rule (c);
the diagonal base with private bank and pull coordinates (every
coordinate of every $S_l$ beyond $s_{\max}(l)$ is available); the
design factors of the configuration, tuning and determinantal arguments;
the {\L}ojasiewicz-adjusted value of $b(w)$ (Proposition
\ref{prop:II-subwindows}(d)); the rate objects (O5)--(O9); the weight
conditions (SF$_\iota$), (b$'$), (GM) for $l\ge2$, (Z0), (FD)
(Proposition~\ref{prop:II-features}); and, for infinite supports, the
doubly exponential base, the depth $s_{\rm t}(l)$ with the factor
$B^\star(l)$, the window factor $\Omega(w)$ (Proposition
\ref{prop:II-RSwindow}) and the rate objects (O10)--(O12).

(b) Two features of earlier designs are deliberately \emph{absent}: an
explosive window function (replaced by the pigeonhole sub-windows) and
lacunary signature sets (incompatible with bounded gaps).  Results that
depend on these features are not claimed for $T_{\mathrm{final}}$; no
result of this part uses them.

(c) The fresh-dominated targets serve only a negative statement (dead
zones forced by positive mismatch) and, below, the fact that
$s_{\max}(l)\to\infty$ (Lemma~\ref{lem:II-scales}(e)) and the stages of the
modified design $D^{U1'}$ (Definition~\ref{def:II-DU1}).  We do not claim
that they can be removed from the schedule.
\end{remark}

%%%%%%%%%%%%%%%%%%%%%%%%%%%%%%%%%%%%%%%%%%%%%%%%%%%%%%%%%%%%%%%%%%%%%%%%

\section{Clean sub-windows and the exactifying companion}\label{sec:II-exact}

In this section we prove the second master theorem: for the operator
$T_{\mathrm{final}}$ and a first row $f$ with finite base support, recovery
of every mate of $f$ follows from two structural conditions at one clean
sub-window of infinitely many levels (Theorem~\ref{thm:II-master2}).  The
proof combines four ingredients.  First, at a clean sub-window the
switching amplitudes of a two-sided decomposition are pinned
\emph{diagonally}: no product of rooms and no slaving between carriers
occurs, because the pinning sets exclude all coarse target coordinates
(Subsection~\ref{subsec:II-pin}).  Second, with the diagonal base
$U_{\mathrm{final}}$, masses placed far out on signature sets act on the
values of single carriers: a \emph{far pull} lowers one value and a
\emph{private bank} raises one value at first order, while all other coarse
values move only at second order; together they tune finitely many values
exactly and two-sidedly (Subsection~\ref{subsec:II-tools}).  Third, every
block has, at every clean sub-window of large level, a peak with robust
margin, and a bank at a far signature coordinate of that peak raises the
block threshold by a prescribed amount; this \emph{threshold buffer}
dominates every drift caused by the other moves
(Subsection~\ref{subsec:II-donors}).  Fourth, one companion per clean
sub-window carries all exactifications at once, at cost $o(T_{\rm lo}^2)$,
and exact $d$-neutral two-piece data are transplanted to it
(Subsections~\ref{subsec:II-companion} and~\ref{subsec:II-transplant}).
Theorem~\ref{thm:II-E} then transfers recovery from the companions to $f$.

\subsection*{Standing assumptions and notation}
Throughout this section $T=T_{\mathrm{final}}$, $U=U_{\mathrm{final}}$
(Section~\ref{sec:II-Tfinal}), $N\ge1$, $I=\{1,\dots,N\}$, $p=p_N$, and
$f\in S_{p^*}$ has forced data $(\xi,q_0,a,w,F,z,\zh,e,\nu)$
[I,~Proposition~1.9] with $F=\supp a$ \emph{finite}.  We use the notation
of Sections~\ref{sec:II-companions} and~\ref{sec:II-Tfinal}.  In
particular, for a block $m$,
\[
 \hat\zeta_m:=R_m^{**}\zh,\qquad A_m:=|\hat\zeta_m|_m=\sigma_m/q_0,\qquad
 \bar\vartheta_m:=\bar\vartheta(\hat\zeta_m)=A_mM_m/C_m ,
\]
so that $\hat\zeta_m(k)=\lambda_{k,m}u_{k,m}(\zh)$.  The \emph{coarse
carriers} of level $l$ are the elements of $\mathcal C(l):=\{l''\le l:
m(l'')\le N\}$; for a carrier $l''=\mathfrak j(k,m)$ we write $k(l'')=k$,
$m(l'')=m$ and use $u_{l''}=u_{k,m}$, $\Phi_{l''}=\Phi_m(k)$,
$\lambda_{l''}=m\Phi_{l''}$ (Theorem~\ref{thm:II-lemmaB}(c)).  Carriers
$l''>l$ are \emph{fine}.  The relative position
$\varrho_{l''}=\varrho_{l''}(f)$ of Definition~\ref{def:II-rates}(O4) equals
$\varrho_{k}(\hat\zeta_m)=m|u_{l''}(\zh)|/(\Phi_{l''}\bar\vartheta_m)$
(Subsection~\ref{subsec:II-threshold}); hence $k(l'')$ is a peak iff
$\varrho_{l''}\ge1$, the margin of a peak is
$\mu_{k,m}=q_0\Phi_{l''}\bar\vartheta_m(\varrho_{l''}-1)/m$, the gap of a strict
non-peak is $\gap_m(k)=M_m(1-\varrho_{l''})$, and $|w_m(k)|=M_m\min(1,
\varrho_{l''})$ (Lemma~\ref{lem:II-T}(b)).  Recall $|u_{l''}(\zh)|\le1$.
Given a sign $\epsilon_{l''}$ we write $\mathrm{val}_{l''}:=\epsilon_{l''}
u_{l''}(\zh)$, and $\mathrm{val}'_{l''}:=\epsilon_{l''}u_{l''}(\zh')$ at
another first row $f'$.

For a sub-window $w$ of level $l$ we abbreviate
\begin{gather*}
 b_w:=b(w),\quad u_w:=u(w),\quad \Lambda_w:=T_{\rm lo}(w)^3,\quad
 \beta_w:=b_w+\mathrm{Des}(l)^2\Lambda_w^2,\\
 D_{\rm r}:=D_{\rm room}(l),\quad \mathrm{Des}:=\mathrm{Des}(l).
\end{gather*}
For $l''\le l$ put $S^{\rm nat}_{l''}(l):=S_{l''}\setminus(F\cup T(l))$,
$\mathfrak m_{l''}:=\|v_{l''}1_{S^{\rm nat}_{l''}(l)}\|_1$ and
\[
 r^{\rm nat}_{l''}:=\min_{\sigma=\pm1}\sum_{s\in S^{\rm nat}_{l''}(l)}
 v_{l''}(s)(1+\sigma z_s),
\]
so that the value of the room object (O1) of $l''$ is $r^{\rm nat}_{l''}/
\mathfrak m_{l''}$.  An \emph{$f$-constant} is a number depending only on
$f$ (and on $N$ and the fixed design), not on $l$, $w$ or the scale; $C_f$
and $c_f>0$ denote $f$-constants whose values may change from line to line.
$l_f$ denotes a level threshold depending only on $f$; it is enlarged
finitely many times below, each time by a condition that depends only on
$f$.  Statements ``for $l\ge l_f$'' refer to sub-windows of level $l$.

\begin{lemma}[Scale relations]\label{lem:II-scales}
Let $w$ be a sub-window of level $l$ and $t\in\mathcal W(w)$.
\begin{enumerate}
\item[(a)] $b_w\le T_{\rm lo}(w)^4/(l\,\mathrm{Des})\le t^4/(l\,\mathrm{Des})$
and $\sum_{l'>l}\lambda_{l'}\le b_w^2/3$.
\item[(b)] $T_{\rm lo}(w)\le(u_w/(4\mathrm{Des}))^{40}$; hence
$\mathrm{Des}^kT_{\rm lo}(w)\le u_w/8$ for $0\le k\le40$, and
$\mathrm{Des}^k\beta_w\le u_w\Lambda_w/(4l)$ for $0\le k\le38$.
\item[(c)] For $l''\le l$: $\lambda_{l''}\ge\Phi_{l''}\ge1/D_{\rm r}$; if
$l\ge\max F$, then $\mathfrak m_{l''}\ge\mathfrak m^{\rm nat}_{l''}(l)\ge
1/D_{\rm r}$.  Each of $l$, $D_{\rm r}$, $2^{G(l)}$, $1+|T(l)|$,
$G^{**}(l)$, $H_{\rm tune}(l)$, $H_{\rm conf}(l)$ is at most
$\mathrm{Des}^{1/6}$.
\item[(d)] For $l''\le l$ let $s'_{l''}(l):=\min(S_{l''}\cap(s_{\max}(l),
\infty))$.  Then $s'_{l''}(l)\le s_{\rm b}(l)\le s_{\rm t}(l)$, and for every
$c\le l$ and $s'\in S_c$ with $s'\le s_{\rm t}(l)$:
$(\mu^\star_{s'})^2v_c(s')\ge\mathrm{Des}^{-1/6}$ and
$(\mu^\star_{s'}v_c(s'))^{-2}\le\mathrm{Des}^{1/3}$.
\item[(e)] $s_{\max}(l)\to\infty$, $u(w)\to0$ and $b(w)\to0$ as
$l\to\infty$ \textup{(}uniformly over the sub-windows of level $l$\textup{)}.
\end{enumerate}
\end{lemma}

\begin{proof}
(a) $b_w\le\eta(w)$ (Definition~\ref{def:II-Tfinal}(6)) and $t\ge
T_{\rm lo}(w)$; the second claim is Proposition~\ref{prop:II-subwindows}(c).
(b) $T_{\rm hi}(w)\le1$, $n(w)\ge Q(w)\ge16$ and $2^{-Q}\le Q^{-2}$ give
$T_{\rm lo}(w)\le2^{-n(w)}\le Q(w)^{-2}=(u_w/(4\mathrm{Des}))^{2\omega(l)+40}$.
Hence $\mathrm{Des}^kT_{\rm lo}\le4^{-40}u_w^{40}\le u_w/8$ for $k\le40$.
By (a), $\mathrm{Des}^kb_w\le(\Lambda_w/l)\mathrm{Des}^{k-1}T_{\rm lo}\le
u_w\Lambda_w/(8l)$ for $k\le41$, and $\mathrm{Des}^{k+2}\Lambda_w^2=
\Lambda_wT_{\rm lo}^2\,\mathrm{Des}^{k+2}T_{\rm lo}\le u_w\Lambda_w/(8l)$ for
$k\le38$, because $T_{\rm lo}\le T_{\rm hi}\le2^{-l^3}/l$ gives
$T_{\rm lo}^2\le1/l$.
(c) $\Phi_{l''}=2^{-m-k}c_{l''}\le\lambda_{l''}$, and $1/\Phi_{l''}$,
$1/\mathfrak m^{\rm nat}_{l''}(l)$ are summands of $D_{\rm room}(l)$.  If
$l\ge\max F$, then $S_{l''}\setminus([1,l]\cup T(l))\subseteq
S^{\rm nat}_{l''}(l)$.  The listed numbers are factors of the bracket in
$\mathrm{Des}(l)$ or are dominated by such factors ($G^{**}(l)\le
H_{\rm conf}(l)$ by Definition~\ref{def:II-pattern} below), and
$\mathrm{Des}(l)$ is at least the sixth power of the bracket.
(d) Proposition~\ref{prop:II-bank}(a)--(c).
(e) At the stages $l'=\mathfrak j(2^r,m')$ the targets are fresh-dominated
(Proposition~\ref{prop:II-features}(g)) with $j_{l'}$ larger than every
earlier target coordinate, and $s_{\max}$ is nondecreasing.  For $w$ of
level $l\ge2$, $u(w)=b(w^-)\le T_{\rm lo}(w^-)^4\le2^{-4(l-1)^3}$, and
$b(w)<u(w)$.
\end{proof}

In particular we may and do assume $l_f\ge\max F$ and
$s_{\max}(l_f)\ge\max F$, so that for $l\ge l_f$
\begin{equation}\label{eq:II-Finside}
 F\cup T(l)\subseteq[1,s_{\max}(l)],\qquad
 S_{l''}\setminus[1,s_{\max}(l)]\subseteq S^{\rm nat}_{l''}(l)
 \quad(l''\le l).
\end{equation}

\subsection{Classification and pinning at a clean sub-window}
\label{subsec:II-pin}

Let $w$ be a sub-window of level $l\ge l_f$ which is clean for $f$
(Definition~\ref{def:II-clean}); by Corollary~\ref{cor:II-clean}(b), every
$l''\in\mathcal C(l)$ has $\varrho_{l''}\in[0,b_w]\cup[u_w,1-u_w]\cup[1-b_w,
1+b_w]\cup[1+u_w,\infty)$, and every room object and every target room
object of level $l$ is tiny or robust.

\begin{definition}[Classification at a clean sub-window]\label{def:II-class}
Let $w$ be clean for $f$, of level $l\ge l_f$.
\begin{enumerate}
\item[(a)] A carrier $l''\in\mathcal C(l)$ is of \emph{class R} if
$r^{\rm nat}_{l''}\ge u_w\mathfrak m_{l''}$ and of \emph{class G} if
$r^{\rm nat}_{l''}\le b_w\mathfrak m_{l''}$.  For $l''$ of class G,
$\epsilon_{l''}\in\{\pm1\}$ is the sign with $r^{\rm nat}_{l''}=\sum_{s\in
S^{\rm nat}_{l''}(l)}v_{l''}(s)(1-\epsilon_{l''}z_s)$ (it is unique,
because the two signed rooms add up to $2\mathfrak m_{l''}>2b_w
\mathfrak m_{l''}$), and the \emph{$d$-weight} of $l''$ is
\[
 q_{l''}:=\frac{\epsilon_{l''}\Phi_{l''}w_m(k)}{mC_m}\qquad
 (l''=\mathfrak j(k,m)).
\]
$l''$ is of \emph{swallowing type} if $\epsilon_{l''}\sgn w_m(k)=1$, of
\emph{anti type} if $\epsilon_{l''}\sgn w_m(k)=-1$, and \emph{$d$-neutral}
if $w_m(k)=0$ (then $\varrho_{l''}=0$).  $G(w)$ is the set of class-G
carriers.
\item[(b)] For a block $m$, $\Sigma_m(w)$ is the set of $l''\in G(w)$ of
block $m$ that are swallowing-type peaks, or strict non-peaks with
$\varrho_{l''}>b_w$ (hence $\varrho_{l''}\ge u_w$) that are not of anti type with
$\varrho_{l''}\ge1-b_w$.  Block $m$ is \emph{$\sigma$-one-signed at $w$}
($\sigma\in\{\pm1\}$) if $\sgn q_{l''}=\sigma$ for every $l''\in\Sigma_m(w)$.
\item[(c)] The \emph{dropped set} $P(w)\subseteq G(w)$ consists of
(P-i) anti-type peaks; (P-ii) anti-type strict non-peaks with $\varrho\ge
1-b_w$; (P-iii) swallowing-type peaks with $\varrho\ge1+u_w$; (P-iv) the sets
$\Sigma_m(w)$ of the blocks $m$ that are one-signed at $w$.  The \emph{kept
set} is $\mathrm{Kp}(w):=G(w)\setminus P(w)$; it consists of (K1)
swallowing-type strict non-peaks with $\varrho\in[u_w,1-u_w]$; (K2) anti-type
strict non-peaks with $\varrho\in[u_w,1-u_w]$; (K3) carriers with $\varrho\le b_w$;
(K4) swallowing-type carriers with $\varrho\in[1-b_w,1+b_w]$ (strict non-peaks,
peaks with small margin, degenerate peaks); in a one-signed block only
(K3) carriers are kept.
\item[(d)] A target coordinate $j\in T(l)\setminus F$ is
\emph{contact-like} if $1-|z_j|\le b_w$ and \emph{free-robust} if
$1-|z_j|\ge u_w$.
\end{enumerate}
\end{definition}

Every carrier of $\mathcal C(l)$ is of class R or of class G, and every
target coordinate off $F$ is contact-like or free-robust, because $w$ is
clean.  In the remainder of this subsection we fix, in addition, $g\in
\Cset(f)$, $\eta\le\eta_*$ with $\eta_\Gamma(\eta)\le1$
[I,~Lemma~8.7], a scale $t\in\mathcal W(w)$ with $t\le\min(t_\eta,1)$
[I,~Lemma~8.6], and a two-sided decomposition $(B_\pm,\Theta_\pm)$ of $g$
at scale $t$ [I,~Lemma~8.5].  As in [I,~Section~8.3] and
[I,~(15)] we write $\Delta\theta_{l''}:=\lambda_{l''}(\Theta_{+}-
\Theta_-)_{m(l'')}(k(l''))$ (and $\Delta\theta_{l''}:=0$ if $m(l'')>N$),
$\Delta B:=B_+-B_-=-\sum_{l''}\Delta\theta_{l''}u_{l''}$, and, with
$d_{\pm,m}$ and $\omega_{\pm,m}$ as in [I,~Lemma~8.8],
$\Delta d_m:=d_{+,m}-d_{-,m}$.  For $l''\in G(w)$ the \emph{switching
amplitude} is $\tau_{l''}:=-\epsilon_{l''}\Delta\theta_{l''}$.  By
[I,~Lemma~8.10] and Lemma~\ref{lem:II-scales}(a),
\begin{equation}\label{eq:II-fine}
 |\Delta\theta_{l''}|\le6\lambda_{l''}/t\quad(\text{all }l''),\qquad
 \sum_{l'>l}|\Delta\theta_{l'}|\le\frac{2b_w^2}t\le t^7 .
\end{equation}

\begin{lemma}[Diagonal pinning]\label{lem:II-diagpin}
Put $K_g:=(1/q_0+1)D_{\rm r}/u_w$.  Then
\begin{enumerate}
\item[(a)] $\sum_{l''\text{ of class R}}|\Delta\theta_{l''}|\le K_gt$;
\item[(b)] $\sum_{l''\in G(w)}(\tau_{l''})_-\le K_gt$, and
$e_0:=\Delta B1_{F^c}-\sum_{l''\in G(w)}\epsilon_{l''}\tau_{l''}u_{l''}
1_{F^c}$ satisfies $\|e_0\|_1\le K_gt+2b_w^2/t\le(K_g+1)t$.
\end{enumerate}
\end{lemma}

\begin{proof}
Let $l''\in\mathcal C(l)$ and $s\in S^{\rm nat}_{l''}(l)$.  For $l'\le l$,
$l'\ne l''$, we have $h_{l'}(s)=0$ (disjoint signature sets) and
$y_{l'}(s)=0$ ($s\notin T(l)$), so $u_{l'}(s)=0$; also $u_{l''}(s)=
v_{l''}(s)$, and $u_{l'}(s)=y_{l'}(s)/n_{l'}$ for $l'>l$.  By
[I,~(15)],
\[
 -\Delta B(s)=\Delta\theta_{l''}v_{l''}(s)+r(s),\qquad r(s):=\sum_{l'>l}
 \Delta\theta_{l'}y_{l'}(s)/n_{l'} .
\]
The sets $S^{\rm nat}_{l''}(l)$, $l''\in\mathcal C(l)$, are pairwise disjoint
subsets of $F^c$, so the numbers
\[
 E_{l''}:=\sum_{s\in S^{\rm nat}_{l''}(l)}\phi_{z_s}(\Delta B(s))
\]
satisfy $\sum_{l''}E_{l''}\le t/q_0$ [I,~Lemma~8.24], and
$\sum_{l''}\sum_s|r(s)|\le\frac43\sum_{l'>l}|\Delta\theta_{l'}|\le
\frac{8b_w^2}{3t}$ by \eqref{eq:II-fine} and $\|y_{l'}\|_1/n_{l'}\le\frac43$.
By [I,~Lemma~8.23(c),(d)],
\[
 E_{l''}\ge|\Delta\theta_{l''}|\sum_{s\in S^{\rm nat}_{l''}(l)}v_{l''}(s)
 \big(1+z_s\sgn\Delta\theta_{l''}\big)-2\sum_s|r(s)|
 \ge r^{\rm nat}_{l''}|\Delta\theta_{l''}|-2\sum_s|r(s)| .
\]
(a) For $l''$ of class R, $r^{\rm nat}_{l''}\ge u_w\mathfrak m_{l''}\ge
u_w/D_{\rm r}$ (Lemma~\ref{lem:II-scales}(c)); sum over class R and use
$16b_w^2/(3t)\le t$.
(b) For $l''\in G(w)$ we have $-\Delta\theta_{l''}v_{l''}(s)=
\epsilon_{l''}\tau_{l''}v_{l''}(s)$, and by [I,~Lemma~8.23(d)],
$\phi_{z_s}(\epsilon_{l''}\tau_{l''}v_{l''}(s))=|\tau_{l''}|v_{l''}(s)
(1-\epsilon_{l''}z_s\sgn\tau_{l''})$, which is
$(\tau_{l''})_-v_{l''}(s)(1+\epsilon_{l''}z_s)$ when $\tau_{l''}<0$.  Summing
over $s$, $(\tau_{l''})_-(2\mathfrak m_{l''}-r^{\rm nat}_{l''})\le E_{l''}
+2\sum_s|r(s)|$, and $2\mathfrak m_{l''}-r^{\rm nat}_{l''}\ge\mathfrak
m_{l''}\ge1/D_{\rm r}$.  Finally $e_0=-\sum\Delta\theta_{l''}u_{l''}1_{F^c}$
over class R and fine carriers, $\|u_{l''}\|_1\le q^*(u_{l''})=1$; use (a)
and \eqref{eq:II-fine}.
\end{proof}

No triangular unrolling and no product of rooms enter: the pinning sets
$S^{\rm nat}_{l''}(l)$ meet no coarse target, so each coarse amplitude is
compared only with its own signature and with fine carriers.

\begin{lemma}[Peak relations]\label{lem:II-peakrel}
Let $l''\in\mathcal C(l)$ with $k:=k(l'')\in P_m$, $\varsigma:=\sgn w_m(k)$,
$\mu:=\mu_{k,m}$, and $Y:=|\omega_{+,m}(k)|+|\omega_{-,m}(k)|$.  Then
$Y=-\varsigma\Delta\theta_{l''}/\lambda_{l''}-\Delta d_mM_m\ge0$, and
$Y\le t/(\lambda_{l''}\mu)$ if $\mu>0$.  Consequently:
\begin{enumerate}
\item[(a)] if $l''$ is of class R, then $\Delta d_mM_m\le|\Delta\theta_{l''}|/
\lambda_{l''}$, and if $\mu>0$ also $\Delta d_mM_m\ge-|\Delta\theta_{l''}|/
\lambda_{l''}-t/(\lambda_{l''}\mu)$;
\item[(b)] if $l''\in G(w)$ is of anti type, then $\tau_{l''}\le
-\lambda_{l''}\Delta d_mM_m$, hence $\Delta d_mM_m\le(\tau_{l''})_-/
\lambda_{l''}$;
\item[(c)] if $l''\in G(w)$ is of swallowing type, then $\tau_{l''}\ge
\lambda_{l''}\Delta d_mM_m$, and $\tau_{l''}\le\lambda_{l''}\Delta d_mM_m+t/\mu$
if $\mu>0$.
\end{enumerate}
If $\varrho_{l''}\ge1+u_w$, then $1/\mu\le mD_{\rm r}/(q_0\bar\vartheta_mu_w)$.
\end{lemma}

\begin{proof}
By [I,~(16)] and [I,~Lemma~8.27], $Y=-\varsigma\Delta\Theta_m(k)-\Delta d_m
M_m\ge0$ and $Y\le t/(\sigma_m|\alpha_m(k)|)=t/(\lambda_{l''}\mu)$ by
[I,~(2)]; and $\Delta\Theta_m(k)=\Delta\theta_{l''}/\lambda_{l''}$.  (a)
$-\varsigma\Delta\theta_{l''}\le|\Delta\theta_{l''}|$.  (b), (c):
$\Delta\theta_{l''}=-\epsilon_{l''}\tau_{l''}$, so $-\varsigma\Delta
\theta_{l''}/\lambda_{l''}=\varsigma\epsilon_{l''}\tau_{l''}/\lambda_{l''}$,
with $\varsigma\epsilon_{l''}=-1$ in (b) and $=+1$ in (c).  Finally
$\mu=q_0\Phi_{l''}\bar\vartheta_m(\varrho_{l''}-1)/m\ge q_0\bar\vartheta_mu_w/(mD_{\rm r})$.
\end{proof}

\begin{definition}[Shift sources]\label{def:II-sources}
Let $w$ be clean of level $l\ge l_f$ and $m\le N$.  An \emph{upper source}
of block $m$ at $w$ is one of: (U1) a peak $k(l'')$ of block $m$ with $l''$
of class R; (U2) an anti-type peak in $G(w)$ of block $m$; (U3) every
$l''\in G(w)$ of block $m$ with $q_{l''}<0$ is an anti-type peak, an
anti-type strict non-peak with $\varrho_{l''}\ge1-b_w$, or has
$\varrho_{l''}\le b_w$.  A \emph{lower source} is one of: (L1) a peak of block
$m$ of class R with $\varrho\ge1+u_w$; (L2) a swallowing-type peak in $G(w)$ of
block $m$ with $\varrho\ge1+u_w$; (L3) every $l''\in G(w)$ of block $m$ with
$q_{l''}>0$ has $\varrho_{l''}\le b_w$.
\end{definition}

A block that is $(+1)$-one-signed at $w$ satisfies (U3), and a block that
is $(-1)$-one-signed at $w$ satisfies (L3): if $\sgn q=+1$ on
$\Sigma_m(w)$, every $l''\in G(w)$ of block $m$ with $q_{l''}<0$ lies
outside $\Sigma_m(w)$, i.e.\ is an anti-type peak, an anti-type strict
non-peak with $\varrho\ge1-b_w$, or has $\varrho\le b_w$; if $\sgn q=-1$ on
$\Sigma_m(w)$, every carrier with $q>0$ lies outside $\Sigma_m(w)$ and is
not a peak (swallowing-type peaks belong to $\Sigma_m(w)$), so it has
$\varrho\le b_w$.

\begin{lemma}[Sources bound the shift on one side]\label{lem:II-sources}
There is $K_d=C_fD_{\rm r}^2/u_w$ such that: if block $m$ has an upper
source at $w$, then $\Delta d_mM_m\le K_dt$; if it has a lower source, then
$\Delta d_mM_m\ge-K_dt$.  Each half uses only its own source.
\end{lemma}

\begin{proof}
We use [I,~(17)]: $\Delta d_mM_m=\frac1{mC_m}\sum_k\Phi_m(k)w_m(k)
\Delta\theta_{\mathfrak j(k,m)}+r_m$ with $|r_m|\le2t/\sigma_m$.  For
$l''\in G(w)$ the summand equals $-q_{l''}\tau_{l''}$; for class R and fine
carriers it is at most $|\Delta\theta_{l''}|/(mC_m)$ in modulus.  Recall
$|q_{l''}|\le\Phi_{l''}M_m/(mC_m)\le1/C_m$, and $|q_{l''}|=\Phi_{l''}M_m
\varrho_{l''}/(mC_m)\le b_w/C_m$ when $\varrho_{l''}\le b_w$.

\emph{Upper bounds.}  (U1): Lemma~\ref{lem:II-peakrel}(a),
Lemma~\ref{lem:II-diagpin}(a) and $\lambda\ge1/D_{\rm r}$ give $\Delta d_m
M_m\le D_{\rm r}K_gt$.  (U2): Lemma~\ref{lem:II-peakrel}(b) and
Lemma~\ref{lem:II-diagpin}(b) give the same bound.  (U3): In
[I,~(17)], the class R and fine terms are at most $(K_gt+t^7)/(mC_m)$
(Lemma~\ref{lem:II-diagpin}(a), \eqref{eq:II-fine}); carriers with $q>0$
contribute $-q\tau\le q(\tau)_-$, in total at most $K_gt/C_m$; carriers
with $q=0$ contribute $0$.  For $q<0$ the contribution is $|q|\tau$, and by
(U3) the carrier is of one of three kinds.  At an anti-type peak,
$\tau\le-\lambda\Delta d_mM_m$ (Lemma~\ref{lem:II-peakrel}(b)).  At an
anti-type strict non-peak $k$ with $\varrho\ge1-b_w$, put $\varsigma:=\sgn
w_m(k)$; by [I,~Lemma~8.8(f)] and $|d_\pm t|\le\frac12$,
$\varsigma(\omega_{-,m}-\omega_{+,m})(k)\ge-3\gap_m(k)/t$, while
$\omega_{-,m}-\omega_{+,m}=-\Delta\Theta_m-\Delta d_mw_m$ gives
$\varsigma(\omega_{-,m}-\omega_{+,m})(k)=-\tau/\lambda-\Delta d_m|w_m(k)|$;
so $\tau\le\lambda(3\gap_m(k)/t-\Delta d_m|w_m(k)|)$ with $\gap_m(k)=M_m(1-
\varrho)\le b_w$.  At a carrier with $\varrho\le b_w$, $|q|\tau\le(b_w/C_m)
6\lambda/t$.  Collecting, with $\sum\lambda\le1$ and $b_w\le t^4$,
\[
 \Delta d_mM_m\Big(1+\sum_{\rm a.p.}|q|\lambda+\sum_{\rm a.n.}|q|\lambda
 \frac{|w_m(k)|}{M_m}\Big)\le\frac{K_gt+t^7}{mC_m}+\frac{K_gt}{C_m}+
 \frac{9t^3}{C_m}+\frac{2t}{\sigma_m},
\]
the sums running over the anti-type peaks and the anti-type near-threshold
strict non-peaks.  If $\Delta d_m\ge0$ the bracket is $\ge1$; otherwise the
bound is trivial.

\emph{Lower bounds.}  (L1): Lemma~\ref{lem:II-peakrel}(a) with
$1/\mu\le mD_{\rm r}/(q_0\bar\vartheta_mu_w)$: $\Delta d_mM_m\ge-(D_{\rm r}K_g+
mD_{\rm r}^2/(q_0\bar\vartheta_mu_w))t$.  (L2): Lemma~\ref{lem:II-peakrel}(c) gives
$\Delta d_mM_m\ge\tau/\lambda-t/(\lambda\mu)\ge-(\tau)_-/\lambda-
t/(\lambda\mu)$, and the same bound follows.  (L3): carriers with $q>0$
contribute $-q\tau\ge-(b_w/C_m)6\lambda/t$, carriers with $q<0$ contribute
$|q|\tau\ge-|q|(\tau)_-$, in total $\ge-K_gt/C_m$; with the class R and
fine terms and $r_m$, $\Delta d_mM_m\ge-C_fK_gt$.  In all cases the constant
is at most $C_fD_{\rm r}^2/u_w$.
\end{proof}

\begin{lemma}[Robust-margin peaks]\label{lem:II-donorpeak}
There is $l_f$ such that for every sub-window $w$ of level $l\ge l_f$ which
is clean for $f$ and every block $m\le N$ there is $c\in\mathcal C(l)$ of
block $m$ such that $k(c)$ is a peak with $\varrho_c\ge1+u_w$; in particular
its margin is at least $q_0\bar\vartheta_mu_w/(mD_{\rm r})$.  Every such $c$ is of
class R, or belongs to $P(w)$ (types (P-i) or (P-iii)).
\end{lemma}

\begin{proof}
By [I,~Lemma~1.7], $\sum_{k\in P_m}|\alpha_m(k)|=1$, and by [I,~(2)],
$\sigma_m|\alpha_m(k)|=\lambda_{k,m}\mu_{k,m}$.  A peak $k=k(l'')$ with
$l''\in\mathcal C(l)$ and $\varrho_{l''}\le1+b_w$ has $\mu_{k,m}\le q_0
\Phi_{l''}\bar\vartheta_mb_w/m$, hence $|\alpha_m(k)|\le q_0\bar\vartheta_mb_w\Phi_{l''}^2/
\sigma_m$, and the total of these is at most $(M_m/C_m)b_w$, since
$q_0\bar\vartheta_m/\sigma_m=M_m/C_m$ and $\sum_k\Phi_m(k)^2\le1$.  A peak of a
fine carrier has $\mu_{k,m}\le|u_{k,m}(\xi)|\le q_0$, and these contribute
at most $q_0\sum_{l'>l}\lambda_{l'}/\sigma_m\le q_0b_w^2/\sigma_m$.  By
Lemma~\ref{lem:II-scales}(e) we may take $l_f$ so large that $(M_m/C_m)b_w+
q_0b_w^2/\sigma_m<1$ for every $m\le N$ and every sub-window of level
$\ge l_f$.  Then some coarse peak of block $m$ has $\varrho>1+b_w$, hence
$\varrho\ge1+u_w$, since $w$ is clean.  If $c$ is of class G, it is an
anti-type peak (P-i) or a swallowing-type peak with $\varrho\ge1+u_w$ (P-iii).
\end{proof}

Since $\|\alpha_m\|_1=1$, tiny-margin peaks and fine peaks carry only
$O(b_w)$ of the $\alpha$-mass; this is why a robust-margin peak, and hence
a \emph{threshold donor} (Lemma~\ref{lem:II-donorraise}), exists in every
block.

\begin{definition}[Shift pattern of $f$ at $w$; condition (SH$_w$)]
\label{def:II-shiftpattern}
Let $w$ be clean of level $l\ge l_f$.  Let $I_{\rm up}(w)$ (resp.\
$I_{\rm lo}(w)$) be the set of blocks $m\le N$ without an upper (resp.\
lower) source at $w$; $P^\star(w)$ the set of $l''\in\mathcal C(l)$ of
blocks in $I_{\rm up}(w)\cup I_{\rm lo}(w)$ such that $k(l'')$ is a peak
with $\varrho_{l''}\ge1+u_w$; $\mathrm{vs}_{l''}:=\sgn w_{m(l'')}(k(l''))$;
$B_f(w):=G(w)\setminus P^\star(w)$ with the signs $\epsilon_{l''}$ of
Definition~\ref{def:II-class}.  The \emph{shift pattern of $f$ at $w$} is
$\pi(w):=(I_{\rm up}(w),I_{\rm lo}(w),P^\star(w),\mathrm{vs},B_f(w),
\epsilon)$, and $c_{\pi(w)}(f)$ is the value at $f$ of the corresponding
rate object (O6) (Definition~\ref{def:II-rates}).  We say that
\textup{(SH$_w$)} holds if $I_{\rm up}(w)\cup I_{\rm lo}(w)=\emptyset$ or
$c_{\pi(w)}(f)\ge u_w$.
\end{definition}

By Lemma~\ref{lem:II-shiftpin}(a) below, $I_{\rm up}(w)\cap
I_{\rm lo}(w)=\emptyset$ and $P^\star(w)\subseteq G(w)$, so $\pi(w)$ is a
shift pattern of level $l$ in the sense of (O6) (every block $m\le N$
contains a coarse carrier by Lemma~\ref{lem:II-donorpeak}).  As $w$ is
clean, (SH$_w$) fails exactly when $I_{\rm up}(w)\cup I_{\rm lo}(w)\ne
\emptyset$ and $c_{\pi(w)}(f)\le b_w$.

\begin{lemma}[Shift pinning at a clean sub-window]\label{lem:II-shiftpin}
Let $w$ be clean of level $l\ge l_f$.
\begin{enumerate}
\item[(a)] $I_{\rm up}(w)\cap I_{\rm lo}(w)=\emptyset$, and $P^\star(w)
\subseteq G(w)$.
\item[(b)] If \textup{(SH$_w$)} holds, then $|\Delta d_m|M_m\le K'_dt$ for every
block $m$, with $K'_d:=C_fl\,D_{\rm r}^2u_w^{-2}$.
\end{enumerate}
\end{lemma}

\begin{proof}
(a) Let $m\in I_{\rm up}(w)\cap I_{\rm lo}(w)$.  Then (U1), (U2) and (L2)
fail for $m$, so every coarse peak of block $m$ is a swallowing-type peak in
$G(w)$ with $\varrho\le1+b_w$; this contradicts
Lemma~\ref{lem:II-donorpeak}.  A class-R peak with $\varrho\ge1+u_w$ is a
source of both kinds ((U1) and (L1)), so blocks in $I_{\rm up}(w)\cup
I_{\rm lo}(w)$ contain none, and $P^\star(w)\subseteq G(w)$.

(b) If $I_{\rm up}(w)\cup I_{\rm lo}(w)=\emptyset$, apply
Lemma~\ref{lem:II-sources}.  Otherwise $c:=c_{\pi(w)}(f)\ge u_w$.  For
$l''\in P^\star(w)$ of block $m$, Lemma~\ref{lem:II-peakrel} gives
$-\Delta\theta_{l''}=\mathrm{vs}_{l''}\lambda_{l''}(\Delta d_mM_m+Y_{l''})$
with $0\le\lambda_{l''}Y_{l''}\le t/\mu_{k(l''),m}\le tmD_{\rm r}/(q_0
\bar\vartheta_mu_w)$.  For $m\in I_{\rm up}(w)$ the block has a lower source, so
$\Delta d_mM_m\ge-K_dt$ (Lemma~\ref{lem:II-sources}); put $\delta_m:=
(\Delta d_m)_+M_m\ge0$.  For $m\in I_{\rm lo}(w)$ put $\delta_m:=
-(\Delta d_m)_-M_m\le0$; then $(\Delta d_m)_+M_m\le K_dt$.  In both cases
$|\Delta d_mM_m-\delta_m|\le K_dt$.  With $\Pi_m:=\sum_{l''\in P^\star(w),\,
m(l'')=m}\mathrm{vs}_{l''}\lambda_{l''}u_{l''}1_{F^c}$ (so $\|\Pi_m\|_1\le
\sum\lambda\le\frac13$) and $-\Delta\theta_{l''}=\epsilon_{l''}
\tau_{l''}=\epsilon_{l''}(\tau_{l''})_+-\epsilon_{l''}(\tau_{l''})_-$ on
$B_f(w)$, [I,~(15)] gives
\[
 \Delta B1_{F^c}=\sum_m\delta_m\Pi_m+\sum_{l''\in B_f(w)}(\tau_{l''})_+
 \epsilon_{l''}u_{l''}1_{F^c}+R,
\]
where, by Lemma~\ref{lem:II-diagpin} and \eqref{eq:II-fine},
\[
 \|R\|_1\le K_dt+\frac{l\,tmD_{\rm r}}{q_0\bar\vartheta_mu_w}+\sum_{G(w)}
 (\tau)_-+\sum_{\rm class\,R}|\Delta\theta|+\sum_{l'>l}|\Delta\theta_{l'}|
 \le C_flD_{\rm r}^2t/u_w .
\]
The vector $x:=((\tau_{l''})_+)_{l''\in B_f(w)}$ is admissible in the
infimum defining $c(\delta;\pi(w))$, and $\delta$ lies in the sign cone of
(O6).  By [I,~Lemma~8.23(c)] and [I,~Lemma~8.24],
$c(\delta;\pi(w))\le\sum_{j\notin F}\phi_{z_j}(\Delta B(j))+2\|R\|_1\le
t/q_0+2\|R\|_1$.  As $c(\cdot;\pi(w))$ is positively homogeneous,
$c(\delta;\pi(w))\ge c\,\|\delta\|_1\ge u_w\|\delta\|_1$.  Hence
$\|\delta\|_1\le C_flD_{\rm r}^2t/u_w^2$ and $|\Delta d_m|M_m\le|\delta_m|+
K_dt\le C_flD_{\rm r}^2t/u_w^2$.
\end{proof}

\begin{lemma}[Pinned carriers]\label{lem:II-pinned}
Let $w$ be clean of level $l\ge l_f$, and suppose $|\Delta d_m|M_m\le K't$
for every block $m$, with some $K'\ge1$.  Put $K_P:=K_g+K'+C_fD_{\rm r}/u_w$
and $K_O:=C_fD_{\rm r}(K'+K_g+lK_P)/u_w$.  Then $|\tau_{l''}|\le K_Pt$ for
$l''$ of type (P-i), (P-ii) or (P-iii), $|\tau_{l''}|\le K_Ot$ for $l''$ of
type (P-iv), and $|\Delta\theta_{l''}|\le K_gt$ for $l''$ of class R.  Under
\textup{(SH$_w$)} we may take $K'=K'_d$, and then
\[
 K_{\rm pin}:=\max(K_g,K_P,K_O,K'_d)\le C_fl^2D_{\rm r}^3u_w^{-3} .
\]
\end{lemma}

\begin{proof}
$(\tau)_-\le K_gt$ on $G(w)$ by Lemma~\ref{lem:II-diagpin}(b).  (P-i):
$\tau\le-\lambda\Delta d_mM_m\le K't$ (Lemma~\ref{lem:II-peakrel}(b),
$\lambda\le1$).  (P-ii): as in the proof of Lemma~\ref{lem:II-sources},
$\tau\le\lambda(3\gap/t+|\Delta d_m|M_m)\le3b_w/t+K't\le(1+K')t$.  (P-iii):
$\tau\le\lambda\Delta d_mM_m+t/\mu\le K't+tmD_{\rm r}/(q_0\bar\vartheta_mu_w)$
(Lemma~\ref{lem:II-peakrel}(c)).  Class R: Lemma~\ref{lem:II-diagpin}(a).

(P-iv): let block $m$ be $\sigma$-one-signed.  In [I,~(17)] the summands
of $l''\in\Sigma_m(w)$ are $-q_{l''}\tau_{l''}$ with $\sgn q_{l''}=\sigma$,
so
\begin{multline*}
 \sum_{\Sigma_m(w)}|q|(\tau)_+\le\mathcal A_m:=|\Delta d_m|M_m+\frac{2t}
 {\sigma_m}\\+\Big|\sum_{l''\notin\Sigma_m(w)}\frac{\Phi_{l''}w_m(k(l''))
 \Delta\theta_{l''}}{mC_m}\Big|+\sum_{\Sigma_m(w)}|q|(\tau)_- .
\end{multline*}
The carriers of block $m$ outside $\Sigma_m(w)$ are of class R or fine
(total $\le(K_gt+t^7)/(mC_m)$), of types (P-i), (P-ii) (at most $l$
carriers, each $\le K_Pt/C_m$), or in $G(w)$ with $\varrho\le b_w$
($|q\tau|\le(b_w/C_m)6\lambda/t$, in total $\le t^3/C_m$).  So $\mathcal
A_m\le C_f(K'+K_g+lK_P)t$.  On $\Sigma_m(w)$, $|q|\ge\Phi M_mu_w/(mC_m)\ge
M_mu_w/(mC_mD_{\rm r})$ (at peaks $|w|=M_m$, at strict non-peaks
$|w|=M_m\varrho\ge M_mu_w$).  Hence $(\tau)_+\le K_Ot$ after enlarging $C_f$.

Under (SH$_w$), Lemma~\ref{lem:II-shiftpin}(b) gives $K'=K'_d$; then
$K_P\le C_flD_{\rm r}^2u_w^{-2}$ and $K_O\le C_fl^2D_{\rm r}^3u_w^{-3}$.
\end{proof}

By Lemma~\ref{lem:II-donorpeak} and Lemma~\ref{lem:II-pinned}, every
robust-margin peak of Lemma~\ref{lem:II-donorpeak} is pinned:
$|\Delta\theta_c|\le K_{\rm pin}t$ under (SH$_w$).  The anti-type
near-threshold strict non-peaks (P-ii) are pinned like anti-type peaks;
their tiny gaps are therefore not rates that need exactification.  Their
swallowing-type counterparts, the (K4) carriers, are not pinned.

\subsection*{Patterns, rays and condition (VR$_w$)}

\begin{definition}[The pattern of $f$ at $w$; ray components; (VR$_w$)]
\label{def:II-pattern}
Let $w$ be clean of level $l\ge l_f$.  The \emph{pattern of $f$ at $w$} is
$\kappa(w):=(G(w),P(w),\epsilon,F\cap T(l),\mathrm{type}_w)$ with
$\mathrm{type}_w(j):=\sgn z_j$ for contact-like and $\mathrm{type}_w(j):=0$
for free-robust $j\in T(l)\setminus F$; it is a pattern of level $l$
(Definition~\ref{def:II-constants}(a)), with kept set $\mathrm{Kp}(w)$.
For $r\in\mathrm{Ext}(\kappa(w))$ and a block $m$ meeting $\supp r$, the
\emph{$d$-component} of $r$ in block $m$ at a first row $f'$ is
\[
 D_{r,m}(f'):=\sum_{l''\in\supp r,\ m(l'')=m}r(l'')\,\epsilon_{l''}
 u_{l''}(\zh'),
\]
and the value of the rate object (O5) of $(\kappa(w),r,m)$ at $f$ is
$|D_{r,m}(f)|/\Phi_{\max}(r,m)$.  The component is \emph{tiny} or
\emph{robust} according to this value.  Condition \textup{(VR$_w$)} holds if
every $r\in\mathrm{Ext}(\kappa(w))$ has at most one robust component.

For box radii $\beta\in(0,\infty)^{E}$ the box $Z_\kappa(\beta)=\{\tau\in
C(\kappa):\tau_{l''}\le\beta_{l''}\ (l''\in\mathrm{Kp}(\kappa))\}$ and the
constant $G^{**}(l)$ (the maximum over all patterns $\kappa$ of level $l$
of the $\ell_1$-Hoffman constant of the system defining $Z_\kappa(\beta)$,
which does not depend on $\beta$) are those of
Definition~\ref{def:II-constants}(d).  The system is a configuration system
of level $l$, so $G^{**}(l)\le H_{\rm conf}(l)$.
\end{definition}

\begin{remark}[Components in one-signed blocks]\label{rem:II-onesigned}
(a) If block $m$ is one-signed at $w$, every kept carrier of block $m$ has
$\varrho\le b_w$, so $|\epsilon_{l''}u_{l''}(\zh)|=\varrho\Phi_{l''}\bar\vartheta_m/m\le
b_w\Phi_{l''}\bar\vartheta_m/m$, and every component in block $m$ has rate at most
$b_w\bar\vartheta_m/m$ (as $\|r\|_1=1$).  Since $b_w\le T_{\rm lo}(w)^4\le(u_w/(4\mathrm{Des}))^{160}$
(Lemma~\ref{lem:II-scales}(a),(b)), $b_w/u_w\le u_w^{159}\to0$, while
$\bar\vartheta_m/m$ is an $f$-constant, the rate is $<u_w$ for $l\ge l_f$, hence tiny: one-signed blocks carry no robust components.
(b) If $N=1$, or if every extreme ray of $C(\kappa(w))$ is supported in
one block, (VR$_w$) holds.
(c) For a kept carrier that is a strict non-peak, $q_{l''}=\epsilon_{l''}
\Phi_{l''}w_m(k)/(mC_m)=\mathrm{val}_{l''}/A_m$ by the clamp formula
($w_m(k)=C_m\hat\zeta_m(k)/(\Phi_{l''}^2A_m)$, [I,~Lemma~1.7]); so
$D_{r,m}/A_m$ is the $d$-sum of the switching direction $r$ in block $m$.
For a (K4) carrier that is a peak, $q_{l''}=\mathrm{val}_{l''}/(\varrho_{l''}
A_m)$ with $\varrho_{l''}\in[1,1+b_w]$.
(d) $C(\kappa)\subseteq[0,\infty)^E$, so for $P'\subseteq E$ the set
$C(\kappa)\cap\{\tau_{l''}=0,\ l''\in P'\}$ is a face of $C(\kappa)$, and
the extreme rays of a face are extreme rays of $C(\kappa)$: dropping more
carriers can only remove rays.
\end{remark}

\subsection{Far pulls and private banks}\label{subsec:II-tools}

For $A\in\ell_1\setminus\{0\}$ and $z'\in B_{\ell_\infty}$ with $z'_j=\sgn A_j$
on $\supp A$ we write $f[A,z']:=f[A/q^*(A),z']$ for the first row with
prescribed data (Subsection~\ref{subsec:II-comp}); its forced data are
$e'=U^*A/\|U^*A\|_H$ (independent of the normalization), $\zh'=z'+Ue'$, and
$w'_m=J_m(R_m^{**}\zh')$.  In this subsection $(A^0,z^0)$ is such a pair,
$\nu_0:=\|U^*A^0\|_H$ and $e^0:=U^*A^0/\nu_0$.

\begin{definition}[Banks and far pulls]\label{def:II-bankpull}
Let $J$ be a finite set disjoint from $\supp A^0$ and $(c_j)_{j\in J}$ real,
$A:=A^0+\sum_{j\in J}c_je^*_j$.  A coordinate $j\in J$ with $|z^0_j|=1$ and
$\sgn c_j=z^0_j$ (the sign vector is not changed at $j$) is a \emph{bank} of
mass $|c_j|$.  A coordinate $j\in J\cap S_{l''}$ with $z^0_j=\epsilon$, at
which $c_j=-\epsilon\mathfrak p$ with $\mathfrak p>0$ and the sign vector is
changed to $-\epsilon$, is a \emph{far pull} (of carrier $l''$, with sign
$\epsilon$) of mass $\mathfrak p$.  In both cases the resulting pair $(A,z)$
satisfies $z_j=\sgn A_j$ on $\supp A$, so $f[A,z]$ is defined.
\end{definition}

\begin{lemma}[Masses off the support]\label{lem:II-bankformula}
In the situation of Definition~\textup{\ref{def:II-bankpull}} put
$X:=\sum_{j\in J}c_j\mu^\star_jk_j$ and $e^A:=U^*A/\|U^*A\|_H$.  Then
$\|U^*A\|_H^2=\nu_0^2+\|X\|^2$ and, for every $u\in\ell_1$,
\begin{equation}\label{eq:II-bankformula}
 \ip{U^*u}{e^A}=\frac{\nu_0\ip{U^*u}{e^0}+\sum_{j\in J}c_j(\mu^\star_j)^2
 u(j)}{(\nu_0^2+\|X\|^2)^{1/2}} .
\end{equation}
Consequently:
\textup{(i)} if $u$ vanishes on $J$, then
\[
 |\ip{U^*u}{e^A}-\ip{U^*u}{e^0}|\le\|U\|\|u\|_1\|X\|^2/(2\nu_0^2);
\]
\textup{(ii)} in general $\ip{U^*u}{e^A}-\ip{U^*u}{e^0}=\sum_{j\in J}c_j
(\mu^\star_j)^2u(j)/\nu_0+R_u$ with $|R_u|\le(\|U\|\|u\|_1+\sum_{j\in J}
|c_j|(\mu^\star_j)^2|u(j)|/\nu_0)\|X\|^2/(2\nu_0^2)$;
\textup{(iii)} $\|e^A-e^0\|_H\le2\|X\|/\nu_0$.
\end{lemma}

\begin{proof}
Lemma~\ref{lem:II-base}(b) with $a:=A^0$ gives $U^*A=U^*A^0+X$, $X\perp
U^*A^0$, the Pythagoras identity and $\ip{U^*u}{X}=\sum_jc_j(\mu^\star_j)^2
u(j)$; divide by $\|U^*A\|$.  Put $x:=\|X\|^2/\nu_0^2$, $p:=\ip{U^*u}{e^0}$
and $s:=\sum_jc_j(\mu^\star_j)^2u(j)$; then $\ip{U^*u}{e^A}=(p+s/\nu_0)
(1+x)^{-1/2}$ and $0\le1-(1+x)^{-1/2}\le x/2$, while $|p|\le\|U^*u\|\le
\|U\|\|u\|_1$.  This gives (i) ($s=0$) and (ii).  (iii):
$\|x'/\|x'\|-y/\|y\|\|\le2\|x'-y\|/\|y\|$ with $x'=U^*A$, $y=U^*A^0$.
\end{proof}

For a carrier $l''$ and a coordinate $j\in S_{l''}$ beyond $s_{\max}(l)$,
$l\ge l''$, the only carrier $\le l$ whose vector is nonzero at $j$ is
$l''$, with $u_{l''}(j)=v_{l''}(j)$ (Theorem~\ref{thm:II-lemmaB}(d)).
Hence a pull at $j$ changes $\epsilon u_{l''}(\zh)$ by $-2v_{l''}(j)$ through
the sign vector, a bank at $j$ changes it by $+\mathfrak b(\mu^\star_j)^2
v_{l''}(j)/\nu_0$ at first order through the Hilbert part, and by
\eqref{eq:II-bankformula} every other carrier $\le l$ moves only at second
order.  This is the diagonal-base mechanism used below.

\begin{lemma}[Exact two-sided tuning]\label{lem:II-TU}
Let $l\ge l_f$, let $(A^2,z^2)$ be as above with $F^2:=\supp A^2$ finite,
$\nu_2:=\|U^*A^2\|_H\ge\nu$ and $\|A^2\|_1\le2$, let $L_0\subseteq\mathcal C(l)$
and signs $(\epsilon_{l''})_{l''\in L_0}$, and assume:
\textup{(T1)} $z^2=\epsilon_{l''}$ on $S_{l''}\setminus[1,s_{\max}(l)]$ for
every $l''\in L_0$; \textup{(T2)} $F^2\cap S_{l''}\cap(s_{\max}(l),\infty)=
\emptyset$ for $l''\in L_0$.  Let $x\in\R^{L_0}$, $\varepsilon_x:=\|x\|_\infty$.
There are $c_T=c_f\mathrm{Des}^{-1}$ and $C_T=C_f\mathrm{Des}$ such that if
$\varepsilon_x\le c_T$ the following holds.  For $l''\in L_0$ let
$j'_{l''}:=s'_{l''}(l)$ \textup{(}the \emph{bank coordinate}\textup{)}, let
$j_{l''}\in S_{l''}$ be the largest element with $v_{l''}(j_{l''})\ge
\varepsilon_x$ \textup{(}the \emph{pull coordinate}; if $x=0$ there are no
moves\textup{)}, and $\mathfrak p_{l''}:=24\lambda_{l''}v_{l''}(j_{l''})$.  Then
$j_{l''}>j'_{l''}$, $v_{l''}(j_{l''})\in[\varepsilon_x,2^{G(l)}\varepsilon_x)$,
and there is $\mathfrak n\in(0,\infty)^{L_0}$ such that the row
$f^\#:=f[A^\#,z^\#]$ with
\[
 A^\#:=A^2-\sum_{l''\in L_0}\epsilon_{l''}\mathfrak p_{l''}e^*_{j_{l''}}+
 \sum_{l''\in L_0}\mathfrak n_{l''}\epsilon_{l''}e^*_{j'_{l''}},\qquad
 z^\#_{j_{l''}}:=-\epsilon_{l''},\quad z^\#:=z^2\text{ elsewhere},
\]
satisfies:
\begin{enumerate}
\item[(a)] $\mathrm{val}^\#_{l''}=\mathrm{val}^{(2)}_{l''}+x_{l''}$ exactly for
every $l''\in L_0$, where $\mathrm{val}^{(2)}$ refers to $f[A^2,z^2]$;
\item[(b)] $|u_k(\zh^\#)-u_k(\zh^{(2)})|\le C_T\varepsilon_x^2$ for every
carrier $k\le l$ with $k\notin L_0$;
\item[(c)] $\|\mathfrak n\|_\infty\le C_T\varepsilon_x$, $\mathfrak p_{l''}\le
6\cdot2^{G(l)}\varepsilon_x$, and the Hilbert vector
\[
 X_{\rm tu}:=\sum_{l''\in L_0}\big(-\epsilon_{l''}\mathfrak p_{l''}
 \mu^\star_{j_{l''}}k_{j_{l''}}+\mathfrak n_{l''}\epsilon_{l''}
 \mu^\star_{j'_{l''}}k_{j'_{l''}}\big)
\]
has $\|X_{\rm tu}\|^2\le C_T\varepsilon_x^2$;
\item[(d)] for $l''\in L_0$, $|u_{l''}(\zh^\#)-u_{l''}(\zh^{(2)})|\le
\varepsilon_x$, and $(A^\#,z^\#)$ is admissible: the pulls carry masses of
the new sign $-\epsilon_{l''}=z^\#_{j_{l''}}$, and the banks carry masses of
sign $\epsilon_{l''}=z^\#_{j'_{l''}}$.
\end{enumerate}
\end{lemma}

\begin{proof}
\emph{Coordinates.}  By (T1), (T2) and \eqref{eq:II-Finside}, all
$j_{l''},j'_{l''}$ lie outside $F^2$, are pairwise distinct (distinct
signature sets) and carry $z^2=\epsilon_{l''}$.  The values $v_{l''}(s)=
\delta_{l''}2^{-s}/n_{l''}$, $s\in S_{l''}$, decrease by the factor
$2^{-G_{l''}}\ge2^{-G(l)}$ between consecutive elements, so the largest
$s$ with $v_{l''}(s)\ge\varepsilon_x$ has $v_{l''}(s)<2^{G(l)}\varepsilon_x$.
By Lemma~\ref{lem:II-scales}(c),(d), $v_{l''}(j'_{l''})\ge
\mathrm{Des}^{-1/6}$ and $2^{G(l)}\le\mathrm{Des}^{1/6}$; so for
$\varepsilon_x<\mathrm{Des}^{-1/3}$ we get $v_{l''}(j'_{l''})>2^{G(l)}
\varepsilon_x$, hence $j_{l''}>j'_{l''}$.  For a carrier $k\le l$ and
$l''\in L_0$, $u_k$ vanishes at $j_{l''}$ and $j'_{l''}$ unless $k=l''$.
The admissibility statement in (d) is clear.

\emph{Step 1 (pulls).}  Let $A^p:=A^2-\sum\epsilon_{l''}\mathfrak p_{l''}
e^*_{j_{l''}}$, $z^p:=z^\#$, $\nu_p:=\|U^*A^p\|\ge\nu_2\ge\nu$, and
$X_p:=\sum-\epsilon_{l''}\mathfrak p_{l''}\mu^\star_{j_{l''}}k_{j_{l''}}$, so
$\|X_p\|^2\le\sum\mathfrak p^2\le l(24\cdot2^{G(l)}\varepsilon_x)^2$.  For
$l''\in L_0$, the sign change at $j_{l''}$ changes $\mathrm{val}_{l''}$ by
$-2v_{l''}(j_{l''})$, and Lemma~\ref{lem:II-bankformula}(ii) with base
$A^2$ gives
\[
 \mathrm{val}^p_{l''}-\mathrm{val}^{(2)}_{l''}=-2v_{l''}(j_{l''})-
 \mathfrak p_{l''}(\mu^\star_{j_{l''}})^2v_{l''}(j_{l''})/\nu_2+R_{l''},
 \qquad|R_{l''}|\le C\|X_p\|^2/\nu^2 .
\]

\emph{Step 2 (banks; explicit solution).}  For $\mathfrak n\in[0,\infty)^{L_0}$
let $A(\mathfrak n):=A^p+\sum\mathfrak n_{l''}\epsilon_{l''}e^*_{j'_{l''}}$,
$\mathcal N(\mathfrak n):=(\nu_p^2+\sum_{l''}\mathfrak n_{l''}^2
(\mu^\star_{j'_{l''}})^2)^{1/2}$, and write $s'_{l''}:=\mu^\star_{j'_{l''}}$,
$v'_{l''}:=v_{l''}(j'_{l''})$, $\varphi_{l''}:=\epsilon_{l''}\ip{U^*u_{l''}}
{e^p}$ ($|\varphi_{l''}|\le\|U\|$).  By \eqref{eq:II-bankformula} with base
$A^p$ (the sign vector is not changed at banks),
\[
 \mathrm{val}_{l''}(\mathfrak n)-\mathrm{val}^p_{l''}=\frac{\mathfrak n_{l''}
 s'^2_{l''}v'_{l''}+\varphi_{l''}\nu_p}{\mathcal N(\mathfrak n)}-
 \varphi_{l''} .
\]
The required increments $\Delta_{l''}:=x_{l''}+\mathrm{val}^{(2)}_{l''}-
\mathrm{val}^p_{l''}$ satisfy $\varepsilon_x/2\le\Delta_{l''}\le4\cdot
2^{G(l)}\varepsilon_x$ for $\varepsilon_x\le c_T$ (Step~1, $v(j)\ge
\varepsilon_x\ge|x_{l''}|$, and $\mathfrak p(\mu^\star_j)^2v(j)\le24v(j)^2$).
For $y\ge\nu_p$ put
\[
 \mathfrak n_{l''}(y):=\frac{\Delta_{l''}y+\varphi_{l''}(y-\nu_p)}
 {s'^2_{l''}v'_{l''}},\qquad\mathcal M(y):=\Big(\nu_p^2+\sum_{l''}s'^2_{l''}
 \mathfrak n_{l''}(y)^2\Big)^{1/2} .
\]
If $\mathcal M(y)=y$, then $\mathcal N(\mathfrak n(y))=y$ and
$\mathfrak n(y)$ solves all equations $\mathrm{val}_{l''}(\mathfrak n)-
\mathrm{val}^p_{l''}=\Delta_{l''}$ simultaneously.  By
Lemma~\ref{lem:II-scales}(c),(d), $(s'_{l''}v'_{l''})^{-2}\le
\mathrm{Des}^{1/3}$ and $(s'^2_{l''}v'_{l''})^{-1}\le\mathrm{Des}^{1/6}$.
Put $\kappa_T:=2(\mathcal M(\nu_p)-\nu_p)$; since $\mathcal M(\nu_p)^2-\nu_p^2
=\sum\Delta_{l''}^2\nu_p^2/(s'_{l''}v'_{l''})^2$,
$\kappa_T\le C_f\mathrm{Des}\,\varepsilon_x^2$ (we use $\nu_p\ge\nu$,
$|L_0|\le l$ and $l\,4^{G(l)}\le\mathrm{Des}^{1/2}$).  On $\mathcal I:=[\nu_p,
\nu_p+\kappa_T]$,
\begin{multline*}
 |\mathcal M'(y)|=\Big|\sum_{l''}\frac{\mathfrak n_{l''}(y)(\Delta_{l''}+
 \varphi_{l''})}{v'_{l''}\mathcal M(y)}\Big|\\\le\sum_{l''}\frac{(2\Delta_{l''}
 +\kappa_T)(\Delta_{l''}+\|U\|)}{(s'_{l''}v'_{l''})^2\nu}\le C_f
 \mathrm{Des}^{2/3}\varepsilon_x\le\tfrac12
\end{multline*}
for $\varepsilon_x\le c_T$ (we used $|L_0|\le l\le\mathrm{Des}^{1/6}$,
$2^{G(l)}\le\mathrm{Des}^{1/6}$, $\nu_p\le1$).  Hence $\nu_p\le\mathcal M(y)
\le\mathcal M(\nu_p)+(y-\nu_p)/2\le\nu_p+\kappa_T$ on $\mathcal I$, and the
continuous self-map $\mathcal M$ of $\mathcal I$ has a fixed point $y^*$.
Put $\mathfrak n:=\mathfrak n(y^*)$; then $\mathfrak n_{l''}\ge(\Delta_{l''}
\nu-\|U\|\kappa_T)/(s'^2_{l''}v'_{l''})>0$ and $\mathfrak n_{l''}\le
\mathrm{Des}^{1/6}(4\cdot2^{G(l)}\varepsilon_x\cdot2+\kappa_T)\le C_f
\mathrm{Des}^{1/3}\varepsilon_x$.  By construction
$\mathrm{val}^\#_{l''}=\mathrm{val}^p_{l''}+\Delta_{l''}=
\mathrm{val}^{(2)}_{l''}+x_{l''}$: this is (a), and (d) follows.

(c) $\mathfrak p_{l''}=24\lambda_{l''}v_{l''}(j_{l''})\le24\cdot2^{G(l)}
\varepsilon_x/4$, and $\|X_{\rm tu}\|^2\le\sum\mathfrak p^2+\sum\mathfrak n^2
\le C_f\mathrm{Des}\,\varepsilon_x^2$.  (b) For $k\le l$, $k\notin L_0$,
$u_k$ vanishes at all pull and bank coordinates, so $u_k(z^\#-z^2)=0$, and
Lemma~\ref{lem:II-bankformula}(i) with base $A^2$ and $J$ the set of all pull
and bank coordinates gives $|u_k(\zh^\#)-u_k(\zh^{(2)})|\le\|U\|
\|X_{\rm tu}\|^2/(2\nu^2)$.
\end{proof}

The contraction constant of the fixed point contains the bank efficiency
$(\mu^\star_{j'}v_{l''}(j'))^{-2}$, not only $(\mu^\star_{j'})^{-2}
v_{l''}(j')^{-1}$; both are design numbers of level $l$
(Lemma~\ref{lem:II-scales}(d)), and this is the only place where the base
enters the constants of the tuning.  The pulls lower and the banks raise
the tuned values; since every pull lowers its value by
$2v_{l''}(j_{l''})\ge2\varepsilon_x$, the required bank increments are
positive whatever the signs of the $x_{l''}$: the tuning is two-sided.

\subsection{Donor peaks, bank donors and the exactifying companion}
\label{subsec:II-donors}\label{subsec:II-companion}

\begin{definition}[The exactifying companion $f^\#_w$]\label{def:II-companionw}
Let $w$ be clean for $f$, of level $l\ge l_f$, and $\Lambda:=\Lambda_w=
T_{\rm lo}(w)^3$.  The following moves are made in the order listed.
\begin{enumerate}
\item[(C1)] \emph{Close class-G rooms}: $z^1_s:=\epsilon_{l''}$ for $s\in
S^{\rm nat}_{l''}(l)$, $l''\in G(w)$.
\item[(C2)] \emph{Close tiny target rooms}: $z^1_j:=\sgn z_j$ for $j\in
T(l)\setminus F$ with $0<1-|z_j|\le b_w$; $z^1:=z$ at all other coordinates.
\item[(C3)] \emph{Donor raise.}  Let $I_D(w)$ be the set of blocks
containing a (K4) carrier.  For $m\in I_D(w)$ let $c_m$ be the least
carrier given by Lemma~\textup{\ref{lem:II-donorpeak}} for block $m$ (the
\emph{donor}), $\mathrm{vs}_m:=\sgn w_m(k(c_m))=\sgn u_{c_m}(\zh)$, and
$s_m:=s'_{c_m}(l)=\min(S_{c_m}\cap(s_{\max}(l),\infty))$.
 \begin{enumerate}
 \item[(a)] If $\lambda_{c_m}v_{c_m}(s_m)(1-\mathrm{vs}_mz^1_{s_m})\ge
 \Lambda$, put $z^2_{s_m}:=z^1_{s_m}+\eta_m\mathrm{vs}_m$ with
 $\eta_m:=\Lambda/(\lambda_{c_m}v_{c_m}(s_m))$ (a $z$-move, no mass).
 \item[(b)] Otherwise put $z^2_{s_m}:=\mathrm{vs}_m$ and place at $s_m$ a
 bank of mass $\mathfrak d_m:=2\nu\Lambda/(\lambda_{c_m}(\mu^\star_{s_m})^2
 v_{c_m}(s_m))$ with sign $\mathrm{vs}_m$.
 \end{enumerate}
 Put $A^2:=a+\sum_{m\text{ in case (b)}}\mathfrak d_m\mathrm{vs}_me^*_{s_m}$,
 $z^2:=z^1$ at all other coordinates, and $f^{(2)}:=f[A^2,z^2]$.
\item[(C4)] \emph{Exact tuning of tiny ray components.}  Let
$\Sigma_{\rm t}(w)$ be the set of tiny components $(r,m)$ of $\kappa(w)$.  If
$\Sigma_{\rm t}(w)=\emptyset$, put $f^\#_w:=f^{(2)}$.  Otherwise let
$L_0:=\bigcup_{(r,m)\in\Sigma_{\rm t}(w)}\{l''\in\supp r:m(l'')=m\}$,
$R:=R_{\Sigma_{\rm t}(w)}$ (Definition~\ref{def:II-constants}(b)),
$V^{(2)}:=(D_{r,m}(f^{(2)}))_{(r,m)\in\Sigma_{\rm t}(w)}$ and $x:=-R^+
V^{(2)}$.  Then $x$ is supported in $L_0$ (it lies in the row space of
$R$), and $Rx=-V^{(2)}$, because $V^{(2)}=R\,\mathrm{val}^{(2)}$ with
$\mathrm{val}^{(2)}:=(\epsilon_{l''}u_{l''}(\zh^{(2)}))_{l''\le l}$ lies in
the range of $R$.  Let $f^\#_w$ be the row given by
Lemma~\textup{\ref{lem:II-TU}} at $(A^2,z^2)$ with $L_0$, the signs
$\epsilon_{l''}$, and $x|_{L_0}$.
\end{enumerate}
We write $f^\#=f^\#_w=f[A^\#,z^\#]$, $a^\#:=A^\#/q^*(A^\#)$, $F^\#:=
\supp a^\#=F\cup\mathrm{Ba}\cup\mathrm{Pu}$, where $\mathrm{Ba}$ is the set
of bank coordinates (the $s_m$ of case (b) and the $j'_{l''}$, $l''\in L_0$)
and $\mathrm{Pu}$ the set of pull coordinates $j_{l''}$, $l''\in L_0$.
Forced data of $f^\#$ carry the superscript $\#$, those of $f^{(2)}$ the
superscript $(2)$.
\end{definition}

The moves act on pairwise disjoint sets of coordinates, except that (C3)
may act after (C1) on the coordinate $s_m$ when $c_m\in G(w)$: (C1) acts on
the sets $S^{\rm nat}_{l''}(l)$, $l''\in G(w)$, which are disjoint and avoid
$F\cup T(l)$; (C2) on $T(l)$; (C3) on $s_m\in S_{c_m}$ beyond $s_{\max}(l)$;
(C4) on $S_{l''}\setminus[1,s_{\max}(l)]$ with $l''\in L_0\subseteq
\mathrm{Kp}(w)$, while $c_m\notin\mathrm{Kp}(w)$ by
Lemma~\ref{lem:II-donorpeak}.  In case (a) of (C3),
$\mathrm{vs}_mz^1_{s_m}\le\mathrm{vs}_mz^2_{s_m}=\mathrm{vs}_mz^1_{s_m}+\eta_m
\le1$, so $|z^2|\le1$; all moves are off $F$, and $(A^2,z^2)$ satisfies
$z^2=\sgn A^2$ on $\supp A^2$.  The companion depends on $f$ and $w$ only,
not on a mate or a scale.

\begin{lemma}[Donor raise]\label{lem:II-donorraise}
Let $l\ge l_f$ and $m\in I_D(w)$, $c:=c_m$.
\begin{enumerate}
\item[(i)] $\mathfrak d_m\le\mathrm{Des}^{1/3}\Lambda_w/8$.
\item[(ii)] $\mathrm{vs}_mu_c(\zh^{(2)})-\mathrm{vs}_mu_c(\zh)\in[\Lambda_w/
\lambda_c-2b_w-C_f\mathrm{Des}^2\Lambda_w^2,\,3\Lambda_w/\lambda_c+2b_w+C_f
\mathrm{Des}^2\Lambda_w^2]$.
\item[(iii)] For every carrier $k\le l$ which is not a donor, (C3) changes
$u_k(\zh)$ by at most $C_f\mathrm{Des}^2\Lambda_w^2$, and by $0$ if no bank
is placed.
\end{enumerate}
\end{lemma}

\begin{proof}
The only carrier $\le l$ whose vector meets $s_m$ is $c$, with
$u_c(s_m)=v_c(s_m)$.  (i) $s_m\le s_{\rm t}(l)$, so
$(\mu^\star_{s_m})^2v_c(s_m)\ge\mathrm{Des}^{-1/6}$ and $1/\lambda_c\le
D_{\rm r}\le\mathrm{Des}^{1/6}$ (Lemma~\ref{lem:II-scales}); $\nu\le
\|U\|=\frac1{16}$.  (iii) In case (a) only $z_{s_m}$ moves and $u_k(s_m)=0$.
Banks: Lemma~\ref{lem:II-bankformula}(i) with base $(a,z^1)$ ($\nu_0=\nu$)
and $J$ the set of donor bank coordinates, $\|X\|^2\le\sum_m\mathfrak d_m^2
\le N\mathrm{Des}^{2/3}\Lambda^2$.  (ii) (C1) changes $u_c(\zh)$ by at most
$r^{\rm nat}_c\le b_w$ (only if $c\in G(w)$), and (C2) by at most $b_w\|u_c\|_1
\le b_w$.  In case (a), (C3) raises $\mathrm{vs}_mu_c$ by $v_c(s_m)\eta_m=
\Lambda/\lambda_c$ exactly, plus second-order effects of the banks of other
blocks (bounded as in (iii)).  In case (b) the closing raises
$\mathrm{vs}_mu_c$ by $v_c(s_m)(1-\mathrm{vs}_mz^1_{s_m})\in[0,\Lambda/
\lambda_c)$, and by Lemma~\ref{lem:II-bankformula}(ii) the banks raise it by
$\mathfrak d_m(\mu^\star_{s_m})^2v_c(s_m)/\nu+R=2\Lambda/\lambda_c+R$ with
$|R|\le C_f\|X\|^2\le C_f\mathrm{Des}^2\Lambda^2$.
\end{proof}

\begin{lemma}[Cost and data of the companion]\label{lem:II-compcost}
For $l\ge l_f$ and every sub-window $w$ of level $l$ that is clean for $f$:
\begin{enumerate}
\item[(a)] in (C4) the hypotheses of Lemma~\textup{\ref{lem:II-TU}} hold, and
$\varepsilon_x\le C_f\mathrm{Des}^{1/6}\beta_w\le c_T$;
\item[(b)] for every $l''\in\mathcal C(l)$ that is not a donor,
\[
 u_{l''}(\zh^\#)-u_{l''}(\zh)=\Delta^{\rm C1}_{l''}+\Delta^{\rm C2}_{l''}+
 x_{l''}1[l''\in L_0]+\Delta^{\rm H}_{l''},
\]
where $\Delta^{\rm C1}_{l''}:=\sum_{s\in S^{\rm nat}_{l''}(l)}v_{l''}(s)
(\epsilon_{l''}-z_s)$ if $l''\in G(w)$ and $:=0$ otherwise
($|\Delta^{\rm C1}_{l''}|=r^{\rm nat}_{l''}\le b_w$),
$|\Delta^{\rm C2}_{l''}|\le b_w$, and $|\Delta^{\rm H}_{l''}|\le C_f(
\mathrm{Des}^2\Lambda_w^2+\mathrm{Des}^{4/3}\beta_w^2)$; in particular
$|u_{l''}(\zh^\#)-u_{l''}(\zh)|\le C_f\mathrm{Des}^{1/6}\beta_w$;
\item[(c)] for every donor, $\mathrm{vs}_mu_{c_m}(\zh^\#)-\mathrm{vs}_m
u_{c_m}(\zh)\in[\Lambda_w/(2\lambda_{c_m}),4\Lambda_w/\lambda_{c_m}]$;
\item[(d)] the total mass placed off $F$ is at most $C_f\mathrm{Des}^{1/3}
\Lambda_w$; $\|a^\#-a\|_1+\|e^\#-e\|_H+|\nu^\#-\nu|\le C_f\mathrm{Des}^{1/3}
\Lambda_w$, and $q^*(A^\#)\le2$;
\item[(e)] with $\delta:=\zh^\#-\zh$ and $\theta_w:=C_f\mathrm{Des}\,
T_{\rm lo}(w)\log(e/T_{\rm lo}(w))$:
$\max_m\Delta_m(\delta)\le C_f\Lambda_w$, $c(\delta)\le\theta_wT_{\rm lo}(w)^2$,
$p^*(f^\#_w-f)\le\theta_wT_{\rm lo}(w)^2$, and for every block $m$
\begin{multline*}
 |C^\#_m-C_m|+|M^\#_m-M_m|+|\sigma^\#_m-\sigma_m|+|q^\#_0-q_0|+|A^\#_m-A_m|\\
 +\|D_m(w^\#_m-w_m)\|_2+\sum_k\lambda_{k,m}|w^\#_m(k)-w_m(k)|\le\theta_w
 T_{\rm lo}(w)^2 ;
\end{multline*}
moreover $\theta_w\le C_fu_w^{20}\to0$ as $l\to\infty$.
\end{enumerate}
\end{lemma}

\begin{proof}
(a) (T1): by (C1) and \eqref{eq:II-Finside}, $z^2=\epsilon_{l''}$ on
$S_{l''}\setminus[1,s_{\max}(l)]\subseteq S^{\rm nat}_{l''}(l)$ for $l''\in
L_0\subseteq G(w)$; (C2) acts inside $[1,s_{\max}(l)]$ and (C3) on the sets
$S_{c_m}$, $c_m\notin L_0$.  (T2): $\supp A^2\subseteq F\cup\{s_m\}$.  Also
$\|A^2\|_1\le1+N\mathrm{Des}^{1/3}\Lambda\le2$ and $\nu_2\ge\nu$
(Lemma~\ref{lem:II-bankformula}).  For $l''\in L_0$ (a kept carrier, not a
donor), the value moves from $f$ to $f^{(2)}$ by at most $r^{\rm nat}_{l''}+
b_w+C_f\mathrm{Des}^2\Lambda^2$ (as in the proof of
Lemma~\ref{lem:II-donorraise}); since $\|r\|_1=1$ and every component of
$\Sigma_{\rm t}(w)$ is tiny at $f$,
$|V^{(2)}_{(r,m)}|\le b_w\Phi_{\max}(r,m)+2b_w+C_f\mathrm{Des}^2\Lambda^2\le
C_f\beta_w$.  Hence $\varepsilon_x\le\|x\|_2\le\|R^+\|_2(\#\Sigma_{\rm t})^{1/2}
\max|V^{(2)}|\le H_{\rm tune}(l)C_f\beta_w\le C_f\mathrm{Des}^{1/6}\beta_w$,
and $C_f\mathrm{Des}^{1/6}\beta_w\le c_f\mathrm{Des}^{-1}$ for $l\ge l_f$ by
Lemma~\ref{lem:II-scales}(b).

(b) $\supp u_{l''}\subseteq\supp y_{l''}\cup S_{l''}\subseteq T(l)\cup
S_{l''}$, and the sets $S^{\rm nat}_{l'}(l)$, $l'\ne l''$, avoid it.  So
(C1) changes $u_{l''}(\zh)$ only through $S^{\rm nat}_{l''}(l)$, and only if
$l''\in G(w)$; there the change is $\Delta^{\rm C1}_{l''}$, of modulus
$\sum v_{l''}(s)|\epsilon_{l''}-z_s|=r^{\rm nat}_{l''}\le b_w\mathfrak m_{l''}
\le b_w$.  (C2) changes it by at most $\sum_{j\in T(l)}|u_{l''}(j)|b_w\le
b_w$.  (C3) contributes at most $C_f\mathrm{Des}^2\Lambda^2$
(Lemma~\ref{lem:II-donorraise}(iii)), and (C4) contributes $x_{l''}$ exactly
if $l''\in L_0$ (Lemma~\ref{lem:II-TU}(a)) and at most $C_T\varepsilon_x^2\le
C_f\mathrm{Des}^{4/3}\beta_w^2$ otherwise (Lemma~\ref{lem:II-TU}(b)).

(c) Lemma~\ref{lem:II-donorraise}(ii) and Lemma~\ref{lem:II-TU}(b) ($c_m\notin
L_0$), together with $b_w+\mathrm{Des}^2\Lambda^2+\mathrm{Des}^{4/3}
\beta_w^2\ll\Lambda\le\Lambda/\lambda_{c_m}$ (Lemma~\ref{lem:II-scales}(b)).

(d) The masses are $\sum_m\mathfrak d_m\le N\mathrm{Des}^{1/3}\Lambda/8$,
$\sum\mathfrak p_{l''}+\sum\mathfrak n_{l''}\le C_f\mathrm{Des}^{4/3}
\varepsilon_x\ll\Lambda$.  With $y$ the total mass, $q^*(A^\#-a)\le
(1+\|U\|)y$, so $|q^*(A^\#)-1|\le2y$ and $\|a^\#-a\|_1\le4y$;
Lemma~\ref{lem:II-bankformula}(iii) with base $a$ gives $\|e^\#-e\|\le2
\|X\|/\nu\le y/(8\nu)$; and $\nu^\#=\|U^*A^\#\|/q^*(A^\#)$.

(e) Let $k$ be a carrier.  If $k\le l$ is not a donor, $|u_k(\delta)|\le
C_f\mathrm{Des}^{1/6}\beta_w$ by (b) (for $k\notin\mathcal C(l)$ there is
nothing to prove).  For a donor, $|u_k(\delta)|\le4\Lambda/\lambda_k$ by (c).
For $k>l$, $|u_k(\delta)|\le\|u_k\|_1\|z^\#-z\|_\infty+\|e^\#-e\|\le3$, and
$\sum_{k>l}\lambda_k\le b_w^2/3$.  Hence $\Delta_m(\delta)=\sum_k\lambda_k
|u_k(\delta)|\le C_f\mathrm{Des}^{1/6}\beta_w+4N\Lambda+b_w^2\le C_f\Lambda$,
and $\sum_k\min(\lambda_k,|u_k(\delta)|)\le l\,C_f\mathrm{Des}^{1/6}\beta_w+
4ND_{\rm r}\Lambda+b_w^2\le C_f\mathrm{Des}^{1/6}\Lambda$.  As $x\log(e/x)$ is
increasing on $(0,1]$, $c(\delta)\le C_f\mathrm{Des}^{1/6}\Lambda\log(e/
\Lambda)\le\theta_wT_{\rm lo}^2$ (note $\log(e/\Lambda)\le3\log(e/
T_{\rm lo})$).  For $l\ge l_f$, $\max_m\Delta_m(\delta)\le c_f$, and
Lemma~\ref{lem:II-cost} gives $p^*(f^\#-f)\le(1+\|U\|)(\|a^\#-a\|_1+C_f
c(\delta))\le\theta_wT_{\rm lo}^2$ and the block data bounds; the bound for
$\sum_k\lambda_{k,m}|w^\#_m(k)-w_m(k)|$ is the one obtained in the proof of
Lemma~\ref{lem:II-cost}(ii), and $|A^\#_m-A_m|\le\Delta_m(\delta)$.
Finally $T\log(e/T)\le2T^{1/2}$ for $T\in(0,1]$ and $T_{\rm lo}\le
(u_w/(4\mathrm{Des}))^{40}$ (Lemma~\ref{lem:II-scales}(b)) give $\theta_w\le
C_f\mathrm{Des}\,T_{\rm lo}^{1/2}\le C_fu_w^{20}$, and $u_w\to0$
(Lemma~\ref{lem:II-scales}(e)).
\end{proof}

The donor raise ($\Lambda_w=T_{\rm lo}^3$) dominates all other changes of
coarse values, which are of order $\beta_w\le T_{\rm lo}^4/l$ (closings and
tuning) or of second order ($\mathrm{Des}^2\Lambda_w^2$); this
\emph{threshold buffer} makes the status of every kept carrier stable.

\begin{lemma}[Status table; stability under the threshold drift]
\label{lem:II-status}
For $l\ge l_f$ and every sub-window $w$ of level $l$ that is clean for $f$,
with $f^\#=f^\#_w$, $\varrho^\#$ the relative positions at $f^\#$ and
$\bar\vartheta^\#_m:=\bar\vartheta(\hat\zeta^\#_m)$:
\begin{enumerate}
\item[(a)] for every block $m$, $E_m:=\sum_{k\ne k(c_m)}|\hat\zeta^\#_m(k)-
\hat\zeta_m(k)|\le C_f\mathrm{Des}^{1/6}\beta_w$ \textup{(}the sum runs over
all $k$ if $m\notin I_D(w)$\textup{)};
\item[(b)] $c_f\Lambda_w\le\bar\vartheta^\#_m-\bar\vartheta_m\le C_f\Lambda_w$ for $m\in
I_D(w)$, and $|\bar\vartheta^\#_m-\bar\vartheta_m|\le C_f\mathrm{Des}^{1/6}\beta_w$ for
$m\notin I_D(w)$;
\item[(c)] for every $l''\in\mathcal C(l)$: if $\varrho_{l''}\ge1+u_w$ then
$\varrho^\#_{l''}\ge1+u_w/2$; if $\varrho_{l''}\le1-u_w$ then $\varrho^\#_{l''}\le
1-u_w/2$; if $\varrho_{l''}\in[u_w,1-u_w]$ then $\varrho^\#_{l''}\in[u_w/2,
1-u_w/2]$; if $\varrho_{l''}\le b_w$ then $\varrho^\#_{l''}\le\frac12$; and every
\textup{(K4)} carrier has $\varrho^\#_{l''}\le1-c_f\Lambda_w$, so that it is a
strict non-peak of $f^\#$ with $\gap^\#\ge c_fM^\#_m\Lambda_w>0$;
\item[(d)] every kept carrier $l''$ is a strict non-peak of $f^\#$, with
$\gap^\#\ge M^\#_mu_w/2$ in cases \textup{(K1), (K2)} and $\gap^\#\ge
M^\#_m/2$ in case \textup{(K3)}; for \textup{(K1), (K2), (K4)},
$\sgn w^\#_m(k(l''))=\sgn w_m(k(l''))\ne0$; and $q^\#_{l''}:=
\epsilon_{l''}\Phi_{l''}w^\#_m(k(l''))/(mC^\#_m)=\mathrm{val}^\#_{l''}/A^\#_m$;
\item[(e)] contacts of $f$ in $T(l)\setminus F$ are contacts of $f^\#$ with
the same sign, every contact-like coordinate is a contact of $f^\#$ with
$z^\#_j=\sgn z_j$, free-robust coordinates have $z^\#_j=z_j$; for
$l''\in G(w)$, $z^\#=z^{(2)}=\epsilon_{l''}$ on $S^{\rm nat}_{l''}(l)$ except at
the pull coordinate $j_{l''}$ \textup{(}if $l''\in L_0$\textup{)} and at
$s_m$ \textup{(}if $l''=c_m$\textup{)}.
\end{enumerate}
\end{lemma}

\begin{proof}
(a) $\hat\zeta_m(k)=\lambda_ku_k(\zh)$; use Lemma~\ref{lem:II-compcost}(b)
for coarse non-donors, $\sum_{k>l}\lambda_k\cdot3\le b_w^2$ for fine
carriers, and note that block $m$ contains at most one donor.

(b) Let $m\in I_D(w)$, $\zeta:=\hat\zeta_m$, $\zeta':=\hat\zeta^\#_m$,
$c:=k(c_m)$, a peak of $\zeta$.  By Lemma~\ref{lem:II-compcost}(c),
$|\zeta'(c)|=|\zeta(c)|+s$ with the same sign and $s\in[\Lambda/2,4\Lambda]$.
By (a), $E_m\le C_f\mathrm{Des}^{1/6}\beta_w\le A_ms/(8(A_m+\bar\vartheta_m+1))$ for
$l\ge l_f$.  Lemma~\ref{lem:II-T2}(b) gives $\bar\vartheta^\#_m-\bar\vartheta_m\ge\min\{1,
A_ms/(8\phi_m(A_m+\bar\vartheta_m+1))\}\ge c_f\Lambda$, $\phi_m:=\|\Phi_m\|_2^2$, and
$\bar\vartheta^\#_m-\bar\vartheta_m\le C_1(s+E_m)\le C_f\Lambda$, where $C_1,h_0$ are the
constants of Lemma~\ref{lem:II-T2} for $\zeta$ (they depend only on $f$) and
$C_1(s+E_m)<h_0$ for $l\ge l_f$.  For $m\notin I_D(w)$,
Lemma~\ref{lem:II-T2}(a) with $\|\zeta'-\zeta\|_1=E_m$.

(c) For every $k$ with $\zeta(k)\ne0$, $\varrho_k(\zeta')=(|\zeta'(k)|/
|\zeta(k)|)\varrho_k(\zeta)\bar\vartheta_m/\bar\vartheta^\#_m$, and $\varrho_{l''}=
\varrho_{k(l'')}(\hat\zeta_m)$.  If $\varrho_{l''}\ge u_w$ and $l''$ is not a
donor, then $|\zeta(k)|=\varrho\Phi_{l''}^2\bar\vartheta_m\ge u_w\bar\vartheta_m/D_{\rm r}^2$ and
$|\zeta'(k)-\zeta(k)|\le C_f\mathrm{Des}^{1/6}\beta_w$, so the first factor is
$1+O(\mathrm{Des}^{1/2}\beta_w/u_w)$; by (b), $\bar\vartheta_m/\bar\vartheta^\#_m\in[1-C_f
\Lambda,1+C_f\mathrm{Des}^{1/6}\beta_w]$.  Since $\Lambda+\mathrm{Des}^{1/2}
\beta_w/u_w\le u_w^{2}$ for $l\ge l_f$ (Lemma~\ref{lem:II-scales}(b)),
$\varrho^\#/\varrho\in[1-u_w/4,1+u_w/4]$, which gives the three robust cases
(for $u_w\le\frac14$, $(1+u)(1-u/4)\ge1+u/2$ and $(1-u)(1+u/4)\le1-u/2$).  For
a donor, $|\zeta'(c)|\ge|\zeta(c)|$ and $\bar\vartheta_m/\bar\vartheta^\#_m\ge1-C_f\Lambda$,
so $\varrho^\#\ge(1+u_w)(1-C_f\Lambda)\ge1+u_w/2$.  If $\varrho_{l''}\le b_w$,
$\varrho^\#_{l''}=m|u_{l''}(\zh^\#)|/(\Phi_{l''}\bar\vartheta^\#_m)\le\varrho_{l''}\bar\vartheta_m/
\bar\vartheta^\#_m+C_fD_{\rm r}\mathrm{Des}^{1/6}\beta_w\le\frac12$.  For a (K4)
carrier, $m\in I_D(w)$, $\varrho\le1+b_w$, the first factor is at most $1+C_f
\mathrm{Des}^{1/2}\beta_w$ and $\bar\vartheta_m/\bar\vartheta^\#_m\le1/(1+c_f\Lambda/
\bar\vartheta_m)\le1-c'_f\Lambda$; as $b_w+C_f\mathrm{Des}^{1/2}\beta_w\le c'_f
\Lambda/4$ for $l\ge l_f$, $\varrho^\#\le1-c'_f\Lambda/2$.  Finally $\gap^\#=
M^\#_m(1-\varrho^\#)$ at strict non-peaks.

(d) Kept carriers are of types (K1)--(K4) and are strict non-peaks of
$f^\#$ by (c).  At a strict non-peak $k$ of $f^\#$, [I,~Lemma~1.7] gives
$w^\#_m(k)=C^\#_m\hat\zeta^\#_m(k)/(\Phi_m(k)^2A^\#_m)$, so $\sgn w^\#_m(k)=
\sgn u_{l''}(\zh^\#)$ and $q^\#_{l''}=\epsilon_{l''}\lambda_{l''}u_{l''}(\zh^\#)
/(m\Phi_{l''}A^\#_m)=\mathrm{val}^\#_{l''}/A^\#_m$.  For (K1), (K2), (K4),
$|u_{l''}(\zh)|=\varrho\Phi_{l''}\bar\vartheta_m/m\ge u_w\bar\vartheta_m/(mD_{\rm r})$ exceeds
the change $C_f\mathrm{Des}^{1/6}\beta_w$, so the sign is that of
$u_{l''}(\zh)$, i.e.\ of $w_m(k)$.

(e) (C2) closes in the direction of $\sgn z_j$, and no other move acts on
$T(l)$; (C1) sets $z^1=\epsilon_{l''}$ on $S^{\rm nat}_{l''}(l)$, and later
moves act only on $s_m$ and on pull coordinates (banks do not change $z$).
\end{proof}

\begin{lemma}[Exactness, no slaving, robust rates]\label{lem:II-exactrobust}
For $l\ge l_f$ and every sub-window $w$ of level $l$ that is clean for $f$:
\begin{enumerate}
\item[(a)] \textup{(No slaving.)}  The change of the value of a coarse
non-donor carrier from $f$ to $f^\#_w$ is the sum of its own closing, the
target closings, its own tuning and a second-order remainder
(Lemma~\ref{lem:II-compcost}(b)); it does not depend on the closing of the
room of any other carrier.
\item[(b)] Every class-G room on $S^{\rm nat}_{l''}(l)$ is exactly closed at
$f^\#_w$ off the pull and donor coordinates, every contact-like target
coordinate is an exact contact, and the zero-cost cone of the pattern read
at $f^\#_w$ is $C(\kappa(w))$ \textup{(}the types on $T(l)\setminus F$ are
the same, Lemma~\ref{lem:II-status}(e)\textup{)}.
\item[(c)] For every tiny component $(r,m)$ of $\kappa(w)$,
$D_{r,m}(f^\#_w)=0$ exactly.  For every robust component,
$|D_{r,m}(f^\#_w)|\ge u_w\Phi_{\max}(r,m)/2$.
\item[(d)] Free-robust target coordinates keep $z^\#_j=z_j$ (room $\ge u_w$),
and robust relative positions stay robust in the sense of
Lemma~\ref{lem:II-status}(c).
\end{enumerate}
\end{lemma}

\begin{proof}
(a), (b), (d): Lemma~\ref{lem:II-compcost}(b) and
Lemma~\ref{lem:II-status}(c),(e).  (c) If $(r,m)\in\Sigma_{\rm t}(w)$, every
$l''\in\supp r$ with $m(l'')=m$ lies in $L_0$, so by
Lemma~\ref{lem:II-TU}(a), $D_{r,m}(f^\#)=\sum r(l'')(\mathrm{val}^{(2)}_{l''}+
x_{l''})=V^{(2)}_{(r,m)}+(Rx)_{(r,m)}=0$.  For a robust component,
$|D_{r,m}(f^\#)-D_{r,m}(f)|\le\max_{l''\in\supp r}|\mathrm{val}^\#_{l''}-
\mathrm{val}_{l''}|\le C_f\mathrm{Des}^{1/6}\beta_w$ (donors are not kept,
hence not in $\supp r$), while $u_w\Phi_{\max}(r,m)\ge u_w/D_{\rm r}\ge2C_f
\mathrm{Des}^{1/6}\beta_w$ for $l\ge l_f$.
\end{proof}

Thus closing one object never reopens or moves another exactified object,
the exactified set is closed under the moves, and no robust object moves by
more than $C_f\mathrm{Des}^{1/6}\beta_w$.  We do not claim that the room of
a donor stays robust at $f^\#_w$ (in case (b) of (C3) the closing at $s_m$
may remove most of it); it is not used: donors are pinned at $f$
(Lemma~\ref{lem:II-pinned}), are not kept, and the data constructed below
vanish at $s_m$.

\begin{lemma}[Inward push of near-threshold coordinates]\label{lem:II-IP}
Fix one block and drop its index.  Let $\zeta\in\ell_1\setminus\{0\}$,
$\bar\vartheta:=\bar\vartheta(\zeta)$, $\mathsf a:=|\zeta|$, $0<b\le\frac12$,
let $\mathcal N$ be a nonempty set of coordinates with $|\mathsf v_k(\zeta)-
\bar\vartheta|\le\bar\vartheta b$ for $k\in\mathcal N$, and let $c\ge4b(\mathsf
a+\bar\vartheta)$.  Let $\zeta'\in\ell_1$ satisfy, for $k\in\mathcal N$,
$\sgn\zeta'(k)=\sgn\zeta(k)$ and $|\zeta'(k)|=|\zeta(k)|-c_k\Phi(k)^2$ with
$c_k\in[c,\bar\vartheta/2]$, and $E:=\sum_{k\notin\mathcal N}|\zeta'(k)-\zeta(k)|
\le\bar\vartheta c\sum_{k\in\mathcal N}\Phi(k)^2/(8(\mathsf a+\bar\vartheta))$.
Then $\bar\vartheta(\zeta')>\bar\vartheta$, and $\mathsf v_k(\zeta')\le
\bar\vartheta(\zeta')-(c-\bar\vartheta b)$ for every $k\in\mathcal N$: every
pushed coordinate is a strict non-peak of $\zeta'$.
\end{lemma}

\begin{proof}
Write $\mathsf A'=\mathsf A(\cdot;\zeta')$, $\mathsf B'=\mathsf B(\cdot;\zeta')$,
$\Phi_{\mathcal N}^2:=\sum_{k\in\mathcal N}\Phi(k)^2$.  For $k\in\mathcal N$,
$|\zeta'(k)|-\bar\vartheta\Phi(k)^2=\Phi(k)^2(\mathsf v_k-c_k-\bar\vartheta)<0$,
so these terms of $\mathsf A'(\bar\vartheta)$ vanish, while their terms in
$\mathsf A(\bar\vartheta;\zeta)$ are at most $\bar\vartheta b\Phi(k)^2$; the
other terms change by at most $|\zeta'(k)-\zeta(k)|$.  With
Lemma~\ref{lem:II-T}(b), $\mathsf A'(\bar\vartheta)\ge\mathsf a-y$,
$y:=\bar\vartheta b\Phi_{\mathcal N}^2+E$.  In $\mathsf B'(\bar\vartheta)$ the
terms with $k\notin\mathcal N$ change by at most $2\bar\vartheta|\zeta'(k)-
\zeta(k)|$; for $k\in\mathcal N$, $\min(\bar\vartheta,\mathsf v_k)\ge
\bar\vartheta(1-b)$ and $0\le\mathsf v'_k=\mathsf v_k-c_k\le\bar\vartheta(1+b)-
c_k$, so the term drops by at least $\Phi(k)^2[\bar\vartheta^2(1-b)^2-
(\bar\vartheta(1+b)-c_k)^2]=\Phi(k)^2(c_k-2\bar\vartheta b)(2\bar\vartheta-c_k)
\ge\Phi(k)^2\bar\vartheta(c-2\bar\vartheta b)$.  Hence $\mathsf B'(\bar\vartheta)
\le\mathsf a^2-\bar\vartheta(c-2\bar\vartheta b)\Phi_{\mathcal N}^2+2\bar\vartheta
E$, and, since $\mathsf A'(\bar\vartheta)^2\ge\mathsf a^2-2\mathsf ay$ (trivial
if $\mathsf a<y$),
\[
 \Xi(\bar\vartheta;\zeta')\ge\bar\vartheta\Phi_{\mathcal N}^2\big(c-2b(\mathsf a+
 \bar\vartheta)\big)-2(\mathsf a+\bar\vartheta)E\ge\tfrac12\bar\vartheta c
 \Phi_{\mathcal N}^2-\tfrac14\bar\vartheta c\Phi_{\mathcal N}^2>0 .
\]
By Lemma~\ref{lem:II-T}(d), $\bar\vartheta(\zeta')>\bar\vartheta$.  Finally
$\mathsf v_k(\zeta')=\mathsf v_k-c_k\le\bar\vartheta(1+b)-c<\bar\vartheta(\zeta')-
(c-\bar\vartheta b)$.
\end{proof}

Lemma~\ref{lem:II-IP} is not used in the construction of $f^\#_w$: a push
that dominates the threshold drift would move the $d$-components of the
rays through the pushed carriers by amounts between $b_w$ and $u_w$ and
destroy the tiny/robust dichotomy on which (C4) rests, whereas the donor
raise rescales the $d$-coefficients of untouched carriers of a block by one
common factor (Corollary~\ref{cor:II-T3}).  Per carrier, a far pull at
$j\in S_{l''}$ realizes a push of $\zeta(k(l''))$ by
$2\lambda_{l''}v_{l''}(j)$, and by bounded gaps every size between $c$ and
$2^{G(l)}c$ is available; Lemma~\ref{lem:II-IP} is the appropriate tool when
no tiny $d$-component passes through the pushed carriers.

\subsection{Transplant to exact $d$-neutral two-piece data}
\label{subsec:II-transplant}

\begin{lemma}[Split with respect to an arbitrary sign vector]
\label{lem:II-splitz}
Let $(B_\pm,\Theta_\pm)$ be a two-sided decomposition of $g\in\Cset(f)$ at
some scale $t>0$, let $z'\in B_{\ell_\infty}$, $K':=\{j\notin F:|z'_j|=1\}$,
and suppose $\Delta B1_{F^c}=V+e$ with $e\in\ell_1$ and $V\in\ell_1$
supported in $K'$ and $z'$-signed.  Then there are $\chi:K'\to[0,1]$ and
$\mathfrak e_\pm\in\ell_1(F^c)$ with $B_+1_{F^c}=\chi V+\mathfrak e_+$,
$B_-1_{F^c}=-(1-\chi)V+\mathfrak e_-$ and
\[
 \|\mathfrak e_\pm\|_1\le\|e\|_1+\sum_{j\notin F}\big(\phi_{z'_j}(B_+(j))+
 \phi_{-z'_j}(B_-(j))\big).
\]
\end{lemma}

\begin{proof}
The proof of [I,~Lemma~8.26], read with $z'$ in place of $z$, gives the
pointwise bound $|\mathfrak e_\pm(j)|\le|e(j)|+\phi_{z'_j}(B_+(j))+
\phi_{-z'_j}(B_-(j))$ for $j\notin F$; it uses only
[I,~Lemma~8.23(a),(b)].  Sum over $j\notin F$.
\end{proof}

\begin{lemma}[Ray removal]\label{lem:II-rayremoval}
Let $E$ be a finite set, $r_1,\dots,r_p\in[0,\infty)^E$ with $\|r_i\|_1=1$,
$\mu\in[0,\infty)^p$, $\tau:=\sum_i\mu_ir_i$, $Q:\R^E\to\R$ linear,
$Q_i:=Q(r_i)$ and $Q_{\min}:=\min\{|Q_i|:Q_i\ne0\}$ \textup{(}$:=\infty$ if all
$Q_i=0$\textup{)}.  Then there is $\mu'$ with $0\le\mu'_i\le\mu_i$ such that
$\tau':=\sum_i\mu'_ir_i$ satisfies $Q(\tau')=0$, $0\le\tau'\le\tau$
coordinatewise, and $\|\tau-\tau'\|_1=\sum_i(\mu_i-\mu'_i)\le|Q(\tau)|/
Q_{\min}$.
\end{lemma}

\begin{proof}
If $Q(\tau)=0$ take $\mu'=\mu$.  Let $Q(\tau)=\mathsf e>0$ (the case $<0$ is
symmetric).  Then $S_+:=\sum_{Q_i>0}\mu_iQ_i\ge\mathsf e$; put $\vartheta:=
\mathsf e/S_+\in(0,1]$, $\mu'_i:=(1-\vartheta)\mu_i$ if $Q_i>0$ and $\mu'_i:=
\mu_i$ otherwise.  Then $Q(\tau')=\mathsf e-\vartheta S_+=0$, and, since the
$r_i$ are nonnegative with $\|r_i\|_1=1$, $\|\tau-\tau'\|_1=\vartheta
\sum_{Q_i>0}\mu_i\le\vartheta S_+/Q_{\min}=\mathsf e/Q_{\min}$.
\end{proof}

In the next proposition the two-piece data of [I,~Definition~7.7] and
the coordinate kinds [1]--[3] of Subsection~\ref{subsec:II-U} are read at
the companion $f^\#=f^\#_w$, whose base support is $F^\#=F\cup\mathrm{Ba}
\cup\mathrm{Pu}$; $\Gamma^\#_w$ denotes $\Gamma_w$ computed at $f^\#$.

\begin{proposition}[Transplant at the exactifying companion]
\label{prop:II-transplant}
Let $w$ be a sub-window of level $l\ge l_f$ that is clean for $f$ and at
which \textup{(SH$_w$)} and \textup{(VR$_w$)} hold, and let $f^\#=f^\#_w$.  Let
$g\in\Cset(f)$, $\eta\le\eta_*$ with $\eta_\Gamma(\eta)\le1$, $t\in\mathcal
W(w)$ with $t\le\min(t_\eta,1)$, and let $(B_\pm,\Theta_\pm)$ be a two-sided
decomposition of $g$ at scale $t$.  Then there is $g_t\in X^*$ carrying
$d$-neutral two-piece data $(b^\pm,\omega^\pm)$ at $f^\#$ such that:
\begin{enumerate}
\item[(i)] $b^\pm(\xi^\#)=0$, $t\|b^\pm\|_1\le A_0$ with an $f$-constant $A_0$,
and $\Gamma^\#_w(b^\pm,\omega^\pm)\le\big(\sqrt{1+\eta_\Gamma(\eta)+C_f
\theta_w}+K_wt\big)^2$;
\item[(ii)] $p^*(g-g_t)\le K_wt$ with $K_w:=C_f\mathrm{Des}(l)\,u_w^{-4}$;
\item[(iii)] every $k\in\supp\omega^\pm_m$ is, at $f^\#$, of kind
\textup{[1]}, of kind \textup{[2]} with $\gamma_B:=\gamma(w):=u_w\min_mM_m/4$
and $A_2:=22$, or of kind \textup{[3]} for the side of the pair with
$A_2=22$;
\item[(iv)] $z^\#_ib^+(i)\ge0\ge z^\#_ib^-(i)$ for $i\in\mathrm{Ba}$, and
$t|b^\pm(i)|\le|a^\#(i)|$ for $i\in\mathrm{Pu}$.
\end{enumerate}
\end{proposition}

\begin{proof}
Write $u:=u_w$, $b:=b_w$, $\Lambda:=\Lambda_w$, and let $K_{\rm pin}\le
C_fl^2D_{\rm r}^3u^{-3}$ be as in Lemma~\ref{lem:II-pinned}.

\emph{Step 0 (pinning at $f$).}  By (SH$_w$) and
Lemma~\ref{lem:II-shiftpin}, $|\Delta d_m|M_m\le K'_dt$ for every $m$; by
Lemma~\ref{lem:II-pinned}, $|\Delta\theta_{l''}|\le K_{\rm pin}t$ for class R
and $|\tau_{l''}|\le K_{\rm pin}t$ on $P(w)$, and
Lemma~\ref{lem:II-diagpin} holds.

\emph{Step 1 (zero-cost projection).}  Let $\tau:=(\tau_{l''})_{l''\in G(w)}$
and $\beta_{l''}:=12\lambda_{l''}/t$.  For $j\in T(l)\setminus F$,
$\Delta B(j)=L_j(\tau)+e_0(j)$ with $L_j(\tau):=\sum_{G(w)}\epsilon_{l''}
\tau_{l''}u_{l''}(j)$ and $e_0$ as in Lemma~\ref{lem:II-diagpin}(b).  We
estimate the violations of the conditions defining $Z_{\kappa(w)}(\beta)$
by $\tau$.  Nonnegativity on $\mathrm{Kp}(w)$: $\sum(\tau)_-\le K_gt$.  At a
contact-like $j$, the condition is $\sgn z_j\,L_j\ge0$, and $(\sgn z_j\,
L_j(\tau))_-\le(\sgn z_j\,\Delta B(j))_-+|e_0(j)|\le\phi_{z_j}(\Delta B(j))+
|e_0(j)|$ (if $\sgn z_j\,x<0$, then $\phi_{z_j}(x)=(1+|z_j|)|x|$).  At a
free-robust $j$ the condition is $L_j=0$, and $|L_j(\tau)|\le|\Delta B(j)|+
|e_0(j)|\le\phi_{z_j}(\Delta B(j))/u+|e_0(j)|$ by [I,~Lemma~8.23(a)].
Summing with [I,~Lemma~8.24], these violations are at most $(t/q_0)(1+1/u)
+2\|e_0\|_1$.  The equations $\tau=0$ on $P(w)$ are violated by at most
$lK_{\rm pin}t$, and the box conditions not at all ($|\tau|\le6\lambda/t$).
So the total violation is $V_w\le C_flK_{\rm pin}t$.  As $0\in
Z_{\kappa(w)}(\beta)$, Definition~\ref{def:II-pattern} gives $\tau_0\in
Z_{\kappa(w)}(\beta)$ with $\|\tau-\tau_0\|_1\le G^{**}(l)V_w$.

\emph{Step 2 ($d$-rows at $f^\#$).}  For a block $m$ and $\tau'\in\R^{G(w)}$
put
\[
 Q^\#_m(\tau'):=\sum_{l''\in\mathrm{Kp}(w),\,m(l'')=m}q^\#_{l''}\tau'_{l''}=
 \frac1{A^\#_m}\sum_{l''\in\mathrm{Kp}(w),\,m(l'')=m}\mathrm{val}^\#_{l''}
 \tau'_{l''}
\]
(Lemma~\ref{lem:II-status}(d)).  By [I,~(17)],
\[
 \sum_{l''\in G(w),\,m(l'')=m}q_{l''}\tau_{l''}=-\Delta d_mM_m+r_m+
 \frac1{mC_m}\sum\Phi w\Delta\theta,
\]
the last sum over class R and fine carriers of block $m$; hence, with
Step~0, the sums over the kept carriers of block $m$ satisfy
$|\sum q\tau|\le C_flK_{\rm pin}t$ and $|\sum q\tau_0|\le C_f(lK_{\rm pin}t+
G^{**}V_w)$.  At a
kept strict non-peak of $f$, $q=\mathrm{val}/A_m$; at a kept (K4) peak of
$f$, $q=\mathrm{val}/(\varrho A_m)$ with $\varrho\in[1,1+b]$
(Remark~\ref{rem:II-onesigned}(c)).  With Lemma~\ref{lem:II-compcost}(b),
$|q^\#_{l''}-q_{l''}A_m/A^\#_m|\le C_f\mathrm{Des}^{1/6}\beta_w$, and
$\sum\tau_0\le12\sum\lambda/t\le4/t$; since $\mathrm{Des}^{1/6}\beta_w\le t^2$
and $A_m/A^\#_m\le2$,
\[
 |Q^\#_m(\tau_0)|\le K_Qt,\qquad K_Q:=C_f(lK_{\rm pin}+G^{**}(l)V_w/t+1)\le
 C_fG^{**}(l)l^3D_{\rm r}^3u^{-3}.
\]

\emph{Step 3 (exact $d$-repair by ray removal).}  $\tau_0\in C(\kappa(w))$, so
$\tau_0=\sum_{r\in\mathrm{Ext}(\kappa(w))}\mu_rr$ with $\mu_r\ge0$.  Since
$r=0$ on $P(w)$, $Q^\#_m(r)=D_{r,m}(f^\#)/A^\#_m$.  By
Lemma~\ref{lem:II-exactrobust}(c) every component of $r$ at $f^\#$ is $0$
(tiny at $f$, or block not met) or has modulus $\ge u\Phi_{\max}(r,m)/2\ge
u/(2D_{\rm r})$ (robust at $f$); by (VR$_w$), $r$ has at most one nonzero
component at $f^\#$, in a block $m(r)$ ($m(r)$ undefined if there is none).
For each block $m$ apply Lemma~\ref{lem:II-rayremoval} to the rays with
$m(r)=m$, the coefficients $\mu_r$, and $Q:=Q^\#_m$: as rays with
$m(r)\ne m$ have $Q^\#_m(r)=0$, $Q^\#_m(\sum_{m(r)=m}\mu_rr)=Q^\#_m(\tau_0)$,
and we obtain $\mu'_r\in[0,\mu_r]$ with $Q^\#_m(\sum_{m(r)=m}\mu'_rr)=0$ and
$\sum_{m(r)=m}(\mu_r-\mu'_r)\le|Q^\#_m(\tau_0)|\,2D_{\rm r}A^\#_m/u$.  Put
$\mu'_r:=\mu_r$ when $m(r)$ is undefined and $\tau':=\sum_r\mu'_rr$.  Then
$\tau'\in C(\kappa(w))$, $0\le\tau'\le\tau_0$ coordinatewise (so
$\tau'\le\beta$ on $\mathrm{Kp}(w)$ and $\tau'=0$ on $P(w)$),
$Q^\#_{m'}(\tau')=0$ for \emph{every} block $m'$, and $\|\tau_0-\tau'\|_1\le
C_fND_{\rm r}K_Qt/u$.  Put $V':=\sum_{l''\in\mathrm{Kp}(w)}\epsilon_{l''}
\tau'_{l''}u_{l''}1_{F^c}$.  Then
\begin{gather*}
 \|\Delta B1_{F^c}-V'\|_1\le\|e_0\|_1+\|\tau-\tau'\|_1\le K_Ut,\\
 K_U:=K_g+1+G^{**}V_w/t+C_fND_{\rm r}K_Q/u\le C_fG^{**}(l)l^3D_{\rm r}^4u^{-4}.
\end{gather*}
\emph{Zero cost.}  Let $z^{(2)}$ be the sign vector of $f^{(2)}$ and
$K^{(2)}:=\{j\notin F:|z^{(2)}_j|=1\}$.  We claim that $V'$ is
$z^{(2)}$-signed and supported in $K^{(2)}$.  Off $F$, $\supp u_{l''}
\subseteq T(l)\cup S_{l''}$ and $S_{l''}\setminus F\subseteq S^{\rm nat}_{l''}(l)
\cup T(l)$.  On $S^{\rm nat}_{l''}(l)$ only $u_{l''}$ is nonzero among the
vectors of coarse carriers, $V'(s)=\epsilon_{l''}\tau'_{l''}v_{l''}(s)$ with
$\tau'_{l''}\ge0$, and $z^{(2)}_s=\epsilon_{l''}$ by (C1) (for kept $l''$;
the coordinate $s_m$ lies in $S_{c_m}$ and $V'(s_m)=0$, since donors are not
kept).  On $T(l)\setminus F$, $V'(j)=L_j(\tau')$, and the rows of
$\kappa(w)$ give $\sgn z_j\,L_j(\tau')\ge0$ at contact-like $j$, where
$z^{(2)}_j=\sgn z_j$ (C2), and $L_j(\tau')=0$ at free-robust $j$.  This
proves the claim.

\emph{Step 4 (split with respect to $z^{(2)}$).}  We bound
$\Sigma:=\sum_{j\notin F}(\phi_{z^{(2)}_j}(B_+(j))+\phi_{-z^{(2)}_j}(B_-(j)))$.
Where $z^{(2)}_j=z_j$, the terms sum to at most $t/q_0$
[I,~Lemma~8.24].  Where $z^{(2)}_j=\pm1$ and $z^{(2)}_jz_j\ge0$ (the (C2)
coordinates and the (C1) coordinates with $\epsilon_{l''}z_s\ge0$),
$\phi_{z^{(2)}_j}\le2\phi_{z_j}$ and $\phi_{-z^{(2)}_j}\le2\phi_{-z_j}$
(if $z^{(2)}_jx<0$ then $\phi_{z^{(2)}_j}(x)=2|x|\le2\phi_{z_j}(x)$), so these
terms sum to at most $2t/q_0$.  At the remaining modified coordinates (the
\emph{flipped} (C1) coordinates, $\epsilon_{l''}z_s<0$, and the $s_m$),
$\phi_{\pm z^{(2)}_j}(x)\le2|x|$ and $|B_+(j)|+|B_-(j)|\le|\Delta B(j)|+
\phi_{z_j}(B_+(j))+\phi_{-z_j}(B_-(j))$ [I,~Lemma~8.23(b)].  On
$S^{\rm nat}_{l''}(l)$, $\Delta B(s)=\epsilon_{l''}\tau_{l''}v_{l''}(s)-r(s)$
(proof of Lemma~\ref{lem:II-diagpin}), and the flipped coordinates carry
$v_{l''}$-mass at most $\sum v_{l''}(s)(1-\epsilon_{l''}z_s)=r^{\rm nat}_{l''}\le
b$; so $\sum_{\rm flipped}|\Delta B|\le\sum_{l''}(6\lambda_{l''}/t)b+
\frac{8b^2}{3t}\le t^2$.  At $s_m$ only the donor $c_m$ and fine carriers
live, so $|\Delta B(s_m)|\le|\Delta\theta_{c_m}|+t^7\le K_{\rm pin}t+t^7$
(Lemma~\ref{lem:II-pinned}).  Hence $\Sigma\le C(1/q_0+NK_{\rm pin}+1)t$, and
Lemma~\ref{lem:II-splitz} with $z'=z^{(2)}$, $V=V'$ and $e:=\Delta B
1_{F^c}-V'$ gives $\chi:K^{(2)}\to[0,1]$ and $\mathfrak e_\pm$ with
\begin{gather*}
 B_+1_{F^c}=\chi V'+\mathfrak e_+,\qquad B_-1_{F^c}=-(1-\chi)V'+\mathfrak e_-,\\
 \|\mathfrak e_\pm\|_1\le K_St,\qquad K_S:=K_U+C(1/q_0+NK_{\rm pin}+1) .
\end{gather*}

\emph{Step 5 (data at $f^\#$).}  For a block $m$ and a coordinate $k$
define $\omega^+_m(k)$ as follows: $\omega^+_m(k):=\omega_{+,m}(k)$ if
$k=k(l'')$ with $l''\in\mathrm{Kp}(w)$; $\omega^+_m(k):=0$ if $k$ is fine, or
a peak of $f^\#$, or a peak of $f$ and a strict non-peak of $f^\#$;
otherwise ($k$ coarse, not kept, a strict non-peak of $f$ and of $f^\#$)
$\omega^+_m(k):=\max\{-2\gap^\#_m(k)/t,\min\{\omega_{+,m}(k),2\gap^\#_m(k)/t\}\}$
if $\gap^\#_m(k)\ge t^2$ and $:=0$ if $\gap^\#_m(k)<t^2$.  Put
\begin{gather*}
 \omega^-_m:=\omega^+_m+\sum_{l''\in\mathrm{Kp}(w),\,m(l'')=m}
 (\epsilon_{l''}\tau'_{l''}/\lambda_{l''})e_{k(l'')},\\
 b^+:=B_+1_F+\chi V'-\varkappa\,a/a(\zh^\#),\qquad
 \varkappa:=(B_+1_F+\chi V')(\zh^\#),\\
 b^-:=b^+-\sum_{l''\in\mathrm{Kp}(w)}\epsilon_{l''}\tau'_{l''}u_{l''},\qquad
 g_t:=b^++\sum_mR_m^*\big(\omega^+_m-d^\#_m(\omega^+_m)w^\#_m\big).
\end{gather*}
Here $a(\zh^\#)=\|a\|_1+\ip{U^*a}{e^\#}$, and by
\eqref{eq:II-bankformula} with base $a$, $\ip{U^*a}{e^\#}=\nu^2/(\nu^2+
\|X\|^2)^{1/2}\in[\nu-\|X\|^2/(2\nu),\nu]$, so $a(\zh^\#)\in[\frac12,1]$.
Finally, at every kept (K4) coordinate $k$ (with $\varsigma_k:=\sgn
w_m(k)=\sgn w^\#_m(k)$) apply the \emph{shift}: replace $\omega^\pm_m(k)$ by
$\omega^\pm_m(k)+o_k$, where $\varsigma_ko_k:=(-\frac32\gap^\#_m(k)/t-
\varsigma_k\omega^-_m(k))_+$.  The shift does not change $\omega^-_m-
\omega^+_m$.

\emph{Admissibility.}  Off $F^\#$ we have $z^\#=z^{(2)}$, $b^+=\chi V'$ and
$b^-=-(1-\chi)V'$; these are $z^\#$-signed, resp.\ $(-z^\#)$-signed, and
supported in $K^\#$ by Step~3.  The vectors $\omega^\pm_m$ are finitely
supported in the strict non-peak set of $f^\#$ (kept carriers are strict
non-peaks of $f^\#$ by Lemma~\ref{lem:II-status}(d)).  Since $R_m^*e_k=
\lambda_{k,m}u_{k,m}$, $\sum_mR_m^*(\omega^-_m-\omega^+_m)=\sum_{\mathrm{Kp}(w)}
\epsilon\tau'u=b^+-b^-$, and
\[
 d^\#_m(\omega^-_m)-d^\#_m(\omega^+_m)=\sum_{l''\in\mathrm{Kp}(w),\,m(l'')=m}
 \frac{\epsilon_{l''}\tau'_{l''}}{\lambda_{l''}}\,\frac{\Phi_{l''}^2
 w^\#_m(k(l''))}{C^\#_m}=Q^\#_m(\tau')=0 .
\]
So $(b^-,\omega^-)$ also represents $g_t$, and the data are $d$-neutral
two-piece data at $f^\#$.  Since $b^+(\zh^\#)=\varkappa-\varkappa=0$,
$b^+(\xi^\#)=0$, and $b^-(\xi^\#)=g_t(\xi^\#)=b^+(\xi^\#)=0$ by
[I,~Lemma~3.2] at $f^\#$.

\emph{(iv).}  At a pull coordinate $j=j_{l''}$, $V'(j)=\epsilon_{l''}
\tau'_{l''}v_{l''}(j)$, so $t|b^\pm(j)|\le t\tau'_{l''}v_{l''}(j)\le12
\lambda_{l''}v_{l''}(j)=\mathfrak p_{l''}/2\le|a^\#(j)|$, because $q^*(A^\#)
\le2$.  At a tuning bank $j'_{l''}$, $z^{(2)}_{j'}=z^\#_{j'}=\epsilon_{l''}$,
so $b^+(j')=\chi V'(j')$ and $b^-(j')=-(1-\chi)V'(j')$ have the required
signs; at a donor bank $s_m$, $V'(s_m)=0$, so $b^\pm(s_m)=0$.

\emph{Estimates.}  First, $B_+-b^+=\mathfrak e_++\varkappa a/a(\zh^\#)$, and
$\varkappa=(B_+-\mathfrak e_+)(\zh^\#)$ with
\[
 B_+(\zh^\#)=B_+(\zh)+\sum_{j\notin F}B_+(j)(z^\#_j-z_j)+\ip{U^*B_+}{e^\#-e}.
\]
Here $|B_+(\zh)|\le t/(2q_0)$ [I,~Lemma~8.7(a)]; at the (C2) coordinates and
the (C1) coordinates with $\epsilon z_s\ge0$, $|z^\#_j-z_j|=1-|z_j|$ and
$(1-|z_j|)|B_+(j)|\le\phi_{z_j}(B_+(j))$, in total $\le t/(2q_0)$; at
flipped coordinates, at the $s_m$ and at pulls, $|z^\#_j-z_j|\le2$ and
$|B_+(j)|\le|\Delta B(j)|+\phi_{z_j}(B_+(j))+\phi_{-z_j}(B_-(j))$, where
$\sum_{\rm pulls}|\Delta B(j)|\le\sum_{l''}6\lambda_{l''}v_{l''}(j_{l''})/t+t^7
\le6\cdot2^{G(l)}\varepsilon_x/t+t^7\le t^2$; and $|\ip{U^*B_+}{e^\#-e}|\le
\|U\|\|B_+\|_1\|e^\#-e\|\le C_f\mathrm{Des}^{1/3}\Lambda/t\le t$, using
$t\|B_+\|_1\le\eta\le1$ [I,~Lemma~8.6] and Lemma~\ref{lem:II-compcost}(d).
With $\|\zh^\#\|_\infty\le1+\|U\|$, $|\varkappa|\le C(K_S+NK_{\rm pin}+1/q_0
+1)t$, and $\|B_+-b^+\|_1\le K_St+2|\varkappa|$.

Next, for the blocks put $X_m:=\Theta_{+,m}-\omega^+_m+d^\#_m(\omega^+_m)
w^\#_m$, so that $g-g_t=(B_+-b^+)+\sum_mR_m^*X_m$ and $\|R_m^*X_m\|_1\le
\sum_k\lambda_{k,m}|X_m(k)|$.  Drop $m$; since $\Theta_+=\omega_+-d_+w$,
\[
 X=(\omega_+-\omega^+)+(d^\#(\omega^+)-d_+)w+d^\#(\omega^+)(w^\#-w).
\]
Let $\delta^\omega_k:=|\omega_+(k)-\omega^+(k)|$.  Recall from the claim in the
proof of [I,~Proposition~8.15(c)] (which uses only [I,~Lemma~8.8]):
at a peak $k$ of $f$, $\lambda_k|\omega_+(k)|\le\lambda_k(|\omega_+(k)|+
|\omega_-(k)|)\le|\Delta\theta_k|+\lambda_k|\Delta d|M$
([I,~(16)]); at a strict non-peak $k$ of $f$, $|\omega_+(k)|\le2\gap(k)/t+
|\Delta\theta_k|/\lambda_k+|\Delta d|M$, and the clamp of $\omega_+(k)$ with
radius $2\gap(k)/t$ (and with $0$ if $\gap(k)<t^2$) differs from
$\omega_+(k)$ by at most $2t+|\Delta\theta_k|/\lambda_k+|\Delta d|M$.
Replacing $\gap$ by $\gap^\#$ in the clamp changes it by at most $2t+2|\gap^\#(k)
-\gap(k)|/t$.  Coordinates that change status between $f$ and $f^\#$ and are
not kept satisfy $|\varrho-1|\le b$ (Lemma~\ref{lem:II-status}(c)); there
$\omega^+=0$ and the bounds above apply with $\gap\le bM\le b$.  At kept
coordinates $\delta^\omega_k=|o_k|$, and $\lambda_k|o_k|\le|\tau_{l''}-
\tau'_{l''}|+\lambda_k|\Delta d|M$: indeed, by [I,~Lemma~8.8(c),(f)],
$\varsigma_k\omega_-(k)\ge-\frac32\gap(k)/t\ge-\frac32\gap^\#(k)/t$ (as
$\gap(k)\le bM\le\gap^\#(k)$ for (K4)), and $\varsigma_k\omega^-(k)=
\varsigma_k\omega_+(k)+\tau'_{l''}/\lambda_k=\varsigma_k\omega_-(k)+(\tau'_{l''}-
\tau_{l''})/\lambda_k+\Delta d|w(k)|$ (from $\omega_--\omega_+=-\Delta\Theta-
\Delta d\,w$ and $\varsigma_k\epsilon_{l''}=1$).  Summing, with
$\sum\lambda\le1$, Step~0 and Steps~1--3,
\begin{multline*}
 \sum_{\rm coarse}\lambda_k\delta^\omega_k\le K_{\rm pin}t+lK_{\rm pin}t+
 \|\tau-\tau'\|_1+K'_dt+4t+\frac{2b}t\\
 +\frac2t\sum_k\lambda_k|\gap^\#(k)-\gap(k)|\le C_f(K_U+lK_{\rm pin})t,
\end{multline*}
because $\sum_k\lambda_k|\gap^\#-\gap|\le|M^\#-M|+\sum_k\lambda_k|w^\#(k)-w(k)|
\le2\theta_wT_{\rm lo}^2\le2\theta_wt^2$ (Lemma~\ref{lem:II-compcost}(e)).
At fine coordinates $\omega^+=0$ and $\sum_{\rm fine}\lambda_k|\omega_+(k)|
\le(4/t)\sum_{l'>l}\lambda_{l'}\le t^7$ ([I,~Lemma~8.8(a),(e)]).  Next
$|d^\#(\omega^+)-d_+|\le|d^\#(\omega^+)-d(\omega^+)|+\|D(\omega^+-\omega_+)
\|_2+|d(\omega_+)-d_+|$, where $|d(\omega_+)-d_+|\le t/\sigma$
[I,~Lemma~8.8(d)], $\|D(\omega^+-\omega_+)\|_2\le\sum_k\lambda_k\delta^\omega_k+t^7$,
and $|d^\#(\omega)-d(\omega)|\le C_f\|D\omega\|_2(\|D(w^\#-w)\|_2+|C^\#-C|)\le
C_f\theta_wt$ for $\omega=\omega^+$, as $\|D\omega^+\|_2\le\|\Phi\|_2
\|\omega^+\|_\infty\le11/t$ (by (iii) below).  Finally $|d^\#(\omega^+)|\le11/t$
and $\sum_k\lambda_k|w^\#(k)-w(k)|\le\theta_wt^2$.  Altogether
\[
 p^*(g-g_t)\le(1+\|U\|)\|g-g_t\|_1\le C_f(K_U+lK_{\rm pin}+1)t\le K_wt,
\]
since $K_U+lK_{\rm pin}+1\le C_fG^{**}(l)l^3D_{\rm r}^4u^{-4}\le C_f
(lD_{\rm r}G^{**}(l))^4u^{-4}\le C_f\mathrm{Des}(l)^{2/3}u^{-4}$ ($l$,
$D_{\rm r}$ and $G^{**}(l)$ are distinct factors of the bracket of
$\mathrm{Des}(l)$).  This is (ii).

\emph{(i).}  $\|V'\|_1\le\sum\tau'\le4/t$ and $t\|B_+\|_1\le1$, so
$t\|b^+\|_1\le1+4+2t|\varkappa|$ and $t\|b^-\|_1\le t\|b^+\|_1+4$: an
$f$-constant $A_0$ bounds both.  $\sqrt{\Gamma^\#_w}$ is a seminorm, and
$H^\#_m$ vanishes on multiples of $w^\#_m$, so
$\sqrt{\Gamma^\#_w(b^+,\omega^+)}\le\sqrt{\Gamma^\#_w(B_+,\Theta_+)}+
\sqrt{\Gamma^\#_w(B_+-b^+,X)}$.  The second term is at most $C_f(\|B_+-
b^+\|_1+\sum_m\sum_k\lambda_k|X_m(k)|)\le K_wt$, using $q^\#_0h^\#(y)\le
\|U\|^2\|y\|_1^2/\nu^\#$ and $\sigma^\#_mH^\#_m(X)\le(\sum_k\lambda_k|X(k)|)^2/
C^\#_m$.  For the first term, $q_0^\#h^\#(B_+)=q^\#_0\|P^\perp_{e^\#}U^*B_+\|^2/
\nu^\#$ with $\|P^\perp_{e^\#}U^*B_+\|\le\|P^\perp_eU^*B_+\|+2\|e^\#-e\|\|U\|
\|B_+\|_1$, $\|P^\perp_eU^*B_+\|^2=\nu h(B_+)\le2\nu/q_0$
[I,~Lemma~8.7(d)], $\|e^\#-e\|\|B_+\|_1\le C_f\mathrm{Des}^{1/3}\Lambda/t\le
C_f\theta_w$, and $|q^\#_0-q_0|+|\nu^\#-\nu|\le C_f\mathrm{Des}^{1/3}\Lambda$;
so $q^\#_0h^\#(B_+)\le q_0h(B_+)+C_f\theta_w$.  For the blocks,
$H_m(\Theta)=(\|D_m\Theta\|_2^2-\ip{D_m\Theta}{D_mw_m}^2/C_m^2)/C_m$, so
$|H^\#_m(\Theta_+)-H_m(\Theta_+)|\le C_f\|D_m\Theta_+\|_2^2(\|D_m(w^\#_m-w_m)\|_2
+|C^\#_m-C_m|)\le C_ft^{-2}\theta_wT_{\rm lo}^2\le C_f\theta_w$ (as
$\|D_m(t\Theta_+)\|_2\le\eta\le1$), and $|\sigma^\#_m-\sigma_m|\le C_f\theta_w$.
With [I,~Lemma~8.7(d)], $\Gamma^\#_w(B_+,\Theta_+)\le1+\eta_\Gamma(\eta)+
C_f\theta_w$.  The $-$ side is compared with $(B_-,\Theta_-)$: off $F$,
$B_--b^-=\mathfrak e_-$; on $F$, by [I,~(15)],
$(B_--b^-)1_F=\sum_{\mathrm{Kp}(w)}\epsilon(\tau'-\tau)u1_F+\sum'\Delta\theta
u1_F+\varkappa a/a(\zh^\#)$, the sum $\sum'$ over all carriers that are not
kept; and $\Theta_--\omega^-+d^\#(\omega^-)w^\#=X-\Delta\Theta-\sum_{\mathrm{Kp},
m}(\epsilon\tau'/\lambda)e_{k(l'')}$ (as $d^\#(\omega^-)=d^\#(\omega^+)$),
whose last two terms have $\lambda$-weighted $\ell_1$-norm at most
$\|\tau-\tau'\|_1+\sum'|\Delta\theta|$.  All these are at most $K_wt$.

\emph{(iii).}  Clamped coordinates are of kind [1] at $f^\#$.  Kept
carriers of types (K1), (K2) have $\gap^\#\ge M^\#u/2\ge\gamma(w)$ and type
(K3) has $\gap^\#\ge M^\#/2\ge\gamma(w)$ (Lemma~\ref{lem:II-status}(d),
$M^\#_m\ge M_m/2$); and $|\omega^+(k)|=|\omega_+(k)|\le|\Theta_+(k)|+|d_+|M\le
(3+\eta)/t\le4/t$ [I,~Lemma~8.8(a),(e)], $|\omega^-(k)|\le4/t+\tau'/\lambda
\le16/t$: kind [2] with $A_2=22$.  At a (K4) coordinate, after the shift,
$\varsigma_k\omega^-\ge-\frac32\gap^\#/t$, and $\varsigma_k\omega^+=
\varsigma_k\omega^--\tau'/\lambda\le\varsigma_k\omega^-$; before the shift
$\varsigma_k\omega^+=\varsigma_k\omega_+(k)\le\frac32\gap(k)/t\le\frac32
\gap^\#(k)/t$, and the shift moves $\varsigma_k\omega^\pm$ up only when
$\varsigma_k\omega^-<-\frac32\gap^\#/t$, in which case afterwards
$\varsigma_k\omega^+\le-\frac32\gap^\#/t$.  So for the $+$ pair either
$\varsigma_k\omega^+\le0$ (kind [3], side $+$) or $0<\varsigma_k\omega^+\le
\frac32\gap^\#/t$ (kind [1]); for the $-$ pair either $\varsigma_k\omega^-\ge0$
(kind [3], side $-$) or $|\omega^-|\le\frac32\gap^\#/t$ (kind [1]).  Sizes:
$|o_k|\le|\omega^-(k)|+\frac32/t\le17.5/t$, so $|\omega^+|\le21.5/t$ and
$|\omega^-|\le16/t$ after the shift; $A_2=22$ suffices.  Peaks of $f^\#$
carry no data.
\end{proof}

\subsection{Master Theorem II}\label{subsec:II-master2}

\begin{theorem}[Master Theorem II]\label{thm:II-master2}
Let $T=T_{\mathrm{final}}$ with base $U_{\mathrm{final}}$, $N\ge1$, and let
$f\in S_{p_N^*}$ have finite base support $F$.  Suppose that for infinitely
many levels $l$ there is a sub-window $w$ of level $l$ which is clean for $f$
and at which the following two conditions hold:
\begin{enumerate}
\item[\textup{(SH$_w$)}] every block $m\le N$ has an upper and a lower shift
source at $w$ \textup{(}Definition~\textup{\ref{def:II-sources})}, or the
shift cost of the shift pattern of $f$ at $w$ is robust:
$c_{\pi(w)}(f)\ge u(w)$ \textup{(}Definition~\textup{\ref{def:II-shiftpattern})};
\item[\textup{(VR$_w$)}] every extreme ray of the zero-cost cone
$C(\kappa(w))$ of the pattern of $f$ at $w$ has at most one robust
$d$-component \textup{(}Definition~\textup{\ref{def:II-pattern})}.
\end{enumerate}
Then $f\in\Rec$: $(f,g)\in\cl\NA((c_0,p_N),\ell_2^2)$ for every $g\in\Cset(f)$.
\end{theorem}

\begin{proof}
Fix $g\in\Cset(f)$, $\rho\in(0,1)$ and $\eta_0\in(0,2]$ with $\rho^2(1+\eta_0)
\le(1+\rho^2)/2$, and choose $\eta\le\eta_*$ with $\eta_\Gamma(\eta)\le
\min(1,\eta_0/8)$.  Let $l_1<l_2<\cdots$ be levels $\ge l_f$ with sub-windows
$w_j$ of level $l_j$ that are clean for $f$ and satisfy (SH$_{w_j}$) and
(VR$_{w_j}$), and put $f_j:=f^\#_{w_j}$ (Definition~\ref{def:II-companionw}).
The rows $f_j$ depend only on $f$ and on the $w_j$.  Their base supports
$F_j=F\cup\mathrm{Ba}_j\cup\mathrm{Pu}_j$ are finite and contain $F$, and
$p^*(f_j-f)\le\theta_{w_j}T_{\rm lo}(w_j)^2\to0$
(Lemma~\ref{lem:II-compcost}(e)).  Let $A_0$ be the $f$-constant of
Proposition~\ref{prop:II-transplant}(i), and let $t_1,j_0,c_0$ be given by
Lemma~\ref{lem:II-U} for the sequence $(f_j)$, $\varepsilon_{\rm tr}:=
\eta_0/2$ and $A_0$.  Put $T_j:=T_{\rm hi}(w_j)$, $n_j:=n(w_j)$, so that
$\mathcal S(T_j,n_j)$ is the set of dyadic scales of $\mathcal W(w_j)$ and
$T_j2^{-n_j}=T_{\rm lo}(w_j)$; $A_{2,j}:=22$, $\gamma_j:=\gamma(w_j)$,
$K_j:=\max(K_{w_j},1)$, $c_j:=c_0\gamma_j/22$.  Discarding finitely many $j$,
we may assume $T_j\le\min(t_\eta,1)$.  For $t\in\mathcal S(T_j,n_j)$ choose a
two-sided decomposition of $g$ at scale $t$ [I,~Lemma~8.5] and let
$g_{j,t}$ and $(b^\pm_{j,t},\omega^\pm_{j,t})$ be given by
Proposition~\ref{prop:II-transplant} at $w_j$.  We verify the hypotheses of
Theorem~\ref{thm:II-E}.

(E-a): $b^\pm(\xi_j)=0$, $t\|b^\pm\|_1\le A_0$, and
$\Gamma^{(j)}_w\le(\sqrt{1+\eta_0/8+C_f\theta_{w_j}}+K_{w_j}T_j)^2\le1+
\eta_0/2$ for large $j$, because $\theta_{w_j}\to0$ and $K_{w_j}T_j\to0$
(below).  (E-b): Proposition~\ref{prop:II-transplant}(iii) gives the kinds;
for $i\in F_j\setminus F=\mathrm{Ba}_j\cup\mathrm{Pu}_j$,
Proposition~\ref{prop:II-transplant}(iv) gives $z_j(i)b^+(i)\ge0\ge
z_j(i)b^-(i)$ on $\mathrm{Ba}_j$ and $t|b^\pm(i)|\le|a_j(i)|$ on
$\mathrm{Pu}_j$.  (E-c): Proposition~\ref{prop:II-transplant}(ii).  (E-d):
write $u:=u(w_j)$, $\mathrm{Des}:=\mathrm{Des}(l_j)$, $Q:=Q(w_j)\ge
(4\mathrm{Des}/u)^{20}\ge\mathrm{Des}\,u^{-5}$; then $K_jT_j\le C_f
\mathrm{Des}\,u^{-4}2^{-l_j^3}/(l_jQ)\le C_f2^{-l_j^3}/l_j\to0$, and, as
$c_j\ge c_fu$ and $n_j\ge l_j2^{l_j^3}Q$,
$n_jc_j/K_j\ge c_fl_j2^{l_j^3}Qu^5/\mathrm{Des}\ge c_fl_j2^{l_j^3}\to\infty$;
so $n_j\ge48\rho^2K_j/(c_j(1-\rho^2))$ for large $j$.  (E-e):
$\varepsilon_j:=p^*(f_j-f)\le\theta_{w_j}T_{\rm lo}(w_j)^2\le C_fu^{20}
T_{\rm lo}(w_j)^2$, while $\theta_*c_j^2(T_j2^{-n_j})^2\ge c_fu^2
T_{\rm lo}(w_j)^2$; since $u(w_j)\to0$ (Lemma~\ref{lem:II-scales}(e)), (E-e)
holds for large $j$.  After discarding finitely many $j$,
Theorem~\ref{thm:II-E} gives $(f,\rho g)\in\cl\NA((c_0,p_N),\ell_2^2)$.  As
$\rho<1$ was arbitrary and $\cl\NA$ is closed, $(f,g)\in\cl\NA$.
\end{proof}

The order of choices is that of Remark~\ref{rem:II-quantifiers}: the design
and $T_{\mathrm{final}}$; then $N$ and $f$; then the levels, the clean
sub-windows and the companions $f^\#_{w_j}$ (which depend on $f$ only);
then $g$ and $\rho$; then the data at each scale.  Every constant entering
the window arithmetic is an $f$-constant, a design number of level $l$
bounded by a power of $\mathrm{Des}(l)$, or a power of $1/u(w)$ coming from
robust rates.

\begin{corollary}[One block]\label{cor:II-master2N1}
Let $N=1$ and let $f\in S_{p_1^*}$ have finite base support.  If for
infinitely many levels $l$ some sub-window of level $l$ that is clean for $f$
satisfies \textup{(SH$_w$)}, then $f\in\Rec$.
\end{corollary}

\begin{proof}
For $N=1$ every component of every ray lies in the single block, so
(VR$_w$) holds at every $w$ (Remark~\ref{rem:II-onesigned}(b)); apply
Theorem~\ref{thm:II-master2}.
\end{proof}

\begin{corollary}[Every block has a bank donor]\label{cor:II-AC}
Let $N\ge1$, let $f\in S_{p_N^*}$ have finite base support, and let $w$ be
a sub-window of level $l\ge l_f$ that is clean for $f$.  Then every block
$m\le N$ has a donor $c_m$: a coarse peak with $\varrho_{c_m}\ge1+u(w)$, which
is of class R or dropped, and which is pinned when \textup{(SH$_w$)} holds.
For every block $m\in I_D(w)$, the move \textup{(C3)} at $c_m$ raises
$\mathrm{vs}_mu_{c_m}(\zh)$ by an amount in $[\Lambda_w,3\Lambda_w]/
\lambda_{c_m}$ up to $2b_w+C_f\mathrm{Des}(l)^2\Lambda_w^2$, either by a move
of the sign vector alone or, when the room of $S_{c_m}$ at $s_m$ toward
$\mathrm{vs}_m$ is too small, by closing $z_{s_m}$ to $\mathrm{vs}_m$ and
placing a bank of mass at most $\mathrm{Des}(l)^{1/3}\Lambda_w$ at $s_m$; it
changes the value of every other carrier $\le l$ by at most $C_f
\mathrm{Des}(l)^2\Lambda_w^2$.  At $f^\#_w$ the threshold of block $m$ has
risen by at least $c_f\Lambda_w$, and every \textup{(K4)} carrier of block $m$
\textup{(}weak or degenerate swallowing-type peak, or swallowing-type strict
non-peak with tiny gap\textup{)} is a strict non-peak of $f^\#_w$ with gap at
least $c_fM^\#_m\Lambda_w$.  In particular no donor hypothesis is needed in
Theorem~\ref{thm:II-master2}, also when every far signature coordinate of
every robust-margin peak is a contact with the sign of that peak.
\end{corollary}

\begin{proof}
Lemmas~\ref{lem:II-donorpeak}, \ref{lem:II-pinned}, \ref{lem:II-donorraise}
and~\ref{lem:II-status}(b),(c).
\end{proof}

\begin{remark}[The aligned corner]\label{rem:II-AC}
A move of the sign vector at a coordinate $s\in S_c$ raises $|u_c(\zh)|$
only if $z_s\ne\mathrm{vs}_c$.  The configuration in which, for every peak
that could serve as a donor, every far signature coordinate is a contact
with the sign of that peak (the \emph{aligned corner}) therefore defeats
every companion that keeps the base part fixed.  A bank at such a contact
(mass with the contact sign, sign vector unchanged) acts through the
Hilbert part: by \eqref{eq:II-bankformula}, $\mathrm{vs}_cu_c(\zh)$ increases
by $\mathfrak d(\mu^\star_s)^2v_c(s)/\nu$ at first order, while every other
coarse value changes only at second order, because the base is diagonal
($\ip{U^*u_k}{k_s}=\mu^\star_su_k(s)=0$ for $k\ne c$).  Banks are admissible
in Theorem~\ref{thm:II-E}, since only the sign of $a_j$ at the coordinates of
$F_j\setminus F$ enters (Remark~\ref{rem:II-E}(c)); the bank coordinates need
not be contacts of $f$.  The smallness of the bank masses
(Lemma~\ref{lem:II-donorraise}(i)) is genuinely used: it keeps the
second-order effects below the raise.
\end{remark}

\begin{corollary}[Residual configurations]\label{cor:II-residual}
Let $N\ge1$ and let $f\in S_{p_N^*}$ have finite base support.  If
$f\notin\Rec$, then there is $l_0$ such that for every level $l\ge l_0$ and
every sub-window $w$ of level $l$ that is clean for $f$ at least one of the
following holds:
\begin{enumerate}
\item[\textup{(B$_w$)}] some extreme ray $r$ of $C(\kappa(w))$ has robust
$d$-components in two different blocks $m_1\ne m_2$:
$|D_{r,m_i}(f)|\ge u(w)\Phi_{\max}(r,m_i)$, $i=1,2$;
\item[\textup{(C$_w$)}] some block has no upper or no lower shift source at
$w$, and the shift cost of the shift pattern of $f$ at $w$ is tiny:
$c_{\pi(w)}(f)\le b(w)$.
\end{enumerate}
For $N=1$, \textup{(B$_w$)} never occurs.
\end{corollary}

\begin{proof}
At a clean sub-window every rate object is tiny or robust, so (SH$_w$)
fails iff (C$_w$) holds, and (VR$_w$) fails iff (B$_w$) holds.  Clean
sub-windows exist at every level (Theorem~\ref{thm:II-clean}).  If the
conclusion failed, there would be infinitely many levels with a clean
sub-window satisfying (SH$_w$) and (VR$_w$), and
Theorem~\ref{thm:II-master2} would give $f\in\Rec$.  The last sentence is
Remark~\ref{rem:II-onesigned}(b).
\end{proof}

\begin{remark}[Scope and comments]\label{rem:II-master2}
(a) Theorem~\ref{thm:II-master2} imposes no hypothesis on donors, on
compensation or one-signedness of blocks, on nearly neutral carriers, and
no rate or growth condition: rooms of signature sets, target rooms,
margins (including degenerate peaks of either type), gaps, $d$-coefficients
and ray $d$-sums are arbitrary, and so are the contact set, the number of
swallowed carriers and the alignment of the far signature tails.  Every
$f$-dependent rate enters only through the tiny/robust dichotomy at a clean
sub-window: tiny objects are made exactly zero at the companion (closed
rooms and target rooms, neutralized ray components, raised thresholds), and
robust objects enter the constants as powers of $1/u(w)$.

(b) Earlier versions of this argument assumed shift sources in every block,
donors given by moves of the sign vector, every block compensated or
one-signed, and no wrong-sign nearly neutral kept carrier.  The first is the
case $I_{\rm up}(w)\cup I_{\rm lo}(w)=\emptyset$ of (SH$_w$); the second is
removed by Corollary~\ref{cor:II-AC}; the last two are removed by the exact
two-sided tuning (C4) together with ray removal (Step~3 of
Proposition~\ref{prop:II-transplant}), within (VR$_w$).  Rays with robust
components in two blocks that are compensated by fixed repair directions are
not covered by Theorem~\ref{thm:II-master2} as stated.

(c) Item (C$_w$) is not exactified by the companion: the shift cost
$c_{\pi(w)}(f)$ is computed from the decomposition at $f$, and no move at the
companion changes that decomposition.  For item (B$_w$) the obstruction is
not that tiny joint objects of several blocks (smallest singular values or
minors of the ray $d$-matrix, rate objects (O7), (O8)) cannot be made zero by
moving the values: these objects are polynomials in the values with design
coefficients.  The tuning (C4) is \emph{linear} (one least-norm solve), and
it neutralizes single components only.  A nonlinear exactification through
the determinantal data $\delta_{\rm comb}(l)$, $\mathrm{Lip}(l)$, $C_L(l)$,
$N_L(l)$ of Definition~\ref{def:II-constants}(c) is not carried out in this
section; it is the subject of Theorem~\ref{thm:II-B} and
Corollary~\ref{cor:II-B1}.

(d) The base enters the constants of this section only through the bank
efficiency $(\mu^\star_{s'})^2v_c(s')$ at coordinates $s'\le s_{\rm t}(l)$
(Lemma~\ref{lem:II-scales}(d)), through $\|U\|$, and through the identities
of Lemma~\ref{lem:II-base}(b); pull coordinates lie far beyond
$s_{\rm t}(l)$, but their Hilbert footprints are only bounded from above.
\end{remark}

%%%%%%%%%%%%%%%%%%%%%%%%%%%%%%%%%%%%%%%%%%%%%%%%%%%%%%%%%%%%%%%%%%%%%%%%

\section{Determinantal exactification, the shift dichotomy and the residual (C$^*$)}\label{sec:II-Cstar}

This section treats first rows with \emph{finite} base support.  It
removes two of the residual configurations left by the assembly of
Section~\ref{sec:II-exact} (nearly dependent $d$-vectors of several
blocks, and $f$-dependent $d$-row Hoffman constants), reduces the third
one (an unpinned uniform shift of the $d$-coefficients) to a single
residual property (C$^*$), and collects what is proved about (C$^*$).

\smallskip\noindent\emph{Standing conventions.}  Throughout,
$T=T_{\mathrm{final}}$ and $U=U_{\mathrm{final}}$ (Section~\ref{sec:II-Tfinal}),
except in Subsection~\ref{subsec:II-absorb}, where a modification of
$T_{\mathrm{final}}$ is used, and except where a statement says ``any
admissible $T$''.  $N\ge1$, $I=\{1,\dots,N\}$, $p=p_N$, and $f\in
S_{p^*}$ has forced data $(\xi,q_0,a,w,F,z,\zh,e,\nu)$ with $F=\supp a$
finite.  We use the notation of Sections~\ref{sec:II-companions} and
\ref{sec:II-Tfinal}; in particular for a carrier $l=\mathfrak j(k,m)$
\[
 \mathrm{val}_l:=\epsilon_lu_l(\zh),\qquad A_m:=|\hat\zeta_m|_m=\sigma_m/q_0,
 \qquad \bar\vartheta_m:=\bar\vartheta(\hat\zeta_m),\qquad
 \varrho_l:=\varrho_{k}(\hat\zeta_m),
\]
where $\epsilon_l$ is a sign attached to $l$ (the swallowing sign when $l$
is swallowed), and $\varrho_l$ is the relative position (it equals the rate
object (O4) of Definition~\ref{def:II-rates}).  The $d$-weight of
$l$ is $q_l=\epsilon_l\Phi_lw_m(k)/(mC_m)$ (Definition~\ref{def:II-class}); at a strict non-peak,
by [I,~Lemma~1.7] and $\hat\zeta_m(k)=m\Phi_lu_l(\zh)$,
\begin{equation}\label{eq:II-qval}
 q_l=\mathrm{val}_l/A_m .
\end{equation}
For a clean sub-window $w$ of level $l$ (Definition~\ref{def:II-clean}) we
write $b=b(w)$, $u=u(w)$, $\mathrm{Des}=\mathrm{Des}(l)$, and we use the
classification of the coarse carriers (classes R and G, swallowing and
anti type, kept and dropped carriers, the closed pattern $\kappa(w)$), the
upper and lower shift sources, the sets $I_{\rm up}(w)$, $I_{\rm lo}(w)$,
the shift pattern $\pi(w)$ with its shift cost $c_{\pi(w)}(f)$ (rate object
(O6)), the shift condition (SH$_w$), the robust-margin peaks of
Lemma~\ref{lem:II-donorpeak}, the pinning bounds of Lemma~\ref{lem:II-diagpin}, the
assembled companion $f^\#_w$ with its moves (C1)--(C4), the tuning
Lemma~\ref{lem:II-TU}, the status Lemma~\ref{lem:II-status}, the transplant
Proposition~\ref{prop:II-transplant} and Master Theorem~II
(Theorem~\ref{thm:II-master2}), all from Section~\ref{sec:II-exact}.
Two conventions for shifts occur: the \emph{data convention}
$\Delta d_m=d_m(\omega^-_m)-d_m(\omega^+_m)$ of [I,~Definition~7.7] for
two-piece data, and the \emph{decomposition convention}
$\Delta d_m=d_{+,m}-d_{-,m}$ of [I,~Lemma~8.27] for two-sided
decompositions at a scale $t$; for the latter we write
$\Delta\theta_l:=\lambda_l(\Theta_+-\Theta_-)_m(k(l))$ and, for class-G
carriers, $\tau_l:=-\epsilon_l\Delta\theta_l$.  Constants $C_f$ depend on
$f$ (and on $g,\rho$ where said) but not on the level, the sub-window or
the scale; they change from line to line.  $l_f$ is a level threshold
depending on $f$ only.

\smallskip\noindent\emph{Organization.}
Subsection~\ref{subsec:II-hoffman}: Hoffman constants through minors and
simultaneous exactification.  Subsection~\ref{subsec:II-thmB}: exactification
of all tiny minors at a companion (Theorem~\ref{thm:II-B}).
Subsection~\ref{subsec:II-dichotomy}: the shift dichotomy
(Theorem~\ref{thm:II-C1}), recovery with shifted data and Master
Theorem~III$'$ (Theorem~\ref{thm:II-MTIII}), whose only residual is (C$^*$)
(Definition~\ref{def:II-Cstar}).  Subsection~\ref{subsec:II-structure}:
rigidity of shifted data.  Subsection~\ref{subsec:II-SA}: realizations of
(C$^*$)-type resonances at block-tame rows, and genericity facts.
Subsection~\ref{subsec:II-absorb}: absorber pairs and Master
Theorem~IV$'$ under the lever hypothesis (KN).
Subsection~\ref{subsec:II-nested}: nested-tuning rows and further tools.
Subsection~\ref{subsec:II-open}: what remains open.

\subsection{Hoffman constants through minors; simultaneous exactification}\label{subsec:II-hoffman}

For a matrix $A\in\R^{p\times n}$ and index sets $J\subseteq\{1,\dots,p\}$,
$K\subseteq\{1,\dots,n\}$ we write $A_J$ for the rows in $J$ and $A_{J,K}$
for the submatrix.  $J$ is \emph{independent} if $J\ne\emptyset$ and the
rows $A_i$ $(i\in J)$ are linearly independent; $\sigma_{\min}(A_J)>0$ is
then the least singular value of $A_J$.  Put
\[
 \delta(A):=\min\{|\det A_{J,K}|:|J|=|K|\ge1,\ \det A_{J,K}\ne0\}
\]
(the least \emph{nonzero} minor; $\delta(A):=1$ if $A=0$).

\begin{lemma}[Hoffman constants through nonzero minors]\label{lem:II-hoffman}
Let $A\in\R^{p\times n}$, $\beta\in\R^p$, and $\mathcal P:=\{x\in\R^n:Ax\le
\beta\}\ne\emptyset$.  Then for every $x\in\R^n$
\begin{equation}\label{eq:II-hoff1}
 \dist_2(x,\mathcal P)\le\|(Ax-\beta)_+\|_2\cdot\max\{\sigma_{\min}(A_J)^{-1}:
 J\ \text{independent}\}
\end{equation}
\textup{(}the maximum over the empty family is $0$\textup{)}, and for every
independent $J$
\begin{equation}\label{eq:II-hoff2}
 \sigma_{\min}(A_J)\ge\max_{|K|=|J|}|\det A_{J,K}|\,\big/\,\|A_J\|_2^{|J|-1},
\end{equation}
where the maximum is positive.  Consequently
\begin{equation}\label{eq:II-hoff3}
 \begin{gathered}
 \dist_2(x,\mathcal P)\le\max(1,\|A\|_2)^{n-1}\delta(A)^{-1}\|(Ax-\beta)_+\|_2,\\
 \dist_1(x,\mathcal P)\le\sqrt n\max(1,\|A\|_2)^{n-1}\delta(A)^{-1}
 \|(Ax-\beta)_+\|_1 .
 \end{gathered}
\end{equation}
Equality constraints $Ex=e$ are included by writing them as two rows
$Ex\le e$, $-Ex\le-e$; their residual is $|Ex-e|$, and $\delta$ is computed
for the matrix with rows $A$ and $E$.  The bounds do not depend on $\beta$.
\end{lemma}

\begin{proof}
Let $x^*$ be the Euclidean projection of $x$ onto the closed convex set
$\mathcal P$ and assume $x\ne x^*$.  Then $x-x^*$ lies in the normal cone of
$\mathcal P$ at $x^*$, which for a polyhedron is the cone generated by the
rows $A_i^{\rm T}$ with $A_ix^*=\beta_i$ (Farkas' lemma).  By the conic
Carath\'eodory theorem $x-x^*=\sum_{i\in J}v_iA_i^{\rm T}$ with $v_i>0$ and $J$
independent.  As $A_ix^*=\beta_i$ for $i\in J$,
\[
 \|x-x^*\|_2^2=\sum_{i\in J}v_iA_i(x-x^*)=\sum_{i\in J}v_i(A_ix-\beta_i)\le
 \|v\|_2\,\|(A_Jx-\beta_J)_+\|_2 ,
\]
and $\|v\|_2\le\|A_J^{\rm T}v\|_2/\sigma_{\min}(A_J)=\|x-x^*\|_2/
\sigma_{\min}(A_J)$; this gives \eqref{eq:II-hoff1}.  For \eqref{eq:II-hoff2}
let $s_1\ge\dots\ge s_k>0$ ($k=|J|$) be the singular values of $A_J$.  By the
Cauchy--Binet formula $s_1^2\cdots s_k^2=\det(A_JA_J^{\rm T})=\sum_K\det
(A_{J,K})^2\ge\max_K\det(A_{J,K})^2$, which is positive because $A_J$ has
rank $k$; hence $s_k\ge(s_1\cdots s_k)/s_1^{k-1}\ge\max_K|\det A_{J,K}|/
\|A_J\|_2^{k-1}$.  Every positive maximum in \eqref{eq:II-hoff2} is $\ge\delta
(A)$, $\|A_J\|_2\le\|A\|_2$ and $|J|\le n$, which gives the first bound of
\eqref{eq:II-hoff3}; the second follows from $\|y\|_1\le\sqrt n\|y\|_2$ on
$\R^n$ and $\|y\|_2\le\|y\|_1$.  An independent set never contains a row
together with its negative, so for equality rows only $|Ex-e|$ enters.
\end{proof}

Thus a system whose minors are either exactly zero or bounded below has a
controlled Hoffman constant, however degenerate its combinatorics; only
minors that are small \emph{but nonzero} are harmful.

\begin{remark}[Nearly dependent $d$-vectors]\label{rem:II-tworay}
Linear tuning of single ray components (rate objects (O5)) does not control
the $d$-row of several blocks.  Let $\mathcal C=\{\mu_1r_1+\mu_2r_2:\mu\in
[0,\infty)^2\}$ with extreme rays $r_1,r_2$ of disjoint supports and
$\|r_i\|_1=1$, and let the $d$-map $D$ of two blocks take the values
$D(r_1)=(1,1)$ and $D(r_2)=(-1,-1+\delta)$, $0<\delta<1$.  Every component
of $D(r_1),D(r_2)$ has modulus $\ge1-\delta$.  But $\mathcal C\cap\ker D=\{0\}$,
and $\tau:=r_1+r_2$ has $D(\tau)=(0,\delta)$, $\dist_1(\tau,\mathcal C\cap\ker
D)=\|\tau\|_1=2$: the Hoffman constant of the system ``$\tau\in\mathcal C$,
$D\tau=0$'' is at least $2/\delta$.  The responsible quantity is the minor
$\det\begin{pmatrix}1&1\\-1&-1+\delta\end{pmatrix}=\delta$, which is
quadratic in the entries; making it vanish is not a linear problem.
\end{remark}

\begin{lemma}[Simultaneous exactification]\label{lem:II-loj}
Let $\pi_1,\dots,\pi_s$ be real polynomials in $n$ variables and
$\mathcal Q:=[-2,2]^n$.  There are $C_L\ge1$ and an integer $N_L\ge1$,
depending only on $(\pi_1,\dots,\pi_s)$, such that for every
$S\subseteq\{1,\dots,s\}$, every $\beta\in[0,1]$ and every $v\in\mathcal Q$
with $\max_{i\in S}|\pi_i(v)|\le\beta$, at least one of the following holds:
\begin{enumerate}
\item[(i)] $\beta\ge1/C_L$;
\item[(ii)] there is $v'\in\R^n$ with $\pi_i(v')=0$ for all $i\in S$ and
$\|v'-v\|_2\le C_L\beta^{2/N_L}$.
\end{enumerate}
\end{lemma}

\begin{proof}
Fix $S$, put $\Pi_S:=\sum_{i\in S}\pi_i^2$ and $Z_S:=\Pi_S^{-1}(0)$, and let
$g_S(v):=\dist_2(v,Z_S)$ if $Z_S\ne\emptyset$, $g_S:\equiv1$ otherwise.  Both
$\Pi_S$ and $g_S$ are continuous semialgebraic functions on the compact
semialgebraic set $\mathcal Q$, and $\Pi_S^{-1}(0)\cap\mathcal Q\subseteq
g_S^{-1}(0)$.  By the {\L}ojasiewicz inequality \cite[Corollary~2.6.7]{BCR} there are an integer
$N_S\ge1$ and $C_S>0$ with $g_S^{N_S}\le C_S\Pi_S$ on $\mathcal Q$.  If
$Z_S=\emptyset$ this reads $1\le C_S\Pi_S(v)\le C_S|S|\beta^2$, i.e.\ (i) for
every $C_L\ge(C_S|S|)^{1/2}$.  Otherwise a nearest point $v'$ of $Z_S$ has
$\|v'-v\|_2=g_S(v)\le(C_S|S|\beta^2)^{1/N_S}$.  Take $N_L:=\max_SN_S$ and
$C_L:=\max_S\max\{(C_S|S|)^{1/2},(C_S|S|)^{1/N_S},1\}$; for $\beta\le1$,
$\beta^{2/N_S}\le\beta^{2/N_L}$.
\end{proof}

The design constants $N_L(l)$, $C_L(l)$ of
Definition~\ref{def:II-constants}(c) are the constants of
Lemma~\ref{lem:II-loj} for the finite family of determinantal polynomials
$\pi_{\kappa,J,K}$ of level $l$; only their existence is used, and the
sub-window parameter $b(w)$ absorbs the non-explicit exponent
(Proposition~\ref{prop:II-subwindows}(d)).  The exponent cannot be $1$ in
general: for $\pi_1(v)=v_1^2$, $|\pi_1(v)|\le\beta$ allows $|v_1|=\beta^{1/2}$,
while the zero set is $\{v_1=0\}$.

\begin{lemma}[Ray components are minors]\label{lem:II-raycomp}
Let $\mathcal C=\{\tau\in\R^n:A\tau\le0\}$ be a pointed polyhedral cone
\textup{(}equalities written as two rows\textup{)}, $r$ an extreme ray with
$\|r\|_1=1$, and $d\in\R^n$ a row vector.  There is a set $J'$ of $n-1$
linearly independent rows of $A$, tight at $r$, such that
\[
 \det\begin{pmatrix}A_{J'}\\ d\end{pmatrix}=\pm\|c\|_1\,(d\cdot r),\qquad
 c_i:=(-1)^{n+i}\det A_{J',\{1,\dots,n\}\setminus\{i\}},\qquad c\ne0 .
\]
\end{lemma}

\begin{proof}
The rows tight at $r$ have rank $n-1$ since $r$ is extreme in a pointed
cone; choose $J'$ independent among them.  Laplace expansion along the last
row gives $\det[A_{J'};d]=d\cdot c$ for every $d$; for $d$ a row of $A_{J'}$
the determinant vanishes, so $c\in\ker A_{J'}=\R r$, and $c\ne0$ because
$A_{J'}$ has rank $n-1$.  Hence $c=\pm\|c\|_1r$.
\end{proof}

Applied to the zero-cost cone $C(\kappa)$ of a pattern and to the value row
of a block $m$ in the matrix $A_\kappa(v)$ of
Definition~\ref{def:II-constants}(c), Lemma~\ref{lem:II-raycomp} shows that
every ray component $D_{r,m}$ (rate object (O5)) is $\pm\|c\|_1^{-1}$ times
a minor $\pi_{\kappa,J,K}(\mathrm{val})$ (rate object (O8)), where $\|c\|_1$ is
a sum of moduli of minors free of value rows, hence lies between
$\delta_{\rm comb}(l)$ and $(r(l)c'(l))^l$ times a design number.  The
converse fails: in Remark~\ref{rem:II-tworay} all components are robust and
only the $2\times2$ minor is small, and there are three-block configurations
with all ray components and all $2\times2$ value minors robust and a single
small $3\times3$ minor.

\begin{lemma}[Threshold buffer through a peak]\label{lem:II-QB}
Let $\zeta\in\ell_1\setminus\{0\}$ be a vector of one block,
$\bar\vartheta:=\bar\vartheta(\zeta)$, $\mathsf a:=|\zeta|$,
$\phi:=\|\Phi\|_2^2$, and let $c$ be a peak of $\zeta$, $s>0$, and
$\zeta'\in\ell_1$ with $\sgn\zeta'(c)=\sgn\zeta(c)$, $|\zeta'(c)|=|\zeta(c)|+s$
and
\[
 E:=\sum_{k\ne c}|\zeta'(k)-\zeta(k)|\le\frac{\mathsf as}{8(\mathsf a+\bar\vartheta+1)} .
\]
Let $L$ be any set of coordinates, $R(s):=\min\{1,\mathsf as/(8\phi(\mathsf a+
\bar\vartheta+1))\}$ and $X:=\sup_{k\in L\setminus\{c\}}|\zeta'(k)-\zeta(k)|/
\Phi(k)^2$.  Then:
\begin{enumerate}
\item[(a)] $\bar\vartheta(\zeta')\ge\bar\vartheta+R(s)$;
\item[(b)] every $k\in L\setminus\{c\}$ with $\mathsf v_k(\zeta)<\bar\vartheta+
R(s)-X$ is a strict non-peak of $\zeta'$;
\item[(c)] if $k_0$ is a peak of $\zeta$ with $\mathsf v_{k_0}(\zeta)>\bar\vartheta$
\textup{(}$k_0=c$ is allowed\textup{)}, $C_1:=(1+2(\bar\vartheta+1)/\mathsf a)/
\Phi(k_0)^2$, $h_0:=\min(1,\mathsf v_{k_0}(\zeta)-\bar\vartheta)$, and
$C_1(s+E)<h_0$, then $\bar\vartheta(\zeta')\le\bar\vartheta+C_1(s+E)$, and every
$k\in L\setminus\{c\}$ with $\mathsf v_k(\zeta)>\bar\vartheta+C_1(s+E)+X$ is a
peak of $\zeta'$ with $\sgn\zeta'(k)=\sgn\zeta(k)$.
\end{enumerate}
\end{lemma}

\begin{proof}
(a) is the lower bound of Lemma~\ref{lem:II-T2}(b).  (b) $\mathsf
v_k(\zeta')\le\mathsf v_k(\zeta)+X<\bar\vartheta+R(s)\le\bar\vartheta(\zeta')$.
(c) The upper bound of Lemma~\ref{lem:II-T2}(a) applies to every peak $k_0$
with $\mathsf v_{k_0}>\bar\vartheta$, the pushed peak $c$ included, with
$\|\zeta'-\zeta\|_1=s+E$ (the $c$-term of $\mathsf A$ grows by $s$, which the
proof of Lemma~\ref{lem:II-T2}(a) accounts for as part of the $\ell_1$
distance).  Then $\mathsf v_k(\zeta')\ge\mathsf v_k(\zeta)-X>\bar\vartheta
(\zeta')$, and $|\zeta'(k)-\zeta(k)|\le X\Phi(k)^2<\mathsf v_k(\zeta)\Phi(k)^2=
|\zeta(k)|$ gives the sign.
\end{proof}

The point of allowing $k_0=c$ is that a block may have no non-degenerate
peak other than the one that is pushed.  In applications $L$ is the set of
coarse carriers of the block; fine coordinates may change their status, and
enter only through $E$.

\begin{lemma}[Buffer peaks]\label{lem:II-buffer}
There is $l_f$ such that for every level $l\ge l_f$, every sub-window $w$
of level $l$ and every block $m$ there is a coarse peak $c=c_m$
\textup{(}a carrier $\le l$\textup{)} with $|\alpha_m(c)|\ge1/(2l)$.  Its
relative margin satisfies
\[
 \varrho_c-1=\frac{A_m|\alpha_m(c)|}{\Phi_c^2\bar\vartheta_m}\ge
 \frac{A_m}{2l\,\Phi_c^2\bar\vartheta_m}\ge\frac{A_m}{2l\,\bar\vartheta_m} .
\]
\end{lemma}

\begin{proof}
By [I,~Lemma~1.7], $\sum_{k\in P_m}|\alpha_m(k)|=1$, and by [I,~(2)],
$\sigma_m|\alpha_m(k)|=\lambda_{k,m}\mu_{k,m}$ with $\mu_{k,m}\le|u_{k,m}(\xi)|
\le q_0\|\zh\|_\infty\le2q_0$.  Hence the peaks of carriers $>l$ carry at
most $2q_0\sum_{l'>l}\lambda_{l'}/\sigma_m\le2q_0T_{\rm lo}(w)^8/\sigma_m$
(Lemma~\ref{lem:II-GW}(GW3)) of the $\alpha$-mass, which is $\le\frac12$ for
$l\ge l_f$; and block $m$ has at most $l$ coarse carriers.  The identity for
the margin is Lemma~\ref{lem:II-T}(b) applied to $\hat\zeta_m$
(${A_m}|\alpha_m(c)|=|\hat\zeta_m(c)|-\bar\vartheta_m\Phi_c^2$), and
$\Phi_c^2\le1$.
\end{proof}

Since $A_m/\bar\vartheta_m$ depends on $f$ only, while $u(w)=b(w^-)\le
T_{\rm lo}(w^-)^4$ for $w\ne(1,1)$ decays faster than any power of $1/l$
(Proposition~\ref{prop:II-subwindows}), the buffer peak is robust
($\varrho_c\ge1+u(w)$) at every sub-window of every large level.  Pushing it outward by a
controlled amount raises the block threshold (Lemma~\ref{lem:II-QB}) and
protects the strict non-peaks of the block; with the diagonal base such a
push is available in either direction on far coordinates of its signature
set, by z-moves (for a class-R peak, as in (C3)) or by pulls and banks (for
a class-G peak after (C1), as in Lemma~\ref{lem:II-TU}), so no donor is
needed for this purpose.

\subsection{Exactification of all tiny minors at a companion}\label{subsec:II-thmB}

Let $w$ be a clean sub-window of level $l\ge l_f$ for $f$ and let
$\kappa=(E,P,\epsilon,F',\mathrm{type})$ be the closed pattern of $w$, with
kept set $\mathrm{Kp}=E\setminus P$ (Definition~\ref{def:II-constants}(a);
$E$ is the set of coarse class-G carriers, $P$ the dropped ones).  The
transplant of Proposition~\ref{prop:II-transplant} projects the switching amplitudes
$\tau=(\tau_{l''})_{l''\in E}$ of the decompositions of a mate at the scales
of $w$ onto the solution set of the following system $\Sigma^\#$, whose
coefficients are taken at a companion $f^\#$ of $f$:
\begin{enumerate}
\item[(Z1)] $\tau_{l''}\ge0$ $(l''\in\mathrm{Kp})$;
\item[(Z2)] $\mathrm{type}(j)L_j(\tau)\ge0$ for $j\in T(l)\setminus F'$ with
$\mathrm{type}(j)=\pm1$, and (Z3) $L_j(\tau)=0$ if $\mathrm{type}(j)=0$, where
$L_j(\tau)=\sum_{l''\in E}\epsilon_{l''}\tau_{l''}u_{l''}(j)$;
\item[(Z4)] $\tau_{l''}=0$ for $l''\in P$ (dropped carriers);
\item[(Z5)] $\sum_{l''\in\mathrm{Kp},\,m(l'')=m}(\mathrm{val}^\#_{l''}/A^\#_m)
\tau_{l''}=0$ for every block $m$ meeting $E$;
\item[(box)] $\tau_{l''}\le12\lambda_{l''}/t$ $(l''\in E)$.
\end{enumerate}
By \eqref{eq:II-qval} the rows (Z5) are the $d$-rows of the kept carriers
at $f^\#$, which are strict non-peaks of $f^\#$ (Lemma~\ref{lem:II-status}).  The
matrix of (Z1)--(Z5) is $A_\kappa(v)$ of Definition~\ref{def:II-constants}(c)
evaluated at $v=\mathrm{val}^\#$, except that the value row of block $m$ is
multiplied by $1/A^\#_m$; the rows (Z4) are unit rows whatever the status of
the dropped carrier at $f^\#$, so the statuses of dropped carriers do not
enter $\Sigma^\#$.  Recall $\eta(w)=T_{\rm lo}(w)^4/(l\,\mathrm{Des})$ and
$\beta(w)=(1+\mathrm{Lip}(l)C_*(l))b(w)$ from
Proposition~\ref{prop:II-subwindows}(d).

\begin{theorem}[Exactification of all tiny minors]\label{thm:II-B}
Let $T=T_{\mathrm{final}}$, $N\ge1$, and let $f\in S_{p^*}$ have $F$ finite, let $w$ be a clean sub-window of level
$l\ge l_f$ for $f$, $\kappa$ its closed pattern, and suppose an assembly is
given in the following form: first z-moves on coarse coordinates which close
the rooms of the class-G carriers on their natural signature sets and the
tiny target rooms on $T(l)$, producing a row $f_1$; then finitely many
further moves $\mathcal M_{\rm ass}$ \textup{(}donor raises by z-moves or
banks; pulls converting tiny-margin peaks\textup{)} at coordinates beyond
$s_{\max}(l)$ of carriers not in $\mathrm{Kp}$.  Assume:
\begin{enumerate}
\item[(A1)] $|\mathrm{val}_{l''}(f_1)-\mathrm{val}_{l''}(f)|\le C_*(l)b(w)$ for
$l''\in\mathrm{Kp}$, and $p^*(f_1-f)\le C_f\mathrm{Des}\,b(w)\log(e/b(w))$;
\item[(A2)] after $\mathcal M_{\rm ass}$ every kept carrier
\textup{(}$l''\in\mathrm{Kp}$\textup{)} is a strict non-peak, every coarse
peak of $f$ with $\varrho\ge1+u(w)$ is a peak with the same sign and
relative margin $\ge u(w)/2$, and every kept carrier either has gap $\ge
c_f\Lambda$
\textup{(}$\Lambda=T_{\rm lo}(w)^3$, in a block with a donor raise\textup{)}
or has $\varrho\le1-u(w)/2$, $\varrho\le C_*(l)b(w)$ or $|\varrho-1|\le
C_*(l)b(w)$ \textup{(}in a block without donor raise\textup{)}.
\end{enumerate}
Then there is a companion $f^\#$ of $f$, obtained from $f_1$ by the moves
$\mathcal M_{\rm ass}$, in each block without donor raise one outward push
of the buffer peak $c_m$ of Lemma~\textup{\ref{lem:II-buffer}}, and the
exact tuning of the values of $\mathrm{Kp}$, such that:
\begin{enumerate}
\item[(a)] $f^\#$ has the carriers $E$, the kept set $\mathrm{Kp}$, the signs
$\epsilon$ and the contacts on $T(l)$ of $\kappa$, and $\supp a^\#$ is
finite; every kept carrier is a strict non-peak of $f^\#$, every coarse peak
of $f$ with $\varrho\ge1+u(w)$ is a peak of $f^\#$ with the same sign, and
every coarse carrier with $|\varrho-1|\le C_*(l)b(w)$ at $f$ \textup{(}kept
or dropped; in particular every tiny-margin peak\textup{)} is a strict
non-peak of $f^\#$.  Nothing else is asserted about dropped carriers; they
enter $\Sigma^\#$ only through the unit rows \textup{(Z4)};
\item[(b)] $\mathrm{val}_{l''}(f^\#)=\mathrm{val}_{l''}(f_1)+x_{l''}$ exactly for
$l''\in\mathrm{Kp}$, with $\|x\|_2\le\eta(w)$; every minor $\pi_{\kappa,J,K}$
which is tiny at $f$ \textup{(}value $\le b(w)$\textup{)} vanishes at
$\mathrm{val}(f^\#)$, and every minor which is robust at $f$ has modulus
$\ge u(w)/2$ at $\mathrm{val}(f^\#)$;
\item[(c)] $p^*(f^\#-f)\le p^*(f_1-f)+C_f(\text{cost of }\mathcal M_{\rm ass})+
C_f\mathrm{Des}^4\eta(w)\log(e/\eta(w))$, which is $o(T_{\rm lo}(w)^2)$
when the cost of $\mathcal M_{\rm ass}$ is;
\item[(d)] the $\ell_1$-Hoffman constant of $\Sigma^\#$ at $f^\#$ \textup{(}all
rows, box rows included\textup{)} satisfies $C^\#_H\le C_f^l\,\mathrm{Des}^2/
u(w)$, for every configuration of the zero-cost cone \textup{(}single- or
multi-block rays, compensated, mixed or one-signed blocks, any face
structure\textup{)}.
\end{enumerate}
\end{theorem}

\begin{proof}
\emph{Step 1 (minors at $f_1$).}  Every minor $\pi_{\kappa,J,K}$ with a value
row is $\mathrm{Lip}(l)$-Lipschitz on $[-2,2]^{[1,l]}$, so by (A1) it moves by
at most $\mathrm{Lip}(l)C_*(l)b$ from $f$ to $f_1$.  Hence tiny minors have
modulus $\le\beta(w)$ at $\mathrm{val}(f_1)$, robust ones $\ge u-\beta(w)\ge
3u/4$ (Proposition~\ref{prop:II-subwindows}(d) and (a)).

\emph{Step 2 ({\L}ojasiewicz point).}  Apply Lemma~\ref{lem:II-loj} with the
family of level $l$, $S$ the tiny minors of $\kappa$, and $v:=\mathrm{val}
(f_1)\in[-2,2]^{[1,l]}$ ($|u_{l''}(\zh)|\le q^*(u_{l''})\,\|\zh\|_\infty\le2$).
Alternative (i) is excluded because $\beta(w)<1/C_L(l)$; so there is $v'$
with $\pi(v')=0$ for every tiny minor and $\|v'-v\|_2\le C_L(l)\beta(w)^{2/N_L
(l)}\le\eta(w)$ (Proposition~\ref{prop:II-subwindows}(d)).  The polynomials of
$\kappa$ involve only the coordinates in $\mathrm{Kp}$; put $x:=(v'-v)
|_{\mathrm{Kp}}$.  Robust minors satisfy $|\pi(v')|\ge3u/4-\mathrm{Lip}(l)\eta(w)
\ge u/2$ (as $\|v'-v\|_\infty\le\|v'-v\|_2\le\eta(w)$ and $\mathrm{Lip}(l)$ is
a constant for the maximum norm), since $\mathrm{Lip}(l)\eta(w)\le T_{\rm lo}(w)^4\le u/8$
(Proposition~\ref{prop:II-subwindows}(e)).

\emph{Step 3 (one combined realization).}  First, in each block $m$
without donor raise, realize an outward push of the buffer peak
$c_m\notin\mathrm{Kp}$ of a size $s\in[s_m,2s_m]$, with $s_m$ the number
fixed in Step~4 from a priori bounds: for a class-G buffer peak by a pull
or bank on far coordinates of $S_{c_m}$ (which satisfy (T1) of
Lemma~\ref{lem:II-TU} after (C1)), for a class-R buffer peak by a z-move
(or a close-and-bank) on far coordinates of its robust room, as in (C3)
(Lemma~\ref{lem:II-donorraise}).  Then apply Lemma~\ref{lem:II-TU}, with
$L_0=\mathrm{Kp}$, at the row obtained from $f_1$ by the moves $\mathcal
M_{\rm ass}$ and the pushes: their masses and z-changes are part of the
base data $(A^2,z^2)$ of that lemma; their coordinates lie in signature
sets $S_{l''}$ with $l''\notin\mathrm{Kp}$, so (T2) is unaffected, and they
are distinct from the pull and bank coordinates of $\mathrm{Kp}$, which lie
in the signature sets $S_{l''}$, $l''\in\mathrm{Kp}$, beyond $s_{\max}(l)$.
The target values are $v'_{l''}$.  Since
Lemma~\ref{lem:II-TU} solves for exact values given all other masses, the
second-order side effects of $\mathcal M_{\rm ass}$ and of the pushes on the
values of $\mathrm{Kp}$ are absorbed, and $\mathrm{val}_{l''}(f^\#)=v'_{l''}$
exactly.  The required increments are $\le2\eta(w)$, far below the range of
Lemma~\ref{lem:II-TU} (an $f$-constant times a negative power of
$\mathrm{Des}$).  Other coarse carriers move by at most $C_f\eta(w)^2$ beyond
the effect of $\mathcal M_{\rm ass}$ and of the pushes; pulls and banks lie
beyond $s_{\max}(l)$, so $T(l)$, the contact pattern on $T(l)$ and all
swallowing signs are unchanged.

\emph{Step 4 (statuses).}  In a block with a donor raise,
Lemma~\ref{lem:II-status} applies verbatim with the additional value moves
$|x|+C_f\eta^2\le2\eta(w)\ll c_f\Lambda$ and $\ll u$, so every kept carrier
stays a strict non-peak and every coarse peak of $f$ with $\varrho\ge1+u$
stays a peak with its sign (Lemma~\ref{lem:II-status}(c),(d)); and every
coarse carrier with $|\varrho-1|\le C_*b$ becomes a strict non-peak, since
the threshold rises by at least $c_f\Lambda$ (Lemma~\ref{lem:II-status}(b))
while its value moves by at most $C_f\mathrm{Des}^{1/6}\beta_w+2\eta\ll
\Lambda$, and $C_*b\ll\Lambda$.  In a block $m$ without donor raise put
$\zeta:=\hat
\zeta_m(f_1)$, $\zeta^\#:=\hat\zeta_m(f^\#)$ (both with the moves of the other
blocks included), and let $X$, $E$ be the quantities of Lemma~\ref{lem:II-QB}
for $L:=$ all coarse carriers of block $m$ (kept or dropped) and the
pushed peak $c=c_m$.  By Step 3 (Lemma~\ref{lem:II-TU}(a),(b) and the
bound $2\eta(w)$ on the increments), $X\le\bar X:=m\cdot2\eta(w)/
\Phi_{\min}(l)\le C\mathrm{Des}\,\eta(w)$ and $E\le\bar E:=\phi\bar X+
2^{-l}\eta(w)$, where $\Phi_{\min}(l)$ is the least weight of a carrier
$\le l$; these a priori bounds do not depend on the push.  With $r_m:=
C_*(l)b(w)\bar\vartheta_m$ the push size of Step~3 is fixed as
\[
 s_m:=\frac{16(\mathsf a+\bar\vartheta+1)}{\mathsf a}\max\{\bar E,\phi(\bar
 X+r_m)\}\qquad(\mathsf a,\bar\vartheta,\phi\text{ of }\zeta),
\]
and the realized push has size $s\in[s_m,2s_m]$.  Then $E\le\mathsf a s/
(8(\mathsf a+\bar\vartheta+1))$ and $R(s)\ge R(s_m)\ge X+r_m$ (for $l\ge l_f$,
$R(s_m)<1$), so Lemma~\ref{lem:II-QB}(a),(b) make every coarse carrier of
block $m$ other than $c_m$ with $\varrho\le1+C_*(l)b$ (kept or dropped; in
particular every tiny-margin peak) a strict non-peak of $\zeta^\#$.  With
$k_0:=c_m$ in Lemma~\ref{lem:II-QB}(c) ($\Phi(c_m)\ge1/\mathrm{Des}$, and
$h_0\ge\min(1,A_m/(2l))$ because Lemma~\ref{lem:II-buffer} gives $\mathsf
v_{c_m}-\bar\vartheta_m=\bar\vartheta_m(\varrho_{c_m}-1)\ge A_m/(2l)$), we
have $C_1(s+E)+X\le C_f\mathrm{Des}^4\eta(w)\le u\bar\vartheta_m/4$; in
particular $C_1(s+E)<h_0$ for $l\ge l_f$, because $C_f\mathrm{Des}^4\eta(w)
\le C_fT_{\rm lo}(w)^4\mathrm{Des}^3/l$ tends to $0$ while $A_m$ is an
$f$-constant.  So every coarse peak of $f$ with $\varrho\ge1+u(w)$ stays a
peak with its sign, and kept carriers with $\varrho\le1-u/2$ stay strict
non-peaks.  Kept near-threshold carriers of such a block become strict
non-peaks with $\gap^\#\ge\gap$, since the threshold does not decrease
(Lemma~\ref{lem:II-QB}(a)).  The buffer peak stays
a robust peak.  The value rows of $\Sigma^\#$ are rescaled by $1/A^\#_m$,
which changes no statement in (b).  This proves (a) and (b).

\emph{Step 5 (cost).}  By Lemma~\ref{lem:II-TU}(c) the pull and bank
masses are $\le C_f\mathrm{Des}\,2^{G(l)}\eta(w)$; with the masses of
$\mathcal M_{\rm ass}$ and of the pushes, Lemma~\ref{lem:II-cost} gives,
exactly as in the proof of Lemma~\ref{lem:II-compcost}(e),
$p^*(f^\#-f_1)\le C_f(\text{cost of }\mathcal M_{\rm ass})+C_f(2\eta)\log(e/
(2\eta))+C_f(s/\lambda_{c_m})\log(e\lambda_{c_m}/s)$ (summed over the
blocks without donor raise, $s\in[s_m,2s_m]$ the realized push), and
$s/\lambda_{c_m}\le2s_m/\lambda_{c_m}\le C_f\mathrm{Des}^4\eta(w)$.  Since $\eta(w)\le T_{\rm lo}(w)^4$,
the added terms are $o(T_{\rm lo}(w)^2)$.

\emph{Step 6 (Hoffman constant).}  The minors of the matrix of $\Sigma^\#$
at $f^\#$ that contain value rows equal the product of the row factors
$1/A^\#_m$ of the value rows used times $\pi_{\kappa,J,K}(\mathrm{val}(f^\#))$;
since $A^\#_m\le m2^{-m}\le\frac12$ \textup{(}$A_m\le\|\hat\zeta_m\|_1\le
m\sum_k\Phi_m(k)\le m2^{-m}$\textup{)}, these factors are $\ge2$, so by (b)
such a minor vanishes or has modulus $\ge u/2$.  Minors without value rows
are design numbers, zero or $\ge\delta_{\rm comb}(l)$.  The entries of the
matrix are bounded by $C_f$ (value entries are $\mathrm{val}/A^\#_m$, and
$A^\#_m$ is within $C_f\,p^*(f^\#-f)$ of $A_m$ by Lemma~\ref{lem:II-cost}).
Lemma~\ref{lem:II-hoffman} with at most $c'(l)$ columns and $r(l)$ rows gives
$C^\#_H\le\sqrt{c'(l)}\,(C_fr(l)c'(l))^{c'(l)}/\min\{\delta_{\rm comb}(l),u/2\}
\le C_f^l\mathrm{Des}^2/u$, using Proposition~\ref{prop:II-bank}(e) and
$c'(l)\le l$.
\end{proof}

The assembly of Section~\ref{sec:II-exact} has the form required by
Theorem~\ref{thm:II-B}: its moves (C1), (C2) produce $f_1$ and satisfy (A1)
by Lemma~\ref{lem:II-compcost}(b) and Lemma~\ref{lem:II-cost} (cost bound),
its donor raise (C3) is $\mathcal M_{\rm ass}$, and (A2) holds by
Lemma~\ref{lem:II-status}(c),(d) for kept carriers and robust peaks; the
statuses of dropped carriers are not needed.  The linear tuning (C4) is replaced by the
combined solve of Step~3, whose target, the {\L}ojasiewicz point, absorbs
the second-order effects of the donor banks on the kept values.

\begin{corollary}[Multi-block rays are not residuals]\label{cor:II-B1}
Let $w$ be a clean sub-window of level $l\ge l_f$ for $f$, and suppose that
the shift bound $|\Delta d_m|M_m\le K_dt$ holds for every block $m$ and
every scale $t\in\mathcal W(w)$, with $K_d\le C_f\mathrm{Des}^{C}/u(w)^{C}$.
Then the conclusions of Proposition~\textup{\ref{prop:II-transplant}} hold at the
companion $f^\#$ of Theorem~\textup{\ref{thm:II-B}} without the hypothesis
\textup{(VR$_w$)}, with a window constant $K_w\le C_f^l\mathrm{Des}^{C'}/
u(w)^{C'}$.  In particular Master Theorem~II
\textup{(}Theorem~\textup{\ref{thm:II-master2})} holds with \textup{(VR$_w$)}
deleted from its hypotheses.
\end{corollary}

\begin{proof}
In Proposition~\ref{prop:II-transplant}, (VR$_w$) is used only to bound the distance
of the projected amplitudes $\tau_0$ to the solution set of the $d$-rows
(blockwise ray removal).  Replace this step by the $\ell_1$-Hoffman
projection of $\tau_0$ onto the solution set of $\Sigma^\#$, box rows
included: Theorem~\ref{thm:II-B}(d) gives $\tau'$ with all (Z5) rows exactly
zero and $\|\tau_0-\tau'\|_1\le C^\#_H\cdot(\text{violations of }\Sigma^\#\text{
at }\tau_0)$.  The violations of (Z1)--(Z4) and (box) are those already
bounded in Proposition~\ref{prop:II-transplant}; the violation of (Z5) is bounded by
the shift term $C_fK_dt$ (through [I,~(17)]) plus $C_f\eta(w)\sum_{l''}
|\tau_{l''}|\le C_f\eta(w)\cdot12/t\le C_fT_{\rm lo}(w)^3$.  To see this, use
the homogeneity of (Z5): multiplying the row of block $m$ by $A^\#_m/A_m\in
[\frac12,2]$ changes neither its solution set nor (up to a factor $\le2$) its
violation, and turns its coefficients into $\mathrm{val}^\#_{l''}/A_m=
q_{l''}+x'_{l''}/A_m$ with $q_{l''}=\mathrm{val}_{l''}/A_m$ the coefficients
at $f$ (Lemma~\ref{lem:II-status}(d)) and $|x'_{l''}|\le C_*(l)b(w)+\eta(w)
\le2\eta(w)$; so $|q^\#_{l''}A^\#_m/A_m-q_{l''}|\le2\eta(w)/A_m\le C_f\eta(w)$,
and the violation is at most $2\big(|\sum_{l''}q_{l''}\tau_{l''}|+C_f\eta(w)
\sum_{l''}|\tau_{l''}|\big)$, where we used $\sum_{l''}\lambda_{l''}\le1$ and
the box rows.  (A bound $|q^\#_{l''}-q_{l''}|\le C\eta(w)$ itself would be
false in blocks with a donor raise, where $|A^\#_m-A_m|$ is of order
$\Lambda_w=T_{\rm lo}(w)^3$.)  The remaining steps
of Proposition~\ref{prop:II-transplant} use only $\|\tau_0-\tau'\|_1$, the sign rows
and the box, not the additional property $\tau'\le\tau_0$ of the ray removal.
All new factors are of the form $C_f^l\mathrm{Des}^{C}u^{-C}$, which the window
arithmetic absorbs (Lemma~\ref{lem:II-GW}(GW2)).  The extra cost
$C_f\mathrm{Des}^4\eta\log(e/\eta)$ of Theorem~\ref{thm:II-B}(c) is
$o(T_{\rm lo}(w)^2)$.
\end{proof}

Consequently the $f$-dependent $d$-row Hoffman constants of mixed or
uncompensated blocks, and nearly dependent $d$-vectors of several blocks
(Remark~\ref{rem:II-tworay}), are not obstructions for $T_{\mathrm{final}}$
and finite $F$.  What Theorem~\ref{thm:II-B} does not give is the shift bound
itself: the violation of (Z5) by the actual amplitudes contains the uniform
shift $\Delta d_m$, which must be pinned at $f$.  This is the subject of the
next subsection.

\subsection{The shift dichotomy and Master Theorem~III$'$}\label{subsec:II-dichotomy}

Fix a clean sub-window $w$ of level $l\ge l_f$ for $f$.  The
\emph{enlarged shift pattern of $f$ at $w$} is the index
$\kappa^{\rm sh}(w)=(\kappa,N_0,R_7,I_{\rm sh},\mathrm{src},G_{\rm pk})$ of
the rate object (O9) of Definition~\ref{def:II-rates} given by: $\kappa=
\kappa(w)$ the closed pattern; $N_0$ the kept carriers with $\varrho\le b$
(nearly neutral; they are exactified to value $0$ at the companion);
$R_7$ the anti-type class-G strict non-peaks with $|\varrho-1|\le b$;
$I_{\rm sh}:=I_{\rm up}(w)\cup I_{\rm lo}(w)$ the source-deficient blocks;
$\mathrm{src}(m)$ equal to ``$\le$'' if block $m\in I_{\rm sh}$ has an
upper source of type (U1), ``$\ge$'' if it has a lower source of type (L1),
and ``free'' otherwise; and $G_{\rm pk}$ the class-G coarse carriers that
are peaks of $f$ with $\varrho\ge1+u(w)$ (tiny-margin peaks, with
$\varrho\in[1,1+b(w)]$, belong to $E\setminus G_{\rm pk}$; by cleanness
every class-G peak is of one of these two kinds).
Since $\kappa^{\rm sh}(w)$ records all window classifications entering the
system $\Sigma^{\rm sh}$ of (O9), the shift-pinning rate $\rho^{\rm
sh}(\kappa^{\rm sh}(w),f)$ is a number depending on $f$ only, and it is
tiny or robust at $w$.  We use the decomposition convention $\delta_m:=
\Delta d_mM_m$ for two-sided decompositions at scales $t\in\mathcal W(w)$, and
put $K_g:=(1/q_0+1)\mathrm{Des}/u$.

\begin{lemma}[Relations satisfied by the decompositions]\label{lem:II-R17}
Let $g\in\Cset(f)$, $t\in\mathcal W(w)$, and let $(B_\pm,\Theta_\pm)$ be a
two-sided decomposition of $g$ at scale $t$ \textup{(}[I,~Lemma~8.5]\textup{)}.
Then the vector $(\delta,\tau)$, $\delta_m=\Delta d_mM_m$, $\tau_{l''}=-
\epsilon_{l''}\Delta\theta_{l''}$ $(l''\in E)$, satisfies the rows
\textup{(S1)}--\textup{(S7)} of $\Sigma^{\rm sh}(\kappa^{\rm sh}(w),f)$ up to a
total $\ell_1$-violation $\mathrm{viol}(\delta,\tau)\le C_f\mathrm{Des}^2K_gt/u$.
\end{lemma}

\begin{proof}
(S1) at strict non-peaks and at peaks (in the form $\tau_{l''}\ge-K_gt$) is
the one-sided pinning of class-G carriers, Lemma~\ref{lem:II-diagpin}.
(S2): by [I,~(16)], at a peak $k=k(l'')$ of block $m$ with $\varsigma_k=
\mathrm{vs}_{l''}$, $-\Delta\theta_{l''}=\mathrm{vs}_{l''}\lambda_{l''}(\delta_m+
e_k)$ with $0\le e_k\le t/(\sigma_m|\alpha_m(k)|)=t/(\lambda_{l''}\mu_k)$; at
$l''\in G_{\rm pk}$ (relative margin $\ge u$) $\mu_k\ge c_f\Phi_{l''}u\ge c_fu/
\mathrm{Des}$, so $|\tau_{l''}-\epsilon_{l''}\mathrm{vs}_{l''}\lambda_{l''}\delta_m|
\le C_f\mathrm{Des}\,t/u$.  (S3): $\Delta B1_{F^c}=V(\tau)+e_0$ with $V(\tau)=
\sum_{l''\in E}\epsilon_{l''}\tau_{l''}u_{l''}1_{F^c}$ and $\|e_0\|_1\le K_gt+
4b^2/t$ (Lemma~\ref{lem:II-diagpin}); the switching budget [I,~Lemma~8.24]
and the cost of the closed rooms (Step~1 of the proof of
Proposition~\ref{prop:II-transplant}) bound
$\sum_{j\notin F}\phi_{z^\#_j}(\Delta B(j))$ by $2t/q_0+CK_gt$, so by
[I,~Lemma~8.23(c)] the cost of $V(\tau)$ is at most $2t/q_0+CK_gt+2\|e_0\|_1$;
at a contact $\phi_z(x)=2(zx)_-$, and at a free coordinate of $\kappa$,
$\phi_{z_j}(x)\ge(1-|z_j|)|x|\ge u|x|$.  (S4): by [I,~(17)], $\delta_m=(mC_m)^{-1}
\sum_k\Phi_m(k)w_m(k)\Delta\theta_{k,m}+r_m$, $|r_m|\le2t/\sigma_m$, and for
$l''\in E$ the term equals $-q_{l''}\tau_{l''}$; class-R and fine carriers
contribute at most $(K_gt+6t^2)/(mC_m)$ (Lemma~\ref{lem:II-diagpin},
Lemma~\ref{lem:II-GW}(GW3)), and carriers of $N_0$ contribute $|q_{l''}
\tau_{l''}|\le b\Phi M6\lambda/(mCt)\le t^3$ by the box bound.  (S5), (S6):
Lemma~\ref{lem:II-sources} (an upper source gives $\delta_m\le K_dt$, a
lower one $\delta_m\ge-K_dt$, $K_d\le C_f\mathrm{Des}^3/u$).  (S7): for
$l''\in R_7$, [I,~Lemma~8.8] gives $\tau_{l''}\le\lambda_{l''}(3\gap/t-\Delta
d_m|w_m(k)|)$ with $\gap\le bM_m$, so the defect is $\le3\lambda b/t\le t^3$.
Summing, and using $t\ge T_{\rm lo}(w)$, $b\le T_{\rm lo}(w)^4$, gives the
bound.
\end{proof}

Any further relation satisfied by the decompositions up to the same order
may be added to $\Sigma^{\rm sh}$; this can only increase $\rho^{\rm sh}$.
$\rho^{\rm sh}(\kappa^{\rm sh},f)=0$ means that $\Sigma^{\rm sh}$ has exact or
asymptotically exact solutions with $\delta\ne0$: a \emph{$d$-constrained
coherent shift resonance}.

\begin{theorem}[Shift dichotomy]\label{thm:II-C1}
Let $T=T_{\mathrm{final}}$, $N\ge1$, and let $f\in S_{p^*}$ have $F$ finite and let $w$ be a clean sub-window of
level $l\ge l_f$ for $f$.  Then exactly one of the following holds.
\begin{enumerate}
\item[(I)] $\rho^{\rm sh}(\kappa^{\rm sh}(w),f)\ge u(w)$.  Then for every
$g\in\Cset(f)$, every scale $t\in\mathcal W(w)$, every two-sided
decomposition at scale $t$ and every block $m$,
\[
 |\Delta d_m|M_m\le K_{\rm sh}t,\qquad K_{\rm sh}:=C_f\mathrm{Des}^2K_g/u(w)^2
 \le C_f\mathrm{Des}^3/u(w)^3 .
\]
\item[(II)] $\rho^{\rm sh}(\kappa^{\rm sh}(w),f)\le b(w)$.
\end{enumerate}
If $I_{\rm up}(w)\cup I_{\rm lo}(w)=\emptyset$, then $\rho^{\rm sh}\ge1$, so
case \textup{(I)} holds.
\end{theorem}

\begin{proof}
The rate object $\rho^{\rm sh}(\kappa^{\rm sh}(w),\cdot)$ belongs to
$\mathcal O(l)$, so its value at $f$ is tiny or robust
(Theorem~\ref{thm:II-clean}).  $\mathrm{viol}$ is positively homogeneous of
degree one, so $\mathrm{viol}(\delta,\tau)\ge\|\delta\|_1\rho^{\rm sh}$ for every
$(\delta,\tau)$.  In case (I), Lemma~\ref{lem:II-R17} gives $\|\delta\|_1\le
C_f\mathrm{Des}^2K_gt/u^2$.  If $I_{\rm sh}=\emptyset$, (S5) forces
$\mathrm{viol}\ge\|\delta\|_1$.
\end{proof}

In case (I) the shift is pinned with a constant of the form $C_f\mathrm{Des}^C
u^{-C}$, which the window arithmetic absorbs (Lemma~\ref{lem:II-GW}(GW2));
every use of a shift bound in Section~\ref{sec:II-exact} holds at $w$ with
this constant.

\smallskip
We use the \emph{combinatorial shift cost} of
Definition~\ref{def:II-constants}(e).  For the enlarged shift pattern
$\kappa^{\rm sh}=\kappa^{\rm sh}(w)$ of $f$ at $w$ take the signs
$\mathrm{vs}_{l''}:=\sgn w_{m(l'')}(k(l''))$ $(l''\in G_{\rm pk})$ of
(O9)(S2), and write $c^{\rm comb}(\delta;\kappa^{\rm sh}):=c^{\rm comb}(\delta;
\kappa^{\rm sh},\mathrm{vs})$, $c^{\rm comb}_*(\kappa^{\rm sh}):=c^{\rm
comb}_*(\kappa^{\rm sh},\mathrm{vs})$, with $G_{\rm np}=E\setminus G_{\rm
pk}$ (dropped carriers included) and $\Pi_m$, $L$, $\mathrm{coef}_{l''}$ as
there.  Since $l\ge l_f\ge\max F$ and $F'=F\cap T(l)$, we have
$F\subseteq[1,l]$, so $1_{F^c}=1$ on $T(l)\setminus F'$ and on every set
$S_{l''}\setminus([1,l]\cup T(l))$; hence these design quantities coincide
with the costs computed with $\Pi_m1_{F^c}$ in place of $\Pi_m$, which is
the form in which they arise from the decompositions below.  By
Definition~\ref{def:II-constants}(e), $c^{\rm comb}_*(\kappa^{\rm sh})>0$
implies $c^{\rm comb}_*(\kappa^{\rm sh})\ge\delta_{\rm sh}(l)$.

\begin{corollary}[Combinatorial shift cost]\label{cor:II-C11}
In the situation of Theorem~\textup{\ref{thm:II-C1}}, write $c_*:=c^{\rm
comb}_*(\kappa^{\rm sh}(w))$.
\begin{enumerate}
\item[(a)] For every $(\delta,\tau)$: if $c_*>0$, then $\|\delta\|_1\le(C\mathrm{Des}/c_*
+2)\,\mathrm{viol}(\delta,\tau)$.  Consequently $c_*>0$ implies $\rho^{\rm sh}\ge
\delta_{\rm sh}(l)/(C\mathrm{Des})>b(w)$, i.e.\ case \textup{(I)}.
\item[(b)] In case \textup{(II)}, $c_*=0$, and the minimum is attained: there
are $\delta$ in the sign-constrained unit sphere of $\R^{I_{\rm sh}}$ and
$x\ge0$ such that $\sum_m\delta_m\Pi_m+\sum_{G_{\rm np}}x\epsilon u$ is
admissible for the contact pattern on $T(l)$ and has no sign defect on the
natural signature sets: an \emph{exact combinatorial coherent resonance}.
\end{enumerate}
\end{corollary}

\begin{proof}
(a) Given $(\delta,\tau)$, take $x:=(\tau)_+$ on $G_{\rm np}$.  On the natural
part of $S_{l''}$, $\epsilon_{l''}\mathrm{coef}_{l''}=x_{l''}\ge0$ for $l''\in
G_{\rm np}$, and for $l''\in G_{\rm pk}$ its sign defect is at most
$(\tau_{l''})_-+|\tau_{l''}-\epsilon\mathrm{vs}\lambda\delta_m|$, i.e.\ (S1) and
(S2) defects.  On $T(l)$, $|L(j)-V_j(\tau)|$ is bounded by $\max|u|$ times
the (S2) defects and $\sum_{G_{\rm np}}(\tau)_-$, and $2(\mathrm{type}(j)V_j)_-$,
$|V_j|$ are (S3) defects.  Hence $c^{\rm comb}(\delta_{I_{\rm sh}};\kappa^{\rm
sh})\le C\mathrm{Des}\,\mathrm{viol}(\delta,\tau)$.  Let $\delta'$ be the
$\ell_1$-nearest point of the sign cone to $\delta_{I_{\rm sh}}$; $\|\delta'-
\delta_{I_{\rm sh}}\|_1$ is at most the (S6) violation, and $c^{\rm comb}(\cdot;
\kappa^{\rm sh})$ is $C\mathrm{Des}$-Lipschitz in $\delta$ ($L$ is linear in
$\delta$ with $\|\Pi_m\|_1\le1$, $\phi_z$ is $2$-Lipschitz, and the defect terms
are $\lambda$-Lipschitz), so $c_*\|\delta'\|_1\le c^{\rm comb}(\delta')\le
C\mathrm{Des}\,\mathrm{viol}$, and $\|\delta\|_1\le\|\delta'\|_1+(\text{(S5) and
(S6) violations})$.  If $c_*>0$ then $c_*\ge\delta_{\rm sh}(l)\ge1/\mathrm{Des}$
(Proposition~\ref{prop:II-bank}(e)) and $b(w)\le T_{\rm lo}(w)^4\le2^{-4\mathrm{
Des}}$.  (b) follows from (a) and compactness.
\end{proof}

So a shift-pinning rate that is positive but smaller than every design
quantity cannot occur at a clean sub-window; only an exact combinatorial
resonance together with a tiny $\rho^{\rm sh}$ blocks the pinning.

\begin{remark}[Two independent pinning criteria]\label{rem:II-twocrit}
Lemma~\ref{lem:II-shiftpin} pins the shift under (SH$_w$), i.e.\ when $c_{\pi(w)}(f)\ge
u(w)$ or $I_{\rm sh}(w)=\emptyset$.  Neither of $c_{\pi(w)}\ge u$ and
$\rho^{\rm sh}\ge u$ is known to imply the other: $c_{\pi(w)}$ uses $z$ on all
of $F^c$ and charges terms proportional to tiny rooms that the rows
(S1)--(S7), written for the closed pattern, do not see (for the actual
decomposition these terms are $\le6b\lambda/t$, but a near-solution of
$\Sigma^{\rm sh}$ may have large $\|\tau\|_1$); conversely $c_{\pi(w)}$ sees
only the peak traces and zero cost, not the $d$-identity (S4), so $\rho^{\rm
sh}$ can be robust while $c_{\pi(w)}$ is tiny.  Both are therefore kept.
\end{remark}

The recovery engines for data that follow an unpinned shift are the
following variants of Theorem~\ref{thm:II-E}.

\begin{theorem}[Windowed recovery with shifted data]\label{thm:II-Eshift}
Let $f$, $g$, $\rho$, $\eta_0$, $f_j$, $(T_j,n_j)$, $K_j$, $A_{2,j}$,
$\gamma_j$ be as in Theorem~\textup{\ref{thm:II-E}}, and suppose all its
hypotheses hold except that the two-piece data $(b^\pm_{j,t},\omega^\pm_{j,t})$
at $f_j$ are not required to be $d$-neutral.
\begin{enumerate}
\item[(a)] If $\Delta d^{(j)}_m:=d^{(j)}_m(\omega^-_m)-d^{(j)}_m(\omega^+_m)\ge0$
for every $m$, $j$ and $t$, then $(f,\rho g)\in\cl\NA((c_0,p),\ell_2^2)$.
\item[(b)] If for every $j$ there is a set $I_-^{(j)}$ of blocks satisfying
the scrambling condition \textup{(SC)} at $f_j$ \textup{(}[I,~Definition~7.16]\textup{)}
such that $\{m:\Delta d^{(j)}_m<0\}\subseteq I_-^{(j)}$ for every $t$, then
$(f,\rho g)\in\cl\NA((c_0,p),\ell_2^2)$.
\end{enumerate}
In both cases, if the hypotheses hold for every $\rho<1$, then $(f,g)\in
\cl\NA((c_0,p),\ell_2^2)$.
\end{theorem}

\begin{proof}
In the proof of Theorem~\ref{thm:II-E}, $d$-neutrality is used only in the
verification of (e) of Theorem~\ref{thm:II-ERT}: (i) it is preserved by the
averaging over $t\in\mathcal S_j$, and (ii) [I,~Corollary~7.21] is applied at
$f_j$ to the averaged data.  Lemma~\ref{lem:II-U} is applied to each side
separately and does not involve $\Delta d$.  (a) $d^{(j)}_m$ is linear, so
the averaged data have $\Delta\bar d^{(j)}_m$ equal to the average of the
$\Delta d^{(j)}_m$, hence $\ge0$; [I,~Corollary~7.21] needs exactly $\Delta
d_m\ge0$ and $\kappa_w\le1$.  (b) The averaged data have $\{m:\Delta\bar
d^{(j)}_m<0\}\subseteq I_-^{(j)}$, and (SC) passes to subsets of $I_-^{(j)}$
(the maximum over fewer blocks is smaller).  The data $\rho\bar D_j$ of
$\rho\bar g_j$ have $\kappa_w\le1$, so [I,~Theorem~7.18] at $f_j$ (finite
$I$, finite $F_j$) gives $(f_j,\rho\bar g_j)\in\cl\NA$ whenever $\rho\bar
g_j\in\Cset(f_j)$.
\end{proof}

\begin{lemma}[Averaging over a subset of scales]\label{lem:II-subset}
Theorems~\textup{\ref{thm:II-ERT}}, \textup{\ref{thm:II-E}} and
\textup{\ref{thm:II-Eshift}} remain valid if, for each $j$, the set
$\mathcal S(T_j,n_j)$ is replaced by a subset $\mathcal S'_j\subseteq\mathcal
S(T_j,n_j)$ with $n'_j:=|\mathcal S'_j|$, all averages are taken over
$\mathcal S'_j$, and the conditions on $n_j$ are replaced by
$n'_j\ge48\rho^2K_j/(c_j(1-\rho^2))$ and $K_jT_j/n'_j\to0$
\textup{(}conditions \textup{(d)} of Theorem~\textup{\ref{thm:II-ERT}} and
\textup{(E-e)} unchanged\textup{)}.
\end{lemma}

\begin{proof}
In the proof of Theorem~\ref{thm:II-ERT} the set of scales enters at three
places.  For $|r|\le r_0$: the set $\mathcal S_r$ of scales $t\in\mathcal
S'_j$ with $ct<\rho|r|$ satisfies $\sum_{\mathcal S_r}t<2\rho|r|/c$ (dyadic
scales), so the averaged error is $\le2\rho^2Kr^2/(cn'_j)\le(1-\rho^2)r^2/24$;
if $\mathcal S_r\ne\emptyset$, then $\rho|r|>c\min\mathcal S'_j\ge2cT_j2^{-n_j}$,
which is all that the use of (d) requires.  For $|r|>r_0$: $\frac1{n'_j}
\sum_{\mathcal S'_j}t\le2T_j/n'_j$, and $2\rho K_jT_j/n'_j\to0$.  Finally
$p^*(\rho\bar g_j-\rho g)\le2\rho K_jT_j/n'_j\to0$.  Averaging over $\mathcal
S'_j$ preserves exactness, side conditions, the sign of $\Delta d_m$ and a
common shift.
\end{proof}

\begin{theorem}[Master Theorem~III$'$]\label{thm:II-MTIII}
Let $T=T_{\mathrm{final}}$ with base $U_{\mathrm{final}}$, let $N\ge1$, and let $f\in S_{p_N^*}$ have finite base support.  Suppose that
for infinitely many levels $l$ some clean sub-window $w$ of level $l$
satisfies
\[
 \rho^{\rm sh}(\kappa^{\rm sh}(w),f)\ge u(w)\qquad\text{or}\qquad
 \text{\textup{(SH$_w$)}}.
\]
Then $f\in\Rec$, i.e.\ $(f,g)\in\cl\NA((c_0,p_N),\ell_2^2)$ for every
$g\in\Cset(f)$.
\end{theorem}

\begin{proof}
Either alternative gives $|\Delta d_m|M_m\le Kt$ for every block $m$ and every
scale $t\in\mathcal W(w)$, with $K\le C_f\mathrm{Des}^3/u^3$
(Theorem~\ref{thm:II-C1}(I)) resp.\ $K\le C_fl\,\mathrm{Des}^3/u^2$
(Lemma~\ref{lem:II-shiftpin}).  Master Theorem~II uses (SH$_w$) only through this
bound, and (VR$_w$) only in the step of Proposition~\ref{prop:II-transplant} that
Corollary~\ref{cor:II-B1} replaces.  Hence the proof of
Theorem~\ref{thm:II-master2} applies along the infinitely many sub-windows in
question, with the companion of Theorem~\ref{thm:II-B} and the window
constants of Corollary~\ref{cor:II-B1}; the remaining window arithmetic is
that of Theorem~\ref{thm:II-master2}.
\end{proof}

\begin{definition}[The residual (C$^*$)]\label{def:II-Cstar}
Let $T=T_{\mathrm{final}}$ and $N\ge1$.  A first row $f\in S_{p_N^*}$ with finite base support is a \emph{(C$^*$) row}
if for all but finitely many levels $l$, \emph{every} clean sub-window $w$
of level $l$ satisfies
\[
 \rho^{\rm sh}(\kappa^{\rm sh}(w),f)\le b(w)\qquad\text{and}\qquad
 c_{\pi(w)}(f)\le b(w);
\]
in particular $I_{\rm up}(w)\cup I_{\rm lo}(w)\ne\emptyset$ at these
sub-windows.
\end{definition}

\begin{corollary}\label{cor:II-Cstarnec}
Let $T=T_{\mathrm{final}}$ and $N\ge1$.  If $f\in S_{p_N^*}$ has finite base support and $f\notin\Rec$, then $f$ is a
\textup{(C$^*$)} row.
\end{corollary}

\begin{proof}
At a clean sub-window both $\rho^{\rm sh}$ and $c_{\pi(w)}$ are tiny or robust,
and (SH$_w$) fails exactly when $I_{\rm sh}(w)\ne\emptyset$ and
$c_{\pi(w)}\le b(w)$.  Apply Theorem~\ref{thm:II-MTIII}.
\end{proof}

In words: a (C$^*$) row has, at every clean sub-window of every large
level, a block lacking an upper or a lower shift source in which the shift
traces of the robust peaks are compensated, exactly at the combinatorial
level (Corollary~\ref{cor:II-C11}(b)) and up to $b(w)$ in the $d$-identity, by
zero-cost switching of class-G strict non-peaks.

\subsection{Structure of shifted data}\label{subsec:II-structure}

By Corollary~\ref{cor:II-Cstarnec} the only first rows with finite base
support that may fail to be in $\Rec$ are (C$^*$) rows.  This subsection
records what exact two-piece data with a nonzero shift look like.  We use
the notation $W_\Delta=\sum_m\Delta d_m\psi_m$, $\psi_m=R_m^*w_m$, and
$\Sigma(\omega)$ of Subsection~\ref{subsec:II-rigid} (data convention).

\begin{lemma}[Shape of a source-deficient block]\label{lem:II-Dprime}
Let $w$ be a clean sub-window of level $l\ge l_f$ for $f$ \textup{(}$F$
finite\textup{)} and $m$ a block.  Block $m$ has a coarse peak with
$\varrho\ge1+u(w)$.  If block $m$ has no upper source at $w$, then every
coarse peak of block $m$ is a class-G peak of swallowing type.  If block $m$
has no lower source at $w$, then every coarse peak of block $m$ with
$\varrho\ge1+u(w)$ is a class-G peak of anti type.
\end{lemma}

\begin{proof}
The first claim is Lemma~\ref{lem:II-donorpeak}.  At a clean sub-window every coarse
carrier is of class R or of class G.  A class-R peak is an upper source
(U1), and a lower source (L1) if $\varrho\ge1+u$; a class-G anti-type peak is
an upper source (U2); a class-G swallowing-type peak with $\varrho\ge1+u$ is a
lower source (L2).  Peaks have $|w_m(k)|=M_m\ne0$, so their type is defined.
\end{proof}

Thus in a source-deficient block ($m\in I_{\rm up}(w)\cup I_{\rm lo}(w)$)
of a (C$^*$) row every robust coarse peak has, at that clean sub-window, a
tiny room on its natural signature set,
and all of them have one type: swallowing type in an upper-deficient block
(``configuration (i)''; there the data shift is negative, $\Delta d_m<0$ in
the data convention) and anti type in a lower-deficient block
(``configuration (ii)'', $\Delta d_m>0$).

\begin{proposition}[Rigidity of shifted data]\label{prop:II-C3}
Let $(b^\pm,\omega^\pm)$ be two-piece data at a first row $f$ with $F$
finite, for some functional.
\begin{enumerate}
\item[(a)] For $j\notin\Sigma(\omega)$: $W_\Delta(j)=0$ if $j$ is a free
coordinate, and $z_jW_\Delta(j)\le0$ if $j$ is a contact.
\item[(b)] Let $l$ be a carrier of block $m$ with $l\notin\Omega$ \textup{(}the
carriers in the supports of $\omega^\pm$\textup{)}, $S_l\cap F=\emptyset$,
$w_m(k(l))\ne0$ and $\Delta d_m\ne0$.  Then there is $s_0$ such that every
$s\in S_l\setminus\Sigma(\omega)$ with $s\ge s_0$ is a contact with
$z_s=-\sgn(\Delta d_m)\sgn w_m(k(l))$.
\item[(c)] For every coordinate $j$ and every block $m$, $u_{k,m}(j)\ne0$ for
infinitely many $k$ \textup{(}any admissible $T$\textup{)}.  Hence at every free
coordinate $j\notin\Sigma(\omega)$ the identity $W_\Delta(j)=0$ of \textup{(a)}
is an infinite series over carriers of every block $m$ with $\Delta d_m\ne0$.
\end{enumerate}
\end{proposition}

\begin{proof}
(a) is Lemma~\ref{lem:II-sandwich}(b),(c).  (b) For $s\in S_l$, $W_\Delta(s)=
\Delta d_m\lambda_lw_m(k(l))v_l(s)+\sum_{l''>l}\Delta d_{m(l'')}\lambda_{l''}
w(l'')y_{l''}(s)/n_{l''}$: earlier carriers do not meet $S_l$ (allowedness
(a) of Definition~\ref{def:II-Tfinal}), and signature sets are disjoint.  A
later carrier $l''$ meets $s\in S_l$ only through its target, and then
$c_{l''}\le\iota_l2^{-2s}c_l\delta_l/2$ by (W4); with $c_{l''+1}\le c_{l''}/4$
the later terms total at most $C\max_{m'}|\Delta d_{m'}|2^{-2s}c_l\delta_l$,
while the own term has modulus $|\Delta d_m||w_m(k(l))|\lambda_l\delta_l2^{-s}/
n_l$.  For $s\ge s_0$ the own term dominates, so $W_\Delta(s)\ne0$ with the
sign of $\Delta d_m w_m(k(l))$, and (a) forces $s$ to be a contact with the
stated sign.  (c) By (T-d), $e_j^*/q^*(e_j^*)\in S_{q^*}$ is a $q^*$-limit of
a subsequence of $(u_{k,m})_k$, so $u_{k,m}(j)\to1/q^*(e_j^*)\ne0$ along it.
\end{proof}

\begin{remark}\label{rem:II-C3}
(a) Proposition~\ref{prop:II-C3}(c) shows that shifted exact data exist only
at rows satisfying one exact cancellation identity per free coordinate
outside the finite set $\Sigma(\omega)$.  That this is a non-generic
property at partial contact is a \emph{heuristic} reading only; it is not
claimed as a theorem.
(b) The identities are void at maximal contact.  There the dichotomy is as
follows.  If $z=\epsilon_0$ is constant on $\N\setminus F$, then $f$ is not a
(C$^*$) row (Proposition~\ref{prop:II-constcontact}).  If $|z|=1$ on
$\N\setminus F$ with signs that are not eventually constant, the cited
results do not exclude (C$^*$): the type of a peak $l$ is $\epsilon_l\sgn
w(k(l))$, and signs on the signature sets can be prescribed carrier by
carrier so that every coarse peak of a block is \emph{self-aligned}
($z=\sgn u_l(\zh)$ on $S_l$, hence swallowing type), while some strict
non-peak of robust relative position is anti-aligned (anti type, $q<0$);
then (U1), (U2), (U3) fail and the block lacks an upper source
(Subsection~\ref{subsec:II-SA}).
\end{remark}

\begin{proposition}[Constant-sign maximal contact]\label{prop:II-constcontact}
Let $T=T_{\mathrm{final}}$, $N\ge1$, $f\in S_{p_N^*}$ with
$F$ finite, and suppose $z_j=\epsilon_0$ for all $j\notin F$, for a fixed
sign $\epsilon_0$.  Then there is $l_1$ such that at every clean sub-window
of every level $l\ge l_1$ every block has an upper source of type
\textup{(U2)} and a lower source of type \textup{(L2)}.  Hence
$I_{\rm up}(w)=I_{\rm lo}(w)=\emptyset$, $\rho^{\rm sh}(\kappa^{\rm sh}(w),f)\ge1$,
$f$ is not a \textup{(C$^*$)} row, and $f\in\Rec$.
\end{proposition}

\begin{proof}
Since $\zh\ne0$ there is $y\in S_{q^*}$ with $y(\zh)=:2c>0$.  By (T-d), every
block $m$ has carriers $k_i\to\infty$ with $q^*(u_{k_i,m}-y)<c$, and carriers
$k'_i\to\infty$ with $q^*(u_{k'_i,m}+y)<c$; then $u_{k_i,m}(\zh)>c$ and
$u_{k'_i,m}(\zh)<-c$, because $|u(\zh)-y(\zh)|\le q^*(u-y)$.  As $\Phi_m(k)\to
0$, their relative positions $\varrho=m|u(\xi)|/(\Phi\cdot m\vartheta_m)$ tend to
infinity; in each block fix one carrier $l_+$ of the first kind and one
carrier $l_-$ of the second kind with $S_{l_\pm}\cap F=\emptyset$ and
$\varrho_{l_\pm}\ge2$ (so both are peaks).  Their signature sets carry $z\equiv
\epsilon_0$, so their rooms (rate object (O1)) vanish and they are of class
G with swallowing sign $\epsilon_0$ at every clean sub-window of every level
$l\ge l_1:=$ the largest of these carriers over all blocks.  Their types are $\epsilon_0\sgn u(\zh)$, i.e.\
one is of swallowing type and one of anti type, and both have $\varrho\ge2\ge
1+u(w)$.  The anti-type one is an upper source (U2), the swallowing-type one
a lower source (L2).  The remaining claims follow from
Theorem~\ref{thm:II-C1} and Theorem~\ref{thm:II-MTIII}.
\end{proof}

\begin{remark}\label{rem:II-constcontact}
The first part of this proof (existence of carriers $l_\pm$ in every block
with $\varrho\ge2$, $S_{l_\pm}\cap F=\emptyset$ and $z\equiv\epsilon_0$ on
their signature sets) uses only (T-d) and the disjointness of the signature
sets, so it holds for every admissible operator with pairwise disjoint
signature sets.  The conclusions about clean sub-windows, sources,
$\rho^{\rm sh}$, (C$^*$) and $\Rec$ use the objects and theorems of
$T_{\mathrm{final}}$.
\end{remark}

\begin{proposition}[Oscillating profiles force shifted data]\label{prop:II-C4}
Let $f'$ be a first row with $F'$ finite, $l$ a peak carrier of block $m$
of $f'$ with $\mathrm{vs}:=\sgn w'_m(k(l))$, and let $h$ carry two-piece data
$(b^\pm,\omega^\pm)$ at $f'$.  Let $S_1,S_2\subseteq S_l\setminus\Sigma(\omega)$ be
finite, nonempty, disjoint sets of contacts with $z'_s=\epsilon_l$ for $s\in
S_1\cup S_2$.  Put $V_i:=\sum_{s\in S_i}v_l(s)$, $A_i(h):=V_i^{-1}\sum_{s\in
S_i}\epsilon_lh(s)$, $d^\pm_{m'}:=d'_{m'}(\omega^\pm_{m'})$ and
\[
 \mathrm{tgt}:=\max_{m',\pm}|d^\pm_{m'}|\sum_{s\in S_1\cup S_2}\sum_{l''>l}
 \lambda_{l''}|u_{l''}(s)| .
\]
If $\epsilon_l\mathrm{vs}=+1$, then $A_i(h)\in[-d^+_m\lambda_lM'_m-\mathrm{tgt}/V_i,\
-d^-_m\lambda_lM'_m+\mathrm{tgt}/V_i]$; if $\epsilon_l\mathrm{vs}=-1$ the same
holds with $d^+_m$ and $d^-_m$ interchanged.  In both cases
\[
 |A_1(h)-A_2(h)|\le|\Delta d_m|\lambda_lM'_m+\mathrm{tgt}\,(1/V_1+1/V_2).
\]
Consequently, if $g$ is a functional and $g_t$ carry two-piece data at $f'$
with $\|g-g_t\|_1\le Kt$, then
\[
 |\Delta d_m(g_t)|\ge\big(|A_1(g)-A_2(g)|-(Kt+\mathrm{tgt})(1/V_1+1/V_2)\big)/
 (\lambda_lM'_m).
\]
\end{proposition}

\begin{proof}
At $s\in S_l\setminus\Sigma(\omega)$ the $+$ representation gives $h(s)=b^+(s)-
\sum_{m'}d^+_{m'}\psi'_{m'}(s)$, and $\psi'_{m'}(s)$ consists of the own term
$\lambda_lw'_m(k(l))v_l(s)=\lambda_l\mathrm{vs}M'_mv_l(s)$ (for $m'=m$) and
target terms of carriers $l''>l$: earlier targets avoid $S_l$ by
allowedness (a), and signature sets are disjoint.  Side admissibility gives
$\epsilon_lb^+(s)=z'_sb^+(s)\ge0$, so $\epsilon_lh(s)\ge-d^+_m\lambda_l\epsilon_l
\mathrm{vs}M'_mv_l(s)-\sum_{m'}|d^+_{m'}|\sum_{l''>l}\lambda_{l''}|u_{l''}(s)|$;
sum over $S_i$ and divide by $V_i$.  The $-$ side gives the other inequality
with $\epsilon_lb^-(s)\le0$.  The last claim follows from $|A_i(g)-A_i(g_t)|\le
\|g-g_t\|_1/V_i$.
\end{proof}

The decompositions of a mate $g$ at a scale $t$ satisfy the same sandwich up
to the budget, so the profile $\epsilon_lg(s)/v_l(s)$ on the signature set of a
robust class-G peak may oscillate inside an interval of length
$|\Delta d_m|\lambda_lM_m$ (decomposition convention).
Proposition~\ref{prop:II-C4} shows that such an oscillation is seen by every
approximating datum at every row $f'$ sharing the peak: $d$-neutral data
($\Delta d=0$) cannot follow a mate whose profile oscillates by more than
$(Kt+\mathrm{tgt})(1/V_1+1/V_2)$.  In case (II) of Theorem~\ref{thm:II-C1}
the exact-data route therefore needs shifted data.

\begin{lemma}[Completion of the fine tail]\label{lem:II-C5}
Let $f^\#\in S_{p^*}$ have finite base support $F^\#$, let $l$ be a level,
$J_{\rm fine}:=\N\setminus(F^\#\cup T(l)\cup\bigcup_{l''\le l}S_{l''})$, and
$\Delta\in\R^I$.  For $z'\in[-1,1]^{J_{\rm fine}}$ let $f(z'):=f[a^\#,z^\#(z')]$,
where $z^\#(z')$ equals $z'$ on $J_{\rm fine}$ and $z^\#$ elsewhere, and put
$W(z'):=\sum_m\Delta_mR_m^*w_m(z')$, $w_m(z')$ the block functionals of $f(z')$.
Then there is $z'^*\in[-1,1]^{J_{\rm fine}}$ with
\[
 W(z'^*)(j)=0\ \text{ if }|z'^*_j|<1,\qquad z'^*_jW(z'^*)(j)\le0\ \text{ if }
 |z'^*_j|=1\qquad(j\in J_{\rm fine}).
\]
Moreover the values $u_{l''}(\zh)$ of all carriers $l''\le l$ are unchanged,
and, with $\varepsilon_l:=\sum_{l''>l}\lambda_{l''}$ and
$\max_m2\varepsilon_l\le c_{f^\#}$, $p^*(f(z'^*)-f^\#)\le C_{f^\#}\varepsilon_l
\log(e/\varepsilon_l)$.
\end{lemma}

\begin{proof}
The data $(a^\#,z^\#(z'))$ are admissible ($J_{\rm fine}\cap F^\#=\emptyset$),
so $f(z')$ is defined (Subsection~\ref{subsec:II-comp}).  The set $\mathcal K:=
[-1,1]^{J_{\rm fine}}$ with the product topology is a compact, convex,
metrizable subset of the locally convex space $\R^{J_{\rm fine}}$.  Define
$\Psi(z')_j:=\mathrm{clamp}_{[-1,1]}(z'_j-W(z')(j))$.  Continuity: $\zh(z')=
z^\#(z')+Ue^\#$ depends coordinatewise continuously on $z'$; $R_m^{**}\zh(z')=
(\lambda_{k,m}u_{k,m}(\zh(z')))_k$ is $\ell_1$-continuous by dominated
convergence ($\sum_k\lambda_{k,m}\|u_{k,m}\|_1<\infty$, $\|\zh(z')\|_\infty\le2$);
it is never $0$; the norming functional of the smooth norm $|\cdot|_m$ is
unique ([I,~Lemma~1.7]), so the duality map is norm-to-weak$^*$ continuous,
and each $w_m(z')(k)$ is continuous and bounded by $1$; finally
$W(z')(j)=\sum_m\Delta_m\sum_k\lambda_{k,m}w_m(z')(k)u_{k,m}(j)$ is continuous by
dominated convergence.  By the Schauder--Tychonoff theorem $\Psi$ has a fixed
point $z'^*$, and $\mathrm{clamp}(z_j-W_j)=z_j$ means $W_j=0$ if $|z_j|<1$,
$W_j\le0$ if $z_j=1$ and $W_j\ge0$ if $z_j=-1$.  Carriers $l''\le l$ have
$\supp u_{l''}\subseteq T(l)\cup S_{l''}$, which misses $J_{\rm fine}$, so their
values are unchanged; with $\delta:=\zh(z'^*)-\zh^\#$ one has $u_{l''}(\delta)=0$
for $l''\le l$ and $|u_{l''}(\delta)|\le2$ otherwise, so $\Delta_m(\delta)\le2
\varepsilon_l$, $c(\delta)\le C\varepsilon_l\log(e/\varepsilon_l)$, and
Lemma~\ref{lem:II-cost} gives the cost bound.
\end{proof}

\begin{remark}[One ray only]\label{rem:II-C5}
$W$ is linear in $\Delta$, so the completion for $\Delta$ also completes
$c\Delta$, $c>0$.  It does not complete two shift vectors $\Delta^{(1)}$,
$\Delta^{(2)}$ spanning a two-dimensional cone at the same time: at a free
coordinate both $W_{\Delta^{(1)}}$ and $W_{\Delta^{(2)}}$ would have to vanish.
If the shift of the decompositions is unpinned, its direction may change
from scale to scale inside a window, while the averaging of
Theorems~\ref{thm:II-E} and \ref{thm:II-Eshift} uses one companion for all
scales; this is resolved in Subsection~\ref{subsec:II-absorb} by normalizing
all scales of a class onto one shift ray (Proposition~\ref{prop:II-ray}).
Lemma~\ref{lem:II-C5} also does not act on the coarse coordinates $T(l)\cup
\bigcup_{l''\le l}S_{l''}$, where the fine contributions to $W_\Delta$ must be
absorbed exactly by the coarse data.
\end{remark}

For one block the existence of shifted exact data is a property of the row
and of the set of switching carriers alone.

\begin{corollary}[One block: shifted exact data]\label{cor:II-A}
Let $I=\{1\}$, $f\in S_{p^*}$ with $F$ finite, $\psi:=R_1^*w_1$, $\Omega\subseteq
Q$ finite, and $\Delta\ne0$.
\begin{enumerate}
\item[(a)] If two-piece data with $\Delta d=\Delta$ and block parts supported
in $\Omega$ exist for some functional, then outside $\Sigma(\omega)$: $\psi$
vanishes at every free coordinate, and $z_j\Delta\psi(j)\le0$ at every
contact.
\item[(b)] Conversely, let $\gamma\in c_{00}$ be supported in $\Omega$ with
$d(\gamma)=\Delta$ \textup{(}possible iff $w(k)\ne0$ for some $k\in\Omega$\textup{)},
and suppose that $v:=R_1^*\gamma-\Delta\psi$ is $z$-signed at every contact
and vanishes at every free coordinate off $F$ \textup{(}outside
$\Sigma(\omega)$ this is the condition of \textup{(a)}; on $\Sigma(\omega)\setminus
F$ it is a condition on $\gamma$\textup{)}.  Then for every $\omega^+$
supported in $\Omega$, every $\chi:K\to[0,1]$ and every $\beta$ supported in
$F$, the pairs $(b^+,\omega^+)$, $(b^-,\omega^-)$ with $\omega^-:=\omega^++\gamma$,
$b^+:=\beta+\chi v1_K$, $b^-:=b^+-v$ are two-piece data with $\Delta d=\Delta$
for $h:=b^++R_1^*(\omega^+-d(\omega^+)w_1)$.
\end{enumerate}
\end{corollary}

\begin{proof}
(a) is Lemma~\ref{lem:II-sandwich} with $W_\Delta=\Delta\psi$.  (b) $d(e_k)=
\Phi(k)^2w(k)/C$, so $d(\gamma)=\Delta$ is solvable iff some $w(k)\ne0$;
$\mathcal V(\gamma)=R_1^*(\gamma-\Delta w_1)=v$, and Lemma~\ref{lem:II-cert}(c)
applies.
\end{proof}

With several blocks $\psi$ is replaced by $W_\Delta$, and blocks with
$\Delta d_m=0$ contribute nothing at the coordinates they own: every
coordinate owned by a carrier of an unshifted block is a \emph{dead zone}
for the shifted blocks, where the sign of $W_\Delta$ is decided by later
carriers.

\begin{remark}[An earlier route to (C$^*$); sketch only]\label{rem:II-C6}
The following route was proposed for (C$^*$) rows and is recorded here
\emph{as a sketch}, not as a result.  Suppose that at infinitely many levels
a clean sub-window $w$ in case (II) has a coarse shifted exact cone at the
companion (the system $\Sigma^\#$ together with the shift columns $\Delta_m
R_m^*w^\#_m$ restricted to coarse coordinates and the coupling $\Delta_m=\sum
q^\#_l\tau_l$) with a Slater point of robust margin, and that every shifted
block has $\Delta_m\ge0$ (data convention), or $\Delta_m<0$ together with
(SC) at the companions.  Then one would transplant with shift columns,
complete the fine tail by Lemma~\ref{lem:II-C5}, absorb the coarse fine
perturbations by the Slater condition, and conclude by
Theorem~\ref{thm:II-Eshift}.  The gaps of this route are: (g1) exact
adjustment of the block scalars entering the coupling; (g2) a conversion
of near coordinates of coarse signature sets into support coordinates
preserving the values; (g3) (SC) at the companions in configuration~(i);
(g4) shift directions that vary with the scale when two or more blocks are
shifted (Remark~\ref{rem:II-C5}); (g5) completion and coarse absorption must
be solved jointly.  The conversion in (g2) is \emph{false} as stated for
super-decreasing signature weights: the achievable signed sums of the near
weights form a Cantor set, and the discrepancy is of the order of the first
converted weight.  Subsection~\ref{subsec:II-absorb} replaces (g2) by
zero-value absorber pairs, and, under the hypotheses of Master
Theorem~IV$'$ (Theorem~\ref{thm:II-MTIV}), closes (g3) in configuration (i),
(g4) and (g5); it reduces (g1) to one scalar per block, which is handled
there by the lever hypothesis (KN); whether a second, independent lever is
necessary is not decided (Remark~\ref{rem:II-onelever}).
\end{remark}

\subsection{Coherent shift resonance at block-tame rows; genericity}\label{subsec:II-SA}

In this subsection $T=T_{\mathrm{final}}$, whose features (SF$_\iota$),
(SF$^*$), (b$'$), (Z0) and the diagonal base are listed in
Proposition~\ref{prop:II-features}; the propositions on genericity hold for
every admissible $T$ with a diagonal base, as stated.  For $p=p_N$ the
\emph{present} carriers are $L_N:=\{l:m(l)\le N\}$.  For a diagonal base
with entries $\mu^\star_j$, $\nu=\|U^*a\|$ and $(Ue)_j=(\mu^\star_j)^2a_j/\nu$,
so $\zh_j=\sgn a_j+(\mu^\star_j)^2a_j/\nu$ on $F$ and $\zh_j=z_j$ off $F$.

\begin{definition}[Owners and the owner recursion]\label{def:II-owner}
The \emph{owner} of a coordinate $j$ is $o_N(j):=\min\{l\in L_N:u_l(j)\ne0\}$
(it exists by Proposition~\ref{prop:II-C3}(c) applied to block $1$); if $j\in
S_l$ with $l\in L_N$, then $o_N(j)=l$ by allowedness (a).  Given $a$ with
finite support $F$ and a sign $\epsilon_l$ for every $l\in L_N$ chosen
along the recursion below, the \emph{owner recursion} assigns $z$ off $F$
as follows.  Initially the coordinates of $F$ are \emph{assigned}.  At step
$l\in L_N$ (in increasing order) let
\[
 A_l:=\sum_{j\in\supp y_l\ \text{assigned}}y_l(j)\zh_j,\qquad
 B_l:=\sum_{j\in\supp y_l\ \text{unassigned}}|y_l(j)| ;
\]
then put $z_j:=\epsilon_l\sgn y_l(j)$ at the unassigned $j\in\supp y_l$ and
$z:=\epsilon_l$ on $S_l$, and declare these coordinates assigned.  The
\emph{owner sign} is $\epsilon_l:=\sgn A_l$ ($+1$ if $A_l=0$).
\end{definition}

At step $l$ the set $S_l$ is unassigned (earlier targets avoid $S_l$ by
allowedness (a), signature sets are disjoint, and $F\cap S_l=\emptyset$ if
$F\subseteq Z_0$); every coordinate off $F$ is eventually assigned; and the
data $(a,z)$ are admissible, so they define a unique first row.  When
$S_l\cap F=\emptyset$ and $z=\epsilon_l$ on $S_l$, the value formula
\begin{equation}\label{eq:II-valueformula}
 n_l\,u_l(\zh)=A_l+\epsilon_l(B_l+\delta_lH_l)
\end{equation}
holds, because $u_l(\zh)=u_l(z)+\ip{U^*u_l}{e}$ and $\ip{U^*u_l}{e}=\sum_{j\in
F}(\mu^\star_j)^2u_l(j)a_j/\nu$.  For the owner sign, $|n_lu_l(\zh)|=|A_l|+B_l+
\delta_lH_l$.  All owner notions are taken relative to $L_N$: with the
recursion over all carriers, a coordinate owned by an absent carrier would
receive that carrier's sign, while every quantity of $p_N$ is dominated
there by the first present carrier.

\smallskip
Let $F_\star:=\{3,5\}\subseteq Z_0$ and recall the target $y^*=(e_3^*-e_5^*)/
q^*(e_3^*-e_5^*)$ of Definition~\ref{def:II-data}(iv); put $c^*:=1/q^*(e_3^*-
e_5^*)$.  Since $y^*$ is supported in $Z_0$ it is allowed at every stage, so
every block has infinitely many carriers with target $y^*$.  For a block $m$
let $\delta^\circ_l:=\delta_lH_l/n_l$ as in Definition~\ref{def:II-Tfinal}.

\begin{theorem}[Self-aligned rows]\label{thm:II-SA}
Let $N\ge1$, $m_0\le N$, and let $l_c$ be the first carrier of block $m_0$.
Let $l_->l_c$ be a carrier of block $m_0$ with target $y^*$ and
$\delta_{l_-}H_{l_-}<c^*\mu^\star_5/2$, and let $0<\eta<\min\{c^*\mu^\star_5/2,\
n_{l_-}\theta_0\Phi_{l_-}/(2m_0)\}$, $\theta_0:=\lambda_{l_c}\delta^\circ_{l_c}/2$.
Then there are $a_3,a_5>0$ such that the first row $f_{\rm SA}\in S_{p_N^*}$
with $a=(a_3e_3^*+a_5e_5^*)/q^*(a_3e_3^*+a_5e_5^*)$ and with $z$ given by the
owner recursion over $L_N$, with owner signs except $\epsilon_{l_-}:=+1$,
has the following properties.
\begin{enumerate}
\item[(a)] Every $l\in L_N$ is exactly swallowed with sign $\epsilon_l$; for
$l\ne l_-$, $\sgn u_l(\zh)=\epsilon_l$ and $|u_l(\zh)|\ge\delta^\circ_l$.
\item[(b)] Every $l\in L_N\setminus\{l_-\}$ is a non-degenerate peak of
swallowing type; if $l$ is not the first carrier of its block, then
$\varrho_l\ge2^{10}$ and $\mu_l\ge q_0\delta^\circ_l/2$.
\item[(c)] $n_{l_-}u_{l_-}(\zh)=-\eta$; $l_-$ is a strict non-peak of anti
type with $\varrho_{l_-}<\frac12$, $\gap\ge M_{m_0}/2$ and $q_{l_-}<0$, and
$u_{l_-}1_{F^c}=v_{l_-}$ is $z$-signed.
\item[(d)] \textup{(Dominance.)}  For each of the vectors
\begin{gather*}
 W:=\sum_{l\in L_N}\lambda_l\epsilon_lu_l1_{F^c},\qquad
 \sum_{l\in L_N,\,l\le L}\lambda_l\epsilon_lu_l1_{F^c}\ \ (L\in\N),\\
 \sum_{l\in L_N}c_l\lambda_l\epsilon_lu_l1_{F^c}\ \ (c_l\in[\tfrac12,1]):
\end{gather*}
at every $j\notin F$ where the sum contains the term of $o_N(j)$, the sum is
nonzero with sign $z_j$.
\item[(e)] Every coordinate off $F$ is a contact; $Q_{m_0}=\{k(l_-)\}$ and
$Q_m=\emptyset$ for $m\ne m_0$; there is no degenerate peak; every block
satisfies $\mathrm{Scr}_m(\Upsilon s)=o(s)$ as $s\to0$ for every
$\Upsilon\ge1$; and $f_{\rm SA}$ satisfies \textup{(BT)}.
\item[(f)] There is $l_1$ such that at every clean sub-window $w$ of every
level $l\ge l_1$: block $m_0$ has no upper source and a lower source
\textup{(L2)}, every other block has both sources, and $c_{\pi(w)}(f_{\rm
SA})=0$.
\end{enumerate}
\end{theorem}

\begin{proof}
\emph{Step 1 (tuning).}  $\supp y^*=F_\star$ and $S_{l_-}$ is unassigned at step
$l_-$, so $n_{l_-}u_{l_-}(\zh)=y^*(\zh_{F_\star})+\delta_{l_-}H_{l_-}$, independently
of the recursion.  With $\zh_j=1+(\mu^\star_j)^2a_j/\nu$ on $F_\star$, $y^*(\zh_{F_\star})
=c^*((\mu^\star_3)^2a_3-(\mu^\star_5)^2a_5)/\nu$ runs continuously over
$(-c^*\mu^\star_5,c^*\mu^\star_3)$ as $a_3/a_5$ runs over $(0,\infty)$.  By the
choice of $l_-$ and $\eta$ and the intermediate value theorem there is a
ratio with $n_{l_-}u_{l_-}(\zh)=-\eta$; fix it.

\emph{Step 2 ((a)).}  $S_l\cap F=\emptyset$ because $F\subseteq Z_0$, so each
$l$ is exactly swallowed with sign $\epsilon_l$, and
\eqref{eq:II-valueformula} with the owner sign gives $|n_lu_l(\zh)|=|A_l|+B_l+
\delta_lH_l\ge\delta_lH_l$, $\sgn u_l(\zh)=\epsilon_l$.

\emph{Step 3 (threshold bound).}  For a block vector $\zeta=\hat\zeta_m$:
$\mathsf A(x;\zeta)\ge\|\zeta\|_1-x\|\Phi_m\|_2^2\ge\|\zeta\|_1-x/4$ and $\mathsf
B(x;\zeta)\le x\|\zeta\|_1$, so $\Xi(\|\zeta\|_1/2;\zeta)\ge(\frac{49}{64}-\frac12)
\|\zeta\|_1^2>0$ and $\bar\vartheta_m>\|\zeta\|_1/2$ (Lemma~\ref{lem:II-T}(a)).
Hence $M_m=\bar\vartheta_m/(\bar\vartheta_m+A_m)>\frac13$ for every block of
every first row, and $\bar\vartheta_{m_0}>\lambda_{l_c}|u_{l_c}(\zh)|/2\ge\theta_0$.

\emph{Step 4 ((b)).}  Let $l\ne l_-$ be a carrier of block $m$ and $x:=\mathsf
v_l=m|u_l(\zh)|/\Phi_l$.  For an earlier carrier $k$ of block $m$,
$\Phi_l/\Phi_k\le c_l/c_{l-1}\le\delta_l2^{-\min S_l-l-10}/(1+\|U\|)\le|u_l(\zh)|
2^{-l-9}/(1+\|U\|)$ by (W5) and $\delta^\circ_l\ge\frac45\delta_l2^{-\min S_l}$;
as $|u_k(\zh)|\le1+\|U\|$, the term $(|\zeta(k)|-x\Phi_k^2)_+=m\Phi_k(|u_k(\zh)|-
|u_l(\zh)|\Phi_k/\Phi_l)_+$ vanishes.  Later carriers contribute at most
$(1+\|U\|)\sum_{l''>l}\lambda_{l''}\le2^{-10}\lambda_l|u_l(\zh)|$ by (SF$_\iota$),
and the own term vanishes at $x=\mathsf v_l$.  So $\mathsf A(\mathsf v_l)\le
2^{-10}m\Phi_l|u_l(\zh)|<m|u_l(\zh)|\le\mathsf B(\mathsf v_l)^{1/2}$, hence
$\Xi(\mathsf v_l)<0$ and $\bar\vartheta_m<\mathsf v_l$: $l$ is a non-degenerate
peak, of swallowing type since $w_m(k(l))=\sgn u_l(\zh)M_m=\epsilon_lM_m$.  If
$l_1$ is the first carrier of block $m$, $\bar\vartheta_m<\mathsf v_{l_1}\le
m(1+\|U\|)/\Phi_{l_1}$, while for later $l$ the second half of (SF$^*$) gives
$\mathsf v_l\ge m\delta^\circ_l/\Phi_l\ge2^{10}m(1+\|U\|)/\Phi_{l_1}$; so
$\varrho_l\ge2^{10}$ and $\mu_l=q_0|u_l(\zh)|(1-1/\varrho_l)\ge q_0\delta^\circ_l/2$.

\emph{Step 5 ((c)).}  $\mathsf v_{l_-}=m_0\eta/(n_{l_-}\Phi_{l_-})<\theta_0/2<
\bar\vartheta_{m_0}/2$: a strict non-peak with $\varrho_{l_-}<\frac12$ and $\gap=
M(1-\varrho_{l_-})\ge M/2$; $w(k(l_-))$ has the sign of $u_{l_-}(\zh)<0=-\epsilon_{l_-}$,
so $l_-$ is of anti type and $q_{l_-}<0$; $u_{l_-}1_{F^c}=v_{l_-}$ lives on
$S_{l_-}$, where $z=+1$.

\emph{Step 6 ((d)).}  Let $j\notin F$ and $o:=o_N(j)$, with $m:=m(o)$,
$k:=k(o)$.  Later carriers meet $j$ only through their targets, with
$|u_{l''}(j)|\le\frac43$.  If $j\in S_o$, then $z_j=\epsilon_o$ and the own
term is $\lambda_o\delta_o2^{-j}/n_o$.  If $j\le\min S_o+m+k+4$, (SF$_\iota$)
bounds the later terms by $\frac43\sum_{l''>o}\lambda_{l''}\le\frac432^{-10}
\lambda_o\delta_o2^{-(m+k+\min S_o)}/n_o\le2^{-5}\lambda_o\delta_o2^{-j}/n_o$;
otherwise (W4) gives $c_{l''}\le\iota_o2^{-2j}c_o\delta_o/2$ for every later
$l''$ with $j\in\supp y_{l''}$, so their total is $\le\frac492^{-2j}c_o\delta_o$
against an own term $\ge\frac452^{-m-k}c_o\delta_o2^{-j}$, ratio $\le\frac592^{m+k-j}
<2^{-5}$.  If $j\notin S_o$, then $j$ was unassigned at step $o$, $z_j=
\epsilon_o\sgn y_o(j)$, and $|\lambda_oy_o(j)/n_o|\ge\lambda_o\underline y_o/n_o\ge
2^9\cdot\frac43\sum_{l''>o}\lambda_{l''}$ by (SF$_\iota$).  In both cases the term
of $o$ has sign $z_j$ and exceeds twice the sum of the moduli of all later
terms, also after multiplication of the terms by factors in $[\frac12,1]$.

\emph{Step 7 ((e)).}  Every coordinate off $F$ is assigned $\pm1$.  By
Steps~4, 5 the only strict non-peak is $k(l_-)$, and there is no degenerate
peak.  For the scrambling sums: finitely many carriers (the first carrier of
each block, and $l_-$) have fixed positive margins or gaps; every other
carrier is a peak with $\mu_l\ge q_0\delta^\circ_l/2$ and contributes to
$\mathrm{Scr}_m(\Upsilon s)$ only if $\delta^\circ_l\le2\Upsilon s/q_0$, in which case
(W5) gives $\Phi_l\le c_l\le C c_{l-1}2^{-l}\Upsilon s/q_0$; if $l(s)$ is the least
contributing carrier, the contributions are $\le\sum_{l\ge l(s)}\Phi_l\le2\Phi_{l(s)}
=o(s)$, since $l(s)\to\infty$.  So (MS) holds, and $f_{\rm SA}$ satisfies (BT).

\emph{Step 8 ((f)).}  All carriers are exactly swallowed, so every room
(rate object (O1)) vanishes and every coarse carrier is of class G.  By (b),
(c), every coarse peak is of swallowing type with $\varrho\ge2^{10}\ge1+u$
(beyond the first carriers), and $l_-$ is a class-G carrier with $q<0$ and
$\varrho_{l_-}\in(0,\frac12)$, which at a clean sub-window of large level is
neither $\le b$ nor $\ge1-b$.  Hence (U1), (U2) fail in every block, (U3)
fails in block $m_0$ and holds vacuously in the other blocks (no carrier
with $q<0$), and (L2) holds in every block.  So $I_{\rm up}(w)=\{m_0\}$,
$I_{\rm lo}(w)=\emptyset$, and in the shift cost $c_{\pi(w)}$ the admissible
$\delta$ is $\delta_{m_0}=1$.  Take $x_{l''}:=\lambda_{l''}$ for $l''\in
B_f(w)$.  Since swallowing-type peaks have $\mathrm{vs}_{l''}=\epsilon_{l''}$,
the vector $\Pi_{m_0}+\sum_{B_f}x\epsilon u1_{F^c}$ is the partial sum of
$W$ over the coarse carriers, which by (d) is $z$-signed wherever nonzero;
there are no free coordinates; so every summand $\phi_{z_j}(\cdot)$ vanishes
and $c_{\pi(w)}(f_{\rm SA})=0$.
\end{proof}

So exact coherent shift resonance occurs at block-tame rows of
$T_{\mathrm{final}}$, and the shift cost of (O6) vanishes there.  By
[I,~Corollary~7.22] these rows are nevertheless in $\Rec$.

\begin{corollary}\label{cor:II-SArec}
The rows $f_{\rm SA}$ of Theorem~\textup{\ref{thm:II-SA}} belong to $\Rec$.
\end{corollary}

\begin{proof}
Theorem~\ref{thm:II-SA}(e) and [I,~Corollary~7.22].
\end{proof}

\begin{proposition}[Explicit coherent-shift mates]\label{prop:II-SAmates}
Let $N=1$, $f=f_{\rm SA}$, $k_-:=k(l_-)$, $A:=|\hat\zeta_1|_1$, $\psi:=R_1^*w_1$, and
\[
 V:=u_{l_-}-\frac{u_{l_-}(\zh)}A\,\psi .
\]
Let $\alpha\in\R$, $\Delta_\alpha>0$, $\omega^+:=\alpha e_{k_-}$, $\omega^-:=(\alpha+
\Delta_\alpha)e_{k_-}$.
\begin{enumerate}
\item[(i)] $\Delta d=d(\omega^-)-d(\omega^+)=\Delta_\alpha\lambda_{l_-}u_{l_-}(\zh)/A<0$.
\item[(ii)] $V1_{F^c}$ is $z$-signed and nonzero at every coordinate off $F$.
\item[(iii)] For every $\chi:F^c\to[0,1]$ and every $\beta^+$ supported in $F$
with $b^+(\xi)=0$, i.e.\ $\beta^+(\xi)=-q_0\Delta_\alpha\lambda_{l_-}\sum_{j
\notin F}\chi_j|V_j|$ \textup{(}such $\beta^+$ exist, since $\xi\ne0$ on
$F$\textup{)}, the pairs $b^+:=\chi\Delta_\alpha\lambda_{l_-}V1_{F^c}+
\beta^+$, $b^-:=b^+-\Delta_\alpha\lambda_{l_-}V$ with $\omega^\pm$ are two-piece
data for $g:=b^++R_1^*(\omega^+-d(\omega^+)w_1)$.
\item[(iv)] $cg\in\Cset(f)$ for all small $c>0$, and $(f,cg)\in\cl\NA((c_0,p_1),
\ell_2^2)$.
\item[(v)] For $s\in S_l$, $l\ne l_-$: $g(s)=\psi(s)\big(-d(\omega^+)+\chi_s
\Delta_\alpha\lambda_{l_-}|u_{l_-}(\zh)|/A\big)$; in particular the profile
$g/\psi$ is not constant on $S_l$ whenever $\chi$ is not.
\end{enumerate}
\end{proposition}

\begin{proof}
(i) Off the peaks, $\Phi(k)^2w(k)/C=\hat\zeta(k)/A$ ([I,~Lemma~1.7]), so $d(e_{k_-})
=\hat\zeta(k_-)/A=\lambda_{l_-}u_{l_-}(\zh)/A$.  (ii) All carriers other than $l_-$
are peaks with $w(k)=\epsilon_kM$, so $\psi=MW'+\lambda_{l_-}w(k_-)u_{l_-}$ with
$W':=\sum_{l\ne l_-}\lambda_l\epsilon_lu_l$, and $w(k_-)=-M\varrho_{l_-}$.  Hence
$V=c_1u_{l_-}+c_2W'$ with $c_2:=M|u_{l_-}(\zh)|/A>0$ and $c_1:=1-\lambda_{l_-}
|u_{l_-}(\zh)|M\varrho_{l_-}/A\in(0,1]$ (as $\lambda|u(\zh)|\le A$ and $M\varrho<1$).
Off $F\cup S_{l_-}$, $V=c_2W$, which is $z$-signed and nonzero by
Theorem~\ref{thm:II-SA}(d); on $S_{l_-}$, $V(s)\ge(c_1-2^{-5}c_2\lambda_{l_-})
v_{l_-}(s)>0$, so $V$ is $z_s$-signed there, since $c_1\ge\frac12$ (as
$\varrho_{l_-}<\frac12$ and $M\le1$) and $2^{-5}c_2\lambda_{l_-}\le\frac1{32}$
(as $\lambda_{l_-}|u_{l_-}(\zh)|\le A$; the factor $2^{-5}$ is from Step~6 of
that proof).  (iii)
$b^+-b^-=\Delta_\alpha\lambda_{l_-}V=R_1^*(\omega^--\omega^+)-\Delta d\,\psi$, since
$R_1^*e_{k_-}=\lambda_{l_-}u_{l_-}$; so both pairs represent $g$.  Off $F$,
$b^+=\chi\Delta_\alpha\lambda_{l_-}V$ is $z$-signed and $b^-=(\chi-1)\Delta_\alpha
\lambda_{l_-}V$ is $(-z)$-signed; every coordinate off $F$ is a contact, and
$\supp\omega^\pm\subseteq Q=\{k_-\}$.  Tangency: off $F$ we have
$\zh_j=z_j$ (the vector $Ue$ is supported in $F$), every coordinate is a
contact and $V1_{F^c}$ is $z$-signed, so $(\chi\Delta_\alpha\lambda_{l_-}V
1_{F^c})(\xi)=q_0\Delta_\alpha\lambda_{l_-}\sum_{j\notin F}\chi_j|V_j|$, and
the normalization of $\beta^+$ gives $b^+(\xi)=0$.  Moreover $V(\zh)=
u_{l_-}(\zh)-(u_{l_-}(\zh)/A)\psi(\zh)=0$ because $\psi(\zh)=w_1(\hat\zeta_1)=A$,
hence $b^-(\xi)=b^+(\xi)=0$; and $R_1^*(\omega^+-d(\omega^+)w_1)(\xi)=0$, so
$g(\xi)=0$.  (iv) [I,~Proposition~7.9] on both sides
gives $p^*(f+rg)\le1+\frac{r^2}2(\kappa_w+1)$ for $|r|\le r_0(g)$.  For $c\le
r_0$ with $c^2(\kappa_w+1)\le\frac34$ and $cp^*(g)\le\frac13$: if $|r|\le1$ then
$p^*(f+rcg)\le1+\frac38r^2\le s(r)$, and if $|r|\ge1$ then $p^*(f+rcg)\le1+
|r|cp^*(g)\le s(r)$, by $s(r)-1\ge\min(r^2,|r|)/3$ ([I,~Lemma~3.6(a)]).  So
$cg\in\Cset(f)$, and Corollary~\ref{cor:II-SArec} applies.  (v) At $s\in S_l$,
$l\ne l_-$, $u_{l_-}(s)=0$ and $R_1^*\omega^+(s)=0$, so $g(s)=\chi_s\Delta_\alpha
\lambda_{l_-}V(s)-d(\omega^+)\psi(s)$ with $V(s)=(|u_{l_-}(\zh)|/A)\psi(s)$.
\end{proof}

These mates switch through the anti-type strict non-peak $k_-$ and carry
the uniform shift through all peaks at every scale: exactly the coherent
shift that windows cannot follow.  The next proposition realizes the
hypotheses of Theorem~\ref{thm:II-NL} (compare Remark~\ref{rem:II-NL}).

\begin{proposition}[Realization of Theorem~\ref{thm:II-NL}]\label{prop:II-NLreal}
Let $N=1$ and $f=f_{\rm SA}$.  There are first rows $f_n\to f$ satisfying
\textup{(BT)} such that, for every carrier $l\ne l_-$, all large $n$ and all
$s_1,s_2\in S_l$, the hypotheses of Theorem~\textup{\ref{thm:II-NL}} hold.
Consequently, for every mate $cg$ of Proposition~\textup{\ref{prop:II-SAmates}}
with $\chi_{s_1}\ne\chi_{s_2}$ and every $\rho\in(0,1]$, $\liminf_n\dist_{p^*}
(\rho cg,\Cset(f_n))>0$, although $f$ and every $f_n$ belong to $\Rec$.
\end{proposition}

\begin{proof}
\emph{Step 1 (the rows).}  By Step~1 of Theorem~\ref{thm:II-SA},
$y^*(\zh_{F_\star})=-(\eta+\delta_{l_-}H_{l_-})<0$.  Choose carriers $o_n\to\infty$,
$o_n>l_-$, with target $y^*$ and $\delta_{o_n}H_{o_n}<(\eta+\delta_{l_-}H_{l_-})/2$.
At $f$, $B_{o_n}=0$ and $A_{o_n}<0$, so $\epsilon_{o_n}=-1$ and $z=-1$ on
$S_{o_n}$.  Let $f_n$ have the same $a$ and $z^{(n)}:=z$, except $z^{(n)}:=+1$ on
$S_{o_n}$ and the coordinates owned by carriers $l>o_n$ are re-assigned by
the \emph{hysteretic} recursion: with the signs $\epsilon_l$ of $f$ as
reference signs, keep $\epsilon_l$ unless $\epsilon_lA_l\le-(B_l+\delta_lH_l)/2$, in
which case $\epsilon_l:=\sgn A_l$.  Coordinates owned by carriers $<o_n$ are
unchanged, so $z^{(n)}\to z$ coordinatewise and $f_n\to f$
(Subsection~\ref{subsec:II-comp}).

\emph{Step 2 ($f_n$ satisfies (BT)).}  Values of carriers $<o_n$ are
unchanged.  $n_{o_n}u_{o_n}(\zh)=A_{o_n}+\delta_{o_n}H_{o_n}\le-(\eta+\delta_{l_-}
H_{l_-})/2$ while $z^{(n)}=+1$ on $S_{o_n}$: $o_n$ is an exactly swallowed
anti-type peak with margin bounded below independently of $n$ (its relative
position tends to infinity, the threshold being bounded).  Every re-run
carrier has $n_l|u_l(\zh)|\ge(B_l+\delta_lH_l)/2$ with sign $\epsilon_l$, hence (Step~4
of Theorem~\ref{thm:II-SA} with $|u_l(\zh)|\ge\delta^\circ_l/2$) is a swallowing-type
peak with margin $\ge q_0\delta^\circ_l/4$.  As $\hat\zeta^{(n)}\to\hat\zeta$ in
$\ell_1$, the threshold converges, and the finitely many carriers whose
status is not robust by Step~4 (the first carrier and $l_-$) keep their
statuses for large $n$.  So $Q^{(n)}=\{k_-\}$, there is no degenerate peak,
and (MS) holds as in Step~7.

\emph{Step 3 (alignment).}  $\Sigma_n=F\cup\supp u_{l_-}=F\cup S_{l_-}$, and every
coordinate off $F$ is a contact of $f_n$.  For $s\in S_c$, $c\ne l_-$ the owner,
$\psi_n(s)=\lambda_cw^{(n)}(c)v_c(s)+\sum_{l''>c}\lambda_{l''}w^{(n)}(l'')y_{l''}(s)/
n_{l''}$, and by Step~6 of Theorem~\ref{thm:II-SA} (which bounds moduli,
$|w|\le1$) the sum is $\le2^{-5}\lambda_cv_c(s)<\lambda_cM^{(n)}v_c(s)$ since
$M^{(n)}>\frac13$.  So $\sgn\psi_n(s)=\sgn w^{(n)}(c)=\sgn u_c(\zh^{(n)})$: the
contacts of $S_c$ are aligned for every swallowing-type peak $c$, and those
of $S_{o_n}$ are anti-aligned ($z^{(n)}=+1$, $w^{(n)}(o_n)=-M^{(n)}$).  All of
them lie outside $\Sigma_n$.  For $s_1,s_2\in S_l$, $l\ne l_-$, and $o_n>l$, the
hypotheses of Theorem~\ref{thm:II-NL} hold.  By
Proposition~\ref{prop:II-SAmates}(v), $\pi_{cg}(s_1)\ne\pi_{cg}(s_2)$ when
$\chi_{s_1}\ne\chi_{s_2}$, and Theorem~\ref{thm:II-NL} applies.
\end{proof}

Hence membership in $\Rec$ cannot be propagated along an arbitrary dense
sequence of block-tame rows: approximants for mates with oscillating
profiles must either bank the oscillation coordinates, as the engineered
approximants of [I,~Section~7] do, or keep every peak of a shifted block
exactly swallowed and aligned outside the data supports, as the hysteretic
re-alignments below do.

\smallskip\noindent\emph{Genericity in fixed-$z$ fibres.}  For the next
results let $T$ be any admissible operator and $U$ diagonal with entries
$\mu_j>0$ ($U^*e_j^*=\mu_jk_j$); for $T=T_{\mathrm{final}}$, $\mu_j=\mu^\star_j$.  Fix a
finite set $F$ with $d:=|F|\ge2$, signs $\sigma\in\{-1,1\}^F$ and $z\in
B_{\ell_\infty}$ with $z=\sigma$ on $F$, and put
\[
 \mathcal A_\sigma:=\{a\in\ell_1:\supp a=F,\ \sgn a=\sigma,\ q^*(a)=1\},\qquad
 f_a:=f[a,z]\quad(a\in\mathcal A_\sigma).
\]
For $a\in\mathcal A_\sigma$ let $\nu(a):=\|U^*a\|$, $b_j(a):=\mu_ja_j/\nu(a)$ and
$\Phi(a):=(\sigma_j+\mu_jb_j(a))_{j\in F}$.  Then $\zh(a)_j=\Phi(a)_j$ on $F$ and
$\zh(a)_j=z_j$ off $F$, so $y(\zh(a))=\ip y{\Phi(a)}$ for $y$ supported in $F$;
and $a\mapsto b(a)$ is a homeomorphism of $\mathcal A_\sigma$ onto the open
orthant $\Omega_\sigma$ of the unit sphere $S^{d-1}$ with signs $\sigma$.

\begin{proposition}[Near-threshold carriers are generic]\label{prop:II-Baire}
In this setting assume that $z$ has a contact off $F$, let $m$ be a block
and $(\varepsilon_l)$ any sequence of positive numbers.  Then
\[
 \{a\in\mathcal A_\sigma:\ |\mathsf v_l(f_a)-\bar\vartheta_m(f_a)|<\varepsilon_l\
 \text{for infinitely many carriers $l$ of block }m\}
\]
\textup{(}$\mathsf v_l=m|u_l(\zh)|/\Phi_l$\textup{)} is a dense $G_\delta$ subset
of $\mathcal A_\sigma$.  If $\varepsilon_l\to0$, the row $f_a$ fails
\textup{(BT)} for every $a$ in this set.
\end{proposition}

\begin{proof}
Work on $\Omega_\sigma$.  For $y\in\ell_1$, $\psi_y(b):=y(\zh)=\sum_{j\notin F}y_jz_j+
\sum_{j\in F}y_j(\sigma_j+\mu_jb_j)$ is affine in $b$; $\hat\zeta_m(b)$ is
$\ell_1$-continuous, so $\bar\vartheta_m(b)$ is continuous and positive.  Fix $L$,
$b_0\in\Omega_\sigma$ and a connected open $O\ni b_0$ with compact closure in
$\Omega_\sigma$.  Choose $n\in\R^F$ with $(n_j\mu_j)_j$ not parallel to $b_0$
($d\ge2$), $c:=-\sum_{j\in F}n_j(\sigma_j+\mu_jb_{0,j})$, a contact $j_1\notin F$,
and $y:=(n+cz_{j_1}e^*_{j_1})/q^*(\cdot)$.  Then $\psi_y(b_0)=0$ and $\psi_y$ has
nonzero differential along the sphere at $b_0$, so $\psi_y(b_\pm)=\pm c_0$ for
some $b_\pm\in O$, $c_0>0$.  By (T-d) there are carriers $l_i\to\infty$ of block
$m$ with $q^*(u_{l_i}-y)\to0$, and $\sup_b|u_{l_i}(\zh(b))-\psi_y(b)|\le q^*(u_{l_i}-
y)$.  For large $i$, $u_{l_i}(\zh(b_\pm))$ has sign $\pm$ and modulus $\ge c_0/2$;
the function $G_i:=\mathsf v_{l_i}-\bar\vartheta_m$ is continuous on $O$, positive at
$b_+$ (as $\Phi_{l_i}\to0$), and equal to $-\bar\vartheta_m<0$ at a zero of
$u_{l_i}(\zh)$ on a path from $b_+$ to $b_-$ in $O$.  Hence $\{b\in O:|G_i(b)|<
\varepsilon_{l_i}\}$ is nonempty and open, and $V_L:=\bigcup_{l\ge L}\{|\mathsf v_l-
\bar\vartheta_m|<\varepsilon_l\}$ is open and dense.  The set in question is
$\bigcap_LV_L$, a dense $G_\delta$ by Baire's theorem.  Infinitely many
carriers with $|\mathsf v_l-\bar\vartheta_m|<\varepsilon_l$ are infinitely many strict
non-peaks ($Q_m$ infinite) or infinitely many peaks with $\mu_l=q_0\Phi_l(\mathsf
v_l-\bar\vartheta_m)/m<s_l:=q_0\Phi_l\varepsilon_l/m$, and $\Phi_l/s_l=m/(q_0
\varepsilon_l)\to\infty$ (when $\varepsilon_l\to0$) contradicts (MS).
\end{proof}

So no admissible design can ensure that only finitely many carriers are
near the threshold, and (BT) is meagre in every fixed-$z$ fibre with a
contact.  Nevertheless (BT) rows are dense among all rows with finite base
support (Proposition~\ref{prop:II-BTdense} below): the owner-type rows
escape Proposition~\ref{prop:II-Baire} because their $z$ depends on $a$.
The next two statements show that (C$^*$) is not stable under tilts of $a$
at fixed $z$.

\begin{proposition}[Tilts create both kinds of robust peaks]\label{prop:II-FZ}
In this setting let $a_0,a_1\in\mathcal A_\sigma$ with $\Phi(a_0)$, $\Phi(a_1)$ not
positively proportional.  Then there is $\eta>0$ such that every block
contains infinitely many carriers $k$ with $u_k(\zh(a_0))\ge\eta/2$ and
$u_k(\zh(a_1))\le-\eta/2$, and infinitely many carriers $k'$ with
$u_{k'}(\zh(a_0))\ge\eta/2$ and $u_{k'}(\zh(a_1))\ge\eta/2$.  Moreover, at most
one $a_1\ne a_0$ has $\Phi(a_1)$ positively proportional to $\Phi(a_0)$.
\end{proposition}

\begin{proof}
$\Phi(a)_j=\sigma_j(1+\mu_j|b_j(a)|)$, so all $\Phi(a)$ lie in the open orthant of
the signs $\sigma$ and $\ip{\Phi(a_0)}{\Phi(a_1)}>0$.  Two nonzero vectors that
are not positively proportional and not opposite are linearly independent,
so there is $c\in\R^F$ with $\ip c{\Phi(a_0)}>0>\ip c{\Phi(a_1)}$.  Put $y_0:=
c/q^*(c)$, $y'_0:=\Phi(a_0)/q^*(\Phi(a_0))$ (supported in $F$) and let $\eta$ be the
least of $\ip{y_0}{\Phi(a_0)}$, $-\ip{y_0}{\Phi(a_1)}$, $\ip{y'_0}{\Phi(a_0)}$,
$\ip{y'_0}{\Phi(a_1)}$.  By (T-d) every block has carriers arbitrarily far out
in every $q^*$-ball around $y_0$ and $y'_0$, and $|u(\zh(a))-y(\zh(a))|\le
q^*(u-y)$.  For the last claim, $\Phi(a_1)=\lambda\Phi(a_0)$ together with
$\|b(a_1)\|_2=1$ is a quadratic equation in $\lambda$ with the root $\lambda=1$.
\end{proof}

\begin{corollary}[Tilts of aligned rows leave (C$^*$)]\label{cor:II-FZ}
Let $T=T_{\mathrm{final}}$, $F$, $\sigma$, $z$ as above, and let $a_0\in\mathcal
A_\sigma$ be such that in every block all but finitely many carriers $k$ are
exactly swallowed with the sign of their value: $z=\sgn u_k(\zh(a_0))$ on
$S_k$.  Then for every $a_1\in\mathcal A_\sigma$ with $\Phi(a_1)$ not positively
proportional to $\Phi(a_0)$, the row $f_{a_1}$ has in every block infinitely
many exactly swallowed robust peaks of anti type and infinitely many of
swallowing type.  Consequently, at every clean sub-window of every
sufficiently large level, every block of $f_{a_1}$ has an upper source
\textup{(U2)} and a lower source \textup{(L2)}, $\rho^{\rm sh}\ge1$, $f_{a_1}$ is
not a \textup{(C$^*$)} row, and $f_{a_1}\in\Rec$.
\end{corollary}

\begin{proof}
For large $i$ the carriers $k_i$ of Proposition~\ref{prop:II-FZ} satisfy
$z=+1$ on $S_{k_i}$ (hypothesis, $u_{k_i}(\zh(a_0))>0$) and $u_{k_i}(\zh(a_1))\le
-\eta/2$: at $f_{a_1}$ they are exactly swallowed (rooms $0$, class G at
every level) with $w=-M$, i.e.\ of anti type, and, as $\Phi_{k_i}\to0$, their
relative positions tend to infinity.  The carriers $k'_i$ are exactly
swallowed robust peaks of swallowing type in the same way.  A fixed carrier
is coarse at all later levels, so for all large levels every clean
sub-window has, in every block, a class-G anti-type peak (U2) and a class-G
swallowing-type peak with $\varrho\ge1+u$ (L2).  Theorem~\ref{thm:II-C1} and
Theorem~\ref{thm:II-MTIII} give the rest.
\end{proof}

\begin{remark}\label{rem:II-FZ}
(a) The hypothesis of Corollary~\ref{cor:II-FZ} (aligned rows) matters.  For a
general (C$^*$) row the argument shows only that the source-deficient blocks
acquire both sources at $f_{a_1}$ (by Lemma~\ref{lem:II-Dprime} all robust coarse
peaks of such a block are of one type); a block that has both sources at
$f_{a_0}$ through carriers of both types may lose one type at $f_{a_1}$.  So the
correct reading is: \emph{aligned} (C$^*$) rows are isolated in their fixed-$z$
fibre, up to at most one other point.
(b) Along $f_{a_1}\to f_{a_0}$ the oscillating coherent-shift mates of $f_{a_0}$
are lost whenever $f_{a_1}$ is block-tame (Proposition~\ref{prop:II-S3}, since
the anti-type peaks provide anti-aligned contacts outside the data
supports).  For non-block-tame $f_{a_1}$ the analogous statement is only
heuristic.
\end{remark}

\smallskip\noindent\emph{Exact data with negative shift.}  We now show that
exact two-piece data with a negative shift in every block are rigid, that
they are recovered when they have no dead zones, and that they exist only on
a meagre set.  In the rest of this subsection $T=T_{\mathrm{final}}$, $p=p_N$
and $F$ is finite.  For two-piece data $(b^\pm,\omega^\pm)$ at $f$ put (data
convention)
\[
 D:=b^+-b^-=\sum_k\gamma_ku_k,\qquad\gamma_k:=\lambda_k\big[(\omega^-_m-\omega^+_m)(k)
 -\Delta d_mw_m(k)\big]\quad(m=m(k))
\]
(Lemma~\ref{lem:II-cert}), $K_\omega$ the finite set of carriers in the supports
of $\omega^\pm$, $E:=\bigcup_{k\in K_\omega}\supp y_k\setminus F$ (finite), and
$C_1:=\max_m|\Delta d_m|$.  A carrier $k$ is \emph{neutral} if $\gamma_k=0$ and
\emph{negligible} if $0<|\gamma_k|<2\iota_kC_1\lambda_k$, with the tolerance
sequence $\iota_k=2^{-k}$ of Definition~\ref{def:II-data}.

\begin{lemma}[Dominance with design thresholds]\label{lem:II-dom}
Let $j\notin F$, $o:=o_N(j)$, and suppose $|\gamma_k|\le C_1\lambda_k$ for every
$k>o$ with $u_k(j)\ne0$.  If $|\gamma_o|\ge\iota_oC_1\lambda_o$, then $\sgn D(j)=
\sgn(\gamma_ou_o(j))$ and $|D(j)|\ge\frac{31}{32}|\gamma_ou_o(j)|$.
\end{lemma}

\begin{proof}
As in Step~6 of Theorem~\ref{thm:II-SA}, with the factor $\iota_o$ carried along:
later carriers meet $j$ only through targets, $|u_k(j)|\le\frac43$.  If $j\in
S_o$ and $j\le\min S_o+m(o)+k(o)+4$, (SF$_\iota$) bounds the later terms by
$\frac43C_1\sum_{k>o}\lambda_k\le\frac432^{-10}\iota_oC_1\lambda_o\delta_o2^{-(m+k+
\min S_o)}/n_o$, at most $2^{-5}$ times the own term $\ge\iota_oC_1\lambda_o\delta_o2^{-j}/
n_o$.  If $j\in S_o$ is larger, (W4) gives $c_k\le\iota_o2^{-2j}c_o\delta_o/2$ for
each later $k$ at $j$, so the later total is $\le\frac49C_1\iota_o2^{-2j}c_o\delta_o$
against an own term $\ge\frac45\iota_oC_12^{-m-k}c_o\delta_o2^{-j}$.  If $j\in\supp y_o
\setminus S_o$, the later total is $\le\frac43C_1\sum_{k>o}\lambda_k\le2^{-9}\iota_oC_1
\lambda_o\underline y_o/n_o\le2^{-9}|\gamma_ou_o(j)|$.
\end{proof}

\begin{proposition}[Exact data with negative shift force owner alignment]\label{prop:II-align}
Let $(b^\pm,\omega^\pm)$ be two-piece data at $f$ with $C_1>0$.  For every
$j\notin F\cup E$ whose owner $o=o_N(j)$ is neither neutral nor negligible,
$D(j)\ne0$, $j$ is a contact, and $z_j=\sgn(\gamma_ou_o(j))$; if moreover $o\notin
K_\omega$, then $z_j=-\sgn(\Delta d_{m(o)})\sgn(u_o(\zh))\sgn(u_o(j))$.
\end{proposition}

\begin{proof}
If $o\notin K_\omega$, a carrier of $K_\omega$ meeting $j$ would have $j\in S_k$
(as $j\notin E$) and would be the owner; so every carrier at $j$ lies outside
$K_\omega$ and has $|\gamma_k|=\lambda_k|\Delta d_m||w_m(k)|\le C_1\lambda_k$.  If $o\in
K_\omega$, then $j\in S_o$ and the later carriers at $j$ are again outside
$K_\omega$.  Lemma~\ref{lem:II-dom} gives $D(j)\ne0$ with the sign of $\gamma_o
u_o(j)$, and Lemma~\ref{lem:II-cert}(b) ($D$ is $z$-admissible off $F$) gives the
rest; for $o\notin K_\omega$, $\gamma_o=-\lambda_o\Delta d_mw_m(k(o))$ and $\sgn w_m(k)=
\sgn u_k(\zh)$.
\end{proof}

Thus exact data with $\Delta d<0$ in a block force $z_j=\sgn u_o(\zh)\sgn u_o(j)$
(owner alignment) at every coordinate outside $F\cup E$ owned by a carrier of
that block that does not carry $\omega$ and is not neutral or negligible.  The
coordinates owned by $K_\omega$-carriers, neutral or negligible carriers are the
\emph{dead zones}.

\begin{proposition}[Exact all-negative data are non-generic]\label{prop:II-R5}
In the setting of Proposition~\textup{\ref{prop:II-Baire}} with
$T=T_{\mathrm{final}}$, the set of $a\in\mathcal A_\sigma$ for which some
$g\in\Cset(f_a)$ carries two-piece data with $\Delta d_m<0$ for every $m\in I$ is
meagre in $\mathcal A_\sigma$.
\end{proposition}

\begin{proof}
For a carrier $l$ of block $m$ let $O_l$ be the set of $b\in\Omega_\sigma$ with
$2\iota_l\bar\vartheta_m\Phi_l/(mM_m)<|u_l(\zh)|<\bar\vartheta_m\Phi_l/(2m)$ such that some
$j\in S_l\setminus F$ has $z_j\ne\sgn u_l(\zh)$; it is open.  \emph{$V^m_L:=
\bigcup_{l\ge L,m(l)=m}O_l$ is dense.}  Given $O$, take $y$ and carriers $l_i$ as in
the proof of Proposition~\ref{prop:II-Baire} (with $S_{l_i}\cap F=\emptyset$).  Let
$\epsilon_i$ be the value of $z$ on $S_{l_i}$ if $z$ is constant there and $\epsilon_i:=1$
otherwise, and let $\kappa_i(b)$ be the midpoint of the interval in the
definition of $O_{l_i}$; it is continuous, positive and $O(\Phi_{l_i})$, and the
interval is nonempty for large $i$ ($\iota_{l_i}\to0$, $M_m>\frac13$).  Then
$h:=u_{l_i}(\zh)+\epsilon_i\kappa_i$ has opposite signs at $b_+$ and $b_-$, hence a
zero $b\in O$, where $|u_{l_i}(\zh)|=\kappa_i(b)$ and $\sgn u_{l_i}(\zh)=-\epsilon_i$; in
both cases some $j\in S_{l_i}$ has $z_j\ne-\epsilon_i$, so $b\in O_{l_i}$.
\emph{Conclusion.}  $\mathcal G:=\bigcap_m\bigcap_LV^m_L$ is a dense $G_\delta$.  Let
$b(a)\in\mathcal G$ and suppose data with all $\Delta d_m<0$ exist; let
$|\Delta d_{m_*}|=C_1$.  Pick $l$ of block $m_*$ with $b\in O_l$, $l\notin K_\omega$,
$S_l\cap(F\cup E)=\emptyset$ (all but finitely many $l$ qualify).  Then $\mathsf
v_l<\bar\vartheta/2$, so $l$ is a strict non-peak with $|w(l)|=M\mathsf v_l/\bar\vartheta>
2\iota_l$, and $|\gamma_l|=\lambda_lC_1|w(l)|>2\iota_lC_1\lambda_l$: $l$ is neither neutral
nor negligible.  For the coordinate $j\in S_l\setminus F$ with $z_j\ne\sgn u_l(\zh)$,
$o_N(j)=l$ and $j\notin F\cup E$, and Proposition~\ref{prop:II-align} gives $z_j=
\sgn u_l(\zh)$ (as $\Delta d_{m_*}<0$ and $u_l(j)>0$), a contradiction.
\end{proof}

\begin{proposition}[Block-tame rows are dense]\label{prop:II-BTdense}
Let $f\in S_{p_N^*}$ have finite base support.  For every $\varepsilon>0$ there is
$f_*\in S_{p_N^*}$ with $a_*=a$, satisfying \textup{(BT)}, with $p^*(f_*-f)<
\varepsilon$.
\end{proposition}

\begin{proof}
\emph{Fine re-alignment.}  For $L\ge1$ let $z^{(L)}:=z$ on $F$ and on the
supports of the carriers $l\le L$, and let $z^{(L)}$ be given on the remaining
coordinates by the owner recursion over $L_N\cap(L,\infty)$ (owner signs).  Put
$f^{(L)}:=f[a,z^{(L)}]$.  With $\delta:=\zh^{(L)}-\zh$, $u_k(\delta)=0$ for $k\le L$ and
$|u_k(\delta)|\le2$ otherwise, so $c(\delta)\le C\varepsilon_L\log(e/\varepsilon_L)$,
$\varepsilon_L:=\sum_{l>L}\lambda_l$, and $p^*(f^{(L)}-f)\le C_f\varepsilon_L\log(e/
\varepsilon_L)$ by Lemma~\ref{lem:II-cost}.  Carriers $\le L$ keep their values;
every carrier $l>L$ with $S_l\cap F=\emptyset$ is an exactly swallowed
non-degenerate swallowing-type peak, with margin $\ge q_0\delta^\circ_l/2$ beyond
the first peak of its block (Steps~2 and~4 of Theorem~\ref{thm:II-SA} use only
the owner rule at $l$, $|u_k(\zh)|\le1+\|U\|$ and (SF$^*$)).

\emph{Levers.}  Let $K_0$ be the finite set of carriers $l\le L$ or with $S_l\cap
F\ne\emptyset$; only they can be degenerate peaks or strict non-peaks.  Choose
carriers $l_a(1)<\dots<l_a(N)$, $l_a(m)$ of block $m$, beyond $K_0$, put $s_m:=
\min S_{l_a(m)}$, and replace $z_{s_m}=\epsilon_{l_a(m)}$ by $t_m\epsilon_{l_a(m)}$,
$t_m\in[\frac34,1]$; afterwards re-run the recursion beyond $l_a(m)$
\emph{hysteretically} (reference signs kept unless $\epsilon_lA_l\le-(B_l+\delta_lH_l)/
2$, then $\epsilon_l:=\sgn A_l$), so that every re-run carrier has $n_l|u_l(\zh)|\ge
(B_l+\delta_lH_l)/2$ and is a peak with margin $\ge q_0\delta^\circ_l/4$.  The lever
changes $|u_{l_a(1)}(\zh)|$ at the exact rate $\delta_{l_a(1)}2^{-s_1}/n_{l_a(1)}$ in
$1-t_1$ and keeps it $\ge\frac34\delta^\circ_{l_a(1)}$; the values of $K_0$ are
unchanged.  Write $\hat\zeta_1(t)=\zeta^{\rm fix}+(\text{own lever term})+\zeta^{\rm
later}(t)$ with $\|\zeta^{\rm later}(t)-\zeta^{\rm later}(1)\|_1\le e_1:=2(1+\|U\|)\sum_{l>
l_a(1)}\lambda_l$, and $\varphi(t):=\bar\vartheta(\zeta^{\rm fix}+\text{own term}+\zeta^{\rm
later}(1))$.  By Lemma~\ref{lem:II-T2}(b) (applied between two values of $t$,
the lever carrier being a peak) $|\varphi(t)-\varphi(t')|\ge c_\Theta\lambda_{l_a(1)}
\delta_{l_a(1)}2^{-s_1}|t-t'|/n_{l_a(1)}$, and by Lemma~\ref{lem:II-T2}(a)
$|\bar\vartheta_1(t)-\varphi(t)|\le L_\Theta e_1$, with $c_\Theta,L_\Theta>0$ uniform over the
family.  Hence for each $k\in K_0$ the set $\{t:\bar\vartheta_1(t)=\mathsf v_k\}$ lies in
an interval of length $\le2L_\Theta e_1n_{l_a(1)}/(c_\Theta\lambda_{l_a(1)}\delta_{l_a(1)}
2^{-s_1})\le(L_\Theta/c_\Theta)2^{-8-k(l_a(1))}$ by (SF$_\iota$).  For $l_a(1)$ large
these $|K_0|$ intervals cover less than a quarter of $[\frac34,1]$; choose $t_1$
outside them, so $\mathrm{dgap}_1:=\min_{k\in K_0}|\bar\vartheta_1-\mathsf v_k|>0$.  Choose
$l_a(2)$ with $2L_\Theta(1+\|U\|)\sum_{l\ge l_a(2)}\lambda_l<\mathrm{dgap}_1/2$, then $t_2$
in the same way, and so on.  The resulting row $f_*$ has no degenerate peak,
$Q_m\subseteq K_0$, and (MS) holds as in Step~7 of Theorem~\ref{thm:II-SA}: it
satisfies (BT).  Its cost is $\le C_f\varepsilon\log(e/\varepsilon)$ with $\varepsilon:=
\sum_{l>\min(L,l_a(1))}\lambda_l$, which is small for $L$ large.
\end{proof}

\begin{lemma}[Transplant of exact data]\label{lem:II-transplant}
Let $f,f^\#\in S_{p^*}$, where $F$ is finite and $F^\#=F\cup B$ with a finite
set $B$ of contacts of $f$ at which $\sgn a^\#_j=z_j$ \textup{(}banks\textup{)},
and $z^\#=z$ on $B$.  Let $(b^\pm,\omega^\pm)$ be two-piece data at $f$ for some
$g$ with $g(\xi)=0$, $D:=b^+-b^-$, and assume:
\begin{enumerate}
\item[(T1)] every $k\in\supp\omega^\pm_m$ is a strict non-peak of $f^\#$;
\item[(T2)] $D^\#:=\sum_mR_m^*(\omega^-_m-\omega^+_m)-\sum_m\Delta d^\#_m\psi^\#_m$,
$\Delta d^\#_m:=d^\#_m(\omega^-_m)-d^\#_m(\omega^+_m)$, vanishes at the free
coordinates of $f^\#$ and satisfies $z^\#_jD^\#(j)\ge0$ at the contacts of
$f^\#$ and at $j\in B$.
\end{enumerate}
Then there are two-piece data $(b^\pm_\#,\omega^\pm)$ at $f^\#$, with
$z_jb^+_\#(j)\ge0\ge z_jb^-_\#(j)$ for $j\in B$, for the functional $G^\#:=b^+_\#+
\sum_mR_m^*(\omega^+_m-d^\#_m(\omega^+_m)w^\#_m)$, such that
\[
 \|b^\pm_\#-b^\pm\|_1\le2\big(\|D^\#-D\|_1+\|D1_{\rm Fl}\|_1+\|D^\#1_{\rm Fl}\|_1\big)
 +C_f|c|,\qquad{\rm Fl}:=\{j:z^\#_jz_j=-1\},
\]
where the rebalancing scalar $c$ satisfies $|c|\le C_f(\|D^\#-D\|_1+\|D1_{\rm
Fl}\|_1+\|D^\#1_{\rm Fl}\|_1+p^*(f^\#-f))$.  Consequently, along rows $f^\#\to f$
for which the bracket tends to $0$, $\Delta d^\#_m\to\Delta d_m$, $\|G^\#-g\|_1\to0$
and the $\Gamma_w$ of the transplanted pairs converge to those at $f$.
\end{lemma}

\begin{proof}
Off $F$, at each $j$ with $D^\#(j)\ne0$ (a contact of $f^\#$ or a bank) write
$D^\#(j)=z^\#_js_j$, $s_j>0$, and put $b^+_\#(j):=z^\#_jx_j$, $b^-_\#(j):=-z^\#_jy_j$
with $x_j,y_j\ge0$, $x_j+y_j=s_j$, choosing $(x_j,y_j)$ nearest to
$(z^\#_jb^+(j),-z^\#_jb^-(j))$; both entries are $\ge0$ when $z^\#_j=z_j$, by side
admissibility at $f$, and then the distance is $\le|D^\#(j)-D(j)|$; at flips it
is $\le|D(j)|+|D^\#(j)|$.  Off $F\cup\supp D^\#$ put $b^\pm_\#:=0$.  On $F$ keep
$b^+$ and put $b^-_\#:=b^+_\#-D^\#$.  Finally add $c\,a1_F$ to both $b^\pm_\#$, with
$c$ chosen such that $b^+_\#(\xi^\#)=0$; this is possible because $a1_F(\xi^\#)=
q^\#_0\,a(\zh^\#)\to q_0a(\zh)=q_0\ne0$, and it keeps $b^+_\#-b^-_\#=D^\#$ and the
contact-like split on $B$.  Then $(b^-_\#,\omega^-)$ represents the same functional
as $(b^+_\#,\omega^+)$, by the definition of $D^\#$; side admissibility holds off
$F^\#$ by construction and by (T2); the block parts are supported in $Q^\#$ by
(T1).  The convergence statements follow from Lemma~\ref{lem:II-cost}
($R_m^*w^\#_m\to R_m^*w_m$ in $\ell_1$, and $d^\#_m\to d_m$ on the fixed finite
supports), and $b^+_\#(\xi^\#)-b^+(\xi)\to0$ with $b^+(\xi)=g(\xi)=0$.
\end{proof}

\noindent\emph{Hypotheses for the recovery of exact negative data.}  Let
$g\in\Cset(f)$ carry two-piece data with $\kappa_w\le1$ and $\Delta d_m<0$ for
\emph{every} block $m$, and assume
\begin{enumerate}
\item[(H1)] no carrier is neutral, and only finitely many carriers (the set
$\mathrm{Neg}$) are negligible;
\item[(H2)] $D(j)\ne0$ for every $j\in E':=E\cup\{j\notin F:\ o_N(j)\in K_\omega\cup
\mathrm{Neg},\ |\gamma_{o_N(j)}u_{o_N(j)}(j)|\le32\sum_{k>o_N(j)}|\gamma_ku_k(j)|\}$.
\end{enumerate}
$E'$ is finite: for $o\in K_\omega\cup\mathrm{Neg}$ the later terms on $S_o\setminus E$
are $\le C2^{-2j}$ by (W4), while $|\gamma_o|v_o(j)\ge c2^{-j}$, and $o$ has finitely
many target coordinates.

\begin{lemma}[Scope]\label{lem:II-scope}
Under \textup{(H1)}, \textup{(H2)}, $|z_j|=1$ for every $j\notin F$.
\end{lemma}

\begin{proof}
Let $j\notin F$.  If $j\in E'$, $D(j)\ne0$ by (H2).  If $o_N(j)\notin K_\omega\cup
\mathrm{Neg}$ and $j\notin E$, $D(j)\ne0$ by (H1) and
Proposition~\ref{prop:II-align}.  If $o_N(j)\in K_\omega\cup\mathrm{Neg}$ and $j\notin
E'$, the owner term exceeds $32$ times the later terms, so $D(j)\ne0$.  By
Lemma~\ref{lem:II-cert}(b), $D$ vanishes at free coordinates.
\end{proof}

So the recovery below concerns rows of maximal contact (typically sign
mixed), exactly where Remark~\ref{rem:II-C3}(b) does not exclude (C$^*$); by
Proposition~\ref{prop:II-R5}, however, rows carrying such data form a meagre
set in every fixed-$z$ fibre.

\begin{lemma}[Fine re-alignment carries exact data]\label{lem:II-finealign}
Assume \textup{(H1)}, \textup{(H2)}.  There is $L_0$ such that for every $L\ge
L_0$ the fine re-alignment $f^{(L)}$ of the proof of
Proposition~\textup{\ref{prop:II-BTdense}} satisfies \textup{(T1)} and
\textup{(T2)} \textup{(}with $B=\emptyset$\textup{)}, and the transplanted data
of Lemma~\textup{\ref{lem:II-transplant}} at $f^{(L)}$ represent functionals
$G_L\to g$ with $\Delta d^{(L)}\to\Delta d$ and $\Gamma_w$ converging, as
$L\to\infty$.
\end{lemma}

\begin{proof}
Choose $L_0$ beyond $K_\omega$, $\mathrm{Neg}$, $E'$ and all carriers $l$ with $S_l\cap
F\ne\emptyset$, and so large that $\iota_{L_0}\le\min_m|\Delta d_m|M_m/(4C_1)$.  Coarse
values are unchanged and $\|\hat\zeta^{(L)}-\hat\zeta\|_1\le2(1+\|U\|)\varepsilon_L$;
$w_m(k)$ is Lipschitz in $\hat\zeta_m$ uniformly in $k\le L$ (it is $\sgn\zeta(k)M$ on
peaks and $M\varrho_k\sgn\zeta(k)$ off peaks, continuous across $\varrho=1$, with
$\bar\vartheta$, $M$, $C$, $|\zeta|$ Lipschitz by Lemmas~\ref{lem:II-T},
\ref{lem:II-T2}); so $\sup_{k\le L}|w^{(L)}(k)-w(k)|\le C_f\varepsilon_L\le\iota_L/4$
for large $L$ by (SF$_\iota$), and $\Delta d^{(L)}\to\Delta d$.  (T1): the
$\omega$-carriers are coarse strict non-peaks with fixed values, and the
threshold converges.  (T2) at $j\notin F$ (all such $j$ are contacts of $f$ by
Lemma~\ref{lem:II-scope}, and of $f^{(L)}$ since the owner recursion assigns
$\pm1$):
(i) $o_N(j)>L$: $z^{(L)}_j=\epsilon^{(L)}_o\sgn u_o(j)$ and $w^{(L)}(o)=\epsilon^{(L)}_oM^{(L)}$,
so $\sgn\gamma^{(L)}_o=\epsilon^{(L)}_o$ and $|\gamma^{(L)}_o|\ge\iota_oC^{(L)}_1\lambda_o$ by the
choice of $L_0$; fine coordinates meet no $K_\omega$-carrier;
Lemma~\ref{lem:II-dom} gives $z^{(L)}_jD^{(L)}(j)>0$.
(ii) $o_N(j)\le L$, $o\notin K_\omega\cup\mathrm{Neg}$, $j\notin E$: $z_j=\sgn(\gamma_o
u_o(j))$ by Proposition~\ref{prop:II-align} and $|\gamma_o|\ge2\iota_oC_1\lambda_o$; the
perturbation of $\gamma_o$ is $\le|\Delta d^{(L)}|\lambda_o\iota_L/4+|\Delta d^{(L)}-\Delta
d|\lambda_o$, so $\gamma^{(L)}_o$ keeps its sign and $|\gamma^{(L)}_o|\ge\iota_oC^{(L)}_1
\lambda_o$; Lemma~\ref{lem:II-dom} applies (it needs only $|\gamma^{(L)}_k|\le C^{(L)}_1
\lambda_k$ for later carriers, true also for the re-aligned fine ones).
(iii) $o_N(j)\in K_\omega\cup\mathrm{Neg}$, $j\notin E'$: $\gamma^{(L)}_o\to\gamma_o\ne0$; on
$S_o$ beyond a finite set the bound $C2^{-2j}$ of the later terms against
$|\gamma_o|v_o(j)\ge c2^{-j}$ gives dominance uniformly in $L$; on the finitely many
remaining coordinates the fine contributions are $\le\frac83C_1\varepsilon_L\to0$ and
the coarse ones converge, so the factor-$32$ dominance at $f$ persists.
(iv) $j\in E'$: $z_jD(j)>0$ by (H2) and $D^{(L)}(j)\to D(j)$.
Finally $\|D^{(L)}-D\|_1\le\frac53\sum_k\lambda_k|\Delta d^{(L)}w^{(L)}(k)-\Delta dw(k)|\to0$
by dominated convergence, and ${\rm Fl}$ consists of fine coordinates, where
$D$ and $D^{(L)}$ are sums over fine carriers, so $\|D1_{\rm Fl}\|_1+\|D^{(L)}
1_{\rm Fl}\|_1\le\frac{10}3C_1\varepsilon_L\to0$.  Lemma~\ref{lem:II-transplant}
applies.
\end{proof}

\begin{lemma}[Bank levers]\label{lem:II-banklever}
Assume \textup{(H1)}, \textup{(H2)}, let $L\ge L_0$ and $f':=f^{(L)}$ with the
transplanted data.  For each block $m$ choose a carrier $l_b(m)>L$ of block $m$,
$l_b(1)<\dots<l_b(N)$, put $s_m:=\min S_{l_b(m)}$ \textup{(}a contact with
$z=\epsilon_{l_b(m)}$\textup{)}, and for $\mu\in[0,\infty)^N$ let $f^\#(\mu)$ have
$a^\#(\mu):=(a+\sum_m\mu_m\epsilon_{l_b(m)}e^*_{s_m})/q^*(\cdot)$, $z$ unchanged up to a
hysteretic re-run of the owner recursion beyond $l_b(1)$.  For suitable small
$\mu$, chosen block after block, $f^\#(\mu)$ has base support $F\cup\{s_1,\dots,
s_N\}$, satisfies \textup{(BT)}, carries transplanted two-piece data with $\Delta
d^\#_m<0$ for every $m$ and with the contact-like split at the banks, and
$p^*(f^\#(\mu)-f')+\|G^\#-G'\|_1\to0$ as $\mu\to0$.
\end{lemma}

\begin{proof}
By Lemma~\ref{lem:II-base}(b), $\nu(\mu)^2=\nu^2+\sum_m(\mu^\star_{s_m}\mu_m)^2$ before
normalization, so $\zh$ moves by $O(|\mu|^2)$ on $F$, and $\zh_{s_m}=\epsilon_{l_b(m)}
(1+(\mu^\star_{s_m})^2\mu_m/\nu)+O(|\mu|^2)$.  Hence $|u_{l_b(m)}(\zh)|$ increases at the
first-order rate $v_{l_b(m)}(s_m)(\mu^\star_{s_m})^2/\nu$; the other carriers containing
$s_m$ are later carriers, of total weight $\le2^{-10}\lambda_{l_b(m)}v_{l_b(m)}(s_m)
2^{-(m+k(l_b(m)))}$ by (SF$_\iota$); carriers depending on $F$ move by
$O(|\mu|^2)$; and forced flips of the hysteretic re-run occur only beyond a
level tending to infinity as $\mu\to0$, with total weight $o(|\mu|)$.  By
Lemma~\ref{lem:II-T2}(b), $\bar\vartheta_m(\mu)-\bar\vartheta_m\ge\kappa_m\mu_m/2$ for
small $\mu_m$, with $\kappa_m>0$, while each coarse $\mathsf v_k$ ($k\le L$, the
only candidates for degeneracy) moves by $O(|\mu|^2)$.  So a small $\mu_1>0$,
avoiding finitely many values, makes block $1$ free of degenerate peaks with
a gap $\mathrm{dgap}_1>0$; then $l_b(2)$, $\mu_2$ are chosen so that block $1$ moves
by less than $\mathrm{dgap}_1/2$, and so on.  (BT): $F^\#$ is finite, $Q^\#_m$ lies
among the coarse carriers, there is no degenerate peak, and (MS) holds by the
margins of the re-run peaks.  (T1) holds since raising the threshold keeps
strict non-peaks; (T2) as in Lemma~\ref{lem:II-finealign}, with
$|\gamma|\le C_1\lambda$ for the re-run carriers; at $s_m$ the owner $l_b(m)$
dominates, so $z_{s_m}D^\#(s_m)>0$ and the split is contact-like there.  The
cost bound follows from Lemma~\ref{lem:II-cost} and
Lemma~\ref{lem:II-transplant}.
\end{proof}

\begin{theorem}[Exact negative data without dead zones are recovered]\label{thm:II-negexact}
Let $N\ge1$, let $f\in S_{p_N^*}$ have finite base support, and let $g\in\Cset(f)$
carry two-piece data with $\kappa_w\le1$ and $\Delta d_m<0$ for every block $m$,
satisfying \textup{(H1)} and \textup{(H2)}.  Then $(f,g)\in\cl\NA((c_0,p_N),
\ell_2^2)$.
\end{theorem}

\begin{proof}
Let $f_i:=f^\#(\mu_i)$ be the rows of Lemma~\ref{lem:II-banklever} built on
$f^{(L_i)}$, with $L_i\to\infty$ and then $\mu_i\to0$ so fast that, by
Lemmas~\ref{lem:II-finealign} and \ref{lem:II-banklever}, $f_i\to f$ and the
transplanted data at $f_i$ represent $G_i\to g$ with $\Delta d^{(i)}_m<0$ and
$\kappa^{(i)}_w\to\kappa_w$.  Every $f_i$ satisfies (BT), hence (SC) for every
set of blocks.  Fix $\rho<1$.  The data have fixed block supports (strict
non-peaks whose gaps at $f_i$ are bounded below), fixed $\|b\|_1$-bounds and
contact-like splits at the banks; Lemma~\ref{lem:II-U}, applied with $t:=t_1$
and constants $A_0,A_2,\gamma_B$ fixed by these data, gives $r_1>0$ and $i_0$ with
$p^*(f_i+rG_i)\le1+\frac{r^2}2(\kappa^{(i)}_w+\varepsilon)$ for $|r|\le r_1$, $i\ge i_0$.
With $\varepsilon$ so small that $\rho^2(1+\varepsilon)<1$, and after decreasing $r_1$
so that $\rho^2(1+\varepsilon)r^2/2\le r^2/2-r^4/8$ for $|r|\le r_1/\rho$, we get
$p^*(f_i+r\rho G_i)\le s(r)$ for $|r|\le r_1/\rho$; for larger
$|r|$, $p^*(f_i+r\rho G_i)\le s(\rho r)+p^*(f_i-f)+|r|\rho p^*(G_i-g)\le s(r)$ for $i$
large, since $s(r)-s(\rho r)\ge(1-\rho^2)\min(r^2,|r|)/3$ ([I,~Lemma~3.6(a)]).  So
$\rho G_i\in\Cset(f_i)$, with two-piece data of $\kappa_w\le1$ and $\Delta d<0$, and
[I,~Theorem~7.18] at $f_i$ gives $(f_i,\rho G_i)\in\cl\NA$.  Let $i\to\infty$,
then $\rho\to1$.
\end{proof}

At the self-aligned rows of Theorem~\ref{thm:II-SA} with finitely many
exceptional carriers tuned through $F$, the hypotheses (H1), (H2) hold for
data supported on the exceptional strict non-peaks with nonzero coefficients
there (every non-exceptional carrier is a robust peak with $|w|=M$, and every
coordinate off $F$ is dominated by its owner); these rows are, however,
block-tame and recovered directly (Corollary~\ref{cor:II-SArec}).  In the
other direction, by Proposition~\ref{prop:II-R5} the exact-data route through
Theorem~\ref{thm:II-negexact} acts only on a meagre set of every fixed-$z$
fibre: the mates of a generic row with finite base support must be handled by
window data, and for these the residual is (C$^*$).

\subsection{Absorber pairs and Master Theorem~IV$'$}\label{subsec:II-absorb}

This subsection constructs exact \emph{shifted} two-piece data at
companions of a row with finite base support, at the scales of a clean
sub-window, under a hypothesis (KN) on the availability of a second lever in
each shifted block.  We use the data convention, and for a block $m$ of a
first row $f'$ we write $\Delta'_m:=\Delta d_mM'_m$.  By
Lemma~\ref{lem:II-cert}, exact data with switching difference $\Delta\omega$ exist
iff the single vector $\mathcal V(\Delta\omega)=\sum_l\gamma_lu_l$ is
$z'$-admissible off $F'$, where $\gamma_l=\lambda_l(\Delta\omega(k(l))-\Delta d_{m(l)}
w'(k(l)))$; by Lemma~\ref{lem:II-kappa}, if $\Delta\omega_m$ is supported in a finite
set $\Omega_m\subseteq Q'_m$, then
\begin{equation}\label{eq:II-X3}
 \Delta'_m\,\varkappa'_m=\sum_{l\in\Omega_m}u_l(\zh')\gamma_l,\qquad
 \varkappa'_m:=\varkappa(\hat\zeta'_m;\Omega_m),
\end{equation}
so the only $f'$-dependent quantities of block $m$ entering the coarse exact
system are the values $u_l(\zh')$, $l\in\Omega_m$, and the single scalar
$\varkappa'_m$.

\subsubsection*{Why one scalar requires two levers}

\begin{proposition}[Status rigidity in ratio space]\label{prop:II-Rkt}
Let $\zeta=\hat\zeta'_m$ be the block vector of a first row, $P$ its peak set,
$\Omega\subseteq Q$ finite, $\varkappa:=\varkappa(\zeta;\Omega)$, and put
$r_l:=u_l(\zh')/\varkappa$ for $l\in\Omega$.
\begin{enumerate}
\item[(i)] After division by $\varkappa_m$, \eqref{eq:II-X3} reads $\Delta'_m=
\sum_{l\in\Omega_m}r_l\gamma_l$: the coarse exact system depends on block $m$ only
through the ratios $r_l$.  By Corollary~\textup{\ref{cor:II-Rkt}} the relative
positions $\varrho_l$, $l\in\Omega$, are functions of $(|r_l|)_{l\in\Omega}$, $\Phi_P^2$
and $s_Q$; more precisely $(a,\mathsf k)$ of Corollary~\textup{\ref{cor:II-Rkt}}
depend smoothly on $(R^2,\Phi_P^2,s_Q)$ near any given point.
\item[(ii)] An outward push of a peak changes $(\bar\vartheta,|\zeta|,\varkappa)$ at the
rates $(C/\Phi_P^2,M,1-X)$ \textup{(}Lemma~\textup{\ref{lem:II-kappader}(b)(i))}; with
the values on $\Omega$ fixed it multiplies all ratios of the block by one
factor.  If the ratios are restored afterwards, every $\varrho_l$, $l\in\Omega$,
returns to its old value up to the change of $\Phi_P^2$ and $s_Q$.
\item[(iii)] A push of an $\Omega$-carrier $d_0$ changes $\varkappa$ only at the rate
$\varrho_{d_0}MX/C$ and $\bar\vartheta$ at the rate $-\varrho_{d_0}M/\Phi_P^2$
\textup{(}Lemma~\textup{\ref{lem:II-kappader}(b)(ii),(iii))}.  So an inward push of
$d_0$, followed by a restoration of $\varkappa$ and of the ratios $r_l$, $l\ne d_0$,
raises $\bar\vartheta$ at first order and lowers every $\varrho_l$, $l\in\Omega\setminus
\{d_0\}$.
\end{enumerate}
\end{proposition}

\begin{proof}
(i) With $\mathsf k=a+\Phi_P^2+s_Q$, the identity $a^2=\Phi_P^2+\mathsf k^2R^2+s_Q$ is
an equation $G(a;R^2,\Phi_P^2,s_Q)=0$ with $\partial_aG=2a-2\mathsf kR^2=2(a-\mathsf k
R^2)>0$, since $a>\mathsf kR\ge\mathsf kR^2$ ($R<1$); the implicit function theorem
gives smooth dependence, and $\varrho_l=\mathsf r_l\mathsf k/\Phi(l)$ with $\mathsf r_l=
m|r_l|$.  (ii) The ratios are $u_l(\zh')/\varkappa$; restoring them restores $R^2$,
hence $(a,\mathsf k)$ and $\varrho_l$ by (i), up to the change of $(\Phi_P^2,s_Q)$.
(iii) The rates are those of Lemma~\ref{lem:II-kappader}; at fixed $\varkappa$ the
relation $\varrho_l=\mathsf r_l\varkappa/(\bar\vartheta\Phi(l))$ shows that raising
$\bar\vartheta$ at fixed ratios lowers $\varrho_l$.
\end{proof}

\begin{remark}[The gap of a one-lever exactification]\label{rem:II-onelever}
Suppose one tries to realize a {\L}ojasiewicz point of the tiny minors of the
coarse shifted system in the variables $(u_l)_{l\in\Omega}$ and $(\varkappa_m)_m$ by the
assembly of Section~\ref{sec:II-exact}, the tuning
Lemma~\ref{lem:II-TU}, and one peak push per block that sets $\varkappa_m$.  The
donor raise (C3) of the assembly raises the donor peak by an amount in $[\Lambda,
3\Lambda]$, $\Lambda=T_{\rm lo}(w)^3$, so it moves $\varkappa_m$ by at least $\frac34\Lambda$,
which is far larger than $\mathrm{Des}\,b(w)$.  If $\varkappa_m$ is pushed back to a
{\L}ojasiewicz value near its value before the raise, then by
Proposition~\ref{prop:II-Rkt}(ii) the threshold returns as well and the raise is
undone, so near-threshold switching carriers lose the strict-non-peak status
that the raise provided.  If instead the {\L}ojasiewicz point is computed
after the raise, the tiny minors are only $\le b+\mathrm{Lip}(l)\Lambda$ there, and
Lemma~\ref{lem:II-loj} gives a distance $C_L(\mathrm{Lip}\,\Lambda)^{2/N_L}$, which
is not $\le T_{\rm lo}(w)^4$ when $N_L>2$.  Thus exactness and threshold
protection cannot be achieved independently by peak pushes; by
Proposition~\ref{prop:II-Rkt}(iii) the missing ingredient is a lever that
moves the threshold without moving $\varkappa$ or the active ratios.  No claim is
made that the conclusion of a one-lever exactification fails.
\end{remark}

\subsubsection*{The design $D^{U1'}$}

\begin{definition}[Absorber targets and the design $D^{U1'}$]\label{def:II-DU1}
\emph{Absorber targets.}  For a coordinate $s$, a sign $\beta\in\{\pm1\}$ and a
stage $p$, an \emph{absorber target} is $y:=x/q^*(x)$ with
\[
 x:=\beta e_s^*+e_{p_0}^*+4\sum_{r=1}^R2^{-r}e_{p_r}^*,
\]
where $p_0<p_1<\dots<p_R$ are fresh odd integers (in no signature set and in no
earlier target) and $R$ is the least integer with $2^{3-R}\le(\mu^\star_{p_0})^2
(c^{\rm low}_p)^2/4$; here $c^{\rm low}_p$ is the minimum of $c_{p-1}/4$,
$b(p-1,M(p-1))^2$, $(\delta^{\max}_pH_p/2)^2$ and the bounds (W4$''$) at $s$,
(W5) with $\delta_p$ replaced by $\delta^{\max}_p/2$, and (W7) below, which do
not depend on $R$ or on $p_0$.  Moreover $p_0$ is chosen so large that
\[
 (\mu^\star_{p_0})^{2^{2^{p_0}}-1}<2^{-R}
\]
(possible: $2^{-R}\ge(\mu^\star_{p_0})^2(c^{\rm low}_p)^2/64$ by the minimality
of $R$, and $c^{\rm low}_p$ does not depend on $p_0$).  Then $|y(s)|\in[\frac18,1]$;
$p_0$ is the \emph{tuning coordinate} and $p_1,\dots,p_R$ the \emph{binary
coordinates}.  Two consecutive carriers of block $1$ with absorber targets for
$(s,+1)$ and $(s,-1)$ form a \emph{biased pair} for $s$.

\emph{Stages.}  $D^{U1'}$ is obtained from $T_{\mathrm{final}}$ by the following
changes of the schedule and of the weights; everything else is kept, the
ladder bijection being determined by the schedule.  The stages are grouped in
rounds $r=1,2,\dots$; round $r$ serves a block $m_r$ from a fixed sequence in
which every block occurs infinitely often, and a target $y$ for $m_r$, namely
the scheduled target of $T_{\mathrm{final}}$ if it satisfies allowedness (a)
and (c) with margin ($\supp y\cap S_{l'}\subseteq[1,p-2|\supp y|-1]$, $p$ the stage
at which $y$ will be placed), and the fresh-dominated target otherwise.  Round
$r$ consists of
\begin{enumerate}
\item[PRE:] for every $s\in\supp y\cap\bigcup_{l'}S_{l'}$ a biased pair for $s$
(block $1$, consecutive stages);
\item[MAIN:] the stage $p$ at which $y$ is placed (block $m_r$), called a
\emph{main stage};
\item[CLUSTER:] a biased pair for every $s\in E^{\rm ch}(p):=T(p)\cup\bigcup_{l''\le p}
(S_{l''}\cap[1,p])$ (block $1$).
\end{enumerate}
\emph{Weights.}  At PRE and MAIN stages, $c_l$ is the minimum of (W1)--(W5),
(W7) of Definition~\ref{def:II-Tfinal}, with (W6) omitted at absorber stages,
and with (W4) strengthened to
\[
 \text{(W4$''$)}\qquad c_l\le\iota_{l'}2^{-2s}c_{l'}\delta_{l'}T_{\rm lo}(l-1,M(l-1))^3/2
 \quad(l'<l,\ s\in\supp y_l\cap S_{l'}) ;
\]
at the $i$-th stage of the cluster of a main stage $p$, $c_{p+i}:=c^{\rm low}_{p+1}
4^{1-i}$, where in (W4$''$) the factor $T_{\rm lo}(\cdot)^3$ is read with the main
stage $p$.  Sub-window data (step (6) of Definition~\ref{def:II-Tfinal}) are
formed at every stage, but windows are used only at main stages.  At a main
stage $L$ the design factor $\mathrm{Des}(L)$ additionally contains $4^{R_{\rm
cl}(L)}$ ($R_{\rm cl}(L):=2|E^{\rm ch}(L)|$), $2^{k(L+R_{\rm cl}(L))}$, the reciprocal
of the least nonzero modulus of a target entry; the rate scheme contains the
minors of Definition~\ref{def:II-Gamma}.  Put $R_\Gamma(L):=2|T(L)|+11L$
(a bound for the number of rows, equations counted twice, of the systems
defining the cones $\Gamma^\#(a)$ of Definition~\ref{def:II-Gamma} at main
stage $L$, after the aggregation in (X1)), $P(L):=2^{R_\Gamma(L)}$ (a bound
for the number of their extreme rays), and, for a sub-window $w$ of a main
stage $L$,
\[
 \bar D(w):=(2L+1)\,3^{2L}\Big(\frac{\mathrm{Des}(L)}{u(w)}\Big)^{C_\Gamma
 (P(L)+2L+2)},
\]
where $C_\Gamma$ is a numerical constant, fixed once and for all, so large
that every bound stated in the form ``$\le C_f\bar D(w)$'' in the second half
of this subsection holds.  Then $\bar D(w)$ is a design number of the
sub-window (it is formed from $L$, $T(L)$, $\mathrm{Des}(L)$ and $u(w)$),
and at main stages $Q(w)$ is replaced by $\max\{Q(w),\bar D(w)^3\}$ (the
value of Definition~\ref{def:II-Tfinal}(6) and $\bar D(w)^3$; both are
known when $Q(w)$ is formed).  The window arguments of
Sections~\ref{sec:II-Tfinal}--\ref{sec:II-Cstar} use only the lower bound
$Q(w)\ge(4\mathrm{Des}(l)/u(w))^{\omega(l)+20}$ and the definitions of $n(w)$,
$T_{\rm hi}(w)$, $T_{\rm lo}(w)$, $\eta(w)$ and $b(w)$ in terms of $Q(w)$,
so they are not affected.
\end{definition}

\begin{lemma}[Properties of $D^{U1'}$]\label{lem:II-DU1}
\textup{(a)} $D^{U1'}$ is admissible \textup{(}conditions \textup{(T-a)}--\textup{(T-d)}
of [I,~Definition~1.3]\textup{)} and does not depend on $N$.
\textup{(b)} At every main stage $L$, $\sum_{l'>L}\lambda_{l'}\le b(L,M(L))^2\le
T_{\rm lo}(w)^8$ for every sub-window $w$ of $L$.
\textup{(c)} Every result of Sections~\ref{sec:II-companions}--\ref{sec:II-exact}
and of Subsections~\ref{subsec:II-hoffman}--\ref{subsec:II-dichotomy} that is used
below holds for $D^{U1'}$ with ``level'' read as ``main stage''.
\end{lemma}

\begin{proof}
(a) (T-a) is immediate.  (T-b), (T-c) follow as for $T_{\mathrm{final}}$, since
that argument uses only allowedness (a) and (W4): absorber targets satisfy
(a) (their coordinates are fresh or the coordinate $s$, which lies in $Z_0$ or
in a signature set of an earlier carrier), cluster weights satisfy (W4) by the
choice of $c^{\rm low}$, and (W4$''$) implies (W4).  (T-d): every target of the
dense family is placed infinitely often in every block, because every block
is served in infinitely many rounds and every request of a round is finite.
The definition is well founded since $R$ is determined from $c^{\rm low}_p$,
which does not depend on $R$.  (W3) at cluster stages: $c_{p+i}\le b(p,M(p))^2\le
2^{-8n(p,M(p))}$, while $(\delta_{p+i}H_{p+i})^2\ge2^{-2(p+i)-6\cdot2^{p+i}-4}$ and
$p+i\le3p+2s_{\max}(p)$, $n(p,M(p))\ge2^{6s_{\max}(p)}$.  Absorbers live in block
$1$, which is present for every $N$.  (b) The cluster weights sum to $\le2c^{\rm
low}_{L+1}$ and the later weights to $\le\frac43c_{L+R_{\rm cl}+1}$, so $\sum_{l'>L}
c_{l'}\le2c_{L+1}\le2b(L,M(L))^2$, and $\lambda_{l'}\le c_{l'}/4$.  (c) The window
arguments use one sub-window's dyadic scales, the design factors of its level,
the box bound (which needs (b)) and $T_{\rm hi}(w')\le T_{\rm lo}(w)$ for later
$w'$ (Lemma~\ref{lem:II-GW}); all hold at main stages.  The non-window
arguments use (T-a)--(T-d), allowedness, bounded gaps and the diagonal base.
The features (W5), (W6), (W7) are lost at absorber stages; they are not used
by the results invoked in this subsection.
\end{proof}

\begin{lemma}[First touchers]\label{lem:II-firsttouch}
Let $w$ be a sub-window of a main stage $L$, let $s_{\rm far}(w)$ be the least $s$
with $2^{-s}\le c_{L+1}^2$, and
\[
 E_c(w):=T(L)\cup\bigcup_{l''\le L}\big(S_{l''}\cap[1,s_{\rm far}(w)]\big).
\]
For every $s\in E_c(w)$, either no carrier at a stage $>L$ meets $s$, or the first
two carriers at stages $>L$ whose vectors meet $s$ form a biased pair for $s$.
\end{lemma}

\begin{proof}
If $s\in E^{\rm ch}(L)$, the cluster of $L$ follows $L$ immediately.  Otherwise $s\in
S_{l''}$ with $l''\le L<s$.  After $L$, the carriers meeting $s$ are absorber pairs
for $s$ (pre-pairs and cluster pairs) and main targets containing $s$, each of
which is preceded by its pre-pair for $s$; fresh-dominated and absorber targets
meet only fresh coordinates and their own coordinate $s'$, and signature sets
are disjoint.  Pairs are placed at consecutive stages.
\end{proof}

\begin{lemma}[Zero-value tuning of an absorber]\label{lem:II-zerovalue}
Let $a$ be an absorber carrier with target $y_a=x/q^*(x)$ for $(s,\beta,p)$, let
$s_{\rm far}(a)$ be the least $s'$ with $2^{-s'}\le\Phi_a^2$, and let $f^{(1)}$ be a
first row with finite base support containing none of $p_0,\dots,p_R$, none of
$S_a\cap[1,s_{\rm far}(a)]$, and not $s$.  Let $0<\theta_a\le\lambda_a$.  Put
$z:=+1$ at $p_0$ and on $S_a$, make $p_r$ $(r\ge1)$ and the coordinates of
$S_a\cap[1,s_{\rm far}(a)]$ support coordinates with masses $\theta_a\sigma_r$
(binary signs $\sigma_r\in\{\pm1\}$, $z_{p_r}:=\sigma_r$), resp.\ $\theta_a$, and
$p_0$ a support coordinate with mass $m_0\in[\theta_a,\theta_a+c_p^2]$.
\begin{enumerate}
\item[(i)] For every $\theta_a\in(0,\lambda_a]$ there are binary signs $\sigma_r$
\textup{(}depending on $\theta_a$, chosen in the proof\textup{)} and $m_0\in
[\theta_a,\theta_a+c_p^2]$ with $u_a(\zh)=0$ exactly; then $a$ is a strict
non-peak with $w(a)=0$ and gap $M$.
\item[(ii)] Finitely many absorbers can be tuned to value $0$ simultaneously
\textup{(}the signs of each absorber chosen as in \textup{(i)} at the row
where every tuning mass equals its $\theta_a$\textup{)}.  Here and in
\textup{(i)} the absorber is at a stage $p$ so late that $2\mu^\star_{p_0}\le
\nu_{(1)}/4$ and $160\lambda_a^2\le\nu_{(1)}^2/2$, where $\nu_{(1)}$ is the
$\nu$ of $f^{(1)}$; this holds for all absorbers at stages beyond an
$f$-dependent stage, which is how the lemma is used.
\item[(iii)] A tuned absorber has no $d$-effect and no effect on $\varkappa$:
$u_a(\zh)\gamma_a=0$ in \eqref{eq:II-X3}, and $\mathsf v_a=0$.
\item[(iv)] The cost is $\le C(\lambda_a+\text{sum of the masses})\log$, and the
values of other carriers move only by the second-order effect $\le C(\sum
\text{masses}\cdot\mu^\star)^2$ of the diagonal base, except for carriers at later
stages whose targets meet the new support coordinates.
\end{enumerate}
\end{lemma}

\begin{proof}
(i) We have $q^*(x)n_au_a(\zh)=\beta\zh_s+\zh_{p_0}+4\sum_r2^{-r}\zh_{p_r}+q^*(x)
\delta_ah_a(\zh)=:\Psi(m_0)$, and at a support coordinate $p$ with mass $m_p$
and sign $z_p$, $\zh_p=z_p(1+(\mu^\star_p)^2m_p/\nu)$, while $\zh_j=z_j$ off the
support (Lemma~\ref{lem:II-base}(b)).  Since $U$ is diagonal, $\nu^2=\sum_j
(\mu^\star_j)^2a_j^2$ depends only on the moduli of the masses, not on the
signs $\sigma_r$; and $s$, $p_0$ and $S_a$ carry no binary mass.  Hence the
quantity
\[
 \mathrm{Tg}:=-\big(\beta\zh_s+\zh_{p_0}+q^*(x)\delta_ah_a(\zh)\big),
\]
evaluated at the row with $m_0=\theta_a$, does not depend on the signs
$\sigma_r$, and at that row $\zh_{p_r}=\sigma_r(1+e_r)$ with $e_r:=(\mu^\star_{p_r})^2
\theta_a/\nu\ge0$.  As $\nu\ge\mu^\star_{p_0}\theta_a$ and $\mu^\star_{p_r}\le
\mu^\star_{p_0+1}=(\mu^\star_{p_0})^{2^{2^{p_0}}}$ for $r\ge1$, we have $e_r\le e:=
(\mu^\star_{p_0})^{2^{2^{p_0}}-1}<2^{-R}$ (Definition~\ref{def:II-DU1}).  Moreover
$|\mathrm{Tg}|\le1+(1+\mu^\star_{p_0})+\frac32(1+\mu^\star_{p_0})\le3.6$, because
$|\zh_s|=|z_s|\le1$ and $q^*(x)\delta_a\|h_a\|_1\le\frac32$ by the choice of
$\delta^{\max}$.  The exact binary sums $4\sum_r2^{-r}\sigma_r$ run through the
odd multiples of $2^{2-R}$ in $(-4,4)$, so there are signs with $|(\mathrm{Tg}-
2^{3-R})-4\sum_r2^{-r}\sigma_r|\le2^{2-R}$; for these signs
$|\sum_r4\cdot2^{-r}e_r\sigma_r|\le4e$, hence
\[
 \Psi(\theta_a)=-\Big(\mathrm{Tg}-\sum_r4\cdot2^{-r}(1+e_r)\sigma_r\Big)\in
 [-3\cdot2^{2-R}-4e,\,-2^{2-R}+4e]\subseteq[-2^{4-R},0) .
\]
Raising $m_0$ from $\theta_a$ by $\Theta\in[0,c_p^2]$ raises $\zh_{p_0}$ by at
least $(\mu^\star_{p_0})^2\Theta/(2\nu)$: $\nu$ is at least the value $\nu_{(1)}$ of
$f^{(1)}$ minus a negligible amount, and $\mu^\star_{p_0}(\theta_a+c_p^2)\le
2\mu^\star_{p_0}\le\nu_{(1)}/4$ for absorbers of late stages, so $d(m_0/\nu)/dm_0
=(\nu^2-(\mu^\star_{p_0}m_0)^2)/\nu^3\ge1/(2\nu)$.  The other terms of $\Psi$ depend
on $m_0$ only through $\nu$: they are sums $\sum_jc_j(1+(\mu^\star_j)^2m_j/\nu)$
over the binary coordinates and $S_a\cap[1,s_{\rm far}(a)]$ with $\sum_j|c_j|
(\mu^\star_j)^2m_j\le5\theta_a$, and $|\Delta(1/\nu)|\le(\mu^\star_{p_0})^2m_0\Theta/
\nu^3$ (as $\Delta(\nu^2)\le2(\mu^\star_{p_0})^2m_0\Theta$); so their total change is
at most $10\theta_am_0/\nu^2$ times the increase $(\mu^\star_{p_0})^2\Theta/(2\nu)$,
and $10\theta_am_0/\nu^2\le160\lambda_a^2/\nu_{(1)}^2\le\frac12$ for absorbers of
late stages.  So $\Psi$ is continuous and
increases by at least $(\mu^\star_{p_0})^2c_p^2/(4\nu)\ge4(\mu^\star_{p_0})^2c_p^2$
($\nu\le\|U\|=\frac1{16}$) over $[\theta_a,\theta_a+c_p^2]$.  Finally $c_p\ge
c^{\rm low}_p$: at a cluster stage $p+i$, $i\ge2$, by the term $c_{p+i-1}/4=
c_{p+i}$ of $c^{\rm low}_{p+i}$, at the first cluster stage by definition, and
at PRE stages because every term of $c^{\rm low}_p$ is at most the corresponding
term of the minimum defining $c_p$ ($\delta_p\ge\delta^{\max}_p/2$).  Hence the
increase is at least $4(\mu^\star_{p_0})^2(c^{\rm low}_p)^2\ge2^{7-R}>2^{4-R}$ by
the choice of $R$, and the intermediate value theorem gives $m_0\in[\theta_a,
\theta_a+c_p^2]$ with $\Psi(m_0)=0$, i.e.\ $u_a(\zh)=0$; by [I,~Lemma~1.7],
$\zeta(a)=0<\bar\vartheta\Phi_a^2$ and $w(a)=0$.  (ii) Choose the signs of every
absorber as in (i) at the row where all tuning masses equal their $\theta_a$.
Each absorber's value depends at first order only on its own $m_0$; the other
masses enter through $\nu$ at second order (bounded as in (i)), so the joint
system is solved by the explicit fixed point of Lemma~\ref{lem:II-TU}, each
coordinate staying in its interval by the margins of (i).  (iii) Lemma~\ref{lem:II-kappa} and the
second expression of $\varkappa$: a strict non-peak with $\mathsf v=0$ contributes
nothing.  (iv) Lemma~\ref{lem:II-cost} and Lemma~\ref{lem:II-base}(b).
\end{proof}

For a biased pair $(a^+,a^-)$ for $s$ and a number $r$ (the \emph{residue} at $s$)
put $c_0:=\lambda_{a^-}2^{-s_{\rm far}(a^-)}$ and
\[
 \gamma_{a^+}:=c_0+\frac{(-r)_+}{|u_{a^+}(s)|},\qquad\gamma_{a^-}:=c_0\frac{|u_{a^+}(s)|}
 {|u_{a^-}(s)|}+\frac{r_+}{|u_{a^-}(s)|} .
\]
Since $u_{a^\pm}(s)$ has the sign $\pm$, $\gamma_{a^+}u_{a^+}(s)+\gamma_{a^-}u_{a^-}(s)=-r$;
both coefficients are $\ge c_0/4$ and Lipschitz in $r$.  On the far parts of
$S_{a^\pm}$ the own term $\gamma_av_a(s')$ dominates the terms of all later carriers
(by (W4$''$) the ratio is $\le C_f2^{s_{\rm far}(a)-s'}T_{\rm lo}^32^{1+k(a)}<\frac12$),
so $z=+1$ there is admissible whatever $r$ is.  A tuned absorber that is not
used to cancel a residue receives the constant coefficient $c_0(a)$; with the
coefficient $0$ the far part of its signature set, where $z=+1$, would carry only
later contributions of uncontrolled sign.

\subsubsection*{The coarse shifted cone, activity classes, and one shift ray}

In the rest of this subsection the design is $D^{U1'}$, $N$ is fixed, $f\in
S_{p_N^*}$ has finite base support, $g\in\Cset(f)$, $\rho<1$, and $w$ is a clean
sub-window of a main stage $L\ge l_f$ for $f$.  $K$ denotes quantities of the
form $C_f(\mathrm{Des}(L)/u(w))^C$ with a numerical constant $C$; for $L\ge
l_f$ every constant of this subsection is at most $C_f\bar D(w)$
(Definition~\ref{def:II-DU1}), and such constants are absorbed directly by
the window, because $T_{\rm hi}(w)\le2^{-L^3}/(LQ(w))$, $n(w)\ge L2^{L^3}Q(w)$
and $Q(w)\ge\bar D(w)^3$.  The
shift components of all data considered below are bounded by an $f$-constant
$\Delta_{\max}$ (Lemma~\ref{lem:II-shift} and the remark following it; a
non-degenerate peak of $f$ persists at the companions).

\begin{definition}[The coarse shifted cone]\label{def:II-Gamma}
Let $f^\#$ be the companion of Theorem~\ref{thm:II-B} for the closed pattern
$\kappa(w)$ (statuses as in Theorem~\ref{thm:II-B}(a)).  Let $\Omega$ be the set of coarse strict non-peaks of
$f^\#$ other than absorbers (kept, clamped and class-R ones; nearly neutral kept
carriers have value $0$ after Theorem~\ref{thm:II-B}), $\Omega_m$ its part in block $m$,
and $P^{\rm c}_m$ the coarse peaks of block $m$, with signs $\mathrm{vs}_l:=\sgn w^\#_m(k(l))$.
For $(\Delta',\gamma)\in\R^I\times\R^\Omega$ put
\[
 \mathcal L(\Delta',\gamma):=\Big(\sum_{l\in\Omega}\gamma_lu_l-\sum_m\Delta'_m\sum_{l\in P^{\rm c}_m}
 \mathrm{vs}_l\lambda_lu_l\Big)\Big|_{E_c(w)\setminus F^\#} .
\]
For an \emph{activity class} $a$ (Lemma~\ref{lem:II-classes}), $\Gamma^\#(a)$ is the
polyhedral cone of all $(\Delta',\gamma)$ with:
\begin{enumerate}
\item[(X1)] $\mathcal L(\Delta',\gamma)$ is $z^\#$-admissible on $E_c(w)\setminus F^\#$;
on each set $S_l\cap(L,s_{\rm far}(w)]$, $l\le L$, where only $l$ lives among the
coarse carriers, the rows of (X1) are positive multiples of one sign row and are
replaced by that row with the coefficient $\|v_l1_{S_l\cap(L,s_{\rm far}(w)]}\|_1$;
\item[(X3)] $\Delta'_m\varkappa^\#_m=\sum_{l\in\Omega_m}u_l(\zh^\#)\gamma_l$ for every $m$, where
$\varkappa^\#_m=\varkappa(\hat\zeta^\#_m;\Omega_m)$; in the \emph{ratio variables} $r_l:=
u_l(\zh^\#)/\varkappa^\#_{m(l)}$ this reads $\Delta'_m=\sum_{l\in\Omega_m}r_l\gamma_l$;
\item[(X4)] $\mathrm{vs}_l\gamma_l/\lambda_l+\Delta'_{m(l)}\ge0$ for every $l\in\Omega$ with
$\gap^\#(k(l))<t^2$ for some $t\in\mathcal W(w)$ (near-threshold carriers);
\item[(X5)] $\gamma_l=0$ and $\Delta'_m=0$ for the components declared inactive in $a$.
\end{enumerate}
The rows other than (X3) have design coefficients (the numbers $\lambda_lu_l(j)$,
the aggregated coefficients, signs, $1/\lambda_l$).  The minors of the matrix of
$\Gamma^\#(a)$ in the ratio variables, and the minors of the matrices $g_J$ formed
by the $\Delta'$-parts of the cofactor vectors of its $(n-1)$-row subsystems, are
polynomials in $r$ with design coefficients; for every pattern and every
class of level $L$ they are added to the rate scheme of $D^{U1'}$, with values
computed at the ratios of $f$.
\end{definition}

For the actual decompositions at a scale $t\in\mathcal W(w)$ we read off the coarse
data (data convention, compare Lemma~\ref{lem:II-R17}): $\Delta'_{\rm dec,m}(t):=
-\Delta d_mM_m$ for $m\in I_{\rm sh}(w)$ and $:=0$ otherwise (for $m\notin I_{\rm sh}(w)$
the shift is $\le K_dt$ by Lemma~\ref{lem:II-sources}), and $\gamma_{\rm dec,l}(t):=
-\Delta\theta_l$ for $l\in\Omega$.

\begin{lemma}[Activity classes]\label{lem:II-classes}
Let $J:=|\Omega|+N$; since $L\ge l_f$, every block $m\le N$ has a carrier
$\le L$, so $N\le L$ and $J\le2L$.  Let $H(L)\le C\mathrm{Des}(L)^C/u(w)$ be
the bound for the $\ell_1$-Hoffman constants of all cones $\Gamma^\#(a)$ at
the exactified point \textup{(}Proposition~\textup{\ref{prop:II-ray}(b)},
Proposition~\textup{\ref{prop:II-KN}}\textup{)}, let $K_*\ge\max\{2\mathrm{Des}
(L),4C_f\mathrm{Des}(L)JH(L)\}$ with $K_*\le(\mathrm{Des}(L)/u(w))^C$, where
$C_f\mathrm{Des}(L)$ bounds the entries of the rows of $\Gamma^\#(a)$, and
put $A_1:=K_*$, $A_{i+1}:=K_*A_i$.  For every scale $t$ one of the $J+1$ intervals $[A_iKt,
A_{i+1}Kt)$, $1\le i\le J+1$, contains the modulus of no component of $X(t):=
(\Delta'_{\rm dec}(t),\gamma_{\rm dec}(t))$; let $i(t)$ be the least such $i$.  The
\emph{class} of $t$ is $a(t):=(i(t),\ \text{the set of components with modulus}\ge
A_{i(t)+1}Kt,\ \text{their signs})$; the other components are \emph{inactive} and
are set to $0$, at an $\ell_1$-cost $\le JA_{i(t)}Kt$.  There are at most $(J+1)3^J$
classes.  Zeroing the inactive components and projecting onto $\Gamma^\#(a(t))$ moves
no active component below $A_{i(t)+1}Kt/2$.  For $L\ge l_f$ the number of
classes, $K_*^{J+2}$, and all constants $A_i$ $(i\le J+2)$ are at most
$\bar D(w)$.
\end{lemma}

\begin{proof}
Pigeonhole ($J$ moduli, $J+1$ disjoint intervals).  Zeroing the inactive
components changes the violation of the rows of $\Gamma^\#(a(t))$ by at
most $C_f\mathrm{Des}(L)\cdot JA_{i(t)}Kt$, so the violation after zeroing is
at most $(1+C_f\mathrm{Des}(L)JA_{i(t)})Kt$ (Proposition~\ref{prop:II-ray}(a)).
The projection moves components by at most $H(L)$ times this violation,
i.e.\ by at most $H(L)(1+C_f\mathrm{Des}(L)JA_{i(t)})Kt\le\frac12K_*A_{i(t)}
Kt=\frac12A_{i(t)+1}Kt$, because $A_{i(t)}\ge1$ and $K_*\ge4C_f\mathrm{Des}
(L)JH(L)$.  The bounds by $\bar D(w)$ follow from $J\le2L$, $(J+1)3^J\le
(2L+1)3^{2L}$ and $K_*\le(\mathrm{Des}/u)^C$, with $C_\Gamma$ large.  The
inactive components have modulus $<A_{i(t)}Kt\le A_{i(t)+1}Kt/K_*$.  An inward row
(X4) coupling a small $\gamma_d$ with a large $\Delta'_m$ would force $\lambda_d\varrho_d
K_*\ge1$, which is excluded by $K_*\ge2\mathrm{Des}(L)\ge2/\lambda_d$.
\end{proof}

A class is \emph{one-signed} if all its active shift components $\Delta'_m$ have
the same sign or none is active, and \emph{mixed} otherwise (possible only for
$N\ge2$).  For a one-signed class with active shifts, \emph{configuration (i)}
means $\Delta'<0$ (the active blocks lack an upper source) and
\emph{configuration (ii)} means $\Delta'>0$.

\begin{proposition}[Projection and one shift ray]\label{prop:II-ray}
Suppose that at the companion $f^\#$ every minor of Definition~\textup{\ref{def:II-Gamma}}
is $0$ or $\ge u/(2\mathrm{Des}^C)$.  Then for every scale $t\in\mathcal W(w)$, with
$a=a(t)$:
\begin{enumerate}
\item[(a)] $X(t)$ violates \textup{(X1)}, \textup{(X3)}, \textup{(X4)} at the
coefficients of $f^\#$ by at most $Kt$, and after the inactive components
are set to $0$ by at most $(1+C_f\mathrm{Des}\,JA_{i(t)})Kt$;
\item[(b)] the $\ell_1$-projection $X_{\rm pr}(t)$ onto $\Gamma^\#(a)$ of $X(t)$
with inactive components set to $0$ satisfies $\|X_{\rm pr}(t)-X(t)\|_1\le
(1+C_f\mathrm{Des}\,JA_{i(t)})Kt$ \textup{(}after enlarging $K$ by the factor
$H(L)$\textup{)}, and no active component is moved below $A_{i(t)+1}Kt/2$
\textup{(}last sentence of Lemma~\textup{\ref{lem:II-classes})};
\item[(c)] \textup{(one ray)} let $r_1,\dots,r_p$ be the extreme rays of $\Gamma^\#(a)$,
$\|r_j\|_1=1$, and $g_j$ their nonzero $\Delta'$-parts; write $\Delta'_{\rm pr}(t)=\sum_j
a_j(t)g_j$ by a fixed rule \textup{(}a basic solution\textup{)}, $a(t)\in[0,A_{\max}]^p$
with $A_{\max}\le K$.  For a cube $\mathcal Q$ of side $\varepsilon$ in $[0,A_{\max}]^p$
with corner $a^{\max}$ \textup{(}coordinatewise maximum\textup{)} and every scale $t$ with
$a(t)\in\mathcal Q$ put
\[
 \Delta'^*:=\sum_ja^{\max}_jg_j,\qquad\mathrm{incr}(t):=\sum_j(a^{\max}_j-a_j(t))r_j\in
 \Gamma^\#(a);
\]
then $\Delta'_{\rm pr}(t)+\Delta'(\mathrm{incr}(t))=\Delta'^*$ and $\|\mathrm{incr}(t)\|_1\le
p\varepsilon$;
\item[(d)] for $\varepsilon:=c_f(u/\mathrm{Des})^C$ the number of pairs \textup{(}class,
cube\textup{)} is at most $D_{\rm cls}(w):=(J+1)3^J(A_{\max}/\varepsilon)^p$, and
$D_{\rm cls}(w)\le\bar D(w)$ for $L\ge l_f$;
\item[(e)] $|\Delta'^*_m|\ge|\Delta'_{\rm pr,m}(t)|$ with the same sign, and the signs
$\epsilon_l\gamma_l\ge0$ of the class-G components are preserved by the increments.
\end{enumerate}
\end{proposition}

\begin{proof}
(a) As in Steps~1--2 of Proposition~\ref{prop:II-transplant}: by the switching budget
[I,~Lemma~8.24], $\sum_{j\notin F}\phi_{z_j}(\Delta B(j))\le t/q_0$ with $\Delta B=
-\sum_l\Delta\theta_lu_l$, and on $E_c(w)$ the coarse part of $\Delta B$ equals
$\mathcal L(X(t))$ up to: the peak errors $e_k\mathrm{vs}\lambda\le C_f\mathrm{Des}\,t/u$
([I,~(16)]; tiny-margin peaks of $f$ are strict non-peaks of $f^\#$, hence in
$\Omega$); fine carriers ($\le3b^2/t$); the difference between $z$ and $z^\#$ on
$E_c(w)$ (closed tiny rooms, as in Proposition~\ref{prop:II-transplant}); and the change
of the coefficients from $f$ to $f^\#$ (values and $\varkappa$ move by $\le C\mathrm{Des}\,b
+\eta(w)$, times $\|\gamma\|_1\le12\sum_l\lambda_l/t$, which is $\le T_{\rm lo}(w)^3$).  (X3) is
[I,~(17)] with the peak traces, as in the proof of Lemma~\ref{lem:II-kappa}.  (X4):
for a near-threshold strict non-peak of $f$, [I,~Lemma~8.8] gives
$\mathrm{vs}\,(\gamma/\lambda)+\Delta'\ge-(4\gap(f)/t+C_fb)$, and at a peak of $f$,
[I,~(16)] gives $\mathrm{vs}(\omega_--\omega_+)(k)=e_k\ge0$; $\gap(f)\le bM$.  A class-R
near-threshold strict non-peak (or a dropped anti-type near-threshold strict
non-peak) of a block pins $\Delta'_{\rm dec}\ge-C_f\mathrm{Des}K_gt$ ([I,~Lemma~8.8] with
$|\Delta\theta|\le K_gt$), so it does not occur in an active block with $\Delta'<0$, and
in an active block with $\Delta'>0$ its row holds with $\gamma=0$.  (b)
Lemma~\ref{lem:II-hoffman}: after the aggregation of (X1) the nonzero minors are
$\ge u/(2\mathrm{Des}^C)$ and the entries are $\le C_f\mathrm{Des}$, so the Hoffman
constant is $\le C\mathrm{Des}^C/u$.  (c) By Lemma~\ref{lem:II-donorpeak} every active block
has a robust coarse peak; it is of class G (a class-R peak is a source of both
kinds), of swallowing type in a block with $\Delta'<0$ and of anti type in a block
with $\Delta'>0$ (Lemma~\ref{lem:II-Dprime}), and (X1) on its near signature set fixes
$\sgn\Delta'_m$ on the whole cone; (X1) on natural signature sets fixes the signs of
the class-G components of $\gamma$, and (X1) at free coordinates forces $\gamma=0$ for
class-R components.  So $\Gamma^\#(a)$ lies in an orthant, is pointed, and is the
cone generated by its extreme rays (Minkowski--Weyl); its $\Delta'$-projection is
the cone generated by $g_1,\dots,g_p$.  A basic representation uses linearly
independent $g_J$ with $a_J=g_J^+\Delta'_{\rm pr}(t)$; by Cauchy--Binet $\sigma_{\min}(g_J)$
is bounded below by a robust minor of $g_J$ over a polynomial bound of its
entries, so $A_{\max}\le C_f(\mathrm{Des}/u)^C\Delta_{\max}$.  $\mathrm{incr}(t)$ is a nonnegative
combination of extreme rays.  (d) Count.  The number $p$ of extreme rays is
at most $P(L)$ (a pointed cone in $\R^n$ cut out by at most $R_\Gamma(L)$
inequalities has at most $\binom{R_\Gamma(L)}{n-1}\le2^{R_\Gamma(L)}$ extreme
rays), so $(A_{\max}/\varepsilon)^p\le(C_f\Delta_{\max}(\mathrm{Des}/u)^{2C})^{P(L)}
\le(\mathrm{Des}/u)^{(2C+1)P(L)}$ as soon as $C_f\Delta_{\max}\le\mathrm{Des}(L)/
u(w)$, which holds for $L\ge l_f$; together with $(J+1)3^J\le(2L+1)3^{2L}$
this gives $D_{\rm cls}(w)\le\bar D(w)$ for $C_\Gamma\ge2C+1$.  (e) follows from the orthant
structure in (c).
\end{proof}

\subsubsection*{Exactification with threshold protection under (KN)}

An $\Omega$-carrier is \emph{near-threshold} if $|\varrho-1|\le b(w)$ at $f$, and
\emph{active in the class $a$} if its component $\gamma_l$ is active.

\begin{definition}[The lever hypothesis (KN$_{w,a}$)]\label{def:II-KN}
\textup{(KN$_{w,a}$)}: every block which is active in $a$ and contains an
$\Omega$-carrier that is near-threshold and active in $a$ also contains an
$\Omega$-carrier $d_0$ that is inactive in $a$ and has $\varrho_{d_0}\ge u(w)$ at $f$.
\end{definition}

\begin{proposition}[Exactification with threshold protection]\label{prop:II-KN}
Let $w$ be a clean sub-window of a main stage $L\ge l_f$ for $f$, and let $a$ be
a class satisfying \textup{(KN$_{w,a}$)}.  Then there is a companion $f^{(1a)}$ of $f$
with:
\begin{enumerate}
\item[(a)] the ratios $r_l$ of the active $\Omega$-carriers equal a point $r'$ at
which every minor of Definition~\textup{\ref{def:II-Gamma}} that is tiny at $f$
vanishes and every robust one is $\ge u/2$;
\item[(b)] every active near-threshold $\Omega$-carrier is a strict non-peak with
gap $\ge c_f\eta(w)\mathrm{Des}(L)M$; every inactive near-threshold $\Omega$-carrier is a
strict non-peak with gap $\ge c_f\eta(w)M$; every other coarse carrier keeps its
robust status;
\item[(c)] $p^*(f^{(1a)}-f)\le C_f\mathrm{Des}\,T_{\rm lo}(w)^4\log(1/T_{\rm lo}(w))=o(T_{\rm
lo}(w)^2)$.
\end{enumerate}
\end{proposition}

\begin{proof}
\emph{Step 0.}  Apply only the closing moves (C1), (C2) of the assembly (no donor
raise in active blocks); by Lemma~\ref{lem:II-status} values move by $\le C_*(l)b$,
and $\varkappa$ by $\le C\mathrm{Des}\,b$ (it is Lipschitz in the block data,
Lemma~\ref{lem:II-kappader}).  \emph{Step 1.}  The tiny minors are $\le\beta(w)$ at
the ratios $r(f_1)$; Lemma~\ref{lem:II-loj} in the ratio variables gives $r'$
with $|r'-r(f_1)|\le C_L\beta^{2/N_L}\le\eta(w)$, and robust minors stay $\ge u/2$.
\emph{Step 2 (realization).}  Apply Lemma~\ref{lem:II-TU} to the active
$\Omega$-carriers with target values $u_l:=r'_l\varkappa_m$ ($\varkappa_m$ held at its current
value).  In every active block $m$ with an active near-threshold carrier, push
the lever $d_0$ of (KN$_{w,a}$) inward (pulls or z-moves on its far signature
coordinates, two-sided by Lemma~\ref{lem:II-TU}) by
$s_0:=C\mathrm{Des}(L)\eta(w)\bar\vartheta_m\Phi_{P_m}^2/(\varrho_{d_0}M_m)$ in units of
$\hat\zeta_m$; then push a peak to restore $\varkappa_m$ exactly (it is continuous,
Lemma~\ref{lem:II-kappader}(a), and the required push is $O(X\cdot s_0)$, with $X\le
C\mathrm{Des}^2\sum_{\rm fine}\Phi^2$); then re-solve the tuning.  The Jacobian of
(ratios, $\varkappa$) with respect to (tuning parameters, pushes) is block
triangular with identity and diagonal blocks, so the explicit fixed point of
Lemma~\ref{lem:II-TU} extended by one scalar equation per block converges.  This
gives (a).  \emph{(b).}  By Proposition~\ref{prop:II-Rkt}(iii) and
Lemma~\ref{lem:II-kappader}, $\bar\vartheta_m$ rises by $\ge c\varrho_{d_0}M s_0/\Phi_P^2\ge
C'\mathrm{Des}\,\eta\bar\vartheta_m$, while the active values move by $\le\eta$ in $u$-units,
i.e.\ by $\le m\eta\mathrm{Des}$ in units of $\mathsf v$; so an active carrier with
$\varrho\le1+b$ ends with $\varrho^\#\le(1+b+C\mathrm{Des}\,\eta)(1-C'\mathrm{Des}\,\eta)<1-
c\mathrm{Des}\,\eta$ for $C'$ large ($b\le\eta$).  Robust carriers move by $O(\mathrm{Des}\,
\eta)\ll u$.  The capacity of $d_0$ suffices: $|\hat\zeta_m(d_0)|=\varrho_{d_0}\bar\vartheta_m
\Phi_{d_0}^2\ge u\bar\vartheta_m/\mathrm{Des}^2>s_0$.  Inactive near-threshold carriers are not
variables of the minors; they are pushed inward individually by $\eta$, and
$\varkappa$ is restored as above.  (c) Lemma~\ref{lem:II-TU} and the cost bound of
Lemma~\ref{lem:II-cost}, with all moves $\le C\mathrm{Des}\,\eta$, $\eta=T_{\rm lo}^4/(L\mathrm{Des})$.
\end{proof}

In a class without active shift no lever is needed: only values enter, and
Theorem~\ref{thm:II-B} applies.

\subsubsection*{The companion and its fine structure}

Fix a class $a$ with active blocks $A\subseteq I_{\rm sh}(w)$, signs $\sigma_m:=\sgn
\Delta'_m$ $(m\in A)$, active $\Omega$-carriers $\Omega_{\rm act}$ and the normalized
reference data $(\Delta'(t),\gamma(t))\in\Gamma^\#(a)$ of
Proposition~\ref{prop:II-ray} for the scales $t$ of one (class, cube) set $\mathcal S$.

\begin{definition}[The companion $f^\#$ of a class]\label{def:II-compclass}
$f^\#$ is obtained from $f$ by the following moves, which act on pairwise
distinct coordinates.
\begin{enumerate}
\item[(1a)] Proposition~\ref{prop:II-KN} (for a class without active shift:
Theorem~\ref{thm:II-B}), producing $f^{(1a)}$ with the exact ratios $r'$ and
$\varkappa'_m:=\varkappa_m(f^{(1a)})$.
\item[(1b)] For every $s\in E_c(w)\setminus F^{(1a)}$ the first pair $P(s)$ of
Lemma~\ref{lem:II-firsttouch}, and all absorbers of the cluster of $L$, are tuned
to value $0$ (Lemma~\ref{lem:II-zerovalue}).
\item[(1c)] $\varkappa_m$ and the ratios of $\Omega_{\rm act}$ are restored to $(\varkappa',r')$
and the tuned absorbers re-tuned to value $0$ (joint explicit fixed point; each
target carrier is controlled at first order by its own masses, cross effects
are of second order; $\varkappa$ is continuous, Lemma~\ref{lem:II-kappader}(a)).
\item[(2)] \emph{Release and fine structure.}  For every coarse carrier $l''$ that is
not an owner candidate below (inactive $\Omega$-carriers, and peaks of inactive
blocks), the far part $S_{l''}\cap(s_{\rm far}(w),\infty)$ is released.  Let $J_{\rm fine}$
consist of the coordinates outside $F^{(1b)}\cup E_c(w)$ not fixed in (1a), (1b),
together with the released ones.  On $J_{\rm fine}$, $z$ is defined by the
recursion of Lemma~\ref{lem:II-recursion} in configuration (i), and by the fixed
point of Lemma~\ref{lem:II-schauder} otherwise.
\end{enumerate}
Step (2) does not change the base part, hence not $e^\#$.  The release changes
the values of the released carriers by $\le4c_{L+1}^2$ (these are not variables
of $\Gamma^\#(a)$), and $\varkappa$ by $\le C\mathrm{Des}\,c_{L+1}^2$.
\end{definition}

The data at $f^\#$ use the switching coefficients $\gamma_l(t)$ on $\Omega_{\rm act}$,
$0$ on the inactive $\Omega$-carriers, the biased coefficients on the used pairs
$P(s)$, and the constant coefficient $c_0(a)$ on the other tuned absorbers.  By
Lemma~\ref{lem:II-cert} the forced coefficient of every other carrier $k$ of
block $m$ is $-\Delta''_m\lambda_kw^\#_m(k)$, where the \emph{true shift} $\Delta''(t)$ is given
by \eqref{eq:II-X3} at $f^\#$; tuned absorbers contribute nothing to it
(Lemma~\ref{lem:II-zerovalue}(iii)).

\begin{lemma}[The true shift]\label{lem:II-trueshift}
\textup{(a)} $|\Delta''_m(t)-\Delta'_m(t)|\le C_f\big(|\varkappa^\#_m-\varkappa'_m|+\|u^\#-u'\|_{\ell_\infty
(\Omega_{\rm act})}\|\gamma(t)\|_1\big)\le C_f\mathrm{Des}\,c_{L+R_{\rm cl}(L)+1}+C_fc_{L+1}^2/t$.  After
\textup{(1c)} the values of $\Omega_{\rm act}$ are not changed by \textup{(2)}, so
$\Delta''(t)=\Delta'(t)\varkappa'_m/\varkappa^\#_m$; for the normalized data of a set $\mathcal S$
this is the same vector $\Delta''^*$ for all $t\in\mathcal S$.
\textup{(b)} $\mathcal L(\Delta''(t),\gamma(t))-\mathcal L(\Delta'(t),\gamma(t))=-\sum_m(\Delta''_m-\Delta'_m)
\mathcal P_m$, $\mathcal P_m:=\sum_{l\in P^{\rm c}_m}\mathrm{vs}_l\lambda_lu_l|_{E_c(w)}$.  On the near
signature sets of the coarse peaks this keeps the sign of the peak trace, so
$\mathcal L(\Delta'',\gamma)$ is admissible there; on $T(L)$ the difference is absorbed by
the cluster pairs (Proposition~\ref{prop:II-absorb}).
\textup{(c)} The signs and the activity of $\Delta''(t)$ are those of $\Delta'(t)$.
\end{lemma}

\begin{proof}
(a) Subtract \eqref{eq:II-X3} at $(u',\varkappa')$ from \eqref{eq:II-X3} at $(u^\#,\varkappa^\#)$ and
divide by $\varkappa^\#\ge A^\#\ge c_f$; $\sum u'\gamma=\Delta'\varkappa'$.  After (1c), $|\varkappa^\#-\varkappa'|$
is bounded by the effect of later moves: only carriers at stages $>L+R_{\rm cl}$ can
change status (cluster carriers are tuned to value $0$ with gap $M$, coarse
carriers have protected margins and gaps), step (2) changes the block data by
$\le C(\sum_{l>L+R_{\rm cl}}\lambda_l+2^{-s_{\rm far}(w)})$, and $\varkappa$ is Lipschitz in the
block data with constant $\le\mathrm{Des}/u$.  The values $u_l$, $l\in\Omega_{\rm act}$, move in
(1b) by second-order effects $\le Cc_{L+1}^2$ and are restored in (1c); their
vectors meet $J_{\rm fine}$ only in their own far signature parts, which are owned
by $l$ with sign fixed since (C1).  (b) $\mathcal L$ is linear in $\Delta'$, and on the
natural signature set of a coarse peak only that peak lives among coarse
carriers; $|\Delta''-\Delta'|<|\Delta'|/2$ for active components by (c).  (c) Active
components have modulus $\ge A_{i+1}Kt/2$, much larger than the bound in (a),
since $c_{L+1}\le T_{\rm lo}(w)^8$.
\end{proof}

\begin{proposition}[Exact absorption on the coarse coordinates]\label{prop:II-absorb}
For $s\in E_c(w)\setminus F^\#$ let the residue $\varrho(s)$ be $\mathcal V(\Delta\omega)(s)-\mathcal L(\Delta'',
\gamma)(s)$ minus the contribution of $P(s)$.  Then:
\begin{enumerate}
\item[(a)] $\varrho(s)$ is the sum of the terms of the carriers after $P(s)$, so $|\varrho(s)|\le
C_f\sum_{k\text{ after }P(s)}\lambda_k$; on $T(L)$ the correction of
Lemma~\textup{\ref{lem:II-trueshift}(b)}, of size $\le C_f\mathrm{Des}\,c_{L+1}$, is added to
it;
\item[(b)] with the biased coefficients for $r:=\varrho(s)$ \textup{(}plus correction\textup{)},
$\mathcal V(\Delta\omega)(s)=\mathcal L(\Delta'',\gamma)(s)$ exactly, which is $z^\#$-admissible;
\item[(c)] \textup{(capacity)} $|\Delta\omega(a)|=|\gamma_a|/\lambda_a\le A_2/t$ for every used
absorber $a$, which is a strict non-peak with $w(a)=0$ and gap $M$ (kind
\textup{[2]});
\item[(d)] at the support coordinates of the absorbers the data satisfy $t|b|\le$
the mass, and the masses cost $\le C_fc_{L+1}T_{\rm hi}(w)+Cc_{L+1}^2$;
\item[(e)] the pairs have no $d$-effect.
\end{enumerate}
\end{proposition}

\begin{proof}
(a) Lemma~\ref{lem:II-firsttouch} and Lemma~\ref{lem:II-cert}(a); the coefficients of
later carriers are $\le C_f\lambda$ ($|\Delta''|\le\Delta_{\max}$, and later absorbers have
$|\gamma|\le C_f\lambda$).  (b) The biased pair.  (c) Cluster pairs have $\lambda_a\ge2^{-1-
k(a)}4^{1-R_{\rm cl}}c^{\rm low}_{L+1}\ge c_{L+1}/\mathrm{Des}(L)$, while the residues are $\le
2c_{L+R_{\rm cl}+1}\le2b(L+R_{\rm cl})^2\ll c_{L+1}/\mathrm{Des}$ plus the correction $\le
C_f\mathrm{Des}\,c_{L+1}$, so $|\gamma_a|/\lambda_a\le C_f\mathrm{Des}^2\ll A_2/T_{\rm hi}(w)$.  A pre-pair
at a stage $p>L$ sees residues only from carriers after it, of total weight
$\le2c_{p+1}\le2b(p,M(p))^2\ll\lambda_a$ (by (W2) and $n(p)\ge\mathrm{Des}(p)\ge2^{1+k(a)}/
c_p$), and $s\notin T(L)$.  (d) $b(p_r)=\gamma_au_a(p_r)$ and $t|\gamma_a|\le C_f\lambda_aT_{\rm hi}\le$
the mass.  (e) Lemma~\ref{lem:II-zerovalue}(iii).
\end{proof}

\begin{lemma}[Ownership and the robust recursion]\label{lem:II-recursion}
Call a carrier an \emph{owner candidate} if its coefficient in $\mathcal V(\Delta\omega)$ is
nonzero with a sign known in advance: \textup{(a)} $l\in\Omega_{\rm act}$ \textup{(}sign
$\epsilon_l$\textup{)}; \textup{(b)} coarse peaks of active blocks \textup{(}coefficient
$-\Delta''_m\mathrm{vs}_l\lambda_l$\textup{)}; \textup{(c)} all tuned absorbers \textup{(}coefficients
$\ge c_0/4>0$\textup{)}; \textup{(d)} carriers $k$ at stages $>L$ of active blocks
other than tuned absorbers \textup{(}coefficient $-\Delta''_m\lambda_kw(k)$, sign
$-\sigma_m\varsigma_k$, $\varsigma_k:=\sgn w(k)$ to be chosen\textup{)}.  For $j\in J_{\rm fine}$
let $\mathrm{own}(j)$ be the least owner candidate meeting $j$; if there is none,
$\mathcal V(\Delta\omega)(j)=0$ and $z_j$ is free.  For carriers $k$ of type \textup{(d)},
in increasing order, let $O(k):=\{j\in J_{\rm fine}:\mathrm{own}(j)=k\}\supseteq S_k$, let
$Y_k$ be $u_k(\zh)$ computed with the coordinates of $\supp u_k\setminus O(k)$ only,
$\varsigma_k:=\sgn Y_k$ \textup{(}$+1$ if $Y_k=0$\textup{)}, and $z_j:=\sgn(-\sigma_m\varsigma_ku_k(j))$ on
$O(k)$; on the sets owned by \textup{(a)}--\textup{(c)} put $z_j:=\sgn(\text{coefficient}
\cdot u_{\rm own}(j))$.  In configuration \textup{(i)} \textup{(}$\sigma\equiv-1$\textup{)}
this gives $z_j=\varsigma_k\sgn u_k(j)$ on $O(k)$ and
\[
 |u_k(\zh^\#)|=|Y_k|+\sum_{j\in O(k)}|u_k(j)|\ge\delta^\circ_k ,
\]
so every carrier of type \textup{(d)} is a peak of sign $\varsigma_k$ with margin
$\ge q^\#_0\delta^\circ_k/2$.
\end{lemma}

\begin{proof}
The recursion is well founded, since $Y_k$ uses only coordinates fixed before
$k$.  On $J_{\rm fine}$, $\zh^\#_j=z_j$ (diagonal base, $J_{\rm fine}\cap F^\#=\emptyset$).
With $\sigma=-1$, $z_ju_k(j)=\varsigma_k|u_k(j)|$ on $O(k)$.  By (W3), $\Phi_k\le c_k\le
(\delta_kH_k)^2$, so $\bar\vartheta\Phi_k/m\ll\delta^\circ_k$, and [I,~Lemma~1.7] gives the
peak and the margin.
\end{proof}

\begin{lemma}[Dominance on the fine coordinates]\label{lem:II-domfine}
In configuration \textup{(i)}, and for the owners of types \textup{(a)}--\textup{(c)} in
every class, for every $j\in J_{\rm fine}$ owned by $k$
\[
 |\mathrm{coef}_k|\,|u_k(j)|\ge2\sum_{k'>k,\ u_{k'}(j)\ne0}|\mathrm{coef}_{k'}|\,|u_{k'}(j)| ;
\]
hence $\sgn\mathcal V(\Delta\omega)(j)=\sgn(\mathrm{coef}_ku_k(j))$ and $z_j\mathcal V(\Delta\omega)(j)>0$.  The
bound holds for every shift and every switching of the class.
\end{lemma}

\begin{proof}
Later coefficients are $\le C_f\lambda$.  (i) $j\in S_k$: later carriers meet $j$ only
through targets, and (W4$''$) gives $c_{k'}\le\iota_k2^{-2j}c_k\delta_kT_{\rm lo}(k-1)^3/2$; their
total is $\le C_f2^{-2j}c_k\delta_kT_{\rm lo}(k-1)^3$, while the own term is $\ge|\mathrm{coef}_k|
\delta_k2^{-j}/n_k$ with $|\mathrm{coef}_k|\ge cA_2KT_{\rm lo}(w)\lambda_k$ for types (a), (b), (d)
and $\ge c_0/4$ for type (c).  With $T_{\rm lo}(k-1)\le T_{\rm lo}(w)$, $\lambda_k=m2^{-m-k(k)}c_k$
and $j\ge3\cdot2^{k(k)}$, the ratio is $<\frac12$; for type (c), $j>s_{\rm far}(a)$ and
$c_0=\lambda_a2^{-s_{\rm far}(a)}$ give a ratio $\le C_f2^{m+k-(j-s_{\rm far}(a))}T_{\rm lo}^2<\frac12$.
(ii) $j\in\supp y_k\setminus S_k$: the own term is $\ge|\mathrm{coef}_k|\underline y_k/n_k$, the later
terms are $\le C_f\sum_{k'>k}\lambda_{k'}\le C_fc_{k+1}\le C_fb(k,M(k))^2\le C_f2^{-8n(k)}$,
and $n(k)\ge\mathrm{Des}(k)\ge1/(\underline y_k\lambda_k)$.  If $k$ is an absorber inside a cluster
or a pre-sequence, only its partner meets $j$ inside the block of consecutive
block-$1$ stages, and $\lambda$ decreases by a factor $\ge8$ between consecutive
block-$1$ stages, so the partner's term is $\le\frac14$ of $k$'s; the first later
non-absorber stage $p'$ has $c_{p'}\le b(p'-1,M(p'-1))^2$, and all later carriers
together give a ratio $\le\frac14+C_f2^{-n(p'-1)}/(T_{\rm lo}(w)\underline y_k\lambda_k)<\frac12$.
Owners of types (a), (b) have no target coordinates in $J_{\rm fine}$.
\end{proof}

\begin{lemma}[Scrambling in configuration (i)]\label{lem:II-SCi}
In configuration \textup{(i)} the set $A$ of active blocks satisfies \textup{(SC)} at
$f^\#$, along any sequence $s_i\downarrow0$.
\end{lemma}

\begin{proof}
Fix $m\in A$ and $\Upsilon\ge1$.  Coarse carriers are finitely many, peaks have relative
margin $\ge u/2$, strict non-peaks have gaps bounded below
(Proposition~\ref{prop:II-KN}(b)), and there is no degenerate peak; tuned
absorbers are finitely many strict non-peaks with gap $M$.  For $s$ below a
positive bound none of them contributes to $\mathrm{Scr}_m(\Upsilon s)$.  Every other
carrier of block $m$ is a peak with margin $\ge q^\#_0\delta^\circ_k/2$
(Lemma~\ref{lem:II-recursion}) and contributes only if $\delta_kH_k\le2n_k\Upsilon s/q^\#_0
=:X$, and then at most $\Phi_k\le c_k\le(\delta_kH_k)^2$ by (W3).  The numbers $\delta_kH_k$
decrease super-exponentially, so the sum of their squares over $\{\delta_kH_k\le X\}$ is
$\le2X^2$.  Hence $\mathrm{Scr}_m(\Upsilon s)\le C(\Upsilon s/q_0)^2=o(s)$.
\end{proof}

\begin{lemma}[Exact completion by a fixed point]\label{lem:II-schauder}
Let the data of the class carry the common true shift $\Delta''^*$ and let $J_{\rm
free}\subseteq J_{\rm fine}$ be the complement of the sets owned by the owners of
types \textup{(a)}--\textup{(c)}.  For $z'\in[-1,1]^{J_{\rm free}}$ let $f(z')$ be the
row with these coordinates replaced, $\Delta''(z')$ the true shift at $f(z')$, and
$\mathcal V(z')$ the fine part of $\mathcal V(\Delta\omega)$ at $f(z')$.  Then $\Psi(z')_j:=\mathrm{clamp}_{
[-1,1]}(z'_j+\mathcal V(z')(j))$ has a fixed point $z'^*$, and at $f^\#:=f(z'^*)$ the
vector $\mathcal V(\Delta\omega)$ is $z^\#$-admissible on $J_{\rm fine}$.  This holds in every
class, one-signed or mixed.
\end{lemma}

\begin{proof}
As in Lemma~\ref{lem:II-C5}: the product topology makes $[-1,1]^{J_{\rm free}}$ a compact
convex metrizable set; $z'\mapsto\mathcal V(z')(j)$ is continuous (dominated convergence,
norm-to-weak$^*$ continuity of the duality maps, continuity of $\varkappa^\#$ and of the
absorber coefficients, which are Lipschitz in the residues); Schauder--Tychonoff
gives $z'^*$, and $\mathrm{clamp}(z_j+\mathcal V_j)=z_j$ means $\mathcal V_j=0$ if $|z_j|<1$ and
$z_j\mathcal V_j\ge0$ if $|z_j|=1$.  On the sets owned by (a)--(c) use
Lemma~\ref{lem:II-domfine}.  Signs of the shifts are not used.
\end{proof}

\begin{lemma}[Cost and statuses]\label{lem:II-coststat}
$p^*(f^\#-f)\le C_f\mathrm{Des}\,T_{\rm lo}(w)^4\log\frac1{T_{\rm lo}(w)}+C_fc_{L+1}\log\frac1{c_{L+1}}
=:\theta_wT_{\rm lo}(w)^2$ with $\theta_w\to0$; every coarse carrier has the status
prescribed by the closed pattern, with the gaps of
Proposition~\textup{\ref{prop:II-KN}(b)}; and every robust minor of $\Gamma^\#(a)$
stays $\ge u/(2\mathrm{Des}^C)$.
\end{lemma}

\begin{proof}
(1a): Proposition~\ref{prop:II-KN}(c).  (1b), (1c): $\lambda_a\le c_{L+1}$ and masses
$\le C_fc_{L+1}T_{\rm hi}$.  (2): with $\delta:=\zh^\#-\zh^{(1c)}$ supported in $J_{\rm fine}$,
coarse carriers meet $J_{\rm fine}$ only beyond $s_{\rm far}(w)$, so $|u_{l''}(\delta)|\le
4\delta_{l''}2^{-s_{\rm far}(w)}\le4c_{L+1}^2$, and carriers at stages $>L$ contribute $\le
\sum\lambda\le c_{L+1}$; Lemma~\ref{lem:II-cost}.  $c_{L+1}\le b(L,M(L))^2\le T_{\rm lo}(w)^8$.
The drifts $Cc_{L+1}$ are $\ll c_f\eta(w)M$.
\end{proof}

\subsubsection*{Mixed classes}

In a mixed class put $A_-:=\{m\in A:\sigma_m=-1\}$ and $A_+:=\{m\in A:\sigma_m=+1\}$, both
nonempty.  [I,~Theorem~7.18] needs (SC) exactly for the blocks with $\Delta d_m<0$,
i.e.\ for $A_-$.  By Lemma~\ref{lem:II-schauder}, exact completions exist in every
class; what may fail in a mixed class is (SC) for $A_-$.  For a carrier $l$ of
block $m$ write $\mathrm{thr}_l:=\bar\vartheta_m\Phi_l/m$, so that $l$ is a peak iff $|u_l(\zh)|\ge
\mathrm{thr}_l$.

\begin{lemma}[States of a negative fine carrier]\label{lem:II-states}
Let $z$ be a point at which $\mathcal V(\Delta\omega)$ is admissible on $J_{\rm fine}$, and $l$ a
carrier of a block $m\in A_-$ at a stage $>L+R_{\rm cl}(L)$ which is not a tuned absorber.
Let $O(l)$ be the set of $j\in\supp u_l\cap J_{\rm fine}$ at which no earlier carrier has a
nonzero coefficient \textup{(}$O(l)\supseteq S_l$\textup{)}, $\mathrm{own}_l:=\sum_{j\in O(l)}
|u_l(j)|\ge\delta^\circ_l$, $v_l:=u_l(\zh)$ and $Y_l:=v_l-\sum_{j\in O(l)}u_l(j)\zh_j$.  Then
exactly one of the following holds:
\begin{enumerate}
\item[(Z$'$)] $|w_m(l)|<M_m/2$, i.e.\ $|v_l|<\mathrm{thr}_l/2$: a strict non-peak with gap
$>M_m/2$;
\item[(R)] $|w_m(l)|\ge M_m/2$ and $\sgn w_m(l)=\sgn Y_l$ \textup{(}or $Y_l=0$\textup{)}: then
$v_l=\sgn(w_m(l))(|Y_l|+\mathrm{own}_l)$, a peak with margin $\ge\mathrm{own}_l/2$ in units of
$|u(\zh)|$;
\item[(W)] $|w_m(l)|\ge M_m/2$ and $\sgn w_m(l)=-\sgn Y_l$: then $|Y_l|<\mathrm{own}_l$ and
$|v_l|=\mathrm{own}_l-|Y_l|$.
\end{enumerate}
\end{lemma}

\begin{proof}
If $|w(l)|\ge M/2$, the coefficient $|\Delta''_m|\lambda_lw(l)$ of $l$ dominates every later
term at every $j\in O(l)$ (the estimates of Lemma~\ref{lem:II-domfine} with $|\mathrm{coef}|
\ge|\Delta''_m|\lambda_lM/2$), so $\mathcal V_j\ne0$ and admissibility forces $z_j=\sgn(w(l))
\sgn(u_l(j))$ on $O(l)$ (as $\Delta''_m<0$ and $\zh_j=z_j$ there).  Hence $v_l=Y_l+
\sgn(w(l))\mathrm{own}_l$, and $\sgn v_l=\sgn w(l)$ by the clamp formula
(Lemma~\ref{lem:II-T}(b)); the two sign cases give (R) and (W).  If $|w(l)|<M/2$, the
clamp formula $|w(l)|=M|v_l|/\mathrm{thr}_l$ gives the bound, and $\gap=M-|w(l)|$.
Finally $\mathrm{own}_l\ge\delta^\circ_l\ge2\mathrm{thr}_l$ by (W3).
\end{proof}

\begin{lemma}[A scrambling criterion that tolerates wrong branches]\label{lem:II-SCband}
Call a carrier $l$ of a block of $A_-$ at a stage $>L+R_{\rm cl}(L)$ \emph{bad} if
$|u_l(\zh^\#)|$ lies in the band $B_l:=(\mathrm{thr}_l(1-\Phi_l^{1/2}/M),\ \mathrm{thr}_l+\Phi_l^{1/2})$.
If at $f^\#$ no such carrier is bad, then $A_-$ satisfies \textup{(SC)} at $f^\#$, along a
sequence common to all blocks of $A_-$.
\end{lemma}

\begin{proof}
Fix $\Upsilon\ge1$.  Coarse carriers and tuned absorbers do not contribute for small
$s$ (proof of Lemma~\ref{lem:II-SCi}).  A peak that is not bad has
$\mu_l=q^\#_0(|u_l(\zh^\#)|-\mathrm{thr}_l)\ge q^\#_0\Phi_l^{1/2}$, and a strict non-peak that is not
bad has $\gap_l=M(1-|u_l(\zh^\#)|/\mathrm{thr}_l)\ge\Phi_l^{1/2}$.  So a contribution to
$\mathrm{Scr}_m(\Upsilon s)$ requires $\Phi_l^{1/2}\le\Upsilon s/q^\#_0$, or
$\Phi_l\gap_l\le\Upsilon s$ with $\gap_l\ge M/2$ (then $\Phi_l\le2\Upsilon s/M$), or $\gap_l<M/2$ and
$\Phi_l^{3/2}\le\Upsilon s$.  In all cases $\Phi_l\le X(s):=\max\{(\Upsilon s/q^\#_0)^2,2\Upsilon s/M,(\Upsilon s)^{2/3}\}$.
Choose main stages $L_i\to\infty$ and $s_i$ with $X(s_i)=(\Phi_{L_i}\Phi_{L_i-1})^{1/2}$.  By (W2)
at $L_i$, $\Phi_{L_i}\le b(L_i-1,M(L_i-1))^2\le T_{\rm lo}^8$ of the preceding stage, which is
$\le\Phi_{L_i-1}^7$ for large $i$ (as $\log_2(1/T_{\rm lo})\ge n\ge\mathrm{Des}\ge1/\Phi_{L_i-1}$);
carriers before $L_i$ have $\Phi_l\ge2^{-N-l}c_{L_i-1}>X(s_i)$, and carriers after $L_i$ have
total weight $\le c_{L_i+1}\ll\Phi_{L_i}$.  So the contributions are $\le2\Phi_{L_i}$, and,
as $s_i\ge X(s_i)^{3/2}/\Upsilon$, $\Phi_{L_i}/s_i\le\Upsilon\Phi_{L_i}^{1/4}\Phi_{L_i-1}^{-3/4}\le\Upsilon\Phi_{L_i-1}\to
0$.
\end{proof}

Carriers in the state (R) are never bad ($\mathrm{own}_l\ge\delta^\circ_l$ while $\mathrm{thr}_l+
\Phi_l^{1/2}\le Cc_l^{1/2}$ and $c_l\le T_{\rm lo}(l-1)^3\ll(\delta^\circ_l)^8$), and (Z$'$) carriers are
not bad; only carriers in the wrong branch (W) with $\mathrm{own}_l-|Y_l|\in B_l$ can be bad.
In configuration (i) the recursion of Lemma~\ref{lem:II-recursion} never produces (W).

\begin{lemma}[The recursion in a mixed class]\label{lem:II-mixed}
Run the recursion of Lemma~\textup{\ref{lem:II-recursion}} in a mixed class with the
rule: a carrier of type \textup{(d)} of a block of $A_-$ takes $\varsigma:=\sgn Y$
\textup{(}self-aligned\textup{)}; a carrier $l$ of type \textup{(d)} of a block of $A_+$ takes
$\varsigma:=\sgn Y_l$ and anti-aligned own coordinates $z_j:=-\varsigma\sgn u_l(j)$ on $O(l)$.
Call such a positive carrier \emph{frustrated} if $|Y_l|<\mathrm{own}_l+\mathrm{thr}_l/2$.  If no
positive carrier is frustrated, the recursion defines $z$ on $J_{\rm fine}$ at which
$\mathcal V(\Delta\omega)$ is admissible on $J_{\rm fine}$ and every negative carrier is in the
state \textup{(R)}; hence $A_-$ satisfies \textup{(SC)} at $f^\#$.
\end{lemma}

\begin{proof}
For a non-frustrated positive carrier, $v_l=\varsigma(|Y_l|-\mathrm{own}_l)$ with $|v_l|\ge\mathrm{thr}_l/2$,
so $|w(l)|\ge M/2$ with sign $\varsigma$; its coefficient $-|\Delta''_m|\lambda_lw(l)$ has the sign
$-\varsigma$ and dominates on $O(l)$, so $\sgn\mathcal V_j=-\varsigma\sgn u_l(j)=z_j$.  Negative
carriers are treated as in Lemma~\ref{lem:II-recursion}; every $j\in J_{\rm fine}$ is
owned or carries $\mathcal V_j=0$.  Lemma~\ref{lem:II-SCband} gives (SC).
\end{proof}

A frustrated positive carrier has no consistent contact state: anti-aligned
own coordinates give the wrong sign of $v_l$, aligned ones the wrong sign of $\mathcal V$
on $O(l)$.  It must take a value of modulus $<\mathrm{thr}_l/2$, which needs free
coordinates or cancellations among later carriers, and downstream negative
carriers may then be forced into (W).

\subsubsection*{Master Theorem~IV$'$}

\begin{theorem}[Master Theorem~IV$'$]\label{thm:II-MTIV}
Let $N\ge1$, let the design be $D^{U1'}$, and let $f\in S_{p_N^*}$ have finite base
support.  Suppose that for infinitely many main stages $L$ some clean sub-window
$w$ of $L$ has at least $n(w)/\bar D(w)$ dyadic scales whose activity class $a$
is of one of the following kinds:
\begin{enumerate}
\item[($\alpha$)] $a$ has no active shift;
\item[($\beta$)] $a$ is one-signed and satisfies \textup{(KN$_{w,a}$)};
\item[($\gamma$)] $a$ is mixed, satisfies \textup{(KN$_{w,a}$)}, and the recursion of
Lemma~\textup{\ref{lem:II-mixed}} meets no frustrated positive carrier.
\end{enumerate}
Then $f\in\Rec$: $(f,g)\in\cl\NA((c_0,p_N),\ell_2^2)$ for every $g\in\Cset(f)$.
\end{theorem}

\begin{proof}
Fix $g$, $\rho<1$ and one such pair $(L,w)$ with $L\ge l_f$.  By pigeonhole over the
$D_{\rm cls}(w)\le\bar D(w)$ pairs (class, cube) of Proposition~\ref{prop:II-ray}(d)
there is a set $\mathcal S$ of at least $n(w)/\bar D(w)^2$ scales with one class
$a$ of the admitted kinds and one cube; $n(w)\ge L2^{L^3}Q(w)$ and $Q(w)\ge\bar D(w)^3$ give $|\mathcal S|\ge
L2^{L^3}\bar D(w)\ge L2^{L^3}(\mathrm{Des}/u)^C$, which exceeds $48\rho^2K/(c(1-\rho^2))$ with $K T_{\rm hi}(w)/|\mathcal S|
\to0$ along the main stages, for $K$ and $c^{-1}$ of the form $C_f(\mathrm{Des}/u)^C$.
Build $f^\#$ by Definition~\ref{def:II-compclass}: in configuration (i) with the
recursion of Lemma~\ref{lem:II-recursion}, in configuration (ii) and in mixed classes
with the fixed point of Lemma~\ref{lem:II-schauder} (in case ($\gamma$) the recursion of
Lemma~\ref{lem:II-mixed}), and for ($\alpha$) with the fixed owners only.

\emph{Data.}  For $t\in\mathcal S$ let $\omega^+$ be the clamped decomposition $\omega_+$ (as in
Proposition~\ref{prop:II-transplant}), let $\Delta\omega(t)$ have the components $\gamma_{\rm pr}(t)+
\gamma(\mathrm{incr}(t))$ on $\Omega_{\rm act}$, the shift $\Delta''^*$ (Lemma~\ref{lem:II-trueshift}), the
absorber coefficients of Proposition~\ref{prop:II-absorb}, and $0$ on inactive
$\Omega$-carriers; $\omega^-:=\omega^++\Delta\omega(t)$; $b^+:=\beta+\chi\mathcal V_{\rm pr}+\chi_{\rm f}\mathcal V_{\rm
fine}$, where $\chi$ is the split of the $+$ side of the decomposition with respect
to the projected coarse vector, $\chi_{\rm f}\in[0,1]$ is arbitrary on the fine contacts,
and $\beta$ is supported in $F^\#$ and chosen with $b^+(\xi^\#)=0$ (possible
since $\xi^\#\ne0$ on $F^\#$; this normalization is needed for tangency, compare
Proposition~\ref{prop:II-SAmates}(iii)); $b^-:=b^+-\mathcal V(\Delta\omega(t))$.  By
Lemma~\ref{lem:II-cert}(c) these are exact two-piece data at $f^\#$ as soon as
$\mathcal V(\Delta\omega(t))$ is $z^\#$-admissible off $F^\#$.  On $E_c(w)$ it equals $\mathcal L(\Delta''^*,
\gamma(t))$ exactly (Proposition~\ref{prop:II-absorb}), which is admissible
(Proposition~\ref{prop:II-ray}, Lemma~\ref{lem:II-trueshift}(b)); on $J_{\rm fine}$ it is
admissible by Lemmas~\ref{lem:II-domfine}, \ref{lem:II-recursion},
\ref{lem:II-schauder} or \ref{lem:II-mixed}; the active far parts of $\Omega$-carriers
are dominated by their own coefficients.  The $+$ side is that of
Proposition~\ref{prop:II-transplant} up to $Kt$ (projection $Kt$, inactive components
$JA_{i(t)}Kt$, fine part and absorber coefficients $\le C_fc_{L+1}$), so $p^*(g-g_t)\le
Kt$.  The $-$ side differs from that of Proposition~\ref{prop:II-transplant} by
$\mathcal V(\mathrm{incr}(t))|_{E_c}$ and $\Delta\omega(\mathrm{incr}(t))$, so, $\sqrt{\Gamma_w}$ being a seminorm,
$\sqrt{\Gamma^\#_w(b^-,\omega^-)}\le\sqrt{\Gamma_w(B_-,\Theta_-)}+Kt+C_fp\varepsilon$; with $\varepsilon=c_f
(u/\mathrm{Des})^C$ small, $\kappa_w\le1+\eta_0/2$.  The block parts are of kinds [1]--[3]:
increments satisfy the inward rows (X4), so near-threshold carriers are of kind
[3]; absorbers are of kind [2] (Proposition~\ref{prop:II-absorb}(c)).  All data of
$\mathcal S$ have the same shift $\Delta''^*$.

\emph{Statuses and (SC).}  Lemma~\ref{lem:II-coststat}.  In configuration (i), $A$
satisfies (SC) at $f^\#$ (Lemma~\ref{lem:II-SCi}); in case ($\gamma$), $A_-$ does
(Lemma~\ref{lem:II-mixed}); in configuration (ii) and in case ($\alpha$), $\Delta d_m\ge0$
for all $m$ and no (SC) is needed.

\emph{Recovery.}  $p^*(f^\#-f)\le\theta_wT_{\rm lo}(w)^2$ with $\theta_w\to0$
(Lemma~\ref{lem:II-coststat}).  Theorem~\ref{thm:II-Eshift} (part (a) or (b)), with
the banked, pulled and absorber supports allowed by Theorem~\ref{thm:II-E}(E-b)
(at absorber coordinates $t|b|\le$ the mass, so no flips), run on the subsets $\mathcal S$
(Lemma~\ref{lem:II-subset}) along the infinitely many sub-windows, gives
$(f,\rho g)\in\cl\NA$; let $\rho\to1$.
\end{proof}

\begin{corollary}[Intrinsic form]\label{cor:II-MTIVintr}
In the setting of Theorem~\textup{\ref{thm:II-MTIV}}, suppose that for infinitely many
main stages $L$ some clean sub-window $w$ of $L$ has $I_{\rm up}(w)=\emptyset$ or $I_{\rm
lo}(w)=\emptyset$, and that \textup{(KN$_{w,a}$)} holds for every activity class $a$ of
$w$.  Then $f\in\Rec$.
\end{corollary}

\begin{proof}
An upper source gives $\Delta'_{\rm dec,m}\ge-K_dt$ and a lower source $\Delta'_{\rm dec,m}\le K_dt$
(Lemma~\ref{lem:II-sources}, data convention).  A block of $I_{\rm up}(w)$ has a lower
source ($I_{\rm up}\cap I_{\rm lo}=\emptyset$, Lemma~\ref{lem:II-shiftpin}), so its active shift can
only be negative; a block of $I_{\rm lo}(w)$ can only have a positive active shift.
Zeroing and projection move components by at most $(1+C_f\mathrm{Des}\,JA_{i(t)})Kt$
and no active component below $A_{i(t)+1}Kt/2$ (Proposition~\ref{prop:II-ray}(b),
Lemma~\ref{lem:II-classes}), and components below $A_1Kt=K_*Kt>(K+K_d)t$ are
inactive; so the sign of an active shift component is that of the
decomposition data.  If $I_{\rm up}(w)$ or $I_{\rm lo}(w)$ is empty, every class of $w$ is one-signed, and
some class carries at least $n(w)/\bar D(w)$ scales (there are at most
$(J+1)3^J\le\bar D(w)$ classes).
\end{proof}

\begin{remark}[What is not proved]\label{rem:II-notMTIV}
(a) The statement of Theorem~\ref{thm:II-MTIV} \emph{without} (KN$_{w,a}$) in case ($\beta$)
is not proved; in particular it is \emph{not} proved that every row with finite base
support belongs to $\Rec$ for $N=1$, although for $N=1$ every class is one-signed.
The reason is Remark~\ref{rem:II-onelever}: near-threshold switching carriers of an
active block with negative shift are forced to be active, so (KN) is a genuine
restriction even for one block.  Without a lever, the residual is an active block
whose inactive $\Omega$-carriers all have $\varrho\le b$ and which has an active
near-threshold switching carrier $d$ such that, near the ratios of $f$, the exact zero
set of the tiny minors lies in $\{\varrho_d\ge1\}$.
(b) The hypotheses of Theorem~\ref{thm:II-MTIV} depend on $N$; since
[I,~Lemma~1.5] transfers density and not membership of individual rows, the theorem
gives no statement for Mart\'in's norm $p$.
\end{remark}

\subsection{Aligned rows, nested-tuning rows and further tools}\label{subsec:II-nested}

In this subsection $N=1$, $T=T_{\mathrm{final}}$, and we drop the block index.  We
show that (C$^*$) occurs at explicit rows, among them rows that are not
block-tame, and that the exact-data mates of such rows are recovered.

\begin{definition}[Aligned rows]\label{def:II-aligned}
A first row $f$ is \emph{aligned with exceptional carriers $l_-$ and $\Lambda$} if
$F\subseteq Z_0$ is finite, $z$ is given by the owner recursion
(Definition~\ref{def:II-owner}) with some signs $\epsilon_l$, and:
\begin{enumerate}
\item[(A1)] $l_-$ has its target supported in $F$, $\epsilon_{l_-}=+1$, $u_{l_-}(\zh)<0$,
and $l_-$ is a strict non-peak with $\varrho_{l_-}\in(0,\frac12]$;
\item[(A2)] every carrier $l\ne l_-$ is \emph{aligned}: $\sgn u_l(\zh)=\epsilon_l$;
\item[(A3)] the carriers of the set $\Lambda$ are strict non-peaks with $\varrho\ge\frac12$;
every carrier outside $\Lambda\cup\{l_-\}$ is a peak, and all but finitely many of them
have $\varrho\ge2$.
\end{enumerate}
\end{definition}

Every coordinate off $F$ is then a contact, all carriers are exactly swallowed,
$l_-$ is of anti type ($q_{l_-}<0$) and all other carriers are of swallowing type.
The rows $f_{\rm SA}$ of Theorem~\ref{thm:II-SA} (with $N=1$) are aligned with
$\Lambda=\emptyset$.  The dominance property of Theorem~\ref{thm:II-SA}(d) holds for
every aligned row, since Step~6 of that proof uses only the owner recursion,
(SF$_\iota$) and (W4).

\begin{proposition}[Coherent resonance at aligned rows]\label{prop:II-alignedres}
Let $f$ be aligned with exceptional carriers $l_-$ and $\Lambda$.  There is $l_1$ such that
at every clean sub-window $w$ of every level $l\ge l_1$: every coarse carrier is of
class G, the block has no upper source and has a lower source \textup{(L2)},
\[
 \rho^{\rm sh}(\kappa^{\rm sh}(w),f)=0\qquad\text{and}\qquad c_{\pi(w)}(f)=0 .
\]
In particular $f$ is a \textup{(C$^*$)} row.
\end{proposition}

\begin{proof}
All rooms vanish, so every coarse carrier is of class G and there is no class-R
carrier.  All peaks are of swallowing type, so (U1), (U2) fail; $l_-$ is a class-G
carrier with $q<0$ and $\varrho\in(0,\frac12]$, which for large $l$ is neither $\le b$
nor $\ge1-b$, so (U3) fails; the peaks with $\varrho\ge2\ge1+u$ give (L2).  Hence
$I_{\rm sh}=I_{\rm up}=\{1\}$, $\mathrm{src}(1)$ is ``free'' (no class-R source), the
set $N_0$ is empty ($\Lambda$ and the peaks have $\varrho\ge\frac12$, and $\varrho_{l_-}$ is fixed),
and $R_7=\emptyset$ (the only anti-type carrier is $l_-$).  Take $\delta:=1$,
$\tau_c:=\lambda_c$ for every coarse carrier $c\ne l_-$, and $\tau_{l_-}:=(1+\sum_{c\ne l_-}
q_c\lambda_c)/|q_{l_-}|>0$.  (S1) holds; (S2) holds since $\epsilon_c\mathrm{vs}_c=+1$ at peaks;
(S4): $\delta+\sum_cq_c\tau_c=1+\sum_{c\ne l_-}q_c\lambda_c-(1+\sum_{c\ne l_-}q_c\lambda_c)=0$; (S5),
(S6), (S7) are void.  (S3): $V(\tau)=\sum_{c\ne l_-}\lambda_c\epsilon_cu_c1_{F^c}+\tau_{l_-}u_{l_-}
1_{F^c}$.  A coordinate $j\in T(l)\setminus F$ is owned by a coarse carrier $o$.  If $o\ne
l_-$, the term $\lambda_o\epsilon_ou_o(j)$ has the sign $z_j$ and dominates the later coarse
terms (Step~6 of Theorem~\ref{thm:II-SA}, which bounds moduli), which never include
$l_-$ because $\supp u_{l_-}\subseteq F\cup S_{l_-}$.  If $o=l_-$, then $j\in S_{l_-}$, $z_j=
+1$, and $\tau_{l_-}v_{l_-}(j)$ dominates a fortiori, as $\tau_{l_-}\ge1/|q_{l_-}|\ge\lambda_{l_-}$.
So all contact rows hold and there are no free rows: $\mathrm{viol}(\delta,\tau)=0$ with
$\|\delta\|_1=1$.  For $c_{\pi(w)}$: $I_{\rm up}=\{1\}$ forces $\delta=1$; with $x:=\lambda$ on
$B_f(w)$ the vector $\Pi_1+\sum_{B_f}x\epsilon u1_{F^c}$ is the partial sum of
$\sum_c\lambda_c\epsilon_cu_c1_{F^c}$ over the coarse carriers, which is $z$-signed wherever it
is nonzero (dominance) and vanishes at coordinates touched by no coarse carrier; so
every summand $\phi_{z_j}(\cdot)$ vanishes.
\end{proof}

So (C$^*$) is a property of rows that may well belong to $\Rec$: the rows
$f_{\rm SA}$ are (C$^*$) rows and block-tame.  The following rows are not block-tame.

\begin{proposition}[Nested-tuning rows]\label{prop:II-NT}
Let $f$ be aligned with exceptional carriers $l_-$ and an infinite set $\Lambda=\{l_1<
l_2<\cdots\}$ whose relative gaps $g_i:=1-\varrho_{l_i}$ satisfy $0<g_i<\Phi_{l_{i+1}}/2$, and
assume that $f$ has no degenerate peak.  Then $Q$ is infinite, $f$ does not satisfy
\textup{(BT)}, the block fails \textup{(SC)}: $\mathrm{Scr}(\Upsilon s)\ge s$ for every $\Upsilon\ge1$
and every $0<s\le\Phi_{l_1}$, and $f$ is a \textup{(C$^*$)} row.
\end{proposition}

\begin{proof}
$Q\supseteq\{k(l_i)\}$ is infinite.  For $\Phi_{l_{i+1}}\le s\le\Phi_{l_i}$ the strict non-peak $l_i$
has $\gap=Mg_i<\Phi_{l_{i+1}}\le s$ and contributes $\min(\Phi_{l_i},\Upsilon s)\ge s$ to
$\mathrm{Scr}(\Upsilon s)$; so $\mathrm{Scr}(\Upsilon s)/s\not\to0$.  Proposition~\ref{prop:II-alignedres}.
\end{proof}

\begin{remark}[Existence of nested-tuning rows; sketch only]\label{rem:II-NTexist}
The existence of rows as in Proposition~\ref{prop:II-NT} is supported by the
following construction, which is \emph{a sketch and not a proof}.  Let $F=\{p,p',p''\}
\subseteq Z_0$, $\sigma=(+,+,+)$, and assume the design addition (Z0$+$): for $F_0:=F$ the
target family contains $c/q^*(c)$ for a countable dense set of directions $c\in\R^{F_0}$,
each used infinitely often (in $T_{\mathrm{final}}$ this is compatible with
allowedness, since $Z_0$ meets no signature set).  Parametrize $a$ by $b$ in the
open positive orthant of $S^2$; for $y$ supported in $F$, $y(\zh)=\ip y{\Phi(b)}$.
Non-special carriers follow a \emph{hysteretic} owner rule (reference signs kept
unless $\epsilon_lA_l\le-(B_l+\delta_lH_l)/2$), so they are aligned with $n_l|u_l(\zh)|\ge(B_l+
\delta_lH_l)/2$; their hysteresis margins decrease super-exponentially.  Specials have
targets $c_i/q^*(c_i)$ from (Z0$+$), so $B=0$ and $A=\ip{c_i}{\Phi(b)}/q^*(c_i)$.  One
holds at stage $i$ a closed ball $B_i$, a level $L_i$, specials $l_1<\dots<l_i$, and
numbers $\varepsilon_j,\kappa_i>0$ such that: (J1) all assignments of carriers $\le L_i$ are
constant on $B_i$; (J2) for $j<i$ the gap of $l_j$ relative to the threshold of the
carriers $\le L_i$ lies in $[\varepsilon_j,\Phi_{l_{j+1}}/2-\varepsilon_j]$ on $B_i$, with margin against
the threshold perturbation of later carriers; (J3) $l_-$ and the specials keep their
signs and the sign conditions with margin; (J4) the gap function $G_i$ of $l_i$ takes
values $\ge\kappa_i$ and $\le-\kappa_i$ in the interior of $B_i$.  The step $i\to i+1$ picks a
zero of $G_i$, a direction $c_{i+1}$ nearly orthogonal to $\Phi$ there with
independent tangential gradient, moves along the curve $\{G_i=\Phi_{l_{i+1}}/4\}$ to a
zero of the gap of a late carrier $l_{i+1}$ with that target (intermediate value
theorem; theta is Lipschitz, but not differentiable on $\{\mathsf v=\bar\vartheta\}$, so
the argument must be topological), takes a ball of radius $\sim\Phi_{l_i}\Phi_{l_{i+1}}$
around that point, and only then chooses $L_{i+1}$.  The limit point of the nested
balls gives the row.  (An earlier version of this induction confined the new gap
to an interval bounded away from $0$ and is inconsistent: the region required at
the next stage is empty.)
\end{remark}

\begin{proposition}[Explicit coherent-shift mates of aligned rows]\label{prop:II-NTM}
Let $f$ be aligned with exceptional carriers $l_-$ and $\Lambda$, $\psi:=R^*w$, $k_-:=k(l_-)$.
\begin{enumerate}
\item[(i)] $\psi1_{F^c}$ is $z$-signed and nonzero at every coordinate off $F\cup S_{l_-}$,
and $z\psi<0$ on $S_{l_-}$.
\item[(ii)] Let $\Omega\subseteq Q$ be finite with $k_-\in\Omega$, and assume either $\Omega=\{k_-\}$, or
that the targets of all carriers of $\Omega$ are supported in $F$.  Let $\omega^+$ and
$\Delta\omega$ be supported in $\Omega$ with $\epsilon_k\Delta\omega(k)\ge0$ for $k\in\Omega\setminus\{k_-\}$,
$\Delta:=d(\Delta\omega)<0$ and $\Delta\omega(k_-)>\Delta w(k_-)$; let $\chi:F^c\to[0,1]$ and $\beta$ be
supported in $F$ with $(\beta+\chi v1_{F^c})(\zh)=0$, where $v:=R^*\Delta\omega-\Delta\psi$.  Then
$\omega^-:=\omega^++\Delta\omega$, $b^+:=\chi v1_{F^c}+\beta$, $b^-:=b^+-v$ are two-piece data
\textup{(}$\Delta d=\Delta<0$\textup{)} for $g:=b^++R^*(\omega^+-d(\omega^+)w)$, and $v\ne0$ off $F$.
\item[(iii)] $cg\in\Cset(f)$ for all small $c>0$.
\item[(iv)] For $s\in S_l$ with $l\notin\Omega\cup\{l_-\}$ and $s\notin\Sigma(\omega)$, $g(s)/\psi(s)=-d(\omega^+)
+\chi_s|\Delta|$; the profile is not constant whenever $\chi$ is not.
\end{enumerate}
\end{proposition}

\begin{proof}
(i) At a coordinate owned by $o\ne l_-$, $|w(o)|\ge M/2$ ($M$ at peaks, $M\varrho_o\ge M/2$
on $\Lambda$) and $M>\frac13$, while the later terms of $\psi$ have moduli $\le2^{-5}$ times
the own term computed with $|w|=1$ (Step~6 of Theorem~\ref{thm:II-SA}); so $\sgn\psi(j)=
\sgn(w(o)u_o(j))=\epsilon_o\sgn u_o(j)=z_j$, by (A2).  On $S_{l_-}$, $w(k_-)<0$ and $z=+1$.
(ii) Off $F$, $v=R^*\Delta\omega-\Delta\psi$.  Under either assumption the carriers of $\Omega$
meet $F^c$ only in their own signature sets.  On $S_k$, $k\in\Omega\setminus\{k_-\}$:
$v(j)=\lambda_k(\Delta\omega(k)-\Delta w(k))v_k(j)-\Delta\sum_{l>k}\lambda_lw(l)u_l(j)$, and
$\epsilon_k(\Delta\omega(k)-\Delta w(k))\ge|\Delta||w(k)|\ge|\Delta|M/2$ (as $\sgn w(k)=\epsilon_k$), while the later
terms are $\le|\Delta|2^{-5}\lambda_kv_k(j)$; so $z_jv(j)\ge(M/2-2^{-5})|\Delta|\lambda_kv_k(j)>0$.  On
$S_{l_-}$: $\Delta=d(\Delta\omega)=\sum_{k\in\Omega}\Phi_k^2w(k)\Delta\omega(k)/C$, the terms with $k\ne k_-$
are $\ge0$, and $w(k_-)<0$, so $\Delta\omega(k_-)\ge|\Delta|C/(\Phi_{k_-}^2|w(k_-)|)\gg|\Delta|$; the own
coefficient $\lambda_{l_-}(\Delta\omega(k_-)-\Delta w(k_-))$ is positive and dominates the later terms.  Elsewhere off $F$, $v=-\Delta\psi=|\Delta|\psi$ is $z$-signed and nonzero by (i).
Hence $b^+$ is $z$-signed, $b^-=(\chi-1)v1_{F^c}+(\beta-v1_F)$ is $(-z)$-signed off $F$, there
are no free coordinates, both pairs represent $g$ since $b^+-b^-=v=R^*(\omega^--\omega^+)-
\Delta\psi$, and $g(\xi)=q_0b^+(\zh)=0$.  (iii) As in Proposition~\ref{prop:II-SAmates}(iv).
(iv) At such $s$, $R^*\Delta\omega(s)=0$, so $g(s)=\chi_sv(s)-d(\omega^+)\psi(s)$ and $v(s)=-\Delta\psi(s)$.
\end{proof}

Without one of the two assumptions in (ii) the conclusion can fail: a target
coordinate off $F$ of a carrier of $\Omega$ that is owned by another carrier receives
the term $\lambda_k\Delta\omega(k)y_k(j)/n_k$, which is not controlled by $\Delta$.

\begin{corollary}[Exact-data mates of aligned rows are recovered]\label{cor:II-NTR}
Let $f$ and the data be as in Proposition~\textup{\ref{prop:II-NTM}(ii)}, with
$\kappa_w\le1$ and $\gamma_k:=\lambda_k(\Delta\omega(k)-\Delta w(k))\ne0$ for $k\in\Omega$.  Then the
functional $g$ satisfies $(f,g)\in\cl\NA((c_0,p_1),\ell_2^2)$ whenever $g\in\Cset(f)$.
\end{corollary}

\begin{proof}
We verify (H1), (H2) of Theorem~\ref{thm:II-negexact}.  Every carrier has $|w(k)|\ge
\min(M\varrho_{l_-},M/2)>0$, so no carrier outside $\Omega$ is neutral ($\gamma_k=-\lambda_k\Delta w(k)\ne
0$), and a carrier outside $\Omega$ is negligible iff $|w(k)|<2\iota_k$, which happens for
finitely many $k$ since $\iota_k\to0$; carriers of $\Omega$ are not neutral by assumption.
By Proposition~\ref{prop:II-NTM}(ii), $D=b^+-b^-=v$ is nonzero at every coordinate off
$F$, so (H2) holds.  Theorem~\ref{thm:II-negexact} applies.
\end{proof}

So the dangerous mates of the nested-tuning rows (oscillating profiles on the
signature sets of robust class-G peaks, forced shift $\Delta<0$, no (SC) at $f$) are
recovered, through block-tame companions, not through windows.  Whether every
mate of such a row is recovered is not proved.

\smallskip\noindent\emph{Banks and pinning.}  The following two tools are used in the
discussion of the remaining steps.

\begin{lemma}[Quadratic cost of banks]\label{lem:II-QB2}
Let $f$ have finite base support, let $E$ be a finite set of contacts, $m_j>0$
$(j\in E)$, and let $f^b$ have the data $a^b:=(a+\sum_{j\in E}m_jz_je_j^*)/q^*(\cdot)$ and
$z$ \textup{(}so $F^b=F\cup E$\textup{)}.
\begin{enumerate}
\item[(a)] For $x\in\R$ and $\mathsf m>0$, $2(x-\mathsf m)_+\le(x_+)^2/(2\mathsf m)$.  Consequently the
summand of the base excess of [I,~Lemma~7.1] at $j\in E$ for a perturbation $tB$,
namely $2(-\sgn(a^b_j)tB_j-|a^b_j|)_+$, is at most $(t^2/2)((z_jB_j)_-)^2/|a^b_j|$ for $t>0$;
for a side-$+$ pair it vanishes unless $z_jb^+_j<0$, for a side-$-$ pair unless
$z_jb^-_j>0$, and for the average piece $b^\theta=\frac12(b^++b^-)$ it is at most
$(t^2/2)(b^\theta_j)^2/|a^b_j|$.
\item[(b)] Suppose that at every $j\in E$, $\max(|b^+_j|,|b^-_j|,|b^\theta_j|)\le C_Ee_j$ with
$e_j:=(z_jb^+_j)_-+(z_jb^-_j)_+$, and put $\varepsilon:=\sum_{j\in E}e_j$.  For $\beta_0>0$ let
$m_j:=(C_E^2q_0\varepsilon/\beta_0)e_j$ and $\mathsf m:=\sum_jm_j=C_E^2q_0\varepsilon^2/\beta_0$.  Then
$q_0\sum_{j\in E}(C_Ee_j)^2/|a^b_j|\le\beta_0(1+\mathsf m)$, and $p^*(f^b-f)\le C_f\mathsf m\log(e/\mathsf
m)$.
\end{enumerate}
\end{lemma}

\begin{proof}
(a) $x^2/(2\mathsf m)-2(x-\mathsf m)=(x-2\mathsf m)^2/(2\mathsf m)\ge0$, and the left side vanishes for
$x\le\mathsf m$.  At $j\in E$, $\sgn a^b_j=z_j$, so the flip occurs only when $z_jtB_j<0$.
(b) $|a^b_j|=m_j/q^*(a+\sum m_jz_je^*_j)\ge m_j/(1+\mathsf m)$, so $q_0\sum C_E^2e_j^2/|a^b_j|\le
(1+\mathsf m)q_0C_E^2\sum e_j^2/m_j=(1+\mathsf m)\beta_0$.  The cost: Lemma~\ref{lem:II-cost} with
$\|a^b-a\|_1\le2\mathsf m$, and $\zh$ moves at first order only on $E$ (diagonal base).
\end{proof}

A violation of $\ell_1$-mass $\varepsilon$ on finitely many contacts thus costs a
companion move of order $\varepsilon^2$; the hypothesis $|b^\pm_j|\le C_Ee_j$ holds when the
violations are of fine origin and the split at $j$ is clamped.

\begin{lemma}[Pinned strict non-peaks bound the shift from above]\label{lem:II-Rpin}
Let $f$ be any first row, $m$ a block, $k\in Q_m$ with $w_m(k)\ne0$, and let $(B_\pm,
\Theta_\pm)$ be a two-sided decomposition at scale $t\le\min(t_\eta,1)$ as in
[I,~Lemma~8.8].  Then, in the decomposition convention,
\[
 \Delta d_m|w_m(k)|\le3\gap_m(k)/t+|\Delta\Theta_m(k)| .
\]
\end{lemma}

\begin{proof}
With $\varsigma:=\sgn w_m(k)$, [I,~Lemma~8.8] gives $\varsigma\omega_+(k)\le(1-d_+t)\gap/t$ and
$\varsigma\omega_-(k)\ge-(1+d_-t)\gap/t$; with $|d_\pm t|\le\frac12$, $\varsigma(\omega_+-\omega_-)(k)\le3\gap/t$.
Since $\Theta_\pm=\omega_\pm-d_\pm w$, $\omega_+-\omega_-=\Delta\Theta+\Delta d\,w$, so $\varsigma(\omega_+-\omega_-)(k)=
\varsigma\Delta\Theta(k)+\Delta d|w(k)|$.
\end{proof}

In particular a class-R strict non-peak with $\gap\le t^2$ (for which $|\Delta\theta|\le K_gt$)
pins the upward shift at scale $t$: it acts as an additional upper source.

\begin{proposition}[Exact coupling for one shifted block]\label{prop:II-S1}
Let $w$ be a clean sub-window of level $l\ge l_f$, $I_{\rm sh}(w)=\{m\}$, and let $k$ be a
class-G strict non-peak of block $m$ with $q_k<0$ whose column $\epsilon_ku_k1_{F^c}$ is a
\emph{free ray} of the combinatorial zero-cost rows \textup{(}at every $j\in
T(l)\setminus F$ with $u_k(j)\ne0$, $j$ is a contact whose row is satisfied for
every value $x_k\ge0$ of the coefficient, the other coefficients being fixed,
e.g.\ because the owner of $j$ dominates; in particular this holds if $u_k$
vanishes on $T(l)\setminus F$\textup{)}.  Let $(\delta,x)$ satisfy the combinatorial rows of
Corollary~\textup{\ref{cor:II-C11}} exactly, and let $r:=\delta_m+\sum_lq_lx_l$ be the
coupling residual.  If $r\ge0$, or if $r<0$ and $x_k\ge|r|/|q_k|$, then $(\delta,x+(r/|q_k|)
e_k)$ satisfies the combinatorial rows and the coupling exactly, and
$|r|/|q_k|\le C\mathrm{Des}(l)|r|/u(w)$.  If $r<0$, $x_k<|r|/|q_k|$, and all other carriers in the
support of $x$ have $q\ge0$, then $\delta_m<0$.
\end{proposition}

\begin{proof}
Changing $x_k$ within $[0,\infty)$ keeps the combinatorial rows (free ray), and
changes the coupling by $q_k$ times the change; in the case $r<0$ the new
coefficient $x_k-|r|/|q_k|$ is $\ge0$ by hypothesis.  At a clean sub-window $\varrho_k$ is robust (it is
not nearly neutral, since $q_k\ne0$ is used with $q'_k=q_k$), so $|q_k|=\Phi_k|w(k)|/(mC)\ge
\Phi_kMu/(mC)\ge u/(C\mathrm{Des})$.  In the last case $\delta_m=r-\sum_{l\ne k}q_lx_l-q_kx_k\le r+
|q_k|x_k<r+|r|=0$.
\end{proof}

For one shifted block with such a carrier, the coupling row (the $d$-identity) can
therefore be satisfied exactly inside the exact combinatorial cone with a constant
of the form $\mathrm{Des}/u$; in the last case of the proposition the sign of the shift
is the one excluded, up to $K_dt$, by a lower source (Lemma~\ref{lem:II-sources}),
so that the shift is pinned.  This is the situation of the aligned rows, where $l_-$
has its target in $F$.

\subsection{What remains open for finite base support}\label{subsec:II-open}

For $T_{\mathrm{final}}$ (and $D^{U1'}$) and first rows with finite base support,
the results of this section reduce Lemma~Z ([I,~Problem~8.21]) to the following
question.

\begin{problem}\label{prob:II-Cstar}
Let $N\ge1$.  Is every \textup{(C$^*$)} row $f\in S_{p_N^*}$ \textup{(}Definition~\ref{def:II-Cstar}\textup{)}
in $\Rec$?
\end{problem}

By Corollary~\ref{cor:II-Cstarnec} an affirmative answer gives $\Rec\supseteq\{f:F\text{
finite}\}$ for $p_N$.  (C$^*$) rows exist (Proposition~\ref{prop:II-alignedres}): some of
them are block-tame and therefore in $\Rec$ (Corollary~\ref{cor:II-SArec}); the
nested-tuning rows of Proposition~\ref{prop:II-NT}, whose existence is only sketched
(Remark~\ref{rem:II-NTexist}), are not block-tame, and only their mates with exact data
are known to be recovered (Corollary~\ref{cor:II-NTR}).  What is proved about
(C$^*$) leaves the following precise gaps.  Every item below is open; where a route is
indicated, it is a sketch or a heuristic and is labelled as such.

\begin{remark}[Status coherence without (KN)]\label{rem:II-openKN}
Master Theorem~IV$'$ needs the lever hypothesis (KN$_{w,a}$) in every class with active
shift that is used.  Without it the residual is: an active block whose inactive
$\Omega$-carriers all have $\varrho\le b(w)$, with an active near-threshold switching carrier
$d$, such that the exact zero set $Z$ of the tiny minors of
Definition~\ref{def:II-Gamma}, in the ratio variables and near the ratios of $f$, lies in
$\{\varrho_d\ge1\}$ (Proposition~\ref{prop:II-Rkt}); if $d$ ends as a peak, its excess
$|\omega_+(d)|+|\omega_-(d)|$ must be $O(Kt/\lambda_d)$, which nothing proved implies.
\emph{Sketch:} near-threshold carriers whose inward excess is inactive can be removed
from $\Omega$ with the continuous coefficient $-\Delta'\lambda_d\min(\varrho_d,1)\mathrm{vs}_d$, which needs no
status, so only carriers with active excess need a lever.  \emph{Heuristic:} if the
design were generic in the sense that the map $(\varrho_d)_d$ restricted to every stratum
of every zero set $Z_S$ is transversal to $(1,\dots,1)$, then $Z$ would enter
$\{\varrho_d<1\}$ near $\{\varrho_d=1\}$ with design-controlled depth; this needs a construction
of such a design (only the parameters of stage $l$ are free at stage $l$), quantitative
transversality near singular strata, and a quantitative treatment of the case $\dim
Z_S<|\{d\}|$.  The absorbers and the one-ray normalization are not affected by this
gap; it concerns only the joint exactness of the ratios and the protection of the
thresholds.
\end{remark}

\begin{remark}[Mixed classes (C$_{\rm mix}$)]\label{rem:II-openmix}
In a mixed class (possible only for $N\ge2$) exact completions exist
(Lemma~\ref{lem:II-schauder}); what is missing is a completion at which no carrier of a
block of $A_-$ is bad in the sense of Lemma~\ref{lem:II-SCband}.  Carriers in the state
(R) are never bad, configuration (i) never produces the state (W), and a bad carrier
in the state (W) requires the coincidence $\mathrm{own}_l-|Y_l|\in B_l$, an interval of
length about $\Phi_l^{1/2}$.  Finite lexicographic models show that mixing can force the
state (W) at every exact completion (\emph{numerical evidence only}).  Conceivable
routes, none carried out: a selection among the Schauder fixed points by continuation
from vanishing positive shift; a genericity argument in finitely many coarse
parameters; or tolerating the sign violations of frustrated positive carriers (which
have $\ell_1$-mass $\le C_fc_{L+1}$, of fine origin) by
Lemma~\ref{lem:II-VT}, which would reduce (C$_{\rm mix}$) to item (S2) below
(\emph{sketch} of a reduction).
\end{remark}

\begin{remark}[A second route: exact coupling and uniform composition]\label{rem:II-openS}
A different route to Problem~\ref{prob:II-Cstar} works at the assembled companion of
Section~\ref{sec:II-exact} and tolerates fine-origin violations instead of
absorbing them.  Its proved ingredients are: the exact combinatorial resonance
(Corollary~\ref{cor:II-C11}(b)); the characterization of shifted exact data
(Proposition~\ref{prop:II-C3}, Corollary~\ref{cor:II-A}); the quadratic cost of banks
(Lemma~\ref{lem:II-QB2}); the violation tolerance of a single engineered stage
(Lemma~\ref{lem:II-VT}); and exact coupling for one shifted block with a free-ray
carrier (Proposition~\ref{prop:II-S1}).  The remaining steps are:
\begin{enumerate}
\item[(S1)] \emph{Exact coupling in general}: with two or more shifted blocks, or
without a free-ray carrier with $q<0$, satisfy the coupling rows $\Delta d_m=d^\#_m(\omega^-)-
d^\#_m(\omega^+)$ exactly inside the exact combinatorial cone with constants of the form
$C_f(\mathrm{Des}/u)^C$.  (Proposition~\ref{prop:II-Rkt} and
Remark~\ref{rem:II-onelever} explain why a single push per block does not
obviously suffice; whether two independent levers per block are necessary is
not decided --- \emph{heuristic}.)
\item[(RT$^*$c)] \emph{Coarse exactness} at the companion modulo the coupling rows and
the contributions of carriers beyond the level (class-R and unused class-G strict
non-peaks of shifted blocks must be neutralized as switching carriers, with sizes
controlled by Lemma~\ref{lem:II-Rpin}) --- \emph{sketch}.
\item[(S2)] \emph{Uniform composition}: (S2a$'$) engineered approximants that are
\emph{$d$-consistent}, i.e.\ at which the normalized values of the finitely many coarse
switching carriers are re-tuned exactly to their values at the row (so that $d'=d$ on
every switching vector), with per-piece bounds $T_0\ge c_\flat t$ uniform in the scale
--- \emph{sketch}; (S2b) the averaging of Theorem~\ref{thm:II-E} carried out at one
engineered approximant for all pieces of a window --- \emph{sketch}; (S2c) a
comparison of the violation mass with the late threshold of the companion ---
\emph{heuristic}.
\end{enumerate}
The $d$-consistency in (S2a$'$) appears to be needed for this route
(\emph{heuristic}): without it, the window
masses of an engineered approximant change the coarse block data at first order, and
for window data with switching of size $1/t$ this produces a mismatch between the
$\theta$-piece and the $\pm$ pieces of order $s_1/t^2$ ($s_1$ the stage scale); an upper
bound of this order follows from the proof of [I,~Theorem~7.18], and finite models show
the scaling with equality up to a constant (\emph{numerical evidence}).  Rebalancing such
a mismatch would require transfer peaks of inefficiency $\lesssim t^2$ at weights that
super-fast ladders do not provide (\emph{heuristic}).  This does not affect
Theorems~\ref{thm:II-E} and \ref{thm:II-Eshift}, which apply [I,~Theorem~7.18] at a fixed
row to fixed exact data.
\end{remark}

\begin{remark}[Further open points]\label{rem:II-openmisc}
(a) The general mates of the nested-tuning rows: a route through exact window data at
the row with a shift column, transplant to block-tame companions with inward moves
of near-threshold special carriers, and (Z0$+$), is plausible but is only a
\emph{sketch}; Lemma~Z at these rows is not proved.
(b) The existence of the nested-tuning rows themselves is a \emph{sketch}
(Remark~\ref{rem:II-NTexist}).
(c) Rows that are given by the owner rule above some level and carry an
\emph{exactly} degenerate swallowing-type peak are not known to exist (the threshold has
a jump part from fine sign changes which may skip the degenerate value); this concerns
only the illustration of configurations, not the residual.
(d) The results of this section are proved for $p_N$ with $N$ fixed; their hypotheses
are not independent of $N$, so they give no row-wise statement for Mart\'in's norm $p$.
(e) Rows with infinite base support are not treated here.
\end{remark}

Lemma~Z ([I,~Problem~8.21]) and the density of $\NA((c_0,p_N),\ell_2^2)$ and of
$\NA((c_0,p),\ell_2^2)$ remain open for every admissible operator, including
$T_{\mathrm{final}}$; nothing proved in this section points to a counterexample.

%%%%%%%%%%%%%%%%%%%%%%%%%%%%%%%%%%%%%%%%%%%%%%%%%%%%%%%%%%%%%%%%%%%%%%%%

\section{Infinite base support}\label{sec:II-infinite}

This section treats first rows $f$ whose base support $F=\supp a$ is
\emph{infinite}.  In Part~I this case was open except for a sketch
[I,~Remark~8.46]; the missing step was the transfer of peaks to a first row
with $a\notin c_{00}$.  We prove that step
(Theorem~\ref{thm:II-transferinf}), so that [I,~Remark~8.46] becomes a
theorem (Theorem~\ref{thm:II-Binf}), and then develop the tools that are
specific to infinite base support: at a scale $t$ a support coordinate $j$
behaves like a contact of sign $\sgn a_j$ with a free anti-sign allowance
(\emph{cushion}) of size $|a_j|/t$ on each side, and the cost of exceeding
the cushions is governed by the \emph{flip profile} of the data.  The main
results are the following.
\begin{itemize}
\item Engineered recovery at arbitrary $F$ needs sparsity of the flips only
for the one-sided parts of the data, because the support may be
\emph{raised} where the average data need it
(Theorem~\ref{thm:II-raised}).
\item Finitely many exactly swallowed carriers whose swallowing by the
support is \emph{cushion sparse} are harmless for every signature-ladder
operator (Theorem~\ref{thm:II-R1}, under (W$^\star$), (H2), (H3-inf));
box-dominated support ($F=\N$) is harmless for \emph{every} admissible
operator (Theorem~\ref{thm:II-N}, under (SC$_I$) and (W$_M$)).
\item For the designed operator $T_{\mathrm{final}}$ the raise-transfer
lemma (Lemma~\ref{lem:II-RT}) removes every profile condition: support
swallowing of any profile by finitely many bad carriers is harmless
(Theorems~\ref{thm:II-RS}, \ref{thm:II-RSstar}, \ref{thm:II-NDprime}, under
(W$^\star$), (H2), (H3-inf)), and
infinitely many bad carriers are harmless under an explicit rate condition
(Theorem~\ref{thm:II-RSinf}).
\end{itemize}
The residual at infinite base support is described in
Subsection~\ref{subsec:II-residualinf}.

\smallskip\noindent\emph{Standing conventions.}  $N\ge1$,
$I=\{1,\dots,N\}$, $p=p_N$, and $f\in S_{p^*}$ has forced data
$(\xi,q_0,a,w,F,z,\zh,e,\nu)$ [I,~Proposition~1.9] with $F=\supp a$
\emph{arbitrary} (finite or infinite).  We write $s_j:=\sgn a_j=z_j$ for
$j\in F$ and $\alpha(M):=\sum_{j\in F,\,j>M}|a_j|$.  Subsections
\ref{subsec:II-transferinf}--\ref{subsec:II-raised} and
\ref{subsec:II-N} hold for every admissible operator $T$.  In
Subsections~\ref{subsec:II-room} and~\ref{subsec:II-R1}, $T$ is a
\emph{signature-ladder operator}: either the operator of
[I,~Definition~8.1] with its windows $\mathcal W(l)$, or
$T=T_{\mathrm{final}}$ with a fixed choice of one sub-window $w_l$ of each
level $l$, used as the $l$-th window (Lemma~\ref{lem:II-GW}(b)); in both
cases (T-a)--(T-d), (P1), (P2), (P3) of [I,~Theorem~8.2] and allowedness
(a), (b) of [I,~Definition~8.1] hold, and these are the only properties
of the operator used there.  In Subsections~\ref{subsec:II-RS}
and~\ref{subsec:II-residualinf}, $T=T_{\mathrm{final}}$ and
$U=U_{\mathrm{final}}$ (Section~\ref{sec:II-Tfinal}).  For a first row
$f'$ with prescribed data we use the notation $f[a',z']$ of
Subsection~\ref{subsec:II-comp}; forced data of other first rows carry
primes, superscripts or indices.

For a two-sided decomposition $(B_\pm,\Theta_\pm)$ of a mate at scale $t$
[I,~Lemma~8.5] we use $\Delta B=B_+-B_-$, $\Delta\theta_{k,m}$ and
[I,~(15)], and the flip quantities of [I,~Lemma~8.28],
\[
 \mathfrak f^+_j:=\Big(-s_jB_+(j)-\frac{|a_j|}t\Big)_+,\qquad
 \mathfrak f^-_j:=\Big(s_jB_-(j)-\frac{|a_j|}t\Big)_+\qquad(j\in F),
\]
which satisfy $\sum_{j\in F}\mathfrak f^\pm_j\le t/(4q_0)$ for every $t>0$
and every $F$ [I,~Lemma~8.28].  The \emph{box profile} of $T$ is
$\mathrm{Bx}(j):=\sum_{m\in I}\sum_k\lambda_{k,m}|u_{k,m}(j)|$; since
$\|u_{k,m}\|_1\le1$ and $\sum_k\lambda_{k,m}\le m2^{-m}$, $\|\mathrm{Bx}\|_1
<2$.

\subsection{Transfer peaks at a base vector outside $c_{00}$}
\label{subsec:II-transferinf}

The proof of [I,~Theorem~5.8] uses $a\in c_{00}$ only through the canonical
truncations, which keep the base vector.  For $a\notin c_{00}$ we truncate
the base vector as well.

\begin{definition}[Truncated canonical approximants]\label{def:II-trunc}
For $N'\in\N$ with $a1_{[1,N']}\ne0$ put
\begin{gather*}
 c^a_{N'}:=q^*(a1_{[1,N']}),\qquad a_{N'}:=a1_{[1,N']}/c^a_{N'},\\
 \nu_{N'}:=\|U^*a_{N'}\|_H,\qquad e_{N'}:=U^*a_{N'}/\nu_{N'},\qquad
 z^{N'}:=z1_{[1,N']},\\ \tilde x_{N'}:=z^{N'}+Ue_{N'},\qquad
 x'_{N'}:=\tilde x_{N'}/p(\tilde x_{N'}),\qquad f'_{N'}:=\nabla p(x'_{N'}),
\end{gather*}
and $c_{N'}:=q(x'_{N'})$.
\end{definition}

\begin{lemma}[Truncated canonical approximants]\label{lem:II-trunc}
Let $T$ be admissible and $f\in S_{p^*}$.
\begin{enumerate}
\item[(a)] $f'_{N'}\in\NA((c_0,p),\R)\cap S_{p^*}$, with forced data
$a_{N'}$, $e_{N'}$, $z^{N'}$ and $\zh'_{N'}=\tilde x_{N'}$.
\item[(b)] As $N'\to\infty$: $a_{N'}\to a$ in $\ell_1$, $e_{N'}\to e$,
$f'_{N'}\to f$ in norm, all forced data of $f'_{N'}$ converge to those of
$f$ as in [I,~Proposition~1.11], $x'_{N'}\to\xi$ weak$^*$ and
$c_{N'}\to q_0$.
\item[(c)] Let $P_{N'}^\perp$ be the orthogonal projection onto
$e_{N'}^\perp$ and $\Psi_{N'}(h):=\|e_{N'}+h\|-1-\ip{e_{N'}}h$.  If
$B\in\ell_1$ is supported in $F\cap[1,N']$ and $\sgn(a_{N'}(j)+tB_j)=
\sgn a_{N'}(j)$ for $j\in F\cap[1,N']$, then
\[
 q^*(a_{N'}+tB)=1+tB(\tilde x_{N'})+\nu_{N'}\Psi_{N'}(tU^*B/\nu_{N'}),
\]
and $\Psi_{N'}(h)\le\|P^\perp_{N'}h\|^2/(2(1-\|h\|))$ for $\|h\|\le\frac12$.
\end{enumerate}
\end{lemma}

\begin{proof}
(a) $a_{N'}\in c_{00}$, $q^*(a_{N'})=1$, $z^{N'}\in c_{00}$,
$\|z^{N'}\|_\infty\le1$ and $z^{N'}=z=\sgn a=\sgn a_{N'}$ on
$\supp a_{N'}=F\cap[1,N']$; apply [I,~Proposition~1.10(c)] and
[I,~Proposition~1.9] at $x'_{N'}$.  (b) $c^a_{N'}\to q^*(a)=1$, $U^*$ is
bounded, and $z^{N'}\to z$ coordinatewise; [I,~Proposition~1.12]
((ii)$\Rightarrow$(i)) and [I,~Proposition~1.11] give the convergence of
$f'_{N'}$ and of its forced data, and $c_{N'}=q(x'_{N'})=1/p(\tilde
x_{N'})\to1/p^{**}(\zh)=q_0$.  (c) No sign change on $F\cap[1,N']$ gives
$\|a_{N'}+tB\|_1=\|a_{N'}\|_1+tB(z^{N'})$, and
$\|U^*a_{N'}+tU^*B\|=\nu_{N'}+t\ip{U^*B}{e_{N'}}+\nu_{N'}\Psi_{N'}(tU^*B/
\nu_{N'})$; add, using $\|a_{N'}\|_1+\nu_{N'}=1$ and $B(z^{N'})+
\ip{U^*B}{e_{N'}}=B(\tilde x_{N'})$.  The bound on $\Psi_{N'}$ is
[I,~Lemma~3.3].
\end{proof}

\begin{theorem}[Transfer peaks at an arbitrary base vector]
\label{thm:II-transferinf}
Let $T$ be admissible, $f\in S_{p^*}$ \textup{(}$F$ arbitrary\textup{)},
and let $g=b+\sum_mR_m^*(\omega_m-d_m(\omega_m)w_m)$ with $b\in\ell_1$,
$\supp b\subseteq F$, $b(\xi)=0$,
\[
 \|b/a\|_\infty:=\sup_{j\in F}|b_j|/|a_j|<\infty,
\]
and $\omega_m\in c_{00}$, $\supp\omega_m\subseteq Q_m$, such that
$g\in\Cset(f)$ and $\Gamma_w(g)\le1$.  Then for every $\rho\in(0,1)$ there
are $g'_{N'}\in X^*$ with $g'_{N'}\to g$ and $\rho g'_{N'}\in
\Cset(f'_{N'})$ for all large $N'$, where $f'_{N'}$ are the truncated
canonical approximants.  In particular $(f,g)\in\cl\NA((c_0,p),\ell_2^2)$.
\end{theorem}

\begin{proof}
We follow the proof of [I,~Theorem~5.8], whose only use of $a\in c_{00}$
is that its canonical truncations keep the base vector $a$; we use the
truncated canonical approximants instead, and indicate the changes.
Steps~2 and~3 there (transfer peaks at $\xi$, built from
$\tau_{*,m}=(\hat\pi_m(\xi)/q_0)a-\hat\pi_m$ and the density of the tails of
$(u_{k,m})_k$, and the limiting levels) do not involve the approximants and
hold verbatim for $a\in\ell_1$.  The block estimates of Steps~1, 4, 5 and
Step~6 use only [I,~Lemmas~1.7 and~3.4] at $x'_{N'}$, the convergence of the
forced data (Lemma~\ref{lem:II-trunc}(b)), [I,~Lemma~1.6] and $g\in
\Cset(f)$; they are unchanged.

\emph{Step 1 (base).}  Put $\kappa_{N'}:=(b1_{[1,N']})(\tilde x_{N'})$ and
$b'_{N'}:=b1_{[1,N']}-\kappa_{N'}a_{N'}$.  Since $\tilde x_{N'}\to\zh$
weak$^*$ boundedly (Lemma~\ref{lem:II-trunc}(b)) and $z^{N'}=z$ on
$[1,N']$, $\kappa_{N'}\to b(\zh)=b(\xi)/q_0=0$, so $b'_{N'}\to b$ in
$\ell_1$; $\supp b'_{N'}\subseteq F\cap[1,N']$ and $b'_{N'}(\tilde
x_{N'})=\kappa_{N'}-\kappa_{N'}a_{N'}(\tilde x_{N'})=0$, because
$a_{N'}(\tilde x_{N'})=1$.  For large $N'$, $c^a_{N'}\le2$ and
$|\kappa_{N'}|\le\frac14$; then for $|t|\le r_b:=\min(1,1/(4\|b/a\|_\infty))$
and $j\in F\cap[1,N']$,
\[
 |tb'_{N'}(j)|\le|t||b_j|+|t\kappa_{N'}||a_{N'}(j)|\le\tfrac14|a_j|+
 \tfrac14|a_{N'}(j)|\le\tfrac34|a_{N'}(j)|,
\]
since $|a_j|=c^a_{N'}|a_{N'}(j)|\le2|a_{N'}(j)|$.  So no sign changes
occur, and Lemma~\ref{lem:II-trunc}(c) gives
\begin{gather*}
 q^*(a_{N'}+tb'_{N'})\le1+\frac{t^2}2h_{N'}\Big(1+\frac{2|t|\|U^*b'_{N'}\|}
 {\nu_{N'}}\Big),\\ h_{N'}:=\frac{\|P^\perp_{N'}U^*b'_{N'}\|^2}{\nu_{N'}}
 =\frac{\|P^\perp_{N'}U^*(b1_{[1,N']})\|^2}{\nu_{N'}},
\end{gather*}
and $h_{N'}\to h(b)=H_b$ because $P^\perp_{N'}\to P^\perp$ in norm and
$\nu_{N'}\to\nu$.  This is the base bound of Step~1 of [I,~Theorem~5.8]
with $a$ replaced by $a_{N'}$: $q^*(a_{N'}+tb'_{N'})\le1+\frac{t^2}2H_b
\theta(t)$ with $\theta(t)\to1$ as $t\to0$, uniformly in large $N'$.

\emph{Step 4.}  Replace $\beta_{m,N}$ by $\beta_{m,N'}:=((R_m^*y_{m,N'})
(x'_{N'})/c_{N'})a_{N'}$ and $e_{m,N'}:=R_m^*y_{m,N'}-\beta_{m,N'}$; then
$e_{m,N'}\to e_{\infty,m}$, because $R_m^*y_{m,N'}\to R_m^*y_m$ in $\ell_1$,
$x'_{N'}\to\xi$ weak$^*$ boundedly, $c_{N'}\to q_0$ and $a_{N'}\to a$ in
$\ell_1$.  The bookkeeping identity [I,~(6)] is an identity at $x'_{N'}$
and holds unchanged.

\emph{Step 5.}  With $A(t):=a_{N'}+t\rho b'_{N'}+\sum_m\varepsilon_m
R_m^*y_{m,N'}=(1+E)a_{N'}+t\rho b'_{N'}+\sum_m\varepsilon_me_{m,N'}$ one has,
by homogeneity and Step~1,
\[
 q^*(A(t))\le(1+E)\,q^*\Big(a_{N'}+\frac{t\rho}{1+E}b'_{N'}\Big)+\rho^2t^2
 \sum_m|\tau_m|q^*(e_{m,N'}),
\]
which is the base bound of Step~5 of [I,~Theorem~5.8] with $a$ replaced by
$a_{N'}$.  Put $g'_{N'}:=b'_{N'}+\sum_mR_m^*(\omega_m-d'_{m,N'}w'_{m,N'})$
with $d'_{m,N'}$ as there.  Then $g'_{N'}\to g$ and $g'_{N'}(x'_{N'})=0$
($b'_{N'}(\tilde x_{N'})=0$, and [I,~Lemma~3.2] at $x'_{N'}$, the support of
$\omega_m$ consisting of strict non-peaks of $w'_{m,N'}$ for large $N'$).
Step~6 uses only $p^*(f'_{N'}-f)\to0$ and $p^*(g'_{N'}-g)\to0$.  Hence
$\rho g'_{N'}\in\Cset(f'_{N'})$ for large $N'$.  As $f'_{N'}$ attains its
norm, $(f'_{N'},\rho g'_{N'})$ attains its norm [I,~Proposition~1.13(d)]
and tends to $(f,\rho g)$; let $\rho\to1$.
\end{proof}

\begin{corollary}[Windowed averaging at arbitrary base support]
\label{cor:II-windowedinf}
Let $T$ be admissible and $f\in S_{p^*}$ \textup{(}$F$ arbitrary\textup{)}.
Then [I,~Theorem~8.17] holds for $f$ provided every certificate $c_t$ in its
hypotheses has a base part $b_t$ with $\|b_t/a\|_\infty<\infty$ \textup{(}for
instance $\supp b_t$ finite\textup{)}.
\end{corollary}

\begin{proof}
The proof of [I,~Theorem~8.17] up to its last sentence uses only the local
validity (i), the closeness (ii), $g\in\Cset(f)$, the convexity of $p^*$ and
of $\Gamma_w$, and not the finiteness of $F$.  It produces an averaged
certificate $c$ whose base part is a finite average of the $b_t$, hence has
$\|b/a\|_\infty<\infty$, with $\rho g_c\in\Cset(f)$, $\Gamma_w(\rho c)\le1$
and $\rho g_c\to\rho g$.  Theorem~\ref{thm:II-transferinf} applied to the
mate $\rho g_c$ gives $(f,\rho g_c)\in\cl\NA$, and $\cl\NA$ is closed.
\end{proof}

\begin{remark}\label{rem:II-transferinf}
(a) The bounded ratio $\|b/a\|_\infty<\infty$ is used only to obtain a
no-flip radius $r_b$ independent of $N'$.  (b) The approximants
$f'_{N'}$ do not depend on $g$: all mates with bounded-ratio base parts and
$\Gamma_w\le1$ are recovered simultaneously along one sequence.  (c) Pure
base mates ($g\in\ell_1$ with $q^*(a+tg)\le s(t)$ for all $t$) are
recovered along every truncation of $a$ as well, by a direct computation
with the base identity; we do not need this.
\end{remark}

\subsection{Cushions and flip profiles}\label{subsec:II-cushion}

\begin{definition}[Flip profiles]\label{def:II-flip}
For $\beta\in\ell_1$ with $\beta\ge0$ put
\begin{gather*}
 m_\beta(x):=\sum\{\beta_j:\ j\in F,\ |a_j|<x\beta_j\}\quad(x>0),\\
 \mathrm{Fl}_\beta(r):=\sum_{j\in F}2(r\beta_j-|a_j|)_+\quad(r\ge0),
\end{gather*}
and $\kappa_\beta(x):=m_\beta(x)/x$.  The profile of $\beta$ is
\emph{sparse} if $\kappa_\beta(x)\to0$ as $x\to0^+$, \emph{super-critical}
if $\kappa_\beta(x)\to\infty$, and \emph{critical} if $0<\liminf_{x\to0^+}
\kappa_\beta(x)\le\limsup_{x\to0^+}\kappa_\beta(x)<\infty$.
\end{definition}

For finite $F$ every profile is sparse ($m_\beta(x)=0$ for
$x<\min_F|a_j|/\|\beta\|_\infty$).

\begin{lemma}[Flip calculus]\label{lem:II-flipcalc}
Let $\beta,\beta_1,\beta_2\in\ell_1$ be nonnegative.
\begin{enumerate}
\item[(a)] $\mathrm{Fl}_\beta(r)=2\int_0^rm_\beta(s)\,ds$; it is convex and
nondecreasing in $r$, and $\mathrm{Fl}_\beta(0)=0$.
\item[(b)] $\beta\mapsto\mathrm{Fl}_\beta(r)$ is convex and coordinatewise
nondecreasing; $\beta\le\beta'$ implies $m_\beta\le m_{\beta'}$;
$m_{c\beta}(x)=c\,m_\beta(cx)$ for $c>0$; and
$m_{\beta_1+\beta_2}(x)\le2m_{\beta_1}(2x)+2m_{\beta_2}(2x)$.
\item[(c)] For $b\in\ell_1$ and $r>0$,
$\sum_{j\in F}\big(|a_j+rb_j|-|a_j|-s_jrb_j\big)=\mathrm{Fl}_{(sb)_-}(r)$,
where $(sb)_-:=((s_jb_j)_-)_{j\in F}$; for $r<0$ the left side equals
$\mathrm{Fl}_{(sb)_+}(|r|)$.
\item[(d)] If $m_\beta(x)\le\kappa x$ for $0<x\le x_0$, then
$\mathrm{Fl}_\beta(r)\le\kappa r^2$ for $0\le r\le x_0$; the same with
$\ge$.  In particular $\mathrm{Fl}_\beta(r)=o(r^2)$ iff the profile of $\beta$
is sparse.
\end{enumerate}
\end{lemma}

\begin{proof}
(a) For each $j$, $r\mapsto2(r\beta_j-|a_j|)_+$ is convex, vanishes at $0$
and has derivative $2\beta_j1[r\beta_j>|a_j|]$ off one point; sum by
monotone convergence.  (b) Each summand is convex and nondecreasing in
$\beta_j\ge0$.  The scaling identity is a substitution.  For the sum, split
$\{|a_j|<x(\beta_1+\beta_2)_j\}$ according to which of $\beta_{1,j},
\beta_{2,j}$ is larger: on the first part $(\beta_1+\beta_2)_j\le
2\beta_{1,j}$ and $|a_j|<2x\beta_{1,j}$.  (c) For $x\ne0$ and $y\in\R$,
$|x+y|-|x|-\sgn(x)y=2(-\sgn(x)y-|x|)_+$; apply it with $x=a_j$, $y=rb_j$.
(d) Integrate (a).  If $\mathrm{Fl}_\beta(r)\le\varepsilon r^2$ for $r\le r_0$,
then, $m_\beta$ being nondecreasing, $\varepsilon(2r)^2\ge\mathrm{Fl}_\beta(2r)\ge
2\int_r^{2r}m_\beta\ge2rm_\beta(r)$, i.e.\ $\kappa_\beta(r)\le2\varepsilon$.
\end{proof}

\begin{lemma}[Critical profiles oscillate]\label{lem:II-oscillate}
Let $S=\{s_1<s_2<\cdots\}\subseteq F$ be infinite, $\beta:=v1_S$ with
$v(s)=c2^{-s}$ $(c>0)$, and suppose that $\mathfrak a(s):=|a_s|/v(s)$ is
nonincreasing along $S$ with $\mathfrak a(s_i)\to0$.  Then
$\kappa_\beta(y^+)\ge2\kappa_\beta(y)$ at every point $y=\mathfrak a(s_i)$ with
$\mathfrak a(s_{i+1})<\mathfrak a(s_i)$, where $\kappa_\beta(y^+):=\lim_{x\downarrow
y}\kappa_\beta(x)$.  Consequently $\limsup_{x\to0^+}\kappa_\beta(x)\ge2
\liminf_{x\to0^+}\kappa_\beta(x)$.
\end{lemma}

\begin{proof}
Put $R_i:=\sum_{i'\ge i}v(s_{i'})$.  As $s_{i'+1}\ge s_{i'}+1$,
$R_{i+1}\le2v(s_{i+1})\le v(s_i)$, so $R_i=v(s_i)+R_{i+1}\ge2R_{i+1}$.  At
$y=\mathfrak a(s_i)$ with $\mathfrak a(s_{i+1})<y$, monotonicity gives $m_\beta(y)=
R_{i+1}>0$ and $m_\beta(y^+)\ge R_i$.  Since $\mathfrak a(s_i)\to0$, such points
$y$ accumulate at $0$, and $\kappa_\beta(y^+)\ge2\kappa_\beta(y)\ge2\inf_{x\le
y}\kappa_\beta(x)$ there.
\end{proof}

For signature vectors $v_l=\delta_lh_l/n_l$ of a signature-ladder operator
whose signature sets have bounded gaps (as for $T_{\mathrm{final}}$, $S_l=
\{2^l(2i+1):i\ge1\}$), and $|a_s|\approx v_l(s)^{1+\beta_0}$ on $S_l\cap F$
($\beta_0>0$) with $S_l\setminus F$ finite, the profile of $v_l1_{S_l\cap F}$
is sparse for $\beta_0<1$ (this part holds for every $S_l$), critical for
$\beta_0=1$ and super-critical for $\beta_0>1$; in the critical case the
local constant $\kappa$ oscillates by a factor at least $2$.  For lacunary
$S_l\cap F$ the cases $\beta_0\ge1$ need not be critical or super-critical
(for $\beta_0=1$ and lacunary $S_l\cap F$ one has $\liminf_ym(y)/y=0$).

\begin{lemma}[Cushions]\label{lem:II-cushion}
Let $T$ be admissible, $g\in\Cset(f)$, $t>0$, and let $(B_\pm,\Theta_\pm)$
be a two-sided decomposition of $g$ at scale $t$.  For $j\in F$,
\begin{gather*}
 (s_jB_+(j))_-\le\frac{|a_j|}t+\mathfrak f^+_j,\qquad (s_jB_-(j))_+\le
 \frac{|a_j|}t+\mathfrak f^-_j,\\ (s_j\Delta B(j))_-\le\frac{2|a_j|}t+
 \mathfrak f^+_j+\mathfrak f^-_j,
\end{gather*}
and $\sum_{j\in F}(\mathfrak f^+_j+\mathfrak f^-_j)\le t/(2q_0)$.
\end{lemma}

\begin{proof}
The first two inequalities are the definitions of $\mathfrak f^\pm_j$; the
third follows from $(x-y)_-\le x_-+y_+$; the sum is [I,~Lemma~8.28].
\end{proof}

Thus at scale $t$ a support coordinate is a contact of sign $s_j$ whose
anti-sign use is free up to the \emph{cushion} $|a_j|/t$ on each side.  The
cushion grows as $t\to0$; far support coordinates ($|a_j|\ll t^2$) behave
like contacts, near ones are two-sided.

We next extend the one-sided transfer expansion of
Lemma~\ref{lem:II-U} to arbitrary base supports, allowing flips on the
support whose profile is sparse uniformly along the sequence.  For a first
row $f_j$ we write $s^{(j)}_i:=\sgn a_j(i)$ on $F_j:=\supp a_j$, and use the
coordinate kinds [1], [2], [3] of Subsection~\ref{subsec:II-U}.

\begin{lemma}[One-sided expansion at arbitrary base support]
\label{lem:II-Uinf}
Let $T$ be admissible, $f\in S_{p^*}$, and let $f_j\to f$ in $S_{p^*}$
\textup{(}the supports $F$, $F_j$ arbitrary; the constant sequence $f_j=f$
is allowed\textup{)}.  Let $C_+\ge1$, $C_*\ge0$ and $U_*\in\ell_1$,
$U_*\ge0$, and suppose that
\[
 \bar m(y):=\sup_j\sum\{U_*(i):\ i\in F_j,\ |a_j(i)|<yU_*(i)\}=o(y)\qquad
 (y\to0^+).
\]
For every $\varepsilon_{\rm tr}>0$ and $A_0\ge1$ there are $t_1>0$ and
$c_0\in(0,\frac18]$, depending only on $f$, $\varepsilon_{\rm tr}$, $A_0$,
$C_+$, $C_*$ and $\bar m$, and $j_0\in\N$, depending in addition on the
sequence $(f_j)$, with the following property.  Let
$A_2\ge1$, $\gamma_B\in(0,1]$ and $c_\flat:=c_0\gamma_B/A_2$.  Let $j\ge j_0$,
$t\in(0,t_1]$, and let $(b,\omega)$, $b\in\ell_1$, $\omega_m\in c_{00}$,
satisfy:
\begin{enumerate}
\item[(OS1)] $b(\xi_j)=0$, $t\|b\|_1\le A_0$ and $\Gamma^{(j)}_w(b,\omega)\le2$;
\item[(OS2$'$)] $b(i)=0$ for $i\notin F_j\cup K_j$, $z_j(i)b(i)\ge0$ for
$i\in K_j$, and $(s^{(j)}_ib(i))_-\le C_+|a_j(i)|/t+C_*U_*(i)$ for $i\in F_j$;
\item[(OS3)] every $k\in\supp\omega_m$ is of kind \textup{[1]},
\textup{[2]} \textup{(}with $A_2,\gamma_B$\textup{)} or \textup{[3]}
\textup{(}side $+$\textup{)} at $f_j$.
\end{enumerate}
Then $G:=b+\sum_mR_m^*(\omega_m-d^{(j)}_m(\omega_m)w_{j,m})$ satisfies
\[
 p^*(f_j+rG)\le1+\frac{r^2}2\big(\Gamma^{(j)}_w(b,\omega)+\varepsilon_{\rm tr}
 \big)\qquad(0<r\le c_\flat t).
\]
The same holds for side-$-$ pairs \textup{(}in \textup{(OS2$'$)} and
\textup{[3]} the signs reversed, $(s^{(j)}_ib(i))_+$ in place of
$(s^{(j)}_ib(i))_-$\textup{)} and $-c_\flat t\le r<0$.
\end{lemma}

\begin{proof}
We run the proof of Lemma~\ref{lem:II-U} and record the changes; side $+$.
In Step~0 the number $a_{\min}$ and the condition on
$\min_{i\in F}|a_j(i)|$ are dropped (they served only to exclude flips); all
other limits and transfer data are as there.  Transfer data
[I,~Lemmas~7.4, 7.5] exist at every first row; their construction does not
involve $F$.  In that proof the sequence enters only through the index
$j_0$ beyond which the forced data of $f_j$ are within the factors of Step~0
of those of $f$ and the transfer data persist; every condition on $c_\flat$
and $t_1$ involves only quantities of $f$, $\varepsilon_{\rm tr}$ and $A_0$
(and, below, $C_+$, $C_*$, $\bar m$).

\emph{Step 1 (base).}  By [I,~Lemma~7.1(b)] at $f_j$, $G^{(j)}_b(a_j+rb)=
\Exc^{(j)}(rb)+\nu_j\Psi(rU^*b/\nu_j)$, and by (OS2$'$) the summands of
$\Exc^{(j)}(rb)$ vanish off $F_j$.  On $F_j$, Lemma~\ref{lem:II-flipcalc}(c)
gives the summand $2(r\beta_i-|a_j(i)|)_+$ with $\beta_i:=(s^{(j)}_ib(i))_-$.
Assume $c_\flat C_+\le\frac14$.  Then $r\beta_i\le|a_j(i)|/4+rC_*U_*(i)$; a
positive summand forces $rC_*U_*(i)>\frac34|a_j(i)|$, so $|a_j(i)|<2C_*rU_*(i)$,
and then $2(r\beta_i-|a_j(i)|)_+\le2r\beta_i\le\frac12|a_j(i)|+2rC_*U_*(i)\le
3rC_*U_*(i)$.  Hence
\[
 \Exc^{(j)}(rb)\le\mathrm{Fl}_*(r):=3C_*r\,\bar m(2C_*r),
\]
and $\mathrm{Fl}_*(r)/r^2\to0$ as $r\to0$, independently of $j$, $t$ and
$(b,\omega)$.  Fix $r_1>0$ with $\mathrm{Fl}_*(r)\le\min(1,\varepsilon_{\rm tr}
q_0/20)\,r^2$ for $0<r\le r_1$, and require $t_1\le r_1$ (then $r\le
c_\flat t\le r_1$).  The Hilbert part is estimated as in
Lemma~\ref{lem:II-U}, and
\[
 G^{(j)}_b(a_j+rb)\le\hat G_b:=\frac{r^2}2h^{(j)}(b)(1+K_Ac_\flat)+
 \mathrm{Fl}_*(r).
\]
\emph{Step 2 (blocks)} is unchanged.  \emph{Step 3 (rebalancing).}  Since
$\mathrm{Fl}_*(r)\le r^2$, the bound $|\varepsilon_m|\le K_3r^2$ holds after
enlarging $K_3$ by an additive constant depending only on $f$; the error
term $|q^*(A)-1|+q^*(A-a_j)\le2q^*(rb)\le2(1+\|U\|)c_\flat A_0$ is
unchanged.  \emph{Step 4 (constants).}  The only new contribution to
$\hat\Gamma$ is $q^{(j)}_0\mathrm{Fl}_*(r)\le\frac{r^2}2\cdot
\frac{\varepsilon_{\rm tr}}{10}$ (as $q^{(j)}_0\le1$); absorb it by replacing
$\varepsilon_{\rm tr}/5$ by $\varepsilon_{\rm tr}/10$ in the other conditions
of that step.  The new condition $c_\flat\le1/(4C_+)$ is of the form
$c_\flat\le\kappa_1$, and the new condition $t_1\le r_1$ does not involve
$A_2$ or $\gamma_B$.  The side-$-$ case is the same with $-r,-b,-\omega$.
\end{proof}

With $C_*=0$ (or $U_*=0$) the hypothesis on $\bar m$ is void and no flips
occur.  For the constant sequence $f_j=f$ and $U_*$ with $m_{U_*}(y)=o(y)$
the lemma is the one-sided expansion with sparse flips used in
Subsection~\ref{subsec:II-R1}.

\begin{lemma}[First-order pin and bounded free switching at any $F$]
\label{lem:II-boundedfree}
Let $T$ be admissible, $g\in\Cset(f)$, $\eta\le\eta_*$ with
$\eta_\Gamma(\eta)\le1$, $0<t\le\min(t_\eta,1)$, and let $(B_\pm,\Theta_\pm)$
be a two-sided decomposition of $g$ at scale $t$.
\begin{enumerate}
\item[(a)] $\big|\sum_{m}\sum_k\Delta\theta_{k,m}u_{k,m}(\zh)\big|=
|\Delta B(\zh)|\le t/q_0$.
\item[(b)] For every finite set $\mathcal U$ of pairs $(k,m)$, $m\in I$,
there is $C_{\mathcal U}<\infty$, depending only on $f$ and $\mathcal U$,
with
\[
 \max_{(k,m)\in\mathcal U}|\Delta\theta_{k,m}|\le C_{\mathcal U}\Big(1+
 \Big\|\sum_{(k,m)\notin\mathcal U}\Delta\theta_{k,m}u_{k,m}\Big\|_1\Big).
\]
\end{enumerate}
\end{lemma}

\begin{proof}
(a) is [I,~(15)] together with $|q_0B_\pm(\zh)|\le t/2$
[I,~Lemma~8.7(a)].  (b) Let $P^\perp$ be the projection onto $e^\perp$ and
$\Phi(c):=(P^\perp U^*\sum_{\mathcal U}c_{k,m}u_{k,m},\ (\sum_{\mathcal U}
c_{k,m}u_{k,m})(\zh))$ for $c\in\R^{\mathcal U}$.  If $\Phi(c)=0$, then
$U^*\sum_{\mathcal U}c_{k,m}u_{k,m}=\mu e$ for some $\mu$, hence
$\sum_{\mathcal U}c_{k,m}u_{k,m}=(\mu/\nu)a$ ($U^*$ is injective), hence
$0=(\mu/\nu)a(\zh)=\mu/\nu$; so $\sum_{\mathcal U}c_{k,m}u_{k,m}=0$ and $c=0$
by (T-b).  Thus $\|c\|_\infty\le C_{\mathcal U}\|\Phi(c)\|$.  Apply this to
$c=\Delta\theta|_{\mathcal U}$: by [I,~(15)], $\sum_{\mathcal U}\Delta\theta
_{k,m}u_{k,m}=-\Delta B-R$ with $R:=\sum_{(k,m)\notin\mathcal U}\Delta\theta
_{k,m}u_{k,m}$; $\|P^\perp U^*\Delta B\|\le2(2\nu/q_0)^{1/2}$ by
[I,~Lemma~8.7(d)], $|\Delta B(\zh)|\le t/q_0\le1/q_0$ by (a),
$\|P^\perp U^*R\|\le\|U\|\|R\|_1$ and $|R(\zh)|\le(1+\|U\|)\|R\|_1$.
\end{proof}

Part (b) is [I,~Lemma~8.29] without the hypothesis that $F$ be finite:
the possibility $a\in Y$ is harmless because the $a$-direction is pinned at
first order by (a).

\begin{lemma}[Un-switching pairs]\label{lem:II-unswitch}
Let $T$ be admissible and let $l=(k,m)$ be a carrier with $w_m(k)=0$ (a
\emph{$d$-neutral} carrier; then $k\in Q_m$).  For $c\in\R$ the pair
$P_l(c):=\big(cu_{k,m},\ -(c/\lambda_{k,m})e_k\text{ in block }m\big)$
represents the zero functional, $d_m(e_k)=0$, its base part vanishes at
$\xi$, and
\[
 \Gamma_w(P_l(c))\le c^2C_{\rm un}(l)^2,\qquad C_{\rm un}(l)^2:=
 \frac{q_0\|U\|^2}\nu+\frac{\sigma_m}{m^2C_m}.
\]
Consequently, if $(b,\omega)$ represents $g$, so does $(b,\omega)+P_l(c)$,
and $\sqrt{\Gamma_w((b,\omega)+P_l(c))}\le\sqrt{\Gamma_w(b,\omega)}+|c|
C_{\rm un}(l)$.
\end{lemma}

\begin{proof}
$M_m>0$ and $w_m(k)=0$ give $k\in Q_m$.  $R_m^*e_k=\lambda_{k,m}u_{k,m}$ and
$d_m(e_k)=\Phi_m(k)^2w_m(k)/C_m=0$, so the pair represents $cu_{k,m}-cu_{k,m}
=0$.  By [I,~Lemma~1.7] off $P_m$, $\lambda_{k,m}q_0u_{k,m}(\zh)=\zeta_m(k)=
\sigma_m\Phi_m(k)^2w_m(k)/C_m=0$, so $u_{k,m}(\xi)=0$.  Further
$h(u_{k,m})\le\|U\|^2\|u_{k,m}\|_1^2/\nu\le\|U\|^2/\nu$ and
$H_m(e_k/\lambda_{k,m})\le\Phi_m(k)^2/(\lambda_{k,m}^2C_m)=1/(m^2C_m)$.
$\sqrt{\Gamma_w}$ is a seminorm on pairs.
\end{proof}

\subsection{Engineered recovery with raised support}\label{subsec:II-raised}

\begin{definition}[Two-piece data at arbitrary base support]
\label{def:II-twopieceinf}
Two-piece data for $g\in X^*$ at $f$ are defined by [I,~Definition~7.7]
read without the hypothesis that $F$ be finite: a pair $(b,\omega)$,
$b\in\ell_1$, $\omega_m\in c_{00}$, $\supp\omega_m\subseteq Q_m$, is
side-$+$ \textup{(}side-$-$\textup{)} admissible if it represents $g$,
$b_j=0$ for $j\notin F\cup K$ and $z_jb_j\ge0$ \textup{(}$\le0$\textup{)}
for $j\in K$; no condition is imposed on $F$.  We put $b^\theta:=
\frac12(b^++b^-)$, $v:=b^+-b^-$, and $\Delta d_m$, $\kappa_w$ as in
[I,~Definition~7.7].  The data satisfy \emph{one-sided cushion sparsity}
\textup{(CS-side)} if the profiles of $(sb^+)_-$ and $(sb^-)_+$
\textup{(}restricted to $F$\textup{)} are sparse:
\[
 m_{(sb^+)_-}(x)=o(x)\quad\text{and}\quad m_{(sb^-)_+}(x)=o(x)\qquad
 (x\to0^+).
\]
\end{definition}

No condition is imposed on $b^\theta$; for finite $F$, (CS-side) is void.

\begin{theorem}[Engineered recovery with raised support]\label{thm:II-raised}
Let $T$ be admissible, $f\in S_{p^*}$ \textup{(}$F$ arbitrary\textup{)}, and
let $g\in\Cset(f)$ carry two-piece data $(b^\pm,\omega^\pm)$ satisfying
\textup{(CS-side)}.  Let $\rho\in(0,1)$ with $\rho^2\kappa_w<1$, and assume
that $I_-:=\{m:\Delta d_m<0\}$ satisfies the scrambling condition
\textup{(SC)} of [I,~Definition~7.16] \textup{(}no assumption if
$I_-=\emptyset$\textup{)}.  Then there are norm-attaining first rows
$f'_i\to f$ with finite base supports and $g'_i\in\Cset(f'_i)$ with
$g'_i\to\rho g$.  In particular $(f,\rho g)\in\cl\NA((c_0,p),\ell_2^2)$, and
$(f,g)\in\cl\NA((c_0,p),\ell_2^2)$ if $\kappa_w\le1$.
\end{theorem}

\begin{proof}
We modify the proof of [I,~Theorem~7.18] and its approximants
[I,~Definition~7.12].  Put $\delta:=(1-\rho^2\kappa_w)/2$ and
$m(x):=m_{(sb^+)_-}(x)+m_{(sb^-)_+}(x)$.

\emph{Approximants.}  For a window $N_w$, a scale $s_1\in(0,1]$ and a
cut-off $N''>N_w$ put $\mathsf m_j:=4\rho s_1|b^\theta_j|$ for $j\in
K\cap[1,N_w]$, $\mathsf r_j:=(8\rho s_1|b^\theta_j|-|a_j|)_+$ for $j\in
F\cap[1,N_w]$, and
\[
 a'':=a1_{[1,N_w]}+\sum_{j\in K\cap[1,N_w]}\mathsf m_jz_je^*_j+\sum_{j\in
 F\cap[1,N_w]}\mathsf r_js_je^*_j ,
\]
so that $|a''_j|=\max(|a_j|,8\rho s_1|b^\theta_j|)$ with sign $s_j$ on
$F\cap[1,N_w]$ (the support is \emph{raised} where the average data need
it).  All other definitions ($a'=a''/q^*(a'')$, $e'$, $z'_j=z_j$ for $j\le
N''$ and $z'_j=0$ beyond, $\hat x'$, $x'$, $f'$) are those of
[I,~Definition~7.12].  The window $N_w$ is chosen \emph{after} $s_1$, with
\[
 \textup{(N0)}\qquad\alpha(N_w)\le s_1^2,\qquad\rho\sum_{j\in F,\,j>N_w}
 (s_jv_j)_-\le\delta s_1/32,
\]
and then $N''>N_w$ with [I,~(11), (12)] and $\rho V_{>N''}\le\delta s_1/32$,
where now $V_{>N''}:=\sum_{j\in F\cup K,\,j>N''}|v_j|$.  This is possible
because $v\in\ell_1$.  Then $\supp a'=(F\cap[1,N_w])\cup\{j\in K\cap[1,N_w]:
b^\theta_j\ne0\}$ is finite and $z'=\sgn a'$ there, so $f'$ attains its
norm and has finite base support ([I,~Lemma~7.13(a)]); support coordinates
in $(N_w,N'']$ become contacts of $f'$ with signs $s_j$, those beyond $N''$
free coordinates.  Moreover $\|a''-a\|_1\le s_1^2+12\rho s_1\|b^\theta\|_1$, so
(E1) of [I,~Lemma~7.13] holds with $4$ replaced by $12$, an additional $s_1^2$,
and $K_e$ replaced by $K'_e:=2\|U\|(12\rho\|b^\theta\|_1+1)/\nu$; (E2) is
unchanged, since $\hat x'-\zh=-z1_{(N'',\infty)}+U(e'-e)$.  The remaining
parts (E3)--(E5) of [I,~Lemma~7.13] and [I,~Lemmas~7.14, 7.15, 7.17] use the
approximants only through (E1), (E2) and the forced data of $f'$; they hold
with the adjusted constants.  At late stages $q^*(a'')\le2$, hence
$|a'_j|\ge|a''_j|/2\ge|a_j|/2$ on $F\cap[1,N_w]$.

\emph{Step 0.}  The condition $(T_1)$ ``$\rho T_0(\|b^+\|_\infty+
\|b^-\|_\infty)\le a_{\min}/8$'' is replaced by
\[
 2\rho|\tau|\,m(4\rho|\tau|)\le\min\big(\tfrac{\delta}{16},1\big)\tau^2
 \qquad(0<|\tau|\le T_0),
\]
which holds for small $T_0$ by (CS-side) and
Lemma~\ref{lem:II-flipcalc}(d); the constant $K_\sharp$ is fixed after this
uniform bound, so that the bounds $\hat G_b\le K_\sharp\tau^2$ below include
the flips.  The other conditions and the order of choices are unchanged.

\emph{Steps 1, 2} are verbatim, with $\beta^\theta:=b^\theta1_{[1,N_w]}$ and
$\beta^\pm:=b^\pm1_{[1,N_w]}\pm\frac12v1_{(N_w,\infty)}$;
$g''-g=-b^\theta1_{(N_w,\infty)}+(\text{block terms})\to0$.

\emph{Step 3 (base).}  We claim
\[
 \Exc'(\tau B^\diamond)\le\tfrac12\rho|\tau|V_{>N''}+\rho|\tau|\sum_{j\in
 F,\,j>N_w}(s_jv_j)_-+\mathrm{Fl}(\tau),\qquad\mathrm{Fl}(\tau)\le2\rho|\tau|
 m(4\rho|\tau|),
\]
and $\Exc'(\tau B^\theta)=0$.  On $\supp a'$, $a'_j+\tau B^\diamond_j=
(1-\tau\rho c)a'_j+\tau\rho\beta^\diamond_j$ with $|\tau\rho c|\le\frac12$, and
the summand of $\Exc'$ equals $|x+y|-|x|-\sgn(x)y$ with $x:=(1-\tau\rho
c)a'_j$, $y:=\tau\rho\beta^\diamond_j$.
(i) $j\in F\cap[1,N_w]$: $s_jx\ge|a''_j|/4$.  For $\diamond=\theta$
($|\tau|\le s_1$), $|y|\le\rho s_1|b^\theta_j|\le|a''_j|/8$: no flip.  For
$\diamond=\pm$ ($\sgn\tau=\diamond$) the anti-sign part of $y$ is at most
$\rho|\tau|\beta_j$ with $\beta=(sb^+)_-$, resp.\ $(sb^-)_+$; since
$|a''_j|\ge|a_j|$, a flip needs $|a_j|<4\rho|\tau|\beta_j$ and costs at most
$2\rho|\tau|\beta_j$ (Lemma~\ref{lem:II-flipcalc}(c)); the sum is at most
$2\rho|\tau|m(4\rho|\tau|)$.
(ii) Window contacts with mass, and contacts in $[1,N'']$: verbatim as in
[I,~Theorem~7.18].
(iii) $j\in F\cap(N_w,N'']$: $j\notin\supp a'$, $z'_j=s_j$, $\beta^\theta_j=0$
and $\beta^\pm_j=\pm\frac12v_j$; for $\diamond=\sgn\tau$ the summand is
$2(s_j\tau\rho\beta^\diamond_j)_-=\rho|\tau|(s_jv_j)_-$.
(iv) $j>N''$, $j\in F\cup K$: $z'_j=0$ and the summand is at most
$\frac12\rho|\tau||v_j|$ (and $0$ for $\theta$).
(v) Off $F\cup K$, $b^\pm$ and $v$ vanish.
This proves the claim.  Hence $\hat G_b$ of [I,~Theorem~7.18] acquires the
additional terms $\rho|\tau|\sum_{j\in F,\,j>N_w}(s_jv_j)_-+\mathrm{Fl}(\tau)$.

\emph{Steps 4, 5.}  The rebalancing is unchanged ($\mathrm{Fl}(\tau)\le
\tau^2$).  In Step~5 the first-order base terms are at most
$\frac12\rho|\tau|V_{>N''}+\rho|\tau|\sum_{j>N_w}(s_jv_j)_-+|\tau|o(s_1)\le
|\tau|\delta s_1(\frac1{64}+\frac1{32}+\frac1{64})\le\delta\tau^2/16$ for
$|\tau|>s_1$ at late stages, as in [I,~Theorem~7.18], and the flips
contribute $c'\mathrm{Fl}(\tau)\le\frac{\tau^2}2\cdot\frac\delta8$.  Hence
$p^*(f'+\tau g')\le1+\frac{\tau^2}2(\rho^2\kappa_w+\frac{6\delta}8)\le1+
\frac{\tau^2}2(1-\delta)$ for $0<|\tau|\le T_0$.  \emph{Step 6} is verbatim.
\end{proof}

The raises do for support coordinates what the window masses do for
contacts: they absorb the two-sided use of the average data below scale
$s_1$, at total mass $O(s_1)$ whatever $b^\theta$ is.  For data satisfying
the two-sided sparsity of $|b^\theta|$, $(sb^+)_-$, $(sb^-)_+$, or bounded
ratios $|b^\pm|\le C|a|$ on $F$, the theorem was obtained without raising;
the point here is that $b^\theta$ need not be sparse.

\begin{theorem}[Flip-extended engineered recovery]\label{thm:II-flipext}
Let $T$, $f$, $g$, $(b^\pm,\omega^\pm)$, $\rho$ and $I_-$ be as in
Theorem~\textup{\ref{thm:II-raised}}, but instead of \textup{(CS-side)} and
$\rho^2\kappa_w<1$ assume that for some $t_0>0$
\[
 \rho^2\kappa^{\rm fl}<1,\qquad\kappa^{\rm fl}:=\max_{\diamond=\pm}\Big(
 \Gamma_w(b^\diamond,\omega^\diamond)+\sup_{0<r\le t_0}\mathrm{Fl}_r(b^\diamond)
 \Big),
\]
where $\mathrm{Fl}_r(b^+):=2q_0\mathrm{Fl}_{(sb^+)_-}(r)/r^2$ and
$\mathrm{Fl}_r(b^-):=2q_0\mathrm{Fl}_{(sb^-)_+}(r)/r^2$.  Then
$(f,\rho g)\in\cl\NA((c_0,p),\ell_2^2)$.
\end{theorem}

\begin{proof}
Run the proof of Theorem~\ref{thm:II-raised} with $\delta:=(1-\rho^2
\kappa^{\rm fl})/2$ and $T_0\le t_0/(2\rho)$, without the flip condition of
Step~0.  In Step~3(i), at late stages $s_jx\ge|a_j|/(1+\epsilon)$ with
$\epsilon\to0$ along the construction ($c\to0$, $q^*(a'')\to1$), so
$\mathrm{Fl}(\tau)\le\mathrm{Fl}_\beta(\rho|\tau|(1+\epsilon))\le\frac{\rho^2\tau^2}
{2q_0}(1+\epsilon)^2\sup_{0<r\le t_0}\mathrm{Fl}_r(b^\diamond)$; this is
$O(\tau^2)$, so $K_\sharp$ exists, and its contribution $c'\mathrm{Fl}(\tau)$ to
$\hat\Gamma$ is at most $\frac{\rho^2\tau^2}2(\sup_r\mathrm{Fl}_r(b^\diamond)+
o(1))$ since $c'\to q_0$.  The $\theta$-data carry no flips, and
$\Gamma_w(b^\theta,\omega^\theta)\le\kappa_w\le\kappa^{\rm fl}$ by convexity.
Step~5 then gives $p^*(f'+\tau g')\le1+\frac{\tau^2}2(\rho^2\kappa^{\rm fl}+
\frac{5\delta}8+o(1))\le1+\frac{\tau^2}2(1-\delta)$ at late stages.
\end{proof}

\begin{theorem}[Windowed recovery at arbitrary base support]
\label{thm:II-masterinf}
Let $T$ be admissible, $f\in S_{p^*}$ \textup{(}$F$ arbitrary\textup{)},
$g\in\Cset(f)$, $\rho\in(0,1)$ and $\eta_0\in(0,2]$ with
$\rho^2(1+\eta_0)\le(1+\rho^2)/2$.  Suppose there are $c_\flat>0$, $C_+\ge1$,
$\beta\in\ell_1$ with $\beta\ge0$ and $m_\beta(x)=o(x)$, and triples
$(T_j,n_j,K_j)$ with $T_j\to0$, $n_j\ge24\rho^2K_j/(c_\flat(1-\rho^2))$ and
$K_jT_j/n_j\to0$, such that for every $j$ and every
$t\in\{T_j2^{1-s}:1\le s\le n_j\}$ there are a functional $g_t$ and two-piece
data $(b^\pm,\omega^\pm)$ for $g_t$ with:
\begin{enumerate}
\item[(i)] $p^*(f+rg_t)\le1+\frac{r^2}2(1+\eta_0)$ for $|r|\le c_\flat t$;
\item[(ii)] $p^*(g-g_t)\le K_jt$;
\item[(iii)] $\Gamma_w(b^\pm,\omega^\pm)\le1+\eta_0/2$;
\item[(iv)] $(s_jb^+_j)_-\le C_+|a_j|/t+\beta_j$ and $(s_jb^-_j)_+\le
C_+|a_j|/t+\beta_j$ for $j\in F$;
\item[(v)] either $\Delta d_m\ge0$ for all these data and all $m$, or the
set of blocks $m$ in which some of these data have $\Delta d_m<0$
satisfies \textup{(SC)}.
\end{enumerate}
Then $(f,\rho g)\in\cl\NA((c_0,p),\ell_2^2)$, along norm-attaining first rows
with finite base supports.
\end{theorem}

\begin{proof}
By [I,~Lemma~8.43] (whose proof uses only (i), (ii), $g\in\Cset(f)$ and the
convexity of $p^*$), $\rho\bar g_j\in\Cset(f)$ for large $j$ and $\rho\bar
g_j\to\rho g$, where $\bar g_j$ is the average of the $g_t$ over the $n_j$
scales.  The averages $\bar D_j$ of the data are two-piece data for $\bar
g_j$ (the representations and the side conditions are linear), $\Delta\bar
d_m$ is the average of the $\Delta d_m$, so (v) passes to $\bar D_j$ (a subset
of a set of blocks satisfying (SC) satisfies (SC)), and $\kappa_w(\rho\bar
D_j)\le\rho^2(1+\eta_0/2)<1$ by convexity of $\Gamma_w$.  By convexity of
$x\mapsto x_\pm$, (iv) holds for $\bar D_j$ with $t$ replaced by the least
scale $T'_j:=T_j2^{1-n_j}$, and Lemma~\ref{lem:II-flipcalc}(b) gives
$m(x)\le2m_{C_+|a|/T'_j}(2x)+2m_\beta(2x)=2m_\beta(2x)=o(x)$ for
$x<T'_j/(2C_+)$: (CS-side).  Theorem~\ref{thm:II-raised} applied to the mate
$\rho\bar g_j$ with data $\rho\bar D_j$ gives $(f,\rho'\rho\bar g_j)\in\cl\NA$
for all $\rho'<1$, hence $(f,\rho\bar g_j)\in\cl\NA$; let $j\to\infty$.
\end{proof}

All window theorems of this section are instances of
Theorem~\ref{thm:II-masterinf} (Lemma~\ref{lem:II-Uinf} turns (iv) and the
block size conditions into (i)), applied at $f$ or, through
Theorem~\ref{thm:II-ERT}, at nearby first rows.

\smallskip
We turn to the exact one-sided coefficient.  We say that $f$ is
\emph{block-tame} if every $Q_m$ is finite, there are no degenerate peaks,
and (MS) [I,~Definition~5.5] holds at $\xi$; nothing is assumed about $F$ or
$K$ (for finite $F$ this is (BT) of [I,~Definition~7.8]).  For $g$ with
$g(\xi)=0$ the one-sided coefficients $\gamma^\pm(g)$ are defined as in
[I,~Definition~7.7], with side admissibility as in
Definition~\ref{def:II-twopieceinf}.

\begin{theorem}[Exact one-sided coefficient at arbitrary base support]
\label{thm:II-onesidedinf}
Let $T$ be admissible, $f\in S_{p^*}$ block-tame \textup{(}$F$
arbitrary\textup{)}, and $g\in X^*$ with $g(\xi)=0$.  Then:
\begin{enumerate}
\item[(a)] $\liminf_{t\to0^\pm}2(p^*(f+tg)-1)/t^2\ge\gamma^\pm(g)$;
\item[(b)] if $\gamma^\pm(g)<\infty$, the infimum defining it is attained;
\item[(c)] for every side-$+$ admissible pair $(b,\omega)$ for $g$,
\[
 \limsup_{t\to0^+}\frac{2(p^*(f+tg)-1)}{t^2}\le\Gamma_w(b,\omega)+2q_0
 \limsup_{t\to0^+}\frac{\mathrm{Fl}_{(sb)_-}(t)}{t^2},
\]
and the same for $t\to0^-$, side-$-$ pairs and $(sb)_+$;
\item[(d)] if some side-$\pm$ pair attaining $\gamma^\pm(g)$ has a sparse
profile $(sb)_\mp$, then $\lim_{t\to0^\pm}2(p^*(f+tg)-1)/t^2=\gamma^\pm(g)$;
\item[(e)] every $g\in\Cset(f)$ has two-piece data with
$\kappa_w=\max(\gamma^+(g),\gamma^-(g))\le1$.
\end{enumerate}
\end{theorem}

\begin{proof}
(a), (b): Steps~2--4 of the proof of [I,~Theorem~7.10] use $F$ only through
the test points of Step~1, and the duality of Step~4 only through the cone
$C_+:=\{\chi_c\in c_{00}:\chi_c=0\text{ on }F,\ z_j\chi_{c,j}\le0\text{ on }
K\}$, whose polar condition $\sup_{\chi_c\in C_+}2b_\pi(\chi_c)<\infty$ is
exactly side-$+$ admissibility in the sense of
Definition~\ref{def:II-twopieceinf} (no condition on $F$).  We replace the
test points of Step~1, which used a Gram system on $F$, by the following
ones, valid for every $F$.  Let $k\in H$ with $\ip ke=0$ and $\chi_c\in C_+$,
put $\kappa:=\|k\|^2/(2q_0)$, $k_2:=-\kappa e$, $\chi:=Uk+\chi_c$,
$\chi_2:=Uk_2$, $\eta_t:=\xi+t\chi+t^2\chi_2$, $\mathbf h_t:=q_0e+tk+t^2k_2$
and $Z_t:=\eta_t-U\mathbf h_t$.  Since $\xi=q_0z+q_0Ue$, $Z_t=q_0z+t\chi_c$;
as $\chi_c$ vanishes on $F$, moves contacts inward and is finitely
supported, $\|Z_t\|_\infty\le q_0$ for small $t$.  Moreover
$\|\mathbf h_t\|^2=(q_0-t^2\kappa)^2+t^2\|k\|^2=q_0^2+t^4\kappa^2$.  As
$B_{q^{**}}=B_{\ell_\infty}+U(B_H)$, $q^{**}(\eta_t)\le\max(\|Z_t\|_\infty,
\|\mathbf h_t\|)\le q_0+O(t^4)$, and $a(\eta_t-\xi)=t\ip{U^*a}k+ta(\chi_c)+
t^2\ip{U^*a}{k_2}=-t^2\kappa\nu$.  Hence the base excess satisfies
$E_q(t\chi+t^2\chi_2)\le\frac{t^2}2\frac\nu{q_0}\|k\|^2+O(t^4)$, which is
the bound of Step~1.  The block test of Step~2 uses $\chi_2$ only through
$q(\chi_2)$ and through a term $w(t^2\delta_2)$ that cancels, and accepts
$\chi_2=Uk_2\in c_0$; the rest is unchanged.
(c) In the proof of [I,~Proposition~7.9],
$G_b(a+tb)=\mathrm{Fl}_{(sb)_-}(t)+\nu\Psi(tU^*b/\nu)$ by [I,~Lemma~7.1(b)] and
Lemma~\ref{lem:II-flipcalc}(c) (off $F$ the summands vanish by side
admissibility).  If the $\limsup$ of $\mathrm{Fl}_{(sb)_-}(t)/t^2$ is infinite
there is nothing to prove; otherwise $\hat G_b:=\mathrm{Fl}_{(sb)_-}(t)+
\frac{t^2}2h(b)(1+K_1t)=O(t^2)$, the rebalancing of that proof applies
([I,~Proposition~7.3] and the transfer data hold at every $f$), and the
additional term in $\hat\Gamma$ is $q_0\mathrm{Fl}_{(sb)_-}(t)$.  (d) follows
from (a) and (c), and (e) from (a), (b) and $p^*(f+tg)\le s(t)\le1+t^2/2$.
\end{proof}

\begin{corollary}[Block-tame first rows with infinite support]
\label{cor:II-BTinf}
Let $T$ be admissible and $f\in S_{p^*}$ block-tame \textup{(}$F$
arbitrary\textup{)}, and suppose that
$V_f:=\sum_m\big(\sum_{k\in Q_m}\lambda_{k,m}|u_{k,m}|+|R_m^*w_m|\big)$ has a
sparse profile.  Then $(f,g)\in\cl\NA((c_0,p),\ell_2^2)$ for every
$g\in\Cset(f)$ such that $|g|1_F$ has a sparse profile.
\end{corollary}

\begin{proof}
By Theorem~\ref{thm:II-onesidedinf}(e) there are two-piece data with
$\kappa_w\le1$, $b^\pm=g-\sum_mR_m^*(\omega^\pm_m-d_m(\omega^\pm_m)w_m)$ and
$\omega^\pm_m\in\R^{Q_m}$; hence $|b^\pm-g|\le CV_f$ with
$C:=\max_\pm\max(\|\omega^\pm\|_\infty,\max_m|d_m(\omega^\pm_m)|)$, and
$(sb^+)_-,(sb^-)_+\le|g|+CV_f$ on $F$.  By Lemma~\ref{lem:II-flipcalc}(b) these
profiles are sparse.  For block-tame $f$ every set of blocks satisfies
(SC) (the remark after [I,~Definition~7.16]; its proof concerns the
blocks only).  Theorem~\ref{thm:II-raised} applies for every $\rho<1$.
\end{proof}

\subsection{Window-pinned mates and room off the support}\label{subsec:II-room}

In this subsection $T$ is a signature-ladder operator (see the standing
conventions), with windows $\mathcal W(l)$, $T_{\rm hi}(l)$, $T_{\rm lo}(l)$,
$n^w_l$, carriers $l$, $u_l$, $\lambda_l$, $\Delta\theta_l$ and
$v_l:=\delta_lh_l/n_l$ as in [I,~Section~8].  Definitions [I,~8.3, 8.30,
8.33] are read verbatim for arbitrary $F$; in particular $J_\gamma=\{j\notin
F:|z_j|\le1-\gamma\}$ excludes $F$.  For $f$ with arbitrary $F$, a mate
$g\in\Cset(f)$ is \emph{window-pinned} if the condition of
[I,~Definition~8.30] holds.

\begin{proposition}[The clamped window certificate]\label{prop:II-clamped}
Let $T$ be a signature-ladder operator, $f\in S_{p^*}$ \textup{(}$F$
arbitrary\textup{)}, $g\in\Cset(f)$, $\eta\le\eta_*$ with
$\eta_\Gamma(\eta)\le1$, $l_*$ a window index, $t\in\mathcal W(l_*)$ with
$t\le\min(t_\eta,1)$, and let $(B_\pm,\Theta_\pm)$ be a two-sided
decomposition of $g$ at scale $t$ with
\[
 \textup{(PIN$_K$)}\qquad\sum_l|\Delta\theta_l|\le Kt\qquad(K\ge1).
\]
Let $b^{\rm cl}$ be the clamp of [I,~Lemma~8.28] \textup{(}$b^{\rm cl}:=0$
off $F$\textup{)}, $M_t:=\min\{M:\alpha(M)\le t^2\}$, $F_t:=F\cap[1,M_t]$,
$a_t:=a1_{F_t}$, $\kappa_t:=(b^{\rm cl}1_{F_t})(\zh)/a_t(\zh)$,
$b_t:=b^{\rm cl}1_{F_t}-\kappa_ta_t$, let $\omega^c_m$ be as in
[I,~Definition~8.14], and $c_t:=(b_t,(\omega^c_m)_m)$.  Assume $Kt\le1$ and
$t^2\le1/(2(1+\|U\|))$, and put $K_\kappa:=\frac1{2q_0}+(1+\|U\|)(2K+
\frac3{2q_0}+2)$.  Then:
\begin{enumerate}
\item[(a)] $a_t(\zh)\ge\frac12$; $c_t$ is a balanced finite certificate with
$\supp b_t\subseteq F_t$ finite and $b_t(\xi)=0$; on $\supp\omega^c_m$ one
has $k\in Q_m$, $\gap_m(k)\ge t^2$ and $|\omega^c_m(k)|\le2\gap_m(k)/t$;
\item[(b)] $|\kappa_t|\le2K_\kappa t$, and if $2K_\kappa t^2\le1$ then
$|b_t(j)|\le3|a_j|/t$ for all $j$;
\item[(c)] $\|B_+-b_t\|_1\le(2K+\frac3{2q_0}+2+2K_\kappa)t$ and
$\|g-g_{c_t}\|_1\le K'_2Kt$;
\item[(d)] $\sqrt{\Gamma_w(c_t)}\le\sqrt{1+\eta_\Gamma(\eta)}+K'_4Kt$,
\end{enumerate}
where $K'_2,K'_4$ depend only on $f$, $N$ and the operator.
\end{proposition}

\begin{proof}
[I,~Proposition~8.13(b),(c)] hold with this $K$: their proofs use only
(PIN$_K$), [I,~Lemma~8.7(b)] and [I,~Lemma~8.8].
(a) $|a_t(\zh)-1|=|(a-a_t)(\zh)|\le\|\zh\|_\infty\alpha(M_t)\le(1+\|U\|)t^2
\le\frac12$, since $a(\zh)=1$; so $b_t(\zh)=0$ and $b_t(\xi)=q_0b_t(\zh)=0$.
The statement on $\omega^c_m$ is [I,~Proposition~8.15(a)], whose proof does
not use $F$.
(b) We have $(b^{\rm cl}1_{F_t})(\zh)=B_+(\zh)-(B_+1_{F^c})(\zh)-((B_+-
b^{\rm cl})1_F)(\zh)-(b^{\rm cl}1_{F\setminus F_t})(\zh)$, with
$|B_+(\zh)|\le t/(2q_0)$ [I,~Lemma~8.7(a)], $\|B_+1_{F^c}\|_1\le Kt+t/q_0$
[I,~Proposition~8.13(b)], $\|(B_+-b^{\rm cl})1_F\|_1\le\|\Delta B1_F\|_1+
t/(2q_0)\le Kt+t/(2q_0)$ [I,~Lemma~8.28], and $\|b^{\rm cl}1_{F\setminus
F_t}\|_1\le(2/t)\alpha(M_t)\le2t$.  As $\|\zh\|_\infty\le1+\|U\|$, this gives
$|(b^{\rm cl}1_{F_t})(\zh)|\le K_\kappa t$ and $|\kappa_t|\le2K_\kappa t$.  Then
$|b_t(j)|\le2|a_j|/t+2K_\kappa t|a_j|\le3|a_j|/t$.
(c) $B_+-b_t=B_+1_{F^c}+(B_+-b^{\rm cl})1_F+b^{\rm cl}1_{F\setminus F_t}+
\kappa_ta_t$; use the bounds of (b).  The block part of $g-g_{c_t}$ is bounded
as in the proof of [I,~Proposition~8.15(c)], which uses only
[I,~Lemmas~8.8, 8.10], (P2) and [I,~Proposition~8.13(a),(c)].
(d) As in the proof of [I,~Proposition~8.15(d)]: $\sqrt{\Gamma_w}$ is a
seminorm, $\Gamma_w(B_+,\Theta_+)\le1+\eta_\Gamma(\eta)$, and the remainder
pair has $q_0h(x)\le\|U\|^2\|x\|_1^2/\nu$ and $\sigma_mH_m(X_m)\le(\sum_k
\lambda_{k,m}|X_m(k)|)^2/C_m$, both at most $(\text{const}\cdot Kt)^2$ by (c).
\end{proof}

\begin{theorem}[Window-pinned mates at arbitrary base support]
\label{thm:II-Bstarinf}
Let $T$ be a signature-ladder operator and $f\in S_{p_N^*}$ \textup{(}$F$
arbitrary\textup{)}.  If $g\in\Cset(f)$ is window-pinned at $f$, then
$(f,g)\in\cl\NA((c_0,p_N),\ell_2^2)$.  No condition on the contact set, on
the rooms of the signature sets or on the blocks of $f$ is needed.
\end{theorem}

\begin{proof}
Fix $\rho\in(0,1)$, choose $\eta_0\in(0,2]$ with $\rho^2(1+\eta_0)\le
(1+\rho^2)/2$, $\kappa_0:=\frac12(\sqrt{1+\eta_0/2}-1)$, $\varepsilon_{\rm tr}:=
\eta_0/2$, and $\eta\le\eta_*$ with $\eta_\Gamma(\eta)\le1$ and
$\sqrt{1+\eta_\Gamma(\eta)}\le1+\kappa_0$.  Let $c_\flat,t_1$ be given by
Lemma~\ref{lem:II-Uinf} for the constant sequence $f_j=f$, $U_*=0$, $C_+:=3$,
$A_0:=4$, $A_2:=1$, $\gamma_B:=1$.  Let $(l_j,K_j)$ be the data of window
pinning and $K:=K_j+6\le7K_j$: by [I,~Lemma~8.10] and (P2) the fine carriers
contribute at most $6t^2$, so (PIN$_K$) holds on $\mathcal W(l_j)$.  Since
$K_jT_{\rm hi}(l_j)\to0$, for large $j$ every dyadic scale $t$ of $\mathcal
W(l_j)$ satisfies $t\le\min(t_\eta,t_1,1)$, $Kt\le1$, $t^2\le1/(2(1+\|U\|))$,
$2K_\kappa t^2\le1$ and $K'_4Kt\le\kappa_0$.  The clamped certificate $c_t$
of Proposition~\ref{prop:II-clamped} then has $\Gamma_w(c_t)\le(1+2\kappa_0)^2=
1+\eta_0/2\le2$ and $t\|b_t\|_1\le3\|a\|_1\le A_0$; since $\supp b_t\subseteq F$,
the pair is admissible on both sides, (OS2$'$) holds with $C_+=3$ on both sides,
and every block coordinate is of kind [1].  Lemma~\ref{lem:II-Uinf} gives
$p^*(f+rg_{c_t})\le1+\frac{r^2}2(1+\eta_0)$ for $|r|\le c_\flat t$, and
$p^*(g-g_{c_t})\le(1+\|U\|)K'_2Kt\le K''_jt$ with $K''_j:=7(1+\|U\|)K'_2K_j$.
As $n^w_{l_j}/K''_j\to\infty$ and $K''_jT_{\rm hi}(l_j)/n^w_{l_j}\to0$,
Corollary~\ref{cor:II-windowedinf} (with $t^{(j)}:=T_{\rm hi}(l_j)$,
$n_j:=n^w_{l_j}$) gives $(f,\rho g)\in\cl\NA$.  Let $\rho\to1$.
\end{proof}

\begin{theorem}[Infinite base support with signature room]\label{thm:II-Binf}
Let $T$ be a signature-ladder operator and $N\ge1$.  Every $f\in S_{p_N^*}$
satisfying \textup{(SR)} of [I,~Definition~8.3], with $F$ finite or
infinite, belongs to $\Rec$.  In particular the claim of
[I,~Remark~8.46] holds.
\end{theorem}

\begin{proof}
[I,~Proposition~8.13(a)] holds at arbitrary $F$: its proof uses
[I,~Lemma~8.7(b)] (where $J_\gamma\subseteq F^c$), the pinning inequality
[I,~Lemma~8.11] (which uses only (P1) and the sets $S_l\cap J_\gamma$) and
[I,~Lemma~8.12].  As in the proof of [I,~Theorem~8.18],
$K_1\Lambda_f(l)T_{\rm hi}(l)\to0$ and $n^w_l/\Lambda_f(l)\to\infty$, because
$\Lambda_f(l)\le\varpi_0^{-l^2}\Lambda^\circ(l)$ and (P3).  So every
$g\in\Cset(f)$ is window-pinned with $K_l:=K_1\Lambda_f(l)$, and
Theorem~\ref{thm:II-Bstarinf} applies.
\end{proof}

\begin{theorem}[Sign-mixed room at arbitrary base support]
\label{thm:II-Bpminf}
Let $T$ be a signature-ladder operator and $N\ge1$.  Let $f\in S_{p_N^*}$
\textup{(}$F$ arbitrary\textup{)} satisfy
\[
 \mathfrak r_l:=\min_{\sigma=\pm1}\sum_{s\in S_l\setminus F}v_l(s)
 (1+\sigma z_s)>0\qquad(l\in\mathcal L_N)
\]
and \textup{(W$^\pm$)} of [I,~Definition~8.33].  Then
$f\in\Rec$.
\end{theorem}

\begin{proof}
The proof of [I,~Lemma~8.34] does not use the finiteness of $F$ (the rooms
are computed on $S_l\setminus F$, and [I,~Lemma~8.24] holds at every $F$).
Hence, as in the proof of [I,~Theorem~8.35], every mate is window-pinned
with $K_j:=(1/q_0+22)\Lambda^\pm_f(l_j)$; apply
Theorem~\ref{thm:II-Bstarinf}.
\end{proof}

At infinite $F$ the support itself can provide room at the scales of a
window, through the cushions.  For $M\ge0$ let $F_{\le M}:=F\cap[1,M]$,
\[
 \mathfrak r^{[M]}_l:=\min_{\sigma=\pm1}\sum_{s\in S_l\setminus F_{\le M}}
 v_l(s)(1+\sigma z_s),\qquad\Lambda^c_f(l,M):=\prod_{l''\le l,\ l''\in
 \mathcal L_N}\Big(1+\frac3{\mathfrak r^{[M]}_{l''}}\Big),
\]
($z_s=s_s$ on $F$) and $M(t):=\min\{M:\alpha(M)\le t^2\}$.

\begin{theorem}[Cushion room]\label{thm:II-cushionroom}
Let $T$ be a signature-ladder operator, $N\ge1$, and let $f\in S_{p_N^*}$
satisfy
\[
 \textup{(W$^c$)}\qquad\liminf_{l\to\infty}\frac{\Lambda^c_f(l,M(T_{\rm lo}(l)))}
 {l2^{l^3}\Lambda^\circ(l)}=0
\]
\textup{(}with $\Lambda^c_f(l,M):=\infty$ if some factor has
$\mathfrak r^{[M]}_{l''}=0$\textup{)}.  Then $f\in\Rec$.  Since $\mathfrak r^{[M]}_l\ge\mathfrak r_l$,
\textup{(W$^\pm$)} implies \textup{(W$^c$)}.
\end{theorem}

\begin{proof}
Let $t\in\mathcal W(l_*)$, $t\le\min(t_\eta,1)$, a decomposition at scale
$t$, and $M\ge M(t)$.  Put $E^c_l:=\sum_{s\in S_l\setminus F_{\le M}}
\phi_{z_s}(\Delta B(s))$.  By [I,~Lemma~8.24] off $F$, and on $F\cap
(M,\infty)$ by $\phi_{s_j}(x)=2(s_jx)_-$ [I,~Lemma~8.23(a)] and
Lemma~\ref{lem:II-cushion},
\[
 \sum_lE^c_l\le\frac t{q_0}+2\sum_{j\in F,\,j>M}\Big(\frac{2|a_j|}t+
 \mathfrak f^+_j+\mathfrak f^-_j\Big)\le\Big(\frac2{q_0}+4\Big)t .
\]
For $s\in S_l$, $-\Delta B(s)=\Delta\theta_lv_l(s)+r_l(s)$ with
$r_l(s):=\sum_{l'>l}\Delta\theta_{l'}y_{l'}(s)/n_{l'}$ ((P1), [I,~(15)]), so
$|\Delta\theta_l|v_l(s)(1+z_s\sgn\Delta\theta_l)\le\phi_{z_s}(\Delta B(s))+
2|r_l(s)|$ [I,~Lemma~8.23(c),(d)]; summing over $s\in S_l\setminus F_{\le M}$
gives $\mathfrak r^{[M]}_l|\Delta\theta_l|\le E^c_l+2\sum_{l'>l}\pi_{l',l}
|\Delta\theta_{l'}|$.  Unrolling as in the proof of [I,~Lemma~8.34],
$\sum_l|\Delta\theta_l|\le(2/q_0+26)\Lambda^c_f(l_*,M)t$.  Since $M(\cdot)$
is nonincreasing, $M:=M(T_{\rm lo}(l_*))$ serves all scales of
$\mathcal W(l_*)$.  By (W$^c$) and (P3) every mate is window-pinned;
Theorem~\ref{thm:II-Bstarinf} applies.
\end{proof}

\begin{remark}[Scope of cushion room]\label{rem:II-cushionroom}
If $z$ is constant of modulus $1$ on $S_l\setminus F_{\le M_0}$ for some $l\in\mathcal L_N$
and some $M_0$ (the signature set is swallowed, off a finite set, by contacts
and support coordinates of one sign), then $\mathfrak r^{[M]}_l=0$ for all
$M\ge M_0$, and \textup{(W$^c$)} fails whatever the decay of $a$.  So
Theorem~\ref{thm:II-cushionroom} adds to Theorem~\ref{thm:II-Bpminf} only
for signature sets on which $a$ is sign-mixed (or $z$ is sign-mixed off
$F$), and then only for fast decay of $a$ on those sets.  Signature sets
swallowed by the support are treated in Subsections~\ref{subsec:II-R1}
and~\ref{subsec:II-RS}.
\end{remark}

\subsection{Swallowed carriers and cushion sparsity}\label{subsec:II-R1}

In this subsection $T$ is a signature-ladder operator.  We extend
[I,~Definition~8.37] to arbitrary base support.

\begin{definition}[Bad carriers at arbitrary base support]
\label{def:II-badinf}
Let $f\in S_{p_N^*}$ \textup{(}$F$ arbitrary\textup{)}.  For
$l\in\mathcal L_N$ let $\mathfrak r_l:=\min_{\sigma=\pm1}\sum_{s\in
S_l\setminus F}v_l(s)(1+\sigma z_s)$ (the room off the support; $0$ if
$S_l\subseteq F$), and $\mathcal B:=\{l\in\mathcal L_N:\mathfrak r_l=0\}$
(the \emph{bad} carriers).  For $l\in\mathcal B$ let $\epsilon_l$ be the sign
with $z\equiv\epsilon_l$ on $S_l\setminus F$ ($\epsilon_l:=1$ if
$S_l\subseteq F$).  Put $\mathcal B_K:=\{l\in\mathcal B:S_l\setminus F
\text{ infinite}\}$ (\emph{contact-swallowed}) and $\mathcal B_F:=\mathcal B
\setminus\mathcal B_K$ (\emph{swallowed by the support}: $S_l\setminus F$ is
finite), $\mathcal B_{\rm pk}:=\{l\in\mathcal B:k(l)\in P_{m(l)}\}$ and
$\mathcal B_{\rm np}:=\mathcal B\setminus\mathcal B_{\rm pk}$.  For
$l\in\mathcal B$ let $\tau_l:=-\epsilon_l\Delta\theta_l$ (for a given
two-sided decomposition) and let $q_l=\epsilon_l\Phi_{m(l)}(k(l))
w_{m(l)}(k(l))/(m(l)C_{m(l)})$ be its $d$-weight (as in Definition~\ref{def:II-class}, with the sign $\epsilon_l$ of $l$).  For good $l$ let $S^\star_l:=S_l\setminus
\big(F\cup\bigcup_{l'\in\mathcal B,\,l'>l}\supp y_{l'}\big)$,
$\mathfrak r^\star_l:=\min_{\sigma=\pm1}\sum_{s\in S^\star_l}v_l(s)
(1+\sigma z_s)$ and $\Lambda^\star_f(l):=\prod_{l''\le l,\,l''\in\mathcal L_N
\setminus\mathcal B}(1+3/\mathfrak r^\star_{l''})$.  Finally
$U_{\mathcal B}(j):=\sum_{l\in\mathcal B}|u_l(j)|$ and $m_{\mathcal B}:=
m_{U_{\mathcal B}}$ (profile on $F$, Definition~\ref{def:II-flip}).
Consider:
\begin{enumerate}
\item[(W$^\star$)] $\mathfrak r^\star_l>0$ for every good $l$, and
$\liminf_l\Lambda^\star_f(l)/(l2^{l^3}\Lambda^\circ(l))=0$;
\item[(H2)] every block $m$ has a peak $k$ with $\alpha_m(k)\ne0$ and
$\mathfrak j(k,m)\notin\mathcal B$;
\item[(H3-inf)] no $l\in\mathcal B_K$ has $k(l)$ a degenerate peak with
$\sgn w_{m(l)}(k(l))=\epsilon_l$, and no $l\in\mathcal B_F$ has $k(l)$ a
degenerate peak;
\item[($\mathcal B_{\rm fin}$)] $\mathcal B$ is finite;
\item[(CS$_{\mathcal B}$)] $m_{\mathcal B}(y)=o(y)$ as $y\to0^+$
(\emph{cushion sparsity});
\item[(H4-inf)] $\sup_{l\in\mathcal B}\sup_{j\in F}|u_l(j)|/|a_j|<\infty$
(\emph{cushion domination}).
\end{enumerate}
\end{definition}

For finite $F$, $\mathcal B_F=\emptyset$, (H3-inf) is (H3), and under
($\mathcal B_{\rm fin}$) the product $\Lambda^\star_f$ differs from that of
[I,~Definition~8.37] by a constant factor.  The product is taken over good
carriers only because carriers of $\mathcal B_F$ may have no room at all.

\begin{lemma}[Cushion sparsity]\label{lem:II-R0}
Assume \textup{($\mathcal B_{\rm fin}$)}.  \textup{(a)} \textup{(H4-inf)}
implies \textup{(CS$_{\mathcal B}$)}.  \textup{(b)} \textup{(CS$_{\mathcal
B}$)} holds iff the profile of $v_l1_{S_l\cap F}$ is sparse for every
$l\in\mathcal B$ with $S_l\cap F$ infinite.
\end{lemma}

\begin{proof}
(a) If $|u_l(j)|\le c_{\mathcal B}|a_j|$ on $F$, then $U_{\mathcal B}\le
|\mathcal B|c_{\mathcal B}|a|$ on $F$, and the set defining
$m_{\mathcal B}(y)$ is empty for $y<1/(|\mathcal B|c_{\mathcal B})$.
(b) $T_{\mathcal B}:=\bigcup_{l\in\mathcal B}\supp y_l$ is finite.  By (P1),
$U_{\mathcal B}=v_l$ on $S_l\setminus T_{\mathcal B}$ for $l\in\mathcal B$,
and $U_{\mathcal B}$ vanishes off $T_{\mathcal B}\cup\bigcup_{l\in\mathcal B}
S_l$.  Each of the finitely many $j\in F\cap T_{\mathcal B}$ has $a_j\ne0$
and leaves the sets defining the profiles for small $y$.
\end{proof}

\begin{lemma}[Exact switching at arbitrary base support]
\label{lem:II-exactswitchinf}
Let $f\in S_{p_N^*}$ \textup{(}$F$ arbitrary\textup{)} satisfy
\textup{(W$^\star$)}, \textup{(H2)}, \textup{(H3-inf)} and
\textup{($\mathcal B_{\rm fin}$)}, let $g\in\Cset(f)$, $\eta\le\eta_*$ with
$\eta_\Gamma(\eta)\le1$, and put $K_\star:=(1/q_0+22)\Lambda^\star_f(l_*)$.
There are $C_f$ and $l_f$, depending only on $f$, $N$ and the operator,
such that for $l_*\ge l_f$, $t\in\mathcal W(l_*)$ with $t\le\min(t_\eta,1)$,
and every two-sided decomposition of $g$ at scale $t$:
\begin{enumerate}
\item[(a)] $\sum_{l\notin\mathcal B}|\Delta\theta_l|\le K_\star t$ and
$\sum_{l>l_*}|\Delta\theta_l|\le6t^2$; and [I,~Lemma~8.39](a)--(c) hold;
\item[(b)] there is $\tau'\in\R^{\mathcal B}$ with $\tau'_l\ge0$
\textup{(}$l\in\mathcal B_K$\textup{)}, $\tau'_l=0$ \textup{(}$l\in\mathcal
B_{\rm pk}$\textup{)}, $\sum_{m(l)=m}q_l\tau'_l=0$ for every $m$,
such that $V':=\sum_{l\in\mathcal B}\epsilon_l\tau'_lu_l1_{F^c}$ is
$z$-signed and supported in $K$, $\sum_{l\in\mathcal B}|\tau_l-\tau'_l|\le
C_fK_\star t$ and $\|\Delta B1_{F^c}-V'\|_1\le C_fK_\star t$.
\end{enumerate}
\end{lemma}

\begin{proof}
(a) The good inequalities in the proof of [I,~Lemma~8.38] use the sets
$S^\star_l\subseteq F^c$, the budget [I,~Lemma~8.24] (valid at every $F$) and
(P1); bad carriers never occur on their right sides, because the targets of
all bad carriers are removed from $S^\star_l$ and other signatures vanish
on $S_l$.  Unrolling over the good carriers gives the bound.  The proof of
[I,~Lemma~8.39] uses only [I,~Lemma~8.27], [I,~(17)], (a) and (H2).
(b) Let $l_f\ge\max\mathcal B$ and $l_f\ge\max_ml^\circ_m$
([I,~Lemma~8.39]).  We run case ($\mathcal B_{\rm fin}$) of the proof of
[I,~Lemma~8.40] with two changes: $T_0:=\bigcup_{l\in\mathcal B}\supp y_l\cup
\bigcup_{l\in\mathcal B_F}(S_l\setminus F)$ (finite), and the cone
$\mathcal Z_f$ has no sign row for $l\in\mathcal B_F$.  For $j\notin F\cup
T_0$ the only bad vector that can be nonzero at $j$ is the signature of the
$l\in\mathcal B_K$ with $j\in S_l$, where $z_j=\epsilon_l$; hence the cost
$\mathfrak c(\tau):=\sum_{j\notin F}\phi_{z_j}(L_j(\tau))$ equals
$\sum_{l\in\mathcal B_K}2(\tau_l)_-\mathfrak m_l+\sum_{j\in T_0\setminus F}
\phi_{z_j}(L_j(\tau))$ with $\mathfrak m_l:=\|v_l1_{S_l\setminus(F\cup T_0)}
\|_1>0$.  The bad peaks are controlled by [I,~Lemma~8.39](c) and (H3-inf):
either $\alpha_{m(l)}(k(l))\ne0$, or $l\in\mathcal B_K$ is at a degenerate peak
with $\sgn w(k(l))=-\epsilon_l$, and then $\tau_l\le\lambda_lK_dt$ and
$(\tau_l)_-\le\mathfrak c(\tau)/(2\mathfrak m_l)$.  The rest of the proof
(Hoffman's bound, [I,~(18)]) is unchanged.
\end{proof}

\begin{lemma}[Common shift and deep assignment]\label{lem:II-commonshift}
\textup{(a)} Let $x,y\in\R$ and $A\ge0$ with $x-y\ge-2A$.  Then
$\{\sigma:x-\sigma\ge-A,\ y-\sigma\le A\}=[y-A,x+A]\ne\emptyset$, and its
point $\sigma$ closest to $0$ satisfies $|\sigma|=(-x-A)_++(y-A)_+$.

\textup{(b)} Let $\alpha\ge0$, $A:=2\alpha$, $\varphi^\pm\ge0$, and let
$x,y_0,\xi,e\in\R$ with $x\ge-\alpha-\varphi^+$, $y_0\le\alpha+\varphi^-$ and
$x-y_0=\xi+e$.  Put $\beta:=x-\sigma$, where $\sigma$ is given by \textup{(a)}
for $(x,x-\xi,A)$, if $\xi\ge-2A$, and $\beta:=-A$ if $\xi<-2A$.  Then
\textup{(i)} $|x-\beta|\le\varphi^++\varphi^-+|e|$; \textup{(ii)}
$\beta\ge-A$; \textup{(iii)} $\beta-\xi\le A$ if $\xi\ge-2A$, and
$\beta-\xi=|\xi|-A\le|\xi|$ if $\xi<-2A$; \textup{(iv)} $|\beta|\le A+|\xi|$.
\end{lemma}

\begin{proof}
(a) The interval is nonempty because $y-A\le x+A$; at most one of $y-A>0$,
$x+A<0$ holds.  (b) If $\xi\ge-2A$, then $y:=x-\xi=y_0+e$; (ii), (iii) follow
from (a), and $|\sigma|\le(\alpha+\varphi^+-2\alpha)_++(\alpha+\varphi^-+|e|-
2\alpha)_+\le\varphi^++\varphi^-+|e|$; (iv): $-A\le\beta\le x-(y-A)=\xi+A$.
If $\xi<-2A$, (ii)--(iv) are clear, $x+A\ge\alpha-\varphi^+\ge-\varphi^+$, and
$x=y_0+\xi+e\le\alpha+\varphi^--4\alpha+|e|$, so $x+A\le\varphi^-+|e|$; hence
$|x-\beta|=|x+A|\le\varphi^++\varphi^-+|e|$.
\end{proof}

\begin{theorem}[Cushion-sparse swallowing]\label{thm:II-R1}
Let $T$ be a signature-ladder operator and $N\ge1$.  Every $f\in S_{p_N^*}$
\textup{(}$F$ arbitrary\textup{)} satisfying \textup{(W$^\star$)},
\textup{(H2)}, \textup{(H3-inf)}, \textup{($\mathcal B_{\rm fin}$)} and
\textup{(CS$_{\mathcal B}$)} belongs to $\Rec$.  In particular this holds
under \textup{(H4-inf)} in place of \textup{(CS$_{\mathcal B}$)}.
\end{theorem}

\begin{proof}
Fix $g\in\Cset(f)$ and $\rho\in(0,1)$; choose $\eta_0,\kappa_0,
\varepsilon_{\rm tr},\eta$ as in the proof of Theorem~\ref{thm:II-Bstarinf}.
Let $l_*\ge l_f$, $t\in\mathcal W(l_*)$ with $t\le\min(t_\eta,1)$, and a
two-sided decomposition at scale $t$.  Constants $C_f,C_\tau,K_\kappa,
A^R_0,A_2$ below depend only on $f$, $N$ and the operator; $K^R_2,K^R_4$ are
such constants times $\Lambda^\star_f(l_*)$.

\emph{Step 1 (pinning, exact switching).}  Lemma~\ref{lem:II-exactswitchinf}
gives (a) and $\tau'$, $V'$ as in (b).

\emph{Step 2 (bounded switching).}  Lemma~\ref{lem:II-boundedfree}(b) with
$\mathcal U:=\{(k(l),m(l)):l\in\mathcal B\}$ gives $|\Delta\theta_l|\le
C_{\mathcal U}(1+K_\star t)$ for $l\in\mathcal B$.  Assume $K_\star t\le1$ and
$C_fK_\star t\le1$; then $|\tau'_l|\le2C_{\mathcal U}+1=:C_\tau$.  Put
$X:=\sum_{l\in\mathcal B}\epsilon_l\tau'_lu_l$.  Then $X1_{F^c}=V'$,
$|X_j|\le C_\tau U_{\mathcal B}(j)$, and $\Delta B-X=\sum_{l\in\mathcal B}
\epsilon_l(\tau_l-\tau'_l)u_l-\sum_{l\notin\mathcal B}\Delta\theta_lu_l$, so
$\|\Delta B-X\|_1\le(C_f+1)K_\star t$.

\emph{Step 3 (exact data).}  Let $\chi$, $\mathfrak e_\pm$ be given by
[I,~Lemma~8.26] for $e:=\Delta B1_{F^c}-V'$: $B_+1_{F^c}=\chi V'+
\mathfrak e_+$, $B_-1_{F^c}=-(1-\chi)V'+\mathfrak e_-$, $\|\mathfrak e_\pm\|_1
\le C_fK_\star t+t/q_0$.  For $j\in F$ apply
Lemma~\ref{lem:II-commonshift}(b) with $\alpha:=|a_j|/t$,
$\varphi^\pm:=\mathfrak f^\pm_j$, $x:=s_jB_+(j)$, $y_0:=s_jB_-(j)$, $\xi:=s_jX_j$
and $e:=s_j(\Delta B-X)_j$ (the hypotheses are Lemma~\ref{lem:II-cushion}),
and let $b^+_1(j):=s_j\beta_j$; $b^+_1:=0$ off $F$.  With $A_j:=2|a_j|/t$ and
$D_t:=\{j\in F:s_jX_j<-2A_j\}$ (the \emph{deep set}), the lemma gives
\begin{gather*}
 \|B_+1_F-b^+_1\|_1\le\frac t{2q_0}+(C_f+1)K_\star t,\qquad s_jb^+_1(j)\ge-A_j,
 \\ s_j(b^+_1-X)_j\le\begin{cases}A_j&j\notin D_t,\\|X_j|&j\in D_t,
 \end{cases}
\end{gather*}
and $|b^+_1(j)|\le A_j+|X_j|$.  Put $\kappa:=(b^+_1+\chi V')(\zh)$,
$b^+:=b^+_1+\chi V'-\kappa a$, $b^-:=b^+-X$; in block $m$ let
$\omega^+_m:=\omega^c_m$ [I,~Definition~8.14] off the coordinates $k(l)$,
$l\in\mathcal B_{\rm np}$, $\omega^+_m(k(l)):=\omega_{+,m}(k(l))$ for
$l\in\mathcal B_{\rm np}$, $\omega^-_m:=\omega^+_m+\sum_{l\in\mathcal
B_{\rm np},\,m(l)=m}(\epsilon_l\tau'_l/\lambda_l)e_{k(l)}$, and
$g_t:=b^++\sum_mR_m^*(\omega^+_m-d_m(\omega^+_m)w_m)$.  Then:
(a) Off $F$, $b^+=\chi V'$ and $b^-=-(1-\chi)V'$, so the pairs are side-$+$,
resp.\ side-$-$ admissible; $b^+-b^-=X=\sum_{l\in\mathcal B_{\rm np}}
\epsilon_l\tau'_lu_l=\sum_mR_m^*(\omega^-_m-\omega^+_m)$ (as $\tau'=0$ on
$\mathcal B_{\rm pk}$) and $d_m(\omega^-_m)-d_m(\omega^+_m)=\sum_{m(l)=m}
q_l\tau'_l=0$; so both pairs represent $g_t$ and are $d$-neutral;
$b^+(\zh)=0$ and $b^\pm(\xi)=0$ [I,~Lemma~3.2].
(b) Since $(b^+_1+\chi V')(\zh)=B_+(\zh)-\mathfrak e_+(\zh)-(B_+1_F-b^+_1)(\zh)$,
[I,~Lemma~8.7(a)] gives $|\kappa|\le K_\kappa t$.  If $K_\kappa t^2\le1$, then on
$F$: $(s_jb^+_j)_-\le3|a_j|/t$, and $(s_jb^-_j)_+\le3|a_j|/t$ off $D_t$ and
$\le|X_j|\le C_\tau U_{\mathcal B}(j)$ on $D_t$.
(c) $t\|b^\pm\|_1\le A^R_0:=3+3C_\tau|\mathcal B|$ (use $\sum_FA_j\le2/t$,
$\|U_{\mathcal B}\|_1\le|\mathcal B|$).
(d) $B_+-b^+=\mathfrak e_++(B_+1_F-b^+_1)+\kappa a$; off $F$, $B_--b^-=
\mathfrak e_-$; on $F$, $B_--b^-=(B_+-b^+)-(\Delta B-X)$; all are
$O(K_\star t)$ in $\ell_1$.  The block parts are compared with
$\Theta_\pm$ exactly as in the proof of [I,~Lemma~8.41(b),(c)].  Hence
$\|g-g_t\|_1\le K^R_2t$ and $\Gamma_w(b^\pm,\omega^\pm)\le(\sqrt{1+
\eta_\Gamma(\eta)}+K^R_4t)^2$.
(e) As in [I,~Lemma~8.41(d)], every block coordinate is of kind [1], or of
kind [2] with $\gamma_B:=\gamma_{\mathcal B}>0$ and $A_2:=4+C_\tau/
\lambda_{\mathcal B}$, $\lambda_{\mathcal B}:=\min_{\mathcal B}\lambda_l$.

\emph{Step 4 (local validity).}  If moreover $K^R_4t\le\kappa_0$ and $t\le
t_1$, Lemma~\ref{lem:II-Uinf} (constant sequence, $U_*:=U_{\mathcal B}$,
$C_+:=3$, $C_*:=C_\tau$, $A_0:=A^R_0$, $A_2$, $\gamma_B$; its hypothesis on
$\bar m$ is (CS$_{\mathcal B}$)) applies to the $+$ data for $0<r\le
c_\flat t$ and to the $-$ data for $-c_\flat t\le r<0$; both represent $g_t$
and have $\Gamma_w\le(1+2\kappa_0)^2=1+\eta_0/2\le2$.  So
$p^*(f+rg_t)\le1+\frac{r^2}2(1+\eta_0)$ for $|r|\le c_\flat t$, and
$p^*(g-g_t)\le(1+\|U\|)K^R_2t$.

\emph{Step 5.}  By (W$^\star$) choose window indices $l_j$ with
$\Lambda^\star_f(l_j)/(l_j2^{l_j^3}\Lambda^\circ(l_j))\to0$.  By (P3) all the
smallness conditions above hold on $\mathcal W(l_j)$ for large $j$, and
$K_j:=(1+\|U\|)K^R_2(l_j)$, $n_j:=n^w_{l_j}$, $T_j:=T_{\rm hi}(l_j)$ satisfy
$n_j\ge24\rho^2K_j/(c_\flat(1-\rho^2))$ and $K_jT_j/n_j\to0$.
Theorem~\ref{thm:II-masterinf} applies with $C_+:=3$ and
$\beta:=C_\tau U_{\mathcal B}$, whose profile is sparse by
(CS$_{\mathcal B}$) and Lemma~\ref{lem:II-flipcalc}(b); condition (v) holds
since the data are $d$-neutral.  Hence $(f,\rho g)\in\cl\NA$; let $\rho\to1$.
The last sentence follows from Lemma~\ref{lem:II-R0}(a).
\end{proof}

\begin{remark}[Scope of Theorem~\ref{thm:II-R1}; non-sparse swallowing]
\label{rem:II-R1}
(a) In the model $v_l(s)\approx2^{-s}$ on $S_l$, $|a_s|\approx v_l(s)^{1+\beta_0}$
on $S_l\cap F$, (CS$_{\mathcal B}$) holds iff $\beta_0<1$; ratio profiles
$|a_s|/v_l(s)=1/s$ or $2^{-\sqrt s}$ are covered, whereas (H4-inf) needs
$\inf|a_s|/v_l(s)>0$.
(b) The deep set is where the method of Theorem~\ref{thm:II-R1} departs
from exact cushion bounds: every two-piece data with $b^+-b^-=X$ on $F$
must have $(s_jb^+_j)_->A_j$ or $(s_jb^-_j)_+>A_j$ on $D_t$, and fixed
data whose anti-sign part on $S_l\cap F$ is $|c|v_l$ pay the flip cost
$\mathrm{Fl}_{|c|v_l}(r)\ge r|c|m_{v_l}(r|c|/2)$ at scale $r$ (each $j$ with
$|a_j|<r|c|v_l(j)/2$ contributes at least $r|c|v_l(j)$), which is not
$o(r^2)$ for non-sparse profiles.  This is a statement about fixed data,
not a non-recovery statement.
(c) \emph{Heuristic, not used.}  For critical profiles a window
decomposition at scale $t$ can share the flips between the cushions of both
sides, while fixed data built from it carry them on one side; in the exactly
self-similar model ($|a_s|=v_l(s)^2$, gaps $1$) the resulting loss is the
factor $2$ in the flip coefficient, and averaging over a window does not
remove it.
(d) (\emph{Proved in our sources, not reproduced here.})  Under additional
hypotheses (bounded gaps, allowedness rule (c), $a$
of constant sign on $S_l\cap F$, $\supp u_l\subseteq F$, $d$-neutral bad
carriers, quasi-monotone ratios) super-critical swallowing is harmless for
signature-ladder operators with these features, by un-switching
(Lemma~\ref{lem:II-unswitch}) and Theorem~\ref{thm:II-raised}; and without
profile conditions one obtains recovery of the shrunken fibre $\rho_*\Cset(f)$
for an explicit $\rho_*<1$.  For $T_{\mathrm{final}}$ both statements are
contained in Theorems~\ref{thm:II-RSstar} and~\ref{thm:II-NDprime} below
(whose hypotheses (W$^\star$), (H2), (H3-inf), ($\mathcal B_{\rm fin}$) are
among those of the cited results; the rows of the first statement can fail
(ND$'$) and are then covered by Theorem~\ref{thm:II-NDprime}), and we do not
reproduce the proofs for general signature-ladder operators.  The non-sparse case for operators without a fast-decaying
diagonal base is open.
\end{remark}

\subsection{Box-dominated support}\label{subsec:II-N}

Here $T$ is an arbitrary admissible operator.  Fix an enumeration
$\mathfrak l:\N\times I\to\N$ of the carriers, put $\mathcal C_L:=
\{(k,m):\mathfrak l(k,m)\le L\}$,
$L(t):=\min\{L:\sum_{\mathfrak l(k,m)>L}\lambda_{k,m}\le t^3\}$, and
\[
 M_f(L):=\sum_{m\in I}\max\Big\{\frac1{\mu_{k,m}}:(k,m)\in\mathcal C_L,\
 k\in P_m\Big\}\qquad(\max\emptyset:=0),
\]
the inverse margins of the coarse peaks.

\begin{theorem}[Box-dominated support]\label{thm:II-N}
Let $T$ be admissible, $N\ge1$, and let $f\in S_{p_N^*}$ satisfy
\begin{enumerate}
\item[(BD)] $\mathrm{Bx}(j)\le C_a|a_j|$ for every $j\in\N$;
\item[(SC$_I$)] the set $I$ of all blocks satisfies \textup{(SC)} of
[I,~Definition~7.16];
\item[(W$_M$)] for every $n\in\N$, $\liminf_{T'\to0^+}T'M_f(L(2^{-n}T'))=0$.
\end{enumerate}
Then $f\in\Rec$.  Condition \textup{(BD)} forces $F=\N$; every carrier is
swallowed by the support, and no room, pinning, design feature, or
condition on the mates or on the strict non-peaks is used.
\end{theorem}

\begin{proof}
$F=\N$: for every $j$ some $u_{k,m}$ has $u_{k,m}(j)\ne0$, because
$\{u_{k,m}\}_k$ is dense in $S_{q^*}$; so $\mathrm{Bx}(j)>0$ and $a_j\ne0$.
(SC$_I$) excludes degenerate peaks, so $\mu_{k,m}>0$ at every peak.  Fix
$g\in\Cset(f)$, $\rho\in(0,1)$, and $\eta_0,\kappa_0,\varepsilon_{\rm tr},\eta$
as in the proof of Theorem~\ref{thm:II-Bstarinf}.  An $f$-constant depends
only on $f,g,\rho,N,T$.

\emph{Step 1 (exact data at one scale).}  Let $t\le\min(t_\eta,1)$, $L$ with
$\sum_{\mathfrak l(k,m)>L}\lambda_{k,m}\le t^3$, and $(B_\pm,\Theta_\pm)$ a
two-sided decomposition at scale $t$, with $d_\pm$, $\omega_\pm$ as in
[I,~Lemma~8.8].  Put $\mathcal Q_m:=Q_m\cap\{k:(k,m)\in\mathcal C_L\}$
(finite), $\omega^\pm_m:=\omega_{\pm,m}1_{\mathcal Q_m}$ and
$b^\pm_0:=g-\sum_mR_m^*(\omega^\pm_m-d_m(\omega^\pm_m)w_m)$.  Both pairs
represent $g$.  Drop $m$.
(a) At a peak $k$, $\sigma|\alpha(k)|=\lambda_k\mu_k$ [I,~(2)], so by
[I,~Lemma~8.8(c)] the coarse peaks contribute $\sum\lambda_k|\omega_\pm(k)|\le
(t/2)\max1/\mu_k$; outside $\mathcal C_L$, $|\omega_\pm|\le4/t$ gives
$\lambda_k|\omega_\pm(k)|\le4\lambda_k/t$.  Hence
$E_\pm:=\sum_m\sum_{k\notin\mathcal Q_m}\lambda_{k,m}|\omega_{\pm,m}(k)|\le
\frac t2M_f(L)+4t^2$.
(b) $|d_\pm-d(\omega^\pm)|\le t/\sigma+E_\pm/m$ ([I,~Lemma~8.8(d)] and
$|d(\omega)|\le\|D\omega\|_2\le\sum_k\Phi_k|\omega(k)|$).
(c) $b^\pm_0-B_\pm=\sum_mR_m^*(\omega_{\pm,m}1_{\mathcal Q_m^c})-\sum_m(d_{\pm,m}-
d_m(\omega^\pm_m))R_m^*w_m$, so $\|b^\pm_0-B_\pm\|_1\le K_1t$ with
$K_1:=C_f(1+M_f(L))$, and, if $t\le t_f$ and $t^2M_f(L)\le1$, pointwise
$|(b^\pm_0-B_\pm)_j|\le5\mathrm{Bx}(j)/t\le5C_a|a_j|/t$ (use
$|(R_m^*w_m)_j|\le\sum_k\lambda_{k,m}|u_{k,m}(j)|$ and $|d_\pm-d(\omega^\pm)|
\le1/t$).
(d) $X:=b^+_0-b^-_0=\sum_mR_m^*\big((\omega^-_m-\omega^+_m)-\Delta d^0_mw_m\big)$
with $\Delta d^0_m:=d_m(\omega^-_m)-d_m(\omega^+_m)$; since
$|\omega^\pm|\le4/t$ and $\|\Phi_m\|_2\le\frac12$, $|\Delta d^0_m|\le4/t$ and
$|X_j|\le12\mathrm{Bx}(j)/t\le12C_a|a_j|/t$.

\emph{Step 2 (common shift).}  Let $A_j:=(6C_a+1)|a_j|/t$, $x_j:=s_j
b^+_{0,j}$ and $y_j:=s_jb^-_{0,j}$.  By Lemma~\ref{lem:II-cushion} and (c),
$x_j\ge-A_j-\mathfrak f^+_j$ and $y_j\le A_j+\mathfrak f^-_j$; by (d),
$x_j-y_j\ge-2A_j$.  Let $\sigma_j$ be the point of $[y_j-A_j,x_j+A_j]$ closest
to $0$; by Lemma~\ref{lem:II-commonshift}(a), $|\sigma_j|\le\mathfrak f^+_j+
\mathfrak f^-_j$.  Put $\tilde\sigma:=\sum_js_j\sigma_je^*_j$
($\|\tilde\sigma\|_1\le t/(2q_0)$), $\kappa':=-\tilde\sigma(\zh)$,
$b^\pm:=b^\pm_0-\tilde\sigma-\kappa'a$ and $g_t:=g-\tilde\sigma-\kappa'a$.  Both
pairs $(b^\pm,\omega^\pm)$ represent $g_t$, $b^\pm(\xi)=0$, $b^+-b^-=X$, and for
$t^2\le q_0/(1+\|U\|)$ (so that $|\kappa'|\le1/t$), $(s_jb^+_j)_-\le
(6C_a+2)|a_j|/t$ and $(s_jb^-_j)_+\le(6C_a+2)|a_j|/t$.  Moreover
$p^*(g-g_t)\le(1+\|U\|)(\|\tilde\sigma\|_1+|\kappa'|)\le Kt$ with
$K:=(1+\|U\|)^2/q_0$.  As $F=\N$ there is no side condition off $F$: these
are two-piece data for $g_t$, with $\Delta d^0_m$ of arbitrary sign.

\emph{Step 3 ($\Gamma_w$).}  The pair $(b^+,\omega^+-d(\omega^+)w)$ differs from
$(B_+,\Theta_+)$ by a base part of $\ell_1$-norm $\le(K_1+K)t$ and the block
part $-\omega_+1_{\mathcal Q^c}+(d_+-d(\omega^+))w$, for which
$\sigma_mH_m\le\sigma_m(E_+/m)^2/C_m$.  As $\sqrt{\Gamma_w}$ is a seminorm and
$\Gamma_w(B_+,\Theta_+)\le1+\eta_\Gamma(\eta)$, $\sqrt{\Gamma_w(b^+,
\omega^+)}\le\sqrt{1+\eta_\Gamma(\eta)}+K_4t$ with $K_4:=C'_f(1+M_f(L))$; the
same for the $-$ data.

\emph{Step 4 (local validity).}  $\|B_\pm\|_1\le\|g\|_1+\sum_k\lambda_k
|\Theta_\pm(k)|\le\|g\|_1+6/t$, so $t\|b^\pm\|_1\le A_0:=\|g\|_1+8$ once
$(K_1+K)t^2\le1$.  For $k\in\mathcal Q_m$ and $\varsigma_k:=\sgn w_m(k)$,
[I,~Lemma~8.8(f)] gives $\varsigma_k\omega^+(k)\le\frac32\gap(k)/t$, and
$\varsigma_k\omega^+(k)\ge-4/t$; if $w_m(k)=0$ both signs are allowed and
$|\omega^+(k)|\le\frac32\gap(k)/t$.  So every block coordinate of the $+$ data
is of kind [1] or of kind [3] (side $+$) with $A_2:=4$, and symmetrically for
the $-$ data.  Lemma~\ref{lem:II-Uinf} (constant sequence, $U_*=0$,
$C_+:=6C_a+2$, $A_0$, $A_2=4$, $\gamma_B=1$) gives $f$-constants
$c_\flat,t_1$, independent of $L$, with $p^*(f+rg_t)\le1+\frac{r^2}2
(1+\eta_0)$ for $|r|\le c_\flat t$, provided $K_4t\le\kappa_0$ and $t\le t_1$.

\emph{Step 5 (windows of fixed length).}  Let $n_0\ge24\rho^2K/(c_\flat
(1-\rho^2))$ be an integer.  By (W$_M$) with $n=n_0$ there are $T_i\to0$ with
$T_iM_f(L_i)\to0$, $L_i:=L(2^{-n_0}T_i)$.  Use the scales
$T_i2^{1-s}$, $1\le s\le n_0$, all with the coarse level $L_i$.  For large
$i$ all conditions of Steps 1--4 hold at these scales.
Theorem~\ref{thm:II-masterinf} applies with the triples $(T_i,n_0,K)$,
$C_+:=6C_a+2$ and $\beta:=0$: (v) holds by (SC$_I$).
\end{proof}

\begin{remark}\label{rem:II-N}
(a) Under (BD) every switching, even of box size, is absorbed by the
cushions, so exact two-piece data exist at every scale without pinning;
what two-piece data cannot carry (peak usage, fine carriers, the
$w$-direction of non-$d$-neutral data) is moved into the base at second-order
cost controlled by the margins, and the closeness constant $K$ is an
$f$-constant: windows of fixed length suffice.  Critical mates are covered.
(b) For signature-ladder operators $L(t)\le l$ for $t\in\mathcal W(l)$, so
(W$_M$) follows from $M_f(l)T_{\rm hi}(l)\to0$ along a subsequence.
(c) In this method (BD) cannot be weakened to a sparse profile of
$\mathrm{Bx}$ on $F$: the box-size switching then exceeds the cushions by
$\approx\mathrm{Bx}(j)/t$ on a set whose flip cost at scales $r<t$ is of order
$r^2$ times a factor growing as $t\to0$.  The example ``$F=\N$, $a>0$'' is
covered when $a$ box-dominates; for the designed operator, thin $a$ and
infinitely many swallowed carriers are treated in
Subsection~\ref{subsec:II-RS} (Theorem~\ref{thm:II-RSinf}).
\end{remark}

\subsection{Raise transfer at the designed operator}\label{subsec:II-RS}

From now on $T=T_{\mathrm{final}}$ and $U=U_{\mathrm{final}}$, with diagonal
base $U^*e^*_s=\mu^\star_sk_s$ (Lemma~\ref{lem:II-base}).  For each level we
use one sub-window $w_l$ of level $l$ as the $l$-th window; all results of
the previous subsections hold for this signature-ladder operator.  We use
raise companions (Definition~\ref{def:II-companion}(b)), the raise-transfer
Lemma~\ref{lem:II-RT} with its extension to small corrections of arbitrary
sign on a finite set (Remark~\ref{rem:II-RTcorr}), Corollary~\ref{cor:II-RT}
and Theorem~\ref{thm:II-ERT}.  For a carrier $k$ we write $\gamma_k:=
\ip{U^*u_k}{U^*a}$.  With the diagonal base, the value change of a carrier
under a move of $a$ on $F$ depends only on $u_k|_F$:
$u_k(\zh')-u_k(\zh)=\ip{U^*u_k}{e'-e}$ when $z$ is fixed, and
$\ip{U^*u}{U^*x}=\sum_s(\mu^\star_s)^2u(s)x_s$.

\begin{lemma}[Value-preserving raises]\label{lem:II-VP}
Let $f\in S_{p^*}$ \textup{(}$F$ arbitrary\textup{)}, let $L_0$ be a finite set
of carriers, and assume
\[
 \textup{(ND$'_{L_0}$)}\qquad a|_F\notin\spn\{u_k|_F:\ k\in L_0\}.
\]
There are a finite set $F_0\subseteq F$ and $c_0,C_0>0$ such that for every
raise $\Delta a$ with $\supp\Delta a\cap F_0=\emptyset$ and $\|U^*\Delta a\|_H
\le c_0$ there is $\Delta a''\in\ell_1(F_0)$ with $\|\Delta a''\|_1\le C_0
\|U^*\Delta a\|_H$ and $|\Delta a''_s|\le|a_s|/2$, such that the first row
$f^\#:=f[(a+D)/q^*(a+D),z]$, $D:=\Delta a+\Delta a''$, satisfies
$u_k(\zh^\#)=u_k(\zh)$ for every $k\in L_0$.  One may take $C_0\ge1$, and
the radius can be chosen as $c_0\ge c_f\min\{1,\min_{s\in F_0}|a_s|\}/(1+
C_0)^2$ with $c_f>0$ depending only on $\nu$ (not on $L_0$); if $F_0=
\emptyset$, no correction is needed.
\end{lemma}

\begin{proof}
Let $V_0:=\{c\in\R^{L_0}:u_c|_F=0\}$, $u_c:=\sum_kc_ku_k$.
(i) If $x\in\ell_1$ is supported in $F$ and $c\in V_0$, then $\ip{U^*u_c}
{U^*x}=\sum_s(\mu^\star_s)^2u_c(s)x_s=0$.  Hence for every $A\in\ell_1(F)
\setminus\{0\}$ the vector $(\ip{U^*u_k}{U^*A/\|U^*A\|-e})_{k\in L_0}$ lies in
$V_0^\perp$.
(ii) Define $\Lambda:\ell_1(F)\to\R^{L_0}$ by $\Lambda(x)_k:=\sum_{s\in F}
(\mu^\star_s)^2x_s(u_k(s)-a_s\gamma_k/\nu^2)$.  A vector $c$ annihilates
$\Lambda(\ell_1(F))$ iff $u_c(s)=a_s(\sum_kc_k\gamma_k)/\nu^2$ for all $s\in F$,
i.e.\ iff $u_c|_F\in\R a|_F$; by (ND$'_{L_0}$) this means $u_c|_F=0$, i.e.\
$c\in V_0$.  So $\Lambda(\ell_1(F))=V_0^\perp$, and there is a finite
$F_0\subseteq F$ with $\Lambda(\R^{F_0})=V_0^\perp$; choose $F'_0\subseteq F_0$
with $|F'_0|=\dim V_0^\perp$ such that $\Lambda|_{\R^{F'_0}}$ is a bijection
onto $V_0^\perp$.
(iii) Let $Y_F$ be the closure of $U^*(\ell_1(F))$ in $H$; by (i),
$\ip{U^*u_c}h=0$ for $c\in V_0$ and $h\in Y_F$.  For $y\in Y_F$ and
$x\in\R^{F'_0}$ near $0$ put $e(y,x):=(U^*a+y+U^*x)/\|U^*a+y+U^*x\|$ and
$G(y,x)_k:=\ip{U^*u_k}{e(y,x)-e}$.  Since $e(y,x)-e\in Y_F$, $G$ is a smooth
map from a neighbourhood of $(0,0)$ in $Y_F\times\R^{F'_0}$ to $V_0^\perp$,
$G(0,0)=0$, and $D_xG(0,0)\xi=\ip{U^*u_k}{P^\perp U^*\xi}/\nu=
\Lambda(\xi)_k/\nu$, an isomorphism $\R^{F'_0}\to V_0^\perp$.  The implicit
function theorem gives a neighbourhood of $0$ in $Y_F$ and a smooth map
$y\mapsto x(y)$ with $G(y,x(y))=0$, $x(0)=0$ and $\|x(y)\|_1\le C_0\|y\|$.
For $\|U^*\Delta a\|\le c_0$ small, $|x_s|\le|a_s|/2$ on the finite set $F'_0$.
Put $\Delta a'':=x(U^*\Delta a)$.
(iv) \emph{Quantitative radius.}  On $\{\|v-U^*a\|\le\nu/2\}$ the map
$v\mapsto v/\|v\|$ has first derivatives bounded by $C/\nu$ and second
derivatives bounded by $C/\nu^2$; hence $\|D_yG\|\le C/\nu$, and $D_xG$ is
$C/\nu^2$-Lipschitz in $(y,x)$ for $\|y\|+\|U^*x\|\le\nu/4$.  Let $B$ be the
inverse of $D_xG(0,0)$ on $V_0^\perp$ and enlarge $C_0$ so that $C_0\ge
\max\{1,2\|B\|C/\nu\}$.  Then $x\mapsto x-BG(y,x)$ is a contraction of the
ball $\{\|x\|_1\le C_0\|y\|\}$ into itself as soon as $\|y\|(1+C_0)\le c\nu^2/
(1+\|B\|)$, and its fixed point $x(y)$ satisfies $|x_s|\le|a_s|/2$ on $F'_0$
as soon as $C_0\|y\|\le\min_{F_0}|a_s|/2$.  Since $\|B\|\le C_0\nu/C$, both
conditions hold for $\|y\|\le c_f\min\{1,\min_{F_0}|a_s|\}/(1+C_0)^2$ with
$c_f$ depending only on $\nu$.  Then $a+D$ has support $F$ and sign $z$
there, $e^\#=e(U^*\Delta a,\Delta a'')$ and $u_k(\zh^\#)-u_k(\zh)=
G(U^*\Delta a,\Delta a'')_k=0$.
\end{proof}

If (ND$'_{L_0}$) fails, some nontrivial combination $u_c$ has $u_c|_F=\kappa
a|_F$ with $\kappa\ne0$, and every move of $a$ on $F$ changes $\sum_kc_k
u_k(\zh)$ by exactly $-\kappa\nu\|e'-e\|^2/2$ (a second-order drift of one
sign); this case is treated in Lemma~\ref{lem:II-ND} below.

For $W\in\ell_1$, $W\ge0$, the \emph{raise room} of $W$ is
\begin{gather*}
 M^W(y):=\sum\{\mu^\star_sW_s:\ s\in F,\ |a_s|<yW_s\},\\
 \textup{(RR$_W$)}\quad M^W(y)=O(y^{1+\varepsilon})\ \text{for some }
 \varepsilon>0 .
\end{gather*}
(RR$_W$) holds, for instance, if $|a_s|\ge c(\mu^\star_s)^{1/(1+\varepsilon)}$
on $F$: then $|a_s|<yW_s$ forces $\mu^\star_s<(y\|W\|_\infty/c)^{1+\varepsilon}$.
Since $\log_2(1/\mu^\star_s)=2^{2^s}$, (RR$_W$) holds for every power profile
$|a_s|\approx v_l(s)^{1+\beta_0}$; it fails only for \emph{$\mu$-thin}
supports.  (RR$_W$) is not a cushion condition: it holds for critical and
super-critical profiles.

\begin{theorem}[Support swallowing with raise room]\label{thm:II-RS}
Let $T=T_{\mathrm{final}}$ and $N\ge1$.  Let $f\in S_{p_N^*}$ \textup{(}$F$
arbitrary\textup{)} satisfy \textup{(W$^\star$)}, \textup{(H2)},
\textup{(H3-inf)}, \textup{($\mathcal B_{\rm fin}$)}, and, with
$U_{\rm np}:=\sum_{l\in\mathcal B_{\rm np}}|u_l|$,
\textup{(ND$'_{\mathcal B_{\rm np}}$)} and \textup{(RR$_{U_{\rm np}}$)}.  Then
$f\in\Rec$.  No condition on the profile of $a$ on the swallowed signature
sets \textup{(}sparse, critical, super-critical or mixed\textup{)}, on the
$d$-neutrality of the bad carriers, or on their targets is imposed.
\end{theorem}

\begin{proof}
Fix $g\in\Cset(f)$, $\rho\in(0,1)$ and $\eta_0,\kappa_0,\varepsilon_{\rm tr},
\eta$ as in the proof of Theorem~\ref{thm:II-Bstarinf}, and let
$\theta_*:=(1-\rho^2)/(12\rho^2)$.  The idea is to run the proof of
Theorem~\ref{thm:II-R1} not at $f$ but at a raised companion $f_j$, where
the support dominates the bounded switching, so that the deep set is empty;
Theorem~\ref{thm:II-ERT} transfers the recovery back to $f$.  Let
$C_\tau$, $K_\star(l)$, $K^R_2(l)$, $c_\flat$, $t_1$ be the constants of the
proof of Theorem~\ref{thm:II-R1} (with $U_*=0$ in Step~4), computed at $f$ with
a factor $2$ of room; we show below that they serve at the companions.

\emph{Step 1 (windows).}  By (W$^\star$) choose levels $l_j\to\infty$ with
$\Lambda^\star_f(l_j)/(l_j2^{l_j^3}\Lambda^\circ(l_j))\to0$, and let $w_j$ be
the sub-window used at level $l_j$.  Put $K_j:=2(1+\|U\|)K^R_2(l_j)$
(a constant times $\Lambda^\star_f(l_j)$) and $n_j:=\lceil48\rho^2K_j/
(c_\flat(1-\rho^2))\rceil$ and $\theta_j:=\theta_*/l_j\le\theta_*$.  Choose
$m_j\in\N$ minimal with
\[
 C_RT_j^\varepsilon\log(e/T_j)\le\theta_jc_\flat^24^{-n_j},\qquad
 T_j:=T_{\rm hi}(w_j)2^{-m_j},
\]
where $\varepsilon$ is the exponent of (RR$_{U_{\rm np}}$) and $C_R$ the
constant of Step~2 below.  Then $m_j=O((n_j+\log(l_j/(\theta_*c_\flat^2)))/
\varepsilon)$, and $\log l_j=o(n(w_j))$; also
$n_j=O(\Lambda^\star_f(l_j))=o(n(w_j))$ by (W$^\star$) and (GW1); so for large
$j$ the scales $\mathcal S_j:=\{T_j2^{1-i}:1\le i\le n_j\}$ lie in
$\mathcal W(w_j)$, and $K_jT_j\to0$.

\emph{Step 2 (the companion).}  Put $y_j:=4C_\tau T_j$,
\[
 \Delta a_s:=s_s\big(y_jU_{\rm np}(s)-|a_s|\big)_+\quad(s\in F),
\]
let $\Delta a''$ be the correction of Lemma~\ref{lem:II-VP} for
$L_0:=\mathcal B_{\rm np}$, $D_j:=\Delta a+\Delta a''$, $\lambda_j:=q^*(a+D_j)$
and $f_j:=f[(a+D_j)/\lambda_j,z]$.  The raised set $\{s\in F:|a_s|<y_j
U_{\rm np}(s)\}$ leaves every finite subset of $F$ as $j\to\infty$, so it
misses $F_0$ for large $j$, and $\max_{s\in\supp\Delta a}\mu^\star_s\to0$.
Hence: $\|\Delta a\|_1\le y_j|\mathcal B|\to0$; $\|U^*\Delta a\|\le y_jM^{U_{\rm
np}}(y_j)\le Cy_j^{2+\varepsilon}$ and $\|U^*\Delta a\|=o(\|\Delta a\|_1)$;
$\lambda_j\ge1+\|\Delta a\|_1-\|\Delta a''\|_1-\|U^*D_j\|\ge1$ for large $j$;
$\lambda_j\to1$ and $f_j\to f$ (Lemma~\ref{lem:II-raise}(d); $D_j$ is not a
raise, but the proofs of Lemma~\ref{lem:II-raise}(b), (c), (e) use only that
$z$ is fixed, and (d) follows from $\|D_j\|_1\to0$ and $\lambda_j\to1$); and with
$\varepsilon_e:=\|U\|\|e_j-e\|\le2\|U\|\|U^*D_j\|/\nu\le Cy_j^{2+\varepsilon}$,
Corollary~\ref{cor:II-RT} (with the correction term of
Remark~\ref{rem:II-RTcorr}) and Lemma~\ref{lem:II-cost} give
\begin{gather*}
 p^*(f_j+r\rho g)\le1+\frac{s(\lambda_j\rho r)-1}{\lambda_j}+\varepsilon'_j,\\
 \varepsilon'_j:=2\|U^*D_j\|+2\|\Delta a''\|_1+p^*(L^*(w_j-w))\le C_R
 T_j^{2+\varepsilon}\log(e/T_j),
\end{gather*}
so $\varepsilon'_j\le\theta_jc_\flat^2(T_j2^{-n_j})^2\le\theta_*c_\flat^2
(T_j2^{-n_j})^2$ by the choice of $m_j$.
Moreover: (i) $|a^{(j)}_s|\ge y_jU_{\rm np}(s)/\lambda_j\ge2C_\tau T_jU_{\rm
np}(s)$ for $s\in F\setminus F_0$ (raised or not); (ii) $z$, $F$, $K$, $J$,
the rooms, $\mathcal B$, the signs $\epsilon_l$ and $T_0$ are those of $f$;
the values $u_l(\zh_j)=u_l(\zh)$, $l\in\mathcal B_{\rm np}$, are preserved
(Lemma~\ref{lem:II-VP}), and all block scalars move by $O(\varepsilon_e)$
(Lemma~\ref{lem:II-cost}(i)).  At a strict non-peak, [I,~Lemma~1.7] gives
$q_l=\epsilon_lq_0u_l(\zh)/\sigma_{m(l)}$; hence the $d$-rows of the
exact-switching cone at $f_j$ are positive multiples (one factor per block)
of those at $f$, and since the remaining rows depend only on $z$, $F$, $T_0$,
$\mathcal B$ and the peak statuses, the cone of
Lemma~\ref{lem:II-exactswitchinf} is the same set at $f_j$ and at $f$, with
Hoffman constants within a factor tending to $1$ (the $d$-rows are rescaled
by block factors $q^{(j)}_0\sigma_m/(q_0\sigma^{(j)}_m)\to1$).  The gaps of the bad strict non-peaks and the
margins of the non-degenerate bad peaks and of the peaks of (H2) are
$f$-constants, so these statuses persist for large $j$.

\emph{Step 3 (decompositions at $f_j$).}  Let $g_j:=g-(g(\xi_j)/q^{(j)}_0)
a^{(j)}$; then $g_j(\xi_j)=0$ and $p^*(g-g_j)=O(\varepsilon_e)$, since
$\|\xi_j-\xi\|_\infty=O(\varepsilon_e)$.  By the display of Step~2 with
$\rho=1$ and Corollary~\ref{cor:II-RT}, $p^*(f_j+tg_j)\le1+\frac{t^2}2
(1+\eta'_j)$ for $t\in[T_j2^{-n_j},T_j]$, with $\eta'_j\le(\lambda_j-1)
\lambda_j^2+2\theta_jc_\flat^2+O(\varepsilon_e/t)\to0$: here $\varepsilon'_j
\le\theta_jc_\flat^2t^2$ on these scales by the choice of $\theta_j$ in
Step~1, $\theta_j\to0$, $\lambda_j\to1$, and $\varepsilon_e/t\le Cy_j^{2+
\varepsilon}2^{n_j}/T_j\to0$.  The lemmas of
[I,~Section~8] used in the proof of Theorem~\ref{thm:II-R1} ([I,~Lemmas~8.5--8.8,
8.10, 8.24, 8.26--8.28, 8.38, 8.39]) use the mate only through the budget
$p^*(f'+tg')-1\le t^2/2$ and $g'(\xi')=0$; they hold at $f_j$ for $g_j$ on these
scales with constants changed by factors at most $2$.  The uniform
smallness [I,~Lemma~8.6] holds along $(f_j)$: its proof by contradiction
applies to sequences of pairs $(j_n,t_n)$ with $t_n\to0$, which force
$j_n\to\infty$, and $f_{j_n}\to f$.

\emph{Step 4 (exact data at $f_j$ without deep set).}  Run Steps 1--3 of the
proof of Theorem~\ref{thm:II-R1} at $f_j$ for $g_j$ at a scale $t\in\mathcal
S_j$.  The bounded-switching constant is uniform: the map $\Phi$ of
Lemma~\ref{lem:II-boundedfree}(b) depends continuously on $(e_j,\zh_j)$, so
$|\tau'_l|\le C_\tau$ for large $j$.  Since $\tau'=0$ on $\mathcal B_{\rm
pk}$, $|X_s|\le C_\tau U_{\rm np}(s)\le|a^{(j)}_s|/(2t)$ for $s\in F\setminus
F_0$ by (i) and $t\le T_j$; on the finite set $F_0$, $|a^{(j)}_s|\ge|a_s|/4$
and $|X_s|\le C_\tau|\mathcal B|\le|a_s|/(2T_j)$ for large $j$.  Hence the deep
set is empty at every scale of $\mathcal S_j$.  The bad carriers of
$\mathcal B_K$ at degenerate peaks with $\sgn w(k(l))=-\epsilon_l$ are handled
as follows: their row $\tau_l=0$ of the cone is kept; at $f_j$, $k(l)$ is
either a peak, and then $\tau_l\le\lambda_lK_dt$ by [I,~Lemma~8.39(c)], or a
strict non-peak of gap $\mathfrak g_j=O(\varepsilon_e)$, and then
[I,~Lemma~8.8(f)] and [I,~(16)] give $\tau_l\le\lambda_l(K_dt+3\mathfrak g_j/t)$
and $\lambda_l(|\omega_+(k)|+|\omega_-(k)|)\le|\tau_l|+\lambda_l(K_dt+6\mathfrak
g_j/t)$; in both cases $(\tau_l)_-\le\mathfrak c(\tau)/(2\mathfrak m_l)$.  As
$\mathfrak g_j=o(t^2)$ on $\mathcal S_j$, the violation of the row is
$O(K_\star t)$, and the window certificate vanishes at $k(l)$ (it vanishes at
peaks and at coordinates of gap $<t^2$).  Steps 1--3 of the proof of
Theorem~\ref{thm:II-R1} then give $d$-neutral two-piece data at $f_j$ for a
functional $g_{j,t}$ with $(s^{(j)}b^+)_-\le3|a^{(j)}|/t$ and $(s^{(j)}b^-)_+\le
3|a^{(j)}|/t$ on $F$, $p^*(g_j-g_{j,t})\le\frac12K_jt$, $\Gamma_w\le1+\eta_0/2$
(the budget factor $1+\eta'_j$ of Step~3 enters the bound of
[I,~Lemma~8.7(d)] multiplicatively: apply it to $g_j/(1+\eta'_j)^{1/2}$; this
gives $\Gamma_w\le(1+\eta'_j)(1+\eta_\Gamma(\eta))+o(1)$, which is $\le1+\eta_0/2$
for large $j$ because $\eta'_j\to0$),
and the block size conditions of kinds [1], [2] with uniform $A_2$,
$\gamma_B$.

\emph{Step 5 (local validity).}  Lemma~\ref{lem:II-Uinf} for the sequence
$f_j\to f$ with $U_*=0$, $C_+=3$ gives, with $c_\flat$, $t_1$, $j_0$
independent of $j$, $p^*(f_j+rg_{j,t})\le1+\frac{r^2}2(1+\eta_0)$ for
$|r|\le c_\flat t$, $t\in\mathcal S_j$, $j\ge j_0$.

\emph{Step 6.}  We verify Theorem~\ref{thm:II-ERT} with the rows $f_j$,
$\lambda_j$, $\varepsilon'_j$, $c_j:=c_\flat$, the windows $(T_j,n_j)$ and $K_j$:
(a) is Step~2; (b) holds with $p^*(g-g_{j,t})\le O(\varepsilon_e)+\frac12K_jt\le
K_jt$ and Step~5; (c), (d) hold by Steps 1--2; for (e), the averaged data
over $\mathcal S_j$ are $d$-neutral two-piece data at $f_j$ for $\bar g_j$ with
$\kappa_w\le1+\eta_0/2$ and side parts at most $3|a^{(j)}|/(T_j2^{1-n_j})$ on
$F$ (convexity), so they satisfy (CS-side) at $f_j$, and if $\rho\bar g_j\in
\Cset(f_j)$, Theorem~\ref{thm:II-raised} at $f_j$ (with $I_-=\emptyset$) applied
to $\rho\bar g_j$, whose data $\rho\bar D_j$ have $\kappa_w<1$, gives
$(f_j,\rho\bar g_j)\in\cl\NA$.  Hence $(f,\rho g)\in\cl\NA$; let $\rho\to1$.
\end{proof}

\begin{remark}\label{rem:II-RS}
(a) The $\ell_1$-mass of the raise is about $y_jm_{U_{\rm np}}(y_j)$, of
order $y_j^2$ for critical and much larger for super-critical profiles; for
these profiles it is not $o(T_j^2)$, so Theorem~\ref{thm:II-E} (which needs $p^*(f_j-f)=
o(T_{\rm lo}^2)$) does not apply.  Theorem~\ref{thm:II-ERT} needs only the
transfer error $\varepsilon'_j$, which involves the Hilbert footprint of the
raise and the normer change, but not its $\ell_1$-mass.
(b) At the raised row the deep set is empty.  (\emph{Heuristic:} the
factor-$2$ loss of Remark~\ref{rem:II-R1}(c), which we attribute to deep
assignment, therefore does not arise in this proof.)
(c) The theorem concerns $p_N$.  It does not by itself say anything about
individual first rows of Mart\'in's norm $p$: [I,~Lemma~1.5] transfers density
for $p_N$ (for infinitely many $N$) to $p$, not membership of a row in $\Rec$.
\end{remark}

The same raise recovers \emph{fixed} $d$-neutral data without any sparsity
condition.

\begin{theorem}[Fixed data at raised rows]\label{thm:II-A}
Let $T=T_{\mathrm{final}}$, $f\in S_{p^*}$ \textup{(}$F$ arbitrary\textup{)},
and let $g\in\Cset(f)$ carry $d$-neutral two-piece data $(b^\pm,\omega^\pm)$
with $\rho^2\kappa_w<1$.  Put $W:=(sb^+)_-+(sb^-)_+$ on $F$ and let $L_0$ be
the set of carriers $(k,m)$ with $k\in\supp\omega^+_m\cup\supp\omega^-_m$.  If
\textup{(RR$_W$)} and \textup{(ND$'_{L_0}$)} hold, then
$(f,\rho g)\in\cl\NA((c_0,p),\ell_2^2)$.
\end{theorem}

\begin{proof}
For $y\to0$ let $f_y$ be the raise of $f$ along $W$ at level $y$
($\Delta a_s:=s_s(yW_s-|a_s|)_+$) followed by the correction of
Lemma~\ref{lem:II-VP} for $L_0$; as in Step~2 above, $\lambda_y\ge1$,
$\lambda_y\to1$, $f_y\to f$, $\varepsilon'_y=O(y^{2+\varepsilon}\log(1/y))$,
and $|a^{(y)}_s|\ge yW_s/\lambda_y$ off the fixed finite set $F_0$.
(1) For $\omega$ supported in strict non-peaks of block $m$,
$d_m(\omega)=(mq_0/\sigma_m)\sum_k\Phi_m(k)u_{k,m}(\zh)\omega(k)$ by
[I,~Lemma~1.7]; the values of $L_0$ are preserved and their gaps are
$f$-constants, so at $f_y$ (for small $y$) $d^{(y)}_m=(q^{(y)}_0\sigma_m/(q_0
\sigma^{(y)}_m))d_m$ on these vectors.  Hence the data remain $d$-neutral, and
since $b^+-b^-=\sum_mR_m^*(\omega^-_m-\omega^+_m)$ does not involve the row,
both pairs represent at $f_y$ one functional $G_y$; replacing $b^\pm$ by
$b^\pm-\kappa_ya^{(y)}$, $\kappa_y:=G_y(\xi_y)/q^{(y)}_0$, we get two-piece data
at $f_y$ (same $z$, $F$, $K$) for a functional $g_y$ with $g_y(\xi_y)=0$ and
$p^*(g-g_y)=O(\varepsilon_e\log(e/\varepsilon_e))$ (the normer change), and
$\Gamma^{(y)}_w\to\Gamma_w$.
(2) Fix $\rho_1\in(\rho,1)$ with $\rho_1^2(\kappa_w+2\varepsilon_{\rm tr})<1$.
$(s^{(y)}b^+)_-,(s^{(y)}b^-)_+\le W\le\lambda_y|a^{(y)}|/y\le3|a^{(y)}|/t$ off
$F_0$ for $t\le T:=y/4$ (and on $F_0$ for small $y$, up to the shift
$\kappa_ya^{(y)}$).  Apply Lemma~\ref{lem:II-Uinf} at $f_y$ ($U_*=0$, $C_+=4$,
block coordinates of kind [2] with $A_2:=1$, $\gamma_B$ half the least gap)
to the scaled data $\rho_1(b^\pm-\kappa_ya^{(y)},\omega^\pm)$ of $\rho_1g_y$,
whose coefficient is $\rho_1^2\Gamma^{(y)}_w\le\rho_1^2(\kappa_w+o(1))<1\le2$,
so that hypothesis (OS1) of that lemma holds (the unscaled data may have
$\Gamma_w>2$ when $\rho<1/\sqrt2$).  It gives $p^*(f_y+r\rho_1g_y)\le1+\frac{r^2}2
(\rho_1^2\Gamma^{(y)}_w+\varepsilon_{\rm tr})$ on the respective sides for
$|r|\le c_\flat T$, uniformly in small $y$.
(3) For $|r|\le c_\flat T$, (2) gives $p^*(f_y+r\rho_1g_y)\le1+\frac{r^2}2
(1-\delta)\le s(r)$ for some $\delta>0$ and small $T$; for $|r|>
c_\flat T$, Corollary~\ref{cor:II-RT}, $\varepsilon'_y=o(T^2)$,
$p^*(g-g_y)=o(T)$ and [I,~Lemma~3.6(a)] give $p^*(f_y+r\rho_1g_y)\le s(r)$.
So $\rho_1g_y\in\Cset(f_y)$ for small $y$.
(4) Theorem~\ref{thm:II-raised} at $f_y$ for the mate $\rho_1g_y$ with data
$\rho_1(b^\pm-\kappa_ya^{(y)},\omega^\pm)$ (side parts $\le4|a^{(y)}|/T$, so
(CS-side) holds), with the factor $\rho/\rho_1$, gives $(f_y,\rho g_y)\in
\cl\NA$; and $(f_y,\rho g_y)\to(f,\rho g)$.
\end{proof}

Since every norm-attaining first row has finite base support
[I,~Proposition~1.10(c)], the approximants produced in
Theorems~\ref{thm:II-raised}, \ref{thm:II-RS} and \ref{thm:II-A} have finite
base supports; so for the pairs covered by these theorems the infinite-support
half of Lemma~Z holds (compare Proposition~\ref{prop:II-LSC} below).

\smallskip
The raise room (RR) is removed by a pinning lemma: an unraised support
coordinate that would be deep for the data pins the switching carrier
through the mate's own flip budget.  For a window index $l_*$ let
$T(l_*):=\bigcup_{l''\le l_*}\supp y_{l''}$ and $T_{\mathcal B}:=\bigcup_{l\in
\mathcal B}\supp y_l$.

\begin{lemma}[Pinning by an unraised deep coordinate]\label{lem:II-P}
Let $T$ be a signature-ladder operator, $f'\in S_{p^*}$ a first row with
forced data $(a',z',\dots)$, $F':=\supp a'$, and $g'\in X^*$ with
$g'(\xi')=0$ and $p^*(f'+rg')\le1+(1+\eta')r^2/2$ for $r=\pm t$, where
$\eta'\le1$.  Let $(B_\pm,\Theta_\pm)$ be a two-sided decomposition of $g'$
at $f'$ at scale $t\in\mathcal W(l_*)$, $l\le l_*$ a carrier, $\epsilon_l$ a
sign, $\tau_l:=-\epsilon_l\Delta\theta_l$, $\sigma:=\sgn\tau_l$, and let
$s\in S_l\cap F'$ with $s\notin T(l_*)$ and $\sgn(a'_s)\epsilon_l=-\sigma$.
Then
\[
 |\tau_l|v_l(s)\le\frac{2|a'_s|}t+\varphi_s+|r_s|,\qquad\varphi_s:=
 \mathfrak f^+_s+\mathfrak f^-_s,
\]
where $\sum_{s\in F'}\varphi_s\le(1+\eta')t/(2q'_0)$ and $r_s:=\sum_{l''>l_*,\,
s\in\supp y_{l''}}\Delta\theta_{l''}y_{l''}(s)/n_{l''}$ satisfies
$\sum_s|r_s|\le8t^2$.  Consequently, if $|\tau_l|v_l(s)\ge\frac{16}5|a'_s|/t$,
then $|\tau_l|\le\frac83(\varphi_s+|r_s|)/v_l(s)\le C_qt/v_l(s)$ with
$C_q:=\frac83(2/q'_0+8)$.  If $s\in(S_l\cap T(l_*)\cap F')\setminus
T_{\mathcal B}$, the same holds with $r_s$ replaced by $r_s+r'_s$,
$r'_s:=\sum\Delta\theta_{l''}y_{l''}(s)/n_{l''}$ over the good $l''\le l_*$
with $s\in\supp y_{l''}$; then $|r'_s|\le\frac43\sum_{\rm good}|\Delta
\theta_{l''}|$.
\end{lemma}

\begin{proof}
By (P1), for $s\in S_l$: $u_l(s)=v_l(s)$, $u_{l'}(s)=0$ for $l'<l$, and
$u_{l''}(s)=y_{l''}(s)/n_{l''}$ for $l''>l$; for $l<l''\le l_*$,
$s\in\supp y_{l''}$ would put $s$ into $T(l_*)$.  So by [I,~(15)],
$\Delta B(s)=\epsilon_l\tau_lv_l(s)-r_s$.  The proof of [I,~Lemma~8.28] uses
only the base budget; with the budget $(1+\eta')t^2/2$ it gives
$s_sB_+(s)\ge-|a'_s|/t-\mathfrak f^+_s$, $s_sB_-(s)\le|a'_s|/t+\mathfrak f^-_s$
($s_s:=\sgn a'_s$) and the stated bound for $\sum\varphi_s$; hence
$s_s\Delta B(s)\ge-2|a'_s|/t-\varphi_s$.  As $s_s\epsilon_l=-\sigma$,
$s_s\Delta B(s)=-|\tau_l|v_l(s)-s_sr_s$, which gives the first inequality.
By [I,~Lemma~8.10] and (P2), $\sum_s|r_s|\le\frac43\sum_{l''>l_*}6
\lambda_{l''}/t\le8T_{\rm lo}(l_*)^3/t\le8t^2$.  Under the deep hypothesis
$2|a'_s|/t\le\frac58|\tau_l|v_l(s)$, which gives the second inequality.
For $s\in T(l_*)\setminus T_{\mathcal B}$ the additional terms come from coarse
good targets containing $s$ ($\|y_{l''}\|_1/n_{l''}\le\frac43$).
\end{proof}

\begin{theorem}[Support swallowing of any profile]\label{thm:II-RSstar}
Let $T=T_{\mathrm{final}}$ and $N\ge1$.  Every $f\in S_{p_N^*}$ \textup{(}$F$
arbitrary\textup{)} satisfying \textup{(W$^\star$)}, \textup{(H2)},
\textup{(H3-inf)}, \textup{($\mathcal B_{\rm fin}$)} and
\textup{(ND$'_{\mathcal B_{\rm np}}$)} belongs to $\Rec$.  In particular
$\mu$-thin supports are allowed: no raise-room condition is needed.
\end{theorem}

\begin{proof}
We describe the changes to the proof of Theorem~\ref{thm:II-RS}.  Write
$\mathcal B':=\mathcal B_{\rm np}$, $U_{\rm np}=\sum_{l\in\mathcal B'}|u_l|$,
$\delta_{\mathcal B}:=\min_{l\in\mathcal B'}\delta_l$, and let $H\ge1$ be the
maximum, over $L'\subseteq\mathcal B'$, of the Hoffman constants (for the
violation measure of [I,~(18)] augmented by $\sum_{l\in L'}|\tau_l|$) of the
cones $\mathcal Z_f(L'):=\mathcal Z_f\cap\{\tau_l=0:l\in L'\}$, where
$\mathcal Z_f$ is the cone of Lemma~\ref{lem:II-exactswitchinf}; there are
finitely many, so $H$ is an $f$-constant.  Let $C'_fK_\star t$ bound the
violation of $\mathcal Z_f$ by the actual amplitudes (proof of
Lemma~\ref{lem:II-exactswitchinf}), $C_q:=\frac83(4/q_0+8)$, and
$A:=96(1+\|U\|)\rho^2/(c_\flat(1-\rho^2))$.

\emph{Step 1$'$ (thresholds).}  For a level $l$ and an integer $J$ put
\begin{gather*}
 K_1(l,J):=\tfrac54C_q\big[(1+2K_\star(l))2^{s_{\max}(l)}+2^J\big]/
 \delta_{\mathcal B},\\ K_{i+1}:=4H\big(C'_fK_\star+|\mathcal B'|K_i
 \big)+2K_i\ \ (1\le i\le|\mathcal B'|+1),
\end{gather*}
and $K'(l,J):=C_B(K_\star+K_{|\mathcal B'|+2})$, with an $f$-constant $C_B$
such that $K'(l,J)t$ bounds the closeness constants of Step~3 of the proof
of Theorem~\ref{thm:II-R1} when $C_fK_\star t$ is replaced there by
$K_{|\mathcal B'|+2}t$.  By induction, $K'(l,J)\le C'_B[(1+K_\star(l))
2^{s_{\max}(l)}+2^J]/\delta_{\mathcal B}$ with an $f$-constant $C'_B$: the
dependence on $2^J$ is additive.

\emph{Step 1$''$ (windows and depth).}  Take levels $l_j$ as in
Theorem~\ref{thm:II-RS}, the sub-window $w_j$ of level $l_j$, $T_j:=
T_{\rm hi}(w_j)$ and $L_j:=\log_2(1/T_j)$.  For $n\in\N$ put
$\theta_j:=\theta_*/l_j$, $x_j(n):=\theta_jc_\flat^2T_j4^{-n}/C_e$ ($C_e$ an
$f$-constant fixed in Step~2$'$) and $J_j(n):=\min\{J:\mu^\star_J\le
x_j(n)^2\}$.  By Lemma~\ref{lem:II-base}(d), $2^{J_j(n)}\le2\log_2(2\log_2(1/
x_j(n)))\le2\log_2(2(L_j+2n+\log_2l_j+C))$.  Let $n_j$ be the least $n\ge1$ with $n\ge
AK'(l_j,J_j(n))$; it exists because the right side grows at most
logarithmically in $n$; put $J_j:=J_j(n_j)$.  By Proposition~\ref{prop:II-RSwindow}(a),
$n(w_j)\ge l_j2^{l_j^3}\Lambda^\circ(l_j)\Omega(w_j)$ and $T_j\le2^{-l_j^3}/
(l_j\Omega(w_j)\Lambda^\circ(l_j))$, where $\Omega(w_j)\ge\Omega_K(l_j)\ge
2^{s_{\max}(l_j)}/\delta_{\mathcal B}$ for large $j$; moreover
$L_j\le l_j^3+l_j+Q(w_j)\le n(w_j)$ (Proposition~\ref{prop:II-subwindows}(f)),
$T_jn(w_j)\le2$, and $\log_2n(w_j)/n(w_j)\to0$
(Proposition~\ref{prop:II-RSwindow}(b)).  With (W$^\star$) this gives
$n_j=o(n(w_j))$, $K'(l_j,J_j)T_j\to0$ and $J_j\to\infty$.  So
$\mathcal S_j:=\{T_j2^{1-i}:1\le i\le n_j\}\subseteq\mathcal W(w_j)$.

\emph{Step 2$'$ (raise only the deep part).}  Let $y_j:=4C_\tau T_j$ and
\begin{gather*}
 R_j:=\Big\{s\in F\cap\bigcup_{l\in\mathcal B'}S_l:\ s\notin T_{\mathcal B},\
 s\ge J_j,\ |a_s|<y_jU_{\rm np}(s)\Big\},\\ \Delta a_s:=s_s\big(y_j
 U_{\rm np}(s)-|a_s|\big)_+1_{R_j}(s),
\end{gather*}
followed by the correction $\Delta a''$ of Lemma~\ref{lem:II-VP} for
$L_0:=\mathcal B'$ ($J_j>\max F_0$ for large $j$), and let $f_j$ be the
resulting row.  Then $\|U^*\Delta a\|\le\mu^\star_{J_j}\|\Delta a\|_1\le
\mu^\star_{J_j}y_j|\mathcal B|$, so $\lambda_j\ge1$ for large $j$, and with
$\mu:=\mu^\star_{J_j}\le x^2$, $x:=x_j(n_j)$, Lemma~\ref{lem:II-cost} and
Lemma~\ref{lem:II-VP} give $\varepsilon'_j\le C\mu y_j\log(e/(\mu y_j))$; since
$\mu\log(e/\mu)\le\sqrt\mu\le x$ and $x^2\log(1/y_j)\le x$ for large $j$,
$\varepsilon'_j\le2Cxy_j=8CC_\tau\theta_jc_\flat^2T_j^24^{-n_j}/C_e\le\theta_j
c_\flat^2(T_j2^{-n_j})^2$ once $C_e\ge8CC_\tau$; as in Step~3 of the proof of
Theorem~\ref{thm:II-RS}, the budget factor $1+\eta'_j$ then tends to $1$, and
the bound $\Gamma_w\le1+\eta_0/2$ of the window data holds for large $j$.  Further $\|\Delta a\|_1\le
y_j|\mathcal B|\to0$, $\lambda_j\to1$, $f_j\to f$, and Step~2(ii) of the
proof of Theorem~\ref{thm:II-RS} holds verbatim.

\emph{Step 4$'$ (exact data without deep set).}  Fix $t\in\mathcal S_j$ and a
two-sided decomposition of $g_j$ at $f_j$ at scale $t$, with actual
amplitudes $\tau_l$, $l\in\mathcal B'$.  The $|\mathcal B'|$ numbers $|\tau_l|$
miss one of the $|\mathcal B'|+1$ intervals $(K_it,K_{i+1}t]$; fix such an
$i_0$, let $L':=\{l:|\tau_l|\le K_{i_0}t\}$ (\emph{pinned}) and call the other
$l\in\mathcal B'$ \emph{kept} ($|\tau_l|>K_{i_0+1}t$).  Let $\tau'$ be a nearest
point of $\mathcal Z_f(L')$ to $\tau$ (the amplitudes of $\mathcal
B_{\rm pk}$ are treated as in Theorem~\ref{thm:II-RS}).  Then
$\sum|\tau_l-\tau'_l|\le H(C'_fK_\star+|\mathcal B'|K_{i_0})t\le K_{i_0+1}t/4$,
so $\tau'_l=0$ on $L'$, and for kept $l$: $\sgn\tau'_l=\sgn\tau_l$ and
$\frac34|\tau_l|\le|\tau'_l|\le\frac54|\tau_l|$; also $|\tau'_l|\le C_\tau$.
\emph{Claim}: $X':=\sum_{l\in\mathcal B'}\epsilon_l\tau'_lu_l$ satisfies
$s_jX'_j\ge-4|a^{(j)}_j|/t$ for all $j\in F$ (the deep set is empty).
(i) $j\in T_{\mathcal B}\cap F$ (finite): $|a^{(j)}_j|\ge|a_j|/4$, so
$4|a^{(j)}_j|/t\ge|a_j|/T_j\ge C_\tau|\mathcal B|\ge|X'_j|$ for large $j$.
(ii) $j\notin T_{\mathcal B}\cup\bigcup_{l\in\mathcal B'}S_l$: $X'_j=0$ by (P1).
(iii) $j\in S_l\setminus T_{\mathcal B}$ with $\tau'_l=0$: $X'_j=0$.
(iv) $j\in S_l\setminus T_{\mathcal B}$, $l$ kept: $X'_j=\epsilon_l\tau'_lv_l(j)$.
If $s_j\epsilon_l=\sgn\tau'_l$ there is nothing to prove; otherwise suppose
$|\tau'_l|v_l(j)>4|a^{(j)}_j|/t$.  If $j\ge J_j$, then $|a^{(j)}_j|\ge
y_jv_l(j)/2$ (whether or not $j\in R_j$), so $4|a^{(j)}_j|/t\ge2y_jv_l(j)/T_j=
8C_\tau v_l(j)>|\tau'_l|v_l(j)$, a contradiction.  If $j<J_j$, then
$|\tau_l|v_l(j)\ge\frac45|\tau'_l|v_l(j)>\frac{16}5|a^{(j)}_j|/t$, and
Lemma~\ref{lem:II-P} at $f'=f_j$ gives $|\tau_l|\le C_qt/v_l(j)\le\frac54C_q
2^{J_j}t/\delta_{\mathcal B}$ if $j\notin T(l_j)$, and (with $|r'_j|\le\frac43
K_\star t$ by Lemma~\ref{lem:II-exactswitchinf}(a) at $f_j$) $|\tau_l|\le
C_q(1+2K_\star)t/v_l(j)\le\frac54C_q(1+2K_\star)2^{s_{\max}(l_j)}t/
\delta_{\mathcal B}$ if $j\in T(l_j)\setminus T_{\mathcal B}$ (we used
$v_l(j)\ge\frac45\delta_{\mathcal B}2^{-j}$).  In both cases $|\tau_l|\le K_1t
<K_{i_0+1}t$, contradicting ``kept''.
Hence Steps 1--3 of the proof of Theorem~\ref{thm:II-R1}, run at $f_j$ with
$\tau'$ (closeness $K_{|\mathcal B'|+2}t$ instead of $C_fK_\star t$), give
$d$-neutral two-piece data at $f_j$ for a functional $g_{j,t}$ with
$(s^{(j)}b^+)_-,(s^{(j)}b^-)_+\le3|a^{(j)}|/t$ on $F$,
$p^*(g_j-g_{j,t})\le K'(l_j,J_j)t$ and $\Gamma_w\le(\sqrt{1+\eta_\Gamma(\eta)}+
K'_4t)^2$ with $K'_4=O(K'(l_j,J_j))$.

\emph{Steps 5, 6} are as in Theorem~\ref{thm:II-RS}, with
$K_j:=2(1+\|U\|)K'(l_j,J_j)$, so that $n_j\ge AK'(l_j,J_j)=48\rho^2K_j/
(c_\flat(1-\rho^2))$ by Step~1$''$.
\end{proof}

The pinning constant of Lemma~\ref{lem:II-P} at the depth $J$ is of order
$2^J$, and a raise beyond $J$ has footprint at most $\mu^\star_J$.  The
doubly exponential decay of $\mu^\star$ makes $2^{J}=O(\log\log(1/x))$ when
$\mu^\star_J\le x^2$, which is absorbed by every window; for
$\mu_s=2^{-s^2-1}$ or $\mu_s=2^{-2^s}$ the corresponding fixed point of
Step~1$''$ need not exist (Remark~\ref{rem:II-whybase}).

\begin{lemma}[Structure of the failure of (ND$'$)]\label{lem:II-ND}
Let $T$ be a signature-ladder operator, $f\in S_{p^*}$ \textup{(}$F$
arbitrary\textup{)}, $L_0$ a finite set of carriers, $T_0:=\bigcup_{k\in L_0}
\supp y_k$, and suppose that $\sum_{k\in L_0}c_ku_k|_F=\kappa a|_F$ with
$\kappa\ne0$.  Then $F\setminus T_0\subseteq\bigcup\{S_k:k\in L_0,\ c_k\ne0\}$
and $a_s=(c_k/\kappa)v_k(s)$ for $s\in S_k\cap F\setminus T_0$.
\end{lemma}

\begin{proof}
Let $s\in F\setminus T_0$.  If $s$ lies in no $S_k$, $k\in L_0$, then
$u_k(s)=0$ for all $k\in L_0$ (signatures of $L_0$ and targets of $L_0$ miss
$s$), so $\kappa a_s=0$, contradicting $s\in F$.  If $s\in S_k$, then
$u_{k'}(s)=0$ for $k'\in L_0\setminus\{k\}$ (disjoint signature sets, targets
in $T_0$) and $u_k(s)=v_k(s)$ by (P1); so $\kappa a_s=c_kv_k(s)$, and $a_s\ne0$
forces $c_k\ne0$.
\end{proof}

\begin{theorem}[Removal of (ND$'$)]\label{thm:II-NDprime}
Let $T=T_{\mathrm{final}}$ and $N\ge1$.  Every $f\in S_{p_N^*}$ \textup{(}$F$
arbitrary\textup{)} satisfying \textup{(W$^\star$)}, \textup{(H2)},
\textup{(H3-inf)} and \textup{($\mathcal B_{\rm fin}$)} belongs to $\Rec$.
\end{theorem}

\begin{proof}
If (ND$'_{\mathcal B_{\rm np}}$) holds, this is Theorem~\ref{thm:II-RSstar}.
Otherwise Lemma~\ref{lem:II-ND} with $L_0:=\mathcal B_{\rm np}$ gives
$F\setminus T_{\mathcal B}\subseteq\bigcup_{c_k\ne0}S_k$ with
$|a_s|=|c_k/\kappa|v_k(s)$ for $s\in S_k\cap F\setminus T_{\mathcal B}$.  For
$s\in F\setminus T_{\mathcal B}$ and $l\in\mathcal B$, $u_l(s)=0$ unless
$s\in S_l$ (bad targets lie in $T_{\mathcal B}$); and $s$ lies in exactly one
$S_k$, so $|u_l(s)|/|a_s|\le\max_{c_k\ne0}|\kappa/c_k|$.  On the finite set
$F\cap T_{\mathcal B}$ the ratios are finite.  So (H4-inf) holds, and
Theorem~\ref{thm:II-R1} gives $f\in\Rec$, with no raise.
\end{proof}

Thus (ND$'$) fails only when the support of $a$ is carried, up to finitely
many coordinates, by the signature sets of the kernel carriers with profiles
proportional to the signatures: then every bad-carrier switching on $F$ is
cushion-dominated and no raise is needed.  Together,
Theorems~\ref{thm:II-RSstar} and~\ref{thm:II-NDprime} settle infinite base
support with finitely many bad carriers for $T_{\mathrm{final}}$, under the
same hypotheses (W$^\star$), (H2), (H3-inf) as at finite $F$.

\smallskip
We turn to infinitely many bad carriers.  Then bounded switching
(Lemma~\ref{lem:II-boundedfree}) is not available, and the switching on $F$
is of box size.  For a level $l_*$ put $\mathcal B(l_*):=\mathcal B\cap
[1,l_*]$, $\mathcal B_{\rm np}(l_*)$, $\mathcal B_{\rm pk}(l_*)$ likewise,
$T_{\mathcal B}(l_*):=\bigcup_{l\in\mathcal B(l_*)}\supp y_l\subseteq
[1,s_{\max}(l_*)]$, and $\mathrm{Bx}_*(s):=\sum_{l\in\mathcal B_{\rm np}(l_*)}
\lambda_l|u_l(s)|$.  The rooms $\mathfrak r^\star_l$ of good $l\le l_*$ are
computed on $S_l\setminus(F\cup\bigcup_{l'\in\mathcal B(l_*),\,l'>l}\supp y_{l'})$,
and $K_\star(l_*):=(1/q_0+22)\Lambda^\star_f(l_*)$ with these rooms.  Further
\begin{gather*}
 \gamma_{\mathcal B}(l_*):=\min\big(\min_mM_m,\ \min\{\gap_{m(l)}(k(l)):l\in
 \mathcal B_{\rm np}(l_*)\}\big),\\
 \mu_{\mathcal B}(l_*):=\min\{\mu_{k(l),m(l)}:
 l\in\mathcal B_{\rm pk}(l_*),\ \alpha_{m(l)}(k(l))\ne0\}
\end{gather*}
($\mu_{\mathcal B}(l_*):=1$ if there is no such $l$).  Let $\mathcal Z(l_*)
\subseteq\R^{\mathcal B(l_*)}$ be the cone of Lemma~\ref{lem:II-exactswitchinf}
built from $\mathcal B(l_*)$, with $T_0:=T_{\mathcal B}(l_*)\cup\bigcup_{l\in
\mathcal B_F(l_*)}(S_l\setminus F)$, and for $P\subseteq\mathcal B_{\rm np}(l_*)$
and $\Sigma\subseteq T_{\mathcal B}(l_*)\cap F$ let
\[
 \mathcal Z(P,\Sigma):=\mathcal Z(l_*)\cap\{\tau_l=0:l\in P\}\cap\Big\{s_j
 \sum_{l\in\mathcal B(l_*)}\epsilon_l\tau_lu_l(j)\ge0:\ j\in\Sigma\Big\}.
\]
Let $H(l_*)\ge1$ be a constant such that for all $P$, $\Sigma$ and
$\tau\in\R^{\mathcal B(l_*)}$, $\dist_1(\tau,\mathcal Z(P,\Sigma))$ is at most
$H(l_*)$ times the sum of: the cost $\mathfrak c(\tau)$ of
Lemma~\ref{lem:II-exactswitchinf}, $\sum_{l\in P\cup\mathcal B_{\rm pk}(l_*)}
|\tau_l|$, the moduli of the $d$-sums, and the negative parts of the
$\Sigma$-rows; it exists by Hoffman's theorem (finitely many systems; the
conversion of $\mathfrak c$ into violations involves the factors
$1/(2\mathfrak m_l)$ and $1/(1-|z_j|)$, $j\in T_0\setminus(F\cup K)$).
Let $\psi(J):=4\sum_{l\in\mathcal B}\lambda_l\|u_l1_{[J,\infty)}\|_1$, so that
$\psi(J)\to0$ as $J\to\infty$.

\begin{lemma}[Window data at a deep-raised row]\label{lem:II-Winf}
Let $T=T_{\mathrm{final}}$, $f\in S_{p_N^*}$ \textup{(}$F$ arbitrary,
$\mathcal B$ possibly infinite\textup{)} satisfy \textup{(H2)} and
\textup{(H3-inf)}, let $g\in\Cset(f)$, and let $\eta_0,\kappa_0,
\varepsilon_{\rm tr},\eta$ be as in the proof of Theorem~\ref{thm:II-Bstarinf}.
There are $f$-constants $C'_f,C''_f,K_4$ and a level $l_f$ such that the
following holds for $l_*\ge l_f$ with $\mathfrak r^\star_l>0$ for all good
$l\le l_*$ and \textup{(ND$'_{\mathcal B_{\rm np}(l_*)}$)}.  Let $J$ be an
integer with $J>\max F_0$ \textup{(}$F_0$ of Lemma~\textup{\ref{lem:II-VP}} for
$L_0:=\mathcal B_{\rm np}(l_*)$\textup{)}, let
\[
 R:=\{s\in F:\ s\ge J,\ |a_s|<4\mathrm{Bx}_*(s)\},\qquad\Delta a:=\sum_{s\in
 R}s_s\big(4\mathrm{Bx}_*(s)-|a_s|\big)e^*_s
\]
\textup{(}a deep raise at box level, of mass $\le\psi(J)$ and footprint
$\le\mu^\star_J\psi(J)$\textup{)}, let $\Delta a''$ be the correction of
Lemma~\textup{\ref{lem:II-VP}}, and $f':=f[(a+D)/q^*(a+D),z]$,
$D:=\Delta a+\Delta a''$, with $\lambda':=q^*(a+D)\le2$.  Put
\begin{gather*}
 K^{\rm pk}:=C'_f(1+K_\star(l_*))\big(1+l_*/\mu_{\mathcal B}(l_*)\big),\\
 K_1:=\tfrac54C_q\big[(1+2K_\star(l_*))2^{s_{\max}(l_*)}+2^J\big]/
 \delta_{\min}(l_*)+4+4K^{\rm pk},\\
 K_{i+1}:=4H(l_*)\big(K^{\rm pk}+l_*K_i+s_{\max}(l_*)K_1\big)+2K_i\ \
 (1\le i\le l_*+1),\\ K_{\rm top}:=K_{l_*+2}.
\end{gather*}
Let $t\in\mathcal W(w_{l_*})$ be a dyadic scale with $t\le\min(t_\eta,1)$,
and let $g'$ be a functional with $g'(\xi')=0$ and $p^*(f'+rg')\le
1+(1+\eta')r^2/2$ at $r=\pm t$, $\eta'\le1$.  Assume:
\begin{enumerate}
\item[(S1)] every $j\in T_{\mathcal B}(l_*)\cap F$ with $j<J$ has $|a_j|\le
t^2$ or $|a_j|\ge K_{\rm top}t^2$;
\item[(S2)] $K_{\rm top}t\le\min(1,\kappa_0/K_4)$;
\item[(S3)] at $f'$, the bad strict non-peaks $l\le l_*$ have gaps
$\ge\gamma_{\mathcal B}(l_*)/2$, and the non-degenerate bad peaks $l\le l_*$
are peaks with margins $\ge\mu_{\mathcal B}(l_*)/2$.
\end{enumerate}
Then every two-sided decomposition of $g'$ at $f'$ at scale $t$ yields
$d$-neutral two-piece data $(b^\pm,\omega^\pm)$ at $f'$ for a functional
$g'_t$ with $(s'b^+)_-\le3|a'|/t$ and $(s'b^-)_+\le3|a'|/t$ on $F$,
$p^*(g'-g'_t)\le C''_fK_{\rm top}t$, $\Gamma_w(b^\pm,\omega^\pm)\le
(\sqrt{(1+\eta')(1+\eta_\Gamma(\eta))}+C''_fK_{\rm top}t)^2$, $t\|b^\pm\|_1\le9$, and
block coordinates of kind \textup{[1]} or of kind \textup{[2]} with
$A_2:=12$, $\gamma_B:=\gamma_{\mathcal B}(l_*)/2$.
\end{lemma}

\begin{proof}
(1) \emph{Pinning.}  Lemma~\ref{lem:II-exactswitchinf}(a) at $f'$ at level
$l_*$ with $\mathcal B(l_*)$ in place of $\mathcal B$: the fine bad carriers
$l>l_*$ enter the good inequalities only through their targets, with total
contribution $\le8t^2$ ([I,~Lemma~8.10], (P2)); so $\sum_{\rm good}|\Delta
\theta_l|\le K_\star t$ (up to a factor $2$ for the budget at $f'$).  For a
non-degenerate bad peak $k=k(l)$, $l\le l_*$, [I,~(16)] at $f'$ with
$\Delta\Theta(k)=-\epsilon_l\tau_l/\lambda_l$ gives $\sgn(w'(k))\epsilon_l\tau_l/
\lambda_l=|\omega_+(k)|+|\omega_-(k)|+\Delta d\,M'$, where $0\le|\omega_+(k)|+
|\omega_-(k)|\le t/(\sigma'|\alpha'(k)|)=t/(\lambda_l\mu'_k)$ by
[I,~Lemma~8.8(c)] and [I,~(2)], and $|\Delta d|M'\le K_dt$
([I,~Lemma~8.39(a)] at $f'$); by (S3), $|\tau_l|\le2t/\mu_{\mathcal B}(l_*)+
\lambda_lK_dt$.  Degenerate bad peaks are handled by (H3-inf) as in
Lemma~\ref{lem:II-exactswitchinf}.  Hence the violation of $\mathcal Z(l_*)$ by
the actual amplitudes $\tau=(\tau_l)_{l\in\mathcal B(l_*)}$ is at most
$K^{\rm pk}t$ (the fine bad $d$-terms are $O(t^2)$ and the fine bad targets on
$T_0$ contribute $\le8t^2$).
(2) \emph{Pigeonhole and projection.}  The at most $l_*$ numbers $|\tau_l|$,
$l\in\mathcal B_{\rm np}(l_*)$, miss one of the $l_*+1$ intervals
$(K_it,K_{i+1}t]$; let $i_0$ be such an index, $P:=\{l:|\tau_l|\le K_{i_0}t\}$
(pinned), and call the others kept.  Let $\Sigma:=\{j\in T_{\mathcal B}(l_*)\cap
F:j<J,\ |a_j|\le t^2\}$.  For $j\in\Sigma$, $j\notin R$ and $|a'_j|\le\frac32
|a_j|$, so [I,~Lemma~8.28] at $f'$ and [I,~(15)] give $s_j\sum_{\mathcal B(l_*)}
\epsilon_l\tau_lu_l(j)\ge-3|a_j|/t-\varphi_j-|r_j|\ge-(3+K^{\rm pk})t$, where
$r_j$ collects the good coarse, fine and bad-peak carriers at $j$.  So the
$\Sigma$-rows are violated by at most $s_{\max}(l_*)K_1t$ in total.  Let $\tau'$
be a nearest point of $\mathcal Z(P,\Sigma)$ to $\tau$: $\sum|\tau-\tau'|\le
H(l_*)(K^{\rm pk}+l_*K_{i_0}+s_{\max}(l_*)K_1)t\le K_{i_0+1}t/4$.  Hence
$\tau'=0$ on $P$, and for kept $l$, $\sgn\tau'_l=\sgn\tau_l$ and
$\frac34|\tau_l|\le|\tau'_l|\le\frac54|\tau_l|\le\frac{15}2\lambda_l/t$
([I,~Lemma~8.10]).
(3) \emph{No deep set.}  Let $X':=\sum_{l\in\mathcal B_{\rm np}(l_*)}\epsilon_l
\tau'_lu_l$; then $|X'_j|\le\frac{15}2\mathrm{Bx}_*(j)/t$.  We claim
$s_jX'_j\ge-4|a'_j|/t$ for all $j\in F$.  If $j\ge J$: $|a'_j|\ge2\mathrm{Bx}_*(j)$
(whether or not $j\in R$, as $\lambda'\le2$), so $4|a'_j|/t>|X'_j|$.  If $j<J$
and $j\in S_l\setminus T_{\mathcal B}(l_*)$ for some $l\in\mathcal B_{\rm np}
(l_*)$: $X'_j=\epsilon_l\tau'_lv_l(j)$, which vanishes if $l$ is pinned; if $l$
is kept and $s_jX'_j<-4|a'_j|/t$, then Lemma~\ref{lem:II-P} at $f'$ applies
to the actual $\tau_l$ and gives $|\tau_l|\le K_1t$, contradicting ``kept''.
If $j<J$ and $j\in T_{\mathcal B}(l_*)\cap F$: for $j\in\Sigma$ the $\Sigma$-row
gives $s_jX'_j\ge0$; otherwise $|a_j|\ge K_{\rm top}t^2$ by (S1), and
$s_jX_j(\tau)\ge-2|a'_j|/t-K^{\rm pk}t$, $|X'_j-X_j(\tau)|\le K_{\rm top}t/4$
and $2|a'_j|/t\ge K_{\rm top}t/2\ge(K^{\rm pk}+K_{\rm top}/4)t$ (as
$|a'_j|\ge|a_j|/4$ and $K_{\rm top}\ge K_1\ge4K^{\rm pk}$), so
$s_jX'_j\ge-4|a'_j|/t$.  In all other cases $X'_j=0$.
(4) \emph{Data.}  Steps 2--3 of the proof of Theorem~\ref{thm:II-R1} at $f'$
with $\tau'$ (empty deep set, closeness $K_{\rm top}t$ instead of
$C_fK_\star t$; bounded switching is not used, only $|\tau'_l|\le\frac{15}2
\lambda_l/t$) give the data; the bound for $\Gamma_w$ is [I,~Lemma~8.7(d)]
applied to $g'/(1+\eta')^{1/2}$, whose budget at $r=\pm t$ is $r^2/2$, plus
the closeness term.  Since $\sum_FA_j\le2/t$ and $\|X'\|_1\le\frac{15}2
\sum_l\lambda_l/t\le3/t$, $t\|b^\pm\|_1\le9$; at $k=k(l)$, $|\omega^-(k)|\le
|\omega^+(k)|+|\tau'_l|/\lambda_l\le(4+\frac{15}2)/t$; and (S3) gives the gap
bound for kind [2].
\end{proof}

For a level $l$ let $C_{\rm VP}(l):=C_0\,2^{\max F_0}/\min_{s\in F_0}|a_s|$,
with $F_0$, $C_0\ge1$ the data of Lemma~\ref{lem:II-VP} for
$L_0:=\mathcal B_{\rm np}(l)$, and $C_{\rm VP}(l):=1$ if $F_0=\emptyset$; so
$C_{\rm VP}(l)\ge2^{\max F_0}\ge1$ (as $|a_s|\le1$).  We may and do assume
$\gamma_{\mathcal B}(l),\mu_{\mathcal B}(l)\le1$ (replace them by their minimum
with $1$).  Put
\[
 \Xi_f(l):=\frac{(4H(l)l+3)^{l+2}\,(1+K_\star(l))\,(1+C_{\rm VP}(l))}
 {\gamma_{\mathcal B}(l)\,\mu_{\mathcal B}(l)} .
\]

\begin{theorem}[Infinitely many bad carriers]\label{thm:II-RSinf}
Let $T=T_{\mathrm{final}}$ and $N\ge1$.  Let $f\in S_{p_N^*}$ \textup{(}$F$
arbitrary, $\mathcal B$ possibly infinite\textup{)} satisfy \textup{(H2)},
\textup{(H3-inf)}, $\mathfrak r^\star_l>0$ for every good $l$,
\textup{(ND$'_{\mathcal B_{\rm np}(l)}$)} for every $l$, and
\[
 \textup{(W$_{\rm inf}$)}\qquad\liminf_{l\to\infty}\frac{\Xi_f(l)}
 {l2^{l^3}\Lambda^\circ(l)}=0 .
\]
Then $f\in\Rec$.  For finite $\mathcal B$ the rates $H$, $C_{\rm VP}$,
$\gamma_{\mathcal B}$, $\mu_{\mathcal B}$ are $f$-constants for large $l$, and
\textup{(W$_{\rm inf}$)} implies \textup{(W$^\star$)} \textup{(}because
$\Xi_f(l)\ge1+K_\star(l)\ge\Lambda^\star_f(l)$\textup{)}; finite $\mathcal B$ is
covered under \textup{(W$^\star$)} alone by Theorems~\textup{\ref{thm:II-RSstar}}
and~\textup{\ref{thm:II-NDprime}}.
\end{theorem}

\begin{proof}
We follow the proof of Theorem~\ref{thm:II-RSstar}, with
Lemma~\ref{lem:II-Winf} in place of Step~4$'$.
(i) \emph{Levels.}  By (W$_{\rm inf}$) choose $l_j$ with
$\Xi_f(l_j)/(l_j2^{l_j^3}\Lambda^\circ(l_j))\to0$, and let $w_j$ be the
sub-window of level $l_j$, $L_j:=\log_2(1/T_{\rm hi}(w_j))$.
(ii) \emph{Local validity.}  Lemma~\ref{lem:II-Uinf} with $U_*=0$, $C_+=3$,
$A_0=9$ gives $c_0,t_1$ depending only on $f$; with $A_2=12$ and
$\gamma_B=\gamma_{\mathcal B}(l_j)/2$, $c_\flat(l_j)=c_0\gamma_{\mathcal B}(l_j)/24$.
(iii) \emph{Sub-window position.}  For $n,p\in\N$ consider the scales
$\mathcal S(n,p):=\{T_{\rm hi}(w_j)2^{-p}2^{1-i}:1\le i\le n\}$.  For an integer
$J>\max F_0$ write $K_{\rm top}(J)$ for the constant of
Lemma~\ref{lem:II-Winf} with $l_*=l_j$ and this $J$.  Condition (S1) with
$K_{\rm top}(J)$ at all scales of $\mathcal S(n,p)$ excludes, for each of the
at most $s_{\max}(l_j)$ coordinates $j'\in T_{\mathcal B}(l_j)\cap F$, at most
$n+\log_2K_{\rm top}(J)+2$ values of $p$; so some $p\le P(n,J):=(s_{\max}(l_j)
+2)(n+\log_2K_{\rm top}(J)+2)$ satisfies it at all scales of $\mathcal S(n,p)$.
(iv) \emph{Depth.}  Put $\theta_j:=\theta_*/l_j$ and, for $(n,J)$, in this
order: $K_{\rm top}(J)$, $P(n,J)$, $x(n,J):=\theta_jc_\flat(l_j)^2\gamma_{\mathcal
B}(l_j)\mu_{\mathcal B}(l_j)2^{-2(L_j+P(n,J)+n)}/(C_e(1+C_{\rm VP}(l_j)))$.  For
each $n$ let $J(n)$ be the least $J\ge1+\max F_0$ (or $J\ge1$ if $F_0=\emptyset$)
with $\mu^\star_J\le x(n,J)^2$.  It exists: $\log_2(1/\mu^\star_J)=2^{2^J}$,
while $\log_2(1/x(n,J))$ grows at most linearly in $J$ (since $K_1$, hence
$K_{\rm top}(J)$, depends on $J$ through the summand $2^J$, $\log_2K_{\rm
top}(J)\le J+C_j$, and so $P(n,J)$ is at most linear in $J$).  Moreover
$2^{J(n)}\le\max\{2^{1+\max F_0},\,C\log_2(L_j+n+\log_2 l_j+C_j)\}$ with
$C_j$ independent of $n$, and $2^{1+\max F_0}\le2C_{\rm VP}(l_j)$
(Lemma~\ref{lem:II-base}(d)).  Put $P:=P(n,J(n))$.  The raise of
Lemma~\ref{lem:II-Winf} at depth $J(n)$ has footprint $\le C\mu^\star_{J(n)}\le
Cx(n,J(n))^2$, which is at most the radius $c_f/(1+C_{\rm VP}(l_j))^2$ of
Lemma~\ref{lem:II-VP} for $C_e$ large; with Lemma~\ref{lem:II-VP} and
Lemma~\ref{lem:II-cost} its transfer error is
$\varepsilon'\le C(1+C_{\rm VP})\sqrt{\mu^\star_{J(n)}}\le C(1+C_{\rm VP})x(n,J(n))
\le\theta_jc_\flat^2(T_{\rm lo}')^2$, where $T'_{\rm lo}\ge T_{\rm hi}(w_j)
2^{-P-n}$ is the least scale, once $C_e$ is a large $f$-constant; and the
values, block scalars, gaps and margins move by $\le C\sqrt{\mu^\star_{J(n)}}\le
\frac12\min(\gamma_{\mathcal B},\mu_{\mathcal B})$, which gives (S3) (the values
of $\mathcal B_{\rm np}(l_j)$ are preserved, and the peak margins are
Lipschitz in the values and the block scalars by Lemma~\ref{lem:II-T2}).
Let $n_j$ be the least $n$ with $n\ge AK_{\rm top}(J(n))/c_\flat(l_j)$,
$A:=96(1+\|U\|)C''_f\rho^2/(1-\rho^2)$; it exists because the right side
depends on $n$ only through $2^{J(n)}$, which grows at most logarithmically
in $n$ (this is the fixed point $J\to K_{\rm top}\to n\to P\to x\to J$).
Unrolling $K_{i+1}\le(4H(l_j)l_j+2)K_i+4H(l_j)(K^{\rm pk}+s_{\max}(l_j)K_1)$
gives
\[
 K_{\rm top}\le(1+s_{\max}(l_j))(4H(l_j)l_j+3)^{l_j+1}(K_1+K^{\rm pk})
\]
(the cruder bound without the factor $1+s_{\max}(l_j)$ is false in general,
because of the term $s_{\max}K_1$ of the recursion).  Using this,
$(1+s_{\max})2^{s_{\max}}/\delta_{\min}\le2\Omega_K(l_j)$, Proposition~\ref{prop:II-RSwindow}
and Proposition~\ref{prop:II-subwindows}(f), one obtains
\[
 n_j+P_j\le C\,\Xi_f(l_j)\,\Omega_K(l_j)\,s_{\max}(l_j)\log_2 n(w_j),\qquad
 K_{\rm top}T_{\rm hi}(w_j)\to0,
\]
and, since $l2^{l^3}\Lambda^\circ(l)\Omega_K(l)s_{\max}(l)\log_2n(w)\le
\mathrm{Des}(l)^{1/2}\cdot2\log_2Q(w)\le Q(w)$ and $n(w)\ge l2^{l^3}Q(w)$,
$(n_j+P_j)/n(w_j)\to0$ by (W$_{\rm inf}$).  So the chosen scales lie in
$\mathcal W(w_j)$, and (S1), (S2), (S3) hold for large $j$.
(v) \emph{Conclusion.}  The raise mass is $\le\psi(J(n_j))\to0$, so
$\lambda_j\to1$ and $f_j\to f$; the budget hypothesis of Lemma~\ref{lem:II-Winf}
holds by Corollary~\ref{cor:II-RT}, as in Step~3 of
Theorem~\ref{thm:II-RS}, with $\eta'_j\to0$ by the choice of $\theta_j$.
Lemma~\ref{lem:II-Winf} at $f_j$ for $g_j$ gives exact
$d$-neutral data at every scale with cushion-compatible side parts, and since
$\eta'_j\to0$ and $K_{\rm top}T_{\rm hi}(w_j)\to0$ these data satisfy
$\Gamma_w\le1+\eta_0/2$ on the chosen scales for large $j$;
Lemma~\ref{lem:II-Uinf} gives the pieces, and Theorem~\ref{thm:II-ERT} together
with Theorem~\ref{thm:II-raised} at $f_j$ concludes as in
Theorem~\ref{thm:II-RS}.
\end{proof}

\begin{remark}[The rates in $\Xi_f$]\label{rem:II-RSinf}
$K_\star(l)$ contains the rooms of the good carriers, as in (W$^\star$);
$H(l)$ the Hoffman constants of the growing exact cones (through the
$d$-rows $q_l=\epsilon_lq_0u_l(\zh)/\sigma_m$ and the contact pattern
on $T_{\mathcal B}(l)\setminus F$, and the factors $1/\mathfrak m_l$);
$C_{\rm VP}(l)$ the conditioning of the value-preserving correction;
$\gamma_{\mathcal B}(l)$ the gaps of the bad strict non-peaks; and
$\mu_{\mathcal B}(l)$ the margins of the bad peaks, which are needed both to
control the rows $\tau_l=0$ at bad peaks (through $\sigma|\alpha(k)|=
\lambda_k\mu_k$) and to keep the bad peaks peaks under the raise.  For
finite $\mathcal B$ they are $f$-constants.  For infinite $\mathcal B$ they
are $f$-dependent rates of the kind that the clean sub-windows of
Section~\ref{sec:II-Tfinal} convert into design constants at finite $F$; no
box domination, no raise room and no bounded switching is used.
\end{remark}

If (ND$'$) fails at some level, it fails at all larger levels, and the
value-preserving correction is not needed.

\begin{proposition}[Infinitely many bad carriers without (ND$'$)]
\label{prop:II-kernel}
Let $T=T_{\mathrm{final}}$, $N\ge1$, and let $f\in S_{p_N^*}$ satisfy
\textup{(H2)}, \textup{(H3-inf)} and $\mathfrak r^\star_l>0$ for every good
$l$.  Suppose \textup{(ND$'_{\mathcal B_{\rm np}(l_0)}$)} fails for some
$l_0$, so that by Lemma~\textup{\ref{lem:II-ND}} there are finitely many
\emph{kernel carriers} $k$, numbers $c_k\ne0$ and $\kappa\ne0$ with
$a=(c_k/\kappa)v_k$ on $S_k\cap F\setminus T_{\mathcal B}(l_0)$ and
$F\setminus T_{\mathcal B}(l_0)\subseteq\bigcup_kS_k$.  Put
$m^F_k(l):=\|v_k1_{S_k\cap F\setminus T_{\mathcal B}(l)}\|_1$.  If
\textup{(W$_{\rm inf}$)} holds with $C_{\rm VP}(l)$ replaced by
$1+\max_k1/m^F_k(l)$, then $f\in\Rec$.
\end{proposition}

\begin{proof}
We run the proof of Lemma~\ref{lem:II-Winf} with $f'=f$ (no raise and no
correction: $R:=\emptyset$, $D:=0$; its hypothesis (ND$'$) is used only for
the correction $\Delta a''$, which is not needed here) at levels $l_*\ge l_0$, with
$K_1\ge3C_q(1+2K_\star)/\min_km^F_k(l_*)$ added to its definition, and
show that the deep set is empty.  Write $\hat S_k:=S_k\cap F\setminus
T_{\mathcal B}(l_*)$.  On $\hat S_k$ the sign of $a$ is constant and
$|a_s|=|c_k/\kappa|v_k(s)$.  If the switching of a kernel carrier $k$ is
anti-signed there and $|\tau_k|\ge3|c_k/\kappa|/t$, summing the first
inequality of Lemma~\ref{lem:II-P} over $s\in\hat S_k$ gives
$(|\tau_k|-2|c_k/\kappa|/t)m^F_k\le t/q_0+8t^2+\frac43K_\star t$, hence
$|\tau_k|\le3C_q(1+2K_\star)t/m^F_k\le K_1t$, so $k$ is pinned.  Thus a kept
kernel carrier has $|\tau'_k|\le\frac54|\tau_k|<\frac{15}4|c_k/\kappa|/t$, and on
$\hat S_k$, $s_sX'_s\ge-\frac{15}4|c_k/\kappa|v_k(s)/t>-4|a_s|/t$.  At
$j\in S_k\cap F\cap T_{\mathcal B}(l_*)\setminus T_{\mathcal B}(l_0)$ the
non-kernel part of $X'_j$ is at most $\frac{15}2\sum_{l>k,\,j\in\supp y_l}
\lambda_l|y_l(j)|/(n_lt)\le5\cdot2^{-2j}\delta_k/t$ by
Proposition~\ref{prop:II-features}(d), which is below
$\frac14|c_k/\kappa|v_k(j)/t$ once $2^j\ge20n_k|\kappa/c_k|$.  The finitely many
remaining coordinates of $F$ (including those of $S_l\cap F$ for non-kernel
bad $l$, which lie in $T_{\mathcal B}(l_0)$) have fixed $|a_j|>0$ and are
thick ($|a_j|\ge K_{\rm top}t^2$) at the scales of large levels, where the
argument of step (3) of the proof of Lemma~\ref{lem:II-Winf} applies.  The
data of Lemma~\ref{lem:II-Winf} then live at $f$ itself, and
Theorem~\ref{thm:II-masterinf} concludes, with the window arithmetic of
Theorem~\ref{thm:II-RSinf}.
\end{proof}

When $F\cap S_k$ has bounded gaps beyond $s_{\max}(l)$, $1/m^F_k(l)\le
C_k2^{s_{\max}(l)}/\delta_k$ is absorbed by $\Omega_K(l)$, and no rate is
added.  For very sparse $F\cap S_k$ the rate $1/m^F_k(l)$ is genuine.

\subsection{Deep raises at clean sub-windows, and the residual}
\label{subsec:II-residualinf}

The finite-support master theorems (Sections~\ref{sec:II-exact}
and~\ref{sec:II-Cstar}) work at clean sub-windows with companions that close
rooms, raise thresholds and tune values.  To transport them to infinite base
support one raises \emph{first}, deep and at box level, and then
exactifies at the raised row.  The next two lemmas are the ingredients that
are proved; the transport itself is only sketched
(Remark~\ref{rem:II-Minf}).  Recall the rate objects (O1)--(O12) of
Definition~\ref{def:II-rates}, clean sub-windows
(Definition~\ref{def:II-clean}), and $s_{\rm t}(l)$ of
Definition~\ref{def:II-Tfinal}(5).

\begin{lemma}[A deep raise preserves the classification]\label{lem:II-DR}
Let $T=T_{\mathrm{final}}$, $f\in S_{p_N^*}$ \textup{(}$F$ arbitrary\textup{)},
$l$ a level, $J>s_{\rm t}(l)$, and let $f^r:=f[(a+\Delta a)/\lambda_r,z]$ for a
raise $\Delta a$ supported in $F\cap[J,\infty)$, $\lambda_r:=q^*(a+\Delta a)
\in[1,1+2\|\Delta a\|_1]$.  Then:
\begin{enumerate}
\item[(a)] the values at $f^r$ of the rate objects \textup{(O1)}, \textup{(O2)},
\textup{(O6)} and \textup{(O11)} of level $l$ equal those at $f$, and the
values of \textup{(O10)} are those at $f$ divided by $\lambda_r$;
\item[(b)] every value $u_k(\zh)$ moves by at most $2\|U^*\Delta a\|/\nu$, and
there is $M_f(l)=C_f^{\,l}\,\mathrm{Des}(l)$ \textup{(}$C_f$ an $f$-constant\textup{)}
such that the value of every object of \textup{(O3)}--\textup{(O5)},
\textup{(O7)}, \textup{(O8)}, \textup{(O12)} of level $l$ moves by at most
$M_f(l)\|U^*\Delta a\|$, provided $M_f(l)\|U^*\Delta a\|\le1$.
\end{enumerate}
Consequently, if $w$ is a clean sub-window of level $l$ for $f$,
$M_f(l)\|U^*\Delta a\|\le b(w)^2$ and $\|\Delta a\|_1\le\frac12$, then every object of
\textup{(O1)}--\textup{(O8)}, \textup{(O10)}--\textup{(O12)} that is tiny at $f$
has value $\le2b(w)$ at $f^r$, and every robust one has value $\ge u(w)/2$.
\end{lemma}

\begin{proof}
(a) $z$ and $F=\supp(a+\Delta a)$ are unchanged, so the rooms (O1), the
target rooms (O2) and the shift costs (O6) (which depend only on $F$, on
$z|_{F^c}$ and on design data) are unchanged.  Since $J>s_{\rm t}(l)$, the
coordinates entering (O10), (O11) are not raised, so $a^r_s=a_s/\lambda_r$
there, and the ratios in (O11) are unchanged.
(b) $e^r=U^*(a+\Delta a)/\|U^*(a+\Delta a)\|$ is invariant under
normalization, so $\|e^r-e\|\le2\|U^*\Delta a\|/\nu$ and $|u_k(\zh^r)-
u_k(\zh)|\le\|U^*u_k\|\|e^r-e\|$ with $\|U^*u_k\|\le q^*(u_k)=1$; the block
scalars move by $C_f\|U^*\Delta a\|$ (Lemma~\ref{lem:II-cost}(i)).  (O3),
(O4) are Lipschitz in the values and the block scalars with constants
$\le C_f/\Phi_{\min}(l)^2$, $\Phi_{\min}(l):=\min_{l''\le l}\Phi_{l''}$, and
$1/\Phi_{l''}\le D_{\rm room}(l)\le\mathrm{Des}(l)^{1/6}$; (O5) is linear in the
values divided by $\Phi_{\max}\ge\Phi_{\min}$; the singular values in (O7) are
$1$-Lipschitz in the operator norm (Weyl); the minors (O8) have Lipschitz
constant $\mathrm{Lip}(l)\le\mathrm{Des}(l)$.  For (O12), $\gamma^r_k\,a^r_s/
(\nu^r)^2=(\gamma_k+\ip{U^*u_k}{U^*\Delta a})a_s/(\nu^2+2\ip{U^*a}{U^*\Delta a}+
\|U^*\Delta a\|^2)$ for $s\le s_{\rm t}(l)$: the factor $\lambda_r$ cancels, so the
entries of $\Gamma$ move by $C_f\|U^*\Delta a\|$, and the determinants of order
at most $l$ (entries bounded by $C_f$) move by at most $l!\,C_f^{\,l}$ times
that.  The last statement follows since $b(w)\le1$, $u(w)/b(w)\ge2^{64}$,
and, for (O10), since $\lambda_r\le1+2\|\Delta a\|_1\le2$ by the hypothesis
$\|\Delta a\|_1\le\frac12$ (without it a robust (O10) object would only have
value $\ge u(w)/\lambda_r$).
\end{proof}

Since $\|U^*\Delta a\|\le\mu^\star_J\|\Delta a\|_1$, the requirement
$M_f(l)\|U^*\Delta a\|\le b(w)^2$ holds for a box-level deep raise as in
Lemma~\ref{lem:II-Winf} as soon as $C_f^{\,l}\mathrm{Des}(l)\mu^\star_J\psi(J)\le
b(w)^2$, and there $\|\Delta a\|_1\le\psi(J)\le\frac12$ for large $J$; by Lemma~\ref{lem:II-base}(d) the corresponding $2^J$ is of order
$l\log_2(C_f)+\log_2\log_2(\mathrm{Des}(l)/b(w))$.  No value-preserving
correction is needed: the classification at $w$ survives with factors $2$,
and all exactifications are done at $f^r$ with its own values.  The
shift-pinning rates (O9) are not covered by the lemma: their definition is an
infimum over unbounded $\tau$ with value-dependent coefficients, and we do
not prove their stability.

At infinite base support the exactifying moves of Section~\ref{sec:II-exact}
(banks and pulls) may fall on support coordinates.  The following
first-order calculus of moves on the support replaces them there.

\begin{lemma}[Moves on the support]\label{lem:II-TR}
Let $T=T_{\mathrm{final}}$ and let $f^\circ=f[a^\circ,z^\circ]$ be a first row
with base support $F^\circ$ and $\nu^\circ:=\|U^*a^\circ\|$, $e^\circ$,
$\gamma^\circ_k:=\ip{U^*u_k}{U^*a^\circ}$.  For a finite set $Y\subseteq F^\circ$
and $\Delta\in\R^Y$ with $|\Delta_x|\le|a^\circ_x|/2$ let
$A:=a^\circ+\sum_{x\in Y}\Delta_xe^*_x$ and $f^\Delta:=f[A/q^*(A),z^\circ]$.  Put
$\alpha_x:=(\mu^\star_x)^2a^\circ_x/(\nu^\circ)^2$.
\begin{enumerate}
\item[(a)] For every carrier $k$,
\begin{gather*}
 u_k(\zh^\Delta)-u_k(\zh^\circ)=\frac1{\nu^\circ}\sum_{x\in Y}(\mu^\star_x)^2
 u_k(x)\Delta_x-\frac{\gamma^\circ_k}{\nu^\circ}\sum_{x\in Y}\alpha_x\Delta_x+
 R_k,\\ |R_k|\le\frac{C}{(\nu^\circ)^2}\Big(\sum_{x\in Y}\mu^\star_x|\Delta_x|
 \Big)^2,
\end{gather*}
with an absolute constant $C$, provided $\sum_x\mu^\star_x|\Delta_x|\le
\nu^\circ/2$.
\item[(b)] \textup{(Pairs.)}  Let $l''$ be a carrier, $L\ge l''$ a level and
$s<s'$ in $S_{l''}\cap F^\circ$ beyond $s_{\max}(L)$, and put
$\mathfrak a_x:=a^\circ_x/v_{l''}(x)$.  For $\Delta_{s'}\in\R$ let
$\Delta_s:=-(\alpha_{s'}/\alpha_s)\Delta_{s'}$.  Then $|\Delta_s|\le
(\mu^\star_{s'}/\mu^\star_s)^2(|a^\circ_{s'}|/|a^\circ_s|)|\Delta_{s'}|$, every
carrier $k\le L$ with $k\ne l''$ satisfies $|u_k(\zh^\Delta)-u_k(\zh^\circ)|\le
|R_k|$, and
\[
 u_{l''}(\zh^\Delta)-u_{l''}(\zh^\circ)=\frac{(\mu^\star_{s'})^2v_{l''}(s')}
 {\nu^\circ}\Big(1-\frac{\mathfrak a_{s'}}{\mathfrak a_s}\Big)\Delta_{s'}+R_{l''} .
\]
\item[(c)] \textup{(Single moves are not private.)}  A single move $\Delta_x$
at a coordinate $x$ with $u_k(x)=0$ changes $u_k(\zh)$ at first order by
$-\gamma^\circ_k\alpha_x\Delta_x/\nu^\circ$.
\end{enumerate}
\end{lemma}

\begin{proof}
(a) $U^*A=U^*a^\circ+\sum_x\Delta_x\mu^\star_xk_x$, and $e^\Delta=U^*A/\|U^*A\|$
does not depend on the normalization of $A$; $z^\circ$ is fixed, so
$u_k(\zh^\Delta)-u_k(\zh^\circ)=\ip{U^*u_k}{e^\Delta-e^\circ}$.  The map
$h\mapsto(U^*a^\circ+h)/\|U^*a^\circ+h\|$ has derivative $h\mapsto(h-
\ip h{e^\circ}e^\circ)/\nu^\circ$ at $0$ and second derivative bounded by
$C/(\nu^\circ)^2$ on $\|h\|\le\nu^\circ/2$.  With $h=\sum_x\Delta_x\mu^\star_xk_x$,
$\ip h{e^\circ}=\sum_x\Delta_x(\mu^\star_x)^2a^\circ_x/\nu^\circ$ and
$\ip{U^*u_k}{k_x}=\mu^\star_xu_k(x)$, $\ip{U^*u_k}{e^\circ}=\gamma^\circ_k/
\nu^\circ$, which gives the formula; $\|U^*u_k\|\le1$.
(b) $\alpha_s\Delta_s+\alpha_{s'}\Delta_{s'}=0$ removes the second term of (a).
For $k\le L$, $k\ne l''$, $u_k$ vanishes at $s,s'$ (Theorem~\ref{thm:II-lemmaB}
and (P1): beyond $s_{\max}(L)$ only $l''$ is nonzero on $S_{l''}$ among the
carriers $\le L$).  For $l''$: $(\mu^\star_s)^2v_{l''}(s)\Delta_s=
-(\mu^\star_{s'})^2a^\circ_{s'}v_{l''}(s)\Delta_{s'}/a^\circ_s=-(\mu^\star_{s'})^2
v_{l''}(s')(\mathfrak a_{s'}/\mathfrak a_s)\Delta_{s'}$.  (c) is (a) with $Y=\{x\}$.
\end{proof}

Since $\mu^\star$ decreases, the deeper coordinate $s'$ must do the work:
$|\Delta_s|\le u^{-1}|\Delta_{s'}|$ when $|a^\circ_s|\ge u|a^\circ_{s'}|$, whereas
pairing in the other order would require astronomically large moves.  If
$|a^\circ_s|,|a^\circ_{s'}|\ge u$, the spread $|1-\mathfrak a_{s'}/\mathfrak a_s|\ge u$ and
$s,s'\le s_{\rm t}(L)$, the private effect in (b) is at least
$\mathrm{Des}(L)^{-1}u\,|\Delta_{s'}|$ (Proposition~\ref{prop:II-bank}(c)), so robust
pairs are two-sided private tuning resources with design constants; a single
anchor coordinate $s_0$ with $u_k(s_0)=0$ for the tuned carriers can replace
the partner by choosing $\Delta_{s_0}$ so that $\sum_x\alpha_x\Delta_x=0$.  By (c),
a single robust support move is not an admissible tuning resource: its common
term moves every value at first order.  Where the tuning coordinates are not
in $F$, Lemma~\ref{lem:II-TU} applies verbatim (its proof uses the base support
only through hypothesis (T2), which keeps the bank and pull coordinates off
the support, and through Lemma~\ref{lem:II-bankformula}, which holds for every
$A^0\in\ell_1$).

\begin{proposition}[Reduction through raised companions]\label{prop:II-Minfred}
Let $T$ be admissible, $f\in S_{p^*}$, $g\in\Cset(f)$, $\rho\in(0,1)$ and
$\eta_0$ as in Theorem~\textup{\ref{thm:II-ERT}}.  Suppose there are first rows
$f^{\rm ex}_j$, raises $\Delta a_j$ of $f^{\rm ex}_j$ with $\lambda_j:=
q^*(a^{\rm ex}_j+\Delta a_j)\ge1$, the raise companions $f_j$, and windows
$(T_j,n_j)$, numbers $c_j\in(0,1]$, $K_j\ge0$ and functionals $g_{j,t}$
satisfying \textup{(b)}--\textup{(e)} of Theorem~\textup{\ref{thm:II-ERT}} at
$f_j$, and
\[
 \varepsilon'_j:=p^*(f^{\rm ex}_j-f)+2\|U^*\Delta a_j\|+p^*\big(L^*(w_j-
 w^{\rm ex}_j)\big)\le\theta_*c_j^2(T_j2^{-n_j})^2 .
\]
Then $(f,\rho g)\in\cl\NA((c_0,p),\ell_2^2)$.
\end{proposition}

\begin{proof}
By Lemma~\ref{lem:II-RT} at $f^{\rm ex}_j$ and $p^*(f^{\rm ex}_j+\lambda_jr\rho
g)\le s(\lambda_j\rho r)+p^*(f^{\rm ex}_j-f)$, hypothesis (a) of
Theorem~\ref{thm:II-ERT} holds with this $\varepsilon'_j$ (as
$\lambda_j\ge1$); apply that theorem.
\end{proof}

The finite-support companions of Sections~\ref{sec:II-exact}
and~\ref{sec:II-Cstar} have $p^*(f^{\rm ex}_j-f)=o(T_{\rm lo}^2)$; the
proposition reduces the transport of the finite-support master theorems to
the production of exact window data, cushion-compatible on $F$, at raised
exactifying companions.

\begin{proposition}[The infinite-support half of Lemma Z]\label{prop:II-LSC}
Let $T$ be a signature-ladder operator and $N\ge1$.  The following are
equivalent:
\begin{enumerate}
\item[(i)] $\NA((c_0,p_N),\ell_2^2)$ is dense;
\item[(ii)] Lemma~Z [I,~Problem~8.21] holds at every $f$ with finite base
support, and \textup{(LSC-trunc)} holds: for every $f\in S_{p_N^*}$,
$g\in\Cset(f)$, $\rho<1$ and $\varepsilon>0$ there is a first row $f'$ with
finite base support, $p_N^*(f'-f)<\varepsilon$ and $\dist(\rho g,\Cset(f'))<
\varepsilon$.
\end{enumerate}
\textup{(LSC-trunc)} holds for every pair $(f,g)$ covered by
Theorems~\textup{\ref{thm:II-transferinf}}, \textup{\ref{thm:II-raised}},
\textup{\ref{thm:II-masterinf}}, \textup{\ref{thm:II-Bstarinf}}--\textup{\ref{thm:II-cushionroom}},
\textup{\ref{thm:II-R1}}, \textup{\ref{thm:II-N}}, Corollary~\textup{\ref{cor:II-BTinf}},
and, for $T=T_{\mathrm{final}}$, Theorems~\textup{\ref{thm:II-RS}}--\textup{\ref{thm:II-RSinf}}
and Proposition~\textup{\ref{prop:II-kernel}}.
\end{proposition}

\begin{proof}
For both operators [I,~Theorem~8.20] holds (its proof uses
[I,~Theorem~8.18] and [I,~Lemma~8.4(a)], and Lemma~\ref{lem:II-GW}(b)).
(i)$\Rightarrow$(ii): [I,~Theorem~8.20] gives Lemma~Z everywhere, and
[I,~Lemma~1.16] gives norm-attaining approximants, which have finite base
support [I,~Proposition~1.10(c)], with $\rho g\in\Li\Cset(f_n)$.
(ii)$\Rightarrow$(i): given $(f,g,\rho,\varepsilon)$, (LSC-trunc) gives $f'$ with
finite $F'$ and $g'\in\Cset(f')$ with $p^*(g'-\rho g)<\varepsilon$; Lemma~Z at
$f'$ for $(g',\rho',\varepsilon)$ gives $f''\in\mathcal R_0$ and $g''\in
\Cset(f'')$ with $p^*(f''-f')<\varepsilon$ and $p^*(g''-\rho'g')<\varepsilon$.
Then $p^*(f''-f)<2\varepsilon$ and $p^*(g''-\rho'\rho g)<2\varepsilon$; as
$\rho\rho'$ ranges over $(0,1)$, Lemma~Z holds at $f$, hence everywhere, and
[I,~Theorem~8.20] gives (i).  For the last statement: each of the listed
results produces norm-attaining approximants $f'_i\to f$ (finite base
support) and $g'_i\in\Cset(f'_i)$ with $g'_i\to\rho g$ (Theorem~\ref{thm:II-ERT}
through rows $f_j\to f$, each of which is recovered along norm-attaining
rows by Theorem~\ref{thm:II-raised}).
\end{proof}

\begin{remark}[Truncation is not the obstruction]\label{rem:II-trunc}
In Theorem~\ref{thm:II-raised} the support beyond the window $N_w$ becomes
contacts of the norm-attaining approximant, at a first-order cost controlled
by the tail of the \emph{fixed} vector $v=b^+-b^-$, and the near support is
raised where the average data need it.  So, once exact data with one-sided
cushion sparsity exist (fixed, or window data at $f$, or at a companion with
small transfer error), truncating the support costs nothing; what remains at
infinite $F$ is the production of such data.  This is an interpretation of
the proofs, not a theorem about pairs without exact data.
\end{remark}

\begin{remark}[Transport of Master Theorem~III$'$ to infinite support;
sketch]\label{rem:II-Minf}
\emph{Claim (sketch, not proved here).}  For $T_{\mathrm{final}}$ and every
$N$, a first row $f$ with infinite $F$ belongs to $\Rec$ unless, at all but
finitely many levels, every clean sub-window $w$ exhibits (C$^*_w$)
(Definition~\ref{def:II-Cstar}, read at the deep-raised row), or a tiny
\emph{tuning regularity modulus} (below), or a failure of (H2), or a bad
degenerate peak with the swallowing sign none of whose block's donor peaks
has a tuning resource.  \emph{Intended route}: at a clean $w$, (1) a deep
box-level raise $f\to f^r$ beyond $J>s_{\rm t}(l)$ (Lemma~\ref{lem:II-DR}; no
value-preserving correction); (2) at $f^r$, the moves (C1)--(C4) of
Section~\ref{sec:II-exact} and the exactification of tiny minors of
Section~\ref{sec:II-Cstar}, with banks and pulls replaced by the support moves
of Lemma~\ref{lem:II-TR} where the coordinates lie in $F$; (3) exact
$d$-neutral window data at the companion, cushion-compatible on $F$ by the
argument of Lemma~\ref{lem:II-Winf} (Lemma~\ref{lem:II-P} and the pigeonholes);
(4) Proposition~\ref{prop:II-Minfred} and Theorem~\ref{thm:II-raised}.
\emph{Items not written out}: the transplant of exact data at infinite $F$
(the finite-support transplant copies $B_+1_F$ and uses [I,~Lemma~8.9]; it must
be rerun with the clamped data of Theorem~\ref{thm:II-R1} and
Lemma~\ref{lem:II-P} at the decomposition row); the stability of the shift-pinning
rates (O9) under the deep raise; the exactification of tiny minors with the
additional pinned faces and contact rows of Lemma~\ref{lem:II-Winf} (these are
subsystems, whose minors are rate objects); the joint fixed point of the
support tuning with the threshold buffer; and the raise level $8\mathrm{Bx}_*$
needed because the box rows allow $|\tau'_l|\le12\lambda_l/t$.  \emph{Tuning
regularity modulus}: the first-order resource map (private two-sided columns
of contacts, convertible thin support coordinates, robust pairs and anchors of
Lemma~\ref{lem:II-TR}; one-sided pull columns; and single-support columns
$((\mu^\star_s)^2/\nu)(u_k(s)-\gamma_ka_s/\nu^2)_{k\in L_0}$ at robust support
coordinates) must map the cone of admissible moves onto $\R^{L_0}$ with a right
inverse of norm at most $(\mathrm{Des}(l)/u(w))^C$.  Its modulus is a finite
minimum of determinants in $(a_s)_{s\le s_{\rm t}(l)}$ and values, i.e.\ built
from rate objects; when only single-support columns exist on proportional
profiles $a=c_kv_k$, it reduces to the Sherman--Morrison denominator
$|1-\sum_kc_k\gamma_k/\nu^2|$ of Remark~\ref{rem:II-newrates}, which vanishes
in particular when (ND$'$) fails (then $a|_F=\sum_kc_ku_k|_F$); the converse
is not claimed.  (A description of this residual as the
conjunction of ``every usable coordinate is thick support'' and ``nearly
proportional profile'' alone is too narrow: thick bank coordinates with
off-support pull coordinates of the wrong orientation and tiny spreads also
lack resources.)
\end{remark}

\begin{remark}[The residual at infinite base support]\label{rem:II-residualinf}
(a) \emph{Recovered} (proved above).  \emph{Rows} (every mate is
recovered, i.e.\ the row is in $\Rec$): first rows with signature room (SR),
sign-mixed room or cushion room at any $F$
(Theorems~\ref{thm:II-Binf}--\ref{thm:II-cushionroom}); finitely many
cushion-sparse swallowed carriers for every signature-ladder operator, under
(W$^\star$), (H2), (H3-inf) (Theorem~\ref{thm:II-R1}); box-dominated support
$F=\N$ for every admissible operator under (SC$_I$) and (W$_M$)
(Theorem~\ref{thm:II-N}); for $T_{\mathrm{final}}$: finitely many bad carriers
with any support profile, under (W$^\star$), (H2), (H3-inf)
(Theorems~\ref{thm:II-RSstar}, \ref{thm:II-NDprime}), and infinitely many bad
carriers under (H2), (H3-inf) and (W$_{\rm inf}$) (Theorem~\ref{thm:II-RSinf},
Proposition~\ref{prop:II-kernel}).  \emph{Pairs} (a given mate is recovered):
window-pinned mates at any $F$ (Theorem~\ref{thm:II-Bstarinf}); cushion-sparse
mates of block-tame rows, with $V_f$ of sparse profile
(Corollary~\ref{cor:II-BTinf}); and, for $T_{\mathrm{final}}$, mates $\rho g$
with fixed $d$-neutral data, $\rho^2\kappa_w<1$, (RR$_W$) and (ND$'_{L_0}$)
(Theorem~\ref{thm:II-A}).
(b) \emph{Open for $T_{\mathrm{final}}$ at infinite $F$}, for pairs $(f,g)$
with $g$ not window-pinned: (b1) the rates: rooms of good carriers decaying
too fast for (W$^\star$) (the support-free part of approximate swallowing),
and, for infinitely many bad carriers, failure of (W$_{\rm inf}$), i.e.\ of the
rates $H(l)$, $C_{\rm VP}(l)$, $\gamma_{\mathcal B}(l)$, $\mu_{\mathcal B}(l)$
(and $1/m^F_k(l)$ in Proposition~\ref{prop:II-kernel}); these would become
design constants through the transport of Remark~\ref{rem:II-Minf}; (b2) failure
of (H2), and bad degenerate peaks with the swallowing sign, (H3-inf), whose
treatment needs a donor with a tuning resource; (b3) the finite-support
residual (C$^*$) of Definition~\ref{def:II-Cstar}, read at deep-raised rows; (b4)
a tiny tuning regularity modulus (Remark~\ref{rem:II-Minf}); (b5) the scrambling
condition (SC) for non-$d$-neutral data at raised rows below the scale of the
normer change, which enters only inside (C$^*$).
(c) For signature-ladder operators without a fast-decaying diagonal base,
non-sparse support swallowing (critical and super-critical profiles beyond the
special cases of Remark~\ref{rem:II-R1}(d)) and infinitely many swallowed
carriers without box domination remain open.
(d) No counterexample is known in any of these classes.  (\emph{Interpretation,
not a statement proved here}:) every infinite-support
mechanism examined is either harmless after a deep raise (thin, critical,
super-critical, $\mu$-thin supports, and box-size switching under
(W$_{\rm inf}$)), pinned by the mate's
own flip budget (thin shallow coordinates), cushion-dominated (failure of
(ND$'$)), or a rate phenomenon of the same kind as at finite base support.
\end{remark}

\section{Status and open problems}\label{sec:II-status}

Throughout this section $T=T_{\mathrm{final}}$ with base $U_{\mathrm{final}}$
(Section~\ref{sec:II-Tfinal}), except where another operator is named, and
$N\ge1$.

\subsection{What is proved}
\begin{enumerate}
\item[(a)] $T_{\mathrm{final}}$ is admissible, satisfies the conclusion of
Mart\'in's Lemma~B, does not depend on $N$, and satisfies every hypothesis
of [I,~Section~8] on every choice of sub-windows
(Theorem~\ref{thm:II-lemmaB}, Lemma~\ref{lem:II-GW}).  Hence density of
$\NA((c_0,p_N),\ell_2^2)$ is equivalent to Lemma~Z for $T_{\mathrm{final}}$
and $N$ [I,~Theorem~8.20], and Lemma~Z for infinitely many $N$ gives
density for Mart\'in's norm $p$ built with $T_{\mathrm{final}}$
(Remark~\ref{rem:II-blocks}).
\item[(b)] \emph{Finite base support.}  Every $f\in S_{p_N^*}$ with finite
base support that is not a (C$^*$) row belongs to $\Rec$
(Corollary~\ref{cor:II-Cstarnec}).  Among (C$^*$) rows, the block-tame ones
are in $\Rec$ [I,~Corollary~7.22] (for instance the self-aligned rows,
Corollary~\ref{cor:II-SArec}); so are the rows in $\mathcal R_0\cup\mathcal
R_0^\pm\cup\mathcal R_S$ [I,~Section~8].  For the norm built with the
modified design $D^{U1'}$, the rows with finite base support covered by
Master Theorem~IV$'$ are in $\Rec$ (Theorem~\ref{thm:II-MTIV},
Corollary~\ref{cor:II-MTIVintr}).  Constant-sign maximal contact never
produces a (C$^*$) row (Proposition~\ref{prop:II-constcontact}).
\item[(c)] \emph{Infinite base support.}  The rows listed in
Remark~\ref{rem:II-residualinf}(a) are in $\Rec$ and the pairs listed there
are in $\cl\NA$; for each of them (LSC-trunc) holds
(Proposition~\ref{prop:II-LSC}).
\end{enumerate}

\subsection{The main problems}

\begin{problem}[Lemma Z for the designed operator]\label{prob:II-Z}
Let $N\ge1$ and $T=T_{\mathrm{final}}$.  Is it true that for every $f\in
S_{p_N^*}$, $g\in\Cset(f)$, $\rho<1$ and $\varepsilon>0$ there is
$f'\in\mathcal R_0$ with $p_N^*(f'-f)<\varepsilon$ and $\dist(\rho
g,\Cset(f'))<\varepsilon$?  Equivalently [I,~Theorem~8.20], is
$\NA((c_0,p_N),\ell_2^2)$ dense?
\end{problem}

Problem~\ref{prob:II-Z} is open for every $N$, including $N=1$.  By
Proposition~\ref{prop:II-LSC}, it has a positive answer if and only if
Lemma~Z holds at every first row with finite base support and
(LSC-trunc) holds.  By Corollary~\ref{cor:II-Cstarnec}, [I,~Lemma~1.16] and
[I,~Lemma~8.4(a)], the first condition follows from a positive answer to
Problem~\ref{prob:II-Cstar}; the second is the following problem.

\begin{problem}[Lower semicontinuity through truncations]\label{prob:II-LSC}
Let $N\ge1$ and $T=T_{\mathrm{final}}$.  Does \textup{(LSC-trunc)} of
Proposition~\textup{\ref{prop:II-LSC}} hold, i.e.\ for every $f\in
S_{p_N^*}$, $g\in\Cset(f)$, $\rho<1$ and $\varepsilon>0$, is there a first
row $f'$ with finite base support, $p_N^*(f'-f)<\varepsilon$ and
$\dist(\rho g,\Cset(f'))<\varepsilon$?
\end{problem}

\begin{problem}[Mart\'in's norm]\label{prob:II-p}
Is $\NA((c_0,p),\ell_2^2)$ dense, for Mart\'in's norm $p$ built with
$T_{\mathrm{final}}$, or with any admissible operator?
\end{problem}

A positive answer to Problem~\ref{prob:II-Z} for infinitely many $N$
answers Problem~\ref{prob:II-p} positively for $T_{\mathrm{final}}$
[I,~Lemma~1.5].  The row-wise results of Sections~\ref{sec:II-exact}
to~\ref{sec:II-infinite} do not transfer to $p$, since their hypotheses
depend on $N$ (Remark~\ref{rem:II-openmisc}(d),
Remark~\ref{rem:II-notMTIV}(b)).

\subsection{Finite base support: the residual (C$^*$)}
Problem~\ref{prob:II-Cstar} asks whether every (C$^*$) row is in $\Rec$.
(C$^*$) rows exist (Proposition~\ref{prop:II-alignedres}); the
nested-tuning rows of Proposition~\ref{prop:II-NT} are (C$^*$) rows that are
not block-tame, but their existence is only sketched
(Remark~\ref{rem:II-NTexist}), and only their mates with exact data are
known to be recovered (Corollary~\ref{cor:II-NTR}).  Two routes are
available; each leaves precise gaps.

\emph{Route through absorbers and levers} (Subsection~\ref{subsec:II-absorb}).
Master Theorem~IV$'$ needs a $\varkappa$-neutral lever in every activity
class with active shift, and in mixed classes the absence of frustrated
positive carriers.  The first requirement is not a technicality: a donor
raise moves the block scalar $\varkappa$ by the same order as the raise,
and the exact shifted system depends on a block only through ratios
(Corollary~\ref{cor:II-Rkt}, Proposition~\ref{prop:II-Rkt}), so this route
needs a second, $\varkappa$-neutral lever in every such block
(Remark~\ref{rem:II-onelever}; no claim is made that a one-lever argument
cannot succeed).  An earlier argument in which a donor
raise alone exactifies the block scalars (which would have put every row
with finite base support in $\Rec$ for $N=1$) fails for this reason; that
conclusion is not proved and is not used.

\begin{problem}[Status coherence without a lever]\label{prob:II-KN}
For the design $D^{U1'}$ and $N\ge1$: let $f$ have finite base support and
let $w$ be a clean sub-window with an active block, an active
near-threshold switching carrier $d$, and no inactive $\Omega$-carrier of
robust relative position, and suppose that the exact zero set of the tiny
minors (Definition~\ref{def:II-Gamma}), in the ratio variables and near the
ratios of $f$, lies in $\{\varrho_d\ge1\}$.  Is $f\in\Rec$?  In particular,
can the window data be completed with $d$ a peak whose excess
$|\omega_+(d)|+|\omega_-(d)|$ is $O(Kt/\lambda_d)$, with $K$ of the form
$C_f(\mathrm{Des}/u)^C$, or can a $\varkappa$-neutral lever be supplied by
other means?  (By Proposition~\ref{prop:II-Rkt}, $\varrho_d$ at an exact
completion is a function of the ratio variables up to fine terms, so under
this hypothesis every exact completion near the ratios of $f$ has
$\varrho_d\ge1$ up to fine terms, and $d$ cannot be kept strictly below its
threshold by such a completion.  Remark~\ref{rem:II-openKN} records a
sketch for carriers with inactive excess and a heuristic genericity route.)
\end{problem}

\begin{problem}[Mixed activity classes]\label{prob:II-mix}
For the design $D^{U1'}$ and $N\ge2$: are the rows with finite base
support whose clean sub-windows carry mixed activity classes in $\Rec$?  One
route is an exact completion (Lemma~\ref{lem:II-schauder}) at which no
carrier of a negative block is bad in the sense of
Lemma~\ref{lem:II-SCband} (Remark~\ref{rem:II-openmix}); finite models
indicate that mixing can force the bad state at every exact completion,
which is numerical evidence only and concerns this route, not the
recoverability of the rows.
\end{problem}

\emph{Route through exact coupling and uniform composition}
(Remark~\ref{rem:II-openS}).

\begin{problem}[Exact coupling]\label{prob:II-S1}
With two or more shifted blocks, or without a free-ray carrier with
$q<0$, can the coupling rows $\Delta d_m=d^\#_m(\omega^-)-d^\#_m(\omega^+)$ be
satisfied exactly inside the exact combinatorial cone, with constants of
the form $C_f(\mathrm{Des}/u)^C$?  (One shifted block with a free-ray carrier:
Proposition~\ref{prop:II-S1}.)  Together with this, coarse exactness at the
companion modulo the coupling rows and the contributions of carriers
beyond the level (item (RT$^*$c) of Remark~\ref{rem:II-openS}) is only
sketched.
\end{problem}

\begin{problem}[Uniform composition]\label{prob:II-S2}
Are there $d$-consistent engineered approximants (the normalized values of
the finitely many coarse switching carriers re-tuned exactly to their
values at the row) with per-piece bounds $T_0\ge c_\flat t$ uniform in the
scale, such that the averaging of Theorem~\ref{thm:II-E} can be carried out
at one engineered approximant for all pieces of a window, and the
violation mass stays below the late threshold of the companion?  (Items
(S2a$'$), (S2b) are sketches and (S2c) is a heuristic,
Remark~\ref{rem:II-openS}.  The route uses $d$-consistency because without
it window masses change the coarse block data at first order, which gives a
mismatch of order up to $s_1/t^2$; the upper bound is proved, finite models
indicate that it is attained up to a constant, which is numerical evidence
only, and that the mismatch cannot then be rebalanced is heuristic.)
\end{problem}


\begin{remark}[Further open points at finite base support]\label{rem:II-statusfin}
(a) For the nested-tuning rows of Proposition~\ref{prop:II-NT}, Lemma~Z is
not proved for general mates; a route through exact window data with a
shift column and transplant to block-tame companions is only a sketch
(Remark~\ref{rem:II-openmisc}(a)).  (b) The existence of these rows is a
sketch (Remark~\ref{rem:II-NTexist}).  (c) Earlier sketches of a stable
shifted resonance theorem required, besides Problems~\ref{prob:II-KN}
and~\ref{prob:II-mix}: exact absorption of fine-peak residues on finitely
many coarse free coordinates (now proved by the zero-value absorber pairs
of Lemma~\ref{lem:II-zerovalue} and Proposition~\ref{prop:II-absorb}, for
$D^{U1'}$); the scrambling condition at companions (proved in configuration
(i), Lemma~\ref{lem:II-SCi}); a shift direction constant along a window
(proved through the one-ray normalization, Proposition~\ref{prop:II-ray});
and a joint completion (Lemma~\ref{lem:II-schauder}).  The exactification
of the block scalars $(\vartheta_m,A_m)$ is the part that remains, in the
form of Problem~\ref{prob:II-KN}.
\end{remark}

\subsection{Infinite base support}
For first rows with infinite base support, Problem~\ref{prob:II-LSC} is
open for $T_{\mathrm{final}}$ outside the classes of
Remark~\ref{rem:II-residualinf}(a).  The obstructions are the following.

\begin{problem}[Transport of Master Theorem~III$'$]\label{prob:II-Minf}
Let $N\ge1$ and $T=T_{\mathrm{final}}$.  Is every $f\in S_{p_N^*}$ with
infinite base support in $\Rec$ unless, at all but finitely many levels,
every clean sub-window exhibits (C$^*$) at the deep-raised row, a tiny
tuning regularity modulus (Problem~\ref{prob:II-tuning}), a failure of
(H2), or a degenerate bad peak with the swallowing sign whose block has no
donor with a tuning resource?  The intended route
(Remark~\ref{rem:II-Minf}) combines a deep raise (Lemma~\ref{lem:II-DR}),
the moves of Sections~\ref{sec:II-exact} and~\ref{sec:II-Cstar} with support
moves (Lemma~\ref{lem:II-TR}), exact window data at the companion
(Lemma~\ref{lem:II-Winf}) and Proposition~\ref{prop:II-Minfred}; it is a
sketch.  The items not written out are: the transplant of exact data at
infinite base support, with clamped data at the decomposition row; the
stability of the shift-pinning rates under the deep raise; the
exactification of tiny minors with the additional pinned faces and contact
rows; and the joint fixed point of the support tuning with the threshold
buffer.
\end{problem}

\begin{problem}[Infinitely many bad carriers]\label{prob:II-Winf}
Let $N\ge1$ and $T=T_{\mathrm{final}}$.  Is every $f\in S_{p_N^*}$ satisfying
the hypotheses of Theorem~\ref{thm:II-RSinf} other than
\textup{(W$_{\rm inf}$)} in $\Rec$?  The rate condition involves the
quantities $H(l)$, $C_{\rm VP}(l)$, $\gamma_{\mathcal B}(l)$, the bad-peak
margin rate $1/\mu_{\mathcal B}(l)$, and, when (ND$'$) fails, the kernel rate
$1/m^F_k(l)$ of Proposition~\ref{prop:II-kernel}; the same question arises
for rooms of good carriers that decay too fast for (W$^\star$).  These
rates would become design constants through Problem~\ref{prob:II-Minf}.
\end{problem}

\begin{problem}[Tuning regularity]\label{prob:II-tuning}
At a clean sub-window of a first row with infinite base support, does the
first-order resource map of Remark~\ref{rem:II-Minf} (private two-sided
columns of contacts, convertible thin support coordinates, robust pairs
and anchors of Lemma~\ref{lem:II-TR}, one-sided pull columns, and
single-support columns at robust support coordinates) map the cone of
admissible moves onto $\R^{L_0}$ with a right inverse of norm at most
$(\mathrm{Des}(l)/u(w))^C$, or can the first rows at which it does not be
recovered otherwise?  On proportional profiles the modulus reduces to the
Sherman--Morrison denominator of Remark~\ref{rem:II-newrates}, which
vanishes in particular when (ND$'$) fails (then $a|_F=\sum_kc_ku_k|_F$); the
converse is not claimed.
\end{problem}

\begin{problem}[Condition (H2) and degenerate bad peaks]\label{prob:II-H2}
Are first rows with infinite base support at which (H2) fails, or which
have a degenerate bad peak with the swallowing sign (failure of
(H3-inf)) whose block has no donor peak with a tuning resource, in
$\Rec$?
\end{problem}

\begin{problem}[Scrambling along deep raises]\label{prob:II-SCdeep}
For non-$d$-neutral data at raised rows, does the scrambling condition
(SC) [I,~Definition~7.16] transfer along a deep raise below the scale of
the change of the normer?  This enters only through shifted
(non-$d$-neutral) data, i.e.\ inside (C$^*$) at infinite base support
(Problem~\ref{prob:II-Minf}), and in the extension of Theorem~\ref{thm:II-A}
to non-$d$-neutral fixed data, which is open (Theorem~\ref{thm:II-A} itself
assumes $d$-neutral data and uses no (SC)).
\end{problem}

\begin{remark}[Other operators]\label{rem:II-statusother}
For signature-ladder operators without a fast-decaying diagonal base,
non-sparse support swallowing (critical and super-critical flip profiles
beyond the special cases of Remark~\ref{rem:II-R1}) and infinitely many
swallowed carriers without box domination remain open
(Remark~\ref{rem:II-residualinf}(c)).  For a general admissible operator,
Problem~\ref{prob:II-p} and the density for $p_N$ are open for every $N$;
the resonance obstruction [I,~Theorem~6.6] shows that the intrinsic
mechanisms of Part~I do not suffice in general, and Theorem~\ref{thm:II-NL}
shows that \emph{block-tame} approximating rows for a mate whose profile
takes different values at two coordinates $s_1,s_2$ outside the data
supports must put $s_1$ or $s_2$ into a bank, the support or the support of
a strict non-peak, or be aligned with maximal contact; nothing is proved
about approximating rows that are not block-tame.  (Such mates exist at
self-aligned rows of $T_{\mathrm{final}}$ for $N=1$,
Proposition~\ref{prop:II-NLreal}.)
\end{remark}

\begin{remark}[Evidence]\label{rem:II-evidence}
No counterexample is known in any of the classes above.  Every mechanism
examined at finite base support is either removed at a clean sub-window
or reduced to (C$^*$); every mechanism examined at infinite base support
is harmless after a deep raise, pinned by the flip budget of the mate,
cushion-dominated, or a rate phenomenon of the same kind as at finite base
support (Remark~\ref{rem:II-residualinf}(d)).  This is a description of the
state of the proofs, not evidence for either answer.
\end{remark}

\begin{thebibliography}{99}
\bibitem{BCR}
J.~Bochnak, M.~Coste and M.-F.~Roy, \emph{Real Algebraic Geometry},
Ergebnisse der Mathematik und ihrer Grenzgebiete (3) \textbf{36},
Springer, Berlin, 1998.
\bibitem{DGS}
G.~Debs, G.~Godefroy and J.~Saint Raymond,
\emph{Topological properties of the set of norm-attaining linear
functionals}, Canad. J. Math. \textbf{47} (1995), 318--329.
\bibitem{Hoffman}
A.~J.~Hoffman, \emph{On approximate solutions of systems of linear
inequalities}, J. Research Nat. Bur. Standards \textbf{49} (1952),
263--265.
\bibitem[I]{PartI}
\emph{Recovery of mates on Mart\'in norms and the resonance obstruction},
research note (Part~I of this study), October 2026.
\bibitem{JW}
W.~B.~Johnson and J.~Wolfe, \emph{Norm attaining operators},
Studia Math. \textbf{65} (1979), 7--19.
\bibitem{KLMW}
V.~Kadets, G.~L\'opez, M.~Mart\'in and D.~Werner,
\emph{Norm attaining operators of finite rank}, arXiv:1905.08272.
\bibitem{KLMW20}
V.~Kadets, G.~L\'opez, M.~Mart\'in and D.~Werner,
\emph{Equivalent norms with an extremely nonlineable set of norm attaining
functionals}, J. Inst. Math. Jussieu \textbf{19} (2020), 259--279.
\bibitem{Lindenstrauss}
J.~Lindenstrauss, \emph{On operators which attain their norm}, Israel J.
Math. \textbf{1} (1963), 139--148.
\bibitem{Martin}
M.~Mart\'in,
\emph{A Banach space whose set of norm-attaining functionals is
algebraically trivial}, J. Funct. Anal. \textbf{288} (2025);
arXiv:2406.07273.
\end{thebibliography}
\end{document}
